% ==================================================================
%  素粒子論のための代数トポロジー演習（全200問）
%  本体ファイル。プリアンブルは preamble.tex、各章は chapters/chNN.tex。
%  旧版（純粋数学向けの問題集）は archive/test_math_v1.tex に保存してある。
% ==================================================================
% ==================================================================
%  preamble.tex : 『素粒子論のための代数トポロジー演習』共通プリアンブル
%  ビルド：upLaTeX (jsbook) + dvipdfmx。本体は test.tex。
% ==================================================================
\documentclass[uplatex, dvipdfmx, 11pt]{jsbook}

\usepackage[utf8]{inputenc}
\usepackage[margin=30truemm]{geometry}
% --- 数学 ---
\usepackage{amsmath,amssymb,amsthm,amsfonts}
\usepackage{mathtools}
\usepackage{physics}      % \tr, \Tr, \dd, \pdv など
\usepackage{bm}
\usepackage{tikz-cd}      % 可換図式
\usepackage{slashed}      % ディラック作用素 \slashed{D}
\numberwithin{equation}{chapter}
% --- 図・色 ---
\usepackage{graphicx, color}
\usepackage{tikz}
\usetikzlibrary{intersections,calc,arrows.meta,decorations.markings,decorations.pathreplacing,patterns,3d}
\usepackage{float}
\usepackage{ascmac}
\usepackage{environ}
\usepackage{showkeys}
\parindent = 0pt

% ------------------------------------------------------------------
% 図用のスタイル（辺の同一視の矢印など）
% ------------------------------------------------------------------
\tikzset{
  mid arrow/.style={postaction={decorate},decoration={markings,
      mark=at position .55 with {\arrow{Stealth[length=2.2mm]}}}},
  mid arrows/.style={postaction={decorate},decoration={markings,
      mark=at position .48 with {\arrow{Stealth[length=2.2mm]}},
      mark=at position .60 with {\arrow{Stealth[length=2.2mm]}}}},
  vtx/.style={circle,fill,inner sep=1.1pt},
  hvtx/.style={circle,draw,fill=white,inner sep=1.3pt},
}
% 図の下に付ける説明（solution 環境の中でも使えるよう float は使わない）
\newcommand{\zucap}[1]{\par\nopagebreak{\centering\footnotesize #1\par}\medskip}

% ------------------------------------------------------------------
% 用語（将来の英語化に備えて英語を併記する）
% ------------------------------------------------------------------
\newcommand{\term}[2]{\textbf{#1}（#2）}   % 初出の用語：太字＋（English）
\newcommand{\en}[1]{（#1）}                % 補足の英語

% 計算過程などを書く注釈
\newenvironment{note}{\vspace{1ex}\par\footnotesize\noindent\textbf{【注】}}{\par\vspace{1ex}}

% ------------------------------------------------------------------
% 記号
% ------------------------------------------------------------------
\newcommand{\id}{\mathrm{id}}
\newcommand{\Z}{\mathbb{Z}}
\newcommand{\R}{\mathbb{R}}
\newcommand{\C}{\mathbb{C}}
\newcommand{\Q}{\mathbb{Q}}
\newcommand{\HH}{\mathbb{H}}               % 四元数
\newcommand{\RP}{\mathbb{RP}}
\newcommand{\CP}{\mathbb{CP}}
\DeclareMathOperator{\Hom}{Hom}
\DeclareMathOperator{\Ext}{Ext}
\DeclareMathOperator{\Tor}{Tor}
\DeclareMathOperator{\coker}{coker}
\DeclareMathOperator{\ind}{ind}
\DeclareMathOperator{\sign}{sign}
\DeclareMathOperator{\ch}{ch}
\DeclareMathOperator{\Td}{Td}
\DeclareMathOperator{\Pf}{Pf}
\DeclareMathOperator{\Aut}{Aut}
\DeclareMathOperator{\End}{End}
\DeclareMathOperator{\Ad}{Ad}
\DeclareMathOperator{\ad}{ad}
\DeclareMathOperator{\Spin}{Spin}
\DeclareMathOperator{\Pin}{Pin}
\newcommand{\Ahat}{\hat{A}}

% ------------------------------------------------------------------
% ハイパーリンク
% ------------------------------------------------------------------
\usepackage[dvipdfmx,bookmarks=true,bookmarksnumbered=true,bookmarkstype=toc,
  colorlinks=true,linkcolor=blue,citecolor=blue,urlcolor=blue]{hyperref}
\usepackage{pxjahyper}

% ------------------------------------------------------------------
% 定理系の囲み（tcolorbox）
%   使い方： \begin{mythm}{タイトル}{ラベル} ... \end{mythm}
%   ラベルは thm:ラベル, def:ラベル, ... として参照できる
% ------------------------------------------------------------------
\usepackage[many]{tcolorbox}
\tcbset{
  myboxstyle/.style={
    enhanced, breakable,
    colback=white, coltitle=white, fonttitle=\bfseries,
    boxrule=0.8pt, arc=1.5mm,
    attach boxed title to top left={xshift=3mm, yshift=-2mm},
    top=3mm, before skip=3mm, after skip=3mm,
    description font=\normalfont\bfseries, separator sign={：}
  }
}
\newtcbtheorem[number within=chapter]{mythm}{定理}{myboxstyle,
  colframe=red!75!black,
  boxed title style={colback=red!50,colframe=red!75!black,boxrule=0.8pt,arc=1.5mm}}{thm}
\newtcbtheorem[use counter from=mythm]{mydef}{定義}{myboxstyle,
  colframe=blue!70!black,
  boxed title style={colback=blue!60!gray!80,colframe=blue!70!black,boxrule=0.8pt,arc=1.5mm}}{def}
\newtcbtheorem[use counter from=mythm]{mylem}{補題}{myboxstyle,
  colframe=green!50!black,
  boxed title style={colback=green!50!gray,colframe=green!50!black,boxrule=0.8pt,arc=1.5mm}}{lem}
\newtcbtheorem[use counter from=mythm]{myprop}{命題}{myboxstyle,
  colframe=orange!80!black,
  boxed title style={colback=orange!80!black,colframe=orange!80!black,boxrule=0.8pt,arc=1.5mm}}{prop}
\newtcbtheorem[use counter from=mythm]{myex}{例}{myboxstyle,
  colframe=gray!70!black,
  boxed title style={colback=gray!70!black,colframe=gray!70!black,boxrule=0.8pt,arc=1.5mm}}{ex}
% 物理との対応をまとめる囲み（番号なし）
\newtcolorbox{physbox}[1][]{enhanced, breakable, colback=yellow!4, colframe=orange!60!black,
  boxrule=0.6pt, arc=1.5mm, fonttitle=\bfseries, title={物理との対応#1}}
% 旧版互換：章冒頭の「この章で用いる事項」
\newtcolorbox{summary}{breakable, enhanced, colback=gray!5, colframe=gray!60!black,
  boxrule=0.6pt, arc=1.5mm, fonttitle=\bfseries, title=この章で用いる事項}

% ------------------------------------------------------------------
% 問題と解答
%   problem 環境は本文に置き、solution 環境の中身は \printsolutions で巻末にまとめて出力する
%   注意：solution の中では # と verbatim を使わない（中身をマクロに蓄えるため）
% ------------------------------------------------------------------
\theoremstyle{definition}
\newtheorem{problem}{問題}[chapter]
\renewcommand{\proofname}{証明}
\newcommand{\allsolutions}{}
\newcommand{\currentchaptertitle}{}
\let\origchaptermark\chaptermark
\renewcommand{\chaptermark}[1]{\gdef\currentchaptertitle{#1}\origchaptermark{#1}}
\newcounter{solchapter}
\NewEnviron{solution}{%
  \ifnum\value{chapter}>\value{solchapter}%
    \setcounter{solchapter}{\value{chapter}}%
    \xdef\allsolutions{\unexpanded\expandafter{\allsolutions}%
      \noexpand\solutionchapter{\thechapter}{\unexpanded\expandafter{\currentchaptertitle}}}%
  \fi
  \xdef\allsolutions{\unexpanded\expandafter{\allsolutions}%
    \noexpand\solutionitem{\theproblem}{\unexpanded\expandafter{\BODY}}}%
}
\newcommand{\solutionchapter}[2]{\section*{第#1章\quad #2}}
\newcommand{\solutionitem}[2]{\par\medskip\noindent\textbf{問題 #1 の解答}\par\nopagebreak #2\qed\par}
\newcommand{\printsolutions}{%
  \chapter*{解答}
  \addcontentsline{toc}{chapter}{解答}
  \markboth{解答}{解答}
  \allsolutions
}


\title{
  素粒子論のための代数トポロジー演習\\[1ex]
  \large ゲージ理論・位相欠陥・アノマリー・弦理論のトポロジー\\
  \large （全200問）
}
\author{}
\date{}

\begin{document}

\frontmatter
\maketitle
\tableofcontents

%================================================
\chapter*{はじめに}
\addcontentsline{toc}{chapter}{はじめに}
%================================================

\section*{対象とねらい}
本書は、素粒子論を専攻する修士1年程度の読者を想定した代数トポロジーの演習書である。前提とするのは
\begin{itemize}
  \item 場の量子論の基礎（経路積分、ゲージ理論、自発的対称性の破れ、フェルミオンの経路積分）
  \item リー群・リー代数の初歩（$SU(N)$, $SO(N)$ とその表現）
  \item 一般相対論の初歩（多様体、微分形式、計量）
  \item 学部程度の集合と位相（開集合、連続写像、コンパクト性、連結性）
\end{itemize}
であり、代数トポロジーそのものの予備知識は仮定しない。

素粒子論でトポロジーが必要になるのは、おおよそ次の場面である。
\begin{enumerate}
  \item \textbf{位相欠陥とソリトン}：渦（宇宙ひも）、モノポール、ドメインウォール、インスタントン、スカーミオンの「位相的電荷」は、真空多様体 $G/H$ や群 $G$ のホモトピー群の元である（第1--3章）。
  \item \textbf{ゲージ場の大域構造}：ディラックの量子化条件、$\theta$ 真空、大きなゲージ変換、チャーン・サイモンズ項の量子化は、ファイバー束とその特性類の言葉で理解される（第5--8章）。
  \item \textbf{アノマリー}：カイラル・アノマリー、フェルミオンのゼロモード、ウィッテンの $SU(2)$ アノマリー、アノマリーの降下方程式は、指数定理と特性類によって統一的に記述される（第9章）。
  \item \textbf{コンパクト化}：余剰次元の多様体のベッチ数・ホッジ数が質量ゼロの場の数や世代数を決め、Dブレーンの電荷はコホモロジー（さらにK理論）で分類される（第10, 11章）。
\end{enumerate}
本書では、これらの物理を「動機」と「応用問題」として前面に出しつつ、そのために必要な数学（ホモトピー、ホモロジー、コホモロジー、ファイバー束、特性類、指数定理）を演習を通して身につけることを目標とする。

\section*{本書の構成と使い方}
各章は次の3節からなる。
\begin{enumerate}
  \item \textbf{物理からの動機}：その章の数学がどの物理の問いに答えるのかを述べる。
  \item \textbf{基本事項のまとめ}：定義・定理を囲み（定義・定理・命題・例）で整理する。証明は要点のみとし、計算の途中経過は【注】に書く。
  \item \textbf{問題}：計算問題と証明問題、物理への応用問題からなる。解答は巻末にまとめた。
\end{enumerate}
各問題には難易度を A〜D の4段階で付した。
\begin{itemize}
  \item A：定義の確認や、ほぼ1行で済む基本問題
  \item B：標準的な定理を1つ用いる計算・証明
  \item C：複数の道具を組み合わせる、またはやや長い計算・証明
  \item D：発展的な内容や、長い証明・工夫を要する問題
\end{itemize}
数学的な厳密さについては、「物理の議論に使える程度に正確で、誤りのないこと」を基準とした。点集合位相の細部（商写像の連続性など）は必要最小限にとどめ、そのかわり具体的な計算と物理的な意味づけを厚くしている。

\section*{章の依存関係}
\begin{center}
\begin{tikzpicture}[x=1.0cm,y=0.95cm,
    box/.style={draw,rounded corners=2pt,fill=blue!5,align=center,font=\footnotesize,minimum width=2.9cm,minimum height=0.8cm},
    arr/.style={-{Stealth[length=2mm]},thick,gray!70!black}]
  \node[box] (c1) at (0,6)  {第1章\\ホモトピーとリー群};
  \node[box] (c2) at (0,4.3)  {第2章\\基本群と被覆};
  \node[box] (c3) at (0,2.6)  {第3章\\高次ホモトピーと欠陥};
  \node[box] (c4) at (4.2,6) {第4章\\ホモロジー};
  \node[box] (c5) at (4.2,4.3) {第5章\\微分形式とド・ラーム};
  \node[box] (c6) at (4.2,2.6) {第6章\\コホモロジー環と双対性};
  \node[box] (c7) at (8.4,6) {第7章\\ファイバー束};
  \node[box] (c8) at (8.4,4.3) {第8章\\接続・曲率・特性類};
  \node[box] (c9) at (8.4,2.6) {第9章\\指数定理とアノマリー};
  \node[box] (c10) at (4.2,0.6) {第10章\\コンパクト化};
  \node[box] (c11) at (8.4,0.6) {第11章\\モース理論・SUSY QM・K理論};
  \draw[arr] (c1) -- (c2); \draw[arr] (c2) -- (c3);
  \draw[arr] (c1) -- (c4); \draw[arr] (c4) -- (c5); \draw[arr] (c5) -- (c6);
  \draw[arr] (c2) -- (c7); \draw[arr] (c3) -- (c7);
  \draw[arr] (c7) -- (c8); \draw[arr] (c5) -- (c8); \draw[arr] (c8) -- (c9); \draw[arr] (c6) -- (c9);
  \draw[arr] (c6) -- (c10); \draw[arr] (c9) -- (c10);
  \draw[arr] (c9) -- (c11); \draw[arr] (c4) .. controls (6.3,3.5) and (6.3,1.5) .. (c11);
\end{tikzpicture}
\end{center}

%================================================
\chapter*{記法と規約}
\addcontentsline{toc}{chapter}{記法と規約}
%================================================
本書全体で次の規約を用いる。章ごとに断りなく変更することはない。

\section*{一般}
\begin{itemize}
  \item 自然単位系 $\hbar=c=1$。ミンコフスキー計量は $\eta=\operatorname{diag}(-1,+1,\dots,+1)$（mostly plus）。インスタントンなどはユークリッド計量 $\delta_{\mu\nu}$ で扱う。
  \item $I=[0,1]$。$S^n=\{x\in\R^{n+1}:\lvert x\rvert=1\}$、$D^n=\{x\in\R^n:\lvert x\rvert\le1\}$、$\RP^n=S^n/(x\sim-x)$、$\CP^n=(\C^{n+1}\setminus0)/\C^\times$、$T^n=(S^1)^n$。
  \item 道の積 $f\cdot g$ は $f$ を先にたどる。$\pi_1(X,x_0)$ の積はこれで定める。
  \item 連続写像 $X\to Y$ のホモトピー類全体を $[X,Y]$ と書く（基点付きのときは $[X,Y]_*$）。
  \item ホモロジー・コホモロジーの係数を省略したときは $\Z$ 係数とする。
\end{itemize}

\section*{ゲージ場}
\begin{itemize}
  \item リー代数の生成子はエルミート：$[T^a,T^b]=if^{abc}T^c$。$SU(N)$ の基本表現では $\tr(T^aT^b)=\tfrac12\delta^{ab}$。
  \item ゲージ場は $A=A_\mu dx^\mu=A_\mu^aT^a\,dx^\mu$（エルミート行列値1形式）で、結合定数は $A$ に吸収する。共変微分と場の強さは
  \[
  D=d-iA,\qquad F=dA-iA\wedge A=\tfrac12F_{\mu\nu}dx^\mu\wedge dx^\nu .
  \]
  \item ゲージ変換 $g\colon M\to G$ に対し
  \[
  \psi\mapsto g\psi,\qquad A\mapsto gAg^{-1}+i\,g\,dg^{-1},\qquad F\mapsto gFg^{-1}.
  \]
  \item $U(1)$ の場合、電荷 $q\in\Z$（単位電荷の整数倍）の場の共変微分は $d-iqA$ である。ディラックの量子化条件は、閉曲面 $\Sigma$ 上で $\displaystyle\int_\Sigma F\in2\pi\Z$。
\end{itemize}

\section*{特性類}
\begin{itemize}
  \item エルミートな曲率 $F$ に対し
  \[
  c(F)=\det\Bigl(1+\frac{F}{2\pi}\Bigr),\qquad \ch(F)=\tr\exp\Bigl(\frac{F}{2\pi}\Bigr).
  \]
  特に $c_1=\dfrac{1}{2\pi}\tr F$、$\ch_2=\dfrac{1}{8\pi^2}\tr F\wedge F$、$SU(N)$ では $c_2=-\ch_2$。
  \item 4次元のインスタントン数を
  \[
  k=\frac{1}{8\pi^2}\int_M\tr F\wedge F=\int_M\ch_2(F)\quad(=-\textstyle\int_Mc_2\ \text{for}\ SU(N))
  \]
  で定める。$\tr(F\wedge *F)\ge0$ なので、自己双対 $F=*F$（向きは $dx^1\wedge dx^2\wedge dx^3\wedge dx^4$）なら $k\ge0$。
  \item チャーン・サイモンズ形式 $\displaystyle CS(A)=\frac{1}{8\pi^2}\tr\Bigl(A\wedge dA-\frac{2i}{3}A\wedge A\wedge A\Bigr)$、$dCS(A)=\dfrac{1}{8\pi^2}\tr F\wedge F$。
  \item 写像 $g\colon S^3\to SU(2)$ の巻き付き数 $\displaystyle w(g)=\frac{1}{24\pi^2}\int_{S^3}\tr(g^{-1}dg)^3$。積分の値は $S^3$ の向きに依存するので注意する。$\R^4$ の球の境界としての向き（外向き法線が先）で積分すると、$g(x)=(x_4+i\bm{x}\cdot\bm{\sigma})/\lvert x\rvert$ は $w=-1$ となる（逆向きなら $+1$）。境界としての向きを用い、遠方で $A\to i\,g\,dg^{-1}$ のとき、インスタントン数は $k=-w(g)$ である（BPST インスタントンで $k=+1$）。純ゲージのチャーン・サイモンズ数は $\int CS(i\,g\,dg^{-1})=-w(g)$。どの向きで積分しているかは各章で明記する（第3, 7, 8章）。
\end{itemize}

\section*{指数}
\begin{itemize}
  \item 偶数次元のディラック作用素 $\slashed{D}$ について、正カイラリティ・負カイラリティのゼロモードの数を $n_\pm$ とし、$\ind\slashed{D}=n_+-n_-$。
  \item カイラリティの符号は、アティヤ・シンガーの指数定理が $\displaystyle\ind\slashed{D}=\int_M\ch(F)\,\Ahat(TM)$ の形になるように選ぶ。特に $S^2$ 上で $\int F/2\pi=n$ なら $\ind=n$、平坦な4次元時空（の1点コンパクト化）上で $\ind=k$。
\end{itemize}

\mainmatter
%================================================
\chapter{位相空間・ホモトピー・リー群}\label{ch:01}
% 18問
%================================================

\section{物理からの動機}

\subsection*{場の配位は写像である}
2次元空間の複素スカラー場 $\phi(\bm{x})$ で、$U(1)$ 対称なポテンシャル
\[
V(\phi)=\lambda\bigl(\lvert\phi\rvert^2-v^2\bigr)^2
\]
をもつものを考える。静的な配位のエネルギーは
\[
E[\phi]=\int d^2x\,\Bigl(\lvert\nabla\phi\rvert^2+V(\phi)\Bigr)
\]
である。$E<\infty$ であるためには、遠方 $r\to\infty$ で $V(\phi)\to0$、すなわち $\phi$ が\term{真空多様体}{vacuum manifold}
\[
\mathcal{M}=\{\phi\in\C:\lvert\phi\rvert=v\}\cong S^1
\]
に近づかなければならない。したがって十分大きな半径 $R$ の円周上で $\phi$ は写像
\[
\phi_\infty\colon S^1_\infty\longrightarrow\mathcal{M}\cong S^1,\qquad \theta\longmapsto\phi(R,\theta)
\]
を定める（下図）。
\begin{center}
\begin{tikzpicture}[scale=1.1]
  \draw[-{Stealth[length=2mm]}] (-2.1,0) -- (2.3,0) node[right] {$\mathrm{Re}\,\phi$};
  \draw[-{Stealth[length=2mm]}] (0,-0.3) -- (0,2.4) node[above] {$V$};
  \draw[very thick,blue!70!black,domain=-1.55:1.55,samples=80,smooth]
     plot (\x,{0.9*(\x*\x-1)*(\x*\x-1)});
  \draw[thick,red!70!black,dashed] (0,0) ellipse (1 and 0.22);
  \node[vtx,red!70!black] at (1,0) {}; \node[vtx,red!70!black] at (-1,0) {};
  \node[red!70!black,below right] at (1.05,-0.05) {\footnotesize $\lvert\phi\rvert=v$};
  \node[blue!70!black] at (1.9,1.6) {\footnotesize $V(\phi)$};
  \begin{scope}[xshift=5.6cm,yshift=1cm]
    \fill[gray!25] (0,0) circle (0.35);
    \draw[thick] (0,0) circle (1.5);
    \foreach \a in {0,45,...,315}{
      \draw[-{Stealth[length=1.8mm]},thick,red!70!black] (\a:1.5) -- ($(\a:1.5)+(\a:0.45)$);
    }
    \node at (0,0) {\footnotesize 芯};
    \node at (0,-1.95) {\footnotesize $\phi(R,\theta)=v\,e^{i\theta}$（$n=1$）};
  \end{scope}
\end{tikzpicture}
\zucap{左：メキシカン・ハット型ポテンシャルの断面。底の円周 $\lvert\phi\rvert=v$ が真空多様体。右：遠方の円周上で $\phi$ の位相（矢印の向き）が1回まわる配位。中心付近（芯）では $\phi$ は真空からはずれる。}
\end{center}

時間発展は（エネルギー有限のまま）連続的に起こるので、写像 $\phi_\infty$ も連続的に変形される。連続的な変形で移り合う写像は「同じもの」とみなすのが自然であり、この同値関係が\term{ホモトピー}{homotopy}である。上の例では、$\phi_\infty$ が円周を何回まわるかという整数（\term{巻き付き数}{winding number}、第2章）は連続変形で変わらないので、\textbf{保存量}になる。これが\term{位相的電荷}{topological charge}の原型であり、渦（ボルテックス）や宇宙ひもの安定性を保証する。

\subsection*{一般の状況}
一般に、コンパクトなリー群 $G$ の作用で不変なポテンシャルが、点 $\phi_0$ で最小値をとり、その固定部分群が $H$ であるとすると、真空多様体は $G$ 軌道
\[
\mathcal{M}=G\cdot\phi_0\cong G/H
\]
である（\textbf{軌道・固定部分群定理}、本章）。$d$ 次元空間で有限エネルギーの配位は、遠方の球面 $S^{d-1}_\infty$ から $G/H$ への写像を定め、そのホモトピー類が位相的電荷を与える。どのホモトピー類がありうるかは、$G/H$ の\textbf{ホモトピー群}で決まる（第2, 3章）。

ユークリッド4次元のゲージ理論では、作用有限の条件から、ゲージ場は遠方で純ゲージ $A\to i\,g\,dg^{-1}$ となり、写像 $g\colon S^3_\infty\to G$ が現れる（インスタントン、第3, 8章）。また $\R^n$ に無限遠点を付け加えた空間は $S^n$ と同相なので（立体射影）、「遠方で一定値に近づく配位」は球面 $S^n$ 上の配位とみなせる。

非線形シグマ模型 $\phi\colon(\text{時空})\to\mathcal{M}$ のように、場の値そのものが曲がった空間（\term{標的空間}{target space}）に値をとる理論もある。カイラル・ラグランジアンの標的は $SU(N)$、$O(3)$ 模型の標的は $S^2$、$\CP^{N-1}$ 模型の標的は $\CP^{N-1}$ である。

\begin{physbox}[：本章の辞書]
\begin{center}
\small
\begin{tabular}{p{0.5\linewidth}p{0.4\linewidth}}
物理 & 数学\\ \hline
場の配位 & 写像 $X\to\mathcal{M}$\\
有限エネルギー（有限作用）の境界条件 & 遠方の球面からの写像、あるいは $\R^n\cup\{\infty\}\cong S^n$\\
エネルギー有限の連続的な時間発展 & ホモトピー\\
位相的電荷 & ホモトピー類\\
真空多様体 & 等質空間 $G/H$\\
ゲージ群・対称性の群 & リー群（多様体としての形：$SU(2)\cong S^3$ など）
\end{tabular}
\end{center}
\end{physbox}

\section{基本事項のまとめ}

\subsection{位相空間と商空間}
本書では、点集合位相の細部には立ち入らず、次のことを用いる。

\begin{mydef}{位相空間・連続写像・同相}{ch01-top}
集合 $X$ の部分集合族 $\mathcal{O}$ で、$\emptyset, X\in\mathcal{O}$、任意個の和と有限個の共通部分で閉じたものを\term{位相}{topology}、その元を開集合という。写像 $f\colon X\to Y$ が\term{連続}{continuous}であるとは、開集合の逆像が開集合であることをいう。連続な全単射で逆も連続なものを\term{同相写像}{homeomorphism}といい、$X\cong Y$ と書く。
\end{mydef}
$\R^n$ の部分集合には距離から定まる位相を入れる。\textbf{コンパクト}（有界閉集合）、\textbf{弧状連結}（任意の2点が道で結べる）などの概念は通常どおり用いる。次の事実はしばしば使う。

\begin{mythm}{コンパクトからハウスドルフへの連続全単射}{ch01-cpthaus}
$X$ がコンパクト、$Y$ がハウスドルフ（$\R^N$ の部分集合など）ならば、連続な全単射 $f\colon X\to Y$ は同相写像である。
\end{mythm}

\begin{mydef}{商空間}{ch01-quot}
空間 $X$ 上の同値関係 $\sim$ に対し、$q\colon X\to X/{\sim}$ を自然な射影とし、$U\subset X/{\sim}$ は $q^{-1}(U)$ が開集合のとき開集合と定める。これを\term{商空間}{quotient space}という。$X/{\sim}$ からの写像 $\bar f$ が連続であることは、$\bar f\circ q$ が連続であることと同値である。部分集合 $A\subset X$ を1点につぶした空間を $X/A$ と書く。
\end{mydef}

\begin{myex}{代表的な商空間}{ch01-quotex}
\begin{itemize}
\item $D^n/\partial D^n\cong S^n$（円板の境界を1点につぶすと球面）。
\item $\RP^n=S^n/(x\sim-x)$：原点を通る直線の空間。$\CP^n=S^{2n+1}/(z\sim\lambda z,\ \lvert\lambda\rvert=1)$：複素直線の空間。
\item トーラス $T^2=I^2/((0,y)\sim(1,y),\ (x,0)\sim(x,1))$。
\item $SO(3)\cong D^3/(\partial D^3\text{ の対蹠点を同一視})\cong\RP^3$（問題 \ref{ch01-prob-so3}）。
\end{itemize}
\end{myex}

\begin{myex}{立体射影と1点コンパクト化}{ch01-stereo}
北極 $N=(0,\dots,0,1)\in S^n$ からの\term{立体射影}{stereographic projection}
\[
\sigma\colon S^n\setminus\{N\}\to\R^n,\qquad \sigma(x)=\frac{(x_1,\dots,x_n)}{1-x_{n+1}}
\]
は同相写像で、逆写像は $\sigma^{-1}(y)=\dfrac{(2y,\ \lvert y\rvert^2-1)}{\lvert y\rvert^2+1}$ である。$\lvert y\rvert\to\infty$ が $x\to N$ に対応するので、$\R^n$ に1点 $\infty$ を付け加えた空間（\term{1点コンパクト化}{one-point compactification}）$\R^n\cup\{\infty\}$ は $S^n$ と同相である。
\begin{center}
\begin{tikzpicture}[scale=1.3]
  \draw[thick] (0,0) circle (1);
  \draw[-{Stealth[length=2mm]}] (-1.8,0) -- (3.3,0) node[right] {$\R^n$};
  \node[vtx] at (0,1) {}; \node[above] at (0,1) {$N\leftrightarrow\infty$};
  \draw[red!70!black,thick] (0,1) -- (2.75,0);
  \node[vtx,red!70!black] at (0.643,0.766) {}; \node[above right,red!70!black] at (0.64,0.78) {$x$};
  \node[vtx,red!70!black] at (2.75,0) {}; \node[below,red!70!black] at (2.75,0) {$\sigma(x)$};
  \node[vtx] at (0,-1) {}; \node[below] at (0,-1) {$S\leftrightarrow 0$};
\end{tikzpicture}
\zucap{立体射影。北極 $N$ と $x$ を結ぶ直線が赤道面と交わる点が $\sigma(x)$。北極に近づくほど $\sigma(x)$ は遠方へ行く。}
\end{center}
物理では「$\lvert x\rvert\to\infty$ で方向によらず一定値に近づく場」を、$S^n$ 上の連続な場とみなすことに使う。
\end{myex}

\subsection{ホモトピー}

\begin{mydef}{ホモトピー・ホモトピー同値}{ch01-htpy}
連続写像 $f,g\colon X\to Y$ が\term{ホモトピック}{homotopic}（$f\simeq g$）であるとは、連続写像 $H\colon X\times I\to Y$ で $H(\cdot,0)=f$、$H(\cdot,1)=g$ となるもの（\term{ホモトピー}{homotopy}）があることをいう。
連続写像 $f\colon X\to Y$ に対し $g\colon Y\to X$ で $g\circ f\simeq\id_X$、$f\circ g\simeq\id_Y$ となるものがあるとき、$f$ を\term{ホモトピー同値写像}{homotopy equivalence}、$X$ と $Y$ は\textbf{ホモトピー同値}であるといい $X\simeq Y$ と書く。1点とホモトピー同値な空間を\term{可縮}{contractible}という。
\end{mydef}

\begin{mydef}{変位レトラクト}{ch01-retract}
$A\subset X$ と連続写像 $H\colon X\times I\to X$ で、$H(x,0)=x$、$H(x,1)\in A$、$H(a,t)=a$（$a\in A$）をみたすものがあるとき、$A$ を $X$ の\term{強変位レトラクト}{strong deformation retract}という。このとき包含 $A\hookrightarrow X$ はホモトピー同値であり、$A\simeq X$。
\end{mydef}

\begin{mythm}{ホモトピーの基本性質}{ch01-htpyprop}
\begin{enumerate}
\item $\simeq$ は写像の集合上の同値関係で、合成と両立する：$f\simeq f'$, $g\simeq g'$ なら $g\circ f\simeq g'\circ f'$。同値類の集合を $[X,Y]$ と書く。
\item ホモトピー同値は空間の同値関係である。
\item 凸集合（特に $\R^n$, $D^n$）は可縮である。可縮な空間 $Y$ への写像はすべてホモトピックで、$[X,Y]$ は1点。
\end{enumerate}
\end{mythm}

\begin{myex}{ヒッグス場の零点を除いた空間}{ch01-punct}
$\R^n\setminus\{0\}$ は $H(x,t)=(1-t)x+t\,x/\lvert x\rvert$ により $S^{n-1}$ に強変位レトラクトするので $\R^n\setminus\{0\}\simeq S^{n-1}$（問題 \ref{ch01-prob-punct}）。物理的には、「零点をもたないヒッグス場」の値域は、動径方向を忘れれば真空多様体の形をしている。
\end{myex}

\subsection{CW複体}
多くの空間は、点（0胞体）に線分（1胞体）、円板（2胞体）、…を境界に沿って貼り合わせて作れる。これを\term{CW複体}{CW complex}という。$n$ 胞体 $D^n$ を貼る写像 $\varphi\colon\partial D^n=S^{n-1}\to X^{n-1}$ を\term{接着写像}{attaching map}と呼ぶ。
\begin{itemize}
\item $S^n=e^0\cup e^n$（$D^n$ の境界全体を1点に貼る）。
\item $\RP^n=e^0\cup e^1\cup\dots\cup e^n$（$e^k$ の接着写像は二重被覆 $S^{k-1}\to\RP^{k-1}$）。
\item $\CP^n=e^0\cup e^2\cup\dots\cup e^{2n}$。
\item 種数 $g$ の閉曲面 $\Sigma_g$ は $e^0$、$2g$ 個の $e^1$、1個の $e^2$。
\end{itemize}
ホモロジーの計算（第4章）では、この胞体構造が主役になる。

\subsection{リー群のトポロジー}
ゲージ群や対称性の群は、多様体としての形（トポロジー）をもつ。次は本書全体で繰り返し用いる。

\begin{myprop}{代表的なリー群の形}{ch01-liegroups}
\begin{enumerate}
\item $U(1)=\{e^{i\alpha}\}\cong S^1$。
\item $SU(2)=\{a_0+i\bm{a}\cdot\bm{\sigma}:a_0^2+\lvert\bm{a}\rvert^2=1\}\cong S^3$（単位四元数の群）。
\item $SO(3)\cong\RP^3$。$SU(2)\to SO(3)$ は2対1の準同型（第2章）。
\item $U(N)\cong SU(N)\times U(1)$（位相空間として。群としては $U(N)\cong(SU(N)\times U(1))/\Z_N$）。
\item $GL(n,\R)\simeq O(n)$、$GL(n,\C)\simeq U(n)$（極分解、グラム・シュミット）。
\item $SO(n)$, $SU(n)$, $U(n)$ は弧状連結、$O(n)$ は2つの連結成分（$\det=\pm1$）をもつ。ローレンツ群 $O(3,1)$ は4つの連結成分をもつ。
\end{enumerate}
\end{myprop}

\begin{center}
\begin{tikzpicture}[scale=1.2]
  \shade[ball color=blue!10,opacity=0.6] (0,0) circle (1.5);
  \draw[thick] (0,0) circle (1.5);
  \draw[dashed,gray] (1.5,0) arc (0:180:1.5 and 0.35);
  \draw[gray] (-1.5,0) arc (180:360:1.5 and 0.35);
  \draw[-{Stealth[length=2mm]},thick] (0,0) -- (40:1.05) node[above left] {\footnotesize $\psi\hat{\bm{n}}$};
  \node[vtx] at (0,0) {}; \node[below] at (0,-0.05) {\footnotesize $\mathbf 1$};
  \node[vtx,red!70!black] at (40:1.5) {}; \node[vtx,red!70!black] at (220:1.5) {};
  \node[red!70!black,above right] at (40:1.5) {\footnotesize $\pi\hat{\bm n}$};
  \node[red!70!black,below left] at (220:1.5) {\footnotesize $-\pi\hat{\bm n}$};
  \draw[red!70!black,densely dotted,thick] (40:1.55) .. controls (100:2.4) and (160:2.4) .. (220:1.55);
  \node at (2.6,0.9) {\footnotesize 半径 $\pi$ の球};
\end{tikzpicture}
\zucap{$SO(3)$ の軸・角度表示。回転軸 $\hat{\bm n}$ と回転角 $\psi\in[0,\pi]$ を、球の点 $\psi\hat{\bm n}$ で表す。角度 $\pi$ の回転は $\hat{\bm n}$ と $-\hat{\bm n}$ で同じなので、球面上の対蹠点（赤）を同一視する。}
\end{center}

\subsection{等質空間と真空多様体}

\begin{mythm}{軌道・固定部分群定理}{ch01-orbit}
コンパクトなリー群 $G$ がハウスドルフ空間 $X$ に連続に作用しているとする。$x_0\in X$ の\term{固定部分群}{stabilizer} $H=G_{x_0}=\{g:gx_0=x_0\}$ は閉部分群で、
\[
G/H\longrightarrow G\cdot x_0,\qquad gH\longmapsto g x_0
\]
は同相写像である。
\end{mythm}
\begin{note}
well-defined かつ単射であることは $gx_0=g'x_0\iff g^{-1}g'\in H$ から従う。$G/H$ はコンパクト（$G$ の連続像）で、軌道はハウスドルフなので、定理 \ref{thm:ch01-cpthaus} により同相である。
\end{note}

\begin{myex}{等質空間}{ch01-homog}
\begin{itemize}
\item $S^n=SO(n+1)/SO(n)$、$S^{2n-1}=U(n)/U(n-1)=SU(n)/SU(n-1)$。
\item $\CP^n=SU(n+1)/S(U(n)\times U(1))=U(n+1)/(U(n)\times U(1))$。
\item グラスマン多様体 $\mathrm{Gr}_k(\C^n)=U(n)/(U(k)\times U(n-k))$（$\C^n$ の $k$ 次元部分空間全体）。
\end{itemize}
\end{myex}

\begin{physbox}[：自発的対称性の破れと真空多様体]
\begin{center}
\renewcommand{\arraystretch}{1.2}\small
\begin{tabular}{p{0.36\linewidth}p{0.3\linewidth}l}
模型（破れ方） & $G\to H$ & 真空多様体 $G/H$\\ \hline
複素スカラー（メキシカン・ハット） & $U(1)\to1$ & $S^1$\\
$O(N)$ 模型 & $O(N)\to O(N-1)$ & $S^{N-1}$\\
随伴ヒッグス（ジョージ・グラショウ） & $SU(2)\to U(1)$ & $S^2$\\
電弱理論 & $SU(2)\times U(1)_Y\to U(1)_{\mathrm{em}}$ & $S^3$\\
カイラル対称性の破れ & $SU(N)_L\times SU(N)_R\to SU(N)_V$ & $SU(N)$\\
$SU(3)$ 随伴（$\langle\phi\rangle\propto\operatorname{diag}(1,1,-2)$） & $SU(3)\to S(U(2)\times U(1))$ & $\CP^2$
\end{tabular}
\end{center}
これらは問題で順に確かめる。どの位相欠陥が安定に存在するかは、第3章でこの表のホモトピー群を計算して決める。
\end{physbox}

\subsection{ホップ写像}
第3, 7章で主役になる写像を導入しておく。$S^3=\{z=(z_1,z_2)\in\C^2:\lvert z_1\rvert^2+\lvert z_2\rvert^2=1\}$ に対し
\[
h\colon S^3\to S^2,\qquad h(z)^a=z^\dagger\sigma^az\quad(a=1,2,3)
\]
を\term{ホップ写像}{Hopf map}という。成分で書くと $h(z)=\bigl(2\,\mathrm{Re}(\bar z_1z_2),\ 2\,\mathrm{Im}(\bar z_1z_2),\ \lvert z_1\rvert^2-\lvert z_2\rvert^2\bigr)$ である。量子力学では、スピン $1/2$ の規格化された状態 $z$ に、そのブロッホ・ベクトル $\langle\bm{\sigma}\rangle$ を対応させる写像にほかならない。各点の逆像は円周 $\{e^{i\alpha}z\}$（状態の全体位相）であり（問題 \ref{ch01-prob-hopf}）、$S^3$ は円周の族に分割される。

\section{問題}

\subsection{ホモトピーの基本}

\begin{problem}[難易度 A]\label{ch01-prob-equiv}
連続写像 $X\to Y$ の間の関係 $\simeq$（ホモトピック）が同値関係であること、および $f\simeq f'\colon X\to Y$、$g\simeq g'\colon Y\to Z$ ならば $g\circ f\simeq g'\circ f'$ であることを示せ。
\end{problem}

\begin{solution}
$I=[0,1]$ とする。
反射律：$H(x,t)=f(x)$。
対称律：$H$ が $f\simeq g$ を与えるなら、$H'(x,t)=H(x,1-t)$ が $g\simeq f$ を与える。
推移律：$H$ が $f\simeq g$、$K$ が $g\simeq h$ を与えるとき
\[
L(x,t)=\begin{cases}H(x,2t)&(0\le t\le\tfrac12)\\ K(x,2t-1)&(\tfrac12\le t\le1)\end{cases}
\]
とおく。$t=\tfrac12$ では $H(x,1)=g(x)=K(x,0)$ で一致し、$X\times[0,\tfrac12]$ と $X\times[\tfrac12,1]$ は閉集合なので、貼り合わせの補題により $L$ は連続で、$f\simeq h$ を与える。

合成：$H\colon f\simeq f'$、$K\colon g\simeq g'$ とすると、$g\circ H$ は $g\circ f\simeq g\circ f'$ を、$K(f'(x),t)$ は $g\circ f'\simeq g'\circ f'$ を与える。推移律より $g\circ f\simeq g'\circ f'$。
\end{solution}

\begin{problem}[難易度 A]\label{ch01-prob-convex}
(1) 空でない凸集合 $C\subset\R^n$ は可縮であり、任意の空間 $X$ からの連続写像 $f,g\colon X\to C$ はホモトピックであることを示せ。\\
(2) 可縮な空間は弧状連結であることを示せ。\\
(3) 対称性の破れがなく真空が一意（真空多様体が1点）な理論、あるいは真空多様体が可縮な理論では、遠方の境界条件による位相的電荷（$[S^{d-1},\mathcal{M}]$ の元）は自明であることを説明せよ。
\end{problem}

\begin{solution}
(1) $H(x,t)=(1-t)f(x)+t\,g(x)$ とおく。凸性より $H(x,t)\in C$ で、$H$ は連続、$H(\cdot,0)=f$、$H(\cdot,1)=g$。特に $X=C$、$f=\id_C$、$g$ を点 $c_0$ への定値写像とすれば $\id_C\simeq$（定値写像）であり、$C$ は可縮である。

(2) 可縮なら $H\colon X\times I\to X$ で $H(x,0)=x$、$H(x,1)=x_0$ となるものがある。$\gamma_x(t)=H(x,t)$ は $x$ から $x_0$ への道なので、任意の2点は $x_0$ を経由して道で結べる。

(3) $\mathcal{M}$ が可縮なら、$\id_{\mathcal{M}}\simeq c$（定値写像）。任意の $f\colon S^{d-1}\to\mathcal{M}$ に対し $f=\id\circ f\simeq c\circ f$ で、$c\circ f$ は定値写像である。よってすべての境界条件は互いにホモトピックで $[S^{d-1},\mathcal{M}]$ は1点、すなわち位相的に異なるセクターは存在しない。真空多様体が1点の場合はもちろん自明である。位相欠陥が存在するには、真空多様体が「穴」をもつ（可縮でない）ことが必要である。
\end{solution}

\begin{problem}[難易度 B]\label{ch01-prob-punct}
$n\ge1$ とする。$S^{n-1}$ が $\R^n\setminus\{0\}$ の強変位レトラクトであることを示し、$\R^n\setminus\{0\}\simeq S^{n-1}$ を結論せよ。これを用いて、$n$ 成分の実スカラー場 $\phi^a$（$O(n)$ 模型）で、ある領域 $\Omega$ 上で $\phi\neq0$ であるものは、写像 $\hat\phi=\phi/\lvert\phi\rvert\colon\Omega\to S^{n-1}$ と（ホモトピーの意味で）同じ情報をもつことを説明せよ。
\end{problem}

\begin{solution}
$X=\R^n\setminus\{0\}$ とし
\[
H(x,t)=(1-t)x+t\frac{x}{\lvert x\rvert}=\Bigl((1-t)+\frac{t}{\lvert x\rvert}\Bigr)x
\]
とおく。係数は正（$1-t$ と $t/\lvert x\rvert$ は非負で同時に $0$ にならない）なので $H(x,t)\neq0$、$H$ は $X\times I\to X$ の連続写像。$H(x,0)=x$、$H(x,1)=x/\lvert x\rvert\in S^{n-1}$、$\lvert a\rvert=1$ なら $H(a,t)=a$。よって強変位レトラクトで、$r(x)=x/\lvert x\rvert$ と包含 $i$ は $r\circ i=\id$、$i\circ r\simeq\id_X$ をみたし $X\simeq S^{n-1}$。

物理への応用：$\phi\colon\Omega\to X$ と $\hat\phi=r\circ\phi$ について、$i\circ\hat\phi=i\circ r\circ\phi\simeq\phi$。したがって $\phi\mapsto\hat\phi$ は全単射 $[\Omega,X]\cong[\Omega,S^{n-1}]$ を与える（逆は $\hat\phi\mapsto i\circ\hat\phi$）。つまり零点をもたない場のホモトピー類は、その「向き」$\hat\phi$ だけで決まり、大きさ $\lvert\phi\rvert$ は位相的に無関係である。零点（欠陥の芯）があるところで初めて位相的な情報が失われうる。
\end{solution}

\begin{problem}[難易度 B]\label{ch01-prob-stereo}
(1) 立体射影 $\sigma(x)=(x_1,\dots,x_n)/(1-x_{n+1})$ と $\tau(y)=\bigl(2y,\lvert y\rvert^2-1\bigr)/(\lvert y\rvert^2+1)$ が互いに逆写像であることを確かめ、$S^n\setminus\{N\}\cong\R^n$ を示せ。\\
(2) $\R^n$ 上の連続な場 $\phi\colon\R^n\to\mathcal{M}$ が、$\lvert x\rvert\to\infty$ で方向によらず一様に一定値 $\phi_\infty$ に収束するとする（$\sup_{\lvert x\rvert\ge R}d(\phi(x),\phi_\infty)\to0$）。$\phi$ は連続写像 $S^n\to\mathcal{M}$ に拡張されることを示せ。\\
(3) 4次元ユークリッド空間のスカラー場の配位で作用が有限になるものは、しばしば $S^4$ 上の場とみなされる。(2) の意味を述べよ。
\end{problem}

\begin{solution}
(1) $y=\sigma(x)$ とおくと $\lvert y\rvert^2=\dfrac{1-x_{n+1}^2}{(1-x_{n+1})^2}=\dfrac{1+x_{n+1}}{1-x_{n+1}}$。よって $\lvert y\rvert^2+1=\dfrac{2}{1-x_{n+1}}$、$\lvert y\rvert^2-1=\dfrac{2x_{n+1}}{1-x_{n+1}}$ であり
\[
\tau(\sigma(x))=\Bigl(\frac{2x_i/(1-x_{n+1})}{2/(1-x_{n+1})},\ \frac{2x_{n+1}/(1-x_{n+1})}{2/(1-x_{n+1})}\Bigr)=x .
\]
逆に $x=\tau(y)$ は $\lvert x\rvert^2=\dfrac{4\lvert y\rvert^2+(\lvert y\rvert^2-1)^2}{(\lvert y\rvert^2+1)^2}=1$ をみたし、$1-x_{n+1}=\dfrac{2}{\lvert y\rvert^2+1}\neq0$ なので $x\neq N$、$\sigma(\tau(y))=\dfrac{2y/(\lvert y\rvert^2+1)}{2/(\lvert y\rvert^2+1)}=y$。両者とも連続なので同相である。

(2) $\Phi(x)=\phi(\sigma(x))$（$x\neq N$）、$\Phi(N)=\phi_\infty$ と定める。$N$ 以外での連続性は明らか。$N$ での連続性：$x\to N$ は $\lvert\sigma(x)\rvert^2=(1+x_{n+1})/(1-x_{n+1})\to\infty$ と同値であり、仮定から $d(\Phi(x),\phi_\infty)\le\sup_{\lvert y\rvert\ge\lvert\sigma(x)\rvert}d(\phi(y),\phi_\infty)\to0$。

(3) 作用 $\int d^4x\,[\tfrac12(\partial\phi)^2+V(\phi)]$ が有限であるには、遠方で $\phi$ が真空に近づく必要がある。真空が一様な値 $\phi_\infty$ に（方向によらず）近づくなら、(2) により配位は $S^4=\R^4\cup\{\infty\}$ 上の連続な場とみなせ、そのホモトピー類 $[S^4,\mathcal{M}]$ が位相的電荷になる。なお、ゲージ理論のインスタントンでは「場そのもの」ではなく「ゲージ同値類」が遠方で一定になるだけなので、$S^4$ 上には自明でないゲージ束が現れる（第7章）。
\end{solution}

\subsection{リー群のかたち}

\begin{problem}[難易度 A]\label{ch01-prob-su2}
$2\times2$ 複素行列 $U$ が $U^\dagger U=\mathbf 1$、$\det U=1$ をみたすための必要十分条件は
\[
U=a_0\mathbf 1+i\,\bm a\cdot\bm\sigma,\qquad (a_0,\bm a)\in\R^4,\quad a_0^2+\lvert\bm a\rvert^2=1
\]
と書けることであることを示し、$SU(2)\cong S^3$ を結論せよ。
\end{problem}

\begin{solution}
$U=\begin{pmatrix}\alpha&\beta\\\gamma&\delta\end{pmatrix}$ とする。$\det U=1$ なら $U^{-1}=\begin{pmatrix}\delta&-\beta\\-\gamma&\alpha\end{pmatrix}$ であり、ユニタリ性 $U^{-1}=U^\dagger=\begin{pmatrix}\bar\alpha&\bar\gamma\\\bar\beta&\bar\delta\end{pmatrix}$ と比べて $\delta=\bar\alpha$、$\gamma=-\bar\beta$。さらに $\det U=\lvert\alpha\rvert^2+\lvert\beta\rvert^2=1$。逆にこの形の行列は $SU(2)$ に属する。

$\alpha=a_0+ia_3$、$\beta=a_2+ia_1$ とおくと
\[
a_0\mathbf 1+i(a_1\sigma_1+a_2\sigma_2+a_3\sigma_3)=\begin{pmatrix}a_0+ia_3&a_2+ia_1\\-a_2+ia_1&a_0-ia_3\end{pmatrix}=\begin{pmatrix}\alpha&\beta\\-\bar\beta&\bar\alpha\end{pmatrix}
\]
であり、条件は $a_0^2+a_1^2+a_2^2+a_3^2=1$。対応 $U\leftrightarrow(a_0,\bm a)$ は連続な全単射で逆も連続なので $SU(2)\cong S^3$。

物理的には、$SU(2)$ のゲージ変換や $N_f=2$ のカイラル場 $\Sigma\in SU(2)$ は「3次元球面に値をとる場」であり、写像 $S^3\to S^3$ の巻き付き数がインスタントン数やバリオン数になる（第3章）。
\end{solution}

\begin{problem}[難易度 B]\label{ch01-prob-so3}
$D=\{\bm v\in\R^3:\lvert\bm v\rvert\le\pi\}$ とし、$\bm v=\psi\hat{\bm n}$ に軸 $\hat{\bm n}$ まわりの角度 $\psi$ の回転 $R(\bm v)\in SO(3)$ を対応させる（$R(\bm 0)=\mathbf 1$）。ロドリゲスの公式
\[
R(\bm v)=\mathbf 1+\frac{\sin\psi}{\psi}K+\frac{1-\cos\psi}{\psi^2}K^2,\qquad K_{ij}=-\epsilon_{ijk}v_k,\quad\psi=\lvert\bm v\rvert
\]
を認めて、$R$ が同相写像 $D/(\bm v\sim-\bm v,\ \lvert\bm v\rvert=\pi)\cong SO(3)$ を誘導することを示せ。さらにこの空間が $\RP^3$ と同相であることを示せ。
\end{problem}

\begin{solution}
\textbf{連続性}：$\sin\psi/\psi$ と $(1-\cos\psi)/\psi^2$ は $\psi^2$ の整級数（$\psi=0$ で $1$, $\tfrac12$）なので $R$ は $D$ 上連続（実は滑らか）。

\textbf{全射}：任意の $R\in SO(3)$ は固有値 $1$ の固有ベクトル（軸）をもち、軸に垂直な平面内の回転なので、ある $\hat{\bm n}$ と $\psi\in[0,\pi]$ で $R=R(\psi\hat{\bm n})$ と書ける（必要なら $\hat{\bm n}\to-\hat{\bm n}$ として角度を $[0,\pi]$ に収める）。

\textbf{同一視される点}：$\operatorname{tr}R(\bm v)=1+2\cos\psi$ なので $R$ から $\psi\in[0,\pi]$ が決まる。反対称部分は $\tfrac12(R-R^{\mathsf T})=\dfrac{\sin\psi}{\psi}K$ である（$K^2$ は対称）。$0<\psi<\pi$ なら $\sin\psi\neq0$ なので $K$、したがって $\bm v$ が $R$ から一意に決まる。$\psi=0$ は $R=\mathbf 1$ のみ。$\psi=\pi$ なら $R=\mathbf 1+2\hat K^2=2\hat{\bm n}\hat{\bm n}^{\mathsf T}-\mathbf 1$（$\hat K=K/\pi$, $\hat K^2=\hat{\bm n}\hat{\bm n}^{\mathsf T}-\mathbf 1$）で、$\hat{\bm n}$ は符号を除いて決まる。よって $R(\bm v)=R(\bm w)$ となるのは $\bm v=\bm w$ か、$\lvert\bm v\rvert=\pi$ かつ $\bm w=-\bm v$ のときに限る。

したがって $R$ は商空間 $D/{\sim}$ からの連続な全単射を誘導する。$D/{\sim}$ はコンパクト、$SO(3)\subset\R^9$ はハウスドルフなので同相（定理 \ref{thm:ch01-cpthaus}）。

\textbf{$\RP^3$ との同相}：$S^3$ の上半球 $S^3_+=\{x_4\ge0\}$ は $D^3$ と同相（$(\bm y,\sqrt{1-\lvert\bm y\rvert^2})\leftrightarrow\bm y$）で、$\RP^3=S^3/(x\sim-x)$ の各点は上半球に代表元をもち、二つの代表元をもつのは赤道 $x_4=0$ 上の対蹠点だけである。よって $\RP^3\cong D^3/(\partial D^3\text{ の対蹠点})$（同じくコンパクト→ハウスドルフの議論）。半径を $\pi$ に拡大すれば上の $D/{\sim}$ である。

この表示で、$\mathbf 1$ から $\pi\hat{\bm n}$ へ行き、対蹠点 $-\pi\hat{\bm n}$ から $\mathbf 1$ へ戻る閉曲線（$2\pi$ 回転）が縮められないことが、スピン $1/2$ の符号の起源になる（第2章）。
\end{solution}

\begin{problem}[難易度 B]\label{ch01-prob-un}
$N\ge2$ とする。\\
(1) $D(z)=\operatorname{diag}(z,1,\dots,1)$ として、$\Phi(U)=\bigl(U\,D(\det U)^{-1},\ \det U\bigr)$ が同相写像 $U(N)\cong SU(N)\times U(1)$ を与えることを示せ。\\
(2) $U(N)$ と $SU(N)\times U(1)$ は群として同型でないことを示せ。（ヒント：中心の中で $g^N=1$ をみたす元の個数を比べよ。）\\
(3) 準同型 $SU(N)\times U(1)\to U(N)$、$(V,e^{i\alpha})\mapsto e^{i\alpha}V$ の核を求め、$U(N)\cong(SU(N)\times U(1))/\Z_N$ を示せ。標準模型のゲージ群の「大域構造」がこの種の商で曖昧さをもつことにも注意せよ。
\end{problem}

\begin{solution}
(1) $\det\bigl(U D(\det U)^{-1}\bigr)=\det U\cdot(\det U)^{-1}=1$ なので $\Phi$ は $SU(N)\times U(1)$ に値をとる。$\Psi(V,z)=V\,D(z)$ とおくと $\det\Psi(V,z)=z$ で、$\Psi\circ\Phi(U)=U$、$\Phi\circ\Psi(V,z)=(V,z)$。$\Phi,\Psi$ はともに行列成分の多項式と $z^{-1}=\bar z$ で書けるので連続。よって同相。ただし $\Phi$ は準同型ではない（$D(\det U)$ が $U$ と可換とは限らない）。

(2) $U(N)$ の中心はスカラー行列 $\{e^{i\alpha}\mathbf 1\}\cong U(1)$ であり、その中で $g^N=1$ をみたす元は $e^{2\pi ik/N}\mathbf 1$ の $N$ 個。$SU(N)\times U(1)$ の中心は $Z(SU(N))\times U(1)=\Z_N\times U(1)$ で、$g^N=1$ をみたす元は $N\cdot N=N^2$ 個。群同型は中心を中心に写し、$g^N=1$ をみたす元の個数を保つので、$N\ge2$ では同型でない。

(3) $e^{i\alpha}V=\mathbf 1$ なら $V=e^{-i\alpha}\mathbf 1\in SU(N)$ なので $e^{-iN\alpha}=1$。核は $\{(\omega^{-1}\mathbf 1,\omega):\omega^N=1\}\cong\Z_N$。全射性：$U\in U(N)$ に対し $\det U=e^{iN\alpha}$ となる $\alpha$ をとれば $V=e^{-i\alpha}U\in SU(N)$。よって $U(N)\cong(SU(N)\times U(1))/\Z_N$（同型定理）。

標準模型では、すべての場の電荷がある $\Z_6$ 部分群の作用を自明にするため、ゲージ群は $SU(3)\times SU(2)\times U(1)_Y$ を $\Gamma\subset\Z_6$ で割ったもののどれか（$\Gamma=1,\Z_2,\Z_3,\Z_6$）でありうる。局所的な物理（摂動論）では区別できないが、許される線演算子やモノポールの電荷（第3, 7章）が変わる。
\end{solution}

\begin{problem}[難易度 C]\label{ch01-prob-gl}
(1) $A\in GL(n,\R)$ の列にグラム・シュミットの直交化を施して $A=Q(A)R(A)$（$Q\in O(n)$、$R$ は対角成分が正の上三角行列）と分解する。$Q,R$ が $A$ の連続関数であることを認めて、
\[
H(A,t)=Q(A)\bigl((1-t)R(A)+t\mathbf 1\bigr)
\]
が $O(n)$ を $GL(n,\R)$ の強変位レトラクトとすることを示せ。\\
(2) 同様に $GL(n,\C)\simeq U(n)$ を示し、$GL(n,\C)$ は弧状連結、$GL(n,\R)$ はちょうど2つの弧状連結成分をもつことを示せ（$SO(n)$, $U(n)$ の弧状連結性は用いてよい）。
\end{problem}

\begin{solution}
(1) $(1-t)R(A)+t\mathbf 1$ は対角成分 $(1-t)R_{kk}+t>0$ をもつ上三角行列なので正則であり、$H(A,t)\in GL(n,\R)$。$H$ は連続、$H(A,0)=QR=A$、$H(A,1)=Q(A)\in O(n)$。$A\in O(n)$ なら列はすでに正規直交なのでグラム・シュミットは何もせず $Q(A)=A$、$R(A)=\mathbf 1$、よって $H(A,t)=A$。以上で強変位レトラクトであり、$GL(n,\R)\simeq O(n)$。

(2) 複素内積でグラム・シュミットを行えば $A=Q(A)R(A)$（$Q\in U(n)$、$R$ は対角成分が正の上三角）となり、同じ $H$ が $GL(n,\C)\simeq U(n)$ を与える。

ホモトピー同値写像は弧状連結成分の集合の全単射を誘導する（$f\colon X\to Y$, $g\colon Y\to X$ で $g\circ f\simeq\id$ なら、各点 $x$ と $g(f(x))$ はホモトピーの軌跡で結ばれる）。$U(n)$ は弧状連結なので $GL(n,\C)$ も弧状連結。$O(n)=SO(n)\sqcup R_1SO(n)$（$R_1=\operatorname{diag}(-1,1,\dots,1)$）で、各成分は弧状連結、かつ連続関数 $\det$ が値 $\pm1$ をとるので両者は結べない。よって $O(n)$、したがって $GL(n,\R)$ はちょうど2つの成分をもつ（$\det>0$ と $\det<0$）。

物理的には、$GL(4,\R)$ の2成分は向きを保つ／反転する座標変換に対応し、一般相対論の「一般座標変換」でも向きの反転は連続変形では得られない。
\end{solution}

\begin{problem}[難易度 A]\label{ch01-prob-lorentz}
(1) $O(n)$ は少なくとも2つの弧状連結成分をもつことを示せ。\\
(2) ローレンツ群 $O(3,1)=\{\Lambda:\Lambda^{\mathsf T}\eta\Lambda=\eta\}$（$\eta=\operatorname{diag}(-1,1,1,1)$）について、$\det\Lambda=\pm1$ かつ $\lvert\Lambda^0{}_0\rvert\ge1$ を示し、$O(3,1)$ が少なくとも4つの弧状連結成分をもつことを示せ。（固有直交時間順ローレンツ群 $SO^+(3,1)$ が弧状連結であることを認めれば、ちょうど4つである。）パリティ $P$ や時間反転 $T$ が「連続的な対称性の極限」として得られないことを説明せよ。
\end{problem}

\begin{solution}
(1) $\det\colon O(n)\to\{\pm1\}$ は連続で、$\mathbf 1$ と $\operatorname{diag}(-1,1,\dots,1)$ で異なる値をとる。道 $\gamma$ で両者が結ばれれば $\det\circ\gamma\colon[0,1]\to\{\pm1\}$ は連続で $1$ から $-1$ に変わることになり、中間値の定理に反する。

(2) $\det(\Lambda^{\mathsf T}\eta\Lambda)=\det\eta$ より $(\det\Lambda)^2=1$。$(0,0)$ 成分を見ると $-1=\eta_{00}=\Lambda^\mu{}_0\eta_{\mu\nu}\Lambda^\nu{}_0=-(\Lambda^0{}_0)^2+\sum_i(\Lambda^i{}_0)^2$ なので $(\Lambda^0{}_0)^2=1+\sum_i(\Lambda^i{}_0)^2\ge1$。

したがって連続写像 $\Lambda\mapsto(\det\Lambda,\ \operatorname{sgn}\Lambda^0{}_0)\in\{\pm1\}\times\{\pm1\}$ が定まり（$\Lambda^0{}_0$ は $0$ にならないので符号は連続）、$\mathbf 1$, $P=\operatorname{diag}(1,-1,-1,-1)$, $T=\operatorname{diag}(-1,1,1,1)$, $PT$ がそれぞれ $(1,1)$, $(-1,1)$, $(-1,-1)$, $(1,-1)$ をとる。(1) と同じ議論で4つの元は互いに道で結べず、少なくとも4成分。$SO^+(3,1)$ が弧状連結なら、各成分はその剰余類 $SO^+$, $P\,SO^+$, $T\,SO^+$, $PT\,SO^+$ でちょうど4つ。

連続的なローレンツ変換（ブーストと回転を連続的に合成したもの）は $\mathbf 1$ と道で結ばれているので $SO^+(3,1)$ に属する。$P$ や $T$ は別の成分にあるので、どんな連続変形の極限としても得られない。したがってローレンツ不変性（連続対称性）から $P$ や $T$ の不変性は従わず、実際、弱い相互作用はこれらを破る（$CPT$ 定理は別の論理による）。
\end{solution}

\subsection{等質空間と真空多様体}

\begin{problem}[難易度 B]\label{ch01-prob-sphere}
(1) $n\ge1$ のとき、$SO(n+1)$ は $S^n$ に推移的に作用し、北極 $e_{n+1}$ の固定部分群は $SO(n)$（左上ブロック）であることを示し、$S^n\cong SO(n+1)/SO(n)$ を結論せよ。\\
(2) 同様に $S^{2n-1}\cong SU(n)/SU(n-1)$（$n\ge2$）を示せ。\\
(3) $O(N)$ 模型 $\mathcal L=-\tfrac12\partial_\mu\phi^a\partial^\mu\phi^a-\lambda(\phi^a\phi^a-v^2)^2$ の真空多様体と、対称性の破れ方 $G\to H$ を答えよ。
\end{problem}

\begin{solution}
(1) 単位ベクトル $u\in S^n$ に対し、$u$ を最後の列にもつ正規直交基底 $(u_1,\dots,u_n,u)$ をとる（グラム・シュミット）。行列式が $-1$ なら $u_1$ を $-u_1$ に替えれば（$n\ge1$ なので可能）$g\in SO(n+1)$ で $g e_{n+1}=u$。よって推移的。$ge_{n+1}=e_{n+1}$ なら $g$ の最後の列は $e_{n+1}$、直交性から最後の行も $e_{n+1}^{\mathsf T}$ で、$g=\operatorname{diag}(g',1)$、$g'\in SO(n)$。軌道・固定部分群定理により $S^n\cong SO(n+1)/SO(n)$。

(2) $u\in S^{2n-1}\subset\C^n$ を最後の列とするユニタリ行列を作り（複素グラム・シュミット）、最初の列に位相 $\overline{\det}$ を掛けて行列式を $1$ にすれば $g\in SU(n)$、$ge_n=u$。固定部分群は $\operatorname{diag}(g',1)$、$g'\in SU(n-1)$。よって $S^{2n-1}\cong SU(n)/SU(n-1)$。

(3) ポテンシャルの最小点は $\phi^a\phi^a=v^2$ で、真空多様体は半径 $v$ の $S^{N-1}$。$G=O(N)$ が推移的に作用し、$\phi_0=(0,\dots,0,v)$ の固定部分群は $O(N-1)$。$O(N)\to O(N-1)$ と破れ、$\mathcal M=O(N)/O(N-1)\cong S^{N-1}$（$SO$ で考えても同じ）。南部・ゴールドストーン粒子の数は $\dim\mathcal M=N-1$ である。
\end{solution}

\begin{problem}[難易度 B]\label{ch01-prob-ew}
電弱理論のヒッグス2重項 $H=(h_1,h_2)^{\mathsf T}\in\C^2$ はポテンシャル $V=\lambda\bigl(H^\dagger H-\tfrac{v^2}{2}\bigr)^2$ をもち、$SU(2)_L\times U(1)_Y$ の元 $(U,e^{i\beta})$ は $H\mapsto e^{i\beta/2}UH$（超電荷 $Y=\tfrac12$）と作用する。\\
(1) 真空多様体が $S^3$ であることを示せ。\\
(2) $H_0=(0,v/\sqrt2)^{\mathsf T}$ の固定部分群を求め、それが $Q=T_3+Y$（$T_3=\sigma_3/2$）で生成される $U(1)_{\mathrm{em}}$ であることを示せ。\\
(3) $SU(2)\times U(1)_Y$ のこの作用は忠実でないことを示し、実際に作用している群が $U(2)$ であることを説明せよ。
\end{problem}

\begin{solution}
(1) 最小点は $H^\dagger H=v^2/2$、すなわち $\lvert h_1\rvert^2+\lvert h_2\rvert^2=v^2/2$ で、$\C^2=\R^4$ の半径 $v/\sqrt2$ の球面 $S^3$。$SU(2)$ だけで推移的である：$U=\begin{pmatrix}\alpha&\beta\\-\bar\beta&\bar\alpha\end{pmatrix}$ に対し $UH_0=\tfrac{v}{\sqrt2}(\beta,\bar\alpha)^{\mathsf T}$ で、$(\beta,\bar\alpha)$ は $S^3$ 全体を動く。

(2) $e^{i\beta/2}UH_0=H_0$ は $U e_2=e^{-i\beta/2}e_2$ を意味する。$U$ はユニタリなので $e_1$ も固有ベクトル、$\det U=1$ より $U=\operatorname{diag}(e^{i\beta/2},e^{-i\beta/2})$。したがって固定部分群は
\[
\Bigl\{\bigl(\operatorname{diag}(e^{i\beta/2},e^{-i\beta/2}),\,e^{i\beta}\bigr)\Bigr\}=\{\exp(i\beta(T_3+Y))\}
\]
である。実際 $T_3+Y$ は $H$ の上成分に $\tfrac12+\tfrac12=1$、下成分に $-\tfrac12+\tfrac12=0$ を与えるので、中性成分 $h_2$ を保つ。これが電磁気の $U(1)_{\mathrm{em}}$、$Q=T_3+Y$ である。軌道定理より $\mathcal M\cong(SU(2)\times U(1)_Y)/U(1)_{\mathrm{em}}\cong S^3$。

(3) $(U,e^{i\beta})=(-\mathbf 1,e^{2\pi i})$、すなわち $\beta=2\pi$ で $U=-\mathbf 1$ の元は $H\mapsto e^{i\pi}(-H)=H$ と自明に作用する。（$U(1)_Y$ の元 $e^{i\beta}$ を $\beta\in[0,4\pi)$ でパラメータ化し、$H$ には $e^{i\beta/2}$ で作用すると考えると、核は $\Z_2$。）実際に $H$ に作用する変換の全体は $\{e^{i\beta/2}U\}=U(2)$ であり、$U(2)$ は $\C^2$ の単位球面 $S^3$ に推移的、$e_2$ の固定部分群は $\{\operatorname{diag}(e^{i\gamma},1)\}\cong U(1)$ で $U(2)/U(1)\cong S^3$。

$\mathcal M=S^3$ は単連結で $\pi_2(S^3)=0$ なので、電弱理論には位相的に安定なひもやモノポールはない（第3章）。
\end{solution}

\begin{problem}[難易度 B]\label{ch01-prob-adj}
ジョージ・グラショウ模型では $SU(2)$ の随伴表現のヒッグス場 $\phi=\phi^a\tfrac{\sigma^a}{2}$（$\phi^a\in\R$）が $\phi\mapsto U\phi U^\dagger$ と変換し、$V=\lambda(\phi^a\phi^a-v^2)^2$ である。\\
(1) $U\phi U^\dagger=R(U)^a{}_b\phi^b\tfrac{\sigma^a}{2}$ で $R(U)\in SO(3)$ となることを用い、真空多様体が $S^2$ であることを示せ。\\
(2) $\phi_0=v\tfrac{\sigma^3}{2}$ の固定部分群を求め、$\mathcal M\cong SU(2)/U(1)$ を示せ。破れずに残るゲージ場（光子）はどの生成子に対応するか。
\end{problem}

\begin{solution}
(1) $\phi^a\phi^a=2\tr\phi^2$ は共役で不変、$U\phi U^\dagger$ もトレースゼロのエルミート行列なので $\phi'^a\sigma^a/2$ と書け、$\phi^a\mapsto\phi'^a$ はノルムを保つ線形写像 $R(U)$ である（$SU(2)$ が連結なので $\det R=1$）。最小点は $\phi^a\phi^a=v^2$、すなわち $\R^3$ の半径 $v$ の球面 $S^2$。$SU(2)\to SO(3)$ は全射（第2章）で $SO(3)$ は $S^2$ に推移的なので、真空多様体全体が1つの軌道である。

(2) $U\sigma^3U^\dagger=\sigma^3$ は $U$ が $\sigma^3$ と可換であることを意味し、$U$ は対角：$U=\operatorname{diag}(e^{i\alpha},e^{-i\alpha})=\exp(i\alpha\sigma^3)$。固定部分群は $U(1)=\{\exp(i\alpha\sigma^3)\}$ であり、$\mathcal M\cong SU(2)/U(1)\cong S^2$。残る生成子は $T^3=\sigma^3/2$（$\phi_0$ の方向）で、これに対応するゲージ場が質量ゼロの光子、$T^1,T^2$ に対応する $W^\pm$ は質量 $\propto v$ を得る。$\pi_2(S^2)=\Z$ がトフーフト・ポリャコフ・モノポールの存在を保証する（第3章）。
\end{solution}

\begin{problem}[難易度 C]\label{ch01-prob-cp2}
$SU(3)$ の随伴表現のヒッグス場 $\phi$（$3\times3$ のトレースゼロ・エルミート行列、$\phi\mapsto U\phi U^\dagger$）の真空が $\phi_0=v\operatorname{diag}(1,1,-2)$ の軌道であるとする。\\
(1) 固定部分群が $S(U(2)\times U(1))=\{\operatorname{diag}(A,e^{i\gamma}):A\in U(2),\ \det A\,e^{i\gamma}=1\}$ であり、これが $U(2)$ と同型であることを示せ。\\
(2) 軌道上の行列 $\phi$ に、固有値 $-2v$ の固有空間（$\C^3$ の複素直線）を対応させることで、真空多様体が $\CP^2$ と同相であることを示せ。\\
(3) 真空多様体の実次元と、破れた生成子の数（＝重くなるゲージ場の数）が一致することを確かめよ。
\end{problem}

\begin{solution}
(1) $U\phi_0U^\dagger=\phi_0$ は $U$ が $\operatorname{diag}(1,1,-2)$ と可換であることと同値で、$U$ は固有空間 $\C^2\oplus\C$ を保つブロック対角行列 $\operatorname{diag}(A,e^{i\gamma})$、$A\in U(2)$。$\det U=1$ より $e^{i\gamma}=(\det A)^{-1}$。写像 $A\mapsto\operatorname{diag}(A,(\det A)^{-1})$ は $U(2)$ から固定部分群への群同型である。

(2) 軌道上の $\phi=U\phi_0U^\dagger$ は固有値 $v$（重複度2）と $-2v$（重複度1）をもつ。$-2v$ の固有空間を $L_\phi=U\C e_3$ とすると、$P_L$ をその上への直交射影として
\[
\phi=v(\mathbf 1-P_L)-2vP_L=v(\mathbf 1-3P_L)
\]
と書ける。よって $\phi\mapsto L_\phi$ は軌道から $\CP^2$（$\C^3$ の直線全体）への写像で、$L\mapsto v(\mathbf 1-3P_L)$ が逆を与える。$SU(3)$ は直線全体に推移的に作用する（直線の単位ベクトルを最後の列とする $SU(3)$ の元がとれる）ので全射。$P_L=uu^\dagger$（$u$ は単位ベクトル）は $u$ の位相によらず連続に依存するので、両方向とも連続。よって $\mathcal M\cong\CP^2\cong SU(3)/S(U(2)\times U(1))$。

(3) $\dim_\R\CP^2=4$。一方 $\dim SU(3)-\dim U(2)=8-4=4$ で、破れた生成子は4個（$\operatorname{diag}(1,1,-2)$ と可換でない、$(1,3),(2,3)$ ブロックの非対角生成子 $T^4,\dots,T^7$）。ヒッグス機構でこれらに対応する4つのゲージ場が質量を得る。一般に $\dim G/H=\dim G-\dim H$ である。
\end{solution}

\begin{problem}[難易度 B]\label{ch01-prob-chiral}
$N_f=N$ フレーバーの QCD では、クォーク凝縮 $\langle\bar q_{Ri}q_{Lj}\rangle\propto\Sigma_{ji}$ が $SU(N)_L\times SU(N)_R$ のもとで $\Sigma\mapsto L\Sigma R^\dagger$ と変換し、真空は $\Sigma=\mathbf 1$ にとれる。\\
(1) $\Sigma=\mathbf 1$ の軌道が $SU(N)$ 全体であり、固定部分群が対角部分群 $SU(N)_V=\{(V,V)\}$ であることを示せ。\\
(2) 写像 $(L,R)\mapsto LR^\dagger$ が同相 $(SU(N)_L\times SU(N)_R)/SU(N)_V\cong SU(N)$ を誘導することを示せ。\\
(3) $N=2$ のときカイラル・ラグランジアンの標的空間は何と同相か。南部・ゴールドストーン粒子（パイオン）の数を答えよ。
\end{problem}

\begin{solution}
(1) 軌道は $\{L R^\dagger\}$ であり、$R=\mathbf 1$ ととれば任意の $L\in SU(N)$ が現れるので $SU(N)$ 全体。$L R^\dagger=\mathbf 1\iff L=R$ なので固定部分群は $\{(V,V)\}$。

(2) 軌道・固定部分群定理そのものである：$(L,R)$ と $(L',R')$ が同じ剰余類（$(L',R')=(LV,RV)$）なら $L'R'^\dagger=LVV^\dagger R^\dagger=LR^\dagger$。逆に $LR^\dagger=L'R'^\dagger$ なら $V=L^{-1}L'=R^{-1}R'$。コンパクト→ハウスドルフなので同相。

(3) $N=2$ では標的は $SU(2)\cong S^3$。パイオン場は $\Sigma=\exp(2i\pi^aT^a/f_\pi)$ と書け、$\dim SU(2)=3$ 個の $\pi^\pm,\pi^0$ が現れる。写像 $\R^3\cup\{\infty\}\cong S^3\to S^3$ の巻き付き数がスカーミオン（バリオン）数である（第3章）。
\end{solution}

\begin{problem}[難易度 B]\label{ch01-prob-hopf}
$z\in S^3\subset\C^2$ に対し $n^a(z)=z^\dagger\sigma^az$ とおく。\\
(1) $\sum_a\sigma^a_{ij}\sigma^a_{kl}=2\delta_{il}\delta_{jk}-\delta_{ij}\delta_{kl}$ を用いて $\lvert\bm n\rvert=1$ を示せ。したがってホップ写像 $h\colon S^3\to S^2$, $z\mapsto\bm n$ が定まる。\\
(2) $zz^\dagger=\tfrac12(\mathbf 1+\bm n\cdot\bm\sigma)$ を示し、$h(z)=h(w)\iff w=e^{i\alpha}z$ を導け。$h$ の各点の逆像は円周である。\\
(3) $\CP^1$（$\C^2$ の複素直線全体）が $S^2$ と同相であることを示せ。量子力学では何に対応するか。
\end{problem}

\begin{solution}
(1) $\sum_a(z^\dagger\sigma^az)^2=\sum_a\bar z_iz_j\bar z_kz_l\,\sigma^a_{ij}\sigma^a_{kl}=2(\bar z_iz_i)(\bar z_kz_k)-(\bar z_iz_i)(\bar z_kz_k)=(z^\dagger z)^2=1$。ここで第1項は $\delta_{il}\delta_{jk}$ により $\bar z_iz_j\bar z_jz_i=\lvert z\rvert^4$、第2項は $\delta_{ij}\delta_{kl}$ により $\lvert z\rvert^4$ である。

(2) $2\times2$ エルミート行列は $\tfrac12(c_0\mathbf 1+\bm c\cdot\bm\sigma)$ と書け、$c_0=\tr$, $c^a=\tr(\,\cdot\,\sigma^a)$ である。$P=zz^\dagger$ について $\tr P=z^\dagger z=1$、$\tr(P\sigma^a)=z^\dagger\sigma^az=n^a$ なので $P=\tfrac12(\mathbf 1+\bm n\cdot\bm\sigma)$。したがって $h(z)=h(w)$ なら $zz^\dagger=ww^\dagger$。両辺を $z$ に作用させると $z=w(w^\dagger z)$ で、$\lvert w^\dagger z\rvert=1$（ノルムを比較）なので $z=e^{i\alpha}w$。逆は明らか。よって $h^{-1}(\bm n)=\{e^{i\alpha}z\}\cong S^1$。

(3) $\CP^1=S^3/(z\sim e^{i\alpha}z)$ であり、(2) より $h$ は連続な全単射 $\bar h\colon\CP^1\to S^2$ を誘導する（全射性：任意の $\bm n\in S^2$ に対し射影 $\tfrac12(\mathbf 1+\bm n\cdot\bm\sigma)$ は階数1で、その像の単位ベクトル $z$ が $h(z)=\bm n$ をみたす）。$\CP^1$ はコンパクト、$S^2$ はハウスドルフなので同相。量子力学では、スピン $1/2$ の状態（全体位相を除いた射線）の空間 $\CP^1$ が\term{ブロッホ球}{Bloch sphere} $S^2$ であることに対応し、$\bm n=\langle\bm\sigma\rangle$ がスピンの向きである。このブロッホ球上のベリー位相が単極子の場になること（第8章）は、ホップ写像のねじれ（第3, 7章）の物理的な現れである。
\end{solution}

\subsection{図形とホモトピー型}

\begin{problem}[難易度 C]\label{ch01-prob-torus}
トーラス $T^2$ を正方形 $I^2$ の対辺を同一視した空間とし、$p$ を正方形の中心 $c=(\tfrac12,\tfrac12)$ の像とする。$T^2\setminus\{p\}$ が $S^1\vee S^1$（2つの円周を1点で貼り合わせた8の字）に強変位レトラクトすることを示せ。（物理的には、トーラス上に1つだけ点状の渦や電荷を置いたときの「その点を除いた空間」の形である。）
\end{problem}

\begin{solution}
\begin{center}
\begin{tikzpicture}[scale=2.2]
  \fill[blue!6] (0,0) rectangle (1,1);
  \draw[thick,red!70!black,mid arrow] (0,0) -- (1,0);
  \draw[thick,red!70!black,mid arrow] (0,1) -- (1,1);
  \draw[thick,blue!70!black,mid arrows] (0,0) -- (0,1);
  \draw[thick,blue!70!black,mid arrows] (1,0) -- (1,1);
  \foreach \a in {0,30,...,330}{
    \draw[-{Stealth[length=1.6mm]},gray] ($(0.5,0.5)+(\a:0.08)$) -- ($(0.5,0.5)+(\a:0.3)$);
  }
  \node[hvtx] at (0.5,0.5) {};
  \node[below] at (0.5,0) {$a$};
  \node[left] at (0,0.5) {$b$};
\end{tikzpicture}
\qquad
\begin{tikzpicture}[scale=1.2]
  \draw[thick,red!70!black,mid arrow] (0,0) .. controls (-1.6,1.2) and (-1.6,-1.2) .. (0,0);
  \draw[thick,blue!70!black,mid arrows] (0,0) .. controls (1.6,1.2) and (1.6,-1.2) .. (0,0);
  \node[vtx] at (0,0) {};
  \node at (-1.4,0) {$a$};
  \node at (1.4,0) {$b$};
  \node at (0,-1.1) {$q(\partial I^2)\cong S^1\vee S^1$};
\end{tikzpicture}
\zucap{中心を除いた正方形を放射状に境界へ押し出す（左）。境界の4辺は $a,b$ の2本に、4頂点は1点に同一視され、8の字になる（右）。}
\end{center}
$q\colon I^2\to T^2$ を商写像、$Q=I^2\setminus\{c\}$、$\lVert z\rVert_\infty=\max(\lvert z_1\rvert,\lvert z_2\rvert)$ とする。$I^2=\{\lVert z-c\rVert_\infty\le\tfrac12\}$ で、境界は $\lVert z-c\rVert_\infty=\tfrac12$ の部分である。
\[
\tilde H(z,t)=c+\Bigl((1-t)+\frac{t}{2\lVert z-c\rVert_\infty}\Bigr)(z-c)
\]
とおくと、係数は正なので $\tilde H(z,t)\neq c$、$\lVert\tilde H(z,t)-c\rVert_\infty=(1-t)\lVert z-c\rVert_\infty+\tfrac t2\le\tfrac12$ で $\tilde H(z,t)\in Q$。$\tilde H(z,0)=z$、$\tilde H(z,1)\in\partial I^2$、境界上の点は動かない。

同一視されるのは境界の点だけで、それらは動かないので、$q\circ\tilde H$ は商写像 $q\times\id_I$（$I$ は局所コンパクト）のファイバー上で一定であり、連続写像 $H\colon(T^2\setminus\{p\})\times I\to T^2\setminus\{p\}$ を誘導する。よって $A=q(\partial I^2)$ は強変位レトラクトである。

$A$ は $q([0,1]\times\{0\})$（両端が同一視された線分 $\cong S^1$、図の $a$）と $q(\{0\}\times[0,1])$（$b$）の和で、両者は1点 $q(0,0)$ のみを共有する。よって $A\cong S^1\vee S^1$。
\end{solution}

\begin{problem}[難易度 B]\label{ch01-prob-letters}
下図のように、アルファベットの大文字を（黒丸を頂点とする）有限グラフとみなす。
\begin{center}
\begin{tikzpicture}[scale=0.62,line width=1.1pt]
  \begin{scope}[xshift=0cm] % A
    \draw (0,0) -- (0.2,0.64) -- (0.5,1.6) -- (0.8,0.64) -- (1,0) (0.2,0.64) -- (0.8,0.64);
    \foreach \p in {(0,0),(0.2,0.64),(0.5,1.6),(0.8,0.64),(1,0)} \node[vtx] at \p {};
  \end{scope}
  \begin{scope}[xshift=1.8cm] % B
    \draw (0,0) -- (0,1.6);
    \draw (0,1.6) .. controls (0.9,1.6) and (0.9,0.8) .. (0,0.8);
    \draw (0,0.8) .. controls (1.05,0.8) and (1.05,0) .. (0,0);
    \foreach \p in {(0,0),(0,0.8),(0,1.6)} \node[vtx] at \p {};
  \end{scope}
  \begin{scope}[xshift=3.4cm] % C
    \draw (1.03,1.33) arc (45:315:0.75);
    \foreach \p in {(1.03,1.33),(1.03,0.27)} \node[vtx] at \p {};
  \end{scope}
  \begin{scope}[xshift=5.2cm] % D
    \draw (0,0) -- (0,1.6) .. controls (1.2,1.6) and (1.2,0) .. (0,0);
    \foreach \p in {(0,0),(0,1.6)} \node[vtx] at \p {};
  \end{scope}
  \begin{scope}[xshift=6.9cm] % E
    \draw (0.9,1.6) -- (0,1.6) -- (0,0) -- (0.9,0) (0,0.8) -- (0.7,0.8);
    \foreach \p in {(0,0),(0,0.8),(0,1.6),(0.9,1.6),(0.7,0.8),(0.9,0)} \node[vtx] at \p {};
  \end{scope}
  \begin{scope}[xshift=8.5cm] % O
    \draw (0.5,0.8) ellipse (0.5 and 0.8);
    \foreach \p in {(0.5,0),(0.5,1.6)} \node[vtx] at \p {};
  \end{scope}
  \begin{scope}[xshift=10.1cm] % P
    \draw (0,0) -- (0,1.6) .. controls (0.9,1.6) and (0.9,0.8) .. (0,0.8);
    \foreach \p in {(0,0),(0,0.8),(0,1.6)} \node[vtx] at \p {};
  \end{scope}
  \begin{scope}[xshift=11.6cm] % Q
    \draw (0.5,0.8) ellipse (0.5 and 0.8);
    \draw (0.65,0.45) -- (1.06,0.02);
    \foreach \p in {(0.5,1.6),(0.854,0.234),(0.65,0.45),(1.06,0.02)} \node[vtx] at \p {};
  \end{scope}
  \begin{scope}[xshift=13.3cm] % R
    \draw (0,0) -- (0,1.6) .. controls (0.9,1.6) and (0.9,0.8) .. (0,0.8) -- (0.9,0);
    \foreach \p in {(0,0),(0,0.8),(0,1.6),(0.9,0)} \node[vtx] at \p {};
  \end{scope}
  \begin{scope}[xshift=14.9cm] % X
    \draw (0,0) -- (1,1.6) (0,1.6) -- (1,0);
    \foreach \p in {(0,0),(1,1.6),(0,1.6),(1,0),(0.5,0.8)} \node[vtx] at \p {};
  \end{scope}
\end{tikzpicture}
\zucap{A, B, C, D, E, O, P, Q, R, X をグラフとみなしたもの（Q の尾は円周と1点で交わる）。}
\end{center}
(1) 頂点数 $V$、辺数 $E$ の連結な有限グラフは $E-V+1$ 個の円周の楔和とホモトピー同値であることを示せ。ただし「CW対 $(X,A)$ で $A$ が可縮なら $X\to X/A$ はホモトピー同値」および「有限グラフの木は可縮」は用いてよい。\\
(2) 10文字をホモトピー型で分類せよ。（$E-V+1$ はファインマン図のループ数と同じ量である。第4章。）
\end{problem}

\begin{solution}
(1) $G$ の極大な木 $T$（閉路をもたない連結部分グラフで包含について極大）をとる。$T$ はすべての頂点を含む：含まない頂点があれば、$G$ の連結性から $T$ の頂点と $T$ 外の頂点を結ぶ辺 $e$ があり、$T\cup e$ はより大きな木になって矛盾。$V$ 頂点の木の辺数は $V-1$（次数1の頂点と辺を1つずつ除く帰納法）。$T$ は可縮な部分複体なので $G\simeq G/T$ で、$G/T$ は1頂点と、$T$ に含まれない $E-(V-1)$ 本の辺がループになったグラフ、すなわち $\bigvee^{E-V+1}S^1$。

(2) 各文字の $(V,E,E-V+1)$ は
\begin{center}
\begin{tabular}{c|cccccccccc}
文字 & A & B & C & D & E & O & P & Q & R & X\\ \hline
$V$ & 5 & 3 & 2 & 2 & 6 & 2 & 3 & 4 & 4 & 5\\
$E$ & 5 & 4 & 1 & 2 & 5 & 2 & 3 & 4 & 4 & 4\\
$E-V+1$ & 1 & 2 & 0 & 1 & 0 & 1 & 1 & 1 & 1 & 0
\end{tabular}
\end{center}
（Q は円周が2頂点で2本の弧に分かれ、尾が交点で2本に分かれる。）よって
可縮：C, E, X。$S^1$ とホモトピー同値：A, D, O, P, Q, R。$S^1\vee S^1$ とホモトピー同値：B。
3つの類が互いに異なることは、基本群（$1$, $\Z$, 自由群 $F_2$；第2章）か、$H_1$ の階数 $0,1,2$（第4章）で区別される。
\end{solution}

\begin{problem}[難易度 D]\label{ch01-prob-vortex}
2次元空間の複素スカラー場 $\phi$ が $V=\lambda(\lvert\phi\rvert^2-v^2)^2$ をもつとする。\\
(1) \textbf{大域的 $U(1)$ の場合}：静的エネルギー $E=\int d^2x\,(\lvert\nabla\phi\rvert^2+V)$ の配位で、$r\ge r_0$ で $\phi=v\,e^{in\theta}$（$n\in\Z$）となるものを考える。$n\neq0$ ならば半径 $R$ の円板内のエネルギーが $R\to\infty$ で $2\pi n^2v^2\log(R/r_0)$ のように発散することを示せ。\\
(2) \textbf{ゲージ化した場合}：$U(1)$ ゲージ場 $A$ を導入し、電荷 $q$ の場の共変微分 $D\phi=(d-iqA)\phi$ を用いて $E=\int d^2x\,(\tfrac12F_{12}^2+\lvert D\phi\rvert^2+V)$ とする。$r\ge r_0$ で $\phi=v\,e^{in\theta}$、$A=\tfrac nq\,d\theta$ とすれば、外側の領域からのエネルギーの寄与がゼロになることを示せ。このとき渦を貫く磁束 $\int F$ を求めよ。\\
(3) 連続な時間発展 $\phi(\bm x,t)$ で、各時刻で遠方の円周 $\lvert\bm x\rvert=R$ 上 $\phi\neq0$ であるとき、$\theta\mapsto\phi(R,\theta,t)/\lvert\phi\rvert$ のホモトピー類は $t$ によらないことを説明せよ。\\
(4) 遠方の円周上の写像 $\phi/\lvert\phi\rvert\colon S^1\to S^1$ が定値写像とホモトピックでない（巻き付き数 $n\neq0$）とき、$\phi$ は円板の内部で必ず零点をもつことを示せ。ただし「円板 $D^2$ 全体で定義された連続写像 $D^2\to S^1$ の境界への制限は定値写像とホモトピック」であることを証明の中で示すこと。
\end{problem}

\begin{solution}
(1) 極座標で $\lvert\nabla\phi\rvert^2=\lvert\partial_r\phi\rvert^2+r^{-2}\lvert\partial_\theta\phi\rvert^2$。$r\ge r_0$ では $\partial_r\phi=0$、$\lvert\partial_\theta\phi\rvert=n v$、$V=0$ なので
\[
E(R)\ge\int_{r_0}^R\!\!\int_0^{2\pi}\frac{n^2v^2}{r^2}\,r\,d\theta\,dr=2\pi n^2v^2\log\frac{R}{r_0}\xrightarrow{R\to\infty}\infty .
\]
大域的な渦（例えば超流動の渦）のエネルギーは系の大きさの対数で発散する（デリックの定理と整合）。

(2) $r\ge r_0$ で $d\phi=in\phi\,d\theta$、$qA=n\,d\theta$ なので $D\phi=i(n-n)\phi\,d\theta=0$。また $A=\tfrac nq d\theta$ は $r>0$ で $F=dA=\tfrac nq\,d(d\theta)=0$、$V=0$。よって外側のエネルギー密度はゼロで、エネルギーは芯の付近だけから来て有限になりうる（アブリコソフ・ニールセン・オレセン渦）。
磁束は、内部で $A$ が滑らかに続くとしてストークスの定理より
\[
\int_{D_R}F=\oint_{\lvert x\rvert=R}A=\frac nq\oint d\theta=\frac{2\pi n}{q}.
\]
単位電荷 $q=1$ では $2\pi n$ で、磁束が量子化される（第5章のディラック量子化と同じ構造）。

(3) $\Phi(\theta,t)=\phi(R,\theta,t)/\lvert\phi(R,\theta,t)\rvert$ は $S^1\times[t_0,t_1]$ 上連続（分母は零にならない）なので、$t_0$ での写像と $t_1$ での写像の間のホモトピーそのものである。よってホモトピー類（第2章で整数の巻き付き数と同一視される）は保存される。

(4) まず主張：$g\colon D^2\to S^1$ が連続なら、$H(\theta,t)=g\bigl((1-t)e^{i\theta}\bigr)$ は $S^1\times I\to S^1$ の連続写像で、$H(\cdot,0)=g|_{\partial D^2}$、$H(\cdot,1)=g(0)$（定値）。よって境界への制限は定値写像とホモトピック。
本題：$\phi$ が円板 $\lvert\bm x\rvert\le R$ 上で零点をもたなければ、$g=\phi/\lvert\phi\rvert$ は円板全体で連続な $S^1$ 値写像であり、上の主張より境界の写像は定値写像とホモトピック、すなわち $n=0$。対偶により、$n\neq0$ なら円板内に零点がある。これが渦の「芯」であり、そこでは $V(\phi)=\lambda v^4>0$ のエネルギーが必ず集中する。
\end{solution}

%================================================
\chapter{基本群・被覆空間と量子力学}\label{ch:02}
% 20問
%================================================

\section{物理からの動機}

\subsection*{アハラノフ・ボーム効果}
無限に長いソレノイドの外側では磁場 $\bm B=0$ であるが、ベクトル・ポテンシャルは $0$ にできない。ソレノイドを貫く磁束を $\Phi$ とすると、外側で $A=\dfrac{\Phi}{2\pi}\,d\theta$ ととれ、$F=dA=0$ でありながら $\oint A=\Phi$ である。電荷 $q$ の粒子の経路積分で、経路 $\gamma$ には位相 $\exp\bigl(iq\int_\gamma A\bigr)$ がかかる。ソレノイドの上を通る経路 $\gamma_1$ と下を通る経路 $\gamma_2$ の位相差は
\[
\exp\Bigl(iq\int_{\gamma_1}A-iq\int_{\gamma_2}A\Bigr)=\exp\Bigl(iq\oint_{\gamma_1\cdot\bar\gamma_2}A\Bigr)=e^{iq\Phi}
\]
であり、スクリーン上の干渉縞がずれる（\term{アハラノフ・ボーム効果}{Aharonov--Bohm effect}、AB効果）。$F=0$ の領域では、$\int_\gamma A$ は端点を止めた連続変形で変わらない（ストークスの定理）。したがって位相は経路そのものではなく、経路がソレノイドのまわりを\textbf{何回まわるか}だけで決まる。
\begin{center}
\begin{tikzpicture}[scale=1.1]
  \fill[gray!30] (0,0) circle (0.45);
  \draw[thick] (0,0) circle (0.45);
  \node at (0,0) {\footnotesize $\Phi$};
  \node[vtx] at (-3.5,0) {}; \node[left] at (-3.5,0) {源};
  \draw[very thick] (3.6,-1.6) -- (3.6,1.6); \node[right] at (3.6,1.3) {スクリーン};
  \node[vtx] at (3.6,0.2) {}; \node[right] at (3.6,0.2) {\footnotesize $x$};
  \draw[thick,red!70!black,mid arrow] (-3.5,0) .. controls (-1.8,1.6) and (1.6,1.6) .. (3.6,0.2);
  \draw[thick,blue!70!black,mid arrow] (-3.5,0) .. controls (-1.8,-1.6) and (1.6,-1.4) .. (3.6,0.2);
  \node[red!70!black] at (0,1.45) {$\gamma_1$};
  \node[blue!70!black] at (0,-1.45) {$\gamma_2$};
  \node at (0,-0.75) {\footnotesize $\bm B=0$ の領域で $A=\frac{\Phi}{2\pi}d\theta$};
\end{tikzpicture}
\zucap{AB効果の配置。灰色のソレノイドの中にだけ磁場がある。粒子が通れる空間は $\R^2\setminus D^2\simeq S^1$ で、2つの経路の差 $\gamma_1\cdot\bar\gamma_2$ はソレノイドを1周するループである。}
\end{center}
粒子が動ける空間 $X=\R^2\setminus(\text{ソレノイド})$ の中のループを、連続変形で移り合うものを同一視して分類したものが\term{基本群}{fundamental group} $\pi_1(X)$ である。いまの場合 $\pi_1(X)\cong\Z$ で、整数は巻き付き数である。AB位相は準同型 $\pi_1(X)\to U(1)$、$n\mapsto e^{inq\Phi}$ にほかならない。

\subsection*{本章の主題}
\begin{itemize}
\item \textbf{渦とひも}：3次元空間のひも状の位相欠陥（宇宙ひも、超伝導の磁束管、超流動の渦）は、ひもを囲む円周から真空多様体 $\mathcal M$ への写像で特徴づけられ、$\pi_1(\mathcal M)$ で分類される。$U(1)\to1$ の破れでは $\pi_1(S^1)=\Z$ で、渦は整数の巻き付き数をもつ。$\pi_1$ が $\Z_2$ や非可換群になる例もある。
\item \textbf{スピン}：$SO(3)$ の中で $2\pi$ 回転のループは1点に縮められず、$4\pi$ 回転のループは縮められる。これは $\pi_1(SO(3))=\Z_2$ であり、$SU(2)\to SO(3)$ が\term{普遍被覆}{universal cover}であることの言いかえである。スピン $1/2$ の状態が $2\pi$ 回転で $-1$ 倍になることの位相幾何的な起源である。
\item \textbf{量子統計}：同種粒子の配位空間の基本群は、3次元では $\Z_2$ 程度の情報（置換群 $S_n$）しかもたないが、2次元では組みひも群 $B_n$ になる。前者からボソン・フェルミオン、後者からエニオンが現れる。
\item \textbf{多重連結空間上の経路積分}：配位空間 $Q$ が単連結でないとき、経路積分は $\pi_1(Q)$ の元ごとのセクターに分かれ、
\[
K=\sum_{\alpha\in\pi_1(Q)}\chi(\alpha)\,K_\alpha,\qquad \chi\colon\pi_1(Q)\to U(1)\ \text{は準同型}
\]
の形に書ける（シュルマン、レイドロー・ドウィット）。$\chi$ の選び方が AB 位相、統計の位相、$\theta$ 角などの「量子化の曖昧さ」である。
\end{itemize}
これらすべてを支える数学的道具が\term{被覆空間}{covering space}であり、多価関数（位相 $\arg\phi$、スピノル、エニオンの波動関数）を一価にする「上の空間」として現れる。

\section{基本事項のまとめ}

\subsection{基本群}

\begin{mydef}{道・ループ・基本群}{ch02-pi1}
$X$ の\term{道}{path}とは連続写像 $f\colon I\to X$ のことで、$f(0)=f(1)=x_0$ のとき $x_0$ を基点とする\term{ループ}{loop}という。$f(1)=g(0)$ のとき道の積を
\[
(f\cdot g)(s)=\begin{cases}f(2s)&(0\le s\le\tfrac12)\\ g(2s-1)&(\tfrac12\le s\le1)\end{cases}
\]
で定め（$f$ を先にたどる）、逆の道を $\bar f(s)=f(1-s)$ とする。2つの道が\textbf{端点を止めてホモトピック}（$f\simeq g$ rel $\{0,1\}$）であるとは、$H\colon I\times I\to X$ で $H(s,0)=f(s)$、$H(s,1)=g(s)$、$H(0,t)$、$H(1,t)$ が $t$ によらないものがあることをいう。$x_0$ を基点とするループのこの同値類の集合
\[
\pi_1(X,x_0)
\]
は積 $[f][g]=[f\cdot g]$ で群になる（単位元は定値ループ、逆元は $[\bar f]$）。$\pi_1(X,x_0)=1$ かつ $X$ が弧状連結のとき、$X$ は\term{単連結}{simply connected}であるという。
\end{mydef}

\begin{myprop}{基点の変更}{ch02-basept}
$h$ を $x_0$ から $x_1$ への道とすると、$\beta_h[\gamma]=[h\cdot\gamma\cdot\bar h]$ は群同型 $\pi_1(X,x_1)\to\pi_1(X,x_0)$ である。別の道 $h'$ に対し $\beta_{h'}=c_{g}\circ\beta_h$（$g=[h'\cdot\bar h]$ による共役 $c_g(x)=gxg^{-1}$）。したがって弧状連結な $X$ では $\pi_1(X)$ は同型を除いて基点によらず、$\pi_1$ がアーベル群なら同型は標準的である。また基点を止めない\term{自由ループ}{free loop}のホモトピー類は、$\pi_1(X,x_0)$ の\textbf{共役類}と1対1に対応する。
\end{myprop}

\begin{mythm}{誘導準同型とホモトピー不変性}{ch02-functor}
連続写像 $f\colon(X,x_0)\to(Y,y_0)$ は準同型 $f_*\colon\pi_1(X,x_0)\to\pi_1(Y,y_0)$、$f_*[\gamma]=[f\circ\gamma]$ を誘導し、$(g\circ f)_*=g_*\circ f_*$、$\id_*=\id$ をみたす。ホモトピー同値 $X\simeq Y$ は基本群の同型を誘導する。特に可縮な空間は単連結であり、$\pi_1(X\times Y)\cong\pi_1(X)\times\pi_1(Y)$。
\end{mythm}

\begin{mythm}{円周の基本群}{ch02-pi1S1}
$p\colon\R\to S^1$、$p(s)=e^{2\pi is}$ とする。$1\in S^1$ を基点とするループ $\gamma$ の、$0$ を始点とする持ち上げ $\tilde\gamma$（$p\circ\tilde\gamma=\gamma$）の終点を対応させる写像
\[
\deg\colon\pi_1(S^1,1)\longrightarrow\Z,\qquad [\gamma]\longmapsto\tilde\gamma(1)
\]
は群同型である。$\deg[\gamma]$ を $\gamma$ の\term{巻き付き数}{winding number}（写像度）という。
\end{mythm}
\begin{center}
\begin{tikzpicture}[scale=1.0]
  \draw[thick,gray,domain=-0.25:2.25,samples=200,smooth]
    plot ({1.6*cos(360*\x)}, {0.45*sin(360*\x)+1.1*\x+1.2});
  \draw[very thick,red!75!black,domain=0:2,samples=160,smooth]
    plot ({1.6*cos(360*\x)}, {0.45*sin(360*\x)+1.1*\x+1.2});
  \foreach \n in {0,1,2}{
    \node[vtx] at (1.6,{1.1*\n+1.2}) {};
    \node[right] at (1.65,{1.1*\n+1.2}) {\footnotesize $\n$};
  }
  \node[left] at (-1.7,3.0) {$\R$};
  \draw[thick] (0,-1.3) ellipse (1.6 and 0.45);
  \draw[very thick,red!75!black,mid arrow] (1.6,-1.3) arc (0:350:1.6 and 0.45);
  \node[vtx] at (1.6,-1.3) {};
  \node[right] at (1.65,-1.3) {\footnotesize $1=p(n)$};
  \node[left] at (-1.7,-1.3) {$S^1$};
  \draw[-{Stealth[length=2.2mm]}] (2.6,1.6) -- (2.6,-0.9) node[midway,right] {$p$};
\end{tikzpicture}
\zucap{普遍被覆 $p\colon\R\to S^1$。$\R$ を螺旋として描き、真下の円周へ射影する。$S^1$ を2周するループ $\gamma(s)=e^{4\pi is}$ の $0$ からの持ち上げ $\tilde\gamma(s)=2s$（赤）の終点 $2$ が巻き付き数である。}
\end{center}

物理では巻き付き数を積分で書くことが多い。$\gamma$ が $\C\setminus\{0\}$ の滑らかなループなら
\[
\deg\gamma=\frac{1}{2\pi}\oint_\gamma d\theta=\frac{1}{2\pi i}\oint_\gamma\frac{dz}{z},\qquad d\theta=\frac{x\,dy-y\,dx}{x^2+y^2}
\]
である（問題 \ref{ch02-prob-windint}）。$d\theta$ は閉形式だが完全形式ではなく、この「完全でない閉形式」が第5章のド・ラーム・コホモロジーの出発点になる。

\subsection{ファン・カンペンの定理と計算例}

\begin{mythm}{ザイフェルト・ファン・カンペンの定理}{ch02-vk}
$X=U\cup V$（$U,V$ は開集合）で、$U$, $V$, $U\cap V$ が弧状連結、$x_0\in U\cap V$ とする。$\pi_1(U\cap V)\to\pi_1(U)$, $\pi_1(V)$ を包含が誘導する準同型 $i_1,i_2$ とすると
\[
\pi_1(X,x_0)\cong\pi_1(U,x_0)*_{\pi_1(U\cap V,x_0)}\pi_1(V,x_0)
=\bigl(\pi_1(U)*\pi_1(V)\bigr)\big/\bigl\langle\!\bigl\langle i_1(c)\,i_2(c)^{-1}\bigr\rangle\!\bigr\rangle .
\]
ここで $*$ は自由積、$\langle\!\langle\cdots\rangle\!\rangle$ は生成される正規部分群である。特に $U\cap V$ が単連結なら $\pi_1(X)\cong\pi_1(U)*\pi_1(V)$。
\end{mythm}
\begin{note}
「どのループも $U$ 内の部分と $V$ 内の部分の積に分解でき（ルベーグ数の議論）、分解の仕方の曖昧さは $U\cap V$ のループを $U$ 側で見るか $V$ 側で見るかだけである」というのが証明の要点である。
\end{note}

\begin{myprop}{2胞体の接着}{ch02-cell}
弧状連結な $X$ に、接着写像 $\varphi\colon S^1\to X$ で2胞体を貼った空間 $Y=X\cup_\varphi D^2$ では
\[
\pi_1(Y)\cong\pi_1(X)\big/\langle\!\langle[\varphi]\rangle\!\rangle
\]
（$[\varphi]$ は基点へ道でつないだループの類）。$n\ge3$ の胞体を貼っても $\pi_1$ は変わらない。したがって CW 複体の $\pi_1$ は2骨格で決まり、「1胞体が生成元、2胞体が関係式」を与える。
\end{myprop}

\begin{myex}{基本群の表}{ch02-table}
\begin{center}
\renewcommand{\arraystretch}{1.25}
\begin{tabular}{lll}
空間 & $\pi_1$ & 理由\\ \hline
$\R^n$, $D^n$, 可縮な空間 & $1$ & 可縮\\
$S^1\simeq\R^2\setminus\{0\}\simeq\R^3\setminus(\text{直線})$ & $\Z$ & 定理 \ref{thm:ch02-pi1S1}\\
$S^n$（$n\ge2$）, $\R^n\setminus\{0\}$（$n\ge3$） & $1$ & ファン・カンペン\\
$\R^2\setminus\{k\ \text{点}\}\simeq\bigvee^kS^1$ & 自由群 $F_k$ & ファン・カンペン\\
$T^n$ & $\Z^n$ & 積\\
$\RP^n$（$n\ge2$） & $\Z_2$ & 被覆 $S^n\to\RP^n$\\
$\Sigma_g$（種数 $g$ の閉曲面） & $\langle a_i,b_i\mid\prod_{i=1}^g[a_i,b_i]\rangle$ & 2胞体の接着\\
レンズ空間 $L(N,p)=S^3/\Z_N$ & $\Z_N$ & 被覆 $S^3\to L(N,p)$
\end{tabular}
\end{center}
ここで $[a,b]=aba^{-1}b^{-1}$。
\end{myex}

可換化 $\pi_1(X)^{\mathrm{ab}}=\pi_1/[\pi_1,\pi_1]$ は第4章の1次ホモロジー $H_1(X)$ と一致する（フレヴィッチの定理）。巻き付き数のような「足し算できる」量は可換化しか見ないので、非可換な $\pi_1$ の情報（交換子）は検出できない（問題 \ref{ch02-prob-twopts}）。

\subsection{被覆空間}

\begin{mydef}{被覆空間}{ch02-cov}
連続全射 $p\colon\tilde X\to X$ が\term{被覆写像}{covering map}であるとは、各 $x\in X$ が開近傍 $U$ をもち、$p^{-1}(U)$ が互いに交わらない開集合 $\tilde U_\alpha$ の和で、各 $p|_{\tilde U_\alpha}\colon\tilde U_\alpha\to U$ が同相となることをいう（$U$ は\textbf{均等に覆われる}）。$p^{-1}(x)$ を\term{ファイバー}{fiber}、その濃度（$X$ が連結なら一定）を\textbf{葉数}という。同相 $\tau\colon\tilde X\to\tilde X$ で $p\circ\tau=p$ をみたすものを\term{被覆変換}{deck transformation}といい、その群を $\mathrm{Deck}(\tilde X/X)$ と書く。
\end{mydef}

\begin{mythm}{持ち上げ}{ch02-lift}
$p\colon\tilde X\to X$ を被覆写像、$p(\tilde x_0)=x_0$ とする。
\begin{enumerate}
\item （道の持ち上げ）$x_0$ を始点とする道 $\gamma$ に対し、$\tilde x_0$ を始点とする持ち上げ $\tilde\gamma$（$p\circ\tilde\gamma=\gamma$）がただ1つ存在する。
\item （ホモトピーの持ち上げ）$H\colon I\times I\to X$ と $H(0,0)$ の上の点を与えると、持ち上げ $\tilde H$ がただ1つ存在する。特に端点を止めてホモトピックな道の持ち上げは同じ終点をもつ。
\item $p_*\colon\pi_1(\tilde X,\tilde x_0)\to\pi_1(X,x_0)$ は単射で、$[\gamma]\in p_*\pi_1(\tilde X,\tilde x_0)$ と「$\tilde\gamma$ が閉じる」ことは同値。葉数は指数 $[\pi_1(X):p_*\pi_1(\tilde X)]$ に等しい（$\tilde X$ 連結のとき）。
\item （持ち上げの判定法）$Y$ が弧状連結かつ局所弧状連結のとき、$f\colon(Y,y_0)\to(X,x_0)$ が持ち上げ $\tilde f\colon(Y,y_0)\to(\tilde X,\tilde x_0)$ をもつための必要十分条件は $f_*\pi_1(Y,y_0)\subset p_*\pi_1(\tilde X,\tilde x_0)$。持ち上げは1点での値で一意に決まる。
\end{enumerate}
\end{mythm}

\begin{mythm}{被覆空間のガロア対応}{ch02-galois}
$X$ が弧状連結・局所弧状連結・半局所単連結（多様体や CW 複体はすべてそう）とする。
\begin{enumerate}
\item 単連結な被覆 $\tilde X\to X$（\term{普遍被覆}{universal cover}）が存在し、$\mathrm{Deck}(\tilde X/X)\cong\pi_1(X)$。
\item 連結な被覆の同型類 $\longleftrightarrow$ $\pi_1(X,x_0)$ の部分群の共役類（$\tilde X\mapsto p_*\pi_1(\tilde X)$）。
\item $H=p_*\pi_1(\tilde X,\tilde x_0)$ が正規部分群であること（\textbf{正規被覆}）と、被覆変換群がファイバーに推移的に作用することは同値で、このとき $\mathrm{Deck}\cong\pi_1(X)/H$。一般には $\mathrm{Deck}\cong N(H)/H$（$N(H)$ は正規化群）。
\end{enumerate}
\end{mythm}

\begin{mythm}{自由な作用による商}{ch02-quotient}
群 $G$ が空間 $Y$ に\term{被覆空間作用}{covering space action}をする（各 $y$ に近傍 $U$ があって $g\ne e$ なら $gU\cap U=\emptyset$；有限群のハウスドルフ空間への自由な作用、リー群の離散部分群の作用など）とき、$Y\to Y/G$ は被覆写像で、被覆変換群は $G$ である。さらに $Y$ が単連結なら
\[
\pi_1(Y/G)\cong G .
\]
\end{mythm}

\begin{myex}{被覆の例}{ch02-covex}
\begin{itemize}
\item $\R\to S^1$（$\Z$ が平行移動で作用）、$S^1\to S^1$、$z\mapsto z^n$（$n$ 重、被覆変換群 $\Z_n$）。$S^1$ の連結な被覆はこれらで尽きる（$\Z$ の部分群は $n\Z$）。
\item $\R^n\to T^n$（$\Z^n$）、$S^n\to\RP^n$（$\Z_2=\{\pm1\}$）。
\item $SU(2)\cong S^3\to SO(3)\cong\RP^3$（$\{\pm\mathbf 1\}$）、$SU(N)\to SU(N)/\Z_N$。
\item $S^3\to L(N,p)$（$\Z_N$；問題 \ref{ch02-prob-lens}）。
\end{itemize}
\end{myex}

\subsection{リー群と等質空間の基本群}

\begin{myprop}{リー群の基本群}{ch02-liepi1}
\begin{center}
\renewcommand{\arraystretch}{1.25}
\begin{tabular}{llll}
群 $G$ & $\pi_1(G)$ & 普遍被覆 & 備考\\ \hline
$U(1)$ & $\Z$ & $\R$ & \\
$SU(N)$ & $1$ & $SU(N)$ & 第3章（ファイブレーション）\\
$U(N)$ & $\Z$ & $SU(N)\times\R$ & $\det_*$ が同型\\
$SU(N)/\Z_N$ & $\Z_N$ & $SU(N)$ & \\
$SO(2)$ & $\Z$ & $\R$ & \\
$SO(n)$（$n\ge3$） & $\Z_2$ & $\Spin(n)$ & $\Spin(3)=SU(2)$ など\\
$SO^+(3,1)$ & $\Z_2$ & $SL(2,\C)$ & \\
$Sp(N)$ & $1$ & $Sp(N)$ &
\end{tabular}
\end{center}
連結なリー群の $\pi_1$ は常にアーベル群である。
\end{myprop}

\begin{mythm}{等質空間の基本群}{ch02-GH}
$G$ が連結かつ単連結なリー群、$H\subset G$ が閉部分群なら
\[
\pi_1(G/H)\cong\pi_0(H)=H/H_0\qquad(H_0\ \text{は単位元を含む連結成分}).
\]
$H$ が離散なら $G\to G/H$ は $H$ を被覆変換群とする普遍被覆で、定理 \ref{thm:ch02-quotient} から従う（問題 \ref{ch02-prob-strings}）。一般の場合はファイブレーションの完全列から従う（第3, 7章）。$G$ が単連結でないときは、普遍被覆 $\tilde G$ と $H$ の逆像 $\tilde H$ に置き換えて $\pi_1(G/H)\cong\pi_0(\tilde H)$ とする。
\end{mythm}

\begin{physbox}[：ひも状欠陥の分類]
3次元空間で対称性が $G\to H$ と破れ、真空多様体が $\mathcal M=G/H$ のとき、ひも（線欠陥）を1周する小さな円周上で場は $\mathcal M$ に値をとり、ループ $S^1\to\mathcal M$ が定まる。基点を止めないので、ひもの種類は $\pi_1(\mathcal M)$ の\textbf{共役類}で分類される。
\begin{center}
\renewcommand{\arraystretch}{1.2}\small
\begin{tabular}{lll}
破れ & $\pi_1(\mathcal M)$ & 欠陥\\ \hline
$U(1)\to1$（超伝導） & $\Z$ & ANO 渦、磁束量子化\\
$U(1)\to\Z_N$（電荷 $N$ のヒッグス） & $\Z$ & 磁束 $2\pi/N$ を単位とする渦\\
$SO(3)\to1$ & $\Z_2$ & $\Z_2$ ひも（2本で対消滅）\\
$SO(3)\to D_2$（2軸性ネマティック液晶） & $Q_8$（非可換） & 非可換ひも\\
$SU(2)\times U(1)\to U(1)_{\mathrm{em}}$（電弱） & $1$ & 安定なひもなし
\end{tabular}
\end{center}
$\pi_1$ が非可換だと、2本のひもが交差したとき交換子に当たるひもでつながれ、自由にすり抜けられない。
\end{physbox}

\subsection{量子力学と基本群}

\begin{mythm}{多重連結空間上の経路積分（シュルマン、レイドロー・ドウィット）}{ch02-sld}
配位空間 $Q$ の普遍被覆を $\tilde Q$、被覆変換群を $\pi_1(Q)$ とする。$Q$ 上の伝播関数を $\tilde Q$ 上の伝播関数 $\tilde K$ の和で
\[
K(q,t;q',0)=\sum_{\alpha\in\pi_1(Q)}\chi(\alpha)\,\tilde K(\tilde q,t;\alpha\tilde q',0)
\]
と表すとき（$\tilde q,\tilde q'$ は $q,q'$ の上の点）、合成則 $\int K K=K$ と確率の保存から $\chi$ は1次元ユニタリ表現（準同型 $\chi\colon\pi_1(Q)\to U(1)$）でなければならない。異なる $\chi$ は異なる量子化を与える。
\end{mythm}
同じことを、「波動関数を $\tilde Q$ 上の関数で $\psi(\alpha\tilde q)=\chi(\alpha)\psi(\tilde q)$ をみたすものとする」と言っても、あるいは「$Q$ 上に平坦な $U(1)$ 接続を入れ、そのホロノミーを $\chi$ とする」と言ってもよい（問題 \ref{ch02-prob-wilson}；第7, 8章）。

\begin{myex}{同種粒子の配位空間と組みひも群}{ch02-braid}
$d$ 次元空間の $n$ 個の同種粒子の配位空間は
\[
C_n(\R^d)=\bigl\{(x_1,\dots,x_n)\in(\R^d)^n: x_i\neq x_j\ (i\ne j)\bigr\}\big/S_n
\]
である。
\begin{itemize}
\item $d\ge3$：$\pi_1(C_n(\R^d))\cong S_n$（置換群）。$S_n$ の1次元表現は自明表現と符号表現だけで、ボソンとフェルミオンに対応する。
\item $d=2$：$\pi_1(C_n(\R^2))\cong B_n$（\term{組みひも群}{braid group}）。生成元 $\sigma_1,\dots,\sigma_{n-1}$（$i$ 番目と $i+1$ 番目の粒子を反時計回りに半回転させて入れ替える）と関係式
\[
\sigma_i\sigma_{i+1}\sigma_i=\sigma_{i+1}\sigma_i\sigma_{i+1},\qquad \sigma_i\sigma_j=\sigma_j\sigma_i\ (\lvert i-j\rvert\ge2)
\]
で表示される（アルティン）。$\sigma_i^2\ne1$ である点が $S_n$ と異なる。
\end{itemize}
\end{myex}
\begin{center}
\begin{tikzpicture}[scale=0.9,thick]
  \begin{scope}
    \draw (0,0) -- (0,2); \draw (3,0) -- (3,2);
    \draw (1,0) .. controls (1,1) and (2,1) .. (2,2);
    \draw[white,line width=5pt] (2,0) .. controls (2,1) and (1,1) .. (1,2);
    \draw (2,0) .. controls (2,1) and (1,1) .. (1,2);
    \foreach \x in {0,1,2,3}{\node[vtx] at (\x,0) {}; \node[vtx] at (\x,2) {};}
    \node[below] at (0,-0.1) {\footnotesize $1$}; \node[below] at (1,-0.1) {\footnotesize $i$};
    \node[below] at (2,-0.1) {\footnotesize $i{+}1$}; \node[below] at (3,-0.1) {\footnotesize $n$};
    \node at (1.5,-0.8) {$\sigma_i$};
  \end{scope}
  \begin{scope}[xshift=4.5cm]
    \draw (0,0) -- (0,2); \draw (3,0) -- (3,2);
    \draw (2,0) .. controls (2,1) and (1,1) .. (1,2);
    \draw[white,line width=5pt] (1,0) .. controls (1,1) and (2,1) .. (2,2);
    \draw (1,0) .. controls (1,1) and (2,1) .. (2,2);
    \foreach \x in {0,1,2,3}{\node[vtx] at (\x,0) {}; \node[vtx] at (\x,2) {};}
    \node at (1.5,-0.8) {$\sigma_i^{-1}$};
  \end{scope}
  \begin{scope}[xshift=9cm]
    \draw (0,0) .. controls (0,0.5) and (1,0.5) .. (1,1);
    \draw[white,line width=5pt] (1,0) .. controls (1,0.5) and (0,0.5) .. (0,1);
    \draw (1,0) .. controls (1,0.5) and (0,0.5) .. (0,1);
    \draw (0,1) .. controls (0,1.5) and (1,1.5) .. (1,2);
    \draw[white,line width=5pt] (1,1) .. controls (1,1.5) and (0,1.5) .. (0,2);
    \draw (1,1) .. controls (1,1.5) and (0,1.5) .. (0,2);
    \foreach \x in {0,1}{\node[vtx] at (\x,0) {}; \node[vtx] at (\x,2) {};}
    \node at (0.5,-0.8) {$\sigma_1^2$};
  \end{scope}
\end{tikzpicture}
\zucap{組みひも群の生成元（時間は下から上）。$\sigma_i$ と $\sigma_i^{-1}$ は交差の上下が逆。$\sigma_1^2$ は一方の粒子が他方のまわりを1周する運動で、2次元平面内では縮められないが、3次元では第3の方向へ持ち上げてほどける。}
\end{center}

\begin{physbox}[：本章の辞書]
\begin{center}
\small
\begin{tabular}{p{0.5\linewidth}p{0.4\linewidth}}
物理 & 数学\\ \hline
AB 位相 $e^{iq\Phi n}$ & 準同型 $\pi_1(X)\to U(1)$\\
平坦な接続のウィルソン・ループ & ホロノミー表現 $\pi_1\to G$\\
渦の巻き付き数・ひもの種類 & $\pi_1(\mathcal M)$ の（共役類の）元\\
多価な位相 $\arg\phi$、スピノル、エニオン波動関数 & 被覆空間上の一価関数\\
スピン $1/2$ の $2\pi$ 回転で $-1$ & $\pi_1(SO(3))=\Z_2$、$SU(2)\to SO(3)$\\
ボソン／フェルミオン／エニオン & $S_n$／$B_n$ の1次元表現\\
$\theta$ セクター、量子化の曖昧さ & $\mathrm{Hom}(\pi_1(Q),U(1))$
\end{tabular}
\end{center}
\end{physbox}

\section{問題}

\subsection{基本群の計算}

\begin{problem}[難易度 A]\label{ch02-prob-basept}
$X$ を弧状連結な空間、$h$ を $x_0$ から $x_1$ への道とし、$\beta_h\colon\pi_1(X,x_1)\to\pi_1(X,x_0)$、$\beta_h[\gamma]=[h\cdot\gamma\cdot\bar h]$ とする。\\
(1) $\beta_h$ が well-defined な群同型であることを示せ。\\
(2) $x_0$ から $x_1$ への別の道 $h'$ について、$\beta_{h'}(x)=g\,\beta_h(x)\,g^{-1}$（$g=[h'\cdot\bar h]\in\pi_1(X,x_0)$）を示せ。\\
(3) 3次元空間のひも状欠陥の「種類」を、ひもを囲むループの $\pi_1(\mathcal M,x_0)$ の元で表そうとすると、基点のとり方によって元が共役で変わりうる。ひもの種類として意味をもつのは何か。$\pi_1(\mathcal M)$ がアーベル群のときはどうか。
\end{problem}

\begin{solution}
(1) $\gamma\simeq\gamma'$（rel 端点）なら、ホモトピーの前後に $h$, $\bar h$ をつなげて $h\cdot\gamma\cdot\bar h\simeq h\cdot\gamma'\cdot\bar h$。道の積は（端点を止めたホモトピーの意味で）結合的で、$\bar h\cdot h$ は $x_1$ での定値ループとホモトピックなので
\[
\beta_h[\gamma]\beta_h[\gamma']=[h\cdot\gamma\cdot\bar h\cdot h\cdot\gamma'\cdot\bar h]=[h\cdot\gamma\cdot\gamma'\cdot\bar h]=\beta_h([\gamma][\gamma']).
\]
$\beta_{\bar h}\colon\pi_1(X,x_0)\to\pi_1(X,x_1)$ が逆写像である（$\bar h\cdot h\cdot\gamma\cdot\bar h\cdot h\simeq\gamma$）。よって同型。

(2) $h'\cdot\gamma\cdot\bar h'\simeq(h'\cdot\bar h)\cdot(h\cdot\gamma\cdot\bar h)\cdot(h\cdot\bar h')$ であり、$h\cdot\bar h'$ は $h'\cdot\bar h$ の逆の道なので、$\beta_{h'}[\gamma]=g\,\beta_h[\gamma]\,g^{-1}$。

(3) ひもを囲むループは基点を止めずに連続変形してよい（ひもの周りの小さな円周の選び方、ひもに沿った移動）。基点 $x_0$ と円周上の点を結ぶ道の選び方の曖昧さは (2) により共役で表されるので、意味をもつのは $\pi_1(\mathcal M,x_0)$ の\textbf{共役類}である（自由ループのホモトピー類と共役類は1対1に対応する）。$\pi_1$ がアーベル群なら共役類は1元からなり、ひもの種類は $\pi_1$ の元そのもの（例えば $U(1)$ 渦の巻き付き数 $n\in\Z$）で、基点のとり方によらない。非可換な場合の例は問題 \ref{ch02-prob-strings} で扱う。
\end{solution}

\begin{problem}[難易度 B]\label{ch02-prob-S1}
$p\colon\R\to S^1$、$p(s)=e^{2\pi is}$ とし、道とホモトピーの持ち上げ（定理 \ref{thm:ch02-lift} (1)(2)）を認める。$\deg[\gamma]=\tilde\gamma(1)$（$\tilde\gamma$ は $\gamma$ の $0$ からの持ち上げ）が well-defined な群同型 $\pi_1(S^1,1)\to\Z$ であることを示せ。また、定値でない写像 $S^1\to S^1$、$z\mapsto z^n$ の巻き付き数を求めよ。
\end{problem}

\begin{solution}
$p(\tilde\gamma(1))=\gamma(1)=1$ なので $\tilde\gamma(1)\in p^{-1}(1)=\Z$。

\textbf{well-defined}：$H$ を $\gamma\simeq\gamma'$（rel 端点）とし、$\tilde H(0,0)=0$ となる持ち上げ $\tilde H$ をとる。$t\mapsto\tilde H(0,t)$ は定値の道 $1$ の持ち上げで、離散集合 $\Z$ に値をとる連続関数なので一定値 $0$。同様に $t\mapsto\tilde H(1,t)\in\Z$ も一定。$\tilde H(\cdot,0)$, $\tilde H(\cdot,1)$ はそれぞれ $\gamma$, $\gamma'$ の $0$ からの持ち上げ（一意性）なので $\tilde\gamma(1)=\tilde\gamma'(1)$。

\textbf{準同型}：$m=\tilde\gamma(1)$ とする。$T_m(s)=s+m$ は $p\circ T_m=p$ をみたすので、$\tilde\gamma'+m$ は $\gamma'$ の $m$ からの持ち上げ。よって $\tilde\gamma\cdot(\tilde\gamma'+m)$ は $\gamma\cdot\gamma'$ の $0$ からの持ち上げで、終点は $m+\tilde\gamma'(1)$。

\textbf{全射}：$\gamma_n(s)=e^{2\pi ins}$ の持ち上げは $\tilde\gamma_n(s)=ns$ で、$\deg=n$。

\textbf{単射}：$\tilde\gamma(1)=0$ なら $\tilde\gamma$ は $\R$ のループで、$\R$ は凸なので $G(s,t)=(1-t)\tilde\gamma(s)$ は $\tilde\gamma$ と定値ループの間の rel 端点ホモトピー。$p\circ G$ が $\gamma\simeq$ 定値ループを与える。

以上より $\deg$ は同型である。$z\mapsto z^n$ を $s\mapsto e^{2\pi ins}$ で表すと、持ち上げは $ns$ で巻き付き数は $n$ である。物理的には、$U(1)$ のゲージ変換 $g(\theta)=e^{in\theta}$ の巻き付き数、渦 $\phi=ve^{in\theta}$ の渦度である。
\end{solution}

\begin{problem}[難易度 B]\label{ch02-prob-windint}
(1) $\gamma\colon I\to\C\setminus\{0\}$ を滑らかなループとする。被覆 $\exp\colon\C\to\C\setminus\{0\}$ に関する持ち上げ $\tilde\gamma$（$e^{\tilde\gamma}=\gamma$）を用いて、
\[
\frac{1}{2\pi i}\oint_\gamma\frac{dz}{z}=\frac{1}{2\pi}\oint_\gamma d\theta
\]
が整数で、$\gamma/\lvert\gamma\rvert\colon I\to S^1$ の巻き付き数に等しいことを示せ。\\
(2) 平面上の複素スカラー場 $\phi(z)=(z-1)^2(\bar z+1)$ について、円周 $\lvert z\rvert=R$ 上の $\phi$ の巻き付き数を $R<1$ と $R>1$ で求めよ。また $z=1$, $z=-1$ のまわりの小さな円周上の巻き付き数を求めよ。\\
(3) 連続な場 $\phi\colon\C\to\C$ の零点が有限個 $z_1,\dots,z_k$ で、各零点のまわりの巻き付き数を $n_j$ とする。$\phi$ の零点をすべて内部に含む大きな円周上の巻き付き数は $\sum_jn_j$ に等しいことを示せ。（$\pi_1(\C\setminus\{z_1,\dots,z_k\})$ が、各 $z_j$ のまわりを1周するループ $a_j$ で生成され、大きな円周がある順序の積 $a_{1}a_2\cdots a_k$ とホモトピックであることは認めてよい。）
\end{problem}

\begin{solution}
(1) $\exp\colon\C\to\C^\times$ は被覆写像（$\R\to S^1$ と $\R_{>0}$ の積）なので、持ち上げ $\tilde\gamma=\log\lvert\gamma\rvert+i\theta$（$\theta$ は連続に選んだ偏角）が存在し、$\gamma$ が滑らかなら $\tilde\gamma$ も滑らか。$dz/z$ を $\gamma$ に引き戻すと $\gamma'/\gamma\,ds=\tilde\gamma'\,ds$ なので
\[
\oint_\gamma\frac{dz}{z}=\tilde\gamma(1)-\tilde\gamma(0)=\underbrace{\log\lvert\gamma(1)\rvert-\log\lvert\gamma(0)\rvert}_{=0}+i\bigl(\theta(1)-\theta(0)\bigr).
\]
$e^{i\theta(1)}=e^{i\theta(0)}$ より $\theta(1)-\theta(0)=2\pi n$、$n\in\Z$。$\theta/2\pi$ は $\gamma/\lvert\gamma\rvert$ の $\R\to S^1$ に関する持ち上げ（を平行移動したもの）なので $n$ は巻き付き数である。また $\dfrac{dz}{z}=\dfrac{dr}{r}+i\,d\theta$ で、$d\theta=\operatorname{Im}(dz/z)=(x\,dy-y\,dx)/(x^2+y^2)$。

(2) 偏角は積で足し算になるので、巻き付き数は積について加法的である：$\deg(\phi_1\phi_2)=\deg\phi_1+\deg\phi_2$。また $\deg\bar\phi=-\deg\phi$。円周 $\lvert z\rvert=R$ 上で、$z-a$（$\lvert a\rvert\ne R$）の巻き付き数は $\lvert a\rvert<R$ なら $1$、$\lvert a\rvert>R$ なら $0$（後者は $z-a$ が円板内で零にならないので、$R\to0$ への縮小で定値に変形できる）。よって
\begin{center}
\begin{tabular}{c|ccc}
 & $(z-1)^2$ & $\bar z+1=\overline{z+1}$ & $\phi$\\ \hline
$R<1$ & $0$ & $0$ & $0$\\
$R>1$ & $2$ & $-1$ & $1$
\end{tabular}
\end{center}
$z=1$ の近く（半径 $\epsilon<2$）では $\bar z+1\neq0$ が円板内で零にならないので巻き付き数 $2$、$z=-1$ の近くでは $-1$。すなわち $z=1$ に渦度 $2$ の渦、$z=-1$ に渦度 $-1$ の反渦がある。

(3) $X=\C\setminus\{z_j\}$ とすると $\phi|_X\colon X\to\C^\times$ は準同型 $\phi_*\colon\pi_1(X)\to\pi_1(\C^\times)=\Z$ を誘導し、$\phi_*(a_j)=n_j$（$a_j$ は基点から $z_j$ の近くまで道でつないだ小円周であるが、$\Z$ はアーベル群なので道の選び方によらない）。大きな円周 $C\simeq a_1\cdots a_k$ より $\deg(\phi|_C)=\phi_*(a_1\cdots a_k)=\sum_jn_j$。(2) の例では $2+(-1)=1$ で整合する。物理的には、遠方での巻き付き数（全渦度）は内部の渦の渦度の和で、渦と反渦は対消滅できるが全渦度は変わらない。
\end{solution}

\begin{problem}[難易度 A]\label{ch02-prob-product}
(1) $\pi_1(X\times Y,(x_0,y_0))\cong\pi_1(X,x_0)\times\pi_1(Y,y_0)$ を示せ。\\
(2) $\pi_1(T^n)\cong\Z^n$ を示せ。\\
(3) 10次元時空 $\R^{1,7}\times T^2$ の閉じた弦は、トーラス方向の\term{巻き付き数}{winding number} $(w_1,w_2)\in\Z^2$ をもつ。弦の運動で $(w_1,w_2)$ が保存される理由を述べよ。また $\R^{1,8}\times S^1$ ではどうか。
\end{problem}

\begin{solution}
(1) $p_X,p_Y$ を射影とし、$\Psi[\gamma]=([p_X\circ\gamma],[p_Y\circ\gamma])$ とおく（$(p_X)_*$, $(p_Y)_*$ の組なので準同型）。逆に $X$ のループ $\alpha$ と $Y$ のループ $\beta$ から $(\alpha,\beta)(s)=(\alpha(s),\beta(s))$ が作れ、$\alpha\simeq\alpha'$（$H$）、$\beta\simeq\beta'$（$K$）なら $(H,K)$ が $(\alpha,\beta)\simeq(\alpha',\beta')$ を与えるので、$([\alpha],[\beta])\mapsto[(\alpha,\beta)]$ は well-defined。両者は互いに逆である（ループは成分の組で決まり、ホモトピーも成分ごとに与えられる）。

(2) $T^n=S^1\times\cdots\times S^1$ なので (1) と $\pi_1(S^1)=\Z$ より $\Z^n$。生成元は $j$ 番目の円周を1周するループ。

(3) 閉じた弦の配位は（各時刻で）ループ $S^1\to\R^{1,7}\times T^2$ であり、その自由ホモトピー類は $\pi_1(\R^{1,7}\times T^2)=\pi_1(T^2)=\Z^2$ の共役類、すなわち $\Z^2$ の元（アーベル群）である。弦の時間発展は連続なので、異なる時刻の弦はホモトピックで、巻き付き数 $(w_1,w_2)$ は保存される（弦の分裂・結合があっても全体の和が保存される）。$\R^{1,8}\times S^1$ では $\pi_1=\Z$ で、1つの巻き付き数 $w\in\Z$ が保存される。T双対性（第10章）では巻き付き数と運動量量子数 $n$ が入れ替わる。
\end{solution}

\begin{problem}[難易度 B]\label{ch02-prob-twopts}
$X=\R^2\setminus\{p,q\}$（$p=(-1,0)$, $q=(1,0)$）、基点 $x_0=(0,-1)$ とする。$a$ を $p$ のまわりを、$b$ を $q$ のまわりを反時計回りに1周するループとし（下図左）、$\pi_1(X,x_0)$ が $a,b$ を生成元とする自由群 $F_2$ であることを用いてよい。
\begin{center}
\begin{tikzpicture}[scale=0.85]
  \begin{scope}
    \draw[thick,red!75!black,mid arrow] (0,-1) .. controls (-0.1,0.3) and (-0.5,0.9) .. (-1,0.9)
       .. controls (-1.7,0.9) and (-2,0) .. (-1.7,-0.6) .. controls (-1.3,-1.1) and (-0.4,-1.1) .. (0,-1);
    \draw[thick,blue!75!black,mid arrow] (0,-1) .. controls (0.4,-1.1) and (1.3,-1.1) .. (1.7,-0.6)
       .. controls (2,0) and (1.7,0.9) .. (1,0.9) .. controls (0.5,0.9) and (0.1,0.3) .. (0,-1);
    \node[red!75!black] at (-2.0,0.6) {$a$};
    \node[blue!75!black] at (2.0,0.6) {$b$};
    \node[hvtx] at (-1,0) {}; \node[below] at (-1,0) {\footnotesize $p$};
    \node[hvtx] at (1,0) {};  \node[below] at (1,0) {\footnotesize $q$};
    \node[vtx] at (0,-1) {};  \node[below] at (0,-1) {\footnotesize $x_0$};
  \end{scope}
  \begin{scope}[xshift=5.2cm]
    \draw[thick,mid arrow] (0,-1) .. controls (1.6,-1.05) and (2.2,-0.6) .. (2.2,0)
       .. controls (2.2,0.8) and (1.2,1.1) .. (0,1.1) .. controls (-1.2,1.1) and (-2.2,0.8) .. (-2.2,0)
       .. controls (-2.2,-0.6) and (-1.6,-1.05) .. (0,-1);
    \node at (0,1.4) {$\gamma_1$};
    \node[hvtx] at (-1,0) {}; \node[below] at (-1,0) {\footnotesize $p$};
    \node[hvtx] at (1,0) {};  \node[below] at (1,0) {\footnotesize $q$};
    \node[vtx] at (0,-1) {};  \node[below] at (0,-1) {\footnotesize $x_0$};
  \end{scope}
  \begin{scope}[xshift=10.4cm]
    \draw[thick,mid arrow] (0,-1) .. controls (-0.1,0.3) and (-0.5,0.9) .. (-1,0.9)
       .. controls (-1.7,0.9) and (-2,0) .. (-1.7,-0.6) .. controls (-1.3,-1.1) and (-0.4,-1.1) .. (0,-1);
    \draw[thick,mid arrow] (0,-1) .. controls (0.1,0.3) and (0.5,0.9) .. (1,0.9)
       .. controls (1.7,0.9) and (2,0) .. (1.7,-0.6) .. controls (1.3,-1.1) and (0.4,-1.1) .. (0,-1);
    \node at (0,1.4) {$\gamma_2$（左の輪が先）};
    \node[hvtx] at (-1,0) {}; \node[below] at (-1,0) {\footnotesize $p$};
    \node[hvtx] at (1,0) {};  \node[below] at (1,0) {\footnotesize $q$};
    \node[vtx] at (0,-1) {};  \node[below] at (0,-1) {\footnotesize $x_0$};
  \end{scope}
\end{tikzpicture}
\zucap{左：生成元 $a,b$。中：両方の点を囲む反時計回りのループ $\gamma_1$。右：8の字のループ $\gamma_2$（左の輪をたどってから右の輪をたどる）。}
\end{center}
(1) $[\gamma_1]$, $[\gamma_2]$ を $a,b$ の語で表せ。\\
(2) $p$, $q$ のまわりの巻き付き数の組 $w=(w_p,w_q)\colon\pi_1(X)\to\Z^2$ が準同型であることを認めて、$w(\gamma_1)$, $w(\gamma_2)$ を求めよ。\\
(3) 交換子 $c=aba^{-1}b^{-1}$ は $w(c)=(0,0)$ だが、定値ループとホモトピックでないことを示せ。\\
(4) $p$, $q$ に非可換ゲージ群 $G$ の磁束管（外側では平坦な接続）があり、$a$, $b$ に沿ったホロノミーがそれぞれ $g_p,g_q\in G$ であるとする。ホロノミーが準同型 $\rho\colon\pi_1(X)\to G$ を与えることを認めて、ループ $c$ に沿って運ばれた粒子の内部状態の変化を述べよ。
\end{problem}

\begin{solution}
(1) $\gamma_1$ の上半分を、$p$ と $q$ の間の点 $(0,0)$ を通るように引き下げる（$p,q$ を横切らないのでホモトピー）。
\begin{center}
\begin{tikzpicture}[scale=0.8]
  \draw[thick,blue!75!black,mid arrow] (0,-1) .. controls (0.4,-1.1) and (1.3,-1.1) .. (1.7,-0.6)
     .. controls (2,0) and (1.7,0.9) .. (1,0.9) .. controls (0.5,0.9) and (0.15,0.4) .. (0.08,-0.95);
  \draw[thick,red!75!black,mid arrow] (-0.08,-0.95) .. controls (-0.15,0.4) and (-0.5,0.9) .. (-1,0.9)
     .. controls (-1.7,0.9) and (-2,0) .. (-1.7,-0.6) .. controls (-1.3,-1.1) and (-0.4,-1.1) .. (0,-1);
  \node[hvtx] at (-1,0) {}; \node[hvtx] at (1,0) {};
  \node[vtx] at (0,-1) {};
  \node[blue!75!black] at (2.1,0.5) {$b$};
  \node[red!75!black] at (-2.1,0.5) {$a$};
  \node at (3.6,0) {$\gamma_1\simeq b\cdot a$};
\end{tikzpicture}
\zucap{中央の縦の往復は打ち消し合うので、$\gamma_1$ はまず $b$、次に $a$ をたどるループとホモトピック。}
\end{center}
$\gamma_1$ は $x_0$ から右へ出て $q$ の右側を上り、上を左へ通って $p$ の左側を下り $x_0$ に戻る。中央に「下って上る」往復 $\delta\cdot\bar\delta$（定値の道とホモトピック）を挿入すると、前半は $q$ を反時計回りに1周する $b$、後半は $p$ を反時計回りに1周する $a$ になるので $[\gamma_1]=ba$。
$\gamma_2$ の左の輪は $a$ そのもの。右の輪は $x_0$ から中央を上り、$q$ の上を右へ通って右側を下り、下を通って戻るので、$q$ のまわりを時計回りに1周する $\bar b$。よって $[\gamma_2]=ab^{-1}$。

(2) $w(a)=(1,0)$、$w(b)=(0,1)$（$a$ は半平面 $\{x<1\}$ に含まれ、そこでは $q$ から見た偏角を $(\pi/2,3\pi/2)$ の範囲で連続に選べるので $w_q(a)=0$。$w_p(b)=0$ も同様）。よって $w(\gamma_1)=(1,1)$、$w(\gamma_2)=(1,-1)$。

(3) $\Z^2$ は可換なので $w(c)=w(a)+w(b)-w(a)-w(b)=(0,0)$。一方、自由群の元は既約語で一意に表され、$aba^{-1}b^{-1}$ は隣り合う文字が互いに逆元でない空でない既約語なので $c\neq1$。巻き付き数（第4章の $H_1$ のクラス）は可換化 $\pi_1^{\mathrm{ab}}$ の情報しかもたず、交換子は検出できない。

(4) $\rho(c)=\rho(a)\rho(b)\rho(a)^{-1}\rho(b)^{-1}=g_pg_qg_p^{-1}g_q^{-1}$。（ホロノミーの順序の規約によっては逆準同型になり $\rho(c)=g_p^{-1}g_q^{-1}g_pg_q$ となるが、議論は同じ。）$g_p$ と $g_q$ が可換でなければ $\rho(c)\ne\mathbf 1$ で、$c$ に沿って一巡した粒子の内部状態（色など）は $\rho(c)$ だけ回転し、干渉実験で検出できる。各磁束管のまわりの「巻き付き数」はゼロなのに効果が残る、非可換 AB 効果である。アーベル群 $G=U(1)$ なら $\rho(c)=1$ で、効果は巻き付き数だけで決まる。
\end{solution}

\begin{problem}[難易度 B]\label{ch02-prob-surface}
(1) 種数 $g\ge1$ の向き付け可能閉曲面 $\Sigma_g$ は、1つの0胞体、1胞体 $a_1,b_1,\dots,a_g,b_g$、語 $[a_1,b_1]\cdots[a_g,b_g]$ に沿って貼られた1つの2胞体からなる。$\pi_1(\Sigma_g)$ とその可換化を求め、$g\ne g'$ なら $\Sigma_g\not\simeq\Sigma_{g'}$ を示せ。\\
(2) $\RP^2$ は1つの0胞体、1つの1胞体 $a$、語 $a^2$ に沿って貼られた2胞体からなる。$\pi_1(\RP^2)$ を求めよ。\\
(3) $g=1$ のとき (1) の結果が問題 \ref{ch02-prob-product} の $\pi_1(T^2)=\Z^2$ と一致することを確かめよ。
\end{problem}

\begin{solution}
(1) 1骨格は $2g$ 個の円周の楔和で、ファン・カンペンの定理を繰り返し用いると $\pi_1$ は $a_1,\dots,b_g$ 上の自由群（$\R^2$ から $2g$ 点を除いた空間と同様）。2胞体の接着（命題 \ref{prop:ch02-cell}）により
\[
\pi_1(\Sigma_g)\cong\bigl\langle a_1,b_1,\dots,a_g,b_g\ \big|\ [a_1,b_1]\cdots[a_g,b_g]\bigr\rangle .
\]
可換化は $\Z^{2g}$ を関係子の像で割ったものだが、可換群では $[a_i,b_i]\mapsto a_i+b_i-a_i-b_i=0$ なので $\pi_1(\Sigma_g)^{\mathrm{ab}}\cong\Z^{2g}$。ホモトピー同値なら $\pi_1$、したがって可換化が同型で、$\Z^{2g}/2\Z^{2g}$ の位数 $2^{2g}$ を比べて $g=g'$。

(2) 1骨格は円周で $\pi_1=\langle a\rangle\cong\Z$。2胞体の境界は $a$ を2回まわるので $\pi_1(\RP^2)\cong\langle a\mid a^2\rangle\cong\Z_2$。（定理 \ref{thm:ch02-quotient} を $S^2\to\RP^2$ に用いても同じ結果が得られる。）

(3) $g=1$：$\langle a,b\mid aba^{-1}b^{-1}\rangle$ は $ab=ba$ を課した自由群で、$\Z^2$ である。

物理との関係：閉じた弦の世界面は $\Sigma_g$ で、弦の摂動展開は種数 $g$ の和になる（第10章）。$\RP^2$ は向き付けられない弦の世界面（クロスキャップ）に現れるほか、3次元の2粒子系の相対配位空間でもある（問題 \ref{ch02-prob-stat}）。
\end{solution}

\begin{problem}[難易度 C]\label{ch02-prob-lens}
$N\ge2$、$p$ を $N$ と互いに素な整数、$\omega=e^{2\pi i/N}$ とする。$S^3=\{(z_1,z_2)\in\C^2:\lvert z_1\rvert^2+\lvert z_2\rvert^2=1\}$ に、$\Z_N$ の生成元を $\tau(z_1,z_2)=(\omega z_1,\omega^pz_2)$ として作用させる。\\
(1) この作用が自由（$\tau^k\ne\id$ なら固定点をもたない）であることを示し、\term{レンズ空間}{lens space} $L(N,p)=S^3/\Z_N$ の基本群を求めよ。\\
(2) $L(2,1)$ は何と同相か。$SO(3)$ の基本群と比べよ。\\
(3) $\C^2$ への作用 $(z_1,z_2)\mapsto(\omega z_1,\omega^{-1}z_2)$ による\term{オービフォールド}{orbifold} $\C^2/\Z_N$ は原点を固定点にもち、作用は自由でない。$\C^2/\Z_N$ と $(\C^2\setminus\{0\})/\Z_N$ の基本群を求めよ。（後者は特異点の「遠方から見た形」で、第10章の弦理論のコンパクト化で使う。）
\end{problem}

\begin{solution}
(1) $1\le k\le N-1$ で $\tau^k(z)=z$ とする。$z_1\ne0$ なら $\omega^k=1$ となり矛盾。$z_1=0$ なら $\lvert z_2\rvert=1$ で $\omega^{kp}=1$、すなわち $N\mid kp$。$\gcd(N,p)=1$ より $N\mid k$ となり矛盾。よって作用は自由である。有限群のハウスドルフ空間への自由な作用は被覆空間作用である（$y$ の像 $gy$ たちは互いに異なる有限個の点なので、互いに交わらない近傍で分離できる）。$S^3$ は単連結なので、定理 \ref{thm:ch02-quotient} より
\[
\pi_1(L(N,p))\cong\Z_N .
\]
$\Z_N$ の生成元は、$(1,0)$ から $(\omega,0)$ へ至る道 $s\mapsto(e^{2\pi is/N},0)$ の像（$L(N,p)$ のループ）で表される。

(2) $N=2$、$p=1$ では $\tau=-\id$ で、$L(2,1)=S^3/\{\pm1\}=\RP^3\cong SO(3)$（第1章）。$\pi_1=\Z_2$ は $\pi_1(SO(3))=\Z_2$（問題 \ref{ch02-prob-su2so3}）と一致する。実際、$S^3\cong SU(2)$ とみなせば $\tau=-\mathbf 1$ の作用による商が $SU(2)/\{\pm\mathbf 1\}=SO(3)$ である。

(3) $\C^2/\Z_N$ は可縮である：$H([z],t)=[(1-t)z]$ は $\Z_N$ の作用（線形）と可換なので商空間上の連続写像を定め、恒等写像と原点への定値写像を結ぶ。よって $\pi_1(\C^2/\Z_N)=1$。一方、$\C^2\setminus\{0\}$ は $S^3$ に強変位レトラクトし（$z\mapsto z/\lvert z\rvert$ への直線ホモトピーは作用と可換）、$(\C^2\setminus\{0\})/\Z_N\simeq S^3/\Z_N=L(N,-1)$。$\gcd(N,-1)=1$ なので (1) より $\pi_1\cong\Z_N$。すなわち特異点を除くと、原点を囲む $S^3/\Z_N$ に $\Z_N$ だけ「ねじれた」ループが残る。このループに沿ったウィルソン線（問題 \ref{ch02-prob-wilson}）や弦の「ねじれたセクター」がオービフォールドの物理に現れる。
\end{solution}

\subsection{被覆空間と持ち上げ}

\begin{problem}[難易度 C]\label{ch02-prob-sqrt}
持ち上げの判定法（定理 \ref{thm:ch02-lift} (4)）を用いて、次を示せ。\\
(1) $\Omega$ が単連結で局所弧状連結（例えば $\R^n$ の凸開集合、$\R^3\setminus\{0\}$）なら、連続な $\phi\colon\Omega\to\C^\times$ は連続な位相をもつ：$\phi=\lvert\phi\rvert e^{i\alpha}$ となる連続関数 $\alpha\colon\Omega\to\R$ がある。特に $\R^3$ の点欠陥は複素スカラー場の位相によっては作れない。\\
(2) $\Omega=\{1<\lvert z\rvert<2\}$ とし、$\phi\colon\Omega\to\C^\times$ の円周 $\lvert z\rvert=3/2$ 上の巻き付き数を $n$ とする。連続な平方根 $\psi$（$\psi^2=\phi$）が存在するための必要十分条件は $n$ が偶数であることを示せ。\\
(3) 秩序変数が $\Psi=e^{i\alpha}\,\hat{\bm d}$（$\alpha\in\R/2\pi\Z$、$\hat{\bm d}=(\cos\beta,\sin\beta)\in S^1$）の形で、$(\alpha,\beta)$ と $(\alpha+\pi,\beta+\pi)$ が同じ物理状態を表す系（スピン1ボース凝縮の極性相などで、$\hat{\bm d}$ を面内に制限した模型）を考える。真空多様体を $\mathcal M=\R^2/\Lambda$（$\Lambda$ は $(2\pi,0)$, $(0,2\pi)$, $(\pi,\pi)$ で生成される格子）と表し、$\pi_1(\mathcal M)$ を求めよ。位相 $\alpha$ が $\pi$ だけ（すなわち「半分だけ」）まわる渦が位相的に安定に存在することを説明せよ。
\end{problem}

\begin{solution}
(1) $\C^\times\cong\R_{>0}\times S^1$ なので $\hat\phi=\phi/\lvert\phi\rvert\colon\Omega\to S^1$ を $p\colon\R\to S^1$、$p(s)=e^{is}$ に持ち上げればよい。$\Omega$ は単連結なので $\hat\phi_*\pi_1(\Omega)=1\subset p_*\pi_1(\R)$ で、判定法から持ち上げ $\alpha$（$e^{i\alpha}=\hat\phi$）が存在する。$\R^3\setminus\{0\}$ は単連結（$\simeq S^2$）なので、点欠陥のまわりで $\phi$ は常に連続な位相をもち、$\hat\phi$ は $\alpha$ の直線ホモトピーで定値に変形できる。これは $[S^2,S^1]=0$ を意味し、$U(1)$ の破れは3次元で安定な点欠陥（モノポール）を与えない。モノポールには $\pi_2(\mathcal M)\ne0$ が必要である（第3章）。

(2) $q\colon\C^\times\to\C^\times$、$q(w)=w^2$ は2重被覆で、$q_*\colon\pi_1(\C^\times)=\Z\to\Z$ は2倍写像なので $q_*\pi_1=2\Z$。$\Omega\simeq S^1$ で $\pi_1(\Omega)=\Z$ は円周 $\lvert z\rvert=3/2$ で生成され、$\phi_*(1)=n$。判定法から、$\psi$（$q\circ\psi=\phi$）が存在 $\iff n\Z\subset2\Z\iff n$ が偶数。（具体的には、$n$ が偶数なら $\phi=z^{n}\cdot\phi_0$ で $\phi_0$ の巻き付き数は $0$、$\phi_0$ は (1) と同様に連続な対数をもつので $\psi=z^{n/2}e^{\frac12\log\phi_0}$。$n$ が奇数なら $\psi$ の巻き付き数 $m$ は $2m=n$ をみたさねばならず不可能。）

(3) $(\alpha,\beta)\in\R^2$ が状態を表し、$(2\pi,0)$, $(0,2\pi)$（各角度の周期性）と $(\pi,\pi)$（同一視）の平行移動で移り合うものが同じ状態である。これらが生成する格子 $\Lambda$ は
\[
\Lambda=\{(\pi m,\pi m')\colon m\equiv m'\ (\mathrm{mod}\ 2)\}=\Z(\pi,\pi)\oplus\Z(2\pi,0)
\]
であり（$(0,2\pi)=2(\pi,\pi)-(2\pi,0)$）、$\Lambda\cong\Z^2$ は $\R^2$ に平行移動で被覆空間作用をする。$\R^2$ は単連結なので定理 \ref{thm:ch02-quotient} より
\[
\pi_1(\mathcal M)\cong\Lambda\cong\Z^2 .
\]
$\mathcal M$ のループは、$\R^2$ で $(0,0)$ から $\lambda\in\Lambda$ へ至る道の像として表され、その類が $\lambda$ である。$\lambda=(\pi,\pi)$ に対応する渦では、渦を1周すると位相 $\alpha$ が $\pi$ だけ増え、同時に $\hat{\bm d}$ が $\pi$ 回転する。$e^{i\alpha}$ も $\hat{\bm d}$ も単独では符号を変えるが、積 $\Psi$ は一価である。$(\pi,\pi)\ne0$ は $\pi_1(\mathcal M)$ の自明でない元なので、この\term{半整数渦}{half-quantum vortex}は連続変形で消せない。(2) と同じく、「$e^{i\alpha}$ だけを一価に定義しようとすると平方根の障害に出会う」ことが、半整数の巻き付きを許す仕組みである。
\end{solution}

\begin{problem}[難易度 C]\label{ch02-prob-figure8}
$X=S^1\vee S^1$ の2つの円周を $a,b$、$\pi_1(X,x_0)=F_2=\langle a,b\rangle$ とする。\\
(1) $X$ の連結な2重被覆は同型を除いてちょうど3つあることを示し、それぞれを図示せよ。（指数2の部分群は正規であることを用いてよい。）\\
(2) 下図のグラフ $\tilde X$（頂点 $v_0,v_1,v_2$）から $X$ への写像 $p$ を、頂点をすべて $x_0$ に、実線（赤）の辺を $a$ に、破線（青）の辺を $b$ に、矢印の向きを保って写すものとする。$p$ が3重被覆であることを確かめ、$H=p_*\pi_1(\tilde X,v_0)$ の生成系を求めよ。$p$ は正規被覆か。被覆変換群を求めよ。
\begin{center}
\begin{tikzpicture}[scale=1.0,thick]
  \tikzset{aedge/.style={red!75!black,mid arrow},bedge/.style={blue!75!black,dashed,mid arrow}}
  \coordinate (V0) at (0,0); \coordinate (V1) at (3,0); \coordinate (V2) at (1.5,2.3);
  \draw[aedge] (V0) to[bend right=15] (V1);
  \draw[aedge] (V1) to[bend right=15] (V2);
  \draw[aedge] (V2) to[bend right=15] (V0);
  \draw[bedge] (V1) to[bend right=55] (V2);
  \draw[bedge] (V2) to[bend right=25] (V1);
  \draw[bedge] (V0) .. controls (-1.5,0.8) and (-1.5,-0.8) .. (V0);
  \node[vtx] at (V0) {}; \node[vtx] at (V1) {}; \node[vtx] at (V2) {};
  \node[below] at (0,-0.1) {$v_0$}; \node[below right] at (3,-0.05) {$v_1$}; \node[above] at (1.5,2.4) {$v_2$};
\end{tikzpicture}
\zucap{$a$ 辺：$v_0\to v_1\to v_2\to v_0$。$b$ 辺：$v_0$ でのループ、$v_1\to v_2$、$v_2\to v_1$。}
\end{center}
(3) 2重被覆の分類が、$X$ 上の「$\Z_2$ のひねり（各ループを1周したときに波動関数の符号が変わるかどうか）」の選び方と対応することを説明せよ。
\end{problem}

\begin{solution}
(1) ガロア対応（定理 \ref{thm:ch02-galois}）により、連結な2重被覆は指数2の部分群の共役類に対応する。指数2の部分群 $H$ は正規なので、全射準同型 $\psi\colon F_2\to\Z_2$ の核と1対1に対応する（$\Z_2$ への全射は核で決まる）。$F_2$ は自由なので $\psi$ は $(\psi(a),\psi(b))\in\Z_2^2$ を任意に指定して決まり、全射なのは $(1,0),(0,1),(1,1)$ の3通り。よって連結な2重被覆はちょうど3つ。
\begin{center}
\begin{tikzpicture}[scale=0.85,thick]
  \tikzset{aedge/.style={red!75!black,mid arrow},bedge/.style={blue!75!black,dashed,mid arrow}}
  \begin{scope}
    \draw[aedge] (0,0) to[bend left=40] (2,0);
    \draw[aedge] (2,0) to[bend left=40] (0,0);
    \draw[bedge] (0,0) .. controls (-1.3,0.9) and (-1.3,-0.9) .. (0,0);
    \draw[bedge] (2,0) .. controls (3.3,-0.9) and (3.3,0.9) .. (2,0);
    \node[vtx] at (0,0) {}; \node[vtx] at (2,0) {};
    \node[below] at (0,-0.1) {\footnotesize $v_0$}; \node[below] at (2,-0.1) {\footnotesize $v_1$};
    \node at (1,-1.2) {$(\psi(a),\psi(b))=(1,0)$};
  \end{scope}
  \begin{scope}[xshift=5.4cm]
    \draw[bedge] (0,0) to[bend left=40] (2,0);
    \draw[bedge] (2,0) to[bend left=40] (0,0);
    \draw[aedge] (0,0) .. controls (-1.3,0.9) and (-1.3,-0.9) .. (0,0);
    \draw[aedge] (2,0) .. controls (3.3,-0.9) and (3.3,0.9) .. (2,0);
    \node[vtx] at (0,0) {}; \node[vtx] at (2,0) {};
    \node[below] at (0,-0.1) {\footnotesize $v_0$}; \node[below] at (2,-0.1) {\footnotesize $v_1$};
    \node at (1,-1.2) {$(0,1)$};
  \end{scope}
  \begin{scope}[xshift=10.6cm]
    \draw[aedge] (0,0) to[bend left=30] (2,0);
    \draw[aedge] (2,0) to[bend left=30] (0,0);
    \draw[bedge] (0,0) to[bend left=80] (2,0);
    \draw[bedge] (2,0) to[bend left=80] (0,0);
    \node[vtx] at (0,0) {}; \node[vtx] at (2,0) {};
    \node[left] at (-0.05,0) {\footnotesize $v_0$}; \node[right] at (2.05,0) {\footnotesize $v_1$};
    \node at (1,-1.2) {$(1,1)$};
  \end{scope}
\end{tikzpicture}
\zucap{連結な2重被覆の3つ。各頂点から出る $a$ 辺・入る $a$ 辺・出る $b$ 辺・入る $b$ 辺がちょうど1本ずつあることが、被覆であるための条件である。}
\end{center}
各場合、$\psi$ の値が $1$ の文字の辺だけが2頂点を入れ替えるので、$v_0$ からの持ち上げが閉じるのは $\psi[\gamma]=0$ のときであり、$p_*\pi_1(\tilde X,v_0)=\ker\psi$ が確かめられる。

(2) \textbf{被覆であること}：各頂点で、出る $a$ 辺・入る $a$ 辺・出る $b$ 辺・入る $b$ 辺がちょうど1本ずつある（$v_0$ の $b$ ループは出る辺と入る辺を1本ずつ与える）。よって $x_0$ の小さな近傍（4方向に出る短い線分の和）の逆像は、各頂点のまわりの同じ形の近傍3つの非交和に同相に写り、辺の内部の点でも同様。3重被覆である。

\textbf{生成系}：道の持ち上げは「各頂点で、読んでいる文字の辺を矢印の向き（逆元なら逆向き）にたどる」ことで得られる。極大な木 $T$ として $a$ 辺 $v_0\to v_1$, $v_1\to v_2$ をとると、$T$ 外の各辺 $e$ に対し「$T$ 内で $v_0$ から $e$ の始点へ、$e$ を通り、$T$ 内で $v_0$ へ戻る」ループが $\pi_1(\tilde X,v_0)$ の自由な生成系になる（第1章：グラフは $E-V+1$ 個の円周の楔和）。
\begin{center}
\begin{tabular}{lll}
$T$ 外の辺 & 対応するループ & $H$ の元\\ \hline
$a\colon v_2\to v_0$ & $v_0\xrightarrow{a}v_1\xrightarrow{a}v_2\xrightarrow{a}v_0$ & $a^3$\\
$b\colon v_0\to v_0$ & $v_0\xrightarrow{b}v_0$ & $b$\\
$b\colon v_1\to v_2$ & $v_0\xrightarrow{a}v_1\xrightarrow{b}v_2\xrightarrow{a^{-1}}v_1\xrightarrow{a^{-1}}v_0$ & $aba^{-2}$\\
$b\colon v_2\to v_1$ & $v_0\xrightarrow{a}v_1\xrightarrow{a}v_2\xrightarrow{b}v_1\xrightarrow{a^{-1}}v_0$ & $a^2ba^{-1}$
\end{tabular}
\end{center}
$p_*$ は単射なので $H=\langle a^3,\ b,\ aba^{-2},\ a^2ba^{-1}\rangle$（階数 $E-V+1=6-3+1=4$）。

\textbf{非正規性}：$b\in H$ だが、$aba^{-1}$ の $v_0$ からの持ち上げは $v_0\xrightarrow{a}v_1\xrightarrow{b}v_2\xrightarrow{a^{-1}}v_1$ で閉じないので $aba^{-1}\notin H$。よって $aHa^{-1}\ne H$ で、正規被覆ではない。

\textbf{被覆変換群}：被覆変換 $\tau$ は $b$ ループ（始点と終点が同じ $b$ 辺）を $b$ ループに写す。$b$ ループをもつ頂点は $v_0$ だけなので $\tau(v_0)=v_0$、持ち上げの一意性より $\tau=\id$。被覆変換群は自明で、$N(H)/H$ とも整合する（$H\subset N(H)\subsetneq F_2$ で指数3は素数なので $N(H)=H$）。

(3) 連結な2重被覆 $\tilde X\to X$ 上の関数で、2つのシートを入れ替える被覆変換のもとで符号を変えるもの（「反周期的な」関数）は、$X$ 上の「ループ $\gamma$ を1周すると $(-1)^{\psi[\gamma]}$ 倍になる多価関数」と同じものである。したがって (1) の3通りは、$a$, $b$ それぞれを1周したときの符号 $((-1)^{\psi(a)},(-1)^{\psi(b)})$ の自明でない選び方、すなわち $\mathrm{Hom}(\pi_1(X),\Z_2)\setminus\{0\}$ に対応する（自明な選び方 $(0,0)$ は非連結な自明被覆 $X\sqcup X$ に対応する）。これは、ねじれた境界条件、$\Z_2$ ウィルソン線、スピン構造（第7章）の選び方と同じ構造である。
\end{solution}

\subsection{リー群の基本群とスピン}

\begin{problem}[難易度 B]\label{ch02-prob-su2so3}
$U\in SU(2)$ に対し $R(U)_{ij}=\tfrac12\tr(\sigma_iU\sigma_jU^\dagger)$ とおく。\\
(1) $U\sigma_jU^\dagger=\sigma_iR(U)_{ij}$ を示し、$R(U)\in SO(3)$、$R(UV)=R(U)R(V)$ を示せ。\\
(2) $U=\exp(-i\psi\,\sigma_3/2)$ のとき $R(U)$ は $z$ 軸まわりの角度 $\psi$ の（反時計回りの）回転であることを示し、$R\colon SU(2)\to SO(3)$ が全射であることを結論せよ。\\
(3) $\ker R=\{\pm\mathbf 1\}$ を示し、$SO(3)\cong SU(2)/\Z_2$、$\pi_1(SO(3))\cong\Z_2$ を結論せよ。
\end{problem}

\begin{solution}
(1) $U\sigma_jU^\dagger$ はトレースゼロのエルミート行列なので $\sum_ic_i\sigma_i$（$c_i\in\R$）と書け、$\tr(\sigma_i\sigma_k)=2\delta_{ik}$ より $c_i=\tfrac12\tr(\sigma_iU\sigma_jU^\dagger)=R(U)_{ij}$。$R(U)$ は実である（$c_i$ が実）。$\bm x\cdot\bm\sigma\mapsto U(\bm x\cdot\bm\sigma)U^\dagger=(R\bm x)\cdot\bm\sigma$ であり、$\det(\bm x\cdot\bm\sigma)=-\lvert\bm x\rvert^2$ は共役で不変なので $\lvert R\bm x\rvert=\lvert\bm x\rvert$、$R\in O(3)$。$SU(2)$ は連結で $\det R$ は連続な $\pm1$ 値、$R(\mathbf 1)=\mathbf 1$ なので $\det R=1$。また
\[
UV\sigma_jV^\dagger U^\dagger=U\sigma_kU^\dagger R(V)_{kj}=\sigma_iR(U)_{ik}R(V)_{kj}
\]
より $R(UV)=R(U)R(V)$。

(2) $\sigma_1\sigma_3=-\sigma_3\sigma_1$ より $\sigma_1e^{i\psi\sigma_3/2}=e^{-i\psi\sigma_3/2}\sigma_1$ なので
\[
U\sigma_1U^\dagger=e^{-i\psi\sigma_3}\sigma_1=(\cos\psi-i\sin\psi\,\sigma_3)\sigma_1=\cos\psi\,\sigma_1+\sin\psi\,\sigma_2
\]
（$\sigma_3\sigma_1=i\sigma_2$）。同様に $U\sigma_2U^\dagger=-\sin\psi\,\sigma_1+\cos\psi\,\sigma_2$、$U\sigma_3U^\dagger=\sigma_3$。よって $R(U)$ は $z$ 軸まわりの角度 $\psi$ の回転。任意の軸 $\hat{\bm n}$ には $\hat{\bm n}=R(W)\hat{\bm z}$ となる $W$ をとれば（$z$ 軸と $y$ 軸まわりの回転の合成で任意の方向に行ける）、$W e^{-i\psi\sigma_3/2}W^\dagger=e^{-i\psi\,\hat{\bm n}\cdot\bm\sigma/2}$ が $\hat{\bm n}$ まわりの角度 $\psi$ の回転に写る。すべての回転はある軸まわりの回転なので $R$ は全射。

(3) $R(U)=\mathbf 1$ なら $U$ はすべての $\sigma_j$ と可換で、シューアの補題（あるいは直接計算）から $U=\lambda\mathbf 1$、$\det U=\lambda^2=1$ より $U=\pm\mathbf 1$。準同型定理より $SO(3)\cong SU(2)/\{\pm\mathbf 1\}$（コンパクト→ハウスドルフの連続全単射なので同相でもある）。$\{\pm\mathbf 1\}$ は $SU(2)\cong S^3$ に自由に作用し、$S^3$ は単連結なので、定理 \ref{thm:ch02-quotient} より $\pi_1(SO(3))\cong\Z_2$。$SU(2)$ が $SO(3)$ の普遍被覆である。
\end{solution}

\begin{problem}[難易度 B]\label{ch02-prob-belt}
$R_z(\psi)$ を $z$ 軸まわりの角度 $\psi$ の回転とし、$SO(3)$ のループ $\gamma_m(t)=R_z(2\pi mt)$（$t\in I$、$m\in\Z$）を考える。\\
(1) 問題 \ref{ch02-prob-su2so3} の被覆 $SU(2)\to SO(3)$ を用いて、$\gamma_m$ の $\mathbf 1$ からの持ち上げを求め、$\gamma_1$（$2\pi$ 回転）は1点に縮められず、$\gamma_2$（$4\pi$ 回転）は縮められることを示せ。\\
(2) $SO(3)$ を半径 $\pi$ の球の境界の対蹠点を同一視したもの（第1章）として、$\gamma_1$ と $\gamma_2$ を図示し、$\gamma_2$ を縮める変形を図で説明せよ。\\
(3) 皿を手のひらに載せて腕をねじらずに1回転させると腕がねじれるが、同じ向きにもう1回転させると元に戻る（\term{ディラックのベルト・トリック}{Dirac belt trick}）。これが (1) とどう関係するか説明せよ。
\end{problem}

\begin{solution}
(1) $U_m(t)=\exp(-i\pi mt\,\sigma_3)$ は $R(U_m(t))=R_z(2\pi mt)$ をみたし（問題 \ref{ch02-prob-su2so3} (2) で $\psi=2\pi mt$）、$U_m(0)=\mathbf 1$ なので $\gamma_m$ の持ち上げである。終点は $U_m(1)=e^{-i\pi m\sigma_3}=(-1)^m\mathbf 1$。定理 \ref{thm:ch02-lift} (3) より、$\gamma_m$ が縮められる（$p_*\pi_1(SU(2))=1$ に属する）$\iff$ 持ち上げが閉じる $\iff m$ が偶数。よって $\gamma_1$ は縮められず、$\gamma_2$ は縮められる。より具体的に、$\gamma_2$ の持ち上げ $U_2(t)=\cos(2\pi t)-i\sin(2\pi t)\sigma_3$ は $S^3$ の大円で、$S^3$ 上で1点に縮められ、その像が $\gamma_2$ の縮約を与える。$[\gamma_1]$ が $\pi_1(SO(3))=\Z_2$ の生成元で、$[\gamma_m]=m\bmod2$。

(2) 軸・角度表示で、$\psi\in[0,\pi]$ では $R_z(\psi)$ は点 $\psi\hat{\bm z}$、$\psi\in[\pi,2\pi]$ では $R_z(\psi)=R_{-z}(2\pi-\psi)$ は点 $-(2\pi-\psi)\hat{\bm z}$ である。よって $\gamma_1$ は中心から北極 $N=\pi\hat{\bm z}$ へ上り、対蹠点の南極 $S$ に「跳んで」中心へ戻る。
\begin{center}
\begin{tikzpicture}[scale=1.0]
  \begin{scope}
    \draw[thick] (0,0) circle (1.4);
    \draw[very thick,red!70!black,mid arrow] (0,0) -- (0,1.4);
    \draw[very thick,red!70!black,mid arrow] (0,-1.4) -- (0,0);
    \node[vtx] at (0,0) {}; \node[right] at (0,0) {\footnotesize $\mathbf 1$};
    \node[vtx,red!70!black] at (0,1.4) {}; \node[above] at (0,1.4) {\footnotesize $N$};
    \node[vtx,red!70!black] at (0,-1.4) {}; \node[below] at (0,-1.4) {\footnotesize $S$};
    \node at (0,-2.2) {(a) $\gamma_1$：$2\pi$ 回転};
  \end{scope}
  \begin{scope}[xshift=4.2cm]
    \draw[thick] (0,0) circle (1.4);
    \draw[very thick,red!70!black,mid arrow] (0,0) -- (0,1.4);
    \draw[very thick,red!70!black,mid arrow] (0,-1.4) -- (0,0);
    \draw[very thick,blue!70!black,mid arrow] (0,0) -- (40:1.4);
    \draw[very thick,blue!70!black,mid arrow] (220:1.4) -- (0,0);
    \node[vtx] at (0,0) {};
    \node[vtx,blue!70!black] at (40:1.4) {}; \node[above right] at (40:1.4) {\footnotesize $Q$};
    \node[vtx,blue!70!black] at (220:1.4) {}; \node[below left] at (220:1.4) {\footnotesize $-Q$};
    \draw[-{Stealth[length=2mm]},gray,thick] (32:1.7) arc (32:-5:1.7);
    \draw[-{Stealth[length=2mm]},gray,thick] (212:1.7) arc (212:175:1.7);
    \node at (0,-2.2) {(b) $\gamma_2$ の2周目を傾ける};
  \end{scope}
  \begin{scope}[xshift=8.4cm]
    \draw[thick] (0,0) circle (1.4);
    \draw[very thick,red!70!black,mid arrow] (-0.07,0) -- (-0.07,1.4);
    \draw[very thick,red!70!black,mid arrow] (-0.07,-1.4) -- (-0.07,0);
    \draw[very thick,blue!70!black,mid arrow] (0.07,0) -- (0.07,-1.4);
    \draw[very thick,blue!70!black,mid arrow] (0.07,1.4) -- (0.07,0);
    \node[vtx] at (0,0) {}; \node[vtx,red!70!black] at (0,1.4) {}; \node[above] at (0,1.4) {\footnotesize $N=-Q$};
    \node[vtx,red!70!black] at (0,-1.4) {}; \node[below] at (0,-1.4) {\footnotesize $S=Q$};
    \node at (0,-2.2) {(c) $Q\to S$：往復に縮む};
  \end{scope}
\end{tikzpicture}
\zucap{$SO(3)$ を半径 $\pi$ の球（境界の対蹠点を同一視）で表す。(a) $2\pi$ 回転は中心→$N$、$S$→中心。対蹠点の組 $\{N,S\}$ は境界上で動かせないので縮まない。(b) $4\pi$ 回転の2周目の直径を $Q,-Q$ へ傾ける。(c) $Q$ を $S$ まで（$-Q$ を $N$ まで）動かすと、ループは赤（1周目）と青（2周目）が同じ直径を逆向きにたどる往復になり、1点に縮む。}
\end{center}
$\gamma_2$ は直径を2回たどる：中心→$N\equiv S$→中心→$N\equiv S$→中心。2周目の直径を（端点の対蹠関係を保ったまま）別の直径 $Q,-Q$ へ傾けたループ「中心→$N\equiv S$→中心→$Q\equiv-Q$→中心」は $\gamma_2$ とホモトピックである。$Q$ を境界に沿って $S$ まで動かす（$-Q$ は $N$ へ）と、ループは「中心→$N\equiv S$→中心→$S\equiv N$→中心」となる。この中の「$S$→中心→$S$」は往復で $S$ に縮み、残る「中心→$N$→中心」も往復なので、ループ全体が1点に縮む。$\gamma_1$ では「跳び」が1回しかないので、このような相殺ができない。

(3) 腕の各点（肩から手のひらまで）での向きの変化を $SO(3)$ の元と見ると、腕は $\mathbf 1$（肩）から皿の向き $R$ への $SO(3)$ の道である。皿を $2\pi$ 回転させる操作の各段階で腕の配位が決まるので、操作全体は「肩から手のひらまでの道」の1パラメータ族、すなわちループ $\gamma_1$ を縮めようとするホモトピーに相当し、(1) よりそれは不可能（腕はねじれたまま）。$4\pi$ 回転では $\gamma_2$ が縮められるので、腕を連続に動かしてねじれのない状態に戻せる。スピン $1/2$ の状態が $2\pi$ 回転で $-1$ 倍、$4\pi$ 回転で元に戻るのは、この $\Z_2$ の現れである（問題 \ref{ch02-prob-spinor}）。
\end{solution}

\begin{problem}[難易度 A]\label{ch02-prob-spinor}
スピン $j$ の状態に対する $z$ 軸まわりの角度 $\psi$ の回転は $\mathcal U(\psi)=\exp(-i\psi J_z)$ で表される。\\
(1) $\mathcal U(2\pi)=(-1)^{2j}$ を示せ。\\
(2) スピン上向きの中性子ビームを2つに分け、一方だけを磁場中でスピン歳差により $z$ 軸まわりに角度 $\psi$ 回転させてから再び重ねる。干渉強度が $\psi$ について周期 $4\pi$ であることを示せ。\\
(3) 物理的状態は射線（全体位相を除いた状態ベクトル）なので、$\mathcal U(2\pi)=-1$ は回転群 $SO(3)$ の\textbf{射影表現}である。この符号が $\pi_1(SO(3))=\Z_2$ と関係する理由を述べよ。また、整数スピンと半整数スピンの状態の重ね合わせを考えると何が起こるか。
\end{problem}

\begin{solution}
(1) $J_z$ の固有値は $m\in\{-j,-j+1,\dots,j\}$ で、$e^{-2\pi im}$ は $m$ が整数なら $1$、半整数なら $-1$。$2j$ と $2m$ の偶奇は一致するので $\mathcal U(2\pi)=(-1)^{2j}$。$j=1/2$ では $\mathcal U(\psi)=e^{-i\psi\sigma_3/2}$、$\mathcal U(2\pi)=-\mathbf 1$。

(2) 出口での状態は（規格化定数を除いて）$\tfrac12\bigl(\mathbf 1+e^{-i\psi\sigma_3/2}\bigr)\lvert\uparrow\rangle=\tfrac12(1+e^{-i\psi/2})\lvert\uparrow\rangle$ なので、強度は
\[
I(\psi)\propto\Bigl\lvert\tfrac12(1+e^{-i\psi/2})\Bigr\rvert^2=\frac{1+\cos(\psi/2)}{2}.
\]
$\psi=2\pi$ で $0$（完全に打ち消し合う）、周期は $4\pi$。これは中性子干渉実験（1975年）で観測された。

(3) $SO(3)$ の各元に $\mathcal U$ を連続に割り当てようとすると、ループ $\gamma_1$（$2\pi$ 回転）に沿って $\mathcal U$ を連続に追いかけた結果が $\mathcal U(2\pi)=(-1)^{2j}$ である。もし $\gamma_1$ が縮められれば、連続性から $\mathcal U(2\pi)$ は $\mathcal U(0)=\mathbf 1$ に連続変形でき、離散的な値 $\pm1$ をとるので $+1$ でなければならない。$\pi_1(SO(3))=\Z_2$ だからこそ $-1$ が許され、$4\pi$ 回転（$\gamma_2\simeq$ 定値）では必ず $+1$ である。言いかえると、半整数スピンは $SO(3)$ ではなく普遍被覆 $SU(2)$ の真の表現である。整数スピン $\lvert a\rangle$ と半整数スピン $\lvert b\rangle$ の重ね合わせ $\lvert a\rangle+\lvert b\rangle$ は $2\pi$ 回転で $\lvert a\rangle-\lvert b\rangle$ になり、物理的に同じ状態であるべきなのに異なる射線になる。したがってこのような重ね合わせは観測できない（ボソン・フェルミオンの\term{超選択則}{superselection rule}）。
\end{solution}

\begin{problem}[難易度 B]\label{ch02-prob-spin4}
$\R^4$ の点 $x=(x_1,x_2,x_3,x_4)$ に $2\times2$ 行列 $X=x_4\mathbf 1+i\,x_k\sigma_k$ を対応させる（$\det X=\lvert x\rvert^2$）。\\
(1) $(q_L,q_R)\in SU(2)\times SU(2)$ に対し $X\mapsto q_LXq_R^\dagger$ が $SO(4)$ の元 $\Lambda(q_L,q_R)$ を定め、$\Lambda$ が準同型であることを示せ。\\
(2) $\ker\Lambda=\{(\mathbf 1,\mathbf 1),(-\mathbf 1,-\mathbf 1)\}$ を示せ。\\
(3) $\Lambda$ が全射であることを（$SO(4)$ の連結性と次元を用いて）示し、$\pi_1(SO(4))\cong\Z_2$、$\Spin(4)\cong SU(2)\times SU(2)$ を結論せよ。ユークリッド化したローレンツ群の表現 $(j_L,j_R)$ との関係を述べよ。
\end{problem}

\begin{solution}
(1) $X$ の形の行列全体 $\mathcal Q=\{a\mathbf 1+i\bm b\cdot\bm\sigma\}$ は $\R_{\ge0}\cdot SU(2)$ に等しく（第1章の $SU(2)$ の表示）、$SU(2)$ の元を左右から掛けても $\mathcal Q$ に留まる。$X\mapsto q_LXq_R^\dagger$ は $\R$ 線形で $\det(q_LXq_R^\dagger)=\det X$、すなわち $\lvert x\rvert^2$ を保つので $\Lambda\in O(4)$。$SU(2)\times SU(2)$ は連結で $\Lambda(\mathbf 1,\mathbf 1)=\mathbf 1$ なので $\det\Lambda=1$。$(q_Lq_L')X(q_Rq_R')^\dagger=q_L(q_L'Xq_R'^\dagger)q_R^\dagger$ より準同型。

(2) $q_LXq_R^\dagger=X$ がすべての $X$ で成り立つとする。$X=\mathbf 1$ とおいて $q_L=q_R=:q$。すると $qXq^\dagger=X$ がすべての $X$、特に $X=i\sigma_k$ で成り立ち、$q$ はすべての $\sigma_k$ と可換なので $q=\pm\mathbf 1$。逆に $(-\mathbf 1,-\mathbf 1)$ は自明に作用する。

(3) 核が離散なので微分 $d\Lambda\colon\mathfrak{su}(2)\oplus\mathfrak{su}(2)\to\mathfrak{so}(4)$ は単射で、両辺とも6次元なので同型。逆関数定理により像は $\mathbf 1$ の近傍を含み、したがって開部分群である。開部分群は（剰余類の和として補集合も開なので）閉でもあり、$SO(4)$ は連結（第1章）なので像は $SO(4)$ 全体。よって $SO(4)\cong(SU(2)\times SU(2))/\Z_2$。$SU(2)\times SU(2)\cong S^3\times S^3$ は単連結で、$\Z_2=\{\pm(\mathbf 1,\mathbf 1)\}$ は自由に作用するので $\pi_1(SO(4))\cong\Z_2$、普遍被覆 $\Spin(4)=SU(2)\times SU(2)$。

物理では、4次元ユークリッド空間の回転群の表現は $SU(2)_L\times SU(2)_R$ の表現 $(j_L,j_R)$ で分類され、$(\tfrac12,0)$ と $(0,\tfrac12)$ が左巻き・右巻きのワイル・スピノル、$(\tfrac12,\tfrac12)$ がベクトル $X$ そのものである。$(-\mathbf 1,-\mathbf 1)$ は $(-1)^{2(j_L+j_R)}$ で作用し、$j_L+j_R$ が半整数の表現（スピノル）は $SO(4)$ の真の表現ではない。また2形式は $(1,0)\oplus(0,1)$ に分かれ、自己双対・反自己双対2形式に対応する（インスタントン、第3, 8章）。
\end{solution}

\begin{problem}[難易度 B]\label{ch02-prob-un}
$N\ge2$ とし、$\pi_1(SU(N))=1$（第3章）を認める。\\
(1) $\det_*\colon\pi_1(U(N))\to\pi_1(U(1))=\Z$ が同型であることを示し、生成元を与えよ。ループ $t\mapsto e^{2\pi it}\mathbf 1_N$ は $\pi_1(U(N))$ のどの元か。\\
(2) $\Z_N=\{\omega^k\mathbf 1:\omega=e^{2\pi i/N}\}$ を $SU(N)$ の中心とする。$\pi_1(SU(N)/\Z_N)\cong\Z_N$ を示し、生成元を $SU(N)$ の道として具体的に書け。\\
(3) 随伴表現の場（グルーオン）だけを含む $SU(N)$ ゲージ理論では、中心 $\Z_N$ は自明に作用するので、実質的なゲージ群は $SU(N)/\Z_N$ である。このとき $\Z_N$ に値をとる「中心磁束」をもつ渦（ひも）がありうることを説明せよ。基本表現のクォークを加えると何が変わるか。
\end{problem}

\begin{solution}
(1) 第1章より同相 $\Phi\colon U(N)\to SU(N)\times U(1)$、$\Phi(U)=(U\,D(\det U)^{-1},\det U)$（$D(z)=\operatorname{diag}(z,1,\dots,1)$）があり、$\det=p_2\circ\Phi$。よって $\pi_1(U(N))\cong\pi_1(SU(N))\times\pi_1(U(1))=1\times\Z$ で、この同型のもとで $\det_*$ は第2成分への射影、すなわち同型である。生成元は $\Phi^{-1}(\mathbf 1,e^{2\pi it})=D(e^{2\pi it})=\operatorname{diag}(e^{2\pi it},1,\dots,1)$。$\det(e^{2\pi it}\mathbf 1_N)=e^{2\pi iNt}$ なので、$e^{2\pi it}\mathbf 1_N$ は生成元の $N$ 倍である。

(2) $\Z_N$ は有限群で、$SU(N)$ に左からの掛け算で自由に作用する（$\omega^kU=U\Rightarrow\omega^k=1$）。$SU(N)$ は単連結なので、定理 \ref{thm:ch02-quotient} より $\pi_1(SU(N)/\Z_N)\cong\Z_N$。$k\in\Z_N$ に対応する元は、$\mathbf 1$ から $\omega^k\mathbf 1$ への $SU(N)$ の道の像である。$k=1$ には
\[
u(t)=\operatorname{diag}\bigl(e^{2\pi it/N},\dots,e^{2\pi it/N},e^{-2\pi i(N-1)t/N}\bigr)\quad(t\in I)
\]
をとればよい：$\det u(t)=e^{2\pi i t(N-1)/N}e^{-2\pi i(N-1)t/N}=1$、$u(0)=\mathbf 1$、$u(1)=\operatorname{diag}(\omega,\dots,\omega,\omega^{-(N-1)})=\omega\mathbf 1$（$\omega^{-(N-1)}=\omega\cdot\omega^{-N}=\omega$）。$N=2$ では $SU(2)/\Z_2=SO(3)$ で問題 \ref{ch02-prob-su2so3} と一致する。

(3) ひもを囲む円周上で、ゲージ変換が「$\mathbf 1$ から中心元 $\omega^k$ へ行く $SU(N)$ の道」になっている配位は、$SU(N)/\Z_N$ ではループであり、(2) の $\pi_1=\Z_N$ の非自明な元を表す。随伴表現の場にとってはこのゲージ変換は一価なので、こうした配位は許され、$\Z_N$ の電荷をもつ渦（\term{中心渦}{center vortex}、トフーフト磁束）として現れる。これらはウィルソン・ループに $\omega^k$ の位相を与え、閉じ込めの描像で重要である。基本表現のクォーク $\psi\mapsto g\psi$ に対しては $\omega^k\ne1$ が自明に作用しないので、ゲージ群は $SU(N)$ のままで $\pi_1=1$、中心渦は位相的に安定でなくなる（ひもの切断）。
\end{solution}

\subsection{位相欠陥：渦とひも}

\begin{problem}[難易度 D]\label{ch02-prob-strings}
(1) $G$ を連結かつ単連結なリー群、$H\subset G$ を離散部分群とする。$H$ は $G$ に右からの掛け算で被覆空間作用をすること（離散部分群なので認めてよい）を用いて、$\pi_1(G/H)\cong H$ を示せ。\\
(2) $U(1)$ ゲージ理論で電荷 $N$ のヒッグス場 $\phi\mapsto e^{iN\lambda}\phi$（$e^{i\lambda}\in U(1)$）が凝縮すると、$U(1)\to\Z_N$ と破れる。$U(1)$ の普遍被覆 $\R$ を用いて $\pi_1(\mathcal M)$ を求め、基本的な渦の磁束が $2\pi/N$ であることを示せ。（単位電荷の粒子の AB 位相は $e^{2\pi i/N}$ になる。）\\
(3) 2つの随伴表現ヒッグス場が平行でない期待値をもち、$SO(3)$ が完全に破れる場合の $\pi_1(\mathcal M)$ を求めよ。2本の基本的なひもを合わせるとどうなるか。\\
(4) \textbf{2軸性ネマティック液晶}：分子の向きは直交する3つの軸（向きは区別しない）で表され、$\mathcal M=SO(3)/D_2$、$D_2=\{\mathbf 1,R_x(\pi),R_y(\pi),R_z(\pi)\}$ である。$\pi_1(\mathcal M)$ が四元数群 $Q_8=\{\pm\mathbf 1,\pm i\sigma_1,\pm i\sigma_2,\pm i\sigma_3\}$ であることを示し、共役類を求めよ。ひもの種類はいくつあるか。\\
(5) (4) で $i\sigma_1$ 型と $i\sigma_2$ 型のひもが交差してすり抜けようとすると、両者をつなぐひもが残る。そのひもの種類を交換子から求めよ。
\end{problem}

\begin{solution}
(1) 被覆空間作用なので $G\to G/H$ は被覆写像で、被覆変換群は $H$（$gH$ の点を $g\mapsto gh$ で動かす）。$G$ は単連結なので定理 \ref{thm:ch02-quotient} より $\pi_1(G/H)\cong H$。具体的には、$h\in H$ に $\mathbf 1$ から $h$ への $G$ の道の像（$G/H$ のループ）が対応し、積は道の連結 $\tilde\gamma_{h_1}\cdot(h_1\tilde\gamma_{h_2})$ で $h_1h_2$ に対応する。

(2) $U(1)=\R/2\pi\Z$ を $e^{i\lambda}$（$\lambda\in\R$）で表す。普遍被覆 $\tilde G=\R$ で、ヒッグス場の固定部分群は $e^{iN\lambda}=1$、すなわち $\tilde H=\{\lambda\in\tfrac{2\pi}{N}\Z\}\cong\Z$。$\mathcal M=\R/\tfrac{2\pi}{N}\Z\cong S^1$ で、(1) より $\pi_1(\mathcal M)\cong\tilde H\cong\Z$。生成元は $\lambda$ が $0$ から $2\pi/N$ まで動く道で、渦を囲む円周上でゲージ変換が $\lambda(\theta)=\theta/N$ のように巻くことに対応する。遠方で $D\phi=(d-iNA)\phi=0$ と $\phi=v\,e^{iN\lambda(\theta)}=ve^{i\theta}$ から $NA=d\theta$、すなわち $A=d\theta/N$ で、磁束 $\oint A=2\pi/N$。単位電荷の粒子の AB 位相は $e^{i\oint A}=e^{2\pi i/N}$ で、遠方からは電荷 $\bmod N$ だけが見える（離散ゲージ理論 $\Z_N$）。

(3) $SO(3)$ を普遍被覆 $SU(2)$ で考えると、$SO(3)$ の固定部分群 $\{1\}$ の逆像は $\tilde H=\{\pm\mathbf 1\}$。$\mathcal M=SU(2)/\{\pm\mathbf 1\}=SO(3)$ で $\pi_1(\mathcal M)\cong\Z_2$（問題 \ref{ch02-prob-su2so3}）。基本的なひもの類は $-\mathbf 1$ で、$(-\mathbf 1)^2=\mathbf 1$ なので2本のひもを合わせると位相的に自明になり、対消滅できる（\term{$\Z_2$ ひも}{$\Z_2$ string}）。

(4) 同様に $SU(2)$ に持ち上げる。$R_k(\pi)$（$k$ 軸まわりの $\pi$ 回転）の逆像は $\pm\exp(-i\pi\sigma_k/2)=\mp i\sigma_k$ なので、$D_2$ の逆像は $\tilde H=\{\pm\mathbf 1,\pm i\sigma_1,\pm i\sigma_2,\pm i\sigma_3\}=Q_8$（$(i\sigma_1)(i\sigma_2)=-\sigma_1\sigma_2=-i\sigma_3$ などで閉じている）。$\mathcal M=SO(3)/D_2=SU(2)/Q_8$ で、(1) より $\pi_1(\mathcal M)\cong Q_8$。

共役類：$\pm\mathbf 1$ は中心なので $\{\mathbf 1\}$, $\{-\mathbf 1\}$。$(i\sigma_2)(i\sigma_1)(i\sigma_2)^{-1}=(i\sigma_2)(i\sigma_1)(-i\sigma_2)=i\,\sigma_2\sigma_1\sigma_2=-i\sigma_1$ なので $\{\pm i\sigma_1\}$、同様に $\{\pm i\sigma_2\}$, $\{\pm i\sigma_3\}$。共役類は5つで、自明なものを除いて4種類のひもがある（$\pm i\sigma_k$ は $k$ 軸以外の2軸が $\pi$ だけ回る「半整数」の回位、$-\mathbf 1$ は $2\pi$ 回位）。

(5) 2本のひもが交差してすり抜けると、交換子 $c=g_1g_2g_1^{-1}g_2^{-1}$ の類をもつひもが両者をつなぐ。$g_1=i\sigma_1$, $g_2=i\sigma_2$ のとき $g_1^{-1}=-i\sigma_1$, $g_2^{-1}=-i\sigma_2$ で
\[
c=(i\sigma_1)(i\sigma_2)(-i\sigma_1)(-i\sigma_2)=(\sigma_1\sigma_2)^2=(i\sigma_3)^2=-\mathbf 1 .
\]
$-\mathbf 1\ne\mathbf 1$ なので、2本のひもは自由にすり抜けられず、$-\mathbf 1$ 型（$2\pi$ 回位）のひもでつながれる。$\pi_1$ が可換なら交換子は常に自明で、このような現象は起こらない（問題 \ref{ch02-prob-twopts} の非可換 AB 効果と同じ構造）。
\end{solution}

\subsection{アハラノフ・ボーム効果と経路積分}

\begin{problem}[難易度 B]\label{ch02-prob-ring}
半径 $R$ の円周上に束縛された質量 $m$、電荷 $q$ の粒子があり、円周の内部を磁束 $\Phi$ が貫いている（円周上で $\bm B=0$）。$A=\frac{\Phi}{2\pi}d\theta$ とする。\\
(1) ハミルトニアン $H=\dfrac{1}{2mR^2}\bigl(-i\partial_\theta-\tfrac{q\Phi}{2\pi}\bigr)^2$ のエネルギー準位が
\[
E_n=\frac{1}{2mR^2}\Bigl(n-\frac{q\Phi}{2\pi}\Bigr)^2\qquad(n\in\Z)
\]
であることを示せ。\\
(2) スペクトル全体は $\Phi\to\Phi+2\pi/q$ で不変であることを示し、基底エネルギーを $\Phi$ の関数として図示せよ。\\
(3) ゲージ変換 $\psi\mapsto e^{iq\lambda}\psi$、$A\mapsto A+d\lambda$ で $A$ を消そうとして $\lambda=-\frac{\Phi}{2\pi}\theta$ とすると何が問題か。$q\Phi\in2\pi\Z$ のときはどうか。\\
(4) 円周上で $\bm B=0$ なのにエネルギーが $\Phi$ に依存することは、何を意味するか。
\end{problem}

\begin{solution}
(1) 共変微分 $D_\theta=\partial_\theta-iqA_\theta=\partial_\theta-iq\Phi/2\pi$ より、運動エネルギー $-\frac{1}{2mR^2}D_\theta^2=\frac{1}{2mR^2}(-iD_\theta)^2$ が与えられたハミルトニアンである。波動関数は円周上の一価関数なので $\psi_n=e^{in\theta}/\sqrt{2\pi}$（$n\in\Z$）が完全系をなし、$(-i\partial_\theta-\tfrac{q\Phi}{2\pi})\psi_n=(n-\tfrac{q\Phi}{2\pi})\psi_n$。よって $E_n$ が得られる。

(2) $\Phi\to\Phi+2\pi/q$ で $E_n\to\frac{1}{2mR^2}(n-1-\frac{q\Phi}{2\pi})^2=E_{n-1}$ なので、スペクトル（集合）は不変。基底エネルギー $E_0(\Phi)=\frac{1}{2mR^2}\min_n(n-\alpha)^2$（$\alpha=q\Phi/2\pi$）は放物線を周期的につないだ形で、$\alpha\in\Z+\tfrac12$ で2つの準位 $n=\alpha\pm\tfrac12$ が交差する。
\begin{center}
\begin{tikzpicture}[xscale=1.3,yscale=2.4]
  \draw[-{Stealth[length=2mm]}] (-2.3,0) -- (2.5,0) node[right] {\footnotesize $\alpha=\frac{q\Phi}{2\pi}$};
  \draw[-{Stealth[length=2mm]}] (0,-0.05) -- (0,0.75) node[above] {\footnotesize $2mR^2E$};
  \foreach \n in {-2,-1,0,1,2}{
    \draw[gray,domain={\n-1.0}:{\n+1.0},samples=40] plot (\x,{(\x-\n)*(\x-\n)*0.6});
  }
  \foreach \n in {-2,-1,0,1,2}{
    \draw[very thick,red!70!black,domain={\n-0.5}:{\n+0.5},samples=30] plot (\x,{(\x-\n)*(\x-\n)*0.6});
    \node[below] at (\n,0) {\footnotesize $\n$};
  }
\end{tikzpicture}
\zucap{各準位 $(n-\alpha)^2$（灰色）と基底エネルギー（赤）。周期1（$\Phi$ について周期 $2\pi/q$）で、半整数で準位が交差する。}
\end{center}

(3) $\lambda=-\frac{\Phi}{2\pi}\theta$ なら $A+d\lambda=0$ となるが、$\psi'=e^{-iq\Phi\theta/2\pi}\psi$ は $\theta\to\theta+2\pi$ で $e^{-iq\Phi}$ 倍になり、一価でない。すなわち $A$ を消す代わりに、波動関数の境界条件が $\psi'(\theta+2\pi)=e^{-iq\Phi}\psi'(\theta)$ にねじれる（このとき $\psi'_n\propto e^{i(n-\alpha)\theta}$ で同じスペクトルを与える）。$g=e^{iq\lambda}$ が円周から $U(1)$ への写像として well-defined なのは $q\Phi\in2\pi\Z$ のときに限り、このとき $g=e^{-ik\theta}$（$k=q\Phi/2\pi\in\Z$）は巻き付き数 $-k$ の「大きなゲージ変換」で、$A$ は真にゲージ同値なものとして消え、スペクトルは $\Phi=0$ のものと一致する（(2) の周期性）。

(4) 円周上の局所的な場の強さ $F=0$ だけでは物理は決まらず、ゲージ不変量 $e^{iq\oint A}=e^{iq\Phi}$（ホロノミー、ウィルソン・ループ）が観測可能である。円周は単連結でない（$\pi_1=\Z$）ので、$F=0$ でも $A$ は大域的に純ゲージとは限らない。これが AB 効果の本質であり、持続電流 $I=-\partial E_0/\partial\Phi$ として実際に観測される。
\end{solution}

\begin{problem}[難易度 D]\label{ch02-prob-pathint}
前問の系を経路積分で扱う。慣性モーメントを $\mathcal I=mR^2$、$\alpha=q\Phi/2\pi$ とし、ユークリッド作用を $S_E=\int_0^\tau d\tau'\bigl[\tfrac{\mathcal I}{2}\dot\theta^2-i\alpha\dot\theta\bigr]$ とする。\\
(1) 円周 $S^1$ の普遍被覆 $\R$ への持ち上げを用いると、$\theta'$ から $\theta$ への経路は巻き付き数 $n\in\Z$（$\R$ 上で $\theta'$ から $\theta+2\pi n$ へ行く）で分類される。熱核 $K(\theta,\theta';\tau)=\langle\theta|e^{-H\tau}|\theta'\rangle$ が
\[
K(\theta,\theta';\tau)=\sum_{n\in\Z}e^{i\alpha(\theta-\theta'+2\pi n)}\sqrt{\frac{\mathcal I}{2\pi\tau}}\exp\Bigl(-\frac{\mathcal I(\theta-\theta'+2\pi n)^2}{2\tau}\Bigr)
\]
と書けることを説明せよ（$\R$ 上の自由粒子の熱核 $\sqrt{\mathcal I/2\pi\tau}\,e^{-\mathcal I x^2/2\tau}$ は既知としてよい）。\\
(2) ポアソンの和公式 $\sum_ng(x+2\pi n)=\frac{1}{2\pi}\sum_k\hat g(k)e^{ikx}$（$\hat g(k)=\int_\R g(y)e^{-iky}dy$）を用いて
\[
K(\theta,\theta';\tau)=\frac{1}{2\pi}\sum_{k\in\Z}e^{ik(\theta-\theta')}\exp\Bigl(-\frac{(k-\alpha)^2}{2\mathcal I}\tau\Bigr)
\]
を示し、前問のスペクトルを再現せよ。\\
(3) 一般に、多重連結な配位空間 $Q$ 上で各ホモトピー類の寄与に重み $\chi$ をつけるとき、伝播関数の合成則から $\chi(n+m)=\chi(n)\chi(m)$、確率の保存から $\lvert\chi\rvert=1$ が要求されることを説明し、$\pi_1(S^1)=\Z$ の場合の可能な重みをすべて求めよ。$\theta$ 角との類似を述べよ。
\end{problem}

\begin{solution}
(1) 円周上の経路 $\theta(\tau')$ は $\R$ への持ち上げ $\tilde\theta(\tau')$ をもち、$\tilde\theta(0)=\theta'$ と固定すると終点は $\theta+2\pi n$ のいずれか（定理 \ref{thm:ch02-lift}）で、$n$ は経路の端点を止めたホモトピー類を区別する。円周上の経路積分は $\R$ 上の経路積分を、終点 $\theta+2\pi n$ について足し上げたものになる。運動項 $\int\frac{\mathcal I}{2}\dot\theta^2$ は $\R$ 上の自由粒子の作用で、その経路積分が熱核 $\sqrt{\mathcal I/2\pi\tau}\,e^{-\mathcal I(\theta+2\pi n-\theta')^2/2\tau}$ を与える。一方、項 $i\alpha\int\dot\theta\,d\tau'=i\alpha(\theta+2\pi n-\theta')$ は経路の端点だけで決まる（位相的な項）ので、各セクターで一定の位相 $e^{i\alpha(\theta-\theta'+2\pi n)}$ として外に出る（$e^{-S_E}$ の中で $-S_E\ni+i\alpha\Delta\tilde\theta$）。以上より与式。

\begin{note}
実時間の作用 $S=\int dt\,[\tfrac{\mathcal I}{2}\dot\theta^2+\alpha\dot\theta]$（$qA_\theta\dot\theta=\alpha\dot\theta$）で $t=-i\tau$ とおくと、$iS=\int d\tau\,[-\tfrac{\mathcal I}{2}(\partial_\tau\theta)^2+i\alpha\partial_\tau\theta]=-S_E$。位相的な項は1階微分なのでウィック回転しても虚数のまま残る。
\end{note}

(2) $g(y)=e^{i\alpha y}\sqrt{\mathcal I/2\pi\tau}\,e^{-\mathcal Iy^2/2\tau}$、$x=\theta-\theta'$ とおくと $K=\sum_ng(x+2\pi n)$。ガウス積分により
\[
\hat g(k)=\sqrt{\frac{\mathcal I}{2\pi\tau}}\int_\R e^{-\mathcal Iy^2/2\tau}e^{-i(k-\alpha)y}dy=\exp\Bigl(-\frac{(k-\alpha)^2\tau}{2\mathcal I}\Bigr).
\]
ポアソンの和公式より与式が従う。一方、スペクトル分解 $K=\sum_k\psi_k(\theta)\overline{\psi_k(\theta')}e^{-E_k\tau}$ と比べると、$\psi_k=e^{ik\theta}/\sqrt{2\pi}$、$E_k=(k-\alpha)^2/2\mathcal I=\frac{1}{2mR^2}(k-\frac{q\Phi}{2\pi})^2$ で、前問の結果と一致する。巻き付き数の和（経路積分）と角運動量の和（ハミルトニアン）がポアソン双対で結ばれている。

(3) 時間 $\tau_1+\tau_2$ の伝播関数は、中間点 $\theta''$ で積分した $\int d\theta''K(\theta,\theta'';\tau_2)K(\theta'',\theta';\tau_1)$ に等しくなければならない。巻き付き数 $m$ の経路と $n$ の経路をつなぐと巻き付き数 $m+n$ の経路になり、逆に巻き付き数 $m+n$ の経路はこのように分解される。したがって各セクターに重み $\chi(n)$ をつけた $K=\sum_n\chi(n)K_n$ が合成則をみたすには $\chi(m)\chi(n)=\chi(m+n)$ が必要である。また $\chi$ の絶対値が1でないと、$K$ が与える時間発展は（実時間で）ユニタリにならない（例えば $\lvert\chi(1)\rvert>1$ なら大きな巻き付き数の寄与が指数的に増える）。よって $\chi\colon\Z\to U(1)$ は準同型で、$\chi(n)=\chi(1)^n=e^{i\vartheta n}$、$\vartheta\in[0,2\pi)$ が可能な重みのすべてである。(1) と比べると $\vartheta=2\pi\alpha=q\Phi$：重みの選択は、円周を貫く磁束（AB位相）を選ぶことと同じである。

QCD の $\theta$ 真空（第8章）でも、ユークリッド配位はインスタントン数 $k\in\Z$ で分類され、経路積分に $e^{ik\theta}$ の重みをつける自由度がある。真空のエネルギーが $\theta$ について周期 $2\pi$ であることは、(2) の $E_0(\alpha)$ の周期性と同じ構造である。
\end{solution}

\begin{problem}[難易度 C]\label{ch02-prob-wilson}
多様体 $M$ 上の $U(1)$ 接続 $A$ が平坦（$F=dA=0$）であるとし、電荷 $q$ の粒子に対するウィルソン・ループを $W(\gamma)=\exp\bigl(iq\oint_\gamma A\bigr)$ とする。\\
(1) $x_0$ を基点とするループ $\gamma_0\simeq\gamma_1$（rel 端点）なら $W(\gamma_0)=W(\gamma_1)$ であることを、ホモトピー $H\colon I\times I\to M$（滑らかとしてよい）とストークスの定理を用いて示し、$W$ が準同型 $\pi_1(M,x_0)\to U(1)$ を定めることを示せ。また $W$ がゲージ変換 $A\mapsto A+d\lambda$ で不変であることを確かめよ。\\
(2) $M=T^2=\R^2/(L_1\Z\oplus L_2\Z)$ 上の定数接続 $A=\frac{a_1}{L_1}dx+\frac{a_2}{L_2}dy$（$q=1$）について、$W$ を $\pi_1(T^2)=\Z^2$ の生成元上で求めよ。ゲージ変換 $g=\exp\bigl(2\pi i(\tfrac{m_1x}{L_1}+\tfrac{m_2y}{L_2})\bigr)$（$m_i\in\Z$）による $(a_1,a_2)$ の変化を求め、平坦接続のゲージ同値類の空間（モジュライ空間）が $\mathrm{Hom}(\Z^2,U(1))\cong T^2$ であることを説明せよ。\\
(3) $M$ がレンズ空間 $L(N,p)$ のとき、平坦 $U(1)$ 接続のホロノミーとしてありうるものを求めよ。$F=0$ でありながら観測可能な効果をもつことを説明せよ。\\
(4) 非可換群 $G$ の平坦接続では、ホロノミーが準同型 $\pi_1(M,x_0)\to G$ を定め、基点の変更やゲージ変換で全体の共役だけ変わる。このことを (1) の類推から説明せよ（詳しくは第8章）。
\end{problem}

\begin{solution}
(1) $H(s,0)=\gamma_0(s)$, $H(s,1)=\gamma_1(s)$, $H(0,t)=H(1,t)=x_0$ とする。正方形 $I^2$ 上でストークスの定理を $H^*A$ に適用すると
\[
\int_{I^2}H^*F=\int_{I^2}d(H^*A)=\oint_{\partial I^2}H^*A=\int_{\gamma_0}A+\int_{\text{右辺}}A-\int_{\gamma_1}A-\int_{\text{左辺}}A .
\]
左右の辺は定値の道 $x_0$ なので寄与はゼロ、$F=0$ より左辺もゼロ。よって $\oint_{\gamma_0}A=\oint_{\gamma_1}A$ で、$W$ は $\pi_1(M,x_0)$ 上で well-defined。積については $\oint_{\gamma\cdot\delta}A=\oint_\gamma A+\oint_\delta A$ なので $W(\gamma\cdot\delta)=W(\gamma)W(\delta)$、準同型である。ゲージ変換では $\oint_\gamma d\lambda=0$（ループ上の完全形式の積分；$\lambda$ が一価な関数のとき）なので不変。$\lambda$ が $U(1)$ 値 $e^{i\lambda}$ としてのみ一価な場合（大きなゲージ変換）は $\oint d\lambda\in2\pi\Z$ で、$q\in\Z$ なら $W$ はやはり不変。

(2) 生成元は $\gamma_1(s)=(L_1s,0)$, $\gamma_2(s)=(0,L_2s)$ で、$W(\gamma_j)=e^{ia_j}$。$g$ は $T^2$ 上で一価（$x\to x+L_1$ で $e^{2\pi im_1}=1$）で、$A\mapsto A+i\,g\,dg^{-1}=A+2\pi(\tfrac{m_1}{L_1}dx+\tfrac{m_2}{L_2}dy)$。すなわち $a_j\mapsto a_j+2\pi m_j$。小さなゲージ変換（一価な $\lambda$）では定数接続の $a_j$ は変わらない。よって定数接続のゲージ同値類は $(a_1,a_2)\in(\R/2\pi\Z)^2$。（任意の平坦接続は $A=A_{\text{定数}}+d\lambda$ と書けることも示せる：調和形式の議論、第5章。）一方、ホロノミー $(e^{ia_1},e^{ia_2})$ は準同型 $\Z^2\to U(1)$ を定め、これはすべての準同型を尽くす。モジュライ空間 $\cong\mathrm{Hom}(\Z^2,U(1))=U(1)^2\cong T^2$。トーラス・コンパクト化のウィルソン線（第10章）はこの自由度である。

(3) $\pi_1(L(N,p))\cong\Z_N$（問題 \ref{ch02-prob-lens}）なので、ホロノミーは準同型 $\Z_N\to U(1)$、すなわち生成元の像 $e^{2\pi ik/N}$（$k=0,\dots,N-1$）で決まる $N$ 通りである。$\Z_N$ は有限なので連続的な変形の余地がなく、$k\neq0$ の接続は $F=0$ であるにもかかわらず自明な接続とゲージ同値でない。電荷 $q$ の粒子が生成元のループを一周すると $e^{2\pi ikq/N}$ の位相を得て、干渉やエネルギー準位（前問と同様）に効果が現れる。オービフォールド・コンパクト化では、こうした離散的なウィルソン線がゲージ対称性を破る手段になる（第10章）。

(4) 非可換の場合、ホロノミーは経路順序積 $P\exp(i\int_\gamma A)$ で、ストークスの定理の代わりに「平坦なら平行移動は端点を止めた変形で変わらない」ことが（非可換ストークスの定理から）示せる。経路をつなぐとホロノミーは積になる（向きの規約により準同型または逆準同型）。基点を $x_0$ から $x_1$ に動かすと、道 $h$ に沿った平行移動 $U_h$ による共役 $U_h(\cdot)U_h^{-1}$ だけ変わり、ゲージ変換 $g$ では $x_0$ での値 $g(x_0)$ による共役で変わる。したがってゲージ不変な情報は $\mathrm{Hom}(\pi_1(M),G)/G$（共役による商）であり、平坦 $G$ 接続のモジュライ空間になる。$\tr$ をとったウィルソン・ループ $\tr P\exp(i\oint A)$ は共役で不変なので、ゲージ不変な観測量である（第8章）。
\end{solution}

\subsection{量子統計と組みひも群}

\begin{problem}[難易度 C]\label{ch02-prob-stat}
$d$ 次元空間の2つの同種粒子の配位空間を $C_2(\R^d)=\{(x_1,x_2)\in\R^d\times\R^d:x_1\ne x_2\}/S_2$ とする（$S_2$ は入れ替え）。\\
(1) 重心 $X=\tfrac12(x_1+x_2)$ と相対座標 $\bm r=x_1-x_2$ を用いて、$C_2(\R^d)\cong\R^d\times P_d$、$P_d=(\R^d\setminus\{0\})/(\bm r\sim-\bm r)$ を示し、$P_d\simeq\RP^{d-1}$ を示せ。\\
(2) $d=1,2,3$ について $\pi_1(C_2(\R^d))$ を求めよ（$d\ge3$ でも同様）。\\
(3) 粒子の交換を表すループ $\varepsilon$（相対座標 $\bm r$ を、ある平面内で角度 $\pi$ だけ回転させる道の像）が、$d\ge3$ では $\varepsilon^2\simeq$ 定値、$d=2$ では $\pi_1$ の生成元であることを示せ。これより、交換の位相 $\chi(\varepsilon)$ は $d\ge3$ で $\pm1$（ボソン・フェルミオン）、$d=2$ では任意の $e^{i\vartheta}$（\term{エニオン}{anyon}）であることを結論せよ。\\
(4) $d=2$ で、相対運動の波動関数を極座標 $\bm r=(r\cos\varphi,r\sin\varphi)$ の普遍被覆上の関数とし、$\psi(r,\varphi+\pi)=e^{i\vartheta}\psi(r,\varphi)$ を課す。相対角運動量 $\ell$（$\psi\propto e^{i\ell\varphi}$）の許される値を求め、ボソン・フェルミオンの場合と比べよ。
\end{problem}

\begin{solution}
(1) $(x_1,x_2)\mapsto(X,\bm r)$ は $\R^d\times\R^d$ の線形同型で、対角集合 $x_1=x_2$ は $\bm r=0$ に、入れ替えは $(X,\bm r)\mapsto(X,-\bm r)$ に対応する。よって $C_2(\R^d)\cong\R^d\times(\R^d\setminus\{0\})/\pm=\R^d\times P_d$（商の連続性は直積の商として確かめられる）。$\R^d$ は可縮なので $C_2\simeq P_d$。$\R^d\setminus\{0\}$ は $S^{d-1}$ に $H(\bm r,t)=(1-t)\bm r+t\bm r/\lvert\bm r\rvert$ で強変位レトラクトし、$H(-\bm r,t)=-H(\bm r,t)$ なので商に降りて $P_d\simeq S^{d-1}/\pm=\RP^{d-1}$。

(2) $d=1$：$\RP^0$ は1点で $\pi_1=1$。（$P_1=(\R\setminus\{0\})/\pm\cong(0,\infty)$。1次元では粒子は互いを通り抜けられず、交換というループが存在しない。）
$d=2$：$\RP^1\cong S^1$（$S^1\to\RP^1$ は $z\mapsto z^2$ と同一視できる）なので $\pi_1=\Z$。
$d\ge3$：$S^{d-1}$（$d-1\ge2$）は単連結で、$\pm1$ は自由に作用するので $\pi_1(\RP^{d-1})\cong\Z_2$（定理 \ref{thm:ch02-quotient}）。

(3) $\varepsilon$ の $S^{d-1}$ への持ち上げは $\hat{\bm r}$ から $-\hat{\bm r}$ への半円で、閉じないので $\varepsilon\ne1$（$d\ge2$）。
$d\ge3$：$\pi_1=\Z_2$ の唯一の非自明元が $\varepsilon$ で、$\varepsilon^2\simeq$ 定値。（$\varepsilon^2$ の持ち上げは大円全体で、$S^{d-1}$（$d\ge3$）上で縮められる。物理的には、相対座標の $2\pi$ 回転は3次元では第3の方向へ持ち上げてほどける。）
$d=2$：$\RP^1\cong S^1$ で、$S^1\to\RP^1$ の2重被覆に関して $\varepsilon$ の持ち上げは半周なので、$\RP^1$ を「ちょうど1周」するループであり、生成元である。$\varepsilon^2$ は相対座標が $2\pi$ 回転するループで、生成元の2倍であり縮まない。

シュルマン・レイドロー・ドウィットの定理より、量子化は1次元ユニタリ表現 $\chi\colon\pi_1\to U(1)$ で分類される。$d\ge3$ では $\chi(\varepsilon)^2=\chi(\varepsilon^2)=1$ より $\chi(\varepsilon)=\pm1$。$+1$ がボソン、$-1$ がフェルミオンである。$d=2$ では $\pi_1=\Z$ が自由に生成されるので $\chi(\varepsilon)=e^{i\vartheta}$ は任意で、交換の位相が任意のエニオンが許される。

(4) 交換 $\varphi\to\varphi+\pi$ に位相 $e^{i\vartheta}$ を課すと、$e^{i\ell(\varphi+\pi)}=e^{i\vartheta}e^{i\ell\varphi}$ より $e^{i\pi\ell}=e^{i\vartheta}$、すなわち
\[
\ell\in\frac{\vartheta}{\pi}+2\Z .
\]
ボソン（$\vartheta=0$）では $\ell$ は偶数、フェルミオン（$\vartheta=\pi$）では奇数で、相対波動関数の対称・反対称性と一致する。エニオンでは $\ell$ が整数でなくなり、角運動量の値がずれる。これは、相対座標の原点に磁束 $\Phi$ を置いた AB 系（問題 \ref{ch02-prob-ring}）で角運動量が $n\to n-q\Phi/2\pi$ とずれるのと同じ構造で、エニオンは「電荷と磁束の複合粒子」とみなせる（チャーン・サイモンズ理論、第8章）。
\end{solution}

\begin{problem}[難易度 C]\label{ch02-prob-braid}
$n\ge2$ とし、$\pi_1(C_n(\R^2))\cong B_n=\langle\sigma_1,\dots,\sigma_{n-1}\mid\sigma_i\sigma_{i+1}\sigma_i=\sigma_{i+1}\sigma_i\sigma_{i+1},\ \sigma_i\sigma_j=\sigma_j\sigma_i\ (\lvert i-j\rvert\ge2)\rangle$、および $d\ge3$ で $\pi_1(C_n(\R^d))\cong S_n$（$\sigma_i\mapsto$ 互換 $(i\ i{+}1)$ による）を認める。\\
(1) $B_n$ の1次元ユニタリ表現 $\chi$ はすべて $\chi(\sigma_i)=e^{i\vartheta}$（$i$ によらない）の形であることを示せ。\\
(2) $\sigma_i\mapsto1$ が準同型 $B_n\to\Z$ を定めることを示し、$B_n$ の可換化が $\Z$ であることを示せ。\\
(3) $d\ge3$ では1次元ユニタリ表現は $\vartheta\in\{0,\pi\}$ に限ることを示せ。\\
(4) 球面 $S^2$ 上の $n$ 粒子では、さらに関係式 $\sigma_1\sigma_2\cdots\sigma_{n-1}\sigma_{n-1}\cdots\sigma_2\sigma_1=1$ が加わる（1つの粒子が他のすべての粒子のまわりを1周するループが、球面の裏側を通って縮められる）。このとき許される $\vartheta$ を求めよ。
\end{problem}

\begin{solution}
(1) $U(1)$ は可換なので、関係式 $\sigma_i\sigma_{i+1}\sigma_i=\sigma_{i+1}\sigma_i\sigma_{i+1}$ から $\chi(\sigma_i)^2\chi(\sigma_{i+1})=\chi(\sigma_i)\chi(\sigma_{i+1})^2$、すなわち $\chi(\sigma_i)=\chi(\sigma_{i+1})$。よってすべて同じ値 $e^{i\vartheta}$。逆に任意の $\vartheta$ で、$\sigma_i\mapsto e^{i\vartheta}$ はすべての関係式（両辺で各生成元の個数が等しい）を保つので表現になる。

(2) 各関係式の両辺は生成元を同じ個数（3個ずつ、2個ずつ）含むので、$\sigma_i\mapsto1$ は自由群からの準同型として関係子を $0$ に送り、準同型 $\ell\colon B_n\to\Z$（「交差の符号付き個数」）を誘導する。可換化 $B_n^{\mathrm{ab}}$ では (1) と同じ議論ですべての $\sigma_i$ が等しくなるので、$B_n^{\mathrm{ab}}$ は1つの元 $\sigma$ で生成される巡回群で、$\ell$ がそれを $\Z$ に全射に写すので $B_n^{\mathrm{ab}}\cong\Z$。特に $B_n$ の元 $\sigma_1^k$（$k\ne0$）は単位元でない。（$n=2$ では $B_2\cong\Z$。$\sigma_1^2$ は図の通り、一方の粒子が他方を1周する運動である。）

(3) $d\ge3$ では $\pi_1=S_n$ で、互換は $\sigma_i^2=1$ をみたす。1次元表現は $\chi(\sigma_i)^2=1$、かつ (1) と同じ関係式（$S_n$ でも組みひも関係式は成り立つ）からすべて等しいので、$\chi(\sigma_i)=e^{i\vartheta}$、$\vartheta\in\{0,\pi\}$：自明表現（ボソン）と符号表現（フェルミオン）。2次元では $\sigma_i^2\ne1$（(2) より $\ell(\sigma_i^2)=2\ne0$）なので、この制限がない。

(4) 1次元表現では関係式の左辺は $e^{i\vartheta\cdot2(n-1)}$ なので $e^{2i(n-1)\vartheta}=1$、すなわち
\[
\vartheta\in\frac{\pi}{n-1}\Z .
\]
平面上では任意だった統計の角度が、球面上では粒子数に依存して量子化される。$n=2$ では $\vartheta\in\pi\Z$ でボソン・フェルミオンに限る。このように多体系の統計は空間のトポロジーに敏感で、トーラス上ではさらに基底状態の縮退が生じる（トポロジカル秩序）。
\end{solution}

%================================================
\chapter{高次ホモトピー群と位相欠陥}\label{ch:03}
% 20問
%================================================

\section{物理からの動機}

\subsection*{欠陥を「囲む球面」}
第1章で見たように、対称性 $G$ が $H$ に自発的に破れると、真空の全体は等質空間 $\mathcal{M}=G/H$（真空多様体）になる。宇宙の温度が下がって相転移が起こるとき、因果的に結びついていない領域どうしは真空 $\mathcal M$ のどの点を選ぶかを独立に決めるので、それらの境目では場が真空からはずれた「傷」が残りうる（\term{キブル機構}{Kibble mechanism}）。この傷を\term{位相欠陥}{topological defect}という。

3次元空間の中の欠陥を、それを取り囲む小さな球面で囲んでみる（下図）。
\begin{itemize}
\item 面状の欠陥（\term{ドメイン・ウォール}{domain wall}）は、壁の両側の2点 $S^0$ で囲まれる。
\item 線状の欠陥（\term{ひも}{string}、渦糸）は、それを1周する円周 $S^1$ で囲まれる。
\item 点状の欠陥（\term{モノポール}{monopole}）は、それを包む球面 $S^2$ で囲まれる。
\end{itemize}
欠陥から離れたところでは場は真空 $\mathcal M$ の値をとるので、囲む球面上の場は写像 $S^{k-1}\to\mathcal{M}$ を定める（$k$ は欠陥の\term{余次元}{codimension}）。この写像が連続変形で定値写像に縮められないとき、欠陥はほどけない。こうして「$\mathcal M$ の中の $(k-1)$ 次元球面のホモトピー類」、すなわち\textbf{高次ホモトピー群} $\pi_{k-1}(\mathcal M)$ が欠陥を分類する。

\begin{center}
\begin{tikzpicture}[scale=0.95]
  % (a) wall
  \begin{scope}
    \fill[blue!12] (-1.1,-1.3) -- (0.5,-0.7) -- (0.5,1.5) -- (-1.1,0.9) -- cycle;
    \draw[blue!60!black,thick] (-1.1,-1.3) -- (0.5,-0.7) -- (0.5,1.5) -- (-1.1,0.9) -- cycle;
    \node[vtx,red!70!black] at (-1.25,0.25) {}; \node[red!70!black,left] at (-1.25,0.25) {\footnotesize $-v$};
    \node[vtx,red!70!black] at (0.9,0.0) {}; \node[red!70!black,right] at (0.9,0.0) {\footnotesize $+v$};
    \draw[red!70!black,densely dotted] (-1.25,0.25) -- (0.9,0.0);
    \node at (-0.3,-2.0) {\small (a) 壁：$S^0$，$\pi_0(\mathcal M)$};
  \end{scope}
  % (b) string
  \begin{scope}[xshift=4.3cm]
    \draw[ultra thick,blue!60!black] (0,-1.4) -- (0,-0.32);
    \draw[ultra thick,blue!60!black] (0,-0.12) -- (0,1.6);
    \draw[thick,red!70!black] (0,0) ellipse (1.0 and 0.35);
    \draw[ultra thick,blue!60!black] (0,0.36) -- (0,0.5);
    \draw[-{Stealth[length=2mm]},red!70!black,thick] (0.35,-0.33) -- (0.45,-0.31);
    \node[red!70!black] at (1.35,0.25) {\footnotesize $S^1$};
    \node at (0,-2.0) {\small (b) ひも：$S^1$，$\pi_1(\mathcal M)$};
  \end{scope}
  % (c) monopole
  \begin{scope}[xshift=8.4cm]
    \shade[ball color=red!8,opacity=0.7] (0,0) circle (1.1);
    \draw[thick,red!70!black] (0,0) circle (1.1);
    \draw[red!70!black,dashed] (1.1,0) arc (0:180:1.1 and 0.3);
    \draw[red!70!black] (-1.1,0) arc (180:360:1.1 and 0.3);
    \node[circle,fill=blue!60!black,inner sep=2pt] at (0,0) {};
    \node[red!70!black] at (1.35,0.8) {\footnotesize $S^2$};
    \node at (0,-2.0) {\small (c) モノポール：$S^2$，$\pi_2(\mathcal M)$};
  \end{scope}
\end{tikzpicture}
\zucap{3次元空間の位相欠陥（青）と、それを囲む球面（赤）。囲む球面上の場の値は真空多様体 $\mathcal M$ への写像を与え、そのホモトピー類が欠陥の位相的電荷である。}
\end{center}

\subsection*{渦とハリネズミ}
最も簡単な例を2つ挙げる。$U(1)\to1$ の破れ（$\mathcal M=S^1$）では、渦を1周すると位相 $\arg\phi$ が $2\pi n$ だけ回る。$n\in\Z=\pi_1(S^1)$ が渦の巻き数である（第2章）。

\begin{center}
\begin{tikzpicture}[scale=0.55]
  \foreach \s/\xs/\lab in {1/0/{$n=+1$（渦）},-1/9/{$n=-1$（反渦）}}{
    \begin{scope}[xshift=\xs cm]
      \foreach \x in {-3,-2,...,3}{
        \foreach \y in {-3,-2,...,3}{
          \pgfmathsetmacro\rr{sqrt(\x*\x+\y*\y)}
          \ifdim \rr pt>0.1pt
            \pgfmathsetmacro\ang{\s*atan2(\y,\x)}
            \draw[-{Stealth[length=1.4mm]},thick,blue!60!black] ($(\x,\y)-0.35*({cos(\ang)},{sin(\ang)})$) -- ($(\x,\y)+0.35*({cos(\ang)},{sin(\ang)})$);
          \fi
        }
      }
      \node[circle,fill=red!70!black,inner sep=2.2pt] at (0,0) {};
      \draw[red!70!black,dashed] (0,0) circle (2.5);
      \node at (0,-4.1) {\small \lab};
    \end{scope}
  }
\end{tikzpicture}
\zucap{2次元の $U(1)$ 渦。矢印は各点での $\phi/\lvert\phi\rvert\in S^1$ の向き（位相）。破線の円周を1周すると、矢印は左では正の向きに1回、右では負の向きに1回まわる。中心（赤）では $\phi=0$ でなければならない。}
\end{center}

$SU(2)$ の随伴表現のヒッグス場が $U(1)$ に破れる場合（ジョージ・グラショウ模型、$\mathcal M=S^2$）には、遠方で $\phi^a\to v\,x^a/r$ となる\term{ハリネズミ}{hedgehog}配位が可能である。遠方の球面 $S^2_\infty$ から真空の球面 $S^2$ への写像は恒等写像であり、写像度1をもつ。これが\term{トフーフト・ポリャコフ・モノポール}{'t Hooft--Polyakov monopole}であり、その磁荷は写像度で決まる。

\begin{center}
\begin{tikzpicture}[scale=1.0]
  \shade[ball color=blue!8,opacity=0.6] (0,0) circle (1.6);
  \draw[thick] (0,0) circle (1.6);
  \draw[dashed,gray] (1.6,0) arc (0:180:1.6 and 0.4);
  \draw[gray] (-1.6,0) arc (180:360:1.6 and 0.4);
  \foreach \a in {0,30,...,330}{
    \draw[-{Stealth[length=2mm]},thick,red!70!black] (\a:1.6) -- (\a:2.25);
    \draw[-{Stealth[length=1.5mm]},red!50!black] (\a+15:0.55) -- (\a+15:0.95);
  }
  \node[circle,fill=black,inner sep=1.5pt] at (0,0) {};
  \node at (0,-2.65) {\small $\phi^a(\bm x)\to v\,\hat x^a$};
  \begin{scope}[xshift=6cm]
    \shade[ball color=red!8,opacity=0.6] (0,0) circle (1.2);
    \draw[thick] (0,0) circle (1.2);
    \draw[dashed,gray] (1.2,0) arc (0:180:1.2 and 0.3);
    \draw[gray] (-1.2,0) arc (180:360:1.2 and 0.3);
    \node at (0,-2.0) {\small 真空多様体 $\{\phi^a\phi^a=v^2\}\cong S^2$};
  \end{scope}
  \draw[-{Stealth[length=2.5mm]},thick] (2.6,0.4) to[bend left=20] node[above] {\footnotesize $\hat{\bm x}\mapsto\bm\phi/v$（写像度 1）} (4.6,0.4);
\end{tikzpicture}
\zucap{ハリネズミ配位。空間の各方向 $\hat{\bm x}$ に、内部空間の同じ方向を向いたヒッグス場を対応させる。中心では $\phi=0$。}
\end{center}

\subsection*{欠陥以外の位相的な対象}
\begin{itemize}
\item \textbf{テクスチャー・スカーミオン}：空間全体で場が真空にあり、遠方で一定値に近づく配位は $\R^3\cup\{\infty\}\cong S^3\to\mathcal M$ を定め、$\pi_3(\mathcal M)$ で分類される。カイラル・ラグランジアン（$\mathcal M=SU(2)\cong S^3$）では、これが\term{スカーミオン}{Skyrmion}であり、その巻き付き数はバリオン数と解釈される。
\item \textbf{インスタントン}：4次元ユークリッドのゲージ理論で作用が有限の配位は、遠方で $A\to i\,g\,dg^{-1}$（純ゲージ）となり、$g\colon S^3_\infty\to G$ が現れる。そのホモトピー類 $\pi_3(G)\cong\Z$ がインスタントン数を決める（詳細は第8章）。量子論では、真空が $\pi_3(G)$ の元で番号づけられ、$\theta$ 真空が生じる。
\item \textbf{2次元のソリトン・ホップ粒子}：2次元の $O(3)$ 模型の場 $\bm n\colon\R^2\cup\{\infty\}\to S^2$ は $\pi_2(S^2)=\Z$ で分類され（\term{ランプ}{lump}、磁性体のスカーミオン）、3次元では $\bm n\colon S^3\to S^2$ が $\pi_3(S^2)=\Z$（ホップ写像）で分類される（\term{ホップ粒子}{hopfion}）。
\item \textbf{大域的アノマリー}：$\pi_4(SU(2))=\Z_2$ は、ウィッテンの $SU(2)$ アノマリー（第9章）の源である。
\end{itemize}
本章では、これらを統一的に扱う道具として高次ホモトピー群 $\pi_n$ を導入し、ファイブレーションの長完全列を使って $\pi_n(G/H)$ を計算する方法を学ぶ。

\section{基本事項のまとめ}

\subsection{高次ホモトピー群}

\begin{mydef}{ホモトピー群}{ch03-pin}
$n\ge1$ とし、$I^n=[0,1]^n$、$\partial I^n$ をその境界とする。基点 $x_0\in X$ に対し、連続写像 $f\colon(I^n,\partial I^n)\to(X,x_0)$（境界を $x_0$ に写す写像）を、境界を $x_0$ に固定したままのホモトピーで同一視した類の集合を $\pi_n(X,x_0)$ と書く。積を
\[
(f+g)(s_1,\dots,s_n)=\begin{cases}f(2s_1,s_2,\dots,s_n)&(s_1\le\tfrac12)\\ g(2s_1-1,s_2,\dots,s_n)&(s_1\ge\tfrac12)\end{cases}
\]
で定めると $\pi_n(X,x_0)$ は群になり、これを $n$ 次の\term{ホモトピー群}{homotopy group}という。$n=1$ は基本群（第2章）である。また $\pi_0(X)$ は弧状連結成分の集合を表す（一般には群ではない）。
\end{mydef}
$I^n/\partial I^n\cong S^n$ なので、$\pi_n(X,x_0)$ は基点付き写像 $(S^n,*)\to(X,x_0)$ の基点付きホモトピー類の集合 $[S^n,X]_*$ と同じである。物理で現れる「遠方で一定値 $x_0$ に近づく $\R^n$ 上の場」は、まさにこの元である（第1章の立体射影）。

\begin{mythm}{$n\ge2$ で $\pi_n$ は可換群}{ch03-abelian}
$n\ge2$ のとき $\pi_n(X,x_0)$ は可換群である。
\end{mythm}
証明の要点は、$n\ge2$ では「別の座標方向」に $f$ と $g$ を並べ替える余地があることである（\term{エックマン・ヒルトンの議論}{Eckmann--Hilton argument}、問題 \ref{ch03-prob-abelian}）。下図のように、境界（基点に写る部分）を広げて $f$ と $g$ をそれぞれ小さな箱に縮め、箱どうしをすれ違わせれば $f+g\simeq g+f$ が得られる。$n=1$ では箱が1次元の線分上に並ぶのですれ違えない。

\begin{center}
\begin{tikzpicture}[scale=0.9]
  \foreach \i/\lab in {0/1,1/2,2/3,3/4}{
    \draw[thick] (3.3*\i,0) rectangle +(2.4,2.4);
  }
  % 1
  \fill[blue!15] (0,0) rectangle (1.2,2.4); \fill[red!15] (1.2,0) rectangle (2.4,2.4);
  \draw (1.2,0) -- (1.2,2.4); \node at (0.6,1.2) {$f$}; \node at (1.8,1.2) {$g$};
  % 2
  \fill[blue!15] (3.3+0.2,0.25) rectangle (3.3+1.0,1.05); \node at (3.3+0.6,0.65) {$f$};
  \fill[red!15] (3.3+1.4,1.35) rectangle (3.3+2.2,2.15); \node at (3.3+1.8,1.75) {$g$};
  % 3
  \fill[blue!15] (6.6+1.4,0.25) rectangle (6.6+2.2,1.05); \node at (6.6+1.8,0.65) {$f$};
  \fill[red!15] (6.6+0.2,1.35) rectangle (6.6+1.0,2.15); \node at (6.6+0.6,1.75) {$g$};
  \draw[-{Stealth[length=1.8mm]},thick,blue!60!black] (6.6+0.6,0.65) -- (6.6+1.5,0.65);
  \draw[-{Stealth[length=1.8mm]},thick,red!60!black] (6.6+1.8,1.75) -- (6.6+0.9,1.75);
  % 4
  \fill[red!15] (9.9,0) rectangle (9.9+1.2,2.4); \fill[blue!15] (9.9+1.2,0) rectangle (9.9+2.4,2.4);
  \draw (9.9+1.2,0) -- (9.9+1.2,2.4); \node at (9.9+0.6,1.2) {$g$}; \node at (9.9+1.8,1.2) {$f$};
  \foreach \i in {0,1,2}{ \node at (3.3*\i+2.85,1.2) {$\simeq$}; }
\end{tikzpicture}
\zucap{$\pi_2$ の可換性（エックマン・ヒルトン）。白い部分はすべて基点 $x_0$ に写る。$f$ と $g$ の箱を縦方向にずらしてから横方向にすれ違わせると、$f+g\simeq g+f$。}
\end{center}

\begin{myprop}{ホモトピー群の基本性質}{ch03-pinprop}
\begin{enumerate}
\item \textbf{関手性}：$\phi\colon X\to Y$（$\phi(x_0)=y_0$）は準同型 $\phi_*\colon\pi_n(X,x_0)\to\pi_n(Y,y_0)$ を誘導し、$(\psi\circ\phi)_*=\psi_*\phi_*$。ホモトピー同値は同型を誘導する。特に可縮な空間では $\pi_n=0$。
\item \textbf{積}：$\pi_n(X\times Y)\cong\pi_n(X)\times\pi_n(Y)$。
\item \textbf{被覆}：$p\colon\tilde X\to X$ が被覆写像なら、$n\ge2$ で $p_*\colon\pi_n(\tilde X)\to\pi_n(X)$ は同型。特に $\pi_n(S^1)=\pi_n(\R)=0$（$n\ge2$）、$\pi_n(\RP^m)\cong\pi_n(S^m)$（$n\ge2$）、$\pi_n(SO(3))\cong\pi_n(SU(2))$（$n\ge2$）。
\item \textbf{基点の取り替えと $\pi_1$ の作用}：$X$ が弧状連結なら $\pi_n(X,x_0)$ の同型類は基点によらない。ただし同型は基点を結ぶ道のホモトピー類に依存し、特に閉じた道により $\pi_1(X,x_0)$ が $\pi_n(X,x_0)$ に作用する。基点を固定しない写像 $S^n\to X$ のホモトピー類 $[S^n,X]$ は $\pi_n(X)$ をこの作用で割った集合と一致する。$X$ が単連結なら $[S^n,X]=\pi_n(X)$。
\end{enumerate}
\end{myprop}
最後の性質は物理的に重要である。基点付きの写像は「遠方で決まった値に近づく配位」に、基点なしの写像は「欠陥を囲む球面上の場」（どの値を基準にするかは決まっていない）に対応する。$\pi_1(\mathcal M)\neq0$ の場合、点欠陥の電荷は $\pi_2(\mathcal M)$ の元そのものではなく、ひもの周りを回すことによる作用の軌道で決まる（問題 \ref{ch03-prob-nematic}）。

\subsection{球面のホモトピー群と写像度}

\begin{mythm}{球面のホモトピー群（低次）}{ch03-spheres}
\begin{enumerate}
\item $k<n$ ならば $\pi_k(S^n)=0$。
\item $\pi_n(S^n)\cong\Z$ であり、同型は写像度で与えられる（ホップの定理）。
\item $\pi_3(S^2)\cong\Z$ であり、ホップ写像 $h\colon S^3\to S^2$ が生成元である。
\end{enumerate}
\end{mythm}
(1) の直観的な理由：$k<n$ なら $f\colon S^k\to S^n$ は（滑らかに近似すれば）全射になりえず、1点 $p$ を取りこぼす。$S^n\setminus\{p\}\cong\R^n$ は可縮なので $f$ は定値写像に縮む。(3) は本章で長完全列から導く。より高い次数については次の表を引用する（計算には本書の範囲を超える道具を要する）。
\begin{center}
\renewcommand{\arraystretch}{1.15}
\begin{tabular}{c|ccccc}
 & $\pi_1$ & $\pi_2$ & $\pi_3$ & $\pi_4$ & $\pi_5$\\ \hline
$S^1$ & $\Z$ & $0$ & $0$ & $0$ & $0$\\
$S^2$ & $0$ & $\Z$ & $\Z$ & $\Z_2$ & $\Z_2$\\
$S^3$ & $0$ & $0$ & $\Z$ & $\Z_2$ & $\Z_2$\\
$S^4$ & $0$ & $0$ & $0$ & $\Z$ & $\Z_2$
\end{tabular}
\end{center}

\begin{mydef}{写像度}{ch03-degree}
$M,N$ を向きのついた連結な閉 $n$ 次元多様体、$f\colon M\to N$ を滑らかな写像とする。\term{正則値}{regular value} $y\in N$（$f^{-1}(y)$ のすべての点で微分 $df$ が同型）に対し
\[
\deg f=\sum_{x\in f^{-1}(y)}\operatorname{sign}\det(df_x)
\]
を\term{写像度}{degree}という。これは $y$ の選び方によらない整数で、ホモトピー不変である。$\int_N\Omega\neq0$ をみたす任意の $n$ 形式 $\Omega$ に対し
\[
\deg f=\frac{\int_Mf^*\Omega}{\int_N\Omega}.
\]
\end{mydef}

\begin{myex}{$S^2$ への写像の写像度}{ch03-deg2}
$S^2\subset\R^3$ に外向き法線で向きを入れると、面積形式は
\[
\omega=\tfrac12\epsilon_{abc}\,n^a\,dn^b\wedge dn^c=\sin\theta\,d\theta\wedge d\varphi,\qquad\int_{S^2}\omega=4\pi .
\]
したがって $\bm n\colon\R^2\cup\{\infty\}\to S^2$（$\R^2$ の向きは $dx^1\wedge dx^2$）の写像度は
\begin{equation}
Q=\frac{1}{8\pi}\int\epsilon_{abc}\,n^a\,dn^b\wedge dn^c=\frac1{4\pi}\int d^2x\;\bm n\cdot(\partial_1\bm n\times\partial_2\bm n)\label{eq:ch03-Q}
\end{equation}
である（問題 \ref{ch03-prob-degint}）。被積分関数は\term{位相的電荷密度}{topological charge density}と呼ばれる。
\end{myex}

\subsection{ファイブレーションの長完全列}
高次ホモトピー群の計算の主な道具は次の定理である。ファイバー束の一般論は第7章で扱い、ここでは道具として認める。

\begin{mythm}{ホモトピー長完全列}{ch03-les}
$p\colon E\to B$ をファイバー束（より一般にファイブレーション）、$B$ を弧状連結、$F=p^{-1}(b_0)$ をファイバー、$e_0\in F$ とする。このとき完全列
\[
\begin{aligned}\cdots\to\pi_n(F)\xrightarrow{i_*}\pi_n(E)\xrightarrow{p_*}\pi_n(B)&\xrightarrow{\partial}\pi_{n-1}(F)\xrightarrow{i_*}\pi_{n-1}(E)\\&\to\cdots\to\pi_0(F)\to\pi_0(E)\to0\end{aligned}
\]
がある。ここで\term{完全}{exact}とは、各所で「前の写像の像＝次の写像の核」となることをいう。
\end{mythm}
\begin{note}
境界写像 $\partial$ の作り方：$[f]\in\pi_n(B)$ を $f\colon(I^n,\partial I^n)\to(B,b_0)$ で表す。$I^n=I^{n-1}\times I$ と見て、$I^{n-1}\times\{0\}\cup\partial I^{n-1}\times I$ を $e_0$ に写すように $f$ を $\tilde f\colon I^n\to E$ に持ち上げる（ホモトピー持ち上げ性質）。$p\circ\tilde f=f$ なので上面 $I^{n-1}\times\{1\}$ は $F$ に写り、これが $\pi_{n-1}(F)$ の元 $\partial[f]$ を与える。直観的には「$B$ の球面を $E$ に持ち上げようとして閉じきらない分が、ファイバーの中の1つ低い次元の球面として残る」。
\end{note}

$G$ をリー群、$H$ を閉部分群とすると、$H\to G\to G/H$ はファイバー束である。コンパクト・リー群について $\pi_2(G)=0$ が成り立つ（下の表）ので、次を得る。

\begin{mythm}{等質空間のホモトピー群}{ch03-GH}
$G$ を連結なコンパクト・リー群、$H$ を閉部分群とすると、完全列
\[
0=\pi_2(G)\to\pi_2(G/H)\xrightarrow{\partial}\pi_1(H)\xrightarrow{i_*}\pi_1(G)\to\pi_1(G/H)\xrightarrow{\partial}\pi_0(H)\to0
\]
がある。したがって
\[
\pi_2(G/H)\cong\ker\bigl(\pi_1(H)\to\pi_1(G)\bigr).
\]
特に $G$ が単連結なら $\pi_2(G/H)\cong\pi_1(H)$、$\pi_1(G/H)\cong\pi_0(H)$ である。また $\pi_2(H)=0$ より $\pi_3(G/H)\cong\coker\bigl(\pi_3(H)\to\pi_3(G)\bigr)$。
\end{mythm}
$G$ が単連結でないときは、普遍被覆 $\tilde G$ と $H$ の逆像 $\tilde H$ をとれば $G/H=\tilde G/\tilde H$ なので、単連結の場合に帰着できる。物理的には、$\pi_2(G/H)\cong\pi_1(H)$ は「モノポールの磁荷は、破れずに残った群 $H$ のループで測られる」ことを意味する（$H$ の $U(1)$ 因子が磁荷の量子化を与える）。

\begin{myex}{ホップ・ファイブレーションと複素射影空間}{ch03-hopf}
第1章のホップ写像 $h(z)^a=z^\dagger\sigma^az$ は $S^1\to S^3\xrightarrow{h}S^2$ というファイバー束であり（ファイバーは状態の全体位相 $\{e^{i\alpha}z\}$）、長完全列は
\begin{center}
\begin{tikzcd}[column sep=small]
\pi_3(S^1)=0\arrow[r]&\pi_3(S^3)\arrow[r,"h_*"]&\pi_3(S^2)\arrow[r]&\pi_2(S^1)=0,\\
\pi_2(S^3)=0\arrow[r]&\pi_2(S^2)\arrow[r,"\partial"]&\pi_1(S^1)\arrow[r]&\pi_1(S^3)=0
\end{tikzcd}
\end{center}
となる。よって $\pi_3(S^2)\cong\pi_3(S^3)\cong\Z$（生成元は $h$）、$\pi_2(S^2)\cong\pi_1(S^1)\cong\Z$。同様に $S^1\to S^{2n+1}\to\CP^n$ から、$\pi_2(\CP^n)\cong\Z$、$\pi_1(\CP^n)=0$、$k\ge3$ で $\pi_k(\CP^n)\cong\pi_k(S^{2n+1})$。
\end{myex}

\subsection{リー群のホモトピー群}
\begin{myprop}{コンパクト・リー群のホモトピー群}{ch03-liegroups}
\begin{enumerate}
\item 任意のコンパクト・リー群 $G$ について $\pi_2(G)=0$（ボットの定理。本書では証明しない）。
\item コンパクトな単純リー群 $G$ について $\pi_3(G)\cong\Z$。$SU(2)\subset G$ を（指数1の）適当な埋め込みで選ぶと、包含が $\pi_3$ の同型を誘導する。
\item 主なリー群について次の表が成り立つ。
\end{enumerate}
\begin{center}
\renewcommand{\arraystretch}{1.15}
\begin{tabular}{l|cccc}
$G$ & $\pi_1$ & $\pi_2$ & $\pi_3$ & $\pi_4$\\ \hline
$U(1)$ & $\Z$ & $0$ & $0$ & $0$\\
$SU(2)\cong Sp(1)\cong S^3$ & $0$ & $0$ & $\Z$ & $\Z_2$\\
$SU(N)$ ($N\ge3$) & $0$ & $0$ & $\Z$ & $0$\\
$U(N)$ ($N\ge3$) & $\Z$ & $0$ & $\Z$ & $0$\\
$Sp(N)$ ($N\ge2$) & $0$ & $0$ & $\Z$ & $\Z_2$\\
$SO(3)$ & $\Z_2$ & $0$ & $\Z$ & $\Z_2$\\
$SO(4)$ & $\Z_2$ & $0$ & $\Z\oplus\Z$ & $\Z_2\oplus\Z_2$\\
$SO(5)$ & $\Z_2$ & $0$ & $\Z$ & $\Z_2$\\
$SO(N)$ ($N\ge6$) & $\Z_2$ & $0$ & $\Z$ & $0$
\end{tabular}
\end{center}
\end{myprop}
$\pi_3$ の生成元は、写像 $g\colon S^3\to G$ の\term{巻き付き数}{winding number}
\begin{equation}
w(g)=\frac{1}{24\pi^2}\int_{S^3}\tr\bigl(g^{-1}dg\bigr)^3\label{eq:ch03-w}
\end{equation}
で検出される（$SU(N)$ の基本表現でトレースをとる。$S^3$ の向きは恒等写像 $S^3\cong SU(2)$ で $w=1$ となるように選ぶ。記法と規約を参照）。$w$ は整数で、各点ごとの積について加法的 $w(gh)=w(g)+w(h)$ である（問題 \ref{ch03-prob-wadd}）。$\pi_4(SU(2))=\Z_2$ は第9章でウィッテンのアノマリーに用いる。$\pi_k$ の周期性（ボット周期性）は第11章で扱う。

\subsection{位相欠陥の分類}

\begin{mythm}{欠陥の分類}{ch03-defects}
$d$ 次元空間で、真空多様体 $\mathcal M$ をもつ理論を考える。
\begin{enumerate}
\item 余次元 $k$ の欠陥（$d=3$ では $k=1,2,3$ がそれぞれ壁・ひも・点）は、それを囲む $S^{k-1}$ 上の場のホモトピー類、すなわち $[S^{k-1},\mathcal M]$（$\mathcal M$ が単連結なら $\pi_{k-1}(\mathcal M)$）で分類される。$\pi_{k-1}(\mathcal M)\neq0$ が位相的に安定な欠陥の存在の必要条件である。
\item 特異点をもたず遠方で一定値に近づく配位（テクスチャー、スカーミオン）は $\pi_d(\mathcal M)$ で分類される。
\item 4次元ユークリッドのゲージ理論で作用有限の配位は、遠方の純ゲージ $A\to i\,g\,dg^{-1}$ を通して $\pi_3(G)$ で分類される。
\end{enumerate}
\end{mythm}
ここで「位相的に安定」とは、エネルギー有限のまま連続変形でほどけないという意味であり、エネルギー的に安定な静的解が実際に存在するかどうかは力学の問題である（デリックの定理、問題 \ref{ch03-prob-skyrmion}）。

\begin{physbox}[：対称性の破れと位相欠陥]
\begin{center}
\renewcommand{\arraystretch}{1.2}\small
\begin{tabular}{p{0.29\linewidth}p{0.15\linewidth}p{0.17\linewidth}p{0.25\linewidth}}
模型（破れ方） & $\mathcal M$ & 非自明な $\pi_k$ & 位相的対象（3次元空間）\\ \hline
実スカラー $\phi^4$：$\Z_2\to1$ & $\{\pm v\}$ & $\pi_0$ & ドメイン・ウォール\\
複素スカラー：$U(1)\to1$ & $S^1$ & $\pi_1=\Z$ & ひも（ANO 渦、宇宙ひも）\\
ジョージ・グラショウ：$SU(2)\to U(1)$（随伴） & $S^2$ & $\pi_2=\pi_3=\Z$ & モノポール\\
電弱：$SU(2)_L\times U(1)_Y\to U(1)_{\mathrm{em}}$ & $S^3$ & $\pi_3=\Z$ のみ & 安定な欠陥なし（テクスチャー）\\
カイラル：$SU(2)_L\times SU(2)_R\to SU(2)_V$ & $SU(2)\cong S^3$ & $\pi_3=\Z$ & スカーミオン（バリオン）\\
$SU(5)$ 大統一：$\to(SU(3)\times SU(2)\times U(1))/\Z_6$ & 12次元 & $\pi_2=\Z$ & 大統一モノポール\\
一軸ネマチック液晶 & $\RP^2$ & $\pi_1=\Z_2$, $\pi_2=\Z$ & 回位、点欠陥\\
$O(3)$ 模型（$\bm n\in S^2$） & $S^2$ & $\pi_2,\pi_3$ & 2次元ランプ、3次元ホップ粒子\\
4次元ユークリッド $SU(N)$ ゲージ理論 & --- & $\pi_3(G)=\Z$ & インスタントン、$\theta$ 真空
\end{tabular}
\end{center}
\end{physbox}

\subsection{トフーフト・ポリャコフ・モノポールの磁荷}
本書の規約（$D=d-iA$、生成子 $T^a=\sigma^a/2$、結合定数は $A$ に吸収）で、$SU(2)$ 随伴ヒッグス場 $\phi=\phi^aT^a$ の共変微分と場の強さは
\[
(D_\mu\phi)^a=\partial_\mu\phi^a+\epsilon^{abc}A^b_\mu\phi^c,\qquad F^a_{\mu\nu}=\partial_\mu A^a_\nu-\partial_\nu A^a_\mu+\epsilon^{abc}A^b_\mu A^c_\nu
\]
である（$D\phi=d\phi-i[A,\phi]$ と $[T^a,T^b]=i\epsilon^{abc}T^c$ から）。$\hat\phi=\phi/\lvert\phi\rvert$ の定義される領域で、\term{トフーフトの場の強さ}{'t Hooft field strength}
\[
\mathcal F_{\mu\nu}=\hat\phi^aF^a_{\mu\nu}-\epsilon^{abc}\hat\phi^aD_\mu\hat\phi^bD_\nu\hat\phi^c
\]
はゲージ不変であり、真空（$D\hat\phi=0$）では破れずに残る $U(1)$（$\hat\phi\cdot T$ で生成される）の場の強さに一致する。直接計算により
\begin{equation}
\mathcal F=d(\hat\phi^aA^a)-\tfrac12\epsilon^{abc}\hat\phi^a\,d\hat\phi^b\wedge d\hat\phi^c\label{eq:ch03-thooft}
\end{equation}
が示せる。

\begin{myprop}{磁荷は写像度で決まる}{ch03-magcharge}
遠方の球面 $S^2_\infty$（外向き法線で向きづける）上で $\phi$ が真空にあり $D\hat\phi\to0$ とすると
\[
\int_{S^2_\infty}\mathcal F=-4\pi\,\deg\bigl(\hat\phi\colon S^2_\infty\to S^2\bigr).
\]
特にハリネズミ $\hat\phi^a=\hat x^a$ では磁場は $\bm B=-\hat{\bm x}/r^2$ で、磁束は $-4\pi$。
\end{myprop}
\begin{note}
(\ref{eq:ch03-thooft}) の第1項は閉曲面上で積分すると消え、第2項は $-\omega$ の引き戻しなので、例 \ref{ex:ch03-deg2} から $-4\pi\deg\hat\phi$ を得る。ハリネズミの場合を直接確かめる：$D_i\hat\phi=0$ となる漸近形は $A^a_i=\epsilon_{aij}x^j/r^2$ であり、これから $\hat\phi^aF^a_{ij}=-2\epsilon_{aij}x^a/r^3+\epsilon_{ijk}x^k/r^3=-\epsilon_{ijk}x^k/r^3$、$B_k=\tfrac12\epsilon_{kij}\mathcal F_{ij}=-x^k/r^3$。符号は $D=d-iA$ という規約による（$D=d+iA$ なら逆符号）。
\end{note}
$U(1)$ の電荷は $\hat\phi\cdot T$ の固有値で測られ、$W^\pm$ ボソンは電荷 $\pm1$、もし基本表現（2重項）の場があれば電荷 $\pm\tfrac12$ をもつ。ディラックの条件 $q\int\mathcal F\in2\pi\Z$ は $q=\tfrac12$ でもみたされる（$\tfrac12\cdot4\pi=2\pi$）。結合定数 $e$ をあらわに書けば磁荷は $4\pi/e$ の整数倍である。ディラック・モノポールとの関係は第5章、モノポール束としての解釈は第7章で扱う。

\section{問題}

\subsection{ホモトピー群と写像度}

\begin{problem}[難易度 B]\label{ch03-prob-abelian}
$n\ge2$ とし、$f,g\colon(I^n,\partial I^n)\to(X,x_0)$ に対し、第 $i$ 座標方向に並べた和 $f+_ig$（$i=1,2$）を定義 \ref{def:ch03-pin} と同様に定める。定値写像を $c$ とする。\\
(1) 写像として $(f+_1g)+_2(h+_1k)=(f+_2h)+_1(g+_2k)$ が成り立つことを確かめよ。\\
(2) (1) を用いて、ホモトピー類について $[f]+_1[g]=[f]+_2[g]=[g]+_1[f]$ を示し、$\pi_n(X,x_0)$ が可換であることを結論せよ。\\
(3) $n=1$ でこの議論が使えない理由を述べよ。
\end{problem}

\begin{solution}
(1) 両辺とも、$I^n$ を $s_1\le\frac12$ / $s_1\ge\frac12$ と $s_2\le\frac12$ / $s_2\ge\frac12$ で4つに分け、
左下 $(s_1,s_2\le\frac12)$ に $f$、右下に $g$、左上に $h$、右上に $k$ を（それぞれの小区画を $I^n$ に拡大して）置いた写像である。実際、左辺で $s_2\le\frac12$ の部分は $(f+_1g)(s_1,2s_2,\dots)$ で、$s_1\le\frac12$ なら $f(2s_1,2s_2,\dots)$。右辺で $s_1\le\frac12$ の部分は $(f+_2h)(2s_1,s_2,\dots)$ で、$s_2\le\frac12$ なら $f(2s_1,2s_2,\dots)$。他の区画も同様に一致する。

(2) $[c]$ は $+_1$, $+_2$ どちらについても単位元である（$f+_ic\simeq f$ は区画を伸縮するホモトピーで示せる）。(1) で $h=k=c$ や $f=k=c$ などを代入すると
\begin{align*}
[f]+_1[g]&=([f]+_2[c])+_1([c]+_2[g])=([f]+_1[c])+_2([c]+_1[g])=[f]+_2[g],\\
[f]+_1[g]&=([c]+_2[f])+_1([g]+_2[c])=([c]+_1[g])+_2([f]+_1[c])=[g]+_2[f].
\end{align*}
第1式から2つの演算は一致し、それを第2式に使えば $[f]+_1[g]=[g]+_1[f]$。よって $\pi_n(X,x_0)$ は可換である。図（本章の基本事項）は、この議論を「箱をすれ違わせる」連続変形として描いたものである。

(3) $n=1$ では座標が1つしかなく、$+_2$ が定義できない。幾何学的には、区間 $I$ 上に並んだ $f$ と $g$ の区画は、端点を基点に固定したまま互いに追い越すことができない。実際、基本群は一般に非可換である（例：8の字の基本群は自由群 $F_2$、第2章）。物理的には、2次元以上の「欠陥の配置空間」では欠陥どうしが空間の余分な方向を使って避け合えるが、ひも（$\pi_1$）はそうできない、という違いに対応する。
\end{solution}

\begin{problem}[難易度 B]\label{ch03-prob-degint}
単位ベクトル場 $\bm n=(n^1,n^2,n^3)$、$\lvert\bm n\rvert=1$ について考える。\\
(1) 2形式 $\omega=\frac12\epsilon_{abc}n^a\,dn^b\wedge dn^c$ を $S^2$ の極座標 $\bm n=(\sin\theta\cos\varphi,\sin\theta\sin\varphi,\cos\theta)$ で表すと $\sin\theta\,d\theta\wedge d\varphi$ になることを示せ。\\
(2) $\bm n\colon\R^2\to S^2$ が遠方で一定値に近づくとき、$Q=\frac1{4\pi}\int d^2x\,\bm n\cdot(\partial_1\bm n\times\partial_2\bm n)$ は写像度に等しいことを説明せよ。\\
(3) 微小変形 $\bm n\to\bm n+\delta\bm n$（$\bm n\cdot\delta\bm n=0$、$\delta\bm n$ は遠方で $0$）に対し、$q=\bm n\cdot(\partial_1\bm n\times\partial_2\bm n)$ の変化が全微分
\[
\delta q=\partial_1\bigl(\bm n\cdot(\delta\bm n\times\partial_2\bm n)\bigr)+\partial_2\bigl(\bm n\cdot(\partial_1\bm n\times\delta\bm n)\bigr)
\]
になることを示し、$Q$ が連続変形で不変であることを直接確かめよ。
\end{problem}

\begin{solution}
(1) $\bm e_r=\bm n$、$\bm e_\theta=\partial_\theta\bm n$、$\bm e_\varphi=\partial_\varphi\bm n/\sin\theta$ は右手系の正規直交基底（$\bm e_r\times\bm e_\theta=\bm e_\varphi$）。$d\bm n=\bm e_\theta\,d\theta+\sin\theta\,\bm e_\varphi\,d\varphi$ なので、外積とウェッジ積を組み合わせた $(d\bm n\times d\bm n)^a=\epsilon_{abc}dn^b\wedge dn^c$ は
\[
d\bm n\times d\bm n=\sin\theta\,(\bm e_\theta\times\bm e_\varphi)\,d\theta\wedge d\varphi+\sin\theta\,(\bm e_\varphi\times\bm e_\theta)\,d\varphi\wedge d\theta=2\sin\theta\,\bm e_r\,d\theta\wedge d\varphi .
\]
よって $\omega=\frac12\bm n\cdot(d\bm n\times d\bm n)=\sin\theta\,d\theta\wedge d\varphi$。これは外向き法線で向きづけた $S^2$ の面積形式で、$\int_{S^2}\omega=4\pi$。

(2) 立体射影により $\R^2\cup\{\infty\}\cong S^2$ であり、遠方で一定値に近づく $\bm n$ は連続写像 $S^2\to S^2$ に拡張される（第1章）。$\bm n^*\omega=\frac12\epsilon_{abc}n^a(\partial_in^b\partial_jn^c)dx^i\wedge dx^j=\bm n\cdot(\partial_1\bm n\times\partial_2\bm n)\,dx^1\wedge dx^2$ なので、定義 \ref{def:ch03-degree} の積分公式から $Q=\frac{1}{4\pi}\int\bm n^*\omega=\deg\bm n$。（$\infty$ の1点は測度ゼロで積分に寄与しない。）

(3) $q=\epsilon_{abc}n^a\partial_1n^b\partial_2n^c$ を変分すると3項が出る。\\
第1項 $\epsilon_{abc}\delta n^a\partial_1n^b\partial_2n^c=\delta\bm n\cdot(\partial_1\bm n\times\partial_2\bm n)$：$\partial_i\bm n\perp\bm n$ なので $\partial_1\bm n\times\partial_2\bm n\parallel\bm n$ で、$\delta\bm n\perp\bm n$ より $0$。\\
第2・3項は部分積分の形に書き換える：
\begin{align*}
\epsilon_{abc}n^a\partial_1\delta n^b\partial_2n^c&=\partial_1(\epsilon_{abc}n^a\delta n^b\partial_2n^c)-\epsilon_{abc}\partial_1n^a\delta n^b\partial_2n^c-\epsilon_{abc}n^a\delta n^b\partial_1\partial_2n^c,\\
\epsilon_{abc}n^a\partial_1n^b\partial_2\delta n^c&=\partial_2(\epsilon_{abc}n^a\partial_1n^b\delta n^c)-\epsilon_{abc}\partial_2n^a\partial_1n^b\delta n^c-\epsilon_{abc}n^a\partial_1\partial_2n^b\delta n^c .
\end{align*}
2階微分の項の和は $-\epsilon_{abc}n^a(\delta n^bX^c+X^b\delta n^c)=0$（$X=\partial_1\partial_2\bm n$、$b,c$ について反対称）。残りの2項はどちらも $\delta\bm n\cdot(\partial_2\bm n\times\partial_1\bm n)$ に等しく、第1項と同じ理由で $0$。したがって
\[
\delta q=\partial_1\bigl(\bm n\cdot(\delta\bm n\times\partial_2\bm n)\bigr)+\partial_2\bigl(\bm n\cdot(\partial_1\bm n\times\delta\bm n)\bigr),
\]
$\delta\bm n$ が遠方で消えれば $\delta Q=0$。$Q$ は局所的な変形で変わらない。$Q$ が整数であること（(2)）と合わせると、$Q$ は連続変形で不変な位相的電荷である。
\end{solution}

\begin{problem}[難易度 A]\label{ch03-prob-degex}
(1) ハリネズミ写像 $\bm n(\bm x)=\bm x/\lvert\bm x\rvert$ を $S^2_R=\{\lvert\bm x\rvert=R\}$（外向き法線で向きづける）から $S^2$ への写像とみて、その写像度が $1$ であることを示せ。また $\bm n(\bm x)=-\bm x/\lvert\bm x\rvert$ の写像度を求めよ。\\
(2) $n\in\Z$ に対し、極座標で $f_n(\theta,\varphi)=(\theta,n\varphi)$ と定める写像 $f_n\colon S^2\to S^2$ が連続であることを確かめ、その写像度が $n$ であることを、正則値の逆像を数える方法と積分公式の両方で示せ。
\end{problem}

\begin{solution}
(1) $\bm x\mapsto\bm x/R$ は向きを保つ微分同相 $S^2_R\to S^2$ なので、写像度は $1$（任意の点の逆像は1点で、そこで向きを保つ）。対蹠写像 $\bm x\mapsto-\bm x$ は $\R^3$ の向きを反転し（$\det(-\mathbf 1_3)=-1$）、外向き法線を外向き法線に写すので、$S^2$ の向きを反転する。よって $-\bm x/\lvert\bm x\rvert$ の写像度は $-1$。一般に $S^n$ の対蹠写像の写像度は $(-1)^{n+1}$ である。

(2) 連続性：$\theta=0,\pi$（極）では $\varphi$ によらず極に写るので、$f_n$ は極でも連続である。\\
逆像を数える方法：$n\ge1$ のとき、極でない点 $(\theta_0,\varphi_0)$ の逆像は $(\theta_0,(\varphi_0+2\pi j)/n)$、$j=0,\dots,n-1$ の $n$ 点。局所座標 $(\theta,\varphi)$ でヤコビ行列は $\operatorname{diag}(1,n)$ で行列式 $n>0$、どの点でも向きを保つので $\deg f_n=n$。$n\le-1$ では $\lvert n\rvert$ 点で各々行列式 $n<0$ なので $\deg f_n=-\lvert n\rvert=n$。$n=0$ では像は経線1本で全射でないので、像にない点を正則値にとれば逆像は空で $\deg f_0=0$。\\
積分公式：$f_n^*(\sin\theta\,d\theta\wedge d\varphi)=\sin\theta\,d\theta\wedge n\,d\varphi$ なので
\[
\deg f_n=\frac{1}{4\pi}\int_0^\pi\!\!\int_0^{2\pi}n\sin\theta\,d\varphi\,d\theta=\frac{n\cdot4\pi}{4\pi}=n .
\]
同じ結果は、$\bm n=(\sin\theta\cos n\varphi,\sin\theta\sin n\varphi,\cos\theta)$ について $\partial_\theta\bm n\times\partial_\varphi\bm n=n\sin\theta\,\bm n$ を確かめても得られる。物理的には、赤道に沿って $\bm n$ の水平成分が $n$ 回まわる「巻き数 $n$ のハリネズミ」であり、ジョージ・グラショウ模型では磁荷 $n$ 単位のモノポールの漸近形を与える。
\end{solution}

\begin{problem}[難易度 A]\label{ch03-prob-table}
命題 \ref{prop:ch03-pinprop} と定理 \ref{thm:ch03-spheres} を用いて次を求めよ。\\
(1) $\pi_n(S^1)$（$n\ge2$）、$\pi_n(T^2)$（$n\ge2$）。\\
(2) $\pi_2(\RP^2)$、$\pi_3(SO(3))$、$\pi_2(S^2\times S^1)$。\\
(3) $\pi_1(S^3)$、$\pi_2(S^3)$、$\pi_3(SU(2)\times SU(2))$。\\
(4) 2次元トーラス上のスカラー場 $\phi\colon T^2\to S^1$ は渦をもちうるが、$S^1$ 値の場で「点状のソリトン」（$\pi_2$ で分類される対象）がありえないのはなぜか。
\end{problem}

\begin{solution}
(1) $\R\to S^1$ は被覆で $\R$ は可縮なので、$n\ge2$ で $\pi_n(S^1)\cong\pi_n(\R)=0$。$T^2=S^1\times S^1$ なので $\pi_n(T^2)=\pi_n(S^1)^2=0$（$n\ge2$）。（$\R^2\to T^2$ も被覆である。）

(2) $S^2\to\RP^2$ は2重被覆なので $\pi_2(\RP^2)\cong\pi_2(S^2)\cong\Z$。$SU(2)\to SO(3)$ も2重被覆で $\pi_3(SO(3))\cong\pi_3(S^3)\cong\Z$。積の公式より $\pi_2(S^2\times S^1)\cong\pi_2(S^2)\oplus\pi_2(S^1)=\Z$。

(3) $k<n$ で $\pi_k(S^n)=0$ なので $\pi_1(S^3)=\pi_2(S^3)=0$。$\pi_3(SU(2)\times SU(2))\cong\Z\oplus\Z$。$SO(4)=(SU(2)\times SU(2))/\Z_2$ なので被覆により $\pi_3(SO(4))\cong\Z\oplus\Z$ でもある（命題 \ref{prop:ch03-liegroups} の表）。カイラル対称性 $SU(2)_L\times SU(2)_R$ のゲージ変換を考えると、左右それぞれに独立な巻き付き数がある。

(4) $\pi_2(S^1)=0$ なので、点を囲む $S^1$ 値の場（2次元では点を囲むのは $S^1$ で、これは渦を与える $\pi_1$）は別として、「空間全体を覆う」タイプの $S^2\to S^1$ の配位はすべて自明である。より具体的には、$S^1$ 値の場の位相的電荷は各ループに沿った巻き数（$\pi_1$ や $H^1$ の情報）だけであり、$\pi_2$ 型の電荷は存在しない。点状のソリトン（ランプ）を作るには、$\pi_2\neq0$ の標的（例えば $S^2$）が必要である。
\end{solution}

\subsection{長完全列と真空多様体}

\begin{problem}[難易度 B]\label{ch03-prob-su2u1}
(1) $G=SU(2)$、$H=U(1)=\{\exp(i\alpha\sigma^3)\}$ として、定理 \ref{thm:ch03-GH} の完全列から $\pi_2(SU(2)/U(1))\cong\Z$、$\pi_1(SU(2)/U(1))=0$ を導け。\\
(2) $U(\theta,\varphi)=e^{-i\varphi\sigma^3/2}e^{-i\theta\sigma^2/2}e^{i\varphi\sigma^3/2}$ とおく。$U\sigma^3U^\dagger=\hat{\bm x}\cdot\bm\sigma$（$\hat{\bm x}$ は極座標 $(\theta,\varphi)$ の単位ベクトル）を示せ。$U$ は北極 $\theta=0$ では $\varphi$ によらないが、南極 $\theta=\pi$ では $U=-i\sigma^2e^{i\varphi\sigma^3}$ となることを示し、境界写像 $\partial$ がハリネズミを $\pi_1(H)$ の生成元に写すことを説明せよ。\\
(3) 同じ $S^2$ を $SO(3)/SO(2)$ と見ると、$\pi_2(S^2)\cong\ker\bigl(\pi_1(SO(2))\to\pi_1(SO(3))\bigr)=2\Z$ となる。(2) との関係を説明せよ。
\end{problem}

\begin{solution}
(1) 定理 \ref{thm:ch03-GH} の完全列は
\[
0=\pi_2(SU(2))\to\pi_2(G/H)\xrightarrow{\partial}\pi_1(U(1))\to\pi_1(SU(2))=0\to\pi_1(G/H)\to\pi_0(U(1))=0 .
\]
よって $\partial$ は同型 $\pi_2(G/H)\cong\pi_1(U(1))\cong\Z$、$\pi_1(G/H)=0$。もちろん $G/H\cong S^2$（第1章）と整合する。

(2) $e^{-i\alpha X}Ye^{i\alpha X}$ の展開と $[\sigma^2,\sigma^3]=2i\sigma^1$、$[\sigma^3,\sigma^1]=2i\sigma^2$ を用いると
\[
e^{-i\theta\sigma^2/2}\sigma^3e^{i\theta\sigma^2/2}=\sigma^3\cos\theta+\sigma^1\sin\theta,\qquad e^{-i\varphi\sigma^3/2}\sigma^1e^{i\varphi\sigma^3/2}=\sigma^1\cos\varphi+\sigma^2\sin\varphi .
\]
$e^{i\varphi\sigma^3/2}$ は $\sigma^3$ と可換なので
\[
U\sigma^3U^\dagger=\sin\theta\cos\varphi\,\sigma^1+\sin\theta\sin\varphi\,\sigma^2+\cos\theta\,\sigma^3=\hat{\bm x}\cdot\bm\sigma .
\]
つまり $U$ はハリネズミ $\hat{\bm x}\mapsto\hat{\bm x}\cdot\bm\sigma\in G/H$ の $G$ への持ち上げである。$\theta=0$ では $U=e^{-i\varphi\sigma^3/2}e^{i\varphi\sigma^3/2}=\mathbf 1$。$\theta=\pi$ では $e^{-i\pi\sigma^2/2}=-i\sigma^2$ で、$\sigma^2$ と $\sigma^3$ が反可換なので $e^{-i\varphi\sigma^3/2}\sigma^2=\sigma^2e^{i\varphi\sigma^3/2}$。よって
\[
U(\pi,\varphi)=-i\sigma^2\,e^{i\varphi\sigma^3/2}e^{i\varphi\sigma^3/2}=-i\sigma^2e^{i\varphi\sigma^3}.
\]
北極（基点）付近では持ち上げは一意だが、南極では $\varphi$ に依存してしまい、そのずれは $H$ のループ $\varphi\mapsto e^{i\varphi\sigma^3}$（$\varphi\in[0,2\pi]$）である。これは $H=\{e^{i\alpha\sigma^3}\}$ をちょうど1周するので $\pi_1(H)$ の生成元である。境界写像 $\partial$ の構成（基本事項の【注】）はまさに「持ち上げの閉じきらない分」を取り出すものなので、$\partial[\text{ハリネズミ}]=\pm1$（符号は向きの約束による）。

(3) $SO(3)$ では $\pi_1(SO(2))=\Z\to\pi_1(SO(3))=\Z_2$ は2で割った余りをとる写像（$z$ 軸まわりの $2\pi$ 回転は $SO(3)$ の非自明なループ）なので、核は $2\Z$。一方 (2) のループ $e^{i\varphi\sigma^3}$ を $SO(3)$ に写すと、随伴作用は $z$ 軸まわりの角度 $2\varphi$ の回転なので、$\varphi\in[0,2\pi]$ で $SO(2)$ を2周する。したがってハリネズミ（$\pi_2(S^2)$ の生成元）は $\pi_1(SO(2))$ では $2$ に対応し、像が $2\Z$ になることと整合する。物理的には、$SO(2)$ のループを単位とすると最小のモノポールは「2単位」の磁荷をもつ。これは命題 \ref{prop:ch03-magcharge} の磁束 $4\pi$ が、整数電荷の場（随伴表現など）に対するディラックの最小値 $2\pi$ の2倍であることに対応する。奇数の元は $\pi_1(SO(3))=\Z_2$ の非自明な元に写り、滑らかな配位としては実現されない（ゲージ束のトポロジーとしての解釈は第7章）。
\end{solution}

\begin{problem}[難易度 A]\label{ch03-prob-ew}
電弱理論の真空多様体は $\mathcal M\cong(SU(2)_L\times U(1)_Y)/U(1)_{\mathrm{em}}\cong S^3$ である（第1章）。\\
(1) $\pi_0(\mathcal M)$、$\pi_1(\mathcal M)$、$\pi_2(\mathcal M)$、$\pi_3(\mathcal M)$ を求め、電弱理論には位相的に安定なドメイン・ウォール、ひも、モノポールがないことを結論せよ。\\
(2) 定理 \ref{thm:ch03-GH} の最後の主張を用いて、$\pi_3(\mathcal M)$ を $G=SU(2)\times U(1)$、$H=U(1)$ から計算せよ。
\end{problem}

\begin{solution}
(1) $S^3$ は弧状連結なので $\pi_0$ は1点、$k<3$ で $\pi_k(S^3)=0$ なので $\pi_1=\pi_2=0$、$\pi_3(S^3)=\Z$。定理 \ref{thm:ch03-defects} より、壁（$\pi_0$）、ひも（$\pi_1$）、モノポール（$\pi_2$）はいずれも位相的に安定ではない。電弱理論で議論される $Z$ ひもなどは、存在するとしても位相的には保護されず、力学的な理由でのみ準安定になりうる。$\pi_3=\Z$ はテクスチャー（3次元空間全体の配位）を許すが、ゲージ理論ではゲージ場が巻き付きをほどくことができ、安定なソリトンにはならない（大きなゲージ変換との関係は第8章）。

(2) $\pi_2(H)=0$ なので完全列 $\pi_3(H)\to\pi_3(G)\to\pi_3(G/H)\to\pi_2(H)=0$ より $\pi_3(G/H)\cong\coker(\pi_3(H)\to\pi_3(G))$。$\pi_3(U(1))=0$、$\pi_3(SU(2)\times U(1))=\pi_3(SU(2))\oplus\pi_3(U(1))=\Z$ なので $\pi_3(\mathcal M)\cong\Z$。(1) と一致する。
\end{solution}

\begin{problem}[難易度 C]\label{ch03-prob-gut}
$SU(5)$ 大統一理論で、随伴ヒッグスが $\operatorname{diag}(2,2,2,-3,-3)$ に比例する期待値をもつと、残る群は
\[
H=\{\operatorname{diag}(A,B):A\in U(3),\ B\in U(2),\ \det A\det B=1\}
\]
となる。\\
(1) $\Phi(U_3,U_2,e^{i\alpha})=\operatorname{diag}(e^{2i\alpha}U_3,\ e^{-3i\alpha}U_2)$ が全射準同型 $SU(3)\times SU(2)\times U(1)\to H$ を与え、その核が $\Z_6$ であることを示せ。すなわち $H\cong(SU(3)\times SU(2)\times U(1))/\Z_6$。\\
(2) $\pi_1(H)\cong\Z$ を示し、$\pi_2(SU(5)/H)\cong\Z$ を導け（大統一モノポール）。\\
(3) $SU(5)$ のループ $\gamma(t)=\operatorname{diag}(1,1,e^{2\pi it},1,e^{-2\pi it})$（$t\in[0,1]$）が $\pi_1(H)$ の生成元であること、および超電荷のループ $\alpha\colon0\to2\pi$（$U_3=U_2=\mathbf 1$）が生成元の6倍であることを示せ。この結果は最小の大統一モノポールの磁荷について何を意味するか。
\end{problem}

\begin{solution}
(1) $\Phi$ は準同型。$\det(e^{2i\alpha}U_3)\det(e^{-3i\alpha}U_2)=e^{6i\alpha}e^{-6i\alpha}=1$ なので像は $H$ に入る。全射：$\operatorname{diag}(A,B)\in H$ に対し $\det B=e^{-6i\alpha}$ となる $\alpha$ を選び、$U_2=e^{3i\alpha}B$、$U_3=e^{-2i\alpha}A$ とおくと $\det U_2=e^{6i\alpha}\det B=1$、$\det U_3=e^{-6i\alpha}\det A=\det B\det A=1$。核：$e^{2i\alpha}U_3=\mathbf 1$、$e^{-3i\alpha}U_2=\mathbf 1$ より $U_3=e^{-2i\alpha}\mathbf 1_3\in SU(3)$ で $e^{-6i\alpha}=1$、$U_2=e^{3i\alpha}\mathbf 1_2\in SU(2)$ で $e^{6i\alpha}=1$。よって $\alpha\in\frac{2\pi}{6}\Z$ で、核は $\{(e^{-2i\alpha}\mathbf 1,e^{3i\alpha}\mathbf 1,e^{i\alpha}):e^{6i\alpha}=1\}\cong\Z_6$。

(2) 準同型 $\det_2\colon H\to U(1)$, $\operatorname{diag}(A,B)\mapsto\det B$ は全射で、核は $\det A=\det B=1$ の部分 $SU(3)\times SU(2)$。リー群の全射準同型はファイバー束なので長完全列
\[
\pi_1(SU(3)\times SU(2))=0\to\pi_1(H)\to\pi_1(U(1))=\Z\to\pi_0(SU(3)\times SU(2))=0
\]
より $\pi_1(H)\cong\Z$（同型は $\det B$ の巻き数）。$SU(5)$ は単連結で $\pi_2(SU(5))=0$ なので、定理 \ref{thm:ch03-GH} より $\pi_2(SU(5)/H)\cong\pi_1(H)\cong\Z$。

(3) $\gamma(t)$ では $B=\operatorname{diag}(1,e^{-2\pi it})$、$\det B=e^{-2\pi it}$ はちょうど1回（負の向きに）まわるので、$\gamma$ は $\pi_1(H)$ の生成元（の逆）である。一方、超電荷のループ $\Phi(\mathbf 1,\mathbf 1,e^{i\alpha})=\operatorname{diag}(e^{2i\alpha}\mathbf 1_3,e^{-3i\alpha}\mathbf 1_2)$ では $\det B=e^{-6i\alpha}$ が $\alpha\in[0,2\pi]$ で6回まわるので、生成元の6倍（符号を除く）。

別の見方：$\gamma(t)=\Phi(U_3(t),U_2(t),e^{i\pi t/3})$ と書ける。実際 $U_3(t)=\exp(2\pi it\operatorname{diag}(-\frac13,-\frac13,\frac23))$、$U_2(t)=\exp(2\pi it\operatorname{diag}(\frac12,-\frac12))$ とすると
\[
e^{2\pi it/3}U_3(t)=\operatorname{diag}(1,1,e^{2\pi it}),\qquad e^{-i\pi t}U_2(t)=\operatorname{diag}(1,e^{-2\pi it}).
\]
$t=1$ で $(U_3,U_2,e^{i\alpha})=(e^{-2\pi i/3}\mathbf 1,-\mathbf 1,e^{i\pi/3})$ は $\Z_6$ の生成元であり、$\gamma$ は $SU(3)\times SU(2)\times U(1)$ では閉じず、$\Z_6$ の分だけずれた道である。超電荷の $\frac16$ 周、色の中心 $\Z_3$、弱アイソスピンの中心 $\Z_2$ を組み合わせて初めて $H$ の中で閉じる。

物理的意味：モノポールの磁荷は $\pi_2(SU(5)/H)\cong\pi_1(H)$ の元で、最小のモノポールは $\gamma$ に対応する。その生成子 $\operatorname{diag}(0,0,1,0,-1)$ は超電荷 $Y=\operatorname{diag}(-\frac13,-\frac13,-\frac13,\frac12,\frac12)$、色の生成子 $\operatorname{diag}(-\frac13,-\frac13,\frac23,0,0)$、弱アイソスピン $\operatorname{diag}(0,0,0,\frac12,-\frac12)$ の組み合わせ $-Y+(\text{色})+(T_3)$ であり、最小の大統一モノポールは純粋に $U(1)_Y$ の磁荷ではなく、色磁荷と弱磁荷も帯びる。超電荷だけで測れば、$Y$ 方向の磁束は「$U(1)_Y$ のディラック・モノポールとして素朴に期待される最小値」の $\frac16$ であり、これは分数電荷 $Y\in\frac16\Z$ をもつクォークとディラック条件が整合するための仕組みである。キブル機構でこのモノポールが大量に生成されることが宇宙論の\term{モノポール問題}{monopole problem}である。
\end{solution}

\begin{problem}[難易度 C]\label{ch03-prob-nematic}
ネマチック液晶の配向は、棒状分子の向き $\bm n\in S^2$ を $\bm n\sim-\bm n$ で同一視したもの（\term{ダイレクター}{director}）で表される。オーダー・パラメータは $Q_{ij}=n_in_j-\frac13\delta_{ij}$ で、$SO(3)$ は $Q\mapsto RQR^{\mathsf T}$ と作用する。\\
(1) $Q_0=\operatorname{diag}(-\frac13,-\frac13,\frac23)$ の固定部分群 $H\subset SO(3)$ が $O(2)$ と同型で2つの連結成分をもつこと、および真空多様体が $\RP^2$ であることを示せ。\\
(2) $SU(2)\to SO(3)$ による $H$ の逆像 $\tilde H$ を用いて、$\pi_1(\RP^2)\cong\Z_2$、$\pi_2(\RP^2)\cong\Z$ を導け。物理的には何が存在するか。\\
(3) $\pi_1(\RP^2)$ の非自明な元は $\pi_2(\RP^2)\cong\Z$ に $N\mapsto-N$ と作用することを説明し、点欠陥の電荷は $\lvert N\rvert$ でしか意味をもたないことを論ぜよ。\\
(4) \term{二軸性ネマチック}{biaxial nematic}では $Q_0=\operatorname{diag}(a,b,c)$（相異なる、和は $0$）である。$\pi_1(\mathcal M)$ が四元数群 $Q_8=\{\pm1,\pm i\sigma^1,\pm i\sigma^2,\pm i\sigma^3\}$ となることを示し、線欠陥の種類が何通りあるかを答えよ。
\end{problem}

\begin{solution}
(1) $RQ_0R^{\mathsf T}=Q_0$ は $R$ が $Q_0$ と可換、すなわち固有空間 $\R^2\oplus\R e_3$ を保つことと同値で、$R=\operatorname{diag}(r,\epsilon)$、$r\in O(2)$、$\epsilon=\pm1$。$\det R=1$ より $\epsilon=\det r$。$r\mapsto\operatorname{diag}(r,\det r)$ は同型 $O(2)\cong H$ で、$\det r=\pm1$ の2成分をもつ。軌道は $\{n n^{\mathsf T}-\frac13\mathbf 1:\bm n\in S^2\}$ で、$nn^{\mathsf T}$ は $\pm\bm n$ を同一視したもの（直線への射影）なので $\mathcal M\cong S^2/(\bm n\sim-\bm n)=\RP^2$。軌道・固定部分群定理より $\mathcal M\cong SO(3)/O(2)$。

(2) $\tilde H$ の単位元成分は $\{e^{i\alpha\sigma^3}\}\cong U(1)$ で、もう1つの成分は $i\sigma^1e^{i\alpha\sigma^3}$（$x$ 軸まわりの $\pi$ 回転の持ち上げ、$\bm n\to-\bm n$ に対応）。$\mathcal M=SU(2)/\tilde H$ で $SU(2)$ は単連結なので、定理 \ref{thm:ch03-GH} より
\[
\pi_1(\RP^2)\cong\pi_0(\tilde H)\cong\Z_2,\qquad\pi_2(\RP^2)\cong\pi_1(\tilde H)\cong\pi_1(U(1))\cong\Z .
\]
（被覆 $S^2\to\RP^2$ からも $\pi_2(\RP^2)\cong\pi_2(S^2)$ が得られる。）物理的には、$\pi_1=\Z_2$ はダイレクターが1周で $\bm n\to-\bm n$ と半回転する線欠陥（\term{回位}{disclination}、強さ $\frac12$）を、$\pi_2=\Z$ はハリネズミ型の点欠陥を与える。回位は2本合わさると消えうる（$\Z_2$）。

(3) $\pi_1(\RP^2)$ の非自明なループを $S^2$ に持ち上げると、$\bm n$ から $-\bm n$ への道になる。点欠陥を囲む球面上の配位をこのループに沿って運ぶと、持ち上げた写像 $S^2\to S^2$ は対蹠写像と合成されたものに変わる。対蹠写像の写像度は $-1$（問題 \ref{ch03-prob-degex}）なので、作用は $N\mapsto-N$。基点を決めない欠陥の分類は $[S^2,\RP^2]=\pi_2(\RP^2)/\pi_1(\RP^2)=\Z/(N\sim-N)$、すなわち $\lvert N\rvert\in\{0,1,2,\dots\}$ である。実際、ハリネズミ $\bm n=\hat{\bm x}$ と $\bm n=-\hat{\bm x}$ はダイレクターとしては同一の配位である。物理的には、点欠陥を回位のまわりに1周させると電荷の符号が反転する。

(4) 固定部分群は、$Q_0$ の3つの（相異なる固有値の）固有軸をそれぞれ保つ回転、すなわち $D_2=\{\mathbf 1,\operatorname{diag}(1,-1,-1),\operatorname{diag}(-1,1,-1),\operatorname{diag}(-1,-1,1)\}\cong\Z_2\times\Z_2$（各座標軸まわりの $\pi$ 回転）。$x^k$ 軸まわりの $\pi$ 回転の $SU(2)$ での持ち上げは $\pm e^{-i\pi\sigma^k/2}=\mp i\sigma^k$ なので、$\tilde H=Q_8$。$SU(2)$ は単連結なので $\pi_1(\mathcal M)\cong\pi_0(\tilde H)=Q_8$（$\tilde H$ は離散）。これは非可換であり（$(i\sigma^1)(i\sigma^2)=-(i\sigma^2)(i\sigma^1)$）、線欠陥は基点を固定しないので共役類で分類される。$Q_8$ の共役類は $\{1\}$, $\{-1\}$, $\{\pm i\sigma^1\}$, $\{\pm i\sigma^2\}$, $\{\pm i\sigma^3\}$ の5つなので、自明なものを含めて5種類（非自明なものは4種類）。非可換性の帰結として、$i\sigma^1$ 型と $i\sigma^2$ 型の線欠陥は互いをすり抜けられず、交差させると交換子 $-1$ に対応する線欠陥でつながれる。
\end{solution}

\begin{problem}[難易度 B]\label{ch03-prob-kink}
1次元空間の実スカラー場で、作用 $S=\int d^2x\,\bigl[-\frac12\partial_\mu\phi\,\partial^\mu\phi-V(\phi)\bigr]$、$V=\frac\lambda4(\phi^2-v^2)^2$ をもつものを考える（計量は mostly plus）。\\
(1) 真空多様体 $\mathcal M$ と $\pi_0(\mathcal M)$ を求め、静的エネルギー $E=\int dx\,[\frac12\phi'^2+V]$ が有限の配位が4つのセクターに分かれることを示せ。\\
(2) $E=\int dx\,\frac12\bigl(\phi'\mp\sqrt{2V}\bigr)^2\pm\int dx\,\phi'\sqrt{2V}$ と書き直して、$\phi(\mp\infty)=\mp v$ のセクターで $E\ge\frac{2\sqrt2}{3}\sqrt\lambda\,v^3$ を示せ。等号を与える解（\term{キンク}{kink}）を求め、真空のまわりの粒子の質量 $m$ を用いて $\phi=v\tanh(mx/2)$、壁の張力が $\sigma=m^3/(3\lambda)$ となることを示せ。\\
(3) 3次元空間ではこの解は何を表すか。
\begin{center}
\begin{tikzpicture}[scale=0.9]
  \draw[-{Stealth[length=2mm]}] (-3.6,0) -- (3.8,0) node[right] {$x$};
  \draw[-{Stealth[length=2mm]}] (0,-1.5) -- (0,1.7) node[above] {$\phi$};
  \draw[dashed,gray] (-3.5,1) -- (3.5,1) node[right] {\footnotesize $+v$};
  \draw[dashed,gray] (-3.5,-1) -- (3.5,-1) node[right] {\footnotesize $-v$};
  \draw[very thick,blue!70!black,domain=-3.5:3.5,samples=80,smooth] plot (\x,{tanh(1.2*\x)});
  \fill[red!40,opacity=0.5,domain=-2.5:2.5,samples=80,smooth] plot (\x,{1.2/(cosh(1.2*\x)^4)}) -- (2.5,0) -- (-2.5,0) -- cycle;
  \node[red!60!black] at (1.5,0.75) {\footnotesize エネルギー密度};
  \node[blue!70!black] at (-2.3,-0.65) {\footnotesize $v\tanh(mx/2)$};
\end{tikzpicture}
\zucap{キンク解（青）とそのエネルギー密度（赤、$\propto\cosh^{-4}(mx/2)$）。場は $x\to\mp\infty$ で異なる真空 $\mp v$ に近づき、その間でエネルギーが局在する。}
\end{center}
\end{problem}

\begin{solution}
(1) $V=0$ となるのは $\phi=\pm v$ なので $\mathcal M=\{-v,+v\}$、$\pi_0(\mathcal M)$ は2点からなる。$E<\infty$ には $x\to\pm\infty$ で $V\to0$ が必要で、$\phi(\pm\infty)\in\{\pm v\}$。左右の値の組 $(\phi(-\infty),\phi(+\infty))$ は連続な時間発展で変わらない（変えるには無限遠までの無限大のエネルギーが必要）ので、4つのセクター $(\pm v,\pm v)$ に分かれる。$(-v,+v)$ がキンク、$(+v,-v)$ が反キンク、残り2つが真空セクターである。

(2) $\frac12(\phi'\mp\sqrt{2V})^2=\frac12\phi'^2+V\mp\phi'\sqrt{2V}$ なので書き直しは恒等式。$\sqrt{2V}=\sqrt{\lambda/2}\,\lvert v^2-\phi^2\rvert$ で、$W(\phi)=\sqrt{\lambda/2}(v^2\phi-\phi^3/3)$ とおくと、$\lvert\phi\rvert\le v$ の範囲で $\phi'\sqrt{2V}=\frac{d}{dx}W(\phi(x))$。キンク・セクター（$\phi(-\infty)=-v$、$\phi(\infty)=v$）で上の符号をとると
\[
E\ge W(v)-W(-v)=\sqrt{\frac\lambda2}\cdot2\Bigl(v^3-\frac{v^3}3\Bigr)=\sqrt{\frac\lambda2}\cdot\frac{4v^3}{3}=\frac{2\sqrt2}{3}\sqrt\lambda\,v^3 .
\]
（$\lvert\phi\rvert>v$ となる区間があっても、そこでは被積分関数がさらに大きくなるだけで下限は変わらない。）等号は $\phi'=\sqrt{2V}=\sqrt{\lambda/2}(v^2-\phi^2)$ のとき。$\phi=v\tanh(kx)$ とおくと $\phi'=vk(1-\tanh^2kx)=\frac kv(v^2-\phi^2)$ なので $k=v\sqrt{\lambda/2}$。一方 $V''(\phi)=\lambda(3\phi^2-v^2)$ より $m^2=V''(v)=2\lambda v^2$、$m=v\sqrt{2\lambda}=2k$。よって $\phi=v\tanh(mx/2)$。張力（エネルギー）は $\sigma=\frac{2\sqrt2}{3}\sqrt\lambda v^3$ で、$m^3=2\sqrt2\lambda^{3/2}v^3$ を使うと $\sigma=\frac{m^3}{3\lambda}$。弱結合 $\lambda\ll1$ で $\sigma\gg m$ となり、キンクは素粒子より重い古典的な対象である。

(3) 3次元空間では $\phi(x,y,z)=v\tanh(mx/2)$ は $x=0$ の平面に局在した\term{ドメイン・ウォール}{domain wall}であり、単位面積あたりのエネルギーが $\sigma$ である。壁は2つの真空 $\pm v$ の領域を隔て、壁を横切る線分の両端点（$S^0$）が $\mathcal M$ の異なる成分に写ることが位相的安定性の理由である（$\pi_0(\mathcal M)$）。宇宙論では、離散対称性の自発的破れがあると壁が生じ、そのエネルギーが宇宙を支配しうる（ドメイン・ウォール問題）。
\end{solution}

\begin{problem}[難易度 B]\label{ch03-prob-sun}
(1) $SU(N)$ は $\C^N$ の単位球面 $S^{2N-1}$ に推移的に作用し、$e_N=(0,\dots,0,1)^{\mathsf T}$ の固定部分群が $SU(N-1)$ であることを示せ（$N\ge2$）。\\
(2) ファイバー束 $SU(N-1)\to SU(N)\to S^{2N-1}$ の長完全列と $SU(2)\cong S^3$ から、帰納法で $\pi_1(SU(N))=\pi_2(SU(N))=0$、$\pi_3(SU(N))\cong\Z$ を示せ。\\
(3) 左上の $2\times2$ ブロックへの埋め込み $SU(2)\hookrightarrow SU(N)$ が $\pi_3$ の同型を誘導することを示し、ゲージ理論での意味（$SU(2)$ のインスタントンを $SU(N)$ に埋め込むこと）を述べよ。
\end{problem}

\begin{solution}
(1) 単位ベクトル $u$ に対し、$u$ を最後の列とする正規直交基底をグラム・シュミットで作ってユニタリ行列 $U$ とし、第1列に位相を掛けて $\det U=1$ にすれば $Ue_N=u$。よって推移的。$Ue_N=e_N$ なら最後の列は $e_N$、ユニタリ性から最後の行も $e_N^{\mathsf T}$ で $U=\operatorname{diag}(U',1)$、$\det U'=1$。固定部分群は $SU(N-1)$ であり、$SU(N)/SU(N-1)\cong S^{2N-1}$。

(2) 長完全列の一部
\[
\pi_{k+1}(S^{2N-1})\to\pi_k(SU(N-1))\xrightarrow{i_*}\pi_k(SU(N))\to\pi_k(S^{2N-1}).
\]
$N\ge3$ では $2N-1\ge5$ なので $k\le3$ に対し $\pi_k(S^{2N-1})=\pi_{k+1}(S^{2N-1})=0$（$k+1\le4<5$）。よって $i_*$ は $k\le3$ で同型 $\pi_k(SU(N-1))\cong\pi_k(SU(N))$。$SU(2)\cong S^3$ から $\pi_1=\pi_2=0$、$\pi_3=\Z$ で、帰納法により任意の $N\ge2$ で同じ。

(3) (2) の同型はすべて左上ブロックへの包含 $SU(N-1)\hookrightarrow SU(N)$ が誘導するものなので、合成 $SU(2)\hookrightarrow SU(N)$ も $\pi_3$ の同型を誘導する。また $g$ を左上ブロックに埋め込んでもトレース $\tr(g^{-1}dg)^3$ は変わらない（残りのブロックは単位行列で $g^{-1}dg$ の寄与ゼロ）ので、巻き付き数 $w$ も保たれる。ゲージ理論では、$SU(2)$ のインスタントン数1の配位（BPST インスタントン、第8章）を $SU(N)$ に埋め込むと、$SU(N)$ のインスタントン数1の配位が得られ、それが $\pi_3(SU(N))$ の生成元を実現する。$\pi_3(SU(N))=\Z$ であって $\Z^{N-1}$ などではないので、$SU(N)$ でもインスタントン数は1つの整数であり、$\theta$ パラメータも1つである。
\end{solution}

\begin{problem}[難易度 C]\label{ch03-prob-so3}
$p\colon SU(2)\to SO(3)$ を2重被覆（随伴表現）とする。\\
(1) ファイブレーション $\Z_2\to SU(2)\to SO(3)$ の長完全列から、$n\ge2$ で $p_*\colon\pi_n(SU(2))\to\pi_n(SO(3))$ が同型であることを示せ。特に $\pi_3(SO(3))\cong\Z$ で、$[p]$ が生成元である。また、任意の $f\colon S^3\to SO(3)$ は $f=p\circ g$ と持ち上がることを説明せよ。\\
(2) $p$ を向きのついた3次元多様体の間の写像 $S^3\to SO(3)\cong\RP^3$ とみると、その写像度は $2$ であることを示せ。「$\pi_3$ の生成元」と「写像度1」が一致しないのはなぜか。\\
(3) 随伴表現の生成子 $(T^a_{\mathrm{adj}})_{bc}=-i\epsilon_{abc}$ でトレースをとった量 $w_{\mathrm{adj}}(f)=\frac{1}{24\pi^2}\int\tr_{\mathrm{adj}}(f^{-1}df)^3$ について、$w_{\mathrm{adj}}(p\circ g)=4\,w(g)$ を示せ。（ヒント：表現 $R$ について $\tr_R(T^a[T^b,T^c])=if^{bca}T(R)$、$\tr_R(T^aT^b)=T(R)\delta^{ab}$。）
\end{problem}

\begin{solution}
(1) 長完全列 $\pi_n(\Z_2)\to\pi_n(SU(2))\xrightarrow{p_*}\pi_n(SO(3))\to\pi_{n-1}(\Z_2)$ で、離散空間 $\Z_2$ の $\pi_k$（$k\ge1$）は自明（連結な $S^k$ からの基点付き写像は定値）。$n\ge2$ なら両端が $0$ で $p_*$ は同型。$\pi_3(SU(2))$ の生成元は $\id$ なので $p_*[\id]=[p]$ が $\pi_3(SO(3))$ の生成元。持ち上げ：$S^3$ は単連結なので、被覆空間の持ち上げ定理（第2章）により $f\colon S^3\to SO(3)$ は $g\colon S^3\to SU(2)$ に持ち上がり $f=p\circ g$。

(2) $SO(3)$ の任意の点 $R$ の逆像は $\{U,-U\}$ の2点。$p$ は局所微分同相で、$p(-U)=p(U)$ より $p=p\circ(-\mathbf 1)$ であり、$U\mapsto-U$ は $\R^4$ の $-\mathbf 1_4$（$\det=+1$）の制限なので $S^3$ の向きを保つ。よって2点での局所的な符号は等しく、$SO(3)$ の向きを $p$ が向きを保つように選べば $\deg p=1+1=2$。$\pi_3$ の同型は群の同型であるが、写像度は $H_3$ への写像（フレヴィッツ写像 $\pi_3(SO(3))\to H_3(SO(3))\cong\Z$）で測った量であり、$SO(3)$ が単連結でないためこれが同型ではなく2倍写像になる。単連結な $S^3$ では両者は一致する。

(3) 1形式 $\alpha=g^{-1}dg=\alpha^aT^a$（$T^a=\sigma^a/2$）は、随伴表現では $p\circ g$ について $(p\circ g)^{-1}d(p\circ g)=\alpha^aT^a_{\mathrm{adj}}$ となる（$p$ は準同型で、微分は随伴表現の生成子への写像）。任意の表現で、$\alpha^a\wedge\alpha^b\wedge\alpha^c$ は $a,b,c$ について完全反対称なので
\[
\tr_R(\alpha^3)=\alpha^a\alpha^b\alpha^c\tr_R(T^aT^bT^c)=\tfrac12\alpha^a\alpha^b\alpha^c\tr_R(T^a[T^b,T^c])=\tfrac{i}{2}f^{abc}T(R)\,\alpha^a\alpha^b\alpha^c .
\]
（$[T^b,T^c]=if^{bcd}T^d$ と $f^{bca}=f^{abc}$ を用いた。）よって $\tr_R\alpha^3$ は $T(R)$ に比例する。$SU(2)$ では $f^{abc}=\epsilon^{abc}$、基本表現 $T(\text{fund})=\frac12$、随伴表現 $\tr(T^a_{\mathrm{adj}}T^b_{\mathrm{adj}})=\epsilon_{acd}\epsilon_{bcd}=2\delta^{ab}$ より $T(\mathrm{adj})=2$。比は $4$ なので $w_{\mathrm{adj}}(p\circ g)=4w(g)$。\\
まとめると、同じ写像 $p$ について「$\pi_3$ での番号」は1、「多様体間の写像度」は2、「随伴表現のトレースで規格化した巻き付き数」は4である。$SO(3)$ ゲージ理論のインスタントン数を随伴表現で規格化するときには、この因子に注意しなければならない。
\end{solution}

\subsection{2次元のソリトン}

\begin{problem}[難易度 B]\label{ch03-prob-bogo}
2次元空間の $O(3)$ 非線形シグマ模型 $\mathcal L=-\frac12\partial_\mu\bm n\cdot\partial^\mu\bm n$（$\lvert\bm n\rvert=1$）の静的エネルギーは $E=\frac12\int d^2x\,(\partial_i\bm n)\cdot(\partial_i\bm n)$ である。遠方で $\bm n$ は一定値に近づくとし、$Q$ を (\ref{eq:ch03-Q}) で定める。\\
(1) 恒等式
\[
\bigl(\partial_i\bm n\pm\epsilon_{ij}\,\bm n\times\partial_j\bm n\bigr)\cdot\bigl(\partial_i\bm n\pm\epsilon_{ik}\,\bm n\times\partial_k\bm n\bigr)=2(\partial_i\bm n)^2\mp4\,\bm n\cdot(\partial_1\bm n\times\partial_2\bm n)
\]
を示し、\term{ボゴモルニ不等式}{Bogomol'nyi bound} $E\ge4\pi\lvert Q\rvert$ を導け。\\
(2) $Q>0$ で等号が成り立つための条件（BPS 方程式）を書け。\\
(3) スケール変換 $\bm n(\bm x)\mapsto\bm n(\bm x/\lambda)$ で $E$ と $Q$ はどう変わるか。
\end{problem}

\begin{solution}
(1) $\lvert\bm n\rvert=1$ より $\bm n\cdot\partial_j\bm n=0$。\\
2乗の項：$(\epsilon_{ij}\bm n\times\partial_j\bm n)\cdot(\epsilon_{ik}\bm n\times\partial_k\bm n)=\delta_{jk}(\bm n\times\partial_j\bm n)\cdot(\bm n\times\partial_k\bm n)=\sum_j\lvert\bm n\times\partial_j\bm n\rvert^2=\sum_j\lvert\partial_j\bm n\rvert^2$（$\bm n\perp\partial_j\bm n$、$\lvert\bm n\rvert=1$）。ここで2次元の $\epsilon_{ij}\epsilon_{ik}=\delta_{jk}$ を用いた。\\
交差項：$2\epsilon_{ij}\,\partial_i\bm n\cdot(\bm n\times\partial_j\bm n)=2\epsilon_{ij}\,\bm n\cdot(\partial_j\bm n\times\partial_i\bm n)=-2\epsilon_{ij}\,\bm n\cdot(\partial_i\bm n\times\partial_j\bm n)=-4\,\bm n\cdot(\partial_1\bm n\times\partial_2\bm n)$。\\
よって左辺は $(\partial\bm n)^2+(\partial\bm n)^2\mp4q$（$q=\bm n\cdot(\partial_1\bm n\times\partial_2\bm n)$）。これが非負なので $(\partial_i\bm n)^2\ge\pm2q$。積分すると
\[
E=\frac12\int(\partial_i\bm n)^2\ge\pm\int q\,d^2x=\pm4\pi Q .
\]
両方の符号で成り立つので $E\ge4\pi\lvert Q\rvert$。

(2) $Q>0$ では上の符号で等号：$\partial_i\bm n=-\epsilon_{ij}\,\bm n\times\partial_j\bm n$、すなわち
\[
\partial_1\bm n=-\bm n\times\partial_2\bm n,\qquad\partial_2\bm n=\bm n\times\partial_1\bm n .
\]
（第2式は第1式に $\bm n\times$ を作用させたものと同値。）$Q<0$ では符号が逆。この1階の方程式の解は自動的に2階の運動方程式をみたし、各位相セクターでエネルギー最小である。

(3) $\bm n_\lambda(\bm x)=\bm n(\bm x/\lambda)$ とすると $\partial_i\bm n_\lambda=\lambda^{-1}(\partial_i\bm n)(\bm x/\lambda)$、$d^2x=\lambda^2d^2y$ なので $E[\bm n_\lambda]=E[\bm n]$。$Q$ も不変（写像度だから）。2次元の $O(3)$ 模型はスケール不変であり、デリックの議論（問題 \ref{ch03-prob-skyrmion}）による崩壊は起こらないが、解の大きさは任意（零モード）になる。なお同じ式は2次元ユークリッドの作用 $S=\frac12\int d^2x(\partial_\mu\bm n)^2$ にも当てはまり、この場合 BPS 解は2次元の\textbf{インスタントン}である。
\end{solution}

\begin{problem}[難易度 C]\label{ch03-prob-lump}
問題 \ref{ch03-prob-bogo} の模型を、南極からの立体射影座標
\[
w=\frac{n^1+in^2}{1+n^3},\qquad\bm n=\frac{\bigl(2\,\mathrm{Re}\,w,\ 2\,\mathrm{Im}\,w,\ 1-\lvert w\rvert^2\bigr)}{1+\lvert w\rvert^2}
\]
で書く。$z=x^1+ix^2$、$\partial_z=\frac12(\partial_1-i\partial_2)$、$\partial_{\bar z}=\frac12(\partial_1+i\partial_2)$ とする。\\
(1) $w$ は $S^2$ の向きを保つ座標であり、球面の計量が $ds^2=\dfrac{4\lvert dw\rvert^2}{(1+\lvert w\rvert^2)^2}$ であることを認めて、
\[
(\partial_i\bm n)^2=\frac{8\bigl(\lvert\partial_zw\rvert^2+\lvert\partial_{\bar z}w\rvert^2\bigr)}{(1+\lvert w\rvert^2)^2},\qquad q=\bm n\cdot(\partial_1\bm n\times\partial_2\bm n)=\frac{4\bigl(\lvert\partial_zw\rvert^2-\lvert\partial_{\bar z}w\rvert^2\bigr)}{(1+\lvert w\rvert^2)^2}
\]
を示し、$E=4\pi Q+8\int d^2x\,\lvert\partial_{\bar z}w\rvert^2/(1+\lvert w\rvert^2)^2$ を導け。$Q>0$ の BPS 解は正則関数 $w(z)$ で与えられる。\\
(2) $w(z)=(z-a)/\lambda$（$a\in\C$、$\lambda\in\C^\times$）のエネルギー密度を求め、$E=4\pi$、$Q=1$ を確かめよ。この解（\term{ランプ}{lump}）の4つの実パラメータの意味を述べよ。\\
(3) 有理関数 $w=P(z)/R(z)$（$P,R$ は共通根のない多項式、$\deg P>\deg R$）の位相的電荷が $Q=\deg P$ であることを示せ。
\end{problem}

\begin{solution}
(1) $w(\bm n)$ が向きを保つことは、$w=z$ に対応する $\bm n$ が原点で $\partial_1\bm n=(2,0,0)$、$\partial_2\bm n=(0,2,0)$、$\bm n=(0,0,1)$ をみたし $q=4>0$ となることから確かめられる。\\
$(\partial_i\bm n)^2$ は球面の計量の引き戻しのトレースなので
\[
(\partial_i\bm n)^2=\frac{4(\lvert\partial_1w\rvert^2+\lvert\partial_2w\rvert^2)}{(1+\lvert w\rvert^2)^2}.
\]
$\partial_1=\partial_z+\partial_{\bar z}$、$\partial_2=i(\partial_z-\partial_{\bar z})$ より $\lvert\partial_1w\rvert^2+\lvert\partial_2w\rvert^2=\lvert a+b\rvert^2+\lvert a-b\rvert^2=2(\lvert a\rvert^2+\lvert b\rvert^2)$（$a=\partial_zw$、$b=\partial_{\bar z}w$）。これで第1式を得る。\\
$q\,dx^1\wedge dx^2$ は面積形式の引き戻しである。座標 $w=u+iv$ での面積形式は（計量から）$\frac{4}{(1+\lvert w\rvert^2)^2}du\wedge dv$（向きを保つ座標なので正符号）。$du\wedge dv=\frac{i}{2}dw\wedge d\bar w$ で、$dw=a\,dz+b\,d\bar z$、$d\bar w=\bar b\,dz+\bar a\,d\bar z$ より $dw\wedge d\bar w=(\lvert a\rvert^2-\lvert b\rvert^2)dz\wedge d\bar z$。$\frac i2dz\wedge d\bar z=dx^1\wedge dx^2$ なので $du\wedge dv=(\lvert a\rvert^2-\lvert b\rvert^2)dx^1\wedge dx^2$。これで第2式を得る。\\
差をとると $(\partial_i\bm n)^2-2q=\frac{16\lvert\partial_{\bar z}w\rvert^2}{(1+\lvert w\rvert^2)^2}$ なので
\[
E=\frac12\int(\partial_i\bm n)^2=\int q+8\int\frac{\lvert\partial_{\bar z}w\rvert^2}{(1+\lvert w\rvert^2)^2}=4\pi Q+8\int\frac{\lvert\partial_{\bar z}w\rvert^2}{(1+\lvert w\rvert^2)^2}\ge4\pi Q .
\]
等号は $\partial_{\bar z}w=0$、すなわち $w$ が正則（コーシー・リーマン方程式）のとき。これは問題 \ref{ch03-prob-bogo} の BPS 方程式を複素座標で書いたものである。

(2) $\partial_zw=1/\lambda$、$\partial_{\bar z}w=0$。$\rho=\lvert z-a\rvert$ として
\[
\mathcal E=\frac12(\partial_i\bm n)^2=\frac{4/\lvert\lambda\rvert^2}{(1+\rho^2/\lvert\lambda\rvert^2)^2}=\frac{4\lvert\lambda\rvert^2}{(\lvert\lambda\rvert^2+\rho^2)^2},
\]
\[
E=\int_0^\infty2\pi\rho\,d\rho\,\frac{4\lvert\lambda\rvert^2}{(\lvert\lambda\rvert^2+\rho^2)^2}=4\pi\lvert\lambda\rvert^2\Bigl[-\frac{1}{\lvert\lambda\rvert^2+\rho^2}\Bigr]_0^\infty=4\pi .
\]
BPS なので $Q=E/4\pi=1$。（あるいは $w$ は $\C\cup\{\infty\}$ の1対1正則写像で、向きを保つので写像度1。）パラメータ：$a$（2実数）はランプの位置、$\lvert\lambda\rvert$ は大きさ、$\arg\lambda$ は内部空間の $n^3$ 軸まわりの回転角。$\bm n$ は $z=a$ で北極 $(0,0,1)$、遠方で南極 $(0,0,-1)$ をとる。大きさ $\lvert\lambda\rvert$ が任意であることはスケール不変性（問題 \ref{ch03-prob-bogo} (3)）の現れである。

(3) $\deg P=N>\deg R$ なので $w(\infty)=\infty$ で、$w$ は $\C\cup\{\infty\}=S^2$ からそれ自身への正則写像に拡張される。一般の値 $c$ に対し $w(z)=c$ は $P(z)-cR(z)=0$ という $N$ 次方程式で、一般の $c$ では相異なる $N$ 個の解をもつ。正則写像は各点で $\lvert\partial_zw\rvert^2>0$（臨界点以外）であり向きを保つので、どの逆像も $+1$ を寄与し $Q=\deg w=N$。（$w$ の座標の向きは (1) で確かめた。）したがって正則な有理関数は任意の $Q>0$ の BPS 解を与え、それらは $Q$ 個のランプの重ね合わせとみなせる。
\end{solution}

\begin{problem}[難易度 C]\label{ch03-prob-cpn}
$\CP^{N-1}$ 模型は $z\in\C^N$、$z^\dagger z=1$ と補助的な $U(1)$ ゲージ場で書かれる。$A=-i\,z^\dagger dz$、$D_\mu z=\partial_\mu z-iA_\mu z$、2次元ユークリッドの作用を $S=\int d^2x\,(D_\mu z)^\dagger(D_\mu z)$ とする。\\
(1) ファイバー束 $S^1\to S^{2N-1}\to\CP^{N-1}$ の長完全列から $\pi_1(\CP^{N-1})=0$、$\pi_2(\CP^{N-1})\cong\Z$ を示せ。\\
(2) $A$ が実の1形式であること、$z\to e^{i\lambda(x)}z$ のもとで $A\to A+d\lambda$、$D_\mu z\to e^{i\lambda}D_\mu z$ と変換することを示せ。\\
(3) $(D_1z)^\dagger D_2z-(D_2z)^\dagger D_1z=iF_{12}$（$F=dA$）を示し、$S\ge2\pi\lvert Q\rvert$、$Q=\frac1{2\pi}\int F$ を導け。等号条件は何か。\\
(4) $N=2$、$z=(1,x^1+ix^2)^{\mathsf T}/\sqrt{1+\lvert x\rvert^2}$ について $A=\frac{r^2}{1+r^2}d\varphi$、$Q=1$ を示し、$Q$ が遠方の円周上の $U(1)$ の巻き数 $\frac1{2\pi}\oint A$ に等しいことを、(1) の境界写像 $\partial\colon\pi_2(\CP^{N-1})\to\pi_1(S^1)$ と関連づけて説明せよ。
\end{problem}

\begin{solution}
(1) 完全列 $\pi_2(S^{2N-1})\to\pi_2(\CP^{N-1})\xrightarrow{\partial}\pi_1(S^1)\to\pi_1(S^{2N-1})\to\pi_1(\CP^{N-1})\to\pi_0(S^1)$。$N\ge2$ では $2N-1\ge3$ なので $\pi_1(S^{2N-1})=\pi_2(S^{2N-1})=0$、$\pi_0(S^1)$ は1点。よって $\pi_2(\CP^{N-1})\cong\pi_1(S^1)=\Z$、$\pi_1(\CP^{N-1})=0$。

(2) $z^\dagger z=1$ を微分すると $dz^\dagger z+z^\dagger dz=0$、すなわち $z^\dagger dz=-\overline{z^\dagger dz}$ は純虚数で、$A=-iz^\dagger dz$ は実。$z\to e^{i\lambda}z$ では $z^\dagger dz\to z^\dagger(dz+i\,d\lambda\,z)=z^\dagger dz+i\,d\lambda$ なので $A\to A+d\lambda$。$Dz\to e^{i\lambda}(dz+i\,d\lambda\,z)-i(A+d\lambda)e^{i\lambda}z=e^{i\lambda}Dz$。本書の規約 $D=d-iA$、$A\to A+d\lambda$（電荷1）と一致する。

(3) $z^\dagger\partial_iz=iA_i$、$(\partial_iz)^\dagger z=-iA_i$ を用いると
\[
(D_1z)^\dagger D_2z=(\partial_1z)^\dagger\partial_2z-iA_2(\partial_1z)^\dagger z+iA_1z^\dagger\partial_2z+A_1A_2=(\partial_1z)^\dagger\partial_2z-A_1A_2 .
\]
$1\leftrightarrow2$ との差をとると $A_1A_2$ が消え、$(D_1z)^\dagger D_2z-(D_2z)^\dagger D_1z=(\partial_1z)^\dagger\partial_2z-(\partial_2z)^\dagger\partial_1z$。一方 $F_{12}=\partial_1A_2-\partial_2A_1=-i\bigl(\partial_1(z^\dagger\partial_2z)-\partial_2(z^\dagger\partial_1z)\bigr)=-i\bigl((\partial_1z)^\dagger\partial_2z-(\partial_2z)^\dagger\partial_1z\bigr)$。よって差は $iF_{12}$。\\
次に
\[
\lvert D_1z\pm iD_2z\rvert^2=\lvert D_1z\rvert^2+\lvert D_2z\rvert^2\pm i\bigl((D_1z)^\dagger D_2z-(D_2z)^\dagger D_1z\bigr)=\lvert D_\mu z\rvert^2\mp F_{12}\ge0 .
\]
積分して $S\ge\pm\int F_{12}d^2x=\pm2\pi Q$、よって $S\ge2\pi\lvert Q\rvert$。$Q>0$ の等号条件は $D_1z+iD_2z=0$。実際、$D_iz=(\mathbf 1-zz^\dagger)\partial_iz$ と書け、$z=u/\lvert u\rvert$（$u$ は $x^1+ix^2$ の正則関数を成分とするベクトル）とすると $\partial_iz$ の $z$ 方向の成分は射影で消えるので $(D_1+iD_2)z=(\mathbf 1-zz^\dagger)(\partial_1+i\partial_2)u/\lvert u\rvert=0$。正則な $u$ が BPS インスタントンを与える。

(4) $w=x^1+ix^2=re^{i\varphi}$、$\rho=1+r^2$、$z=(1,w)^{\mathsf T}\rho^{-1/2}$ とすると
\[
z^\dagger dz=\frac{\bar w\,dw}{\rho}+\rho^{1/2}d(\rho^{-1/2})=\frac{\bar w\,dw}{\rho}-\frac{d\rho}{2\rho}=\frac{\bar w\,dw-w\,d\bar w}{2\rho}
\]
（$d\rho=\bar w\,dw+w\,d\bar w$）。$\bar w\,dw-w\,d\bar w=2ir^2d\varphi$ なので $A=\frac{r^2}{1+r^2}d\varphi$。$F=dA=\frac{2r}{(1+r^2)^2}dr\wedge d\varphi$ で、$r\,dr\wedge d\varphi=dx^1\wedge dx^2$ だから
\[
Q=\frac{1}{2\pi}\int F=\int_0^\infty\frac{2r\,dr}{(1+r^2)^2}=1 .
\]
（$n^a=z^\dagger\sigma^az$ とすると $n^1+in^2=2\bar z_1z_2=2w/(1+r^2)$、$n^3=(1-r^2)/(1+r^2)$ で、これは問題 \ref{ch03-prob-lump} の $w=z$ の解そのもの。$\CP^1$ 模型と $O(3)$ 模型は同じであり、$Q$ は $\bm n$ の写像度に一致する。）\\
ストークスの定理より $\int_{D_R}F=\oint_{\lvert x\rvert=R}A\to\oint d\varphi=2\pi$。遠方で $Dz\to0$ なので $A$ は純ゲージで、$z\to e^{i\varphi}z_\infty$（$z_\infty$ は定ベクトル）の形になる。つまり遠方の円周上で $z$ は1点 $[z_\infty]\in\CP^1$ のファイバー $S^1$ を1周している。これは境界写像 $\partial$ の構成そのものである：$\R^2\cup\{\infty\}\to\CP^{N-1}$ を $S^{2N-1}$ に持ち上げると（$z$ がその持ち上げ）、遠方の円周上でファイバーに沿ったループが残り、その巻き数 $\frac1{2\pi}\oint A=Q$ が $\partial[\,\cdot\,]\in\pi_1(S^1)=\Z$ である。
\end{solution}

\subsection{$\pi_3$ と巻き付き数：インスタントンとスカーミオン}

\begin{problem}[難易度 A]\label{ch03-prob-wid}
$SU(2)$ の点を $g=a_0+i\bm a\cdot\bm\sigma$（$a_0^2+\lvert\bm a\rvert^2=1$）と書く。\\
(1) 単位元の近くで $g^{-1}dg=i\sigma^k\,da^k$ であることから、単位元で $\tr(g^{-1}dg)^3=12\,da^1\wedge da^2\wedge da^3$ を示せ。\\
(2) $\tr(g^{-1}dg)^3$ が左移動 $g\mapsto hg$ で不変であることを用いて、$\tr(g^{-1}dg)^3=12\,\mathrm{vol}_{S^3}$ を示せ（$\mathrm{vol}_{S^3}$ は単位 $S^3$ の体積形式で、単位元での向きを $da^1\wedge da^2\wedge da^3$ とする）。$\mathrm{Vol}(S^3)=2\pi^2$ を用いて、恒等写像 $S^3\to SU(2)$ の巻き付き数が $w=1$ であることを確かめよ。\\
(3) $g\mapsto g^{-1}$ の巻き付き数を求めよ。
\end{problem}

\begin{solution}
(1) 単位元 $a_0=1$, $\bm a=0$ では $da_0=0$（$a_0=\sqrt{1-\lvert\bm a\rvert^2}$ の微分は $\bm a=0$ で消える）なので $dg=i\sigma^kda^k$、$g^{-1}=\mathbf 1$。
\[
\tr(g^{-1}dg)^3=i^3\tr(\sigma^i\sigma^j\sigma^k)\,da^i\wedge da^j\wedge da^k=-i\cdot2i\,\epsilon_{ijk}\,da^i\wedge da^j\wedge da^k=2\cdot6\,da^1\wedge da^2\wedge da^3 .
\]
ここで $\tr(\sigma^i\sigma^j\sigma^k)=2i\epsilon_{ijk}$（反対称部分のみ寄与）と $\epsilon_{ijk}da^i\wedge da^j\wedge da^k=6\,da^1\wedge da^2\wedge da^3$ を用いた。

(2) $(hg)^{-1}d(hg)=g^{-1}h^{-1}h\,dg=g^{-1}dg$ なので $\tr(g^{-1}dg)^3$ は左不変。同様に $\tr$ の巡回性から右移動 $g\mapsto gh$ でも不変（$h^{-1}(g^{-1}dg)h$ になるだけ）。$S^3$ の回転群 $SO(4)$ の単位成分は左右の移動 $g\mapsto h_1gh_2^{-1}$ で尽くされ、単位元の体積形式 $da^1\wedge da^2\wedge da^3$ を各点に運んだものが $\mathrm{vol}_{S^3}$ なので、$\tr(g^{-1}dg)^3=12\,\mathrm{vol}_{S^3}$。よって
\[
w(\id)=\frac{1}{24\pi^2}\cdot12\cdot2\pi^2=1 .
\]
本書の規約（恒等写像で $w=1$）は、$S^3$ の向きを単位元で $da^1\wedge da^2\wedge da^3$ ととることに相当する。

(3) $\iota(g)=g^{-1}$ では $\iota^{-1}d\iota=g\,d(g^{-1})=-dg\,g^{-1}$。$\tr(-dg\,g^{-1})^3=-\tr(g\,(g^{-1}dg)^3g^{-1})=-\tr(g^{-1}dg)^3$ なので $w(\iota)=-1$。（$g^{-1}=a_0-i\bm a\cdot\bm\sigma$ は $\bm a\to-\bm a$ という $\R^3$ の向き反転であることからもわかる。）一般に $w(g^{-1})=-w(g)$（各点ごとの逆元）。
\end{solution}

\begin{problem}[難易度 B]\label{ch03-prob-wadd}
$g,h\colon S^3\to SU(N)$ とし、各点ごとの積 $(gh)(x)=g(x)h(x)$ を考える。\\
(1) $\alpha=g^{-1}dg$、$\beta=dh\,h^{-1}$ とおくと $(gh)^{-1}d(gh)=h^{-1}(\alpha+\beta)h$、$d\alpha=-\alpha\wedge\alpha$、$d\beta=\beta\wedge\beta$ であることを示せ。\\
(2) $\tr(\alpha+\beta)^3=\tr\alpha^3+\tr\beta^3-3\,d\tr(\alpha\wedge\beta)$ を示し、$w(gh)=w(g)+w(h)$ を導け。\\
(3) 各点ごとの $n$ 乗 $g^n$ について $w(g^n)=n\,w(g)$ を示し、任意の整数 $n$ に対し $w=n$ の写像 $S^3\to SU(2)$ を構成せよ。
\end{problem}

\begin{solution}
(1) $(gh)^{-1}d(gh)=h^{-1}g^{-1}(dg\,h+g\,dh)=h^{-1}\alpha h+h^{-1}dh=h^{-1}(\alpha+dh\,h^{-1})h$。$d(g^{-1})=-g^{-1}dg\,g^{-1}$ より $d\alpha=d(g^{-1})\wedge dg=-g^{-1}dg\wedge g^{-1}dg=-\alpha\wedge\alpha$。$d\beta=-dh\wedge d(h^{-1})=dh\wedge h^{-1}dh\,h^{-1}=\beta\wedge\beta$。

(2) 行列値の $p$ 形式 $\mu$ と $q$ 形式 $\nu$ について $\tr(\mu\wedge\nu)=(-1)^{pq}\tr(\nu\wedge\mu)$。これより $\tr(\alpha\beta\alpha)=\tr(\beta\alpha\alpha)=\tr(\alpha\alpha\beta)$（1形式を2形式の反対側へ移すと符号 $(-1)^2=1$）、同様に $\beta$ が2つの項も等しい。よって
\[
\tr(\alpha+\beta)^3=\tr\alpha^3+3\tr(\alpha^2\beta)+3\tr(\alpha\beta^2)+\tr\beta^3 .
\]
一方 $d\tr(\alpha\beta)=\tr(d\alpha\,\beta)-\tr(\alpha\,d\beta)=-\tr(\alpha^2\beta)-\tr(\alpha\beta^2)$。したがって $\tr(\alpha+\beta)^3=\tr\alpha^3+\tr\beta^3-3d\tr(\alpha\beta)$。また $\tr\beta^3=\tr(h\,(h^{-1}dh)^3h^{-1})=\tr(h^{-1}dh)^3$、共役 $h^{-1}(\cdot)h$ はトレースを変えない。$S^3$ は閉多様体なので完全形式の積分は $0$ で、
\[
w(gh)=\frac{1}{24\pi^2}\int\bigl[\tr\alpha^3+\tr(h^{-1}dh)^3\bigr]=w(g)+w(h).
\]
これは $w$ が群準同型 $\pi_3(SU(N))\to\Z$ を与えることを示す（各点ごとの積は $\pi_3$ の和とホモトピックである）。

(3) 帰納法で $w(g^n)=w(g^{n-1})+w(g)=nw(g)$（$n\ge1$）。$w(g^0)=w(\text{定値})=0$、$w(g^{-n})=-w(g^n)$（問題 \ref{ch03-prob-wid} (3)）。恒等写像 $\id\colon S^3\to SU(2)$ をとれば、$x\mapsto x^n$（四元数の $n$ 乗）が $w=n$ の写像である。物理的には、巻き付き数 $n$ のゲージ変換やスカーミオン数 $n$ の配位の具体的な代表元を与える。
\end{solution}

\begin{problem}[難易度 D]\label{ch03-prob-skyrmion}
$SU(2)$ カイラル場 $U\colon\R^3\to SU(2)$、$U\to\mathbf 1$（$\lvert\bm x\rvert\to\infty$）を考える。$\R^3\cup\{\infty\}\cong S^3$（$\R^3$ の向きは $dx^1\wedge dx^2\wedge dx^3$）とみて、\term{バリオン数}{baryon number}を $B=-w(U)=-\frac{1}{24\pi^2}\int_{\R^3}\tr(U^{-1}dU)^3$ で定める。\\
(1) \term{ハリネズミ仮定}{hedgehog ansatz} $U=\exp\bigl(if(r)\hat{\bm x}\cdot\bm\sigma\bigr)=\cos f+i\sin f\,\hat{\bm x}\cdot\bm\sigma$（$f(0)\in\pi\Z$、$f(\infty)=0$）について $B=\bigl(f(0)-f(\infty)\bigr)/\pi$ を示せ。特に $f(0)=\pi$ なら $B=1$。（ヒント：$S^3$ の点を $\cos\chi+i\sin\chi\,\hat{\bm m}\cdot\bm\sigma$ と書くと、正規化された体積形式は $\frac{1}{2\pi^2}\sin^2\chi\,d\chi\wedge d\Omega_{\hat{\bm m}}$。ただし向きは $\chi\to0$ で $da^1\wedge da^2\wedge da^3$ と一致するようにとる。）\\
(2) バリオン数カレント $B^\mu\propto\epsilon^{\mu\nu\rho\sigma}\tr(L_\nu L_\rho L_\sigma)$（$L_\mu=U^{-1}\partial_\mu U$）が運動方程式を使わずに保存することを、微分形式 $d\tr(U^{-1}dU)^3=0$ として示せ。\\
(3) スキルム模型の静的エネルギー（スキルム単位）
\[
E=E_2+E_4,\qquad E_2=-\frac12\int d^3x\,\tr(L_iL_i),\qquad E_4=-\frac1{16}\int d^3x\,\tr\bigl([L_i,L_j]^2\bigr)
\]
について、$E_2,E_4\ge0$ を示し、スケール変換 $U_\lambda(\bm x)=U(\bm x/\lambda)$ のもとでの振る舞いから、静的解では $E_2=E_4$ であることを導け。$E_4$ がないと何が起こるか（\term{デリックの定理}{Derrick's theorem}）。
\end{problem}

\begin{solution}
(1) $U$ は写像 $\bm x\mapsto(\chi,\hat{\bm m})=(f(r),\hat{\bm x})$ である。$w(U)$ は正規化体積形式の引き戻しの積分なので
\[
w(U)=\frac{1}{2\pi^2}\int\sin^2f(r)\,f'(r)\,dr\wedge d\Omega_{\hat{\bm x}} .
\]
$\chi\to0$ では $a^k\approx\chi\,\hat m^k$ で、極座標 $(\chi,\hat{\bm m})$ に対する $d\chi\wedge d\Omega$ は $da^1\wedge da^2\wedge da^3$ と同じ向き（$\R^3$ の通常の極座標）なので、ヒントの体積形式の符号は正しい。$\R^3$ でも $dr\wedge d\Omega$ は $dx^1\wedge dx^2\wedge dx^3$ と同じ向きなので、
\[
w(U)=\frac{4\pi}{2\pi^2}\int_0^\infty\sin^2f\,f'\,dr=\frac2\pi\Bigl[\frac{f-\sin f\cos f}{2}\Bigr]_{r=0}^{r=\infty}=\frac{f(\infty)-f(0)}{\pi}
\]
（$f(0),f(\infty)\in\pi\Z$ で $\sin f\cos f$ の項は消える）。よって $B=-w(U)=(f(0)-f(\infty))/\pi$、$f(0)=\pi$、$f(\infty)=0$ で $B=1$。\\
符号の解釈：$r\to\infty$ では $U\approx\mathbf 1+if(r)\hat{\bm x}\cdot\bm\sigma$、$f$ は減少関数なので、写像 $\bm x\mapsto f(r)\hat{\bm x}$ は動径方向を反転し向きを逆にする（ヤコビアン $f'(f/r)^2<0$）。$f$ が単調なら任意の正則値の逆像は1点で符号は負なので、写像度（$=w$）は $-1$ である。$B=-w$ という符号の約束は、この標準的なスカーミオンが $B=+1$ となるように選んだものである（スキルム模型の文献の標準的な定義と一致する）。

(2) $L=U^{-1}dU$ は $dL=-L\wedge L$ をみたす（問題 \ref{ch03-prob-wadd}）。
\[
d\tr L^3=\tr(dL\,L^2)-\tr(L\,dL\,L)+\tr(L^2dL)=-3\tr L^4 ,
\]
（各項は巡回性で $-\tr L^4$ に等しい。）一方、奇数次形式 $L$ と3形式 $L^3$ について $\tr(L\wedge L^3)=(-1)^{1\cdot3}\tr(L^3\wedge L)=-\tr L^4$ なので $\tr L^4=0$。よって $d\tr L^3=0$。4次元で $B^\mu$ に対応する3形式 $\tr L^3$ が閉形式であることは $\partial_\mu B^\mu=0$ と同値であり、これは運動方程式によらない（\textbf{位相的カレント}）。$B$ が整数であることと合わせ、バリオン数は厳密に保存される（古典論の範囲で）。

(3) $L_i$ は反エルミート（$U^{-1}\partial_iU$ で $U$ はユニタリ）なので $L_i=iX_i$、$X_i$ エルミートとして $-\tr(L_iL_i)=\tr(X_iX_i)\ge0$。$[L_i,L_j]$ も反エルミートなので $-\tr([L_i,L_j]^2)\ge0$。よって $E_2,E_4\ge0$。\\
$U_\lambda(\bm x)=U(\bm x/\lambda)$ では $L_i\to\lambda^{-1}L_i$、$d^3x\to\lambda^3d^3y$ なので
\[
E(\lambda)=\lambda E_2+\lambda^{-1}E_4 .
\]
静的解はエネルギーの停留点なので、特にこの1パラメータ変形について $\frac{dE}{d\lambda}\big|_{\lambda=1}=E_2-E_4=0$。$E_4=0$（シグマ模型項だけ）だと $E(\lambda)=\lambda E_2$ は $\lambda\to0$ で減少し、配位は縮んで消えてしまう：位相的に $B\neq0$ でも有限サイズの静的解は存在しない（デリックの定理）。スキルム項 $E_4$ が縮みを妨げ、$E_2=E_4$ でつり合う。なお、スキルム模型では $E\ge12\pi^2\lvert B\rvert$（ファデーエフ・ボゴモルニ不等式）が成り立ち、ハリネズミの $B=1$ 解は数値的に $E\approx1.23\times12\pi^2$ である（ここでは証明しない）。
\end{solution}

\begin{problem}[難易度 B]\label{ch03-prob-theta}
4次元ユークリッド空間の $SU(N)$ ゲージ場で、遠方の3次元球面 $S^3_\infty$（$\R^4$ の外向き法線で向きづける）上で $A\to i\,g\,dg^{-1}$ となるものを考える。$CS(A)$、$k$ は記法と規約のとおりとする。\\
(1) 純ゲージ $A=i\,g\,dg^{-1}$ について $CS(A)=-\frac{1}{24\pi^2}\tr(g^{-1}dg)^3$ を示せ。\\
(2) $dCS=\frac{1}{8\pi^2}\tr F\wedge F$ とストークスの定理から $k=-w(g)=w(g^{-1})$ を導け。特に $k$ は整数である。\\
(3) 3次元空間 $\R^3\cup\{\infty\}\cong S^3$ 上の古典的真空 $A=i\,g\,dg^{-1}$（$g\to\mathbf 1$ at $\infty$）に $N_{\mathrm{CS}}=\int_{\R^3}CS(A)=-w(g)$ を対応させる。巻き付き数 $w(h)=-1$ のゲージ変換 $h$ で真空を移すと $N_{\mathrm{CS}}$ が1増えることを示せ。この操作の量子論での演算子を $T$ とし、$N_{\mathrm{CS}}=n$ の状態を $\ket{n}$ とすれば $T\ket n=\ket{n+1}$ である。\\
(4) $\ket\theta=\sum_ne^{in\theta}\ket n$ が $T$ の固有状態であることを示し、ハミルトニアンが $T$ と可換なら $\bra{\theta'}H\ket\theta=0$（$\theta\neq\theta'$）であることを示せ。
\end{problem}

\begin{solution}
(1) $F=dA-iA\wedge A=0$（純ゲージ）なので $dA=iA^2$。よって $\tr(A\,dA-\frac{2i}{3}A^3)=\tr(iA^3-\frac{2i}3A^3)=\frac i3\tr A^3$。$d(g^{-1})=-g^{-1}dg\,g^{-1}$ より $A=i\,g\,dg^{-1}=-i\,dg\,g^{-1}$ で、$A^3=(-i)^3(dg\,g^{-1})^3=i(dg\,g^{-1})^3$。$\tr(dg\,g^{-1})^3=\tr(g(g^{-1}dg)^3g^{-1})=\tr(g^{-1}dg)^3$ なので
\[
CS(A)=\frac{1}{8\pi^2}\cdot\frac i3\cdot i\,\tr(g^{-1}dg)^3=-\frac{1}{24\pi^2}\tr(g^{-1}dg)^3 .
\]

(2) $A$ は $\R^4$ 全体で滑らかなので $CS(A)$ は大域的な3形式で、半径 $R$ の球 $B_R$ でストークスの定理（境界は外向き法線で向きづける）を用いると
\[
k=\lim_{R\to\infty}\frac{1}{8\pi^2}\int_{B_R}\tr F\wedge F=\lim_{R\to\infty}\int_{S^3_R}CS(A)=-\frac{1}{24\pi^2}\int_{S^3_\infty}\tr(g^{-1}dg)^3=-w(g).
\]
$w(g^{-1})=-w(g)$（問題 \ref{ch03-prob-wid}）なので $k=w(g^{-1})$。$w$ は整数なので $k\in\Z$。つまりインスタントン数は、遠方の純ゲージを与える写像 $S^3_\infty\to G$ のホモトピー類（$\pi_3(G)=\Z$）で決まる。符号は規約 $A\to i\,g\,dg^{-1}$ と向きの選び方による。BPST 解の具体形は第8章で扱う。

(3) ゲージ変換 $h$ で $A\mapsto hAh^{-1}+i\,h\,dh^{-1}$。$A=i\,g\,dg^{-1}$ なら
\[
hAh^{-1}+i\,h\,dh^{-1}=i\,h g\,dg^{-1}h^{-1}+i\,h\,dh^{-1}=i\,(hg)\,d\bigl(g^{-1}h^{-1}\bigr)=i\,(hg)\,d(hg)^{-1}.
\]
よって $g\mapsto hg$ で、問題 \ref{ch03-prob-wadd} より $N_{\mathrm{CS}}=-w(hg)=-w(h)-w(g)=N_{\mathrm{CS}}+1$。（$w(h)=-1$ の例は問題 \ref{ch03-prob-skyrmion} のハリネズミ $h=e^{if(r)\hat{\bm x}\cdot\bm\sigma}$、$f(0)=\pi$。）古典的真空は $\pi_3(G)=\Z$ で番号づけられる無限個の位相的に異なる真空に分かれ、$T$ はそれらを1つずつずらす「大きなゲージ変換」である。インスタントンはこれらの間のトンネル効果を表し、$N_{\mathrm{CS}}$ を $\lvert k\rvert$ だけ変える（向きの規約を含めた詳細は第8章）。

(4) $T\ket\theta=\sum_ne^{in\theta}\ket{n+1}=\sum_me^{i(m-1)\theta}\ket m=e^{-i\theta}\ket\theta$。$T$ はユニタリで $[H,T]=0$ なら
\[
\bra{\theta'}H\ket\theta=\bra{\theta'}T^\dagger HT\ket\theta=e^{i\theta'}e^{-i\theta}\bra{\theta'}H\ket\theta ,
\]
よって $\theta\neq\theta'$（mod $2\pi$）なら $\bra{\theta'}H\ket\theta=0$。$\theta$ は時間発展で保存され、異なる $\theta$ の真空は別々の理論（超選択セクター）とみなせる。経路積分では $\theta$ 真空の間の振幅にインスタントン数 $k$ のセクターが位相 $e^{\pm ik\theta}$ で寄与し、これがラグランジアンの $\theta$ 項 $\frac{\theta}{8\pi^2}\tr F\wedge F$ と同等になる（符号の規約を含め第8章）。
\end{solution}

\subsection{ホップ写像と $\pi_4$}

\begin{problem}[難易度 D]\label{ch03-prob-hopf}
$S^3=\{(z_1,z_2)\in\C^2:\lvert z_1\rvert^2+\lvert z_2\rvert^2=1\}$、$z_1=x_1+ix_2$、$z_2=x_3+ix_4$ とし、$S^3$ には単位球体 $B^4\subset\R^4$（向き $dx_1\wedge dx_2\wedge dx_3\wedge dx_4$）の境界としての向きを入れる。ホップ写像を $h(z)^a=z^\dagger\sigma^az$ とする。\\
(1) $h^{-1}(0,0,1)$ と $h^{-1}(0,0,-1)$ を求め、点 $(0,0,0,1)$ からの立体射影 $S^3\setminus\{(0,0,0,1)\}\to\R^3$、$x\mapsto(x_1,x_2,x_3)/(1-x_4)$ による像が、それぞれ単位円周と $x_3$ 軸（と無限遠点）になり、互いに1回絡むことを示せ。\\
(2) $S^2$ の面積形式を $\omega$ とし、$\mathcal A=x_1dx_2-x_2dx_1+x_3dx_4-x_4dx_3$ とおく。$h^*\omega=2\,d\mathcal A$ が点 $(1,0,0,0)$ で成り立つことを確かめよ（$h$ と $\mathcal A$ の $SU(2)$ 不変性により全体で成り立つ）。\\
(3) $\bm n\colon S^3\to S^2$ に対し $F=\bm n^*\omega$ は $S^3$ 上の完全形式 $F=dA$ であり、\term{ホップ不変量}{Hopf invariant} $H(\bm n)=\frac{1}{16\pi^2}\int_{S^3}A\wedge F$ は $A$ の選び方によらないことを示せ。$H(h)=1$ を計算せよ。\\
(4) ホップ・ファイブレーションの長完全列から $\pi_3(S^2)\cong\Z$ を導き、3次元のファデーエフ・ニーミ模型（$\bm n\colon\R^3\to S^2$、遠方で一定）の位相的ソリトン（ホップ粒子）がホップ不変量で分類されることを説明せよ。
\end{problem}

\begin{center}
\begin{tikzpicture}[scale=1.0]
  % left panel: circle and axis
  \draw[thick,blue!70!black] (2,0) arc (0:180:2 and 0.6);
  \draw[thick,red!70!black,preaction={draw,white,line width=4pt}] (0,-2.3) -- (0,2.3);
  \draw[thick,blue!70!black,preaction={draw,white,line width=4pt}] (-2,0) arc (180:360:2 and 0.6);
  \draw[-{Stealth[length=2mm]},red!70!black,thick] (0,1.9) -- (0,2.35);
  \node[blue!70!black] at (2.35,-0.5) {\footnotesize $h^{-1}(N)$};
  \node[red!70!black,right] at (0.05,2.2) {\footnotesize $h^{-1}(S)$（$x_3$ 軸 $\cup\{\infty\}$）};
  % right panel: Hopf link
  \begin{scope}[xshift=6.6cm]
    \draw[very thick,blue!70!black] (0,0) circle (1);
    \draw[very thick,green!50!black] (1.2,0) circle (1);
    \draw[very thick,blue!70!black,preaction={draw,white,line width=4pt}] (40:1) arc (40:66:1);
    \node at (0.6,-1.55) {\footnotesize 任意の2本のファイバー（絡み数 1）};
  \end{scope}
\end{tikzpicture}
\zucap{ホップ・ファイブレーションのファイバーを立体射影で $\R^3$ に写したもの。左：北極と南極の上のファイバーは単位円周と $x_3$ 軸になり、1回絡んでいる。右：任意の2点の上のファイバーも同様に1回絡む（ホップ・リンク）。}
\end{center}

\begin{solution}
(1) $h(z)=(2\,\mathrm{Re}(\bar z_1z_2),2\,\mathrm{Im}(\bar z_1z_2),\lvert z_1\rvert^2-\lvert z_2\rvert^2)$。$h=(0,0,1)$ は $\lvert z_1\rvert=1$、$z_2=0$、すなわち $\{(e^{i\alpha},0)\}$：$x_3=x_4=0$、$x_1^2+x_2^2=1$。立体射影の像は $(x_1,x_2,0)$ で、$x_1x_2$ 平面の単位円周。$h=(0,0,-1)$ は $z_1=0$、$\lvert z_2\rvert=1$：$x_1=x_2=0$、$x_3^2+x_4^2=1$。像は $(0,0,x_3/(1-x_4))$ で、$x_4\to1$ のとき無限遠へ行き、$x_3/(1-x_4)$ は $\R$ 全体をとる。つまり $x_3$ 軸 $\cup\{\infty\}$。$x_3$ 軸は単位円周が張る円板をちょうど1回（原点で）貫くので、2つの閉曲線は1回絡む（絡み数 $\pm1$。符号は向きの約束による。絡み数は第6章で扱う）。

(2) 点 $p=(1,0,0,0)$、すなわち $z=(1,0)$ で $h(p)=(0,0,1)$。接空間は $\partial_{x_2},\partial_{x_3},\partial_{x_4}$ で張られる。$\partial_{x_2}$ はファイバー方向（$e^{i\alpha}z$ の微分 $(i,0)$）で $h$ は一定、よって $h_*\partial_{x_2}=0$。$z_1=1$、$z_2=x_3+ix_4$ の1次で $h\approx(2x_3,2x_4,1)$ なので $h_*\partial_{x_3}=(2,0,0)$、$h_*\partial_{x_4}=(0,2,0)$。北極での $\omega$ は $\omega(u,v)=\bm n\cdot(u\times v)$ なので
\[
h^*\omega(\partial_{x_3},\partial_{x_4})=(0,0,1)\cdot\bigl((2,0,0)\times(0,2,0)\bigr)=4,\quad h^*\omega(\partial_{x_2},\cdot)=0 .
\]
一方 $d\mathcal A=2(dx_1\wedge dx_2+dx_3\wedge dx_4)$ なので $2d\mathcal A(\partial_{x_3},\partial_{x_4})=4$、$2d\mathcal A(\partial_{x_2},\partial_{x_j})=4\,dx_1\wedge dx_2(\partial_{x_2},\partial_{x_j})=0$（$j=3,4$；$\partial_{x_1}$ 成分がないため）。よって $p$ で $h^*\omega=2d\mathcal A$。（$\mathcal A=-iz^\dagger dz$ を $S^3$ に制限したものであり、問題 \ref{ch03-prob-cpn} の $A$ に等しい。$z\mapsto Uz$, $U\in SU(2)$ で $\mathcal A$ は不変、$h$ は $SO(3)$ 回転を受けるだけなので $\omega$ の引き戻しも不変。$SU(2)$ は $S^3$ に推移的なので1点で確かめれば十分。）

(3) $F=\bm n^*\omega$ は閉形式（$S^2$ 上の2形式の引き戻しで、$d\omega=0$）。$S^3$ の2次ド・ラーム・コホモロジーは $0$（第5章）なので $F=dA$ と書ける。$A'=A+\lambda$（$d\lambda=0$）とすると、$H^1(S^3)=0$ より $\lambda=d\mu$ で、
\[
\int\lambda\wedge F=\int d\mu\wedge F=\int d(\mu F)=0 \qquad(dF=0)
\]
なので $H$ は $A$ によらない。$h$ については (2) より $A=2\mathcal A$、$F=2d\mathcal A$ とでき、
\[
H(h)=\frac{4}{16\pi^2}\int_{S^3}\mathcal A\wedge d\mathcal A=\frac{1}{4\pi^2}\int_{B^4}d\mathcal A\wedge d\mathcal A .
\]
$d\mathcal A\wedge d\mathcal A=4(dx_1dx_2+dx_3dx_4)^2=8\,dx_1\wedge dx_2\wedge dx_3\wedge dx_4$ で、$\mathrm{Vol}(B^4)=\pi^2/2$ より $\int_{B^4}=4\pi^2$。よって $H(h)=1$。\\
（解釈：$F/4\pi$ は「$S^2$ の点の逆像」の双対であり、$H$ は2つの点の逆像の絡み数に等しい。(1) と整合する。）

(4) $S^1\to S^3\to S^2$ の長完全列 $\pi_3(S^1)=0\to\pi_3(S^3)\xrightarrow{h_*}\pi_3(S^2)\to\pi_2(S^1)=0$ より $h_*$ は同型、$\pi_3(S^2)\cong\Z$ で生成元は $[h]$（$h_*[\id]=[h]$）。$H$ はホモトピー不変で（連続変形での変化は完全形式の積分になる）、$H(h)=1$ なので、$H$ はこの同型 $\pi_3(S^2)\to\Z$ を与える。\\
ファデーエフ・ニーミ模型では、遠方で一定の $\bm n\colon\R^3\to S^2$ は $S^3\to S^2$ を定め、そのホモトピー類 $\pi_3(S^2)=\Z$ がホップ電荷 $H$ である。$\pi_2$ による分類（点を囲む球面）とは異なり、これは空間全体の結び目構造であり、$\bm n^{-1}(\bm p_1)$ と $\bm n^{-1}(\bm p_2)$ という2本の閉曲線の絡み数が $H$ になる。ホップ粒子は閉じたひも状（結び目状）の配位であり、スキルム項に似た4階の項によりデリックの崩壊が防がれ、エネルギーは $E\ge c\lvert H\rvert^{3/4}$（ヴァクレンコ・カピタンスキー不等式）で下から押さえられる。
\end{solution}

\begin{problem}[難易度 C]\label{ch03-prob-pi4}
$\pi_4(S^3)=\Z_2$ を認める。\\
(1) $\pi_4(SU(2))$、$\pi_4(SO(3))$ を求めよ。\\
(2) $Sp(N)$ は $\HH^N$ の単位球面 $S^{4N-1}$ に推移的に作用し、固定部分群は $Sp(N-1)$ である（認めてよい）。$Sp(1)\cong SU(2)$ を用いて、$N\ge2$ で $\pi_4(Sp(N))\cong\Z_2$ を示せ。\\
(3) $\pi_4(SU(3))=0$ を認めて、$SU(2)\to SU(3)\to S^5$ の長完全列から、境界写像 $\partial\colon\pi_5(S^5)\to\pi_4(SU(2))$ が全射であることを示せ。\\
(4) 4次元時空を $S^4$ にコンパクト化すると、ゲージ変換 $g\colon S^4\to G$ は $\pi_4(G)$ で分類される。$SU(2)$ と $SU(3)$ ではどう違うか。第9章で扱うウィッテンの $SU(2)$ アノマリーとの関係を述べよ。
\end{problem}

\begin{solution}
(1) $SU(2)\cong S^3$ なので $\pi_4(SU(2))\cong\Z_2$。2重被覆 $SU(2)\to SO(3)$ は $n\ge2$ で $\pi_n$ の同型を誘導するので（問題 \ref{ch03-prob-so3}）$\pi_4(SO(3))\cong\Z_2$。

(2) ファイバー束 $Sp(N-1)\to Sp(N)\to S^{4N-1}$ の長完全列
\[
\pi_5(S^{4N-1})\to\pi_4(Sp(N-1))\to\pi_4(Sp(N))\to\pi_4(S^{4N-1}).
\]
$N\ge2$ では $4N-1\ge7$ なので $\pi_5(S^{4N-1})=\pi_4(S^{4N-1})=0$ で、$\pi_4(Sp(N-1))\cong\pi_4(Sp(N))$。帰納法で $\pi_4(Sp(N))\cong\pi_4(Sp(1))=\pi_4(S^3)=\Z_2$。

(3) 完全列
\[
\pi_5(S^5)\xrightarrow{\partial}\pi_4(SU(2))\xrightarrow{i_*}\pi_4(SU(3))\to\pi_4(S^5)
\]
で $\pi_4(SU(3))=0$ なので $i_*=0$、完全性より $\operatorname{im}\partial=\ker i_*=\pi_4(SU(2))$、すなわち $\partial\colon\Z\to\Z_2$ は全射（奇数を非自明な元に写す）。言い換えると、$SU(2)$ の非自明な写像 $S^4\to SU(2)$ は $SU(3)$ に埋め込むと自明になる。

(4) $SU(2)$ では $\pi_4=\Z_2$ なので、時空 $S^4$ 上に「連続変形で恒等変換に戻せない」ゲージ変換が存在する。$SU(3)$（一般に $SU(N)$、$N\ge3$）では $\pi_4=0$ でそのようなものはない。ウィッテンは、$SU(2)$ の2重項のワイル・フェルミオンが奇数個あると、フェルミオンの経路積分（の符号）がこの非自明なゲージ変換で $-1$ 倍になり、理論がゲージ不変でなくなることを示した（\term{ウィッテンの $SU(2)$ アノマリー}{Witten's SU(2) anomaly}、大域的アノマリー）。$\pi_4$ が自明でないことはこのアノマリーの存在の必要条件である。(3) の結果は、$SU(2)\subset SU(3)$ に埋め込むと $\Z_2$ の情報が摂動的な（三角形の）アノマリーに吸収されることと関係する（第9章）。同様に $Sp(N)$ のゲージ理論でも $\pi_4=\Z_2$ により大域的アノマリーが起こりうる。
\end{solution}

%================================================
\chapter{ホモロジー}\label{ch:04}
% 20問
%================================================

\section{物理からの動機}

\subsection*{境界の境界はゼロ}
3次元空間の領域 $V$ を考える。その境界 $\partial V$ は閉曲面であり、閉曲面には境界がない：
\[
\partial(\partial V)=\emptyset .
\]
この自明に見える事実が、電磁気学の基本法則の背後にある。ストークスの定理 $\oint_{\partial S}\bm A\cdot d\bm l=\int_S\bm B\cdot d\bm S$（$\bm B=\nabla\times\bm A$）を $S=\partial V$ に適用すると
\[
\int_{\partial V}\bm B\cdot d\bm S=\oint_{\partial\partial V}\bm A\cdot d\bm l=0
\]
となり、ベクトルポテンシャルで書ける磁場は湧き出しをもたない（$\nabla\cdot\bm B=0$）。微分形式の言葉では $d^2=0$、すなわちビアンキ恒等式 $dF=d(dA)=0$ である（第5章）。同様に、マクスウェル方程式 $\partial_\mu F^{\mu\nu}=J^\nu$ から $\partial_\nu J^\nu=\partial_\nu\partial_\mu F^{\mu\nu}=0$（電荷の保存）が従うのも、「2回微分すると消える」という同じ構造である。

本章では、この「$\partial\partial=0$」を代数的な出発点として、空間の\textbf{ホモロジー}を定義する。$\partial\partial=0$ なので境界はつねに「閉じたもの」（\term{サイクル}{cycle}）であるが、逆に閉じたものが何かの境界になっているとは限らない。境界にならないサイクルを数えるのがホモロジー群
\[
H_k=\frac{\{\text{$k$ 次元の閉じたもの}\}}{\{\text{$k+1$ 次元のものの境界}\}}
\]
である。たとえばトーラスの上で、トーラスを一周する輪はどんな面の境界にもならない。トーラス上の平坦なゲージ場のウィルソン・ループ $\exp(i\oint_\gamma A)$ は、輪 $\gamma$ を面の境界の分だけずらしても変わらないので、$\gamma$ のホモロジー類だけで決まる。独立なウィルソン・ループの数はホモロジーで数えられる。

\subsection*{ガウスの法則と全電荷}
閉じた空間（例えば周期境界条件の箱 $T^3$）ではガウスの法則 $\nabla\cdot\bm E=\rho$ を全空間で積分すると、表面項がないので全電荷が $0$ になる：
\[
Q=\int_{T^3}\rho\,d^3x=\int_{T^3}\nabla\cdot\bm E\,d^3x=0 .
\]
格子上でこれを書くと、電荷分布（各サイトの数）が電場（各リンクの数）の境界になっているという主張になり、$H_0$ が「各連結成分の全電荷」という保存量を与える（問題で扱う）。

\subsection*{格子ゲージ理論}
ウィルソンの格子ゲージ理論では、ゲージ場は格子の\textbf{リンク}（向き付けられた辺）上の群の元 $U_\ell$（平行移動）、場の強さは\textbf{プラケット}（最小の正方形）の縁をまわるホロノミー $U_p$ である。$U(1)$ の場合 $U_\ell=e^{i\theta_\ell}$ と書くと
\[
\theta_p=\theta_1(x)+\theta_2(x+\hat1)-\theta_1(x+\hat2)-\theta_2(x)
\]
はプラケット $p$ の\textbf{境界} $\partial p$ に沿った $\theta$ の和であり、連続極限で $\theta_p\approx a^2F_{12}$（$a$ は格子間隔）となる。サイト・リンク・プラケット・キューブはそれぞれ 0, 1, 2, 3 次元の胞体であり、格子ゲージ理論は\textbf{胞体複体の上の鎖と余鎖}の理論にほかならない。キューブの境界の境界がゼロであることが格子上のビアンキ恒等式になる。
\begin{center}
\begin{tikzpicture}[scale=1.25]
  \draw[gray!50] (-0.6,-0.6) grid (3.6,2.6);
  \fill[blue!10] (1,1) rectangle (2,2);
  \draw[very thick,red!70!black,mid arrow] (1,1) -- (2,1);
  \draw[very thick,red!70!black,mid arrow] (2,1) -- (2,2);
  \draw[very thick,red!70!black,mid arrow] (2,2) -- (1,2);
  \draw[very thick,red!70!black,mid arrow] (1,2) -- (1,1);
  \foreach \x in {0,...,3} \foreach \y in {0,...,2} \node[vtx,gray] at (\x,\y) {};
  \node[vtx] at (1,1) {}; \node[below left] at (1,1) {$x$};
  \node[below] at (1.5,0.95) {\footnotesize $\theta_1(x)$};
  \node[right] at (2.02,1.5) {\footnotesize $\theta_2(x+\hat1)$};
  \node[above] at (1.5,2.02) {\footnotesize $-\theta_1(x+\hat2)$};
  \node[left] at (0.98,1.5) {\footnotesize $-\theta_2(x)$};
  \node at (1.5,1.5) {$p$};
  \draw[-{Stealth[length=2mm]}] (4.1,0) -- (4.8,0) node[right] {$1$};
  \draw[-{Stealth[length=2mm]}] (4.1,0) -- (4.1,0.7) node[above] {$2$};
\end{tikzpicture}
\zucap{2次元正方格子のプラケット $p$ とその境界 $\partial p$（赤）。リンク変数 $\theta_\ell$ を縁に沿って（向きの符号つきで）足したものがプラケット変数 $\theta_p=\theta(\partial p)$ で、$p$ を貫く磁束 $a^2F_{12}$ にあたる。}
\end{center}

\subsection*{ファインマン図のループ数とオイラー標数}
ファインマン図はグラフであり、独立なループ運動量の数 $L=I-V+1$（$I$：内線の数、$V$：頂点の数）は、グラフの1次元ホモロジーの階数（第1ベッチ数）$b_1$ にほかならない。$\hbar$ 展開の次数はこの $L$ で決まる。

より一般に、空間を胞体に分けたとき
\[
\chi=(\text{頂点の数})-(\text{辺の数})+(\text{面の数})-\cdots
\]
は分け方によらない\term{オイラー標数}{Euler characteristic}になる。$\chi$ は曲率の積分（ガウス・ボンネの定理、第8章）、弦の世界面の種数展開 $g_s^{-\chi}$（第10章）、ラージ $N$ 展開の $N^{\chi}$ などに現れる。本章でこれらの計算の基礎を準備する。

\begin{physbox}[：本章の辞書]
\begin{center}
\small
\renewcommand{\arraystretch}{1.15}
\begin{tabular}{p{0.46\linewidth}p{0.40\linewidth}}
物理 & 数学\\ \hline
ビアンキ恒等式 $dF=0$、電荷保存 $\partial_\mu J^\mu=0$ & $\partial\partial=0$（双対的に $d^2=0$）\\
格子のサイト・リンク・プラケット・キューブ & 0, 1, 2, 3 胞体\\
電流の流線、ウィルソン・ループの経路 & 1 鎖、1 サイクル\\
平坦な場でのループの変形（ストークス） & ホモロジー類（境界の分の差は効かない）\\
閉じた空間の全電荷 $=0$ & $H_0$ と $\varepsilon\circ\partial=0$\\
独立なポリヤコフ・ループ、ループ運動量 & 第1ベッチ数 $b_1$\\
離散的なウィルソン・ライン（$W^p=1$） & ホモロジーのねじれ（torsion）\\
種数展開・ラージ $N$ 展開の次数 & オイラー標数 $\chi$\\
ヘッジホッグ・バビー・スキルミオンの電荷 & 写像度
\end{tabular}
\end{center}
\end{physbox}

\section{基本事項のまとめ}

\subsection{鎖複体とホモロジー}

\begin{mydef}{鎖複体・ホモロジー}{ch04-chain}
アーベル群 $C_k$（$k\in\Z$、$k<0$ では $0$）と準同型 $\partial_k\colon C_k\to C_{k-1}$ の列で $\partial_{k}\circ\partial_{k+1}=0$ をみたすものを\term{鎖複体}{chain complex}という。$\partial$ を\term{境界作用素}{boundary operator}と呼び、
\[
Z_k=\ker\partial_k,\qquad B_k=\operatorname{im}\partial_{k+1},\qquad H_k=Z_k/B_k
\]
とおき、$Z_k$ の元を\term{サイクル}{cycle}、$B_k$ の元を\term{バウンダリー}{boundary}、$H_k$ を $k$ 次の\term{ホモロジー群}{homology group}という。$\partial\partial=0$ より $B_k\subset Z_k$ である。$z,z'\in Z_k$ が $z-z'\in B_k$ のとき\term{ホモローグ}{homologous}であるという。
\end{mydef}

鎖複体の間の準同型の族 $f_k\colon C_k\to C'_k$ が $\partial'f_k=f_{k-1}\partial$ をみたすとき\term{鎖写像}{chain map}といい、$f_*[z]=[f(z)]$ によってホモロジーの準同型を誘導する。

本書で現れるホモロジー群はすべて\textbf{有限生成アーベル群}であり、次の構造定理により
\[
H_k\cong\Z^{b_k}\oplus\Z_{m_1}\oplus\cdots\oplus\Z_{m_r}
\]
と書ける（$\Z_m=\Z/m\Z$）。$b_k=\operatorname{rank}H_k$ を\term{ベッチ数}{Betti number}、有限位数の元のなす部分 $\Z_{m_1}\oplus\cdots$ を\term{ねじれ部分}{torsion}という。ねじれは「$z$ 自身は境界でないが $mz$ は境界である」サイクルを表す。係数を $\R$ や $\Q$ にとると（$C_k\otimes\R$ を使うと）ねじれは消え、$\dim H_k(\,\cdot\,;\R)=b_k$ となる（普遍係数定理、第6章）。

\begin{mydef}{完全列}{ch04-exact}
準同型の列 $\cdots\to A\xrightarrow{f}B\xrightarrow{g}C\to\cdots$ が $B$ で\term{完全}{exact}であるとは $\operatorname{im}f=\ker g$ であることをいう。$0\to A\to B\to C\to0$ が完全（\term{短完全列}{short exact sequence}）であることは、$A\to B$ が単射、$B\to C$ が全射、$B/A\cong C$ を意味する。$C$ が自由アーベル群なら短完全列は分裂し $B\cong A\oplus C$。有限生成アーベル群の短完全列では階数が加法的：$\operatorname{rank}B=\operatorname{rank}A+\operatorname{rank}C$。
\end{mydef}

\subsection{単体的ホモロジー}

\begin{mydef}{単体複体と向き}{ch04-simp}
$\R^N$ の中で一般の位置にある $k+1$ 点 $v_0,\dots,v_k$ の凸包を $k$ \term{単体}{simplex} という（0 単体＝点、1 単体＝線分、2 単体＝三角形、3 単体＝四面体）。その頂点の部分集合の凸包を\textbf{面}という。単体の集まりで、各単体の面を含み、2つの単体の共通部分が（空であるか）共通の面であるものを\term{単体複体}{simplicial complex} $K$ という。

頂点の並べ方を偶置換で移り合うものを同じとみなした同値類を単体の\textbf{向き}という。向きの付いた単体を $[v_0,\dots,v_k]$ と書き、
\[
[v_{\pi(0)},\dots,v_{\pi(k)}]=\operatorname{sgn}(\pi)\,[v_0,\dots,v_k]
\]
と約束する。$k$ \term{鎖}{chain} の群 $C_k(K)$ は、向き付けた $k$ 単体（各単体について一方の向きを選ぶ）を基底とする自由アーベル群である。
\end{mydef}

\begin{mydef}{単体的境界作用素}{ch04-bdry}
\[
\partial[v_0,\dots,v_k]=\sum_{i=0}^k(-1)^i[v_0,\dots,\widehat{v_i},\dots,v_k]
\]
（$\widehat{v_i}$ は $v_i$ を除くこと）。例：$\partial[v_0,v_1]=v_1-v_0$、$\partial[v_0,v_1,v_2]=[v_1,v_2]-[v_0,v_2]+[v_0,v_1]$。
\end{mydef}

\begin{mylem}{$\partial\partial=0$}{ch04-dd}
$\partial_{k-1}\circ\partial_k=0$。したがって $(C_*(K),\partial)$ は鎖複体であり、\term{単体的ホモロジー}{simplicial homology} $H_k(K)$ が定まる。
\end{mylem}
証明は問題で行う。三角形の境界 $\partial[v_0,v_1,v_2]$ は $v_0\to v_1\to v_2\to v_0$ と一周する閉じた道であり、その境界 $(v_1-v_0)+(v_2-v_1)+(v_0-v_2)$ は $0$ になる。

\begin{myex}{円周}{ch04-circle}
$K=\partial\Delta^2$（三角形の3辺と3頂点）では $C_1=\Z^3$、$C_0=\Z^3$、$C_2=0$。$z=[v_0,v_1]+[v_1,v_2]+[v_2,v_0]$ は $\partial z=0$ をみたし、$Z_1=\Z z$。境界になるものはないので $H_1(\partial\Delta^2)=\Z z\cong\Z$：円周を一周する輪は境界にならない。$\Delta^2$（内部の2単体も含む）では $z=\partial[v_0,v_1,v_2]$ となり $H_1=0$ である（問題）。
\end{myex}

\term{$\Delta$ 複体}{$\Delta$-complex} では、単体の辺どうしの同一視を許す（例：正方形を対角線で2つの三角形に分け、向かい合う辺を貼るとトーラス）。境界作用素は同じ式で定義され、少ない単体で計算できる。

\subsection{特異ホモロジーと基本性質}

任意の位相空間 $X$ に対しては、標準 $k$ 単体 $\Delta^k$ からの連続写像 $\sigma\colon\Delta^k\to X$（\term{特異単体}{singular simplex}）全体を基底とする自由アーベル群 $C_k(X)$ と、同じ形の境界作用素 $\partial\sigma=\sum_i(-1)^i\sigma|_{[e_0,\dots,\widehat{e_i},\dots,e_k]}$ から\term{特異ホモロジー}{singular homology} $H_k(X)$ を定める。具体的な計算には使わないが、次の性質を保証する役割をもつ。

\begin{mythm}{ホモロジーの基本性質}{ch04-basic}
\begin{enumerate}
\item（関手性）連続写像 $f\colon X\to Y$ は準同型 $f_*\colon H_k(X)\to H_k(Y)$ を誘導し、$(g\circ f)_*=g_*\circ f_*$、$\id_*=\id$。
\item（ホモトピー不変性）$f\simeq g$ なら $f_*=g_*$。特に $X\simeq Y$ なら $H_k(X)\cong H_k(Y)$。可縮な空間では $H_0=\Z$、$H_k=0$（$k\ge1$）。
\item（$H_0$）$H_0(X)\cong\Z^{(\text{弧状連結成分の数})}$。各成分の点が生成元で、同じ成分の2点はホモローグ（道 $\gamma$ を1単体とみなすと $\partial\gamma=$ 終点 $-$ 始点）。
\item（次元）$n$ 次元の胞体複体では $k>n$ で $H_k=0$。
\item（一致）単体複体（$\Delta$ 複体）$K$ の単体的ホモロジーは、その実現 $\lvert K\rvert$ の特異ホモロジーと自然に同型。したがって単体的ホモロジーは単体分割の取り方によらない。
\end{enumerate}
\end{mythm}

\begin{mydef}{被約ホモロジー}{ch04-reduced}
\term{添加写像}{augmentation} $\varepsilon\colon C_0\to\Z$、$\varepsilon(\sum n_iv_i)=\sum n_i$ は $\varepsilon\circ\partial_1=0$ をみたす。$\tilde H_0=\ker\varepsilon/\operatorname{im}\partial_1$、$\tilde H_k=H_k$（$k\ge1$）を\term{被約ホモロジー}{reduced homology}という。$H_0\cong\tilde H_0\oplus\Z$（$X\ne\emptyset$）であり、可縮な空間ではすべての $\tilde H_k$ が $0$。
\end{mydef}

\subsection{完全列・切除・マイヤー・ヴィートリス}

部分空間 $A\subset X$ に対し $C_k(X,A)=C_k(X)/C_k(A)$ は鎖複体になり、\term{相対ホモロジー}{relative homology} $H_k(X,A)$ が定まる。相対サイクルとは「境界が $A$ の中にある鎖」である。

\begin{mythm}{対の長完全列}{ch04-les}
次の列は完全である：
\[
\cdots\to H_k(A)\xrightarrow{i_*}H_k(X)\xrightarrow{j_*}H_k(X,A)\xrightarrow{\ \delta\ }H_{k-1}(A)\to\cdots\to H_0(X,A)\to0 .
\]
連結準同型 $\delta$ は、相対サイクル $c$ に $[\partial c]$（$A$ の中のサイクル）を対応させる。
\end{mythm}

\begin{mythm}{切除定理（認める）}{ch04-excision}
$\bar Z\subset\operatorname{int}A$ なら $H_k(X\setminus Z,A\setminus Z)\cong H_k(X,A)$。とくに $A$ が $X$ の閉部分複体（よい対）なら $H_k(X,A)\cong\tilde H_k(X/A)$。例：$H_k(D^n,S^{n-1})\cong\tilde H_k(S^n)$。
\end{mythm}

\begin{mythm}{マイヤー・ヴィートリス完全列}{ch04-mv}
$X=U\cup V$（$U,V$ は開集合、あるいは $X$ の内部が $X$ を覆う部分複体）のとき、完全列
\[
\cdots\to H_k(U\cap V)\xrightarrow{\ \Phi\ }H_k(U)\oplus H_k(V)\xrightarrow{\ \Psi\ }H_k(X)\xrightarrow{\ \delta\ }H_{k-1}(U\cap V)\to\cdots
\]
がある。ここで $\Phi(c)=(c,-c)$、$\Psi(a,b)=a+b$（包含写像による像）。$\delta$ は、$X$ のサイクル $z=u+v$（$u$ は $U$ の鎖、$v$ は $V$ の鎖）に $[\partial u]=[-\partial v]$（$U\cap V$ のサイクル）を対応させる。$U\cap V\ne\emptyset$ なら被約ホモロジーでも完全。
\end{mythm}
ゲージ理論で空間を2つの座標近傍（パッチ）で覆い、それぞれで場を記述する（ディラック・モノポールのウー・ヤン構成、第5, 7章）のと同じ発想で、ホモロジーを「部品」から組み立てる道具である。

\subsection{胞体ホモロジー}

CW 複体（第1章）$X$ の $k$ 骨格を $X^k$ とする。

\begin{mythm}{胞体ホモロジー}{ch04-cellular}
$C_k^{\mathrm{CW}}=\Z^{(\text{$k$ 胞体の数})}$（各胞体に向きを1つ選んで基底とする）と
\[
d_k(e^k_\alpha)=\sum_\beta d_{\alpha\beta}\,e^{k-1}_\beta
\]
は鎖複体をなし、そのホモロジーは $H_k(X)$ と同型である。係数 $d_{\alpha\beta}$ は、$e^k_\alpha$ の接着写像 $S^{k-1}\to X^{k-1}$ と、$X^{k-1}$ から $e^{k-1}_\beta$ 以外をつぶす写像 $X^{k-1}\to X^{k-1}/(X^{k-1}\setminus e^{k-1}_\beta)\cong S^{k-1}$ の合成の写像度である。
\end{mythm}

実際の計算でよく使うのは次の2つである。
\begin{itemize}
\item $d_1(e^1)=(\text{終点})-(\text{始点})$。
\item 2胞体 $e$ を、辺のラベルの語（例：$aba^{-1}b^{-1}$）で表される多角形として貼るとき、$d_2(e)$ における $1$ 胞体 $a$ の係数は、語に現れる $a$ の\textbf{指数の和}である。
\end{itemize}

\begin{center}
\begin{tikzpicture}[scale=1.15]
  \fill[blue!6] (0,0) rectangle (2,2);
  \draw[thick,mid arrow] (0,0) -- (2,0); \node[below] at (1,0) {$a$};
  \draw[thick,mid arrows] (2,0) -- (2,2); \node[right] at (2,1) {$b$};
  \draw[thick,mid arrow] (0,2) -- (2,2); \node[above] at (1,2) {$a$};
  \draw[thick,mid arrows] (0,0) -- (0,2); \node[left] at (0,1) {$b$};
  \foreach \p in {(0,0),(2,0),(2,2),(0,2)} \node[vtx] at \p {};
  \node at (1,1) {$e$};
  \node at (1,-0.8) {$T^2$：$aba^{-1}b^{-1}$};
  \begin{scope}[xshift=4cm]
  \fill[blue!6] (0,0) rectangle (2,2);
  \draw[thick,mid arrow] (0,0) -- (2,0); \node[below] at (1,0) {$a$};
  \draw[thick,mid arrows] (2,0) -- (2,2); \node[right] at (2,1) {$b$};
  \draw[thick,mid arrow] (2,2) -- (0,2); \node[above] at (1,2) {$a$};
  \draw[thick,mid arrows] (0,0) -- (0,2); \node[left] at (0,1) {$b$};
  \foreach \p in {(0,0),(2,0),(2,2),(0,2)} \node[vtx] at \p {};
  \node at (1,1) {$e$};
  \node at (1,-0.8) {クラインの壺 $K$：$abab^{-1}$};
  \end{scope}
  \begin{scope}[xshift=8cm]
  \fill[blue!6] (0,0) rectangle (2,2);
  \draw[thick,mid arrow] (0,0) -- (2,0); \node[below] at (1,0) {$a$};
  \draw[thick,mid arrows] (2,0) -- (2,2); \node[right] at (2,1) {$b$};
  \draw[thick,mid arrow] (2,2) -- (0,2); \node[above] at (1,2) {$a$};
  \draw[thick,mid arrows] (0,2) -- (0,0); \node[left] at (0,1) {$b$};
  \node[vtx] at (0,0) {}; \node[vtx] at (2,2) {};
  \node[hvtx] at (2,0) {}; \node[hvtx] at (0,2) {};
  \node at (1,1) {$e$};
  \node at (1,-0.8) {$\RP^2$：$abab$};
  \end{scope}
\end{tikzpicture}
\zucap{正方形の辺を矢印の向きがそろうように貼り合わせる。語は左下の頂点から反時計回りに読む。$T^2$ と $K$ では4頂点がすべて1点になり、$\RP^2$ では2点（黒・白）になる。$d_2e$ は順に $a+b-a-b=0$、$a+b+a-b=2a$、$2a+2b$。}
\end{center}

\begin{myex}{胞体ホモロジーの例}{ch04-cellex}
\begin{itemize}
\item $S^n=e^0\cup e^n$（$n\ge1$）：$d=0$ なので $H_0=H_n=\Z$、他は $0$。
\item $T^2$：$C_2=\Z e$、$C_1=\Z a\oplus\Z b$、$C_0=\Z v$、$d_2e=0$、$d_1=0$。$H_*=(\Z,\Z^2,\Z)$。
\item クラインの壺 $K$：$d_2e=2a$、$d_1=0$ から $H_0=\Z$、$H_1=\Z\oplus\Z_2$、$H_2=0$（問題）。
\item $\CP^n=e^0\cup e^2\cup\dots\cup e^{2n}$：奇数次元の胞体がないので $d=0$、$H_{2k}=\Z$（$0\le k\le n$）。
\end{itemize}
\end{myex}

\begin{center}
\small
\renewcommand{\arraystretch}{1.25}
\begin{tabular}{p{0.36\linewidth}|p{0.42\linewidth}|c}
空間 & $H_k$（$k=0,1,2,\dots$） & $\chi$\\ \hline
$S^n$（$n\ge1$） & $\Z$（$k=0,n$）、他は $0$ & $1+(-1)^n$\\
$T^n$ & $\Z^{\binom nk}$ & $0$\\
$\Sigma_g$（種数 $g$ の向き付け可能な閉曲面） & $\Z,\ \Z^{2g},\ \Z$ & $2-2g$\\
$N_h$（$h$ 個の $\RP^2$ の連結和） & $\Z,\ \Z^{h-1}\oplus\Z_2,\ 0$ & $2-h$\\
クラインの壺 $K=N_2$ & $\Z,\ \Z\oplus\Z_2,\ 0$ & $0$\\
$\RP^n$ & $\Z$（$k=0$ と $k=n$ 奇数）、$\Z_2$（$k$ 奇数 $<n$）、他は $0$ & $\frac{1+(-1)^n}{2}$\\
$\CP^n$ & $\Z$（$k=0,2,\dots,2n$）、他は $0$ & $n+1$\\
レンズ空間 $L(p,q)$ & $\Z,\ \Z_p,\ 0,\ \Z$ & $0$
\end{tabular}
\end{center}
$T^n$ の結果はキュネットの公式（第6章）から従う。本章では $T^2$ と、$H_1(T^d)\cong\Z^d$（各座標方向に一周する輪が基底）を用いる。

\subsection{オイラー標数}

\begin{mythm}{オイラー・ポアンカレの公式}{ch04-euler}
有限 CW 複体 $X$ の $k$ 胞体の数を $c_k$ とすると
\[
\chi(X):=\sum_k(-1)^kc_k=\sum_k(-1)^kb_k .
\]
右辺はホモロジーだけで決まるので、$\chi$ は胞体分割によらないホモトピー不変量である。
\end{mythm}
\begin{note}
$Z_k=\ker d_k$、$B_{k-1}=\operatorname{im}d_k$ とすると短完全列 $0\to Z_k\to C_k\to B_{k-1}\to0$、$0\to B_k\to Z_k\to H_k\to0$ から $c_k=\operatorname{rank}Z_k+\operatorname{rank}B_{k-1}=b_k+\operatorname{rank}B_k+\operatorname{rank}B_{k-1}$。$(-1)^k$ をかけて足すと $\operatorname{rank}B_k$ の項が打ち消し合う。
\end{note}
性質：$\chi(X\times Y)=\chi(X)\chi(Y)$、$\chi(X\cup Y)=\chi(X)+\chi(Y)-\chi(X\cap Y)$（部分複体のとき）、$n$ 重被覆 $\tilde X\to X$ では $\chi(\tilde X)=n\chi(X)$。

閉曲面 $\Sigma$ 上のリーマン計量のガウス曲率 $K$ について $\int_\Sigma K\,dA=2\pi\chi(\Sigma)$（ガウス・ボンネの定理、第8章）、閉多様体 $M$ 上の孤立零点をもつベクトル場の零点の指数の和は $\chi(M)$（ポアンカレ・ホップの定理）である。どちらも本章では結果のみを用いる。

\subsection{フレヴィッチの定理・写像度・レフシェッツの定理}

\begin{mythm}{フレヴィッチの定理（1次）}{ch04-hurewicz}
$X$ が弧状連結なら、ループ $\gamma$ をその1サイクルとしてのクラスに写す写像 $h\colon\pi_1(X,x_0)\to H_1(X)$ は全射準同型で、核は交換子部分群 $[\pi_1,\pi_1]$ である。すなわち $H_1(X)\cong\pi_1(X)^{\mathrm{ab}}=\pi_1/[\pi_1,\pi_1]$。
\end{mythm}
例：$\pi_1(S^1\vee S^1)=F_2$（自由群）に対し $H_1=\Z^2$。$\pi_1(\Sigma_g)=\langle a_i,b_i\mid\prod[a_i,b_i]=1\rangle$ に対し $H_1=\Z^{2g}$。$\pi_1$ が可換なら $H_1\cong\pi_1$（例：$\pi_1(SO(3))=\Z_2=H_1(SO(3))$）。

\begin{mydef}{写像度}{ch04-degree}
$n\ge1$ とし、$f\colon S^n\to S^n$ に対し $f_*\colon H_n(S^n)\cong\Z\to\Z$ はある整数 $\deg f$ 倍である。これを $f$ の\term{写像度}{degree}という。
\end{mydef}

\begin{myprop}{写像度の性質}{ch04-degprop}
\begin{enumerate}
\item $\deg(g\circ f)=\deg g\cdot\deg f$、$\deg\id=1$。$f\simeq g$ なら $\deg f=\deg g$。
\item 1つの座標の符号を反転する鏡映は $\deg=-1$。対蹠写像 $x\mapsto-x$ は $n+1$ 個の鏡映の合成なので $\deg=(-1)^{n+1}$。
\item（局所写像度）$f$ が滑らかで $y$ が正則値（$f^{-1}(y)$ の各点で $\det Df\ne0$）なら $\deg f=\sum_{x\in f^{-1}(y)}\operatorname{sign}\det Df_x$。
\item $n=1$ では $\deg f$ は巻き付き数（第2章）に、$n=2$ の $O(3)$ シグマ模型では $\deg\bm n=\frac1{4\pi}\int\bm n\cdot(\partial_1\bm n\times\partial_2\bm n)\,d^2x$ に等しい（第3章）。
\end{enumerate}
\end{myprop}

\begin{mythm}{レフシェッツの不動点定理（認める）}{ch04-lefschetz}
$X$ を有限 CW 複体（コンパクトな多様体など）、$f\colon X\to X$ を連続写像とし、
\[
L(f)=\sum_k(-1)^k\tr\bigl(f_*\colon H_k(X;\Q)\to H_k(X;\Q)\bigr)
\]
を\term{レフシェッツ数}{Lefschetz number}とする（$H_k(X;\Q)\cong\Q^{b_k}$ でのトレース）。$L(f)\ne0$ なら $f$ は不動点をもつ。また $L(\id)=\chi(X)$ であり、$f\simeq\id$ なら $L(f)=\chi(X)$。不動点がすべて孤立していれば、その指数の和は $L(f)$ に等しい（レフシェッツ・ホップ）。
\end{mythm}

\section{問題}

\subsection{鎖複体・グラフ・ファインマン図}

\begin{problem}[難易度 A]\label{ch04-prob-dd}
(1) 単体的境界作用素について $\partial\circ\partial=0$ を示せ。\\
(2) 四面体 $\sigma=[0,1,2,3]$ について $\partial\sigma$ と $\partial\partial\sigma$ を具体的に計算せよ。\\
(3) $\R^3$ の領域 $V$ に対し、ベクトルポテンシャル $\bm A$ で書ける磁場 $\bm B=\nabla\times\bm A$ の $\partial V$ を貫く全磁束が $0$ になることが、(1) の「連続版」であることを説明せよ。
\end{problem}

\begin{solution}
(1) $\sigma=[v_0,\dots,v_k]$ とする。
\[
\begin{aligned}
\partial\partial\sigma&=\sum_{i}(-1)^i\,\partial[v_0,\dots,\widehat{v_i},\dots,v_k]\\
&=\sum_{j<i}(-1)^i(-1)^j[\dots\widehat{v_j}\dots\widehat{v_i}\dots]+\sum_{j>i}(-1)^i(-1)^{j-1}[\dots\widehat{v_i}\dots\widehat{v_j}\dots].
\end{aligned}
\]
第2の和で $v_j$（$j>i$）は $v_i$ を除いた列の $j-1$ 番目にあるので符号が $(-1)^{j-1}$ になる。第2の和で $i$ と $j$ の名前を入れかえると $\sum_{i<j}$ が $\sum_{j<i}$ になり、各項は $-(-1)^{i+j}[\dots\widehat{v_j}\dots\widehat{v_i}\dots]$ となって第1の和と打ち消し合う。よって $\partial\partial\sigma=0$ で、線形性から $\partial\partial=0$。

(2)
\[
\partial[0,1,2,3]=[1,2,3]-[0,2,3]+[0,1,3]-[0,1,2].
\]
各項の境界は
\begin{align*}
\partial[1,2,3]&=[2,3]-[1,3]+[1,2],&\partial[0,2,3]&=[2,3]-[0,3]+[0,2],\\
\partial[0,1,3]&=[1,3]-[0,3]+[0,1],&\partial[0,1,2]&=[1,2]-[0,2]+[0,1].
\end{align*}
符号つきで足すと
\[
\begin{aligned}
\partial\partial\sigma&=\bigl([2,3]-[1,3]+[1,2]\bigr)-\bigl([2,3]-[0,3]+[0,2]\bigr)\\
&\quad+\bigl([1,3]-[0,3]+[0,1]\bigr)-\bigl([1,2]-[0,2]+[0,1]\bigr)=0 .
\end{aligned}
\]
6本の辺がそれぞれちょうど2回、逆符号で現れる。$\partial\sigma$ は四面体の表面（4つの三角形）を外向きに一貫して向き付けたもので、閉曲面だから境界をもたない。

(3) ストークスの定理 $\int_S(\nabla\times\bm A)\cdot d\bm S=\oint_{\partial S}\bm A\cdot d\bm l$ と、ガウスの定理 $\int_V\nabla\cdot\bm B\,d^3x=\int_{\partial V}\bm B\cdot d\bm S$ は、どちらも「積分領域の境界 $\partial$」と「微分 $d$」が互いに随伴であること $\int_{\partial c}\omega=\int_c d\omega$ の例である。$S=\partial V$ とすると
\[
\int_{\partial V}\bm B\cdot d\bm S=\oint_{\partial\partial V}\bm A\cdot d\bm l=0
\]
（$\partial\partial V=\emptyset$）。同じことを微分の側で言えば $\nabla\cdot(\nabla\times\bm A)=0$、すなわち $d^2=0$ である。磁気単極子があれば $\bm A$ は大域的には存在せず、この議論は破れる（ディラック・モノポール、第5章）。
\end{solution}

\begin{problem}[難易度 A]\label{ch04-prob-simplex}
$\Delta^2$ を三角形 $[v_0,v_1,v_2]$（すべての面を含む単体複体）、$\partial\Delta^2$ をその3辺と3頂点からなる部分複体とする。$H_k(\Delta^2)$ と $H_k(\partial\Delta^2)$ を、鎖群と境界作用素を書き下して計算せよ。被約ホモロジーではどうなるか。
\end{problem}

\begin{solution}
辺を $e_0=[v_1,v_2]$、$e_1=[v_0,v_2]$、$e_2=[v_0,v_1]$ とおく。
\[
\partial e_0=v_2-v_1,\quad\partial e_1=v_2-v_0,\quad\partial e_2=v_1-v_0,\qquad\partial[v_0,v_1,v_2]=e_0-e_1+e_2 .
\]

$\partial_1$：$\partial_1(xe_0+ye_1+ze_2)=-(y+z)v_0+(z-x)v_1+(x+y)v_2$。これが $0$ となるのは $z=x$、$y=-x$ のときで、$Z_1=\Z(e_0-e_1+e_2)$。$\operatorname{im}\partial_1$ は $v_1-v_0$、$v_2-v_0$ で生成され、$\operatorname{im}\partial_1=\ker\varepsilon=\{\sum n_iv_i:\sum n_i=0\}$。よって $H_0=C_0/\operatorname{im}\partial_1\cong\Z$（$\varepsilon$ による同型）。

$\partial\Delta^2$：$C_2=0$ なので $H_2=0$、$H_1=Z_1=\Z(e_0-e_1+e_2)\cong\Z$、$H_0\cong\Z$。生成元 $e_0-e_1+e_2=[v_0,v_1]+[v_1,v_2]+[v_2,v_0]$ は三角形の縁を一周するサイクルである。

$\Delta^2$：$\partial_2\colon\Z\to C_1$ は $[v_0,v_1,v_2]\mapsto e_0-e_1+e_2\ne0$ で単射なので $H_2=\ker\partial_2=0$。$B_1=\operatorname{im}\partial_2=\Z(e_0-e_1+e_2)=Z_1$ なので $H_1=0$。$H_0\cong\Z$。

被約ホモロジー：両方とも $\tilde H_0=\ker\varepsilon/\operatorname{im}\partial_1=0$。$\tilde H_*(\Delta^2)=0$（可縮だから当然）、$\tilde H_1(\partial\Delta^2)=\Z$、他は $0$。$\partial\Delta^2$ は円周と同相なので、$\tilde H_k(S^1)=\Z$（$k=1$）、$0$（それ以外）が得られた。縁を一周する輪は、内部の三角形があれば境界になり（「穴がふさがっている」）、なければ境界にならない。
\end{solution}

\begin{problem}[難易度 B]\label{ch04-prob-graph}
有限グラフ $\Gamma$（頂点 $V$ 個、辺 $E$ 個、連結成分 $b_0$ 個、2単体はない）を1次元の単体複体（またはループや多重辺を許す $\Delta$ 複体）とみなす。\\
(1) $H_0(\Gamma)\cong\Z^{b_0}$、$H_1(\Gamma)$ は階数 $E-V+b_0$ の自由アーベル群であることを示せ。\\
(2) $\Gamma$ が連結で、$T\subset\Gamma$ をその全域木（すべての頂点を含む木、辺は $V-1$ 本；木は $H_1=0$ をみたすことを認めてよい）とする。$T$ に含まれない辺 $e$ ごとに、$e$ と $T$ の中の道でできる閉路 $z_e$ を考える。$\{z_e\}$ が $H_1(\Gamma)$ の基底であることを示せ。
\end{problem}

\begin{solution}
(1) $H_0$：同じ連結成分の2頂点 $u,w$ は辺の列で結ばれ、その辺を向きに応じて符号をつけて足した1鎖 $c$ は $\partial c=w-u$ をみたすので $u$ と $w$ はホモローグ。よって $H_0$ は各成分の代表頂点 $u_1,\dots,u_{b_0}$ で生成される。一方、成分ごとの添加写像 $\varepsilon_\alpha(\sum n_vv)=\sum_{v\in\Gamma_\alpha}n_v$ は各辺の境界（同じ成分の頂点の差）で $0$ になるので、$(\varepsilon_1,\dots,\varepsilon_{b_0})\colon H_0\to\Z^{b_0}$ が定まり、$u_\alpha$ を標準基底に写す。したがって $H_0\cong\Z^{b_0}$。

$H_1$：$C_2=0$ なので $H_1=\ker\partial_1$ は自由アーベル群 $C_1$ の部分群であり、自由である。$H_0=C_0/\operatorname{im}\partial_1\cong\Z^{b_0}$ から $\operatorname{rank}\operatorname{im}\partial_1=V-b_0$。完全列 $0\to\ker\partial_1\to C_1\to\operatorname{im}\partial_1\to0$ で階数は加法的なので
\[
b_1=\operatorname{rank}H_1=E-(V-b_0)=E-V+b_0 .
\]
連結なら $b_1=E-V+1$。とくに $\chi(\Gamma)=V-E=b_0-b_1$。

(2) $T$ の外の辺は $E-(V-1)=b_1$ 本ある。$e$ の端点を $T$ の中の（一意な）道で結んで閉路 $z_e$ を作ると、$z_e$ は $e$ の係数が $1$、$T$ の外の他の辺の係数が $0$ のサイクルである。

一次独立性：$\sum_en_ez_e=0$ なら、辺 $e$ の係数を見て $n_e=0$。

生成：任意のサイクル $z$ に対し、$z$ における辺 $e\notin T$ の係数を $n_e$ として $z'=z-\sum_en_ez_e$ とおくと、$z'$ は $T$ の辺だけでできたサイクルである。$H_1(T)=0$ かつ $T$ には2単体がないので $z'=0$。よって $z=\sum n_ez_e$。

物理では、全域木の外の各辺 $e$ が「独立なループ運動量」1つに対応し、$z_e$ に沿ってその運動量が流れる。
\end{solution}

\begin{problem}[難易度 B]\label{ch04-prob-feynman}
連結なファインマン図を、頂点 $V$ 個、内線 $I$ 本、外線 $E$ 本のグラフ $\Gamma$ とみなす（外線は無視するか、外線の先に1価の頂点を付けても $b_1$ は変わらない）。内線 $\ell$ に向きを付け、4元運動量 $k_\ell\in\R^4$ を割り当てる。頂点 $v$ に外から流れ込む外線運動量の和を $p_v$ とする。\\
(1) 各頂点での運動量保存（流入 $=$ 流出）が $\partial k+p=0$ と書けることを示せ。ただし $k=\sum_\ell k_\ell\,\ell\in C_1(\Gamma)\otimes\R^4$、$p=\sum_vp_v\,v\in C_0(\Gamma)\otimes\R^4$。\\
(2) 解 $k$ が存在するための必要十分条件は $\sum_vp_v=0$（全運動量保存）であり、解全体は $4L$ 次元のアフィン空間で $L=I-V+1=b_1(\Gamma)$ となることを示せ（$L$ をループ数という）。\\
(3) 作用を $S/\hbar$ の形で書くと、伝播関数（内線・外線とも）は $\hbar$ に比例し、頂点は $\hbar^{-1}$ に比例する。この規格化で、連結図の寄与が $\hbar^{E+L-1}$ に比例することを示せ。\\
(4) $\phi^4$ 理論で $4V=2I+E$、$\phi^3$ 理論で $3V=2I+E$ を示し、$L$ を $V$ と $E$ で表せ。下図の (a)(b)(c) について $V,I,E,L$ を求めよ。
\begin{center}
\begin{tikzpicture}[scale=0.9]
  % (a) fish
  \node[vtx] (a1) at (0,0) {}; \node[vtx] (a2) at (1.6,0) {};
  \draw[thick] (a1) to[out=60,in=120] (a2);
  \draw[thick] (a1) to[out=-60,in=-120] (a2);
  \draw[thick] (a1) -- ++(150:0.9); \draw[thick] (a1) -- ++(210:0.9);
  \draw[thick] (a2) -- ++(30:0.9); \draw[thick] (a2) -- ++(-30:0.9);
  \node at (0.8,-1.3) {(a) $\phi^4$, 4点1ループ};
  % (b) sunset
  \begin{scope}[xshift=4.6cm]
  \node[vtx] (b1) at (0,0) {}; \node[vtx] (b2) at (1.6,0) {};
  \draw[thick] (b1) to[out=70,in=110] (b2);
  \draw[thick] (b1) -- (b2);
  \draw[thick] (b1) to[out=-70,in=-110] (b2);
  \draw[thick] (b1) -- ++(180:0.9); \draw[thick] (b2) -- ++(0:0.9);
  \node at (0.8,-1.3) {(b) $\phi^4$, サンセット};
  \end{scope}
  % (c) triangle
  \begin{scope}[xshift=9.2cm]
  \node[vtx] (c1) at (0,-0.5) {}; \node[vtx] (c2) at (1.4,-0.5) {}; \node[vtx] (c3) at (0.7,0.7) {};
  \draw[thick] (c1) -- (c2) -- (c3) -- (c1);
  \draw[thick] (c1) -- ++(210:0.8); \draw[thick] (c2) -- ++(-30:0.8); \draw[thick] (c3) -- ++(90:0.7);
  \node at (0.7,-1.3) {(c) $\phi^3$, 三角形};
  \end{scope}
\end{tikzpicture}
\zucap{ファインマン図をグラフとみなす。黒丸が相互作用頂点、太線が伝播関数。ループ数は内線と頂点が作るグラフの第1ベッチ数。}
\end{center}
\end{problem}

\begin{solution}
(1) $\partial\ell=(\text{終点})-(\text{始点})$ なので、$\partial k$ の頂点 $v$ の係数は
\[
(\partial k)_v=\sum_{\ell\text{ が }v\text{ に入る}}k_\ell-\sum_{\ell\text{ が }v\text{ から出る}}k_\ell .
\]
運動量保存「流入の和 $=$ 流出の和」は $(\partial k)_v+p_v=0$ にほかならない。

(2) $\R^4$ の各成分ごとに考えればよい（$C_*\otimes\R^4=(C_*\otimes\R)^{\oplus4}$）。グラフが連結なので、前問より $H_0(\Gamma;\R)=C_0\otimes\R/\operatorname{im}\partial\cong\R$ で、同型は添加写像 $\varepsilon$ で与えられる。すなわち $\operatorname{im}\partial=\ker\varepsilon$。よって $-p\in\operatorname{im}\partial$ となる条件は $\varepsilon(p)=\sum_vp_v=0$。解が1つ $k_0$ あれば、解全体は $k_0+\ker\partial=k_0+Z_1(\Gamma;\R)\otimes\R^4$ であり、$C_2=0$ なので $Z_1=H_1$、その次元は（前問より）$4(I-V+1)=4b_1$。したがって独立な積分変数（ループ運動量）は $L=b_1=I-V+1$ 個である。

(3) 図の値は $\hbar^{I+E}\cdot\hbar^{-V}$ に比例し、
\[
I+E-V=E+(I-V+1)-1=E+L-1 .
\]
外線の数 $E$ を固定すると、$\hbar$ の次数はループ数で決まる：木図（$L=0$）が古典的寄与、$L$ ループが $\hbar^L$ の量子補正である。

(4) 線の端点を数える。$\phi^4$ では各頂点から4本の線の端が出る。内線は両端とも頂点につくので2回、外線は1回数えられ、$4V=2I+E$。$\phi^3$ では $3V=2I+E$。よって
\[
\phi^4:\ L=\frac{4V-E}{2}-V+1=V-\frac E2+1,\qquad\phi^3:\ L=\frac{3V-E}{2}-V+1=\frac{V-E}{2}+1 .
\]
(a) $V=2$、$I=2$、$E=4$：$L=2-2+1=1$。(b) $V=2$、$I=3$、$E=2$：$L=3-2+1=2$（公式でも $2-1+1=2$）。(c) $V=3$、$I=3$、$E=3$：$L=3-3+1=1$（公式でも $(3-3)/2+1=1$）。(b) のサンセット図は3本の内線が2つの独立な閉路を作り、2ループの自己エネルギーである。
\end{solution}

\begin{problem}[難易度 A]\label{ch04-prob-c60}
凸多面体の表面は $S^2$ と同相で、頂点・辺・面が CW 構造を与える。$\chi(S^2)=2$ を用いよ。\\
(1) すべての頂点に辺がちょうど3本集まる多面体で、$n$ 角形の面の数を $f_n$ とする。$\sum_n(6-n)f_n=6\chi=12$ を示せ。\\
(2) 面が五角形と六角形だけなら、五角形はちょうど12個であることを示せ。フラーレン $\mathrm{C}_{60}$（サッカーボール、六角形20個）の頂点・辺・面の数を求めよ。\\
(3) 同じ条件（3価、五角形と六角形のみ）でトーラス $T^2$ を分割するとき、五角形は現れないことを示せ（カーボン・ナノチューブを閉じたトーラスは六角形だけでできる）。
\end{problem}

\begin{solution}
(1) $F=\sum_nf_n$。各辺は2つの面に属するので $\sum_nnf_n=2E$。各頂点に3本の辺が集まり、各辺は2頂点をもつので $3V=2E$。よって
\[
6\chi=6V-6E+6F=4E-6E+6F=6F-2E=\sum_n(6-n)f_n .
\]
$S^2$ では $\chi=2$ なので右辺は $12$。

(2) $f_5$ と $f_6$ だけなら $(6-5)f_5+0\cdot f_6=12$、すなわち $f_5=12$（$f_6$ は任意）。$\mathrm{C}_{60}$ では $F=12+20=32$、$2E=5\cdot12+6\cdot20=180$ より $E=90$、$V=2E/3=60$（炭素原子60個）。確かに $60-90+32=2$。

(3) $\chi(T^2)=0$ なので $f_5=6\chi=0$。六角形だけの3価分割（グラフェンの蜂の巣格子を周期的に閉じたもの）はトーラス上で可能だが、球面上では不可能である。一般に $n=7$ の面（七角形）は曲率の負の欠陥で、$\sum(6-n)f_n=6\chi$ は「曲率欠陥の総和は位相で決まる」というガウス・ボンネの定理（第8章）の離散版である。
\end{solution}

\subsection{格子ゲージ理論と鎖}

以下の3問では、$d$ 次元の周期的な立方格子（各方向に $N\ge3$ サイト）を、サイト（0胞体）、リンク（1胞体）、プラケット（2胞体）、キューブ（3胞体）、…からなる $T^d$ の胞体分割とみなす。$\hat\mu$ は $\mu$ 方向の単位ベクトル、$\ell_\mu(x)$ は $x$ から $x+\hat\mu$ へのリンク、$p_{\mu\nu}(x)$ は $x,x+\hat\mu,x+\hat\mu+\hat\nu,x+\hat\nu$ を頂点とするプラケットで $\partial p_{\mu\nu}(x)=\ell_\mu(x)+\ell_\nu(x+\hat\mu)-\ell_\mu(x+\hat\nu)-\ell_\nu(x)$ となる向きをもつ（$p_{\nu\mu}=-p_{\mu\nu}$）。$U(1)$ リンク変数 $U_\ell=e^{i\theta_\ell}$ は $\ell$ に沿った平行移動（$\theta_\ell\approx\int_\ell A$）とし、$\theta$ を1鎖上の線形関数 $\theta(\sum n_\ell\ell)=\sum n_\ell\theta_\ell$ に拡張する。ゲージ変換 $g(x)=e^{i\alpha(x)}$ は $\theta_\ell\mapsto\theta_\ell+\alpha(\partial\ell)$、すなわち $\theta_{\ell}\mapsto\theta_\ell+\alpha(\text{終点})-\alpha(\text{始点})$ と作用する（$A\mapsto A+d\alpha$ に対応）。

\begin{problem}[難易度 B]\label{ch04-prob-gauss}
有限連結グラフ $\Gamma$（例えば上の格子のサイトとリンク）上の静電場を考える。電場は各リンクの値 $E_\ell\in\R$ で、1鎖 $\mathcal E=\sum_\ell E_\ell\,\ell$ とみなす（リンクの向きを逆にすると符号が変わる）。電荷分布は0鎖 $\rho=\sum_x\rho_x\,x$。ガウスの法則を
\[
\sum_{\ell:\ x\text{ から出る}}E_\ell-\sum_{\ell:\ x\text{ に入る}}E_\ell=\rho_x
\]
とする。\\
(1) ガウスの法則が $\partial\mathcal E=-\rho$ と書けることを示し、全電荷 $Q=\sum_x\rho_x=0$ を導け。\\
(2) 逆に $Q=0$ なら、ガウスの法則をみたす $\mathcal E$ が存在することを示せ。解の不定性は何か。\\
(3) $\Gamma$ を周期的立方格子 $T^d$ とし、$\rho=0$ とする。$\mu$ 方向のリンクのうち $x_\mu=s$ から $x_\mu=s+1$ へ向かうものすべてについての和 $\Phi_\mu(s)=\sum E_{\ell_\mu(x)}$（$x_\mu=s$）は $s$ によらないことを示せ。これは「$T^d$ を一周する電束」で、ガウスの法則だけでは決まらない量である。
\end{problem}

\begin{solution}
(1) $\partial\ell=(\text{終点})-(\text{始点})$ なので $\partial\mathcal E$ の $x$ の係数は（入る $E_\ell$ の和）$-$（出る $E_\ell$ の和）、すなわち $-(\nabla\cdot E)_x$ である。よってガウスの法則は $\partial\mathcal E=-\rho$。$\varepsilon\circ\partial=0$ より
\[
Q=\varepsilon(\rho)=-\varepsilon(\partial\mathcal E)=0 .
\]
これは連続理論で $\int_M\nabla\cdot\bm E=\int_{\partial M}\bm E\cdot d\bm S=0$（$M$ が閉じた空間）としたことの格子版である。

(2) $\Gamma$ が連結なので $H_0(\Gamma;\R)\cong\R$ で、$\operatorname{im}\partial_1=\ker\varepsilon$（問題 \ref{ch04-prob-graph}）。$Q=\varepsilon(\rho)=0$ なら $-\rho\in\operatorname{im}\partial_1$、すなわち $\partial\mathcal E=-\rho$ となる $\mathcal E$ が存在する。電荷が整数で $E_\ell$ も整数（コンパクトな $U(1)$ の電束）としても、$\Z$ 係数で同じ議論が成り立つ。解の不定性は $\ker\partial_1=Z_1(\Gamma)$、すなわち電荷を生まない「閉じた電束」である。連結でない場合は各成分ごとに全電荷が $0$ であることが条件になる。

$Z_1$ はプラケットの境界（局所的な「渦」をもつ電場、光子の横波の自由度）と、空間を一周する電束に分かれる。前者は $B_1=\partial(\text{プラケットの2鎖})$、後者は商 $Z_1/B_1=H_1(T^d)\cong\Z^d$ で数えられる。次の (3) はこの後者の部分を取り出す。

(3) $x_\mu=s+1$ をみたすすべてのサイトについてガウスの法則（$\rho=0$）を足す。$\mu$ 以外の方向のリンクは両端がこの「層」に含まれるので、出る分と入る分で1回ずつ現れて打ち消す。$\mu$ 方向のリンクのうち、層から出るものは $x_\mu=s+1\to s+2$、層に入るものは $x_\mu=s\to s+1$ である。よって
\[
0=\Phi_\mu(s+1)-\Phi_\mu(s),
\]
すなわち $\Phi_\mu$ は $s$ によらない。$\Phi_\mu$ は $\mathcal E$ の $\mu$ 方向の「巻き付き」を測る量で、たとえば $\mu$ 方向に一周する閉路 $P_\mu$ 上に $E=1$ を置いた配位は $\rho=0$ で $\Phi_\mu=1$ をもつ。$\Phi_\mu$ はプラケットの境界を加えても変わらない（$p_{\mu\nu}$ の境界は $\mu$ 方向のリンクを同じ層の中に $+\ell_\mu(x)$、$-\ell_\mu(x+\hat\nu)$ と1本ずつ含み、$\mu$ を含まないプラケットの境界は $\mu$ 方向のリンクを含まない）ので、ホモロジー類 $[\mathcal E]\in H_1(T^d;\R)\cong\R^d$ の成分を与える。
\end{solution}

\begin{problem}[難易度 B]\label{ch04-prob-bianchi}
3次元の周期的立方格子 $T^3$ を考え、$x$ を最も原点側の頂点とするキューブ $c(x)$ に、$\R^3$ の標準的な向き（右手系）を与える。
\begin{center}
\begin{tikzpicture}[scale=1.5]
  \coordinate (A) at (0,0); \coordinate (B) at (2,0); \coordinate (C) at (2,2); \coordinate (D) at (0,2);
  \coordinate (E) at (0.9,0.7); \coordinate (F) at (2.9,0.7); \coordinate (G) at (2.9,2.7); \coordinate (H) at (0.9,2.7);
  \fill[red!12] (D) -- (C) -- (G) -- (H) -- cycle;
  \fill[blue!10] (B) -- (F) -- (G) -- (C) -- cycle;
  \draw[thick] (A) -- (B) -- (C) -- (D) -- cycle;
  \draw[thick] (B) -- (F) -- (G) -- (C);
  \draw[thick] (D) -- (H) -- (G);
  \draw[dashed] (A) -- (E) -- (F); \draw[dashed] (E) -- (H);
  \node[vtx] at (A) {}; \node[below left] at (A) {$x$};
  \draw[-{Stealth[length=2mm]},thick,red!70!black] (1.45,2.35) -- (1.45,3.1) node[above] {\footnotesize $+p_{12}(x+\hat3)$};
  \draw[-{Stealth[length=2mm]},thick,blue!70!black] (2.45,1.35) -- (3.25,1.35) node[right] {\footnotesize $+p_{23}(x+\hat1)$};
  \node[gray] at (1.0,1.0) {\footnotesize $-p_{31}(x)$（手前）};
  \draw[-{Stealth[length=2mm]}] (-1.3,0) -- (-0.6,0) node[right] {\footnotesize $1$};
  \draw[-{Stealth[length=2mm]}] (-1.3,0) -- (-0.85,0.35) node[above right] {\footnotesize $2$};
  \draw[-{Stealth[length=2mm]}] (-1.3,0) -- (-1.3,0.7) node[above] {\footnotesize $3$};
\end{tikzpicture}
\zucap{キューブ $c(x)$。各面は外向きの法線が右手系を与えるように向き付けられる（上面 $x_3=1$ は $+p_{12}$、右面 $x_1=1$ は $+p_{23}$、手前の面 $x_2=0$ は $-p_{31}(x)=p_{13}(x)$）。}
\end{center}
(1) $\partial c(x)=\sum_{\mu}\bigl[p_{\nu\lambda}(x+\hat\mu)-p_{\nu\lambda}(x)\bigr]$（$(\mu,\nu,\lambda)$ は $(1,2,3)$ の巡回置換を走る）とする。リンク $\ell_1(x)$ の係数を調べて $\partial\partial c(x)$ でそれが $0$ になることを確かめよ（他のリンクも同様である）。\\
(2) リンク変数の代表 $\theta_\ell\in\R$ を選び、プラケット変数 $\theta_p=\theta(\partial p)$ とする。$\theta(\partial\partial c)=0$ を、格子上の「$\nabla\cdot\bm B=0$」として解釈せよ。\\
(3) $\theta_p=\bar\theta_p+2\pi n_p$（$\bar\theta_p\in(-\pi,\pi]$、$n_p\in\Z$）と分け、キューブの\textbf{モノポール数} $m_c=\frac1{2\pi}\bar\theta(\partial c)$ を定める（ドグラン・トゥサンの定義）。$m_c\in\Z$ であり、$m_c$ はリンク変数の代表の選び方によらないこと、さらに全モノポール数 $\sum_cm_c$ が $0$ であることを示せ。
\end{problem}

\begin{solution}
(1) $\ell_1(x)$ を含むプラケットは $p_{12}(x)$ と $p_{13}(x)$ で、それぞれの境界に係数 $+1$ で現れる（$\partial p_{\mu\nu}(x)$ の第1項 $\ell_\mu(x)$）。$\partial c(x)$ にはこれらが $-p_{12}(x)$（$\mu=3$ の項）と $-p_{31}(x)=+p_{13}(x)$（$\mu=2$ の項）として現れる。よって $\partial\partial c(x)$ における $\ell_1(x)$ の係数は $-1+1=0$。キューブの12本の辺はそれぞれちょうど2つの面に属し、外向きに向き付けた面の境界では逆向きにたどられるので、すべて打ち消し合う（キューブの表面は閉曲面）。

(2) $\theta$ は線形なので $\theta(\partial\partial c)=0$ は、6つのプラケット変数について
\[
\sum_{p\subset\partial c}[\partial c:p]\,\theta_p=0\qquad(\text{実数として厳密に})
\]
を意味する（$[\partial c:p]=\pm1$ は $\partial c$ における $p$ の係数）。$\theta_p\approx a^2F_{\nu\lambda}$、すなわち $p_{\nu\lambda}$ を貫く磁束 $a^2B_\mu$ なので、これは
\[
\sum_\mu\bigl[B_\mu(x+\hat\mu)-B_\mu(x)\bigr]a^2=0
\]
という差分版の $\nabla\cdot\bm B=0$ である。すなわち「キューブの表面を出る全磁束は $0$」で、$dF=0$（ビアンキ恒等式）の格子版である。$\theta$ が1鎖上の関数、$\theta_p$ がその「$d$」であることから、これは $\partial\partial=0$ の双対 $d^2=0$ にほかならない。

(3) $\theta(\partial c)$ の定義から $\sum_p[\partial c:p]\theta_p=0$ なので
\[
m_c=\frac1{2\pi}\sum_p[\partial c:p]\bar\theta_p=\frac1{2\pi}\sum_p[\partial c:p](\theta_p-2\pi n_p)=-\sum_p[\partial c:p]\,n_p\in\Z .
\]
代表の取りかえ $\theta_\ell\to\theta_\ell+2\pi k_\ell$ では $\theta_p$ が $2\pi$ の整数倍だけ変わるだけなので $\bar\theta_p$ は不変、よって $m_c$ も不変。ゲージ変換でも $\theta_p$ は不変（$\alpha(\partial\partial p)=0$）である。$\bar\theta_p$ は物理的な磁束（$(-\pi,\pi]$ に入る）、$2\pi n_p$ は $p$ を貫く「見えないディラック弦」と解釈され、$m_c\neq0$ のキューブにはディラック弦の端点、すなわち磁気単極子がある。

全モノポール数：周期的格子では各プラケット $p_{\nu\lambda}(y)$ はちょうど2つのキューブ $c(y)$ と $c(y-\hat\mu)$ の面であり、前者では $-p_{\nu\lambda}(y)$、後者では $+p_{\nu\lambda}(y)$ として現れる。したがって $\sum_c\partial c=0$（$T^3$ の基本サイクル $\sum_cc$ はサイクル）であり
\[
\sum_cm_c=\frac1{2\pi}\bar\theta\Bigl(\sum_c\partial c\Bigr)=0 .
\]
閉じた空間では磁気単極子と反単極子が必ず同数ある。これは問題 \ref{ch04-prob-gauss} の全電荷 $0$ と同じ構造である。
\end{solution}

\begin{problem}[難易度 C]\label{ch04-prob-polyakov}
$d$ 次元の周期的立方格子 $T^d$ 上の $U(1)$ 格子ゲージ理論を考える。1サイクル $\gamma$ に対しウィルソン・ループを $W(\gamma)=e^{i\theta(\gamma)}$ とし、$P_\mu$ を原点を通り $\mu$ 方向に一周する閉路（ポリヤコフ・ループ）とする。$H_1(T^d)\cong\Z^d$ で、$[P_1],\dots,[P_d]$ が基底であることは認めてよい。\\
(1) $W(\gamma)$ はゲージ不変であることを示せ。\\
(2) 配位が\textbf{平坦}（すべてのプラケットで $U_p=e^{i\theta_p}=1$）なら、$W(\gamma)$ は $\gamma$ のホモロジー類だけで決まることを示せ。したがって平坦な配位のゲージ不変な情報は $(W(P_1),\dots,W(P_d))\in U(1)^d$ に含まれる。\\
(3) 任意の $(e^{i\varphi_1},\dots,e^{i\varphi_d})\in U(1)^d$ が平坦な配位のポリヤコフ・ループとして実現されることを示せ。\\
(4) ポリヤコフ・ループがすべて等しい2つの平坦な配位は、ゲージ変換で移り合うことを示せ。以上から、平坦な $U(1)$ 配位のゲージ同値類の空間は $U(1)^d=\Hom(H_1(T^d),U(1))$ であり、その次元は $b_1(T^d)=d$ である。
\end{problem}

\begin{solution}
(1) ゲージ変換で $\theta(\gamma)\mapsto\theta(\gamma)+\alpha(\partial\gamma)$。$\gamma$ はサイクルなので $\partial\gamma=0$、よって $\theta(\gamma)$ は不変で $W(\gamma)$ も不変。

(2) $\gamma-\gamma'=\partial c$、$c=\sum_pc_pp$（$c_p\in\Z$）とすると
\[
\theta(\gamma)-\theta(\gamma')=\theta(\partial c)=\sum_pc_p\theta_p\in2\pi\Z
\]
（平坦なので各 $\theta_p\in2\pi\Z$）。よって $W(\gamma)=W(\gamma')$。$H_1(T^d)$ は $[P_\mu]$ で生成されるので、任意のサイクル $\gamma$ は $\gamma\sim\sum_\mu n_\mu P_\mu$ と書け、$W(\gamma)=\prod_\mu W(P_\mu)^{n_\mu}$。すなわち平坦な配位のすべてのウィルソン・ループは $d$ 個のポリヤコフ・ループで決まる。

(3) すべての $\mu$ 方向リンクに $\theta_{\ell_\mu(x)}=\varphi_\mu/N$ を置く。プラケット変数は $\theta_{p_{\mu\nu}(x)}=\varphi_\mu/N+\varphi_\nu/N-\varphi_\mu/N-\varphi_\nu/N=0$ なので平坦であり、$P_\mu$ は $\mu$ 方向リンク $N$ 本からなるので $\theta(P_\mu)=\varphi_\mu$、$W(P_\mu)=e^{i\varphi_\mu}$。これは連続理論の定数ゲージ場 $A_\mu=\varphi_\mu/L_\mu$（$L_\mu$ は周期）の格子版で、局所的には純ゲージだが大域的には自明でない。

(4) 2つの配位 $\theta,\theta'$ の差 $\eta=\theta-\theta'$ を考える（$U(1)$ は可換なので差も配位であり、平坦性とウィルソン・ループは乗法的にふるまう）。$\eta$ は平坦で $e^{i\eta(P_\mu)}=1$。(2) より任意のサイクル $\gamma$ で $e^{i\eta(\gamma)}=1$。

原点 $0$ から各サイト $x$ への道（1鎖 $c_x$ で $\partial c_x=x-0$）を選び、$\alpha(x)=\eta(c_x)$（$\bmod 2\pi$）とおく。別の道 $c'_x$ をとっても $c_x-c'_x$ はサイクルなので $\eta(c_x)-\eta(c'_x)\in2\pi\Z$ で、$e^{i\alpha(x)}$ は道によらない。リンク $\ell$（始点 $y$、終点 $x$）について $c_y+\ell$ は $0$ から $x$ への道だから
\[
\alpha(x)\equiv\eta(c_y)+\eta_\ell=\alpha(y)+\eta_\ell\pmod{2\pi},
\]
すなわち $\eta_\ell\equiv\alpha(\partial\ell)$。よって $\theta_\ell\equiv\theta'_\ell+\alpha(\partial\ell)$ であり、$\theta$ は $\theta'$ をゲージ変換 $e^{i\alpha}$ で移したものである。

以上から、平坦な配位のゲージ同値類は $(W(P_\mu))_\mu$ と1対1に対応し、空間 $U(1)^d$ をなす。これは $\Hom(H_1(T^d),U(1))$（$H_1$ の基底の行き先を自由に指定する）であり、コホモロジー $H^1(T^d;U(1))$ とも書ける（第5, 6章）。有限温度の格子QCDでは時間方向のポリヤコフ・ループが閉じ込めの秩序変数になり、トーラス・コンパクト化ではこれらがゼロモード（ウィルソン・ライン）として低エネルギー理論に残る（第10章）。
\end{solution}

\subsection{胞体ホモロジーと閉曲面}

\begin{problem}[難易度 B]\label{ch04-prob-torus}
(1) 正方形の語 $aba^{-1}b^{-1}$ による $T^2$ の胞体構造（頂点1つ、辺2つ、面1つ）から $H_*(T^2)$ を計算せよ。\\
(2) 下図の $3\times3$ の格子を対角線で分けた三角形分割（周期的に同一視）について、頂点・辺・三角形の数を求め、$\chi(T^2)=0$ を確かめよ。（これが単体複体であることは認めてよい。）\\
(3) (2) で、すべての三角形を反時計回りに向き付けた和 $\Sigma=\sum\sigma$ がサイクルであることを示し、$H_2(T^2)=\Z\Sigma$ であることを示せ。さらに (2) とオイラー・ポアンカレの公式から $b_1=2$ を再確認せよ。
\begin{center}
\begin{tikzpicture}[scale=1.0]
  \fill[blue!6] (0,0) rectangle (3,3);
  \foreach \i in {0,...,3} { \draw (\i,0) -- (\i,3); \draw (0,\i) -- (3,\i); }
  \foreach \i in {0,1,2} \foreach \j in {0,1,2} \draw (\i,\j) -- (\i+1,\j+1);
  \foreach \i in {0,...,3} \foreach \j in {0,...,3} {
    \pgfmathtruncatemacro\lab{mod(\i,3)+3*mod(\j,3)}
    \node[vtx] at (\i,\j) {};
    \node[font=\scriptsize,fill=white,inner sep=0.5pt] at (\i-0.15,\j+0.18) {$\lab$};
  }
  \draw[very thick,red!70!black,mid arrow] (0,0) -- (3,0);
  \draw[very thick,red!70!black,mid arrow] (0,3) -- (3,3);
  \draw[very thick,green!50!black,mid arrows] (0,0) -- (0,3);
  \draw[very thick,green!50!black,mid arrows] (3,0) -- (3,3);
  \draw[-{Stealth[length=2mm]},thick] (4.0,1.2) arc (-60:240:0.35);
  \node at (4.0,0.6) {\footnotesize 向き};
\end{tikzpicture}
\zucap{トーラスの9頂点の三角形分割。同じ番号の頂点、同じ色の縁は同一視される。各正方形を対角線で2つの三角形に分ける。}
\end{center}
\end{problem}

\begin{solution}
(1) 頂点：語 $aba^{-1}b^{-1}$ の正方形の4頂点はすべて同一視される（要約の図）。よって $C_0=\Z v$、$C_1=\Z a\oplus\Z b$、$C_2=\Z e$。$d_1a=v-v=0$、$d_1b=0$。$d_2e$ は語の指数和で $a+b-a-b=0$。すべての $d$ が $0$ なので
\[
H_0=\Z,\qquad H_1=\Z a\oplus\Z b\cong\Z^2,\qquad H_2=\Z e\cong\Z .
\]

(2) 頂点は $3\times3=9$ 個。辺は各頂点から「右」「上」「右上」に1本ずつ出るので $27$ 本。三角形は各正方形に2つで $18$ 個。$\chi=9-27+18=0$。

(3) 下の三角形を $[(i,j),(i+1,j),(i+1,j+1)]$、上の三角形を $[(i,j),(i+1,j+1),(i,j+1)]$（どちらも反時計回り）とする。各辺はちょうど2つの三角形の辺であり、平面上で反時計回りに向き付けた隣り合う2つの三角形は共通の辺を逆向きにたどる。周期的な同一視は平行移動なので向きを保ち、外周の辺についても同じことが成り立つ（例えば下端の辺 $(i,0)\to(i+1,0)$ は下の段の三角形で $+$、上端 $(i,3)\to(i+1,3)$ として上の段の三角形で $-$ に現れる）。よって $\partial\Sigma=0$。

逆に $c=\sum n_\sigma\sigma$ がサイクルなら、辺 $e$ を共有する2つの三角形 $\sigma,\sigma'$ について $\partial c$ の $e$ の係数は $\pm(n_\sigma-n_{\sigma'})$ なので $n_\sigma=n_{\sigma'}$。三角形は辺を通じてすべてつながっているので、全係数が等しく $c=n\Sigma$。$C_3=0$ なので $H_2=Z_2=\Z\Sigma$。

$b_0=1$、$b_2=1$、$\chi=0$ から $b_1=b_0+b_2-\chi=2$ で (1) と一致する。$\Sigma$ はトーラス全体を表す\term{基本類}{fundamental class}で、周期境界条件の格子上での積分 $\int_{T^2}$ はこの2サイクル上の和にあたる。
\end{solution}

\begin{problem}[難易度 B]\label{ch04-prob-klein}
(1) 正方形の語 $abab^{-1}$ によるクラインの壺 $K$ の胞体構造について、頂点がすべて1点に同一視されることを確かめ、$H_*(K)$ と $\chi(K)$ を求めよ。\\
(2) $H_1(K)$ のねじれ元を幾何学的に説明せよ。\\
(3) $K$ の胞体構造上の $U(1)$ 格子ゲージ場（各1胞体に角度 $\theta_a,\theta_b$）で、2胞体のホロノミーが自明（$e^{i\theta(d_2e)}=1$）なものを平坦という。平坦な配位のウィルソン・ループ $W(a)=e^{i\theta_a}$、$W(b)=e^{i\theta_b}$ の取りうる値を求めよ。
\end{problem}

\begin{solution}
(1) 正方形の頂点を左下から反時計回りに $P_0,P_1,P_2,P_3$ とする。語の各文字は辺 $P_0P_1$、$P_1P_2$、$P_2P_3$、$P_3P_0$ に対応し、指数 $+1$ なら $P_k\to P_{k+1}$、$-1$ なら逆向きが文字の向き。$a$：$P_0\to P_1$ と $P_2\to P_3$ より $P_0\sim P_2$、$P_1\sim P_3$。$b$：$P_1\to P_2$ と $P_0\to P_3$（$b^{-1}$ が $P_3\to P_0$）より $P_1\sim P_0$、$P_2\sim P_3$。よって4頂点はすべて同一視され、頂点は1つ $v$。

胞体鎖複体は $C_2=\Z e\to C_1=\Z a\oplus\Z b\to C_0=\Z v$ で、$d_1=0$、$d_2e=a+b+a-b=2a$。$d_2$ は単射なので $H_2=0$。
\[
H_1=\frac{\Z a\oplus\Z b}{\Z\,2a}\cong\Z_2\oplus\Z,\qquad H_0=\Z .
\]
$\chi=1-2+1=0$（$=1-1+0$）。$H_2=0$ は、$K$ が向き付け不可能で基本類をもたないことを表す（トーラスでは $d_2e=0$ だった）。

(2) $a$ は境界ではないサイクルだが、$2a=\partial e$ は2胞体全体の境界である。正方形を貼り合わせたとき、$e$ の縁は $a$ を同じ向きに2回たどり、$b$ は逆向きに1回ずつたどって打ち消す。したがって「$a$ を2周した輪」は曲面全体を張る膜の縁になるが、$a$ を1周した輪はどんな膜の縁にもならない。実係数では $a$ は $0$ になり $b_1=1$ であり、ねじれは実数の（微分形式の）手法では見えない（第5, 6章）。

(3) 平坦条件は $\theta(d_2e)=2\theta_a\in2\pi\Z$、すなわち $W(a)^2=1$、$W(a)=\pm1$。$\theta_b$ は制限されない。よって平坦な配位のウィルソン・ループは $(W(a),W(b))\in\{\pm1\}\times U(1)$。ゲージ変換は頂点が1つなので $\theta(\partial\ell)=0$ で自明に作用し、平坦な配位の同値類の空間は $\Z_2\times U(1)=\Hom(H_1(K),U(1))$ である。ねじれ $\Z_2$ に対応する $W(a)=-1$ は、場の強さが至るところ $0$ なのに物理的に観測できる離散的なウィルソン・ラインであり、連続的に $W(a)=1$ へ変形できない。
\end{solution}

\begin{problem}[難易度 C]\label{ch04-prob-hexagon}
正六角形の辺に、反時計回りに語を読んでラベルと向きを付け、同じラベルの辺を矢印の向きがそろうように貼り合わせて得られる閉曲面を考える。
\begin{center}
\begin{tikzpicture}[scale=1.0]
  \foreach \xs/\word in {0/{0/a/1,1/b/1,2/c/1,3/a/-1,4/b/-1,5/c/-1}, 5.5/{0/a/1,1/b/1,2/c/1,3/a/1,4/b/1,5/c/1}}{
    \begin{scope}[xshift=\xs cm]
      \fill[blue!6] (240:1.5) \foreach \k in {1,...,5} { -- ({240+60*\k}:1.5) } -- cycle;
      \foreach \k/\lab/\s in \word {
        \pgfmathsetmacro\angA{240+60*\k}
        \pgfmathsetmacro\angB{240+60*(\k+1)}
        \ifnum\s=1
          \draw[thick,mid arrow] (\angA:1.5) -- (\angB:1.5);
        \else
          \draw[thick,mid arrow] (\angB:1.5) -- (\angA:1.5);
        \fi
        \node at ({(\angA+\angB)/2}:1.62) [anchor={(\angA+\angB)/2+180}] {$\lab$};
      }
      \node[left] at (240:1.6) {\footnotesize $P_0$};
      \node[right] at (300:1.6) {\footnotesize $P_1$};
    \end{scope}
  }
  \node at (0,-2.2) {(1) $abca^{-1}b^{-1}c^{-1}$};
  \node at (5.5,-2.2) {(2) $abcabc$};
\end{tikzpicture}
\zucap{六角形の頂点を左下から反時計回りに $P_0,\dots,P_5$ とし、$k$ 番目の辺を $P_kP_{k+1}$（$P_6=P_0$）とする。}
\end{center}
それぞれについて、頂点の同値類を求めて胞体ホモロジー $H_*$ とオイラー標数を計算し、閉曲面の分類定理（連結な閉曲面は $\Sigma_g$ か $N_h$ のいずれかに同相で、$H_2\cong\Z$ なら $\Sigma_g$、$H_2=0$ なら $N_h$）を用いて、どの曲面かを決定せよ。
\end{problem}

\begin{solution}
ラベル $x$ の辺が指数 $+1$ ならその辺の $x$ は $P_k\to P_{k+1}$ 向き、$-1$ なら $P_{k+1}\to P_k$ 向きである。同じラベルの辺を貼ると、始点どうし・終点どうしが同一視される。

(1) $a$：$P_0\to P_1$ と $P_4\to P_3$ より $P_0\sim P_4$、$P_1\sim P_3$。$b$：$P_1\to P_2$ と $P_5\to P_4$ より $P_1\sim P_5$、$P_2\sim P_4$。$c$：$P_2\to P_3$ と $P_0\to P_5$ より $P_2\sim P_0$、$P_3\sim P_5$。よって頂点は
\[
v=\{P_0,P_2,P_4\},\qquad w=\{P_1,P_3,P_5\}
\]
の2つである。$a$ は $v\to w$、$b$ は $w\to v$、$c$ は $v\to w$ なので
\[
d_1a=w-v,\qquad d_1b=v-w,\qquad d_1c=w-v,\qquad d_2e=a+b+c-a-b-c=0 .
\]
$H_2=\Z e\cong\Z$。$\operatorname{im}d_1=\Z(w-v)$ なので $H_0\cong\Z$。$d_1(xa+yb+zc)=(x-y+z)(w-v)$ より $\ker d_1$ は階数2で基底 $a+b$、$b+c$。$\operatorname{im}d_2=0$ なので $H_1\cong\Z^2$。$\chi=2-3+1=0$。$H_2\cong\Z$ より向き付け可能で $2-2g=0$、$g=1$：\textbf{トーラス}である。

(2) $a$：$P_0\to P_1$ と $P_3\to P_4$ より $P_0\sim P_3$、$P_1\sim P_4$。$b$：$P_1\to P_2$ と $P_4\to P_5$ より $P_2\sim P_5$。$c$：$P_2\to P_3$ と $P_5\to P_0$ で新しい同一視はない。頂点は $u_0=\{P_0,P_3\}$、$u_1=\{P_1,P_4\}$、$u_2=\{P_2,P_5\}$ の3つで
\[
d_1a=u_1-u_0,\quad d_1b=u_2-u_1,\quad d_1c=u_0-u_2,\qquad d_2e=2a+2b+2c .
\]
$d_2$ は単射なので $H_2=0$。$d_1(xa+yb+zc)=(z-x)u_0+(x-y)u_1+(y-z)u_2$ が $0$ となるのは $x=y=z$ のときで $\ker d_1=\Z(a+b+c)$。よって
\[
H_1=\Z(a+b+c)/\Z\,2(a+b+c)\cong\Z_2 .
\]
$\operatorname{im}d_1$ は階数2なので $H_0\cong\Z$。$\chi=3-3+1=1$。$H_2=0$ より向き付け不可能で $2-h=1$、$h=1$：\textbf{実射影平面} $\RP^2$ である。

\begin{note}
(2) は六角形の中心に関する対蹠点を同一視しているので、$\RP^2=D^2/(x\sim-x,\ x\in\partial D^2)$ の構成そのものである。頂点の同一視は「ある頂点から出発し、隣の辺の貼り合わせ先へ移る」ことを繰り返して追跡するとよい。
\end{note}
\end{solution}

\begin{problem}[難易度 B]\label{ch04-prob-genus}
(1) 種数 $g\ge1$ の向き付け可能閉曲面 $\Sigma_g$ は、$4g$ 角形を語 $a_1b_1a_1^{-1}b_1^{-1}\cdots a_gb_ga_g^{-1}b_g^{-1}$ で貼り合わせて得られ、頂点はすべて1点に同一視される（認めてよい）。$H_*(\Sigma_g)$ と $\chi(\Sigma_g)$ を求めよ。\\
(2) $N_h$（$h\ge1$）は $2h$ 角形を語 $a_1a_1a_2a_2\cdots a_ha_h$ で貼り合わせたもので、頂点は1点になる（認めてよい）。$H_*(N_h)$ と $\chi(N_h)$ を求めよ。\\
(3) 閉じた弦の摂動論で、種数 $g$ の世界面の寄与は弦の結合定数 $g_s$ の $g_s^{-\chi(\Sigma_g)}$ 倍に比例する。$g=0,1,2$ の次数を書き、なぜ $\chi$ が現れるかを、ポリヤコフ作用に付け加える項 $\lambda\cdot\frac1{4\pi}\int\sqrt{h}\,R\,d^2\sigma=\lambda\chi$（ガウス・ボンネ、第8章）から説明せよ（$g_s=e^\lambda$）。向き付け不可能な世界面（$N_h$）をもつ弦理論ではどうなるか。
\end{problem}

\begin{solution}
(1) $C_0=\Z v$、$C_1=\bigoplus_i(\Z a_i\oplus\Z b_i)$、$C_2=\Z e$。頂点が1つなので $d_1=0$。各 $a_i$, $b_i$ は語の中に指数 $+1$ と $-1$ で1回ずつ現れるので $d_2e=0$。よって
\[
H_0=\Z,\qquad H_1=\Z^{2g},\qquad H_2=\Z,\qquad\chi(\Sigma_g)=1-2g+1=2-2g .
\]
$g=0$（$S^2$）でも $\chi=2$ で公式は正しい。

(2) $C_0=\Z v$、$C_1=\Z^h$、$C_2=\Z e$、$d_1=0$、$d_2e=2(a_1+\dots+a_h)$。$d_2$ は単射なので $H_2=0$。$C_1$ の基底を $s=a_1+\dots+a_h,\ a_2,\dots,a_h$ にとりかえると（変換行列は行列式 $1$ の整数行列）
\[
H_1=\frac{\Z s\oplus\Z a_2\oplus\dots\oplus\Z a_h}{\Z\,2s}\cong\Z_2\oplus\Z^{h-1},\qquad H_0=\Z,\qquad\chi(N_h)=1-h+1=2-h .
\]
$h=1$ で $\RP^2$（$H_1=\Z_2$）、$h=2$ でクラインの壺（$H_1=\Z\oplus\Z_2$）となり、問題 \ref{ch04-prob-hexagon}, \ref{ch04-prob-klein} と一致する。

(3) $g=0$：$g_s^{-2}$（球面、樹木レベル）、$g=1$：$g_s^{0}$（トーラス、1ループ）、$g=2$：$g_s^{2}$（2ループ）。一般に $g_s^{2g-2}$。

経路積分の重み $e^{-S}$ に $\lambda\chi$ が加わると、各世界面は $e^{-\lambda\chi}=g_s^{-\chi}$ 倍される。ガウス・ボンネの定理により $\frac1{4\pi}\int\sqrt hR=\chi$ は計量によらない位相不変量なので、この項は世界面の局所的な力学に影響せず、トポロジーに応じた重みだけを与える。$\chi=2-2g$ は取っ手を1つ増やすごとに $2$ 減り、「弦が1回分かれて再びつながる」（相互作用頂点2つ）ごとに $g_s^2$ が付くことと整合する。

向き付け不可能な世界面を許す理論（I 型弦理論など）では、$N_h$ の寄与 $g_s^{-\chi(N_h)}=g_s^{h-2}$ も現れる。例えば $\RP^2$（クロスキャップ1つ）は $g_s^{-1}$ で、球面と（境界のある）円板 $\chi=1$ の間の次数になる。クロスキャップ2つ（クラインの壺）はトーラスと同じ $g_s^0$ である。
\end{solution}

\begin{problem}[難易度 C]\label{ch04-prob-rpn}
$\RP^n$ は各次元 $0\le k\le n$ に1つずつ胞体 $e^k$ をもち、$e^k$ の接着写像は二重被覆 $q\colon S^{k-1}\to\RP^{k-1}$ である。\\
(1) 胞体境界の係数 $d_k$ は、合成 $S^{k-1}\xrightarrow{q}\RP^{k-1}\to\RP^{k-1}/\RP^{k-2}\cong S^{k-1}$ の写像度である。この合成は、赤道 $S^{k-2}$ をつぶすと上半球と下半球それぞれの上で同相写像になり、下半球上のものは上半球上のものに対蹠写像 $A$ を合成したものになる。これから $d_k=\pm(1+\deg A)=\pm(1+(-1)^k)$ を導け（$k\ge1$、$S^{k-1}$ 上の対蹠写像の写像度は $(-1)^k$）。\\
(2) $H_k(\RP^n)$ を求めよ。\\
(3) $SO(3)\cong\RP^3$ のホモロジーを書き、$H_1(SO(3))=\Z_2$ がスピノルの存在（$\pi_1(SO(3))=\Z_2$、第2章）とフレヴィッチの定理で整合することを確かめよ。また、ネマティック液晶の秩序変数空間 $\RP^2$ は $\pi_2(\RP^2)=\Z$（ヘッジホッグ欠陥、第3章）をもつが $H_2(\RP^2)$ はどうか。
\end{problem}

\begin{solution}
(1) 合成 $\psi\colon S^{k-1}\to S^{k-1}$ は、赤道を1点（基点）につぶした $S^{k-1}/S^{k-2}\cong S^{k-1}\vee S^{k-1}$ を経由する。上半球 $D_+$ 上の写像を $h$（$D_+/\partial D_+\cong S^{k-1}$ への同相、写像度 $\pm1$。向きを選んで $+1$ とする）とすると、下半球の点 $x$ は $q(x)=q(-x)$ なので $\psi(x)=h(-x)$、すなわち下半球上では $h\circ A$。写像度は各半球からの寄与の和であり（局所写像度の和、あるいは $S^{k-1}\vee S^{k-1}$ のホモロジー）
\[
\deg\psi=\deg h+\deg(h\circ A)=1+\deg A=1+(-1)^k .
\]
（$A$ は $\R^k$ の $k$ 個の座標の符号反転なので $\deg A=(-1)^k$。）よって $d_k=0$（$k$ 奇数）、$d_k=\pm2$（$k$ 偶数、$k\ge2$）。

(2) 胞体鎖複体は
\[
0\to\Z\xrightarrow{d_n}\cdots\to\Z\xrightarrow{\ 2\ }\Z\xrightarrow{\ 0\ }\Z\xrightarrow{\ 2\ }\Z\xrightarrow{\ 0\ }\Z\to0
\]
（右端から $C_0,C_1,C_2,\dots$。$d_1=0$、$d_2=2$、$d_3=0$、…）。
\begin{itemize}
\item $k=0$：$H_0=\Z/\operatorname{im}d_1=\Z$。
\item $0<k<n$、$k$ 奇数：$\ker d_k=\Z$（$d_k=0$）、$\operatorname{im}d_{k+1}=2\Z$ なので $H_k=\Z_2$。
\item $0<k<n$、$k$ 偶数：$d_k=2$ は単射なので $H_k=0$。
\item $k=n$：$n$ 奇数なら $d_n=0$ で $H_n=\Z$、$n$ 偶数なら $d_n=2$ で $H_n=0$。
\end{itemize}
$\chi(\RP^n)=\sum_{k=0}^n(-1)^k$ は $n$ 偶数で $1$、奇数で $0$。

(3) $H_*(SO(3))=H_*(\RP^3)=(\Z,\Z_2,0,\Z)$。$\pi_1(SO(3))=\Z_2$ は可換なのでフレヴィッチの定理より $H_1=\pi_1=\Z_2$ で一致する。$2\pi$ 回転のループは $H_1$ の生成元で、$4\pi$ 回転は境界（ホモロジー的に自明）である。$H_3=\Z$ は $SO(3)$ が向き付け可能な3次元多様体であることを表す。

$\RP^2$ では $H_2=0$ である。ヘッジホッグを囲む球面からの写像 $f\colon S^2\to\RP^2$ は、$\RP^2$ の $H_2$ が $0$ なのでホモロジーでは自明に見える。$\pi_2$ の情報はホモロジーでは失われうる（$f$ は被覆 $S^2\to\RP^2$ を通して持ち上がり、$S^2\to S^2$ の写像度としては見える）。ホモトピー群とホモロジー群は一般に異なる。
\end{solution}

\begin{problem}[難易度 A]\label{ch04-prob-cpn}
(1) $\CP^n=e^0\cup e^2\cup\dots\cup e^{2n}$ を用いて $H_*(\CP^n)$ と $\chi(\CP^n)$ を求めよ。\\
(2) $\lambda_0=1,\lambda_1,\dots,\lambda_n$ を相異なる絶対値1の複素数とし、$f[z_0:\dots:z_n]=[\lambda_0z_0:\dots:\lambda_nz_n]$ とする（$\CP^{N-1}$ 模型の大域的 $U(1)^{N-1}$ 変換）。$f$ の不動点をすべて求め、その数が $\chi(\CP^n)$ と一致することを確かめよ（$f\simeq\id$ であること、および各不動点の指数が $+1$ であることを認めて、レフシェッツ・ホップの定理と比べよ）。
\end{problem}

\begin{solution}
(1) 胞体は偶数次元にしかないので、$d_k\colon C_k\to C_{k-1}$ は $C_k$ か $C_{k-1}$ のどちらかが $0$ で、すべて $0$。よって $H_{2k}(\CP^n)=\Z$（$0\le k\le n$）、奇数次は $0$。$\chi=n+1$。

(2) $[z]$ が不動点 $\iff$ ある $c\in\C^\times$ で $\lambda_jz_j=cz_j$（すべての $j$）。$z_j\ne0$ となる $j$ では $\lambda_j=c$ であり、$\lambda_j$ は相異なるので $z_j\ne0$ となる $j$ はただ1つ。よって不動点は座標点
\[
[1:0:\dots:0],\ [0:1:0:\dots:0],\ \dots,\ [0:\dots:0:1]
\]
の $n+1$ 個で、$\chi(\CP^n)=n+1$ と一致する。$\lambda_j=e^{i\alpha_j}$ として $\alpha_j\to0$ と連続に動かせば $f\simeq\id$、よって $L(f)=L(\id)=\chi=n+1$ であり、指数 $+1$ の不動点が $n+1$ 個あることと整合する。物理的には、$\CP^{N-1}$ 模型の真空多様体上で、大域的対称性の一般の元で不変な真空はちょうど $N=\chi(\CP^{N-1})$ 個ある（局所化の考え方、第11章）。
\end{solution}

\begin{problem}[難易度 B]\label{ch04-prob-lens}
$p\ge2$、$q$ を $p$ と互いに素な整数とし、$\zeta=e^{2\pi i/p}$ とする。レンズ空間 $L(p,q)$ は $S^3\subset\C^2$ を $(z,w)\mapsto(\zeta z,\zeta^qw)$ で生成される $\Z_p$ の作用（自由な作用）で割った空間である。$L(p,q)$ は各次元 $0\le k\le3$ に1つずつ胞体 $e^k$ をもち、$d_2(e^2)=p\,e^1$ であることを認める。\\
(1) $H_*(L(p,q))$ と $\chi$ を求めよ。$L(2,1)=\RP^3$ の場合と比べよ。\\
(2) フレヴィッチの定理と $\pi_1(L(p,q))=\Z_p$ が整合することを確かめよ。\\
(3) $L(p,q)$ 上の平坦な $U(1)$ ゲージ場のウィルソン・ラインの取りうる値を求めよ（胞体構造上の格子ゲージ場として考えてよい）。
\end{problem}

\begin{solution}
(1) $C_k=\Z e^k$（$0\le k\le3$）。0胞体が1つなので $d_1=0$。$d_2=p$。$d_2\circ d_3=0$ で $d_2$ は単射なので $d_3=0$。胞体鎖複体は
\[
0\to\Z\xrightarrow{\ 0\ }\Z\xrightarrow{\ p\ }\Z\xrightarrow{\ 0\ }\Z\to0
\]
で
\[
H_0=\Z,\qquad H_1=\Z/p\Z=\Z_p,\qquad H_2=\ker(p)=0,\qquad H_3=\Z .
\]
$\chi=1-1+1-1=0$（自由な $\Z_p$ 作用による $p$ 重被覆 $S^3\to L(p,q)$ から $\chi(S^3)=p\,\chi(L)$、$0=p\cdot0$ とも整合）。$p=2$、$q=1$ では作用は $(z,w)\mapsto(-z,-w)$、すなわち対蹠写像で $L(2,1)=\RP^3$、前問の $(\Z,\Z_2,0,\Z)$ と一致する。$H_3=\Z$ はレンズ空間が向き付け可能な閉3次元多様体であることを表す。

(2) $S^3$ は単連結で $\Z_p$ は自由に作用するので $S^3\to L(p,q)$ は普遍被覆、$\pi_1=\Z_p$（第2章）。$\Z_p$ は可換なので $H_1=\pi_1^{\mathrm{ab}}=\Z_p$ で一致する。

(3) ゲージ場は1胞体 $e^1$ に角度 $\theta$ を置いたもの、平坦条件は $\theta(d_2e^2)=p\theta\in2\pi\Z$。よってウィルソン・ライン $W=e^{i\theta}$ は $W^p=1$、すなわち $W\in\{1,\zeta,\dots,\zeta^{p-1}\}$。ゲージ変換は（頂点が1つなので）自明に作用し、平坦な $U(1)$ 場の同値類は $\Hom(H_1,U(1))=\Hom(\Z_p,U(1))\cong\Z_p$ の $p$ 個である。連続的なモジュライはなく、ねじれから来る離散的なウィルソン・ラインだけがある。このような離散ウィルソン・ラインは、$\Z_p$ による商をとるコンパクト化でゲージ対称性を破る道具になる（第10章）。
\end{solution}

\subsection{完全列とマイヤー・ヴィートリス}

\begin{problem}[難易度 B]\label{ch04-prob-mvsphere}
$n\ge1$ とし、$S^n$ を $U=S^n\setminus\{S\}$（南極を除く）、$V=S^n\setminus\{N\}$（北極を除く）で覆う。\\
(1) $U,V$ は可縮で、$U\cap V\simeq S^{n-1}$ であることを説明せよ。\\
(2) 被約ホモロジーのマイヤー・ヴィートリス完全列から $\tilde H_k(S^n)\cong\tilde H_{k-1}(S^{n-1})$ を導き、$\tilde H_k(S^0)$ から出発して $\tilde H_k(S^n)$ を求めよ。\\
(3) この被覆は、ディラック・モノポールのゲージ場を北半球・南半球で別々に書く構成（第5, 7章）と同じである。$H_2(S^2)\cong\Z$ の生成元を、モノポールを囲む閉曲面の言葉で説明せよ。
\end{problem}

\begin{solution}
(1) 立体射影で $U\cong\R^n$、$V\cong\R^n$ で可縮。$U\cap V=S^n\setminus\{N,S\}\cong\R^n\setminus\{0\}\simeq S^{n-1}$（立体射影で $S\mapsto0$、$N\mapsto\infty$、第1章）。

(2) $U\cap V\ne\emptyset$ なので被約ホモロジーの完全列
\[
\tilde H_k(U)\oplus\tilde H_k(V)\to\tilde H_k(S^n)\xrightarrow{\ \delta\ }\tilde H_{k-1}(U\cap V)\to\tilde H_{k-1}(U)\oplus\tilde H_{k-1}(V)
\]
で両端が $0$ なので、$\delta$ は同型 $\tilde H_k(S^n)\cong\tilde H_{k-1}(S^{n-1})$。$S^0$ は2点なので $H_0(S^0)=\Z^2$、$\tilde H_0(S^0)=\Z$、他は $0$。これを $n$ 回くり返すと
\[
\tilde H_k(S^n)\cong\tilde H_{k-n}(S^0)=\begin{cases}\Z&k=n\\0&k\neq n\end{cases}
\]
特に $\tilde H_0(S^n)=0$（$n\ge1$、連結）であり、$H_k(S^n)$ は $k=0,n$ で $\Z$、他は $0$。

(3) $\delta$ の構成を逆にたどると、$H_2(S^2)$ の生成元（基本類 $[S^2]$）は、北半球の鎖 $u$ と南半球の鎖 $v$ の和 $u+v$ で、共通部分（赤道の近く）で $\partial u=-\partial v$ が赤道の円周 $S^1$ を表すものである。モノポールを囲む閉曲面上の磁束 $\int_{S^2}F$ は、北と南のパッチの寄与 $\int_uF+\int_vF=\oint_{\partial u}A_N+\oint_{\partial v}A_S=\oint_{\text{赤道}}(A_N-A_S)$ に分かれ、赤道上の遷移関数（ゲージ変換）の巻き付き数で決まる。ホモロジーの $\delta\colon H_2(S^2)\cong H_1(S^1)$ は、この「2次元の磁束を赤道上の1次元の量に帰着させる」構造の位相的な骨組みである。ディラックの量子化条件 $\int F\in2\pi\Z$ は第5, 7章で扱う。
\end{solution}

\begin{problem}[難易度 C]\label{ch04-prob-mvtorus}
$T^2=S^1\times S^1$（座標 $(\phi_1,\phi_2)$）を、$U=\{\phi_1\ne0\}$、$V=\{\phi_1\ne\pi\}$ で覆う。\\
(1) $U$, $V$ は円筒 $(\text{開区間})\times S^1\simeq S^1$、$U\cap V$ は2つの円筒の非交和 $\simeq S^1\sqcup S^1$ であることを確かめよ。\\
(2) マイヤー・ヴィートリス完全列を書き、写像 $\Phi\colon H_k(U\cap V)\to H_k(U)\oplus H_k(V)$ を $k=0,1$ について具体的に求め、$H_*(T^2)$ を計算せよ。\\
(3) $H_1(T^2)\cong\Z^2$ の2つの生成元がそれぞれ完全列のどこから来るかを説明せよ。
\end{problem}

\begin{solution}
(1) $U=(0,2\pi)\times S^1$、$V=(-\pi,\pi)\times S^1$ はどちらも円筒で、$\phi_2$ 方向の円周 $\beta$ に変位レトラクトする。$U\cap V=\bigl((0,\pi)\sqcup(\pi,2\pi)\bigr)\times S^1$ は2つの円筒 $W_1\sqcup W_2$ で、それぞれ円周 $\beta_1,\beta_2$ に変位レトラクトする。

(2) $H_1(U\cap V)=\Z\beta_1\oplus\Z\beta_2$、$H_1(U)=\Z\beta_U$、$H_1(V)=\Z\beta_V$、$H_0(U\cap V)=\Z p_1\oplus\Z p_2$、$H_0(U)=\Z$、$H_0(V)=\Z$、$H_2$ はいずれも $0$（円筒・円周の次元）。完全列は
\[
0\to H_2(T^2)\xrightarrow{\delta}\Z^2\xrightarrow{\Phi_1}\Z^2\xrightarrow{\Psi_1}H_1(T^2)\xrightarrow{\delta}\Z^2\xrightarrow{\Phi_0}\Z^2\xrightarrow{\Psi_0}H_0(T^2)\to0 .
\]
$\beta_i$ は $U$ でも $V$ でも $\phi_2$ 方向の一周なので $\beta_i\mapsto\beta_U$、$\beta_i\mapsto\beta_V$。よって
\[
\Phi_1(\beta_1)=\Phi_1(\beta_2)=(\beta_U,-\beta_V),\qquad\Phi_0(p_1)=\Phi_0(p_2)=(1,-1).
\]
$\ker\Phi_1=\Z(\beta_1-\beta_2)$、$\operatorname{im}\Phi_1=\Z(\beta_U,-\beta_V)$、$\ker\Phi_0=\Z(p_1-p_2)$、$\operatorname{im}\Phi_0=\Z(1,-1)$。
\begin{itemize}
\item $H_2(T^2)\cong\ker\Phi_1=\Z$。
\item $H_0(T^2)\cong\operatorname{coker}\Phi_0=\Z^2/\Z(1,-1)\cong\Z$。
\item $H_1$：短完全列 $0\to\operatorname{coker}\Phi_1\to H_1(T^2)\to\ker\Phi_0\to0$ で、$\operatorname{coker}\Phi_1=\Z^2/\Z(\beta_U,-\beta_V)\cong\Z$、$\ker\Phi_0\cong\Z$。右端が自由なので分裂し $H_1(T^2)\cong\Z^2$。
\end{itemize}

(3) 一方の生成元は $\operatorname{coker}\Phi_1$ から来る $\Psi_1(\beta_U)=[\beta]$、すなわち $\phi_2$ 方向の円周で、$U$ の中にすでにある輪である。もう一方 $[\alpha]$ は $\phi_1$ 方向に一周する輪で、$U$ にも $V$ にも含まれない。$\alpha=u+v$（$u\subset U$、$v\subset V$ の弧）と分けると $\partial u$ は $W_1$ の点と $W_2$ の点の差になり、$\delta[\alpha]=\pm(p_1-p_2)$ が $\ker\Phi_0$ の生成元になる。つまり $\alpha$ は「2つの重なりを両方通って戻ってくる」ことで生まれる輪である。また $\delta\colon H_2(T^2)\to\ker\Phi_1$ は、基本類 $[T^2]$ を2つの重なりにある円周の差 $\beta_1-\beta_2$ に写す。
\end{solution}

\subsection{フレヴィッチ・写像度・ラージ $N$}

\begin{problem}[難易度 C]\label{ch04-prob-hurewicz}
平面から $n$ 点 $p_1,\dots,p_n$ を除いた $X=\R^2\setminus\{p_1,\dots,p_n\}$ を考える（$\R^3$ から平行な $n$ 本の直線を除いたものとホモトピー同値で、$n$ 本の細いソレノイドの外側の空間）。$X\simeq\bigvee^nS^1$ で $\pi_1(X)=F_n$（$p_j$ をまわるループ $x_j$ で生成される自由群）であることは認めてよい。\\
(1) $H_1(X)$ を求めよ。\\
(2) ソレノイド $j$ の磁束を $\Phi_j$ とする。電荷 $q$ の粒子がループ $\gamma$ に沿って得るアハラノフ・ボーム位相 $\exp(iq\oint_\gamma A)$ が $\gamma$ のホモロジー類だけで決まることを示し、$[\gamma]=\sum_jn_j[x_j]$ のときの値を書け。\\
(3) 非可換ゲージ群 $G=SU(2)$ の「非可換磁束」（平坦な接続で、$x_j$ に沿ったホロノミーが $g_j\in G$）を考える。ループ $\gamma=x_1x_2x_1^{-1}x_2^{-1}$ は $H_1(X)$ で $0$ であることを示し、$g_1=i\sigma_1$、$g_2=i\sigma_3$ のときそのホロノミーが $-\mathbf 1\ne\mathbf 1$ であることを確かめよ。この例から、ホモロジーとホモトピーの情報の違いについて何が言えるか。
\end{problem}

\begin{solution}
(1) フレヴィッチの定理より $H_1(X)\cong F_n^{\mathrm{ab}}=\Z^n$（$[x_1],\dots,[x_n]$ が基底）。直接には、$\bigvee^nS^1$ は頂点1つ・辺 $n$ 本のグラフで $b_1=n-1+1=n$（問題 \ref{ch04-prob-graph}）。

(2) $A$ は $X$ 上で平坦（$F=dA=0$）。ループ $\gamma,\gamma'$ がホモローグで $\gamma-\gamma'=\partial c$（$c$ は $X$ 内の2鎖）なら、ストークスの定理より
\[
\oint_\gamma A-\oint_{\gamma'}A=\int_cF=0 .
\]
したがって位相は $[\gamma]\in H_1(X)$ だけで決まる。$\oint_{x_j}A=\Phi_j$（$x_j$ は $p_j$ を正の向きに1周）なので、$[\gamma]=\sum n_j[x_j]$ なら
\[
\exp\Bigl(iq\oint_\gamma A\Bigr)=\exp\Bigl(iq\sum_jn_j\Phi_j\Bigr).
\]
$n_j$ は $\gamma$ が $p_j$ をまわる巻き付き数である。可換なゲージ理論では、平坦な場の観測量はホモロジーで決まる（より正確には $\Hom(H_1(X),U(1))$ の元）。

(3) $\gamma$ を1鎖とみなすと $x_1+x_2-x_1-x_2=0$ なので $[\gamma]=0\in H_1(X)$。一方、非可換の場合、平坦接続のホロノミーは基本群の準同型 $\pi_1(X)\to G$ を定め（基点での共役を除いて）、$\gamma$ のホロノミーは（道の積の順序の規約によらず）$g_1,g_2$ の交換子 $g_1g_2g_1^{-1}g_2^{-1}$ またはその共役・逆元である。
\[
g_1g_2=(i\sigma_1)(i\sigma_3)=-\sigma_1\sigma_3=i\sigma_2,\qquad g_2g_1=(i\sigma_3)(i\sigma_1)=-\sigma_3\sigma_1=-i\sigma_2 .
\]
よって $g_1g_2g_1^{-1}g_2^{-1}=(g_1g_2)(g_2g_1)^{-1}=(i\sigma_2)(-i\sigma_2)^{-1}$。$(-i\sigma_2)(i\sigma_2)=\sigma_2^2=\mathbf 1$ より $(-i\sigma_2)^{-1}=i\sigma_2$ なので
\[
g_1g_2g_1^{-1}g_2^{-1}=(i\sigma_2)(i\sigma_2)=-\sigma_2^2=-\mathbf 1 .
\]
$-\mathbf 1$ は $SU(2)$ の中心で共役によらないので、ホロノミーは確かに $-\mathbf 1\ne\mathbf 1$。

結論：ホモロジー $H_1=\pi_1^{\mathrm{ab}}$ は交換子の情報を捨てる。可換ゲージ理論の位相（AB 効果）は $H_1$ で完全に記述されるが、非可換な磁束（アリス弦や離散ゲージ理論の非可換エニオンなど）の物理は $\pi_1$ そのものを必要とする。2つの非可換磁束の周りを交換子の形にまわると、どちらのソレノイドも「正味0回」しかまわっていないのに、粒子の内部状態が変わる。
\end{solution}

\begin{problem}[難易度 D]\label{ch04-prob-degree}
(1) $n\ge1$ とする。$S^2=\C\cup\{\infty\}$ 上の写像 $f_n(z)=z^n$（$f_n(\infty)=\infty$）の写像度が $n$ であることを、局所写像度の公式で示せ。$z\mapsto z^{-n}$、$z\mapsto\bar z^n$ の写像度は何か。\\
(2) $f\colon S^2\to S^2$ のレフシェッツ数が $L(f)=1+\deg f$ であることを示し、$\deg f\ne-1$ なら $f$ は不動点をもつことを導け。不動点をもたない写像の例を挙げよ。\\
(3)（毛の生えた球の定理）$S^{2m}$ 上に零点のない連続な接ベクトル場 $v$ は存在しないことを、$F(x,t)=\cos(\pi t)\,x+\sin(\pi t)\,v(x)/\lvert v(x)\rvert$ を用いて示せ。$S^{2m+1}$ では零点のないベクトル場の例を挙げよ。\\
(4) 任意の連続写像 $f\colon\RP^{2m}\to\RP^{2m}$ は不動点をもつことを示せ。一方 $\RP^{2m+1}$ には不動点をもたない写像があることを示せ。
\end{problem}

\begin{solution}
(1) $z$ 平面の向きを $dx\wedge dy$（$z=x+iy$）で定める。正則関数 $f$ の実ヤコビ行列は、$f'=s+it$ と書くと（コーシー・リーマンの関係より）$\begin{pmatrix}s&-t\\t&s\end{pmatrix}$ で、$\det=s^2+t^2=\lvert f'\rvert^2\ge0$。$w=1$ は $f_n$ の正則値で、逆像は $1$ の $n$ 乗根 $z_k=e^{2\pi ik/n}$ の $n$ 個、各点で $f_n'(z_k)=nz_k^{n-1}\ne0$、符号は $+1$。よって $\deg f_n=n$。（$\infty$ の近くでも $\zeta=1/z$ 座標で $\zeta\mapsto\zeta^n$ と正則である。）

$z\mapsto z^{-n}$ も $S^2$ 上の正則写像（$0\leftrightarrow\infty$ を入れかえる）で、$w=1$ の逆像は $n$ 個、すべて符号 $+1$ なので写像度は $+n$。$S^1$ 上では $e^{i\theta}\mapsto e^{-in\theta}$ の巻き付き数は $-n$ だが、$S^2$ 上では正則写像は向きを保つので写像度は正である。$z\mapsto\bar z^n$ は $z\mapsto z^n$ と複素共役（鏡映、写像度 $-1$）の合成なので $\deg=-n$。$O(3)$ シグマ模型では $\bm n$ を立体射影した複素場 $w(z)$ が正則なら $\deg>0$（反正則なら $<0$）のバビー・スキルミオン（ベラヴィン・ポリヤコフ解）である。

(2) $H_0(S^2;\Q)=\Q$ 上で $f_*$ は恒等（連結空間の点を点に写す）、$H_2(S^2;\Q)=\Q$ 上で $\deg f$ 倍、他は $0$。よって $L(f)=1+\deg f$。$\deg f\ne-1$ なら $L(f)\ne0$ で、レフシェッツの定理より不動点がある。例：回転（$\deg=1$）は回転軸の2点、$z\mapsto z^n$ は $0,\infty$ などを固定する。不動点のない写像の例は対蹠写像 $x\mapsto-x$（$z\mapsto-1/\bar z$）で、$\deg=(-1)^3=-1$、$L=0$。

(3) $v$ が零点をもたないとし、$u=v/\lvert v\rvert$ とおく。$u(x)\perp x$、$\lvert u\rvert=\lvert x\rvert=1$ なので
\[
\lvert F(x,t)\rvert^2=\cos^2(\pi t)+\sin^2(\pi t)=1,
\]
$F$ は $S^{2m}\times I\to S^{2m}$ の連続写像で、$F(\cdot,0)=\id$、$F(\cdot,1)=-\id$（対蹠写像）。写像度のホモトピー不変性より $1=\deg(-\id)=(-1)^{2m+1}=-1$、矛盾。よって零点がある。（ポアンカレ・ホップの定理によれば零点の指数の和は $\chi(S^{2m})=2$ である。地球上で風速ゼロの地点が必ずある。）

$S^{2m+1}\subset\C^{m+1}$ では $v(z)=iz$、すなわち $(x_1,y_1,\dots)\mapsto(-y_1,x_1,\dots)$ は $\langle v(z),z\rangle_{\R}=\mathrm{Re}(\bar z\cdot iz)=\mathrm{Re}(i\lvert z\rvert^2)=0$ をみたす接ベクトル場で、$\lvert v\rvert=1$。これは $U(1)$ 作用 $z\mapsto e^{i\alpha}z$ の生成子（ホップ・ファイバーの方向）である。

(4) 問題 \ref{ch04-prob-rpn} より $H_k(\RP^{2m})$ は $k=0$ で $\Z$、$k$ 奇数 $<2m$ で $\Z_2$、他は $0$。$\Q$ 係数では $\Z_2\otimes\Q=0$ なので $H_k(\RP^{2m};\Q)=0$（$k\ge1$）、$H_0=\Q$。よって任意の $f$ について $L(f)=1\ne0$ で、不動点がある。

$\RP^{2m+1}=S^{2m+1}/\pm$ では、$J\colon\R^{2m+2}\to\R^{2m+2}$ を $J^2=-1$ の線形写像（$\C^{m+1}$ の $i$ 倍）とし、$f[x]=[Jx]$ とおく（$J(-x)=-Jx$ なので well-defined）。不動点 $[x]=[Jx]$ は $Jx=\lambda x$（$\lambda=\pm1$）を意味するが、$J^2x=\lambda^2x=x$ と $J^2x=-x$ が矛盾する。よって不動点はない。実際 $f\simeq\id$（$e^{\alpha J}$, $\alpha\in[0,\pi/2]$）で、$L(f)=\chi(\RP^{2m+1})=0$ と整合する。
\end{solution}

\begin{problem}[難易度 D]\label{ch04-prob-thooft}
$N\times N$ エルミート行列場 $\Phi$ の理論で、作用を $S=\frac{N}{\lambda}\int\tr\bigl(\frac12(\partial\Phi)^2+\frac13\Phi^3\bigr)$ とする（$\lambda$ はトフーフト結合定数）。伝播関数は $\langle\Phi_{ij}\Phi_{kl}\rangle\propto\frac{\lambda}{N}\delta_{il}\delta_{jk}$ で、行列の添字の流れを2本の線で表す（2重線表示）。このとき、ファインマン図（真空図）は頂点 $V$、伝播関数 $E$、添字の閉じたループ $F$ をもち、頂点ごとに $N/\lambda$、伝播関数ごとに $\lambda/N$、添字ループごとに（添字の和から）$N$ が付く。2重線図は、頂点を円板、伝播関数を帯としたリボン・グラフで、その境界の各円周（添字ループ）に円板を貼ると向き付け可能な閉曲面 $\Sigma$ の胞体分割（頂点 $V$、辺 $E$、面 $F$）が得られる（認めてよい）。\\
(1) 連結な真空図の $N,\lambda$ 依存性が $N^{\chi(\Sigma)}\lambda^{E-V}=N^{2-2g}\lambda^{L-1}$ となることを示せ（$g$ は $\Sigma$ の種数、$L=E-V+1$）。\\
(2) 3価の頂点2つを3本の伝播関数で結ぶ「$\theta$ 図」には、頂点での伝播関数の巡回順序の取り方により、2つの2重線図がある。半辺の置換を用いて、それぞれの添字ループの数 $F$ を数え、平面図（$g=0$）と非平面図（$g=1$）であることを確かめよ。\\
(3) $N\to\infty$ で $\lambda$ を固定すると平面図だけが残ることを説明し、弦の種数展開（問題 \ref{ch04-prob-genus}）との対応 $g_s\sim1/N$ を述べよ。
\begin{center}
\begin{tikzpicture}[scale=1.0]
  \node[vtx] (l) at (-1.4,0) {}; \node[vtx] (r) at (1.4,0) {};
  \draw[thick] (l) to[out=70,in=110] (r);
  \draw[thick] (l) -- (r);
  \draw[thick] (l) to[out=-70,in=-110] (r);
  \node at (0,-1.4) {$\theta$ 図（1本線）};
  \begin{scope}[xshift=5.2cm]
    \fill[blue!8] (0,0) ellipse (2.0 and 1.15);
    \fill[white] (-1.55,0.12) .. controls (-0.8,0.85) and (0.8,0.85) .. (1.55,0.12) -- cycle;
    \fill[white] (-1.55,-0.12) .. controls (-0.8,-0.85) and (0.8,-0.85) .. (1.55,-0.12) -- cycle;
    \draw[thick,red!70!black] (0,0) ellipse (2.0 and 1.15);
    \draw[thick,green!50!black] (-1.55,0.12) .. controls (-0.8,0.85) and (0.8,0.85) .. (1.55,0.12) -- cycle;
    \draw[thick,blue!70!black] (-1.55,-0.12) .. controls (-0.8,-0.85) and (0.8,-0.85) .. (1.55,-0.12) -- cycle;
    \node at (0,1.45) {\footnotesize 添字ループ1（外側）};
    \node[green!50!black] at (0,0.4) {\footnotesize 2};
    \node[blue!70!black] at (0,-0.4) {\footnotesize 3};
    \node at (0,-1.6) {2重線表示（平面的）：$F=3$};
  \end{scope}
\end{tikzpicture}
\zucap{$\theta$ 図とその平面的な2重線表示。色の付いた帯が伝播関数（2本の添字線にはさまれた部分）、3本の閉曲線が添字ループである。$V-E+F=2-3+3=2=\chi(S^2)$。}
\end{center}
\end{problem}

\begin{solution}
(1) 図の因子は
\[
\Bigl(\frac N\lambda\Bigr)^V\Bigl(\frac\lambda N\Bigr)^EN^F=N^{V-E+F}\lambda^{E-V}.
\]
$V-E+F$ は閉曲面 $\Sigma$ の胞体分割の交代和なので、オイラー・ポアンカレの公式より $\chi(\Sigma)=2-2g$（$\Sigma$ は連結な向き付け可能閉曲面）。また $E-V=L-1$（$L=E-V+1$ はグラフのループ数、問題 \ref{ch04-prob-feynman}）。よって $N^{2-2g}\lambda^{L-1}$。なお $L=E-V+1=F-1+2g$ なので、ループ数は添字ループと種数の両方に分配される。3価頂点では $2E=3V$ なので $\lambda^{E-V}=\lambda^{V/2}$。

(2) 頂点 $A$ の半辺を $a_1,a_2,a_3$、頂点 $B$ の半辺を $b_1,b_2,b_3$ とし、伝播関数 $i$ は $a_i$ と $b_i$ を結ぶとする（対合 $\alpha\colon a_i\leftrightarrow b_i$）。各頂点で半辺を反時計回りに並べた巡回置換を $\sigma$ とすると、リボン・グラフの境界の円周（添字ループ）は置換 $\sigma\circ\alpha$ の巡回に1対1に対応する（帯を渡って反対側の頂点へ行き、そこで次の帯に移る）。

平面図：平面上に描くと、$A$ で $a_1\to a_2\to a_3$ が反時計回りなら $B$ では鏡像の順 $b_1\to b_3\to b_2$。$\sigma\alpha$ を追うと
\[
a_1\mapsto b_3,\ b_3\mapsto a_1;\qquad a_2\mapsto b_1,\ b_1\mapsto a_2;\qquad a_3\mapsto b_2,\ b_2\mapsto a_3
\]
（例：$a_1\xrightarrow{\alpha}b_1\xrightarrow{\sigma}b_3$、$b_3\xrightarrow{\alpha}a_3\xrightarrow{\sigma}a_1$）。巡回は3つで $F=3$、$\chi=2-3+3=2$、$g=0$。図の3つの添字ループに対応する。

非平面図：$B$ でも $b_1\to b_2\to b_3$ とする。$\sigma\alpha$ を順に追うと（例：$a_1\xrightarrow{\alpha}b_1\xrightarrow{\sigma}b_2$）
\[
a_1\mapsto b_2\mapsto a_3\mapsto b_1\mapsto a_2\mapsto b_3\mapsto a_1
\]
となり、長さ6の巡回1つで $F=1$。$\chi=2-3+1=0$、$g=1$：この2重線図はトーラス上にしか交差なしに描けない。$N$ 依存性は平面図が $N^2$、非平面図が $N^0$ で、$1/N^2$ だけ抑制される。（向き付け可能なら $V-E+F$ は偶数なので、$F$ は奇数 $1$ か $3$ に限られ、これで尽きている。）

(3) 固定した $\lambda$ で、真空図の和は
\[
\log Z=\sum_{g\ge0}N^{2-2g}\mathcal F_g(\lambda)=N^2\mathcal F_0(\lambda)+\mathcal F_1(\lambda)+N^{-2}\mathcal F_2(\lambda)+\cdots
\]
と整理され、$N\to\infty$ では種数 $0$、すなわち球面に描ける平面図だけが主要項になる（各 $\mathcal F_g$ はループ数について無限和）。これを閉じた弦の種数展開（問題 \ref{ch04-prob-genus}）$\sum_gg_s^{2g-2}\mathcal F_g^{\mathrm{string}}$ と比べると $g_s\sim1/N$ と同定され、2重線図の曲面 $\Sigma$ が弦の世界面にあたる（トフーフト；AdS/CFT 対応の出発点）。どちらの展開でも次数を決めているのはオイラー標数 $\chi(\Sigma)$ である。
\end{solution}

%================================================
\chapter{微分形式とド・ラーム・コホモロジー}\label{ch:05}
% 18問
%================================================

\section{物理からの動機}

\subsection*{ポテンシャルはいつ存在するか}
マクスウェル方程式の半分 $\nabla\cdot\bm B=0$、$\nabla\times\bm E+\partial_t\bm B=0$ は、微分形式で書くと $dF=0$ の一つの式になる（本章）。$F$ が閉形式であることから「$F=dA$ となるベクトル・ポテンシャル $A$ がある」と結論するのが普通であるが、これは\textbf{局所的に}しか正しくない。$d\omega=0$ をみたす形式（\term{閉形式}{closed form}）が $\omega=d\alpha$（\term{完全形式}{exact form}）と書けるかどうかは、空間の形に依存する。その障害を測るのが
\[
H^p_{\mathrm{dR}}(M)=\frac{\{\text{閉 $p$ 形式}\}}{\{\text{完全 $p$ 形式}\}}
\]
で、これを\term{ド・ラーム・コホモロジー}{de Rham cohomology}という。可縮な領域では $H^p_{\mathrm{dR}}=0$（$p\ge1$）であり（\term{ポアンカレの補題}{Poincar\'e lemma}）、物理で「ポテンシャルが存在する」と言うときに暗黙に使っているのはこの事実である。

\subsection*{ディラック・モノポール}
原点に磁荷 $g$ をおくと、$\R^3\setminus\{0\}$ 上で $\bm B=g\,\hat{\bm r}/4\pi r^2$、すなわち
\[
F=\frac{g}{4\pi}\sin\theta\,d\theta\wedge d\phi ,\qquad dF=0,\qquad \int_{S^2}F=g
\]
である。もし $\R^3\setminus\{0\}$ 全体で $F=dA$ となる $A$ があれば、ストークスの定理により $\int_{S^2}F=\int_{\partial S^2}A=0$ となって矛盾する。したがって $F$ は閉形式だが完全形式ではなく、$[F]\neq0\in H^2_{\mathrm{dR}}(\R^3\setminus\{0\})\cong\R$。実際には、北側と南側の二つの領域でそれぞれポテンシャル $A_N$, $A_S$ をとり、重なりでゲージ変換で貼り合わせる（下図）。貼り合わせが量子力学的に矛盾しないという条件が、\term{ディラックの量子化条件}{Dirac quantization condition} $qg\in2\pi\Z$ である。
\begin{center}
\begin{tikzpicture}[scale=1.15]
  \fill[red!8] (-1.6,0) arc (180:360:1.6) -- cycle;
  \fill[blue!8] (-1.6,0) arc (180:0:1.6) -- cycle;
  \begin{scope}
    \clip (0,0) circle (1.6);
    \fill[violet!18] (-1.7,-0.45) rectangle (1.7,0.45);
  \end{scope}
  \draw[thick] (0,0) circle (1.6);
  \draw[gray] (-1.6,0) arc (180:360:1.6 and 0.35);
  \draw[gray,dashed] (1.6,0) arc (0:180:1.6 and 0.35);
  \foreach \a in {0,30,...,330}{
    \draw[-{Stealth[length=1.6mm]},orange!80!black] (\a:0.25) -- (\a:0.75);
  }
  \node[vtx] at (0,0) {};
  \node at (0.25,0.15) {\footnotesize $g$};
  \draw[ultra thick,red!70!black] (0,-0.1) -- (0,-2.6);
  \node[red!70!black,right] at (0.05,-2.4) {\footnotesize $A_N$ のディラック弦};
  \draw[very thick,blue!70!black,dashed] (0,0.1) -- (0,2.6);
  \node[blue!70!black,right] at (0.05,2.4) {\footnotesize $A_S$ のディラック弦};
  \node[blue!70!black] at (-2.4,1.0) {\footnotesize $U_N$：$A_N$};
  \node[red!70!black] at (-2.4,-1.0) {\footnotesize $U_S$：$A_S$};
  \node[violet] at (2.75,0) {\footnotesize $U_N\cap U_S$};
  \draw[violet,->] (2.1,0) -- (1.65,0.05);
\end{tikzpicture}
\zucap{ディラック・モノポール。$A_N=\frac{g}{4\pi}(1-\cos\theta)d\phi$ は北側 $U_N$（負の $z$ 軸を除く）で、$A_S=-\frac{g}{4\pi}(1+\cos\theta)d\phi$ は南側 $U_S$（正の $z$ 軸を除く）で滑らか。それぞれが特異になる半直線がディラック弦である。赤道付近の重なりで $A_N-A_S=\frac{g}{2\pi}d\phi$。}
\end{center}

\subsection*{アハロノフ・ボーム効果}
無限に長いソレノイドの外側では $\bm B=0$、すなわち $F=dA=0$ であるが、ソレノイドを囲む閉曲線に沿って $\oint A=\Phi$（内部の磁束）$\neq0$ である。外側の領域 $\simeq S^1$ では $A$ は閉形式だが完全形式でなく、その類 $[A]\in H^1_{\mathrm{dR}}$ が電荷 $q$ の粒子の干渉縞をずらす位相 $e^{iq\Phi}$ を決める。\textbf{場の強さが局所的にゼロでも、コホモロジー類として観測可能な情報が残る}。

\subsection*{保存量と位相的項}
電荷 $Q=\int_{S^2}{*F}$ は、源のない領域では $d{*F}=0$ なので、$S^2$ を連続的に変形しても変わらない。一般に「閉形式をサイクルの上で積分した量」はサイクルのホモロジー類と形式のコホモロジー類だけで決まる（ストークスの定理）。また $\theta$ 項 $F\wedge F=d(A\wedge dA)$ のような\term{位相的項}{topological term}は局所的には完全形式なので運動方程式を変えないが、大域的には完全でないことがあり、量子論で意味をもつ（第6, 8章）。

\subsection*{カルツァ・クライン理論のゼロモード}
時空 $M_4\times K$（$K$ はコンパクトな内部空間）上の $p$ 形式ゲージ場を4次元から見ると、質量ゼロ（massless）の場の数は $K$ 上の\term{調和形式}{harmonic form}の数、すなわち $K$ の\term{ベッチ数}{Betti number}で決まる（\textbf{ホッジの定理}）。弦理論のコンパクト化で現れる質量ゼロのスカラー（モジュライ）や $U(1)$ ゲージ場の数は、内部空間のトポロジーで決まるのである（第10章）。

\begin{physbox}[：本章の辞書]
\begin{center}
\small
\begin{tabular}{p{0.46\linewidth}p{0.44\linewidth}}
物理 & 数学\\ \hline
場の強さ $F$、カレント $j$ & 2形式、1形式（${*j}$ は3形式）\\
ビアンキ恒等式、電荷保存 & $d^2=0$\\
ポテンシャルの存在（局所） & ポアンカレの補題\\
モノポール、AB効果 & 閉形式だが完全でない形式、$H^2$, $H^1\neq0$\\
ゲージ変換 $A\to A+d\alpha$ & コホモロジー類を変えない\\
ガウスの法則・保存電荷 & ストークスの定理、周期\\
ディラックの量子化 & 周期 $\int F\in2\pi\Z$（整数類）\\
KK の massless モード & 調和形式、ベッチ数 $b_p$\\
クーロン・ゲージ & ホッジ分解
\end{tabular}
\end{center}
\end{physbox}

\section{基本事項のまとめ}

\subsection{微分形式と外微分}
$M$ を $n$ 次元の滑らかな多様体、$x^\mu$ を局所座標とする。

\begin{mydef}{微分形式}{ch05-forms}
$M$ 上の\term{$p$ 形式}{$p$-form}とは、局所的に
\[
\omega=\frac1{p!}\,\omega_{\mu_1\cdots\mu_p}(x)\,dx^{\mu_1}\wedge\cdots\wedge dx^{\mu_p}
\]
（$\omega_{\mu_1\cdots\mu_p}$ は完全反対称）と書かれ、座標変換で反対称共変テンソルとして変換するものである。全体を $\Omega^p(M)$ と書く。\term{外積}{wedge product} $\wedge$ は $\alpha\wedge\beta=(-1)^{pq}\beta\wedge\alpha$（$\alpha\in\Omega^p$, $\beta\in\Omega^q$）をみたす。\term{外微分}{exterior derivative} $d\colon\Omega^p\to\Omega^{p+1}$ は
\[
d\omega=\frac1{p!}\,\partial_\nu\omega_{\mu_1\cdots\mu_p}\,dx^\nu\wedge dx^{\mu_1}\wedge\cdots\wedge dx^{\mu_p}
\]
で、座標によらない。
\end{mydef}

\begin{myprop}{外微分の性質}{ch05-dprop}
\begin{enumerate}
\item $d^2=0$（偏微分の対称性と $dx^\mu\wedge dx^\nu$ の反対称性による）。
\item ライプニッツ則 $d(\alpha\wedge\beta)=d\alpha\wedge\beta+(-1)^p\alpha\wedge d\beta$（$\alpha\in\Omega^p$）。
\item 滑らかな写像 $f\colon M\to N$ の\term{引き戻し}{pullback} $f^*$ は $f^*(\alpha\wedge\beta)=f^*\alpha\wedge f^*\beta$、$f^*d=d f^*$ をみたす。
\end{enumerate}
\end{myprop}

\begin{myex}{$\R^3$ のベクトル解析}{ch05-r3}
$\R^3$（ユークリッド計量、向き $dx\wedge dy\wedge dz$）では、$d\bm S=(dy\wedge dz,\ dz\wedge dx,\ dx\wedge dy)$、$d^3x=dx\wedge dy\wedge dz$ として
\[
f\ \xrightarrow{\ d\ }\ \nabla f\cdot d\bm x,\qquad
\bm v\cdot d\bm x\ \xrightarrow{\ d\ }\ (\nabla\times\bm v)\cdot d\bm S,\qquad
\bm v\cdot d\bm S\ \xrightarrow{\ d\ }\ (\nabla\cdot\bm v)\,d^3x .
\]
$d^2=0$ は $\nabla\times\nabla f=0$、$\nabla\cdot(\nabla\times\bm v)=0$ にほかならない（問題 \ref{ch05-prob-r3}）。
\end{myex}

\subsection{内部積・リー微分・ストークスの定理}
ベクトル場 $X$ に対し、\term{内部積}{interior product} $\iota_X\colon\Omega^p\to\Omega^{p-1}$ を $(\iota_X\omega)_{\mu_2\cdots\mu_p}=X^\nu\omega_{\nu\mu_2\cdots\mu_p}$ で定める。$X$ の流れ $\varphi_s$ による\term{リー微分}{Lie derivative} $L_X\omega=\frac{d}{ds}\varphi_s^*\omega\big|_{s=0}$ について、次の\textbf{カルタンの公式}が成り立つ：
\[
L_X=d\,\iota_X+\iota_X\,d .
\]
例えば $\R^3$ で $L_{\bm v}(d^3x)=d(\iota_{\bm v}d^3x)=d(\bm v\cdot d\bm S)=(\nabla\cdot\bm v)\,d^3x$（体積の変化率は発散）。

\begin{mythm}{ストークスの定理}{ch05-stokes}
$M$ を境界つきの向きづけられた $n$ 次元多様体、$\omega$ をコンパクト台の $(n-1)$ 形式とすると
\[
\int_M d\omega=\int_{\partial M}\omega .
\]
ただし $\partial M$ の向きは「外向き法線を先頭に置いて $M$ の向きになる」ように定める（\textbf{外向き法線先頭の規約}）。
\end{mythm}
特に、境界のない閉多様体 $M$ 上で完全形式の積分はゼロであり、$\int_M$ は $H^n_{\mathrm{dR}}(M)$ 上の線形写像を定める。

\begin{center}
\begin{tikzpicture}[scale=1.0]
  \fill[blue!8] plot[smooth cycle,tension=0.8] coordinates {(-2.2,0) (-1.4,1.3) (0.2,1.5) (1.6,0.9) (1.9,-0.5) (0.6,-1.4) (-1.3,-1.2)};
  \draw[thick,blue!70!black,mid arrow] plot[smooth cycle,tension=0.8] coordinates {(-2.2,0) (-1.3,-1.2) (0.6,-1.4) (1.9,-0.5) (1.6,0.9) (0.2,1.5) (-1.4,1.3)};
  \fill[white] (-0.3,0.1) circle (0.45);
  \draw[thick,red!70!black,mid arrow] (-0.3,0.1) ++(0:0.45) arc (0:-360:0.45);
  \node[vtx] at (-0.3,0.1) {};
  \node at (-0.3,-0.12) {\scriptsize $0$};
  \node[blue!70!black] at (1.15,0.35) {\small $V$};
  \node[blue!70!black] at (2.25,0.6) {\small $\Sigma_1$};
  \node[red!70!black] at (-0.3,0.82) {\small $S_\varepsilon$};
  \begin{scope}[xshift=4.6cm]
    \fill[green!8] plot[smooth cycle,tension=0.9] coordinates {(-0.9,-0.4) (0.1,0.9) (1.2,0.5) (1.0,-0.8)};
    \draw[thick,green!50!black,mid arrow] plot[smooth cycle,tension=0.9] coordinates {(-0.9,-0.4) (1.0,-0.8) (1.2,0.5) (0.1,0.9)};
    \node[green!50!black] at (1.55,0.75) {\small $\Sigma_2$};
    \node[vtx] at (-1.6,0.6) {};
    \node at (-1.6,0.35) {\scriptsize $0$};
  \end{scope}
\end{tikzpicture}
\zucap{ストークスの定理の使い方（断面図）。左：原点を囲む閉曲面 $\Sigma_1$ と小球面 $S_\varepsilon$ は、原点を含まない領域 $V$ の境界 $\partial V=\Sigma_1-S_\varepsilon$ をなす。右：原点を囲まない $\Sigma_2$ は $\R^3\setminus\{0\}$ 内の領域の境界。閉形式 $F$ の積分は $\int_{\Sigma_1}F=\int_{S_\varepsilon}F$、$\int_{\Sigma_2}F=0$（問題 \ref{ch05-prob-flux}）。}
\end{center}

\subsection{ホッジ作用素}
\begin{mydef}{ホッジ作用素}{ch05-hodge}
向きづけられた $n$ 次元（擬）リーマン多様体 $(M,g)$ の体積形式を $\mathrm{vol}=\sqrt{\lvert g\rvert}\,dx^1\wedge\cdots\wedge dx^n$ とする。$p$ 形式の内積 $\langle\ ,\ \rangle$ を $\langle dx^{\mu_1}\wedge\cdots\wedge dx^{\mu_p},dx^{\nu_1}\wedge\cdots\wedge dx^{\nu_p}\rangle=\det(g^{\mu_i\nu_j})$ で定め、\term{ホッジ作用素}{Hodge star} ${*}\colon\Omega^p\to\Omega^{n-p}$ を
\[
\alpha\wedge{*\beta}=\langle\alpha,\beta\rangle\,\mathrm{vol}
\]
で定める。成分では（$\varepsilon_{12\cdots n}=+1$）
\[
{*}(dx^{\mu_1}\wedge\cdots\wedge dx^{\mu_p})=\frac{\sqrt{\lvert g\rvert}}{(n-p)!}\,g^{\mu_1\nu_1}\cdots g^{\mu_p\nu_p}\,\varepsilon_{\nu_1\cdots\nu_p\rho_1\cdots\rho_{n-p}}\,dx^{\rho_1}\wedge\cdots\wedge dx^{\rho_{n-p}} .
\]
\end{mydef}

\begin{myprop}{${**}$ の符号}{ch05-starstar}
$s=\operatorname{sign}\det g$ とすると、$p$ 形式上で
\[
{**}=(-1)^{p(n-p)}\,s .
\]
ユークリッド計量（$s=+1$）の4次元では2形式上 ${**}=+1$ で、${*}$ の固有値 $\pm1$ に応じて\textbf{自己双対}・\textbf{反自己双対}に分解できる（インスタントン、第8章）。ローレンツ計量（$s=-1$）の4次元では2形式上 ${**}=-1$ であり、実の自己双対形式は存在しない（問題 \ref{ch05-prob-star}）。
\end{myprop}

本書のミンコフスキー空間では $x^\mu=(t,x,y,z)$、$\eta=\operatorname{diag}(-1,1,1,1)$、向きは $dt\wedge dx\wedge dy\wedge dz$（$\varepsilon_{0123}=+1$）とする。このとき
\[
{*}dt=-d^3x,\quad {*}dx^i=-dt\wedge dS_i,\quad {*}(dt\wedge dx^i)=-dS_i,\quad {*}dS_i=dt\wedge dx^i,\quad {*}d^3x=-dt
\]
である。ここで $dS_i=\tfrac12\varepsilon_{ijk}dx^j\wedge dx^k$。

\begin{mydef}{余微分とラプラシアン}{ch05-codiff}
$L^2$ 内積 $(\alpha,\beta)=\int_M\alpha\wedge{*\beta}$ に関する $d$ の形式的随伴を\term{余微分}{codifferential} $\delta$ という。$p$ 形式に対し
\[
\delta=(-1)^p\,{*}^{-1}d\,{*}\qquad\bigl(\text{リーマンの場合}=(-1)^{n(p+1)+1}{*d*}\bigr)
\]
であり、$\delta^2=0$。\term{ラプラシアン}{Laplacian}（ホッジ・ラプラシアン）を $\Delta=d\delta+\delta d$ と定め、$\Delta\omega=0$ をみたす形式を\textbf{調和形式}という。
\end{mydef}
\begin{note}
$\alpha\in\Omega^{p-1}$, $\beta\in\Omega^p$ とすると $d(\alpha\wedge{*\beta})=d\alpha\wedge{*\beta}+(-1)^{p-1}\alpha\wedge d{*\beta}$。閉多様体上で積分すると $(d\alpha,\beta)=(-1)^p\int\alpha\wedge d{*\beta}=\int\alpha\wedge{*}\bigl((-1)^p{*}^{-1}d{*\beta}\bigr)$ となり、上の $\delta$ を得る。平坦な $\R^n$ では関数に対し $\Delta f=\delta df=-\sum_i\partial_i^2f$（正の作用素）、1形式に対し $\delta(A_idx^i)=-\partial_iA_i$ である。ミンコフスキー空間では $\Delta=-\Box$、$\Box=-\partial_t^2+\nabla^2$。
\end{note}

\subsection{形式で書いたマクスウェル方程式}
$A=A_\mu dx^\mu$、$A_\mu=(-\varphi,\bm A)$、$\bm E=-\nabla\varphi-\partial_t\bm A$、$\bm B=\nabla\times\bm A$ とすると、$F=dA=\tfrac12F_{\mu\nu}dx^\mu\wedge dx^\nu$ は
\[
F=-E_i\,dt\wedge dx^i+B_i\,dS_i,\qquad {*F}=B_i\,dt\wedge dx^i+E_i\,dS_i
\]
である。カレント $j=j_\mu dx^\mu=-\rho\,dt+J_i\,dx^i$ に対し ${*j}=\rho\,d^3x-J_i\,dt\wedge dS_i$。

\begin{physbox}[：マクスウェル方程式（mostly plus、本書の規約）]
\[
dF=0\quad\Longleftrightarrow\quad\nabla\cdot\bm B=0,\ \ \nabla\times\bm E+\partial_t\bm B=0,
\]
\[
d{*F}={*j}\ \ (\Longleftrightarrow\ \delta F=j\ \Longleftrightarrow\ \partial_\mu F^{\nu\mu}=j^\nu)\quad\Longleftrightarrow\quad\nabla\cdot\bm E=\rho,\ \ \nabla\times\bm B-\partial_t\bm E=\bm J .
\]
作用は $S=\int\bigl(-\tfrac12F\wedge{*F}+A\wedge{*j}\bigr)$ で、$-\tfrac12F\wedge{*F}=\tfrac12(\bm E^2-\bm B^2)\,d^4x$（$d^4x=dt\wedge d^3x$）。$d^2=0$ から $d{*j}=0$（連続の式）が従う。時刻一定面 $\Sigma_t$ 上で ${*j}|_{\Sigma_t}=\rho\,d^3x$、${*F}|_{\Sigma_t}=\bm E\cdot d\bm S$、$F|_{\Sigma_t}=\bm B\cdot d\bm S$ なので
\[
Q_e=\int_{\Sigma_t}{*j}=\oint_{S^2_\infty}{*F},\qquad Q_m=\oint_{S^2}F .
\]
導出は問題 \ref{ch05-prob-maxwell}。
\end{physbox}

\subsection{ド・ラーム・コホモロジー}
\begin{mydef}{ド・ラーム・コホモロジー}{ch05-dR}
$Z^p(M)=\ker(d\colon\Omega^p\to\Omega^{p+1})$（閉形式）、$B^p(M)=d\Omega^{p-1}(M)$（完全形式）とし、
\[
H^p_{\mathrm{dR}}(M)=Z^p(M)/B^p(M)
\]
と定める。$d(\alpha\wedge\beta)$ の公式から、外積は $H^p\times H^q\to H^{p+q}$ を誘導する（コホモロジー環、第6章）。$b_p(M)=\dim H^p_{\mathrm{dR}}(M)$ をベッチ数という。
\end{mydef}
$H^0_{\mathrm{dR}}(M)$ は局所定数関数の空間で、$\dim H^0=$（連結成分の数）。また $p>n$ では $H^p=0$。ゲージ変換 $A\to A+d\alpha$ は $A$ のコホモロジー類を変えない（$A$ が閉のとき）。

\begin{mythm}{ポアンカレの補題}{ch05-poincare}
$U\subset\R^n$ が原点に関して星形（$x\in U\Rightarrow tx\in U$, $0\le t\le1$）ならば、$p\ge1$ で $H^p_{\mathrm{dR}}(U)=0$。具体的に、オイラー・ベクトル場 $X=x^\mu\partial_\mu$ を用いて
\[
(K\omega)(x)=\int_0^1 t^{p-1}\,(\iota_X\omega)(tx)\,dt\qquad(\omega\in\Omega^p)
\]
（$(\iota_X\omega)(tx)$ は係数を $tx$ で評価し、$dx$ はそのまま）とおくと、$p\ge1$ で
\[
\omega=dK\omega+Kd\omega
\]
が成り立つ。特に $d\omega=0$ なら $\omega=d(K\omega)$（問題 \ref{ch05-prob-poincare}）。
\end{mythm}
$K\omega$ は $\iota_X K\omega=0$、すなわち $x^\mu A_\mu=0$ をみたす。これは場の理論で使われる\textbf{フォック・シュウィンガー・ゲージ}である。

\begin{mythm}{ホモトピー不変性}{ch05-htpy}
滑らかな写像 $f,g\colon M\to N$ がホモトピックならば $f^*=g^*\colon H^p_{\mathrm{dR}}(N)\to H^p_{\mathrm{dR}}(M)$。したがってホモトピー同値な多様体は同じド・ラーム・コホモロジーをもつ。
\end{mythm}
\begin{note}
ホモトピー $H\colon M\times I\to N$ に対し、$\omega\in\Omega^p(N)$ から $(\mathcal K\omega)=\int_0^1\iota_{\partial_t}H^*\omega\,dt$ を作ると $g^*\omega-f^*\omega=d\mathcal K\omega+\mathcal Kd\omega$ となる（ポアンカレの補題と同じ形の\term{ホモトピー作用素}{homotopy operator}）。閉形式なら右辺は完全形式。
\end{note}
例えば $\R^n\setminus\{0\}\simeq S^{n-1}$（第1章）なので $H^*_{\mathrm{dR}}(\R^n\setminus\{0\})\cong H^*_{\mathrm{dR}}(S^{n-1})$、ソレノイドの外側 $\simeq S^1$、など。

\begin{mythm}{マイヤー・ヴィートリス完全列}{ch05-MV}
$M=U\cup V$（$U,V$ 開集合）とする。$\Omega^p(M)\to\Omega^p(U)\oplus\Omega^p(V)$, $\omega\mapsto(\omega|_U,\omega|_V)$ と $\Omega^p(U)\oplus\Omega^p(V)\to\Omega^p(U\cap V)$, $(\alpha,\beta)\mapsto\beta|_{U\cap V}-\alpha|_{U\cap V}$ から、長い完全列
\[
\cdots\to H^p(M)\to H^p(U)\oplus H^p(V)\to H^p(U\cap V)\xrightarrow{\ \partial^*\ }H^{p+1}(M)\to\cdots
\]
が得られる。$\partial^*$ は、1の分割 $\rho_U+\rho_V=1$（$\operatorname{supp}\rho_U\subset U$, $\operatorname{supp}\rho_V\subset V$）を用いて、閉形式 $\omega\in\Omega^p(U\cap V)$ に対し $U$ 上で $d(\rho_V\omega)$、$V$ 上で $-d(\rho_U\omega)$ に等しい $M$ 上の閉形式の類を対応させる（符号は規約による）。
\end{mythm}
第4章のホモロジーのマイヤー・ヴィートリス完全列と同じ構造であり、物理的には「北側・南側のパッチで書いた量を、重なりでの差（遷移関数）で貼り合わせる」操作そのものである。

\begin{myex}{基本的な計算結果}{ch05-results}
\begin{center}
\renewcommand{\arraystretch}{1.25}\small
\begin{tabular}{lll}
$M$ & $H^p_{\mathrm{dR}}(M)$ & 生成元\\ \hline
$\R^n$, 星形領域 & $\R$（$p=0$）、$0$（$p\ge1$） & $1$\\
$S^1$, $\R^2\setminus\{0\}$ & $\R$（$p=0,1$） & $1$, $\frac{1}{2\pi}d\theta$\\
$S^n$ ($n\ge1$) & $\R$（$p=0,n$）、他は $0$ & $1$, 体積形式／$\mathrm{Vol}(S^n)$\\
$\R^3\setminus\{0\}$ & $\R$（$p=0,2$） & $1$, $\frac{1}{4\pi}\sin\theta\,d\theta\wedge d\phi$\\
$\R^2\setminus\{a_1,\dots,a_k\}$ & $\R$（$p=0$）、$\R^k$（$p=1$） & $\frac1{2\pi}d\theta_i$\\
$T^n$ & $\R^{\binom np}$ & $dx^{i_1}\wedge\cdots\wedge dx^{i_p}$\\
$\Sigma_g$（種数 $g$ の閉曲面） & $\R,\ \R^{2g},\ \R$ & 第4章のホモロジーと双対
\end{tabular}
\end{center}
$T^n$ や $S^2\times S^2$ のような積には\textbf{キュネットの公式} $H^k(M\times N)\cong\bigoplus_{p+q=k}H^p(M)\otimes H^q(N)$ を用いる（$\pi_M^*\alpha\wedge\pi_N^*\beta$ が生成する。詳しくは第6章）。
\end{myex}

\subsection{ド・ラームの定理と周期}
滑らかな $p$ 次元サイクル $c$（例えば向きづけられた閉部分多様体）と閉 $p$ 形式 $\omega$ に対し、\term{周期}{period} $\int_c\omega$ はストークスの定理により $[\omega]\in H^p_{\mathrm{dR}}$ と $[c]\in H_p(M)$ だけで決まる：
\[
\int_c(\omega+d\alpha)=\int_c\omega,\qquad \int_{c+\partial b}\omega=\int_c\omega+\int_bd\omega=\int_c\omega .
\]
\begin{mythm}{ド・ラームの定理}{ch05-deRham}
積分のペアリングは同型
\[
H^p_{\mathrm{dR}}(M)\xrightarrow{\ \cong\ }H^p(M;\R)\cong\Hom(H_p(M),\R),\qquad[\omega]\mapsto\Bigl([c]\mapsto\int_c\omega\Bigr)
\]
を与える。特に、\textbf{閉形式が完全形式であるための必要十分条件は、すべての周期がゼロであること}。$b_p(M)$ は第4章のホモロジーのベッチ数（$H_p(M)$ の階数）と一致する。
\end{mythm}
すべての周期が $\Z$ に属する閉形式を\textbf{整数類}を代表するという。ディラックの量子化 $\int_\Sigma F\in2\pi\Z$ は「$F/2\pi$ は整数類」と言い換えられ、これが第1チャーン類 $c_1=[F/2\pi]$ である（第7, 8章）。ねじれ（torsion）はド・ラーム・コホモロジーでは見えない（第6章の普遍係数定理）。

\begin{center}
\begin{tikzpicture}[scale=1.0]
  \draw[thick] (0,0) ellipse (2.6 and 1.25);
  \draw[thick] (-1.1,0.08) .. controls (-0.5,-0.35) and (0.5,-0.35) .. (1.1,0.08);
  \draw[thick] (-0.9,-0.05) .. controls (-0.4,0.25) and (0.4,0.25) .. (0.9,-0.05);
  \draw[very thick,blue!70!black,mid arrow] (0,-0.03) ellipse (1.85 and 0.82);
  \node[blue!70!black] at (0,-1.05) {$a$};
  \draw[very thick,red!70!black,mid arrow] (1.9,-0.02) arc (90:270:0.24 and 0.6);
  \draw[very thick,red!70!black,dashed] (1.9,-1.22) arc (-90:90:0.24 and 0.6);
  \node[red!70!black] at (1.35,-0.75) {$b$};
  \begin{scope}[xshift=5cm,yshift=-1.1cm]
    \draw[thick,blue!70!black,mid arrow] (0,0) -- (2.2,0);
    \draw[thick,blue!70!black,mid arrow] (0,2.2) -- (2.2,2.2);
    \draw[thick,red!70!black,mid arrows] (0,0) -- (0,2.2);
    \draw[thick,red!70!black,mid arrows] (2.2,0) -- (2.2,2.2);
    \node[blue!70!black,below] at (1.1,0) {$a$};
    \node[red!70!black,right] at (2.2,1.1) {$b$};
    \node at (1.1,1.1) {$T^2$};
    \node[below left] at (0,0) {\scriptsize $(0,0)$};
    \node[above right] at (2.2,2.2) {\scriptsize $(1,1)$};
    \node at (1.1,-0.75) {\scriptsize $x$ 方向};
    \node[rotate=90] at (-0.45,1.1) {\scriptsize $y$ 方向};
  \end{scope}
\end{tikzpicture}
\zucap{トーラス $T^2=\R^2/\Z^2$ の基本サイクル $a$（$x$ 方向）と $b$（$y$ 方向）。右は正方形の辺を同一視した表示。閉1形式 $\omega$ の類は二つの周期 $\oint_a\omega$, $\oint_b\omega$ で決まる（問題 \ref{ch05-prob-torus}）。}
\end{center}

\subsection{ディラック・モノポール}
\begin{myex}{ディラック・モノポールのポテンシャル}{ch05-dirac}
$\R^3\setminus\{0\}$ 上の $F=\frac{g}{4\pi}\sin\theta\,d\theta\wedge d\phi$ に対し、
\begin{gather*}
A_N=\frac{g}{4\pi}(1-\cos\theta)\,d\phi\quad(U_N=\R^3\setminus\{\theta=\pi\}),\\
A_S=-\frac{g}{4\pi}(1+\cos\theta)\,d\phi\quad(U_S=\R^3\setminus\{\theta=0\})
\end{gather*}
はそれぞれの領域で滑らかで $dA_N=dA_S=F$ をみたす。重なりで
\[
A_N-A_S=\frac{g}{2\pi}\,d\phi=d\alpha,\qquad\alpha=\frac{g\phi}{2\pi}
\]
であり、これは $U(1)$ ゲージ変換 $A\to A+d\alpha$（本書の規約 $A\to gAg^{-1}+i\,g\,dg^{-1}$ で $g=e^{i\alpha}$）である。電荷 $q$ の波動関数は $\psi_N=e^{iq\alpha}\psi_S$ と変換するので、$e^{iqg\phi/2\pi}$ が $\phi\sim\phi+2\pi$ で一価であることを要求すると
\[
qg\in2\pi\Z .
\]
単位電荷 $q=1$ が存在すれば $g=\int_{S^2}F\in2\pi\Z$ で、本書の規約 $\int_\Sigma F\in2\pi\Z$ と一致する（問題 \ref{ch05-prob-dirac}）。$e^{i\alpha}$ を遷移関数とする $U(1)$ 主束としての定式化は第7章で扱う。
\end{myex}

\subsection{ホッジの定理}
\begin{mythm}{ホッジの定理}{ch05-hodgethm}
$M$ をコンパクトで向きづけられた境界のないリーマン多様体とする。
\begin{enumerate}
\item $\omega$ が調和 $\iff$ $d\omega=0$ かつ $\delta\omega=0$（$(\Delta\omega,\omega)=\lVert d\omega\rVert^2+\lVert\delta\omega\rVert^2$ による）。
\item \textbf{ホッジ分解}：$L^2$ 直交分解
\[
\Omega^p(M)=\mathcal H^p\oplus d\,\Omega^{p-1}(M)\oplus\delta\,\Omega^{p+1}(M)
\]
が成り立つ。ここで $\mathcal H^p$ は調和 $p$ 形式の空間。
\item 各コホモロジー類はただ一つの調和代表元をもち、$\mathcal H^p\cong H^p_{\mathrm{dR}}(M)$。特に $\dim\mathcal H^p=b_p(M)$ で、計量によらない。
\end{enumerate}
\end{mythm}
${*}$ は $\mathcal H^p$ を $\mathcal H^{n-p}$ に同型に写すので $b_p=b_{n-p}$（ポアンカレ双対性、第6章）。調和代表元は、各類の中で $L^2$ ノルム $\lVert\omega\rVert^2$ を最小にするものである：$\lVert h+d\alpha\rVert^2=\lVert h\rVert^2+\lVert d\alpha\rVert^2$。物理的には「固定したフラックス（周期）をもつ配位のうちエネルギー最小のもの」が調和形式である。

\begin{physbox}[：カルツァ・クライン理論とベッチ数]
$M_d\times K$（$K$ はコンパクトな $k$ 次元リーマン多様体）上の $p$ 形式ゲージ場 $C$（作用 $-\tfrac12\int dC\wedge{*dC}$）を
\[
C=\sum_{q}\sum_{i} c^{(i)}_q(x)\wedge\omega^{(i)}_{p-q}(y)
\]
と内部空間 $K$ の固有形式で展開すると、$\Delta_K\omega=m^2\omega$ の固有値 $m^2$ が $d$ 次元の質量の2乗を与える。\textbf{質量ゼロのモードは調和形式}に対応し、$d$ 次元の質量ゼロの $q$ 形式場の数は
\[
b_{p-q}(K)
\]
である（問題 \ref{ch05-prob-kk}, \ref{ch05-prob-kkgen}）。例：$K=S^1$ 上の5次元マクスウェル場 → 光子（$b_0=1$）＋スカラー（$b_1=1$、ウィルソン線）。ベッチ数は位相不変量なので、低エネルギーの場の数は内部空間の計量の詳細によらない。
\end{physbox}

\subsection{位相的項とガウス・ボンネの予告}
$U(1)$ では $F\wedge F=d(A\wedge F)$、非可換では
\[
\frac{1}{8\pi^2}\tr F\wedge F=d\,CS(A),\qquad CS(A)=\frac{1}{8\pi^2}\tr\Bigl(A\wedge dA-\frac{2i}{3}A\wedge A\wedge A\Bigr)
\]
（記法と規約）。このように局所的に完全な形式の積分は場の局所的変分に依存せず、トポロジーだけを拾う。$CS(A)$ のゲージ変換性と大きなゲージ変換、$\theta$ 真空は第8章で扱う。同様に、曲面の曲率2形式 $R$ も局所的には $d\omega$（$\omega$ はスピン接続）であり、$\frac{1}{2\pi}\int_\Sigma R=\chi(\Sigma)$（ガウス・ボンネの定理、第8章）となる。$S^2$ の場合は問題 \ref{ch05-prob-gb} で、モノポールと同じ構造であることを見る。

\section{問題}

\subsection{微分形式とマクスウェル方程式}

\begin{problem}[難易度 A]\label{ch05-prob-r3}
$\R^3$（ユークリッド計量、向き $d^3x=dx\wedge dy\wedge dz$）で、ベクトル場 $\bm v$ に1形式 $\bm v\cdot d\bm x=v_idx^i$ と2形式 $\bm v\cdot d\bm S=\iota_{\bm v}d^3x$ を対応させる。\\
(1) $df=\nabla f\cdot d\bm x$、$d(\bm v\cdot d\bm x)=(\nabla\times\bm v)\cdot d\bm S$、$d(\bm v\cdot d\bm S)=(\nabla\cdot\bm v)\,d^3x$ を示し、$d^2=0$ から $\nabla\times\nabla f=0$、$\nabla\cdot(\nabla\times\bm v)=0$ を導け。\\
(2) ${*}dx^i=dS_i$、${*}dS_i=dx^i$、${*}1=d^3x$ を確かめ、$\delta df=-\nabla^2f$ を示せ。\\
(3) カルタンの公式を用いて $L_{\bm v}d^3x=(\nabla\cdot\bm v)\,d^3x$ を示し、非圧縮流体（$\nabla\cdot\bm v=0$）の流れが体積を保つことの形式による表現であることを説明せよ。
\end{problem}

\begin{solution}
(1) $df=\partial_if\,dx^i$。次に
\[
d(v_idx^i)=\partial_jv_i\,dx^j\wedge dx^i=\sum_{j<i}(\partial_jv_i-\partial_iv_j)\,dx^j\wedge dx^i .
\]
$dx^j\wedge dx^i$ を $dS_k$（$\varepsilon_{kji}=1$）で書くと、例えば $dy\wedge dz=dS_1$ の係数は $\partial_yv_z-\partial_zv_y=(\nabla\times\bm v)_1$。他も同様で $d(\bm v\cdot d\bm x)=(\nabla\times\bm v)\cdot d\bm S$。最後に
\[
d(v_1dy\wedge dz+v_2dz\wedge dx+v_3dx\wedge dy)=(\partial_xv_1+\partial_yv_2+\partial_zv_3)\,dx\wedge dy\wedge dz .
\]
$d^2f=0$ は $(\nabla\times\nabla f)\cdot d\bm S=0$、$d^2(\bm v\cdot d\bm x)=0$ は $\nabla\cdot(\nabla\times\bm v)\,d^3x=0$ を意味する。

(2) ユークリッド計量では $\langle dx^i,dx^j\rangle=\delta^{ij}$。$dx\wedge dy\wedge dz=\mathrm{vol}$ なので $dx\wedge{*dx}=\mathrm{vol}$ をみたすのは ${*dx}=dy\wedge dz=dS_1$（$dx\wedge dS_2=dx\wedge dz\wedge dx=0$ などから他の成分は入らない）。同様に ${*}dS_1$：$dS_1\wedge{*dS_1}=\mathrm{vol}$ より ${*}dS_1=dx$。${*}1=\mathrm{vol}$ は定義そのもの。$n=3$ では ${**}=(-1)^{p(3-p)}=+1$ である。
$df=\partial_if\,dx^i$ に対し $\delta=(-1)^1{*}^{-1}d{*}$ を使うと
\[
{*df}=\partial_if\,dS_i,\quad d{*df}=\nabla^2f\,d^3x,\quad \delta df=-{*}^{-1}(\nabla^2f\,d^3x)=-\nabla^2f .
\]

(3) $d(d^3x)=0$（4形式は $\R^3$ でゼロ）なので $L_{\bm v}d^3x=d\iota_{\bm v}d^3x+\iota_{\bm v}d(d^3x)=d(\bm v\cdot d\bm S)=(\nabla\cdot\bm v)\,d^3x$。流れ $\varphi_s$ で運ばれる領域 $D_s=\varphi_s(D)$ の体積は $\int_{D_s}d^3x=\int_D\varphi_s^*d^3x$ で、その $s$ 微分は $\int_{D_s}L_{\bm v}d^3x=\int_{D_s}\nabla\cdot\bm v\,d^3x$。$\nabla\cdot\bm v=0$ なら体積は時間によらない。ハミルトン力学では位相空間の体積形式 $\omega^n/n!$ についての同じ計算がリウヴィルの定理になる。
\end{solution}

\begin{problem}[難易度 A]\label{ch05-prob-star}
(1) ユークリッド $\R^4$（向き $dx^1\wedge dx^2\wedge dx^3\wedge dx^4$）で ${*}(dx^1\wedge dx^2)=dx^3\wedge dx^4$、${*}(dx^1\wedge dx^3)=-dx^2\wedge dx^4$ を確かめ、2形式上 ${**}=1$ であることを示せ。2形式の空間（6次元）は自己双対・反自己双対の3次元ずつに分かれることを、基底を与えて示せ。\\
(2) ミンコフスキー空間（$\eta=\operatorname{diag}(-1,1,1,1)$、向き $dt\wedge dx\wedge dy\wedge dz$）で ${*}(dt\wedge dx)=-dy\wedge dz$、${*}(dy\wedge dz)=dt\wedge dx$ を示し、2形式上 ${**}=-1$ であることを確かめよ。また1形式上 ${**}=+1$ であることを ${*}dt$ を計算して確かめよ。\\
(3) (2) の状況で、実の2形式 $F\neq0$ が $F={*F}$ をみたすことはないことを示せ。$F=-E_idt\wedge dx^i+B_idS_i$ のとき $F-i{*F}$ が $\bm E+i\bm B$ だけで書け、${*}$ の固有値 $i$ の固有形式であることを示せ。
\end{problem}

\begin{solution}
(1) 定義 $\alpha\wedge{*\beta}=\langle\alpha,\beta\rangle\mathrm{vol}$ を使う。$\langle dx^1\wedge dx^2,dx^1\wedge dx^2\rangle=1$、他の基底との内積は $0$ なので、${*}(dx^1\wedge dx^2)=c\,dx^3\wedge dx^4$ で $dx^1\wedge dx^2\wedge dx^3\wedge dx^4=\mathrm{vol}$ より $c=1$。同様に $dx^1\wedge dx^3\wedge(-dx^2\wedge dx^4)=+dx^1\wedge dx^2\wedge dx^3\wedge dx^4$ なので ${*}(dx^1\wedge dx^3)=-dx^2\wedge dx^4$。一般式 ${**}=(-1)^{2\cdot2}\cdot(+1)=1$。直接にも ${*}(dx^3\wedge dx^4)=dx^1\wedge dx^2$（$dx^3\wedge dx^4\wedge dx^1\wedge dx^2=\mathrm{vol}$）である。
${**}=1$ なので $P_\pm=\frac12(1\pm{*})$ は射影で、$\Omega^2=\Omega^2_+\oplus\Omega^2_-$。基底は
\[
dx^1\wedge dx^2\pm dx^3\wedge dx^4,\qquad dx^1\wedge dx^3\mp dx^2\wedge dx^4,\qquad dx^1\wedge dx^4\pm dx^2\wedge dx^3
\]
（${*}(dx^1\wedge dx^4)=dx^2\wedge dx^3$ を用いた）で、それぞれ3次元。これが $\mathfrak{so}(4)\cong\mathfrak{su}(2)\oplus\mathfrak{su}(2)$ の分解に対応し、インスタントン $F=\pm{*F}$ の舞台になる（第8章）。

(2) $\langle dt\wedge dx,dt\wedge dx\rangle=\eta^{tt}\eta^{xx}=-1$。$\mathrm{vol}=dt\wedge dx\wedge dy\wedge dz$ なので $dt\wedge dx\wedge{*}(dt\wedge dx)=-\mathrm{vol}$ より ${*}(dt\wedge dx)=-dy\wedge dz$。$\langle dy\wedge dz,dy\wedge dz\rangle=1$ と $dy\wedge dz\wedge dt\wedge dx=+\mathrm{vol}$（2形式どうしの交換は符号なし）より ${*}(dy\wedge dz)=dt\wedge dx$。したがって
\[
{**}(dt\wedge dx)={*}(-dy\wedge dz)=-dt\wedge dx ,
\]
他の基底も同様で ${**}=-1=(-1)^{2\cdot2}\cdot(-1)$。
1形式：$\langle dt,dt\rangle=-1$ より $dt\wedge{*dt}=-\mathrm{vol}$、${*dt}=-dx\wedge dy\wedge dz$。また $\langle d^3x,d^3x\rangle=1$ と $d^3x\wedge dt=-dt\wedge d^3x$ より ${*}d^3x=-dt$。よって ${**}dt=-{*}d^3x=dt$、一般式 $(-1)^{1\cdot3}(-1)=+1$ と一致。

(3) $F={*F}$ なら $F={**F}=-F$ で $F=0$。よって実の自己双対2形式はない。${*}$ の固有値は ${*}^2=-1$ より $\pm i$ である。
${*F}=-E_i{*}(dt\wedge dx^i)+B_i{*}dS_i=E_idS_i+B_idt\wedge dx^i$。したがって
\[
F-i{*F}=-(E_i+iB_i)\,dt\wedge dx^i+(B_i-iE_i)\,dS_i=-(E_i+iB_i)\,dt\wedge dx^i-i(E_i+iB_i)\,dS_i ,
\]
は $\bm E+i\bm B$ だけで書ける。さらに ${*}(F-i{*F})={*F}+iF=i(F-i{*F})$。$\bm E+i\bm B$（リーマン・ジルバーシュタイン・ベクトル）は光子の一方のヘリシティに対応し、ローレンツ符号では自己双対性は複素化して初めて意味をもつ。
\end{solution}

\begin{problem}[難易度 B]\label{ch05-prob-maxwell}
ミンコフスキー空間（本書の規約）で $A=A_\mu dx^\mu$、$A_\mu=(-\varphi,\bm A)$、$\bm E=-\nabla\varphi-\partial_t\bm A$、$\bm B=\nabla\times\bm A$ とし、$j=j_\mu dx^\mu$、$j^\mu=(\rho,\bm J)$ とする。\\
(1) $F=dA=-E_idt\wedge dx^i+B_idS_i$ を示せ。\\
(2) $dF=dt\wedge(\nabla\times\bm E+\partial_t\bm B)\cdot d\bm S+(\nabla\cdot\bm B)\,d^3x$ を示し、$dF=0$ が斉次方程式であることを確かめよ。\\
(3) ${*F}$ と ${*j}$ を求め、$d{*F}={*j}$ が $\nabla\cdot\bm E=\rho$、$\nabla\times\bm B-\partial_t\bm E=\bm J$ と同値であることを示せ。\\
(4) $d{*j}=0$ が連続の式 $\partial_t\rho+\nabla\cdot\bm J=0$ であることを示し、それが $d^2=0$ から従うことを説明せよ。\\
(5) $-\tfrac12F\wedge{*F}=\tfrac12(\bm E^2-\bm B^2)\,d^4x$ を示せ。また $S=\int(-\tfrac12F\wedge{*F}+A\wedge{*j})$ を $A$ について変分すると $d{*F}={*j}$ が得られることを示せ。
\end{problem}

\begin{solution}
(1) $F_{\mu\nu}=\partial_\mu A_\nu-\partial_\nu A_\mu$ より $F_{0i}=\partial_tA_i+\partial_i\varphi=-E_i$、$F_{ij}=\partial_iA_j-\partial_jA_i=\varepsilon_{ijk}B_k$。したがって
\[
F=F_{0i}\,dt\wedge dx^i+\tfrac12F_{ij}dx^i\wedge dx^j=-E_i\,dt\wedge dx^i+B_k\,dS_k .
\]

(2) $d(-E_idt\wedge dx^i)=-\partial_jE_i\,dx^j\wedge dt\wedge dx^i=dt\wedge(\partial_jE_i\,dx^j\wedge dx^i)=dt\wedge(\nabla\times\bm E)\cdot d\bm S$（問題 \ref{ch05-prob-r3}(1)）。また
\[
d(B_kdS_k)=\partial_tB_k\,dt\wedge dS_k+(\nabla\cdot\bm B)\,d^3x .
\]
和をとれば主張の式。$dt\wedge dS_k$ と $d^3x$ は独立なので、$dF=0\iff\nabla\cdot\bm B=0$ かつ $\nabla\times\bm E+\partial_t\bm B=0$。

(3) 問題 \ref{ch05-prob-star} の結果 ${*}(dt\wedge dx^i)=-dS_i$、${*}dS_i=dt\wedge dx^i$ から
\[
{*F}=E_i\,dS_i+B_i\,dt\wedge dx^i .
\]
$j_\mu=\eta_{\mu\nu}j^\nu=(-\rho,\bm J)$ と ${*dt}=-d^3x$、${*}dx^i=-dt\wedge dS_i$（$dx\wedge(-dt\wedge dy\wedge dz)=\mathrm{vol}$ で確かめられる）から
\[
{*j}=\rho\,d^3x-J_i\,dt\wedge dS_i .
\]
(2) と同じ計算で（$-E\to B$、$B\to E$ と置き換える）
\[
d{*F}=dt\wedge(\partial_t\bm E-\nabla\times\bm B)\cdot d\bm S+(\nabla\cdot\bm E)\,d^3x .
\]
係数を比べて、$d{*F}={*j}\iff\nabla\cdot\bm E=\rho$ かつ $\partial_t\bm E-\nabla\times\bm B=-\bm J$。

(4) $d(\rho\,d^3x)=\partial_t\rho\,dt\wedge d^3x$、$d(J_idt\wedge dS_i)=\partial_kJ_i\,dx^k\wedge dt\wedge dS_i=-dt\wedge(\nabla\cdot\bm J)\,d^3x$（$dx^k\wedge dS_i=\delta^k_id^3x$）。よって
\[
d{*j}=(\partial_t\rho+\nabla\cdot\bm J)\,dt\wedge d^3x .
\]
$d{*F}={*j}$ が成り立てば $d{*j}=d^2{*F}=0$。すなわち電荷保存はマクスウェル方程式の可積分条件であり、カレントが保存しなければ方程式に解がない。

(5) $\alpha\wedge{*\alpha}=\langle\alpha,\alpha\rangle\mathrm{vol}$ で、$\langle dt\wedge dx^i,dt\wedge dx^j\rangle=-\delta^{ij}$、$\langle dS_i,dS_j\rangle=\delta_{ij}$、交差項はゼロ。よって $\langle F,F\rangle=-\bm E^2+\bm B^2$、$-\frac12F\wedge{*F}=\frac12(\bm E^2-\bm B^2)\,d^4x$。これは $-\frac14F_{\mu\nu}F^{\mu\nu}$ に等しい。
変分：$\alpha\wedge{*\beta}=\beta\wedge{*\alpha}$（内積の対称性）より $\delta(F\wedge{*F})=2\,d\delta A\wedge{*F}$。$d(\delta A\wedge{*F})=d\delta A\wedge{*F}-\delta A\wedge d{*F}$ を使い、表面項を落とすと
\[
\delta S=\int\bigl(-d\delta A\wedge{*F}+\delta A\wedge{*j}\bigr)=\int\delta A\wedge\bigl(-d{*F}+{*j}\bigr).
\]
任意の $\delta A$ に対しゼロなので $d{*F}={*j}$。成分では $\partial_\mu F^{\mu\nu}=-j^\nu$（mostly plus 特有の符号）である。
\end{solution}

\begin{problem}[難易度 B]\label{ch05-prob-bianchi}
(1) $U(1)$ ゲージ場について、$F=dA$ ならば $dF=0$（ビアンキ恒等式）であることを示せ。逆に、$\R^4$ 全体で $dF=0$ ならば $F=dA$ となる $A$ があることを、定理 \ref{thm:ch05-poincare} を用いて述べよ。\\
(2) 非可換ゲージ場 $F=dA-iA\wedge A$ について、$DF:=dF-i(A\wedge F-F\wedge A)=0$ を示せ。\\
(3) 静的な磁場 $\bm B$ が $\R^3\setminus\{0\}$ 上で $\nabla\cdot\bm B=0$ をみたし、原点を囲む球面を通る磁束がゼロでないとする。このとき $\R^3\setminus\{0\}$ 全体で滑らかなベクトル・ポテンシャルは存在しないことを示せ。
\end{problem}

\begin{solution}
(1) $dF=d^2A=0$。成分では $\partial_{[\lambda}F_{\mu\nu]}=0$。逆に $\R^4$ は星形なので、閉2形式 $F$ に対しポアンカレの補題から $A=KF$ をとれば $F=dA$。

(2) 外微分は $d(A\wedge A)=dA\wedge A-A\wedge dA$（$A$ は1形式）なので
\[
dF=-i\,(dA\wedge A-A\wedge dA).
\]
一方 $A\wedge F-F\wedge A=A\wedge dA-iA^3-dA\wedge A+iA^3=A\wedge dA-dA\wedge A$（$A^3=A\wedge A\wedge A$）。よって
\[
DF=-i(dA\wedge A-A\wedge dA)-i(A\wedge dA-dA\wedge A)=0 .
\]
ゲージ変換 $F\to gFg^{-1}$ のもとで $DF\to g(DF)g^{-1}$ と共変である。

(3) $F=\bm B\cdot d\bm S$ とする（問題 \ref{ch05-prob-r3}）。$\nabla\cdot\bm B=0$ は $dF=0$。もし $\R^3\setminus\{0\}$ 全体で $\bm B=\nabla\times\bm A$、すなわち $F=dA$（$A=\bm A\cdot d\bm x$）なら、原点を囲む球面 $S^2$（境界なし）上でストークスの定理より
\[
\int_{S^2}\bm B\cdot d\bm S=\int_{S^2}dA=\int_{\partial S^2}A=0 ,
\]
仮定に反する。よって $[F]\neq0\in H^2_{\mathrm{dR}}(\R^3\setminus\{0\})$ であり、ポテンシャルは局所的にしか存在しない。モノポールの磁荷は「ビアンキ恒等式 $dF=0$ が成り立つ領域でも $F$ が完全とは限らない」ことの表れである。
\end{solution}

\subsection{ド・ラーム・コホモロジーの計算}

\begin{problem}[難易度 A]\label{ch05-prob-dtheta}
$X=\R^2\setminus\{0\}$ 上の1形式 $\omega=\dfrac{-y\,dx+x\,dy}{x^2+y^2}$ を考える。\\
(1) $d\omega=0$ を示し、局所的には $\omega=d\theta$（偏角）であることを確かめよ。\\
(2) 単位円周（反時計回り）に沿った積分を求め、$\omega$ は $X$ 上で完全形式でないことを示せ。\\
(3) 閉曲線 $\gamma\colon[0,2\pi]\to X$ の巻き付き数を $w(\gamma)=\frac1{2\pi}\oint_\gamma\omega$ と定める。$\gamma(t)=e^{2it}\bigl(1+\tfrac12e^{it}\bigr)$（複素数表示）について $w(\gamma)$ を求めよ。\\
(4) $X$ 上の任意の閉1形式 $\eta$ は、$c=\frac1{2\pi}\oint_{S^1}\eta$ と関数 $f$ を用いて $\eta=c\,\omega+df$ と書けることを示せ。ただし「$X$ 内の閉曲線は単位円周の $n$ 周とホモトピックで、ホモトピックな閉曲線に沿った閉1形式の積分は等しい」ことを認めてよい。
\end{problem}

\begin{solution}
(1) $\omega=P\,dx+Q\,dy$、$P=-y/r^2$、$Q=x/r^2$（$r^2=x^2+y^2$）とすると
\[
\partial_xQ=\frac{r^2-2x^2}{r^4}=\frac{y^2-x^2}{r^4},\qquad \partial_yP=-\frac{r^2-2y^2}{r^4}=\frac{y^2-x^2}{r^4},
\]
$d\omega=(\partial_xQ-\partial_yP)dx\wedge dy=0$。$x=r\cos\theta$, $y=r\sin\theta$ とすると $-y\,dx+x\,dy=r^2d\theta$ なので、$\theta$ が一価に定まる領域（例えば負の $x$ 軸を除いた領域）では $\omega=d\theta$。複素数で書けば $\omega=\mathrm{Im}(dz/z)$。

(2) $\gamma(t)=(\cos t,\sin t)$ で $\gamma^*\omega=dt$、$\oint\omega=2\pi$。$\omega=df$ ならば閉曲線上の積分は $f(\gamma(2\pi))-f(\gamma(0))=0$ となるはずなので、$\omega$ は完全でない。したがって $H^1_{\mathrm{dR}}(X)\neq0$。

(3) $\gamma=z_1z_2$、$z_1=e^{2it}$、$z_2=1+\frac12e^{it}$ とすると $dz/z=dz_1/z_1+dz_2/z_2$ で
\[
\oint\omega=\mathrm{Im}\oint\frac{dz_1}{z_1}+\mathrm{Im}\oint\frac{dz_2}{z_2}=4\pi+0 .
\]
第2項：$z_2$ は中心 $1$、半径 $\frac12$ の円を動き原点を囲まない。その円を含む半平面 $\mathrm{Re}\,z>0$ では偏角が一価で $\omega=d\theta$ なので積分はゼロ。よって $w(\gamma)=2$。

(4) $\eta'=\eta-c\,\omega$ は閉で、$\oint_{S^1}\eta'=0$。任意の閉曲線 $\gamma$ は $S^1$ の $n$ 周とホモトピックなので $\oint_\gamma\eta'=n\oint_{S^1}\eta'=0$。そこで基点 $x_0$ を固定して $f(x)=\int_{x_0}^{x}\eta'$（任意の道に沿う）と定めると、二つの道の差は閉曲線なので $f$ は道によらず well-defined で、$df=\eta'$。よって $\eta=c\,\omega+df$、すなわち $H^1_{\mathrm{dR}}(X)=\R[\omega]$ で、類は周期 $\oint_{S^1}\eta=2\pi c$ で決まる。

物理：2次元の渦（第1章）の位相 $\theta(\bm x)$ について、$d\theta$ は芯の外で閉1形式であり、巻き付き数は $\frac1{2\pi}\oint d\theta\in\Z$。超流動の循環の量子化 $\oint\bm v_s\cdot d\bm x=\frac{2\pi n}{m}$ はこの式である。
\end{solution}

\begin{problem}[難易度 C]\label{ch05-prob-poincare}
$U\subset\R^n$ を原点に関して星形な開集合、$X=x^\mu\partial_\mu$ とし、$p\ge1$ の $p$ 形式 $\omega$ に対し定理 \ref{thm:ch05-poincare} の $K\omega$ を考える。\\
(1) $h_t(x)=tx$ とすると $(K\omega)=\int_0^1\frac1t\,h_t^*(\iota_X\omega)\,dt$ と書けることを示せ。\\
(2) $\frac{d}{dt}h_t^*\omega=\frac1t\,h_t^*(L_X\omega)$（$0<t\le1$）を示し、カルタンの公式を用いて $\omega=dK\omega+Kd\omega$ を証明せよ。\\
(3) 一様磁場 $F=B\,dx\wedge dy$（$\R^3$）に対し $A=KF$ を求めよ。一般の一様磁場 $F=\bm B\cdot d\bm S$ では $\bm A=\frac12\bm B\times\bm r$ となることを示し、これが $\bm r\cdot\bm A=0$ と $\nabla\cdot\bm A=0$ をみたすことを確かめよ。\\
(4) ディラック・モノポールの $F$ に対し北側の領域 $U_N=\R^3\setminus\{(0,0,z):z\le0\}$ でポテンシャルが存在することを、ポアンカレの補題から説明せよ（$U_N$ が点 $(0,0,1)$ に関して星形であることを示せ）。
\end{problem}

\begin{solution}
(1) $\iota_X\omega=\frac{1}{(p-1)!}x^\nu\omega_{\nu\mu_2\cdots\mu_p}(x)\,dx^{\mu_2}\wedge\cdots\wedge dx^{\mu_p}$。$h_t^*$ は $x\to tx$、$dx^\mu\to t\,dx^\mu$ と置き換えるので
\[
h_t^*(\iota_X\omega)=\frac{t^p}{(p-1)!}\,x^\nu\omega_{\nu\mu_2\cdots\mu_p}(tx)\,dx^{\mu_2}\wedge\cdots\wedge dx^{\mu_p}=t^p\,(\iota_X\omega)(tx).
\]
$\frac1t$ をかけて積分すれば $K\omega$ の定義式 $\int_0^1t^{p-1}(\iota_X\omega)(tx)\,dt$ に一致する。$p\ge1$ なので被積分関数は $t\to0$ で有界である。

(2) $X$ の流れは $\varphi_s(x)=e^sx$ で、$h_t=\varphi_{\log t}$。$t=e^s$ とすると
\[
\frac{d}{dt}h_t^*\omega=\frac{ds}{dt}\frac{d}{ds}\varphi_s^*\omega=\frac1t\,\varphi_s^*L_X\omega=\frac1t\,h_t^*L_X\omega .
\]
（$\frac{d}{ds}\varphi_s^*\omega=\varphi_s^*L_X\omega$ は流れの群性 $\varphi_{s+u}=\varphi_s\circ\varphi_u$ から従う。）カルタンの公式 $L_X=d\iota_X+\iota_Xd$ と $h_t^*d=dh_t^*$ により
\[
\frac{d}{dt}h_t^*\omega=d\Bigl(\frac1th_t^*\iota_X\omega\Bigr)+\frac1th_t^*\iota_X(d\omega).
\]
$\varepsilon<t\le1$ で積分して $\varepsilon\to0$ とすると、左辺は $h_1^*\omega-h_\varepsilon^*\omega\to\omega-0$（$h_\varepsilon^*\omega$ は係数に $\varepsilon^p$ がかかるので $p\ge1$ でゼロに行く）。右辺は (1) により $dK\omega+K(d\omega)$。（$d\omega$ は $p+1$ 形式なので $K$ の中の冪は $t^{p}$ になる。）以上で $\omega=dK\omega+Kd\omega$。

(3) $\iota_X(dx\wedge dy)=x\,dy-y\,dx$ なので
\[
KF=\int_0^1t\,B\,(x\,dy-y\,dx)\,dt=\frac B2(x\,dy-y\,dx),
\]
対称ゲージである。一般に $F=\frac12\varepsilon_{ijk}B_i\,dx^j\wedge dx^k$ で $\iota_XF=\varepsilon_{ijk}B_ix^j\,dx^k=(\bm B\times\bm r)_k\,dx^k$。よって $\bm A=\frac12\bm B\times\bm r$。検算：$\nabla\times(\bm B\times\bm r)=\bm B(\nabla\cdot\bm r)-(\bm B\cdot\nabla)\bm r=3\bm B-\bm B=2\bm B$ なので $\nabla\times\bm A=\bm B$。$\bm r\cdot(\bm B\times\bm r)=0$（フォック・シュウィンガー・ゲージ、一般に $\iota_XK\omega=0$ は $\iota_X^2=0$ から従う）。$\nabla\cdot(\bm B\times\bm r)=\bm r\cdot(\nabla\times\bm B)-\bm B\cdot(\nabla\times\bm r)=0$。

(4) $p=(0,0,1)$、$x\in U_N$ とし、線分 $(1-s)p+sx$（$0\le s\le1$）を考える。$x$ が $z$ 軸上にないとき、$s>0$ の点の $(x,y)$ 成分は $s(x_1,x_2)\neq0$ で、$s=0$ は $p$ なので負の $z$ 軸に触れない。$x=(0,0,x_3)$、$x_3>0$ のときは線分は正の $z$ 軸上にある。よって $U_N$ は $p$ に関して星形であり、平行移動してポアンカレの補題を使えば、$U_N$ 上の閉2形式 $F$ は完全：$F=dA$ となる $A$ が $U_N$ 上にある。$A_N=\frac{g}{4\pi}(1-\cos\theta)d\phi$ はその一つの例である（問題 \ref{ch05-prob-dirac}）。
\end{solution}

\begin{problem}[難易度 B]\label{ch05-prob-sn}
(1) マイヤー・ヴィートリス完全列を用いて、$n\ge1$ に対し
\[
H^p_{\mathrm{dR}}(S^n)\cong\begin{cases}\R&(p=0,n)\\0&(\text{それ以外})\end{cases}
\]
を示せ。$U=S^n\setminus\{N\}$、$V=S^n\setminus\{S\}$（北極・南極を除く）とし、$U\cap V\simeq S^{n-1}$ を用いよ。\\
(2) 向きづけられた閉多様体 $M$（境界なし、$\dim M=n$）上の体積形式は完全形式でないことを示せ。特に $S^2$ 上の $\sin\theta\,d\theta\wedge d\phi$ は閉だが完全でない。\\
(3) $S^n$ 上の $n$ 形式 $\omega$ が完全であるための必要十分条件は $\int_{S^n}\omega=0$ であることを示せ。電荷密度 $\rho$ をもつ静電場を $S^3$ 型の閉じた宇宙（空間が $S^3$）で考えると、全電荷 $\int\rho\,\mathrm{vol}$ がゼロでなければならないことを説明せよ。
\end{problem}

\begin{solution}
(1) 立体射影で $U\cong V\cong\R^n$ は可縮なので $H^p(U)=H^p(V)=0$（$p\ge1$）、$H^0=\R$。$U\cap V\cong\R^n\setminus\{0\}\simeq S^{n-1}$。

\textbf{$n=1$}：$U\cap V$ は2つの開弧で $H^0(U\cap V)=\R^2$、$H^1(U\cap V)=0$。完全列
\[
0\to H^0(S^1)\to\R\oplus\R\xrightarrow{\ r\ }\R^2\xrightarrow{\partial^*}H^1(S^1)\to0,\qquad r(a,b)=(b-a,\,b-a)
\]
で、$r$ の像は1次元なので $H^1(S^1)\cong\R^2/\R\cong\R$、$H^0(S^1)=\R$（連結）。

\textbf{$n\ge2$}：$p\ge1$ では $H^p(U)\oplus H^p(V)=0$ なので完全列の部分 $0\to H^p(U\cap V)\xrightarrow{\partial^*}H^{p+1}(S^n)\to0$ から $H^{p+1}(S^n)\cong H^p(S^{n-1})$（$p\ge1$）。$p=0$ の部分
\[
0\to H^0(S^n)\to\R\oplus\R\xrightarrow{\ r\ }H^0(U\cap V)=\R\xrightarrow{\partial^*}H^1(S^n)\to0,\qquad r(a,b)=b-a
\]
（$U\cap V$ は連結）で、$(a,b)\mapsto b-a$ は全射なので $H^1(S^n)=0$、$H^0(S^n)=\R$。$p\ge2$ では $H^p(S^n)\cong H^{p-1}(S^{n-1})$ を繰り返す。$p<n$ なら $H^1(S^{n-p+1})=0$（$n-p+1\ge2$）に、$p=n$ なら $H^1(S^1)=\R$ に、$p>n$ なら $H^{p-n+1}(S^1)=0$（次数 $\ge2$）に行き着く。よって主張が成り立つ。

(2) 体積形式 $\mathrm{vol}$ が $d\alpha$ なら、ストークスの定理より $\int_M\mathrm{vol}=\int_{\partial M}\alpha=0$ だが、$\int_M\mathrm{vol}>0$ なので矛盾。$S^2$ では $\int\sin\theta\,d\theta\wedge d\phi=4\pi\neq0$。

(3) ストークスより $\int_{S^n}:H^n(S^n)\to\R$ は well-defined な線形写像で、(2) より $[\mathrm{vol}]\mapsto\mathrm{Vol}(S^n)\neq0$。$H^n(S^n)\cong\R$ は1次元なのでこの写像は同型であり、核はゼロ。すなわち $\int\omega=0\iff[\omega]=0\iff\omega$ は完全。
静電場：空間 $S^3$ 上でガウスの法則は $d({*_3E})=\rho\,\mathrm{vol}$（${*_3E}$ は電場の2形式、問題 \ref{ch05-prob-maxwell} の ${*F}|_{\Sigma_t}=\bm E\cdot d\bm S$）。右辺は完全形式なので $\int_{S^3}\rho\,\mathrm{vol}=0$。閉じた空間では電気力線の逃げ場がなく、正味の電荷を置くことはできない（ガウスの法則の大域的な制約）。逆に $\int\rho=0$ なら (3) により解 $\bm E$ が存在する。
\end{solution}

\begin{center}
\begin{tikzpicture}[scale=1.05]
  \fill[gray!20] (0,1) circle (0.3); \draw[thick] (0,1) circle (0.3);
  \fill[gray!20] (0,-1) circle (0.3); \draw[thick] (0,-1) circle (0.3);
  \node at (0,1) {\scriptsize $\Phi_1$}; \node at (0,-1) {\scriptsize $\Phi_2$};
  \node[right] at (0.3,1.15) {\footnotesize $a_1$};
  \node[right] at (0.3,-1.15) {\footnotesize $a_2$};
  \node[vtx] at (-3.2,0) {}; \node[left] at (-3.2,0) {$S$};
  \draw[thick] (3.4,-2.3) -- (3.4,2.3); \node[right] at (3.4,2.1) {\footnotesize スクリーン};
  \node[vtx] at (3.4,0) {}; \node[right] at (3.4,-0.25) {$P$};
  \draw[thick,blue!70!black,mid arrow] (-3.2,0) .. controls (-1.5,2.2) and (1.5,2.2) .. (3.4,0);
  \draw[thick,green!50!black,mid arrow] (-3.2,0) .. controls (-1.5,0.2) and (1.5,-0.2) .. (3.4,0);
  \draw[thick,red!70!black,mid arrow] (-3.2,0) .. controls (-1.5,-2.2) and (1.5,-2.2) .. (3.4,0);
  \node[blue!70!black] at (-1.3,1.75) {$\gamma_+$};
  \node[green!50!black] at (-1.3,0.35) {$\gamma_0$};
  \node[red!70!black] at (-1.3,-1.75) {$\gamma_-$};
\end{tikzpicture}
\zucap{2本のソレノイド（$z$ 軸に平行、断面は点 $a_1=(0,1)$, $a_2=(0,-1)$ のまわりの灰色の円板、磁束 $\Phi_1,\Phi_2$ は $+z$ 向き）と、源 $S$ からスクリーン上の点 $P$ へ向かう3つの経路。ソレノイドの外側では $F=0$ だが $A$ は完全でない（問題 \ref{ch05-prob-ab2}）。}
\end{center}

\begin{problem}[難易度 B]\label{ch05-prob-ab2}
平面 $M=\R^2\setminus\{a_1,a_2\}$（$a_1=(0,1)$, $a_2=(0,-1)$）を考え、$\theta_i$ を $a_i$ のまわりの偏角、$\omega_i=\frac1{2\pi}d\theta_i$ とする。\\
(1) $U=\{y>-\tfrac12\}\setminus\{a_1\}$、$V=\{y<\tfrac12\}\setminus\{a_2\}$ としてマイヤー・ヴィートリス完全列を使い、$H^1_{\mathrm{dR}}(M)\cong\R^2$ で $[\omega_1],[\omega_2]$ が基底であることを示せ。\\
(2) 上図のように、断面が $a_i$ を中心とする小円板のソレノイドに磁束 $\Phi_i$（$+z$ 向き）が閉じ込められている。ソレノイドの外側では $A$ は閉形式で、ある関数 $\chi$ により $A=\Phi_1\omega_1+\Phi_2\omega_2+d\chi$ と書けることを示せ。\\
(3) 電荷 $q$ の粒子が経路 $\gamma$ に沿って伝わるとき、振幅は位相 $\exp(iq\int_\gamma A)$ を得る。図の経路 $\gamma_+,\gamma_0,\gamma_-$ について、$P$ における振幅の相対位相 $\exp\bigl(iq\int_{\gamma_+}A-iq\int_{\gamma_0}A\bigr)$ などを求めよ。どの経路の組でも干渉縞がずれないための条件を求めよ。
\end{problem}

\begin{solution}
(1) $U\simeq S^1$（$a_1$ を除いた半平面）、$V\simeq S^1$、$U\cup V=M$、$U\cap V=\{-\tfrac12<y<\tfrac12\}$ は帯で可縮（$a_1,a_2$ を含まない）。完全列
\begin{multline*}
0\to H^0(M)\to H^0(U)\oplus H^0(V)\to H^0(U\cap V)\\
\xrightarrow{\partial^*}H^1(M)\to H^1(U)\oplus H^1(V)\to H^1(U\cap V)=0
\end{multline*}
で、$\R\oplus\R\to\R$, $(a,b)\mapsto b-a$ は全射なので $\partial^*=0$。よって $H^1(M)\cong H^1(U)\oplus H^1(V)\cong\R^2$（制限写像が同型）。
$\omega_1$ の $U$ への制限は $H^1(U)$ の生成元（$a_1$ を回る円周での積分が $1$）。$V$ への制限：$\omega_1$ は凸集合 $\{y<\frac12\}$ 全体で滑らかな閉形式なので完全、類はゼロ。同様に $\omega_2\mapsto(0,\text{生成元})$。よって $[\omega_1],[\omega_2]$ は基底。また $a_i$ のまわりの小円 $C_i$（反時計回り）について $\oint_{C_i}\omega_j=\delta_{ij}$。

(2) 外側の領域 $M'$（平面から2つの閉円板を除いたもの）は $M$ にホモトピー同値なので（(1) と同じ議論で）$H^1(M')=\R^2$ で基底 $[\omega_1],[\omega_2]$。外側で $F=dA=0$ なので $A$ は閉。ソレノイドを含めた平面全体で $A$ は滑らかなので、ストークスの定理により $\oint_{C_i}A=\int_{D_i}F=\Phi_i$（$D_i$ は $C_i$ が囲む円板）。$A'=A-\Phi_1\omega_1-\Phi_2\omega_2$ は閉で、$\oint_{C_i}A'=0$。(1) の同型の下で $[A']$ の像は $H^1(U)$, $H^1(V)$ での類で、それらは $C_1$, $C_2$ 上の周期（ゼロ）で決まる（$H^1(S^1)\cong\R$ は周期で検出される、問題 \ref{ch05-prob-dtheta}）。よって $[A']=0$、$A'=d\chi$。

(3) 位相差は閉曲線 $\gamma_+\cdot\gamma_0^{-1}$ 上の積分 $\exp(iq\oint A)$ で、$d\chi$ の寄与は消える。$\gamma_+\cdot\gamma_0^{-1}$ は $S\to$（上側）$\to P\to$（中央）$\to S$ と進み、$a_1$ だけを\textbf{時計回り}に1周する（$a_1$ の巻き付き数 $-1$、$a_2$ は $0$）。よって
\begin{gather*}
e^{iq(\int_{\gamma_+}A-\int_{\gamma_0}A)}=e^{-iq\Phi_1},\qquad
e^{iq(\int_{\gamma_0}A-\int_{\gamma_-}A)}=e^{-iq\Phi_2},\\
e^{iq(\int_{\gamma_+}A-\int_{\gamma_-}A)}=e^{-iq(\Phi_1+\Phi_2)} .
\end{gather*}
一般の2経路の差の閉曲線が $a_i$ を $n_i$ 回（反時計回り）回るなら位相は $e^{iq(n_1\Phi_1+n_2\Phi_2)}$。すべての組で干渉縞がずれない条件は $q\Phi_1\in2\pi\Z$ かつ $q\Phi_2\in2\pi\Z$。$q\Phi_1+q\Phi_2\in2\pi\Z$ だけでは $\gamma_+$ と $\gamma_0$ の干渉が区別できる。粒子が感じるのは $F$ そのものではなく、コホモロジー類 $[A]\in H^1(M;\R)$ を周期で読んだ $e^{iq\Phi_i}$ である。
\end{solution}

\begin{problem}[難易度 B]\label{ch05-prob-torus}
$T^n=\R^n/\Z^n$（座標 $x^i\sim x^i+1$、平坦計量）とする。\\
(1) 定数係数の形式 $dx^{i_1}\wedge\cdots\wedge dx^{i_p}$ は調和であることを示せ。さらに平坦計量のもとで調和形式は定数係数のものに限ることを示し、ホッジの定理から $b_p(T^n)=\binom np$ を結論せよ。オイラー標数 $\chi(T^n)=\sum_p(-1)^pb_p$ を求めよ。\\
(2) $T^2$ の閉1形式 $\omega$ に対し、周期 $\alpha=\oint_a\omega$, $\beta=\oint_b\omega$（$a,b$ は図のサイクル）を用いて $\omega-\alpha\,dx-\beta\,dy$ が完全形式であることを、普遍被覆 $\R^2$ への持ち上げを用いて示せ。\\
(3) $\omega=\bigl(2+2\pi\cos2\pi(x+y)\bigr)dx+\bigl(-1+2\pi\cos2\pi(x+y)\bigr)dy$ が閉であることを確かめ、周期と調和代表元を求めよ。\\
(4) $T^2$ 上の $U(1)$ ゲージ場で $F=0$ のもの（平坦接続）を考える。ゲージ不変な量 $e^{i\oint_aA}$, $e^{i\oint_bA}$ は $A$ の調和部分（定数部分）だけで決まることを説明せよ。
\end{problem}

\begin{solution}
(1) 平坦計量では ${*}$ は定数係数の形式を定数係数の形式に写すので、定数係数の $\omega$ は $d\omega=0$、$d{*\omega}=0$、よって $\delta\omega=0$、$\Delta\omega=0$。
逆に、平坦計量の下で $\omega=\frac1{p!}\omega_I\,dx^I$ に対し $\Delta\omega=\frac1{p!}(-\sum_k\partial_k^2\omega_I)\,dx^I$（ユークリッド空間のホッジ・ラプラシアンは成分ごとに $-\nabla^2$ として働く）。調和なら各成分 $f=\omega_I$ は $T^n$ 上で $\nabla^2f=0$ をみたし、
\[
0=-\int_{T^n}f\,\nabla^2f\,d^nx=\int_{T^n}\lvert\nabla f\rvert^2d^nx
\]
（周期性により表面項なし）から $f$ は定数。よって $\mathcal H^p$ は $\{dx^{i_1}\wedge\cdots\wedge dx^{i_p}\}_{i_1<\cdots<i_p}$ で張られ、ホッジの定理より $b_p=\dim\mathcal H^p=\binom np$。$\chi(T^n)=\sum_p(-1)^p\binom np=(1-1)^n=0$。

(2) $\pi\colon\R^2\to T^2$ で引き戻した $\tilde\omega=\pi^*\omega$ は $\R^2$ 上の閉形式なので、ポアンカレの補題により $\tilde\omega=d\tilde f$。$\tilde\omega'=\tilde\omega-\alpha\,dx-\beta\,dy=d\tilde g$、$\tilde g=\tilde f-\alpha x-\beta y$。$\tilde g(x+1,y)-\tilde g(x,y)$ は線分 $(x,y)\to(x+1,y)$ 上の $\tilde\omega'$ の積分で、その像は $T^2$ の閉曲線で $a$ にホモトピック（$y$ をずらすホモトピー）。閉形式の積分はホモトピー不変なので $\oint_a\omega-\alpha=0$。同様に $y$ 方向も周期的。よって $\tilde g$ は $T^2$ 上の関数 $g$ を定め、$\omega-\alpha\,dx-\beta\,dy=dg$。これは $H^1(T^2)\cong\R^2$ が周期 $(\alpha,\beta)$ で同定されることを示す（ド・ラームの定理の具体例）。

(3) $u=2\pi(x+y)$ とすると $\omega=2\,dx-dy+\cos u\,du=2\,dx-dy+d(\sin u)$。よって $d\omega=0$。周期は
\[
\oint_a\omega=\int_0^1\bigl(2+2\pi\cos2\pi(x+y_0)\bigr)dx=2,\qquad\oint_b\omega=-1,
\]
調和代表元は $2\,dx-dy$（ノルム最小の代表元）。

(4) $F=dA=0$ なので $A$ は閉。(2) より $A=\alpha\,dx+\beta\,dy+dg$ で、$\oint_aA=\alpha$、$\oint_bA=\beta$。$dg$ は小さなゲージ変換 $e^{ig}$ で消せ、周期に寄与しない。したがってウィルソン・ループ $e^{i\alpha}$, $e^{i\beta}$ は定数部分だけで決まる。さらに大きなゲージ変換 $e^{2\pi i(mx+ny)}$（$m,n\in\Z$）は $(\alpha,\beta)\to(\alpha+2\pi m,\beta+2\pi n)$ とずらすので、ゲージ同値類の空間は $\R^2/(2\pi\Z)^2\cong T^2$（問題 \ref{ch05-prob-holonomy}）。
\end{solution}

\subsection{ディラック・モノポールとAB効果}

\begin{problem}[難易度 B]\label{ch05-prob-dirac}
$\R^3\setminus\{0\}$ 上で $F=\frac{g}{4\pi}\sin\theta\,d\theta\wedge d\phi$ とする。\\
(1) デカルト座標で $F=\frac{g}{4\pi r^3}(x\,dy\wedge dz+y\,dz\wedge dx+z\,dx\wedge dy)$、すなわち $\bm B=\frac{g\,\hat{\bm r}}{4\pi r^2}$ であることを示し、$dF=0$ と $\int_{S^2}F=g$ を確かめよ。\\
(2) $(1-\cos\theta)\,d\phi=\dfrac{x\,dy-y\,dx}{r(r+z)}$ を示し、$A_N=\frac{g}{4\pi}(1-\cos\theta)d\phi$ が負の $z$ 軸の外で滑らかであることを示せ。$A_S$ についても同様に述べよ。\\
(3) 負の $z$ 軸を囲む半径 $\varepsilon$ の小円（$\theta$ が $\pi$ に近い緯線、$+\phi$ 向き）に沿って $\oint A_N$ を計算し、$\varepsilon\to0$ の極限を求めよ。これを「ディラック弦に沿って磁束 $g$ が原点に流れ込んでいる」と解釈せよ。\\
(4) $A_N-A_S$ を求め、電荷 $q$ の波動関数の貼り合わせが一価であるための条件として $qg\in2\pi\Z$ を導け。単位電荷 $q=1$ のとき最小の磁荷はいくらか。$\hbar$ と結合定数 $e$ を戻すとどう書けるか。
\end{problem}

\begin{solution}
(1) $\Omega=\frac{1}{r^3}(x\,dy\wedge dz+y\,dz\wedge dx+z\,dx\wedge dy)=\frac1{r^3}\iota_X d^3x$（$X=x^i\partial_i$）とおく。$\iota_X\Omega=\frac1{r^3}\iota_X\iota_Xd^3x=0$ で、$\Omega$ は拡大 $x\to\lambda x$ で不変（係数が $-2$ 次、$dx\wedge dx$ が $+2$ 次）。したがって $\Omega$ は $\theta,\phi$ だけを含み $dr$ を含まない2形式で、$r$ によらない。単位球面上では $\iota_{\hat{\bm r}}d^3x$ は（外向きの向きの）面積要素 $\sin\theta\,d\theta\wedge d\phi$ に等しい。よって $\Omega=\sin\theta\,d\theta\wedge d\phi$（立体角の形式）。$F=\frac{g}{4\pi}\Omega=\bm B\cdot d\bm S$ で $\bm B=\frac{g\bm r}{4\pi r^3}$。
$dF=0$：$\nabla\cdot(\bm r/r^3)=\frac{3}{r^3}-\frac{3r^2}{r^5}=0$（$r\neq0$）。あるいは $d(\sin\theta\,d\theta\wedge d\phi)=0$（$\theta$ だけの関数）。積分 $\int_0^\pi\!\int_0^{2\pi}\frac{g}{4\pi}\sin\theta\,d\phi\,d\theta=g$。

(2) $\cos\theta=z/r$、$d\phi=\frac{x\,dy-y\,dx}{x^2+y^2}$、$x^2+y^2=r^2-z^2=(r-z)(r+z)$ より
\[
(1-\cos\theta)\,d\phi=\frac{r-z}{r}\cdot\frac{x\,dy-y\,dx}{(r-z)(r+z)}=\frac{x\,dy-y\,dx}{r(r+z)} .
\]
分母がゼロになるのは $r+z=0$、すなわち負の $z$ 軸（と原点）だけ。正の $z$ 軸上でも分子・分母は滑らかで分母 $=2r^2\neq0$。よって $A_N$ は $U_N$ で滑らか。同様に $(1+\cos\theta)d\phi=\frac{x\,dy-y\,dx}{r(r-z)}$ なので $A_S=-\frac{g}{4\pi}\frac{x\,dy-y\,dx}{r(r-z)}$ は正の $z$ 軸の外で滑らか。$dA_N=\frac{g}{4\pi}\sin\theta\,d\theta\wedge d\phi=F$、$dA_S=-\frac{g}{4\pi}(-\sin\theta)\,d\theta\wedge d\phi=F$。

(3) $\theta=\theta_0$ の緯線上で $\oint A_N=\frac{g}{4\pi}(1-\cos\theta_0)\cdot2\pi=\frac g2(1-\cos\theta_0)$。$\theta_0\to\pi$（小円が負の $z$ 軸に縮む）で $\oint A_N\to g$。一方、点 $(0,0,-a)$ を中心とする半径 $\varepsilon$ の水平な円板（$+z$ 向き、境界は $+\phi$ 向き）を通るモノポールの磁束は、立体角が $O(\varepsilon^2/a^2)$ なので $\varepsilon\to0$ でゼロになる（しかも下向きの磁場なので符号は負）。境界の周回積分は $g$ に近づくのに、円板を通る $F$ の磁束はゼロに近づく。ストークスの定理と整合させるには、負の $z$ 軸上に $+z$ 向きの磁束 $g$ を運ぶ無限に細いソレノイド（ディラック弦）があると考えればよい：弦が $-\infty$ から原点に磁束 $g$ を運び込み、原点から等方的に湧き出す。$A_N$ を全空間で使うと、弦の分も含めて $d(\text{全磁束})=0$ が保たれる。弦が観測できない条件が次の量子化条件である（AB効果の位相 $e^{iqg}=1$）。

(4) $A_N-A_S=\frac{g}{4\pi}\bigl[(1-\cos\theta)+(1+\cos\theta)\bigr]d\phi=\frac{g}{2\pi}d\phi=d\alpha$、$\alpha=\frac{g\phi}{2\pi}$。本書の規約で $U(1)$ ゲージ変換 $h=e^{i\alpha}$ は $A\to A+i\,h\,dh^{-1}=A+d\alpha$ であり、電荷 $q$ の場は $\psi\to e^{iq\alpha}\psi$（$D=d-iqA$ と共変）。よって重なりで $\psi_N=e^{iqg\phi/2\pi}\psi_S$。$\psi_N,\psi_S$ はそれぞれ一価なので、赤道を1周したとき $e^{iqg}=1$、すなわち $qg\in2\pi\Z$。
$q=1$ なら最小の磁荷は $g=2\pi$、一般に $g=2\pi n$ で $\int_{S^2}F/2\pi=n\in\Z$（第1チャーン数、第7, 8章）。$e$ と $\hbar$ を戻すと（$A\to eA/\hbar$）$eg=2\pi\hbar n$。ヘヴィサイド・ローレンツ単位の慣用形 $\frac{eg}{4\pi\hbar}=\frac n2$ と同じである。
\end{solution}

\begin{problem}[難易度 B]\label{ch05-prob-flux}
(1) $\Sigma$ を $\R^3\setminus\{0\}$ 内の向きづけられた閉曲面（外向き）とし、問題 \ref{ch05-prob-dirac} の $F$ について、$\Sigma$ が原点を囲むなら $\int_\Sigma F=g$、囲まないなら $0$ であることを、ストークスの定理を用いて示せ（基本事項の図を参照）。\\
(2) 電荷分布 $\rho$ が有界領域 $D$ に台をもつとする。$D$ の外の閉曲面 $\Sigma$ に対し $Q(\Sigma)=\int_\Sigma{*F}$（$\Sigma$ は時刻一定面内）は、$\Sigma$ を $D$ と交わらないように連続変形しても変わらないことを示せ。$Q(\Sigma)$ は何に等しいか。\\
(3) 時刻一定面 $\Sigma_t=\{t\}\times\R^3$ 上の全電荷 $Q(t)=\int_{\Sigma_t}{*j}$ が、$\bm J$ が空間遠方で十分速く減衰すれば $t$ によらないことを、$[t_1,t_2]\times\R^3$ 上のストークスの定理で示せ。境界の向きに注意せよ。
\end{problem}

\begin{solution}
(1) 囲まない場合：$\Sigma$ は $\R^3$ の有界領域 $V$ の境界で、$V$ は原点を含まない（$\Sigma\subset\R^3\setminus\{0\}$ でジョルダン・ブラウワーの分離定理により $\R^3\setminus\Sigma$ の有界成分が $V$）。$F$ は $V$ 上滑らかで閉なので $\int_\Sigma F=\int_VdF=0$。
囲む場合：十分小さな $\varepsilon$ で球面 $S_\varepsilon$ は $V$ の内部にあり、$V'=V\setminus B_\varepsilon$ の境界は $\partial V'=\Sigma-S_\varepsilon$（$S_\varepsilon$ の外向き法線は $V'$ にとって内向きなので負号）。$\int_\Sigma F-\int_{S_\varepsilon}F=\int_{V'}dF=0$。$\int_{S_\varepsilon}F=\frac{g}{4\pi}\cdot4\pi=g$（$\varepsilon$ によらない）。よって $\int_\Sigma F=g$。まとめて $\int_\Sigma F=g\cdot(\text{原点を囲む回数})$ で、これは $[\Sigma]\in H_2(\R^3\setminus\{0\})\cong\Z$ と $[F]$ のペアリングである。

(2) $D$ の外 $\R^3\setminus D$ では $\rho=0$ なので、静的とは限らなくても時刻一定面に制限したガウスの法則 $d({*F}|_{\Sigma_t})=\rho\,d^3x$（問題 \ref{ch05-prob-maxwell}(3) の $d^3x$ 成分）から、${*F}|_{\Sigma_t}=\bm E\cdot d\bm S$ は $\R^3\setminus D$ 上の閉2形式。$\Sigma_0$ から $\Sigma_1$ への $D$ を避ける連続変形 $H\colon\Sigma\times[0,1]\to\R^3\setminus D$ があれば、
\[
\int_{\Sigma_1}{*F}-\int_{\Sigma_0}{*F}=\int_{\Sigma\times I}d(H^*{*F})=\int_{\Sigma\times I}H^*d{*F}=0 .
\]
（$\partial(\Sigma\times I)=\Sigma\times\{1\}-\Sigma\times\{0\}$。）$\Sigma$ が $D$ を囲むとき、$\Sigma$ の内部を $V$ として $\int_\Sigma\bm E\cdot d\bm S=\int_V\rho\,d^3x=Q_{\text{内部}}$：囲まれた電荷。

(3) $W=[t_1,t_2]\times B_R$（向き $dt\wedge d^3x$）に $d{*j}=0$ を積分すると $0=\int_{\partial W}{*j}$。外向き法線先頭の規約により、$t=t_2$ の面は外向き法線 $+\partial_t$ で向き $d^3x$（$dt\wedge d^3x$ が正）、$t=t_1$ の面は外向き法線 $-\partial_t$ なので向き $-d^3x$。よって
\[
0=\int_{\{t_2\}\times B_R}\rho\,d^3x-\int_{\{t_1\}\times B_R}\rho\,d^3x+\int_{[t_1,t_2]\times\partial B_R}{*j}.
\]
側面の項は $-\int dt\oint_{\partial B_R}\bm J\cdot d\bm S$ で、$R\to\infty$ で $\bm J$ が十分速く減衰すれば消える。したがって $Q(t_2)=Q(t_1)$。$d{*j}=0$ は「閉形式の積分はサイクルを動かしても不変」という形で保存則を表す。ネーターの保存量一般（$d{*J}=0$）について同じことが言える。
\end{solution}

\begin{problem}[難易度 C]\label{ch05-prob-holonomy}
$M=\R^3\setminus\{x^2+y^2\le r_0^2\}$（無限に長いソレノイドの外側）とし、$M$ 上の $U(1)$ ゲージ場 $A$ が $F=dA=0$ をみたすとする。ソレノイド内部の磁束を $\Phi$ とする（$+z$ 向き）。\\
(1) $M\simeq S^1$ を示し、$H^1_{\mathrm{dR}}(M)=\R[\omega]$、$\omega=\frac1{2\pi}d\phi$ を結論せよ。\\
(2) $M$ 内で自由ホモトピックな閉曲線 $\gamma_0,\gamma_1$ に対し $\oint_{\gamma_0}A=\oint_{\gamma_1}A$ を示せ。ホロノミー $\exp(iq\oint_\gamma A)$ が $[A]\in H^1_{\mathrm{dR}}(M)$ と $\gamma$ の巻き付き数 $n$ だけで決まり、$e^{iqn\Phi}$ に等しいことを示せ。\\
(3) $M$ 上の滑らかな写像 $h\colon M\to U(1)$ によるゲージ変換を考える。$h=e^{i\lambda}$（$\lambda$ は一価関数）の場合と $h=e^{im\phi}$（$m\in\Z$）の場合で $[A]$ がどう変わるかを調べ、「物理的に区別できる平坦な配位」の空間が $\R/2\pi\Z\cong U(1)$ であることを示せ（すべての電荷 $q$ は整数とする）。\\
(4) $M$ を $\R^3$ から $k$ 本の平行なソレノイドを除いたものに替えると、この空間はどうなるか。また、クーパー対（電荷 $2$）だけが関与する超伝導体の干渉実験では、$\Phi$ の周期がいくらになるか。
\end{problem}

\begin{solution}
(1) $r_0'>r_0$ を固定し、$C=\{x^2+y^2=r_0'^2,\ z=0\}$ とする。$\rho=\sqrt{x^2+y^2}>r_0$ として
\[
H\bigl((x,y,z),t\bigr)=\Bigl(\bigl((1-t)+t\,r_0'/\rho\bigr)x,\ \bigl((1-t)+t\,r_0'/\rho\bigr)y,\ (1-t)z\Bigr)
\]
とおくと、動径は $(1-t)\rho+t\,r_0'$ で $\rho$ と $r_0'$ の間にあり $r_0$ より大きいので、点は $M$ 内にとどまる。$H(\cdot,0)=\id$、$H(\cdot,1)\in C$、$C$ 上では $H$ は恒等的なので、$C$ は $M$ の強変位レトラクト。よって $M\simeq S^1$、$H^1_{\mathrm{dR}}(M)\cong H^1_{\mathrm{dR}}(S^1)=\R$。$\oint_C\omega=1$ なので $[\omega]$ が生成元。

(2) 自由ホモトピー $H\colon S^1\times[0,1]\to M$ を（滑らかに近似して）とると、ストークスの定理と $dA=0$ から
\[
\oint_{\gamma_1}A-\oint_{\gamma_0}A=\int_{\partial(S^1\times I)}H^*A=\int_{S^1\times I}H^*dA=0 .
\]
$A$ を $A+d\chi$ に替えても閉曲線上の積分は変わらないので、$\oint_\gamma A$ は $[A]$ と $\gamma$ のホモトピー類（$\pi_1(M)=\Z$ の元、すなわち巻き付き数 $n$）で決まる。$[A]=c[\omega]$ とすると $\oint_C A=c$。一方 $A$ はソレノイドの内部まで滑らかに伸びているので、$C$ が囲む円板に対しストークスの定理から $\oint_CA=\Phi$。よって $c=\Phi$、$\oint_\gamma A=n\Phi$、ホロノミーは $e^{iqn\Phi}$。

(3) 本書の規約で $A\to A+i\,h\,dh^{-1}$。$h=e^{i\lambda}$ では $A\to A+d\lambda$ で $[A]$ は不変。$h=e^{im\phi}$ は $M$ 上で一価で滑らか（$\phi$ 自体は多価だが $e^{im\phi}$ は一価）であり、$i\,h\,dh^{-1}=i\,e^{im\phi}(-im\,d\phi)e^{-im\phi}=m\,d\phi$。よって $A\to A+m\,d\phi$、$[A]=\Phi[\omega]\to(\Phi+2\pi m)[\omega]$。これは $\oint_CA$ を $2\pi m$ ずらすが、整数電荷 $q$ の場は $h^q=e^{iqm\phi}$ と一価に変換されるので許されるゲージ変換であり、ホロノミー $e^{iqn\Phi}$ も不変。
任意の $h\colon M\to U(1)$ は、$h|_C$ の巻き付き数を $m$ とすると $h=e^{im\phi}e^{i\lambda}$（$\lambda$ 一価）と書ける（$h\,e^{-im\phi}$ は巻き付き数ゼロなので対数が一価にとれる）。したがって $[A]$ はゲージ変換で $2\pi\Z[\omega]$ だけずれ、平坦な配位のゲージ同値類は
\[
H^1_{\mathrm{dR}}(M)/2\pi H^1(M;\Z)\cong\R/2\pi\Z\cong U(1),\qquad[A]\mapsto e^{i\Phi}.
\]
すべての観測量（ホロノミー）は $e^{i\Phi}$ の関数で、これが完全な不変量である。

(4) $k$ 本なら $M\simeq\bigvee_kS^1$（$\R^2\setminus k$ 点）で $H^1=\R^k$（問題 \ref{ch05-prob-ab2}）、平坦配位の空間は $\R^k/(2\pi\Z)^k\cong T^k$（各ソレノイドの $e^{i\Phi_i}$）。一般に $H_1$ がねじれを持たないとき $H^1(M;\R)/2\pi H^1(M;\Z)\cong T^{b_1}$。
電荷 $q$ の干渉では位相 $e^{iq\Phi}$ なので $\Phi$ の周期は $2\pi/q$。$q=2$ なら $\pi$。単位を戻すと $2\pi\hbar/(2e)=h/2e$ で、超伝導の磁束量子である。
\end{solution}

\begin{problem}[難易度 C]\label{ch05-prob-ff}
(1) $U(1)$ ゲージ場について $F\wedge F=d(A\wedge dA)$ を示せ。4次元時空の作用に $\theta$ 項 $S_\theta=\frac{\theta}{8\pi^2}\int F\wedge F$ を加えても、運動方程式が変わらないことを示せ。\\
(2) $M=S^2\times S^2$（積の向き）上で、$F=\pi_1^*F_1+\pi_2^*F_2$、$F_i$ は第 $i$ 因子の $S^2$ 上の閉2形式で $\int_{S^2}F_i=2\pi n_i$（$n_i\in\Z$）とする。$\frac{1}{8\pi^2}\int_MF\wedge F=n_1n_2$ を示せ。$n_1n_2\neq0$ ならば $M$ 全体で $F=dA$ となる $A$ は存在しないことを示せ。\\
(3) 非可換の場合、$F=dA-iA\wedge A$ について
\[
\tr F\wedge F=d\,\tr\Bigl(A\wedge dA-\frac{2i}{3}A\wedge A\wedge A\Bigr)
\]
を示せ（行列値形式のトレースの巡回性 $\tr(\alpha\wedge\beta)=(-1)^{pq}\tr(\beta\wedge\alpha)$ を用いよ）。
\end{problem}

\begin{solution}
(1) $d(A\wedge dA)=dA\wedge dA-A\wedge d^2A=F\wedge F$。変分：$\delta(F\wedge F)=2\,d\delta A\wedge F=2\,d(\delta A\wedge F)+2\,\delta A\wedge dF=2\,d(\delta A\wedge F)$（$dF=0$）。よって $\delta S_\theta$ は表面項だけで、局所的な変分（コンパクト台の $\delta A$）ではゼロ。運動方程式は変わらない。ただし $S_\theta$ の値そのものは大域的なトポロジーに依存しうる（(2)）。

(2) $\pi_i^*F_i\wedge\pi_i^*F_i=\pi_i^*(F_i\wedge F_i)=0$（2次元上の4形式）なので
\[
F\wedge F=\pi_1^*F_1\wedge\pi_2^*F_2+\pi_2^*F_2\wedge\pi_1^*F_1=2\,\pi_1^*F_1\wedge\pi_2^*F_2
\]
（2形式どうしは可換）。積の向きでのフビニの定理により $\int_M\pi_1^*F_1\wedge\pi_2^*F_2=\int F_1\int F_2=4\pi^2n_1n_2$。よって
\[
\frac1{8\pi^2}\int_MF\wedge F=\frac{2\cdot4\pi^2n_1n_2}{8\pi^2}=n_1n_2 .
\]
$F=dA$ が $M$ 全体で成り立つなら $F\wedge F=d(A\wedge F)$ で、閉多様体上のストークスにより積分はゼロ。$n_1n_2\neq0$ と矛盾。（実際には $\int_{S^2\times\{pt\}}F=2\pi n_1\neq0$ だけで $F$ が完全でないことがわかる。）$\theta$ 項は $e^{iS_\theta}=e^{i\theta n_1n_2}$ を与え、$\theta\sim\theta+2\pi$。なお $\frac1{8\pi^2}\int F\wedge F=\frac12\int c_1^2$ は一般には半整数になりうる（例えば $\CP^2$。交叉形式、第6章）。

(3) $\tr F\wedge F=\tr(dA\,dA)-i\tr(dA\,A^2)-i\tr(A^2\,dA)-\tr(A^4)$（$\wedge$ を省略）。$A^2$ は2形式なので $\tr(A^2dA)=\tr(dA\,A^2)$。$\tr(A^4)=\tr(A\cdot A^3)=(-1)^{1\cdot3}\tr(A^3\cdot A)=-\tr A^4$ より $\tr A^4=0$。よって
\[
\tr F\wedge F=\tr(dA\,dA)-2i\tr(dA\,A^2).
\]
一方 $d\tr(A\,dA)=\tr(dA\,dA)$、
\[
d\tr(A^3)=\tr(dA\,A^2)-\tr(A\,dA\,A)+\tr(A^2\,dA),
\]
ここで $\tr(A\cdot(dA\,A))=(-1)^{1\cdot3}\tr(dA\,A\cdot A)=-\tr(dA\,A^2)$、$\tr(A^2dA)=\tr(dA\,A^2)$ なので $d\tr A^3=3\tr(dA\,A^2)$。したがって
\[
d\,\tr\Bigl(A\,dA-\frac{2i}3A^3\Bigr)=\tr(dA\,dA)-2i\tr(dA\,A^2)=\tr F\wedge F .
\]
$8\pi^2$ で割れば記法と規約の $dCS(A)=\frac1{8\pi^2}\tr F\wedge F$。$CS(A)$ はゲージ不変でなく、大きなゲージ変換で整数（巻き付き数 $w$）だけ変わることを第8章で示す。
\end{solution}

\begin{problem}[難易度 B]\label{ch05-prob-gb}
単位球面 $S^2$（計量 $d\theta^2+\sin^2\theta\,d\phi^2$、向き $d\theta\wedge d\phi$）上に正規直交な双対枠 $e^1=d\theta$, $e^2=\sin\theta\,d\phi$ をとる（極以外で定義）。接続形式 $\omega^a{}_b=-\omega^b{}_a$ を構造方程式 $de^a+\omega^a{}_b\wedge e^b=0$ で定め、曲率を $R^1{}_2=d\omega^1{}_2$ とする（2次元では $SO(2)$ が可換なので $\omega\wedge\omega$ の項は消える）。\\
(1) $\omega^1{}_2=-\cos\theta\,d\phi$、$R^1{}_2=\sin\theta\,d\theta\wedge d\phi$ を示し、ガウス曲率 $K=1$（$R^1{}_2=K\,e^1\wedge e^2$）を確かめよ。\\
(2) 複素枠 $\epsilon=e^1+ie^2$ について $d\epsilon=i\,\omega^1{}_2\wedge\epsilon$ を示し、枠を回転した $\epsilon'=e^{i\alpha}\epsilon$ の接続形式が $\omega'=\omega^1{}_2+d\alpha$ であることを示せ。\\
(3) $\epsilon_N=e^{i\phi}\epsilon$ は北極の近傍で、$\epsilon_S=e^{-i\phi}\epsilon$ は南極の近傍で滑らかな枠であることを示し、対応する接続形式 $\omega_N=(1-\cos\theta)d\phi$、$\omega_S=-(1+\cos\theta)d\phi$ を求めよ。\\
(4) 上半球と下半球に分けてストークスの定理を用い、$\int_{S^2}R^1{}_2=\oint_{\text{赤道}}(\omega_N-\omega_S)=4\pi$ を示せ。ディラック・モノポールとの対応（$g$ は何に当たるか）を述べ、$\frac1{2\pi}\int R=2=\chi(S^2)$ がガウス・ボンネの定理（第8章）の特別な場合であることを述べよ。
\end{problem}

\begin{solution}
(1) $de^1=0$、$de^2=\cos\theta\,d\theta\wedge d\phi$。$\omega^1{}_2=-\cos\theta\,d\phi$（$\omega^2{}_1=\cos\theta\,d\phi$）とおくと
\[
de^1+\omega^1{}_2\wedge e^2=-\cos\theta\sin\theta\,d\phi\wedge d\phi=0,\qquad
de^2+\omega^2{}_1\wedge e^1=\cos\theta\,d\theta\wedge d\phi+\cos\theta\,d\phi\wedge d\theta=0 .
\]
（反対称な解は一意。）$R^1{}_2=d(-\cos\theta\,d\phi)=\sin\theta\,d\theta\wedge d\phi=e^1\wedge e^2$ なので $K=1$。

(2) $\omega\equiv\omega^1{}_2$ と書く。$de^1=-\omega\wedge e^2$、$de^2=-\omega^2{}_1\wedge e^1=\omega\wedge e^1$ より
\[
d\epsilon=-\omega\wedge e^2+i\,\omega\wedge e^1=i\,\omega\wedge(e^1+ie^2)=i\,\omega\wedge\epsilon .
\]
$d\epsilon'=i\,d\alpha\wedge e^{i\alpha}\epsilon+e^{i\alpha}i\,\omega\wedge\epsilon=i(\omega+d\alpha)\wedge\epsilon'$。よって $\omega'=\omega+d\alpha$ で、$U(1)$ ゲージ変換と同じ形である。曲率 $d\omega'=d\omega$ は不変。

(3) 北極付近で $x=\theta\cos\phi$, $y=\theta\sin\phi$ を座標とすると $dx+i\,dy=d(\theta e^{i\phi})=e^{i\phi}(d\theta+i\theta\,d\phi)$。$\sin\theta=\theta+O(\theta^3)$ なので $\epsilon_N=e^{i\phi}(d\theta+i\sin\theta\,d\phi)$ は $dx+i\,dy$ に滑らかな補正を加えたもので、北極でも滑らかかつ非退化。南極付近では $\theta'=\pi-\theta$, $x'+iy'=\theta'e^{-i\phi}$ として $d(\theta'e^{-i\phi})=-e^{-i\phi}(d\theta+i\theta'd\phi)$ で、$\epsilon_S=e^{-i\phi}\epsilon\approx-(dx'+i\,dy')$ は南極で滑らか。(2) より
\[
\omega_N=-\cos\theta\,d\phi+d\phi=(1-\cos\theta)\,d\phi,\qquad\omega_S=-\cos\theta\,d\phi-d\phi=-(1+\cos\theta)\,d\phi .
\]
いずれも各極で滑らか（$(1\mp\cos\theta)\sim\theta^2/2$ 等が $d\phi$ の特異性を打ち消す、問題 \ref{ch05-prob-dirac}(2) と同じ）。

(4) 上半球 $D_N$ の境界は、外向き法線 $\partial_\theta$ を先頭に置く規約で $+\phi$ 向きの赤道、下半球 $D_S$ の境界は $-\phi$ 向きの赤道。$R=d\omega_N$（$D_N$ 上）$=d\omega_S$（$D_S$ 上）なので
\[
\int_{S^2}R=\int_{D_N}d\omega_N+\int_{D_S}d\omega_S=\oint_{\text{赤道}}(\omega_N-\omega_S)=\oint2\,d\phi=4\pi .
\]
直接計算 $\int\sin\theta\,d\theta\,d\phi=4\pi$ とも一致。ディラック・モノポールの $A_N,A_S$ は $\frac{g}{4\pi}=1$、すなわち $g=4\pi$ の場合と全く同じ式であり、遷移関数は $e^{2i\phi}$（巻き付き数 $2$）。$\frac1{2\pi}\int R=2=\chi(S^2)$。遷移関数の巻き付き数は計量を連続的に変形しても変わらない整数なので、任意の計量で $\frac1{2\pi}\int R=2$ となる。一般の閉曲面で $\frac1{2\pi}\int_\Sigma R=\chi(\Sigma)=2-2g$ となることは第8章で示す。
\end{solution}

\subsection{ホッジ理論とカルツァ・クライン理論}

\begin{problem}[難易度 C]\label{ch05-prob-riemann}
$\Sigma_g$ を種数 $g$ の向きづけられた閉曲面、計量 $h$ をもつとする。$H_1(\Sigma_g;\Z)\cong\Z^{2g}$（第4章）とド・ラームの定理を用いてよい。\\
(1) 2次元では、1形式上の ${*}$ が共形変換 $h\to e^{2\sigma}h$ で不変であることを示せ。局所共形座標（$h=e^{2\sigma}(dx^2+dy^2)$）で ${*}dx=dy$、${*}dy=-dx$、1形式上 ${**}=-1$ を確かめよ。\\
(2) 調和1形式の空間 $\mathcal H^1$ の次元が $2g$ であること、および ${*}$ が $\mathcal H^1$ を保つことを示せ。\\
(3) $z=x+iy$ として ${*}dz=-i\,dz$ を示し、$\omega=f\,dz+k\,d\bar z$ が調和であることと $f$ が正則・$k$ が反正則であることが同値であることを示せ。これより正則1形式の空間の複素次元が $g$ であることを説明せよ。\\
(4) コンパクト・ボソン $X\sim X+2\pi$ の作用 $S=\frac{\kappa}{2}\int_{\Sigma_g}dX\wedge{*dX}$（$\kappa>0$）を考える。$dX$ は周期 $\oint_{c}dX\in2\pi\Z$ をもつ閉1形式である。周期を固定したとき作用を最小にする配位では $dX$ が調和であることを示せ。
\end{problem}

\begin{solution}
(1) 定義式の成分表示で $n=2$, $p=1$：${*}dx^i=\sqrt{h}\,h^{ij}\varepsilon_{jk}dx^k$。$h_{ij}=e^{2\sigma}\delta_{ij}$ なら $\sqrt h=e^{2\sigma}$、$h^{ij}=e^{-2\sigma}\delta^{ij}$ で因子が打ち消し合い ${*}dx^i=\varepsilon_{ik}dx^k$：${*}dx=dy$、${*}dy=-dx$。共形因子によらない。任意の計量は局所的に共形平坦（等温座標の存在）なので、1形式上の ${*}$ は共形類だけで決まる。${**}dx={*}dy=-dx$、一般式 $(-1)^{1\cdot1}=-1$ と一致。

(2) ホッジの定理より $\dim\mathcal H^1=b_1=\dim H^1_{\mathrm{dR}}(\Sigma_g)$。ド・ラームの定理より $H^1_{\mathrm{dR}}\cong\Hom(H_1(\Sigma_g),\R)\cong\R^{2g}$。
2次元の1形式では $\delta=-{*d*}$（リーマンの式 $(-1)^{n(p+1)+1}=(-1)^{5}=-1$）。$\omega$ が調和 $\iff d\omega=0$ かつ $d{*\omega}=0$。このとき $\eta={*\omega}$ について $d\eta=d{*\omega}=0$、$d{*\eta}=d({**\omega})=-d\omega=0$ なので $\eta$ も調和。

(3) ${*}dz={*}dx+i{*}dy=dy-i\,dx=-i(dx+i\,dy)=-i\,dz$。同様に ${*}d\bar z=i\,d\bar z$。$\partial=\partial_z$, $\bar\partial=\partial_{\bar z}$ として
\[
d\omega=\bar\partial f\,d\bar z\wedge dz+\partial k\,dz\wedge d\bar z=(\partial k-\bar\partial f)\,dz\wedge d\bar z ,
\]
\[
{*\omega}=-i f\,dz+i k\,d\bar z,\qquad d{*\omega}=-i\bar\partial f\,d\bar z\wedge dz+i\partial k\,dz\wedge d\bar z=i(\bar\partial f+\partial k)\,dz\wedge d\bar z .
\]
両方ゼロ $\iff\bar\partial f=0$ かつ $\partial k=0$。（共形座標の取り替えは正則なので、この条件は座標によらない。）
したがって $\mathcal H^1\otimes\C=\mathcal H^{1,0}\oplus\mathcal H^{0,1}$（${*}$ の固有値 $-i$, $+i$ の固有空間）で、$\mathcal H^{1,0}$ は正則1形式 $f\,dz$ の空間。複素共役は $\mathcal H^{1,0}$ と $\mathcal H^{0,1}$ を入れ替えるので両者の次元は等しく、$\dim_\C\mathcal H^{1,0}=\frac12\dim_\C(\mathcal H^1\otimes\C)=\frac12\cdot2g=g$。これは弦理論で種数 $g$ の世界面の正則微分の数（周期行列の大きさ）として現れる（第10章）。

(4) 周期を固定すると $[dX]$ が固定され、$dX=h+df$（$h$ は類の調和代表元、$f$ は一価関数）と書ける。$(h,df)=(\delta h,f)=0$ なので
\[
S=\frac\kappa2\lVert h+df\rVert^2=\frac\kappa2\bigl(\lVert h\rVert^2+\lVert df\rVert^2\bigr)\ge\frac\kappa2\lVert h\rVert^2 ,
\]
等号は $df=0$。よって最小作用配位は $dX=h$（調和）。各巻き付きセクター（$H^1(\Sigma_g;2\pi\Z)\cong(2\pi\Z)^{2g}$ の元）に1つずつ古典解があり、分配関数はそれらの和（格子和、テータ関数）を含む。
\end{solution}

\begin{problem}[難易度 D]\label{ch05-prob-coulomb}
$M$ をコンパクトで向きづけられた境界のない3次元リーマン多様体（例えば $T^3=\R^3/L\Z^3$）とする。ホッジの定理を用いてよい。\\
(1) $(\Delta\omega,\omega)=\lVert d\omega\rVert^2+\lVert\delta\omega\rVert^2$ を示し、調和 $\iff$ 閉かつ余閉 を示せ。また $d\Omega^{p-1}$、$\delta\Omega^{p+1}$、$\mathcal H^p$ が互いに直交することを示せ。\\
(2) 任意の1形式 $A$ に対し、関数 $\lambda$ で $\delta(A+d\lambda)=0$（クーロン・ゲージ）となるものが存在し、定数を除いて一意であることを示せ。（関数上 $\Delta=\delta d$ の像は定数の直交補空間であることを用いよ。）\\
(3) クーロン・ゲージで $A=h+\delta\beta$（$h\in\mathcal H^1$）と書ける。$T^3$ 上で $h$ を具体的に書き、小さなゲージ変換では消せないが、大きなゲージ変換 $e^{i\lambda}$, $\lambda=2\pi\,\bm m\cdot\bm x/L$（$\bm m\in\Z^3$）でどう変わるかを述べよ。\\
(4) 空間が $M$ の静電場について、ガウスの法則 $\nabla\cdot\bm E=\rho$ が解をもつための必要十分条件は $\int_M\rho\,\mathrm{vol}=0$ であることを示せ。解 $\bm E=-\nabla\varphi$ の一意性と、一般の解 $\bm E$ の不定性（調和部分）について述べよ。周期境界条件の格子QEDで1個の荷電粒子だけを置けない理由を説明せよ。
\end{problem}

\begin{solution}
(1) $(\Delta\omega,\omega)=(d\delta\omega,\omega)+(\delta d\omega,\omega)=(\delta\omega,\delta\omega)+(d\omega,d\omega)$（$\delta$ は $d$ の随伴）。リーマン計量では $\lVert\cdot\rVert^2\ge0$ なので $\Delta\omega=0\iff d\omega=0$ かつ $\delta\omega=0$（逆は自明）。
直交性：$(d\alpha,\delta\beta)=(d^2\alpha,\beta)=0$。$h$ 調和なら $(h,d\alpha)=(\delta h,\alpha)=0$、$(h,\delta\beta)=(dh,\beta)=0$。

(2) $\delta(A+d\lambda)=\delta A+\Delta\lambda$（関数上 $\Delta=\delta d$）。よって $\Delta\lambda=-\delta A$ を解けばよい。関数上の調和形式は $d\lambda=0$、すなわち定数（$M$ 連結とする）であり、ホッジ分解 $\Omega^0=\mathcal H^0\oplus\delta\Omega^1$ と $\Delta$ の自己共役性から $\operatorname{im}\Delta=(\mathcal H^0)^\perp=\{f:\int_Mf\,\mathrm{vol}=0\}$。$\delta A$ は $(\delta A,1)=(A,d1)=0$ をみたすので像に入り、解 $\lambda$ が存在する。二つの解の差は $\Delta$ の核、すなわち定数（ゲージ変換としては大域的な位相 $e^{ic}$、$A$ を変えない）。$U(1)$ では「グリボフの曖昧さ」は生じない。

(3) ホッジ分解 $A=h+d\lambda'+\delta\beta$ をとると、 $\delta A=\delta d\lambda'=\Delta\lambda'=0$ なので $d\lambda'=0$。よって $A=h+\delta\beta$。$T^3$（平坦）では問題 \ref{ch05-prob-torus} と同様に $h=\sum_ia_i\,dx^i$（$a_i$ 定数）。$h$ はゲージ変換 $A\to A+d\lambda$（$\lambda$ 一価）で変わらない（$d\lambda$ は $\mathcal H^1$ と直交）。しかし $e^{i\lambda}$, $\lambda=2\pi\bm m\cdot\bm x/L$ は $T^3$ 上一価な $U(1)$ 値関数であり（$x^i\to x^i+L$ で $\lambda\to\lambda+2\pi m_i$）、$A\to A+d\lambda$ で $a_i\to a_i+2\pi m_i/L$。物理的な自由度は $e^{ia_iL}$（$i$ 方向に一周するウィルソン・ループ）で、$(\R/\tfrac{2\pi}{L}\Z)^3\cong T^3$ に値をとる。これは $b_1(T^3)=3$ 個のゲージ場のゼロモード（\en{toron}）で、有限体積のQEDでは特別な扱いが必要になる。

(4) $\bm E$ に対応する2形式 $\mathcal E={*_3}(E_idx^i)$ を用いるとガウスの法則は $d\mathcal E=\rho\,\mathrm{vol}$。必要性：ストークスの定理より $\int_M\rho\,\mathrm{vol}=\int_Md\mathcal E=0$（$\partial M=\emptyset$）。十分性：$\int\rho=0$ なら (2) と同じ理由で $\Delta\varphi=\rho$（$\Delta=-\nabla^2$ の平坦の場合、すなわち $-\nabla^2\varphi=\rho$）の解 $\varphi$ が存在し、$\bm E=-\nabla\varphi$ は $\nabla\cdot\bm E=-\nabla^2\varphi=\rho$ をみたす。
一意性：$\varphi$ は定数を除いて一意なので $\bm E=-\nabla\varphi$ は一意。ガウスの法則だけを課した一般の解は、$d\mathcal E'=0$ をみたす $\mathcal E'$ を加える不定性があり、静電場（$\nabla\times\bm E=0$、すなわち $E_idx^i$ も閉）まで課すと差は閉かつ余閉、つまり調和1形式 $\sum c_idx^i$（一様電場、$b_1=3$ 個）だけが残る。これは「トーラスを一周する電気フラックス」のセクターである。
格子QED：周期境界条件の空間はトーラスで閉多様体なので、ガウスの法則（の格子版）から全電荷はゼロでなければならない。1個の電荷を置くと矛盾し、一様な中和背景電荷を加えるか、荷電共役で捻った境界条件（C周期境界条件）などの工夫が必要になる。
\end{solution}

\begin{problem}[難易度 C]\label{ch05-prob-kk}
5次元時空 $M_4\times S^1$（座標 $(x^\mu,y)$、$y\sim y+2\pi R$、計量 $\eta_{\mu\nu}dx^\mu dx^\nu+dy^2$）上のマクスウェル理論
\[
S=-\frac1{2g_5^2}\int F\wedge{*F}=-\frac{1}{4g_5^2}\int d^4x\,dy\,F_{MN}F^{MN}
\]
を考える。\\
(1) $A_\mu(x,y)=\sum_{n\in\Z}A^{(n)}_\mu(x)e^{iny/R}$, $A_y(x,y)=\sum_n\phi^{(n)}(x)e^{iny/R}$ と展開する。$n\neq0$ のモードではゲージ変換で $\phi^{(n)}=0$ にできることを示し、そのとき $A^{(n)}_\mu$ が質量 $\lvert n\rvert/R$ のベクトル場になることを示せ。\\
(2) $n=0$ のモードが4次元の質量ゼロの光子 $A_\mu^{(0)}$ と質量ゼロのスカラー $\phi^{(0)}$ を与えることを示し、4次元の結合定数 $g_4$ を $g_5$, $R$ で表せ。これらが $S^1$ の調和0形式・調和1形式（$b_0=b_1=1$）に対応することを説明せよ。\\
(3) スカラーがウィルソン線 $W(x)=\oint_{S^1}A=2\pi R\,\phi^{(0)}(x)$ と同一視でき、大きなゲージ変換によって $W\sim W+2\pi$ と周期的（コンパクトなスカラー）であることを示せ。
\end{problem}

\begin{solution}
(1) 5次元のゲージ変換 $A\to A+d\lambda$、$\lambda=\sum_n\lambda^{(n)}(x)e^{iny/R}$ で $A_y\to A_y+\partial_y\lambda$、すなわち $\phi^{(n)}\to\phi^{(n)}+\frac{in}{R}\lambda^{(n)}$。$n\neq0$ なら $\lambda^{(n)}=\frac{iR}{n}\phi^{(n)}$ と選べば $\phi^{(n)}=0$（スカラーは「食べられる」）。このとき $F_{\mu y}=\partial_\mu A_y-\partial_yA_\mu$ の $n$ モードは $-\frac{in}{R}A^{(n)}_\mu$。
$F_{MN}F^{MN}=F_{\mu\nu}F^{\mu\nu}+2F_{\mu y}F^{\mu y}$（$g^{yy}=1$）なので、ラグランジアン密度は $-\frac1{4g_5^2}F_{\mu\nu}F^{\mu\nu}-\frac1{2g_5^2}F_{\mu y}F^{\mu}{}_y$。$y$ で積分すると $\int_0^{2\pi R}e^{i(n+m)y/R}dy=2\pi R\,\delta_{n,-m}$、$A^{(-n)}=\overline{A^{(n)}}$（$A$ は実）より、$n$ モードの部分は
\[
\frac{2\pi R}{g_5^2}\Bigl[-\frac12\overline{F^{(n)}_{\mu\nu}}F^{(n)\mu\nu}-\frac{n^2}{R^2}\overline{A^{(n)}_\mu}A^{(n)\mu}\Bigr]\qquad(n>0\text{ と }-n\text{ をまとめた})
\]
となる。mostly plus 計量でのプロカ場（複素）のラグランジアン $-\frac12\lvert F_{\mu\nu}\rvert^2-m^2\lvert A_\mu\rvert^2$ と比べて $m=\lvert n\rvert/R$。正準規格化は全体の因子 $2\pi R/g_5^2$ で場を再定義すればよい。

(2) $n=0$：$\partial_y$ はゼロに作用するので $F_{\mu y}=\partial_\mu\phi^{(0)}$。
\[
S\supset2\pi R\int d^4x\Bigl[-\frac{1}{4g_5^2}F^{(0)}_{\mu\nu}F^{(0)\mu\nu}-\frac{1}{2g_5^2}\partial_\mu\phi^{(0)}\partial^\mu\phi^{(0)}\Bigr].
\]
質量項はなく、光子と質量ゼロのスカラー。$\frac1{g_4^2}=\frac{2\pi R}{g_5^2}$。$\phi^{(0)}$ に対するゲージ変換は $\lambda^{(0)}$ が $y$ によらないので作用せず、物理的なスカラーである（ただし (3) の大きなゲージ変換は除く）。
展開 $A=\sum A^{(n)}_\mu dx^\mu e^{iny/R}+\sum\phi^{(n)}e^{iny/R}dy$ は $A=\sum_q c_q\wedge\omega_{1-q}$ の形で、$S^1$ のラプラシアン $-\partial_y^2$ の固有値 $n^2/R^2$ が質量の2乗。固有値ゼロ（調和形式）は0形式では定数 $1$（$b_0=1$）、1形式では $dy$（$b_1=1$）。前者に掛かる $A^{(0)}_\mu$ が光子、後者に掛かる $\phi^{(0)}$ がスカラー。

(3) $\oint_{S^1}A=\int_0^{2\pi R}A_y\,dy=2\pi R\,\phi^{(0)}$（$n\neq0$ は積分で消える）で、ゲージ不変な量 $e^{iq\oint A}$ はウィルソン線（$S^1$ 方向のホロノミー）。ゲージ変換 $e^{i\lambda}$, $\lambda=my/R$（$m\in\Z$）は $S^1$ 上一価で、$A\to A+\frac mR dy$、$\phi^{(0)}\to\phi^{(0)}+\frac mR$、$W\to W+2\pi m$。よって $W$ は周期 $2\pi$ のコンパクトなスカラー（$H^1(S^1;\R)/2\pi H^1(S^1;\Z)$、問題 \ref{ch05-prob-holonomy}）。古典的には平坦方向（モジュライ）で、$W$ の値に応じて電荷 $q$ の場の KK 質量が $\lvert n-qW/2\pi\rvert/R$ にずれる（$\partial_y-iqA_y$ の固有値から）（ホソタニ機構の原型）。
\end{solution}

\begin{problem}[難易度 D]\label{ch05-prob-kkgen}
$M_d\times K$（$K$ はコンパクトで向きづけられた $k$ 次元リーマン多様体、積計量）上の $p$ 形式ゲージ場 $C$（場の強さ $G=dC$、ゲージ変換 $C\to C+d\Lambda$、作用 $-\frac12\int G\wedge{*G}$）を考える。$K$ のホッジ分解と、積計量では（ユークリッド化して）$\Delta_{M\times K}(\alpha\wedge\beta)=\Delta_M\alpha\wedge\beta+\alpha\wedge\Delta_K\beta$ が成り立つことを用いてよい。\\
(1) $C=\sum_q\sum_ic^{(i)}_q(x)\wedge\omega^{(i)}_{p-q}(y)$（$\omega^{(i)}_{p-q}$ は $\Delta_K$ の固有形式）と展開し、ゲージ $\delta_KC=0$ の下で運動方程式が $(\Delta_M+\Delta_K)C=0$（の各成分）になることを認めて、$d$ 次元で質量ゼロの $q$ 形式場の数が $b_{p-q}(K)$ であることを説明せよ。\\
(2) $K=T^6$、$p=2$（10次元の2形式場 $B$）のとき、4次元の質量ゼロ場（2形式、ベクトル、スカラー）の数を求めよ。\\
(3) $K=S^2\times S^2$（キュネットの公式を用いて $b_p$ を求めよ）、$p=3$ のとき、$d=7$ 次元の質量ゼロ場の数を求めよ。\\
(4) 質量ゼロの場の数が $K$ の計量（大きさや形）によらないのはなぜか。$K$ を $S^6$ に替えると (2) の答えはどうなるか。
\end{problem}

\begin{solution}
(1) $\Delta_K\omega=m^2\omega$ の固有形式に掛かる係数 $c_q$ は $(\Delta_M+m^2)c_q=0$、すなわち（ローレンツ符号に戻して）$(-\Box+m^2)c_q=0$ をみたし、質量 $m$ の $q$ 形式場になる。質量ゼロは $\Delta_K\omega=0$、すなわち調和形式。ホッジの定理により調和 $(p-q)$ 形式の空間の次元は $b_{p-q}(K)$。したがって質量ゼロの $q$ 形式場は $b_{p-q}(K)$ 個（$q=0,\dots,p$）。
補足：$d$ 次元側にも残るゲージ変換 $c_q\to c_q+d\Lambda_{q-1}$（$\Lambda=\Lambda_{q-1}\wedge\omega_{p-q}$、$\omega$ 調和）があり、$c_q$ は $d$ 次元の $q$ 形式ゲージ場として振る舞う。$\delta_K C=0$ は $C$ の内部成分から $d_K$ の像（純ゲージ）を除く条件で、残る非調和モードは質量をもつ。

(2) $b_p(T^6)=\binom6p$：$b_0=1$, $b_1=6$, $b_2=15$。よって4次元で2形式場 $B_{\mu\nu}$ が $1$ 個（4次元では質量ゼロの2形式はスカラー（アクシオン）と双対）、ベクトル場 $B_{\mu i}$ が $6$ 個、スカラー $B_{ij}$ が $15$ 個。最後の15個は $T^6$ 上の定数2形式（$B$ 場のモジュライ）である。

(3) $S^2$ の $b=(1,0,1)$。キュネットより $b_k(S^2\times S^2)=\sum_{p+q=k}b_p(S^2)b_q(S^2)$：
\[
b_0=1,\quad b_1=0,\quad b_2=2,\quad b_3=0,\quad b_4=1 .
\]
$p=3$：質量ゼロの3形式 $b_0=1$ 個、2形式 $b_1=0$ 個、ベクトル $b_2=2$ 個（$S^2$ 各因子の面積形式に対応）、スカラー $b_3=0$ 個。

(4) ベッチ数はトポロジー（ド・ラーム・コホモロジー＝特異コホモロジー）だけで決まり、計量によらないから。計量を変えると調和形式そのもの（どの形式が調和か）は変わるが、その数は変わらず、質量をもつモードの質量だけが変わる。$K=S^6$ なら $b_0=1$, $b_1=b_2=0$ なので、$B_{\mu\nu}$ が1個だけで、ベクトルもスカラーも出ない。弦理論のコンパクト化でモジュライや $U(1)$ ゲージ場の数がホッジ数（$b_p$ の精密化）で決まる、という第10章の議論の原型である。
\end{solution}

%================================================
\chapter{コホモロジー環・ポアンカレ双対・交叉形式}\label{ch:06}
% 18問
%================================================

\section{物理からの動機}

\subsection*{位相項は「流束の積」である}
4次元のゲージ理論には、計量によらない\term{位相項}{topological term}
\[
S_\theta=\theta\,k,\qquad k=\frac{1}{8\pi^2}\int_X\tr F\wedge F
\]
を加えることができる（記法と規約の $k$）。$\R^4$ 上の $SU(N)$ インスタントンでは $k\in\Z$ なので $e^{iS_\theta}$ は $\theta\to\theta+2\pi$ で不変であり、$\theta$ は角度変数である。ところが $U(1)$ ゲージ理論をコンパクトな4次元多様体 $X$ 上で考えると、$F/2\pi$ はそれ自体が整数係数のコホモロジー類 $c\in H^2(X;\Z)$ を代表し（第5章のディラック量子化）、
\[
k=\frac{1}{8\pi^2}\int_X F\wedge F=\frac12\int_X\frac{F}{2\pi}\wedge\frac{F}{2\pi}=\frac12\,Q(c,c)
\]
となる。ここで $Q(c,c)$ は「2次のコホモロジー類どうしの積を $X$ 上で積分した整数」で、$X$ の\textbf{交叉形式}と呼ばれる（本章）。$Q(c,c)$ が常に偶数なら $k\in\Z$ で $\theta$ の周期は $2\pi$、奇数になりうるなら $k\in\tfrac12\Z$ で周期は $4\pi$ になる（問題 \ref{ch06-prob-theta}）。つまり「$\theta$ の周期」という物理量が、時空の\textbf{コホモロジーの積構造}で決まる。

$\int F\wedge F$ のほか、チャーン・サイモンズ項 $\int A\wedge dA$、BF 項 $\int B\wedge F$、離散ゲージ場の位相作用 $\pi\int a\cup a$ など、位相項はほとんどが「流束（コホモロジー類）の積」の形をしている。微分形式の $\wedge$ を、整数係数や $\Z_N$ 係数でも意味をもつように拡張したものが\textbf{カップ積}である。

\subsection*{ねじれと離散ゲージ場}
$\Z_N$ ゲージ場（格子上の $\Z_N$ 値のリンク変数で、各プラケットのホロノミーが $1$ のもの）はゲージ同値類として $H^1(X;\Z_N)$ の元であり、微分形式では表せない。整数係数のコホモロジーの\term{ねじれ部分}{torsion}（有限位数の元）も同様に、ド・ラーム・コホモロジー（第5章）には見えない。こうした「流束はゼロなのにアハラノフ・ボーム位相だけをもつ」配位を扱う道具が\textbf{普遍係数定理}と係数を変えたコホモロジーである。

\subsection*{ポアンカレ双対：ブレーンとサイクル}
弦理論では、Dブレーンが内部空間の部分多様体（サイクル）$C$ に巻き付き、その電荷はホモロジー類 $[C]$ で決まる。ブレーンがつくる場の強さは $C$ の\textbf{ポアンカレ双対}な微分形式 $\eta_C$ で表され、2枚のブレーンの交点の数（\textbf{交叉数}）$\int\eta_{C}\wedge\eta_{C'}$ が交点に局在する massless なカイラル・フェルミオンの数を与える（下図、問題 \ref{ch06-prob-branes}）。

\begin{center}
\begin{tikzpicture}[scale=1.0]
  \draw[thick] (0,0) ellipse (2.6 and 1.3);
  \draw[thick] (-1.0,0.12) .. controls (-0.5,-0.32) and (0.5,-0.32) .. (1.0,0.12);
  \draw[thick] (-0.8,-0.02) .. controls (-0.4,0.25) and (0.4,0.25) .. (0.8,-0.02);
  \draw[very thick,blue!70!black,mid arrow] (0,-0.83) arc (-90:270:1.85 and 0.83);
  \draw[very thick,red!70!black] (0,-0.25) .. controls (-0.25,-0.5) and (-0.25,-1.05) .. (0,-1.3);
  \draw[very thick,red!70!black,dashed] (0,-0.25) .. controls (0.25,-0.5) and (0.25,-1.05) .. (0,-1.3);
  \node[vtx] at (-0.17,-0.83) {};
  \node[blue!70!black] at (2.15,0.55) {$a$};
  \node[red!70!black] at (-0.45,-1.15) {$b$};
  \node at (4.6,0.3) {\footnotesize ブレーン $D_1$ が $a$ に、};
  \node at (4.6,-0.1) {\footnotesize $D_2$ が $b$ に巻き付く};
  \node at (4.6,-0.6) {\footnotesize 交叉数 $a\cdot b=1$};
\end{tikzpicture}
\zucap{トーラス上の2つのサイクル $a$（青）と $b$（赤、点線は裏側）は1点で交わる。$a,b$ に巻き付いたブレーンの交点には、交叉数の数だけカイラルなモードが局在する。}
\end{center}

\subsection*{絡み数：ガウス積分とエニオン}
3次元空間内の互いに交わらない2つの閉曲線 $C_1,C_2$ に対し、ガウスの積分
\[
\mathrm{Lk}(C_1,C_2)=\frac{1}{4\pi}\oint_{C_1}\!\oint_{C_2}\frac{(\bm x_1-\bm x_2)\cdot(d\bm x_1\times d\bm x_2)}{\lvert\bm x_1-\bm x_2\rvert^3}
\]
は整数（\textbf{絡み数}）である。これは「$C_1$ を流れる単位電流がつくる磁場を $C_2$ に沿って周回積分したもの」であり、アンペールの法則により $C_1$ が $C_2$ の張る面を貫く回数に等しい。アーベリアン・チャーン・サイモンズ理論では、電荷 $q_1,q_2$ のウィルソン・ループの期待値が $\exp\bigl(\mp\frac{2\pi i}{k}q_1q_2\,\mathrm{Lk}\bigr)$ となり、これがエニオンの組みひも統計の位相を与える（問題 \ref{ch06-prob-cs}）。
\begin{center}
\begin{tikzpicture}[scale=1.1]
  \draw[very thick,blue!70!black] (0,0) circle (1);
  \draw[very thick,red!70!black] (1.2,0) circle (1);
  % 上の交点：C2 が上
  \draw[white,line width=5pt] ([shift={(1.2,0)}]110:1) arc (110:145:1);
  \draw[very thick,red!70!black] ([shift={(1.2,0)}]105:1) arc (105:150:1);
  % 下の交点：C1 が上
  \draw[white,line width=5pt] (-70:1) arc (-70:-36:1);
  \draw[very thick,blue!70!black] (-75:1) arc (-75:-31:1);
  \draw[very thick,blue!70!black,-{Stealth[length=2.5mm]}] (175:1) arc (175:185:1);
  \draw[very thick,red!70!black,-{Stealth[length=2.5mm]}] ([shift={(1.2,0)}]5:1) arc (5:-5:1);
  \node[blue!70!black] at (-1.35,0.6) {$C_1$};
  \node[red!70!black] at (2.6,0.6) {$C_2$};
  \node at (0.6,-1.5) {\footnotesize ホップ絡み目：$\mathrm{Lk}(C_1,C_2)=\pm1$};
\end{tikzpicture}
\zucap{ホップ絡み目。2つの円周は連続変形でほどけない。符号は曲線の向きで決まる。}
\end{center}

\begin{physbox}[：本章の辞書]
\begin{center}
\small
\begin{tabular}{p{0.47\linewidth}p{0.43\linewidth}}
物理 & 数学\\ \hline
量子化された流束 $[F/2\pi]$ & $H^2(X;\Z)$ の元（の実像）\\
$\Z_N$ ゲージ場（平坦） & $H^1(X;\Z_N)=\Hom(\pi_1(X),\Z_N)$\\
流束ゼロでも AB 位相をもつ $U(1)$ 束 & $H^2(X;\Z)$ のねじれ元\\
位相項 $\int F\wedge F$, $\int A\wedge dA$, $\pi\int a\cup a$ & カップ積（とその積分）\\
ブレーンの電荷・巻き付き & ホモロジー類とポアンカレ双対形式\\
ブレーン交点のカイラル・モードの数 & 交叉数\\
$U(1)$ インスタントン数・$\theta$ の周期 & 4次元多様体の交叉形式 $Q$、偶奇性\\
チャーン・サイモンズのウィルソン・ループ & 絡み数
\end{tabular}
\end{center}
\end{physbox}

\section{基本事項のまとめ}

\subsection{特異コホモロジーと普遍係数定理}

\begin{mydef}{特異コホモロジー}{ch06-coh}
$G$ をアーベル群とする。特異鎖複体 $(C_\bullet(X),\partial)$（第4章）に対し、$C^k(X;G)=\Hom(C_k(X),G)$ を\term{$k$ コチェイン}{cochain}の群とし、\term{余境界}{coboundary} $\delta\colon C^k\to C^{k+1}$ を $(\delta\varphi)(\sigma)=\varphi(\partial\sigma)$ で定める。$\delta^2=0$ であり
\[
H^k(X;G)=\ker\delta/\operatorname{im}\delta
\]
を $G$ 係数の\term{特異コホモロジー}{singular cohomology}という。単体複体・胞体複体についても同様に定義でき、同じ群を与える。$H^k$ は反変的：$f\colon X\to Y$ は $f^*\colon H^k(Y;G)\to H^k(X;G)$ を誘導し、ホモトピー不変である。
\end{mydef}
直感的には、$k$ コチェインは「$k$ 次元の小片それぞれに値を割り当てるもの」で、格子上の $k$ 形式場にあたる。$\delta$ は格子上の外微分であり、$\delta\varphi=0$ が「閉形式」、$\varphi=\delta\lambda$ が「完全形式（ゲージ自由度）」である。

\begin{mythm}{ド・ラームの定理（第5章）}{ch06-derham}
滑らかな多様体 $M$ では、微分形式 $\omega$ に「各滑らかな特異単体上の積分」を対応させる写像が同型 $H^k_{\mathrm{dR}}(M)\cong H^k(M;\R)$ を誘導する。この同型のもとで $\wedge$ は次節のカップ積に対応する。
\end{mythm}

\begin{mythm}{普遍係数定理}{ch06-uct}
任意の空間 $X$ とアーベル群 $G$ に対し、分裂する完全列
\[
0\to\Ext(H_{k-1}(X),G)\to H^k(X;G)\to\Hom(H_k(X),G)\to0
\]
がある。したがって $H^k(X;G)\cong\Hom(H_k(X),G)\oplus\Ext(H_{k-1}(X),G)$（自然でない同型）。
\end{mythm}
有限生成アーベル群に対しては次の表だけで計算できる（$\Hom$, $\Ext$ はともに直和を直和に移す）。
\begin{center}
\renewcommand{\arraystretch}{1.25}
\begin{tabular}{c|cc}
 & $\Hom(\cdot,G)$ & $\Ext(\cdot,G)$\\ \hline
$\Z$ & $G$ & $0$\\
$\Z_n$ & $\{g\in G:ng=0\}$ & $G/nG$
\end{tabular}
\end{center}
特に $\Hom(\Z_n,\Z)=0$、$\Ext(\Z_n,\Z)=\Z_n$、$\Hom(\Z_n,\Z_N)\cong\Ext(\Z_n,\Z_N)\cong\Z_{\gcd(n,N)}$。

\begin{myprop}{整数係数：ねじれは1つ上にずれる}{ch06-shift}
$H_k(X)\cong\Z^{b_k}\oplus T_k$（$T_k$ は有限）なら
\[
H^k(X;\Z)\cong\Z^{b_k}\oplus T_{k-1}.
\]
$b_k$ を\term{ベッティ数}{Betti number}という。体 $\mathbb F$ 係数では $\Ext(\cdot,\mathbb F)=0$ なので $H^k(X;\mathbb F)\cong\Hom(H_k(X),\mathbb F)$ であり、$\mathbb F=\R,\Q$ なら $\dim H^k(X;\R)=b_k$。
\end{myprop}
\begin{note}
$\Ext(\Z_n,\Z)=\Z_n$ は、自由分解 $0\to\Z\xrightarrow{\times n}\Z\to\Z_n\to0$ に $\Hom(\cdot,\Z)$ を施した $\Z\xrightarrow{\times n}\Z$ の余核として求まる。一般に $\Ext(\Z_n,G)=\coker(G\xrightarrow{\times n}G)=G/nG$、$\Hom(\Z_n,G)=\ker(G\xrightarrow{\times n}G)$ である。
\end{note}

\begin{myex}{代表例}{ch06-uctex}
\begin{itemize}
\item $\RP^n$：$H_k=\Z\ (k=0)$, $\Z_2\ (k\text{ 奇},\ 0<k<n)$, $\Z\ (k=n\text{ 奇})$、それ以外 $0$（第4章）。よって $H^k(\RP^n;\Z)=\Z\ (k=0)$, $\Z_2\ (k\text{ 偶},\ 0<k\le n)$, $\Z\ (k=n\text{ 奇})$、それ以外 $0$（問題 \ref{ch06-prob-uct}）。
\item \term{レンズ空間}{lens space} $L(p;q)=S^3/\Z_p$（$(z_1,z_2)\mapsto(\zeta z_1,\zeta^qz_2)$, $\zeta=e^{2\pi i/p}$, $\gcd(p,q)=1$）：$H_*=(\Z,\Z_p,0,\Z)$、$H^*(\,\cdot\,;\Z)=(\Z,0,\Z_p,\Z)$。$L(2;1)=\RP^3$。
\end{itemize}
\end{myex}

\subsection{$\Z_N$ 係数と離散ゲージ場}

\begin{myprop}{$H^1$ と基本群}{ch06-h1}
弧状連結な $X$ とアーベル群 $G$ に対し $H^1(X;G)\cong\Hom(H_1(X),G)\cong\Hom(\pi_1(X),G)$。
\end{myprop}
$\Ext(H_0,G)=\Ext(\Z,G)=0$ なので第1の同型は普遍係数定理から、第2の同型は $H_1=\pi_1^{\mathrm{ab}}$（フレヴィッツの定理、第4章）と $G$ の可換性から従う。

\begin{physbox}[：平坦な $\Z_N$ ゲージ場]
単体複体（格子）$X$ の各向き付き辺 $e$ に $U_e=e^{2\pi ia_e/N}$（$a_e\in\Z_N$）を置く。プラケット（2単体）$\sigma$ のホロノミー $\prod_{e\subset\partial\sigma}U_e=e^{2\pi i(\delta a)(\sigma)/N}$ が $1$ であること（平坦性）は $\delta a=0$、頂点 $v$ でのゲージ変換 $U_e\mapsto e^{2\pi i\lambda_{v'}/N}U_e\,e^{-2\pi i\lambda_v/N}$（$e=[v,v']$）は $a\mapsto a+\delta\lambda$ である。したがって
\[
\{\text{平坦な }\Z_N\text{ ゲージ場}\}/\{\text{ゲージ変換}\}=H^1(X;\Z_N)\cong\Hom(\pi_1(X),\Z_N),
\]
各類はループに沿ったホロノミー（$\Z_N$ のウィルソン・ループ）で完全に指定される。2次元空間 $\Sigma$ 上の $\Z_N$ ゲージ理論（トーリック符号）の基底状態はこの集合で張られ、縮退度は $\lvert H^1(\Sigma;\Z_N)\rvert$ である（問題 \ref{ch06-prob-zn}）。
\end{physbox}

\begin{mythm}{係数の完全列（ボックシュタイン列）}{ch06-bock}
アーベル群の短完全列 $0\to A\to B\to C\to0$ は長完全列
\[
\cdots\to H^k(X;A)\to H^k(X;B)\to H^k(X;C)\xrightarrow{\ \beta\ }H^{k+1}(X;A)\to\cdots
\]
を誘導する。特に $0\to\Z\to\R\xrightarrow{\exp 2\pi i(\cdot)}U(1)\to0$ から
\[
H^1(X;\R)\to H^1(X;U(1))\xrightarrow{\ \beta\ }H^2(X;\Z)\to H^2(X;\R).
\]
$H^1(X;U(1))=\Hom(\pi_1(X),U(1))$ は平坦な $U(1)$ 接続のゲージ同値類（ホロノミー）であり、$\beta$ はその下にある直線束のチャーン類 $c_1$ を与える（符号は規約による；直線束の分類は第7章）。
\end{mythm}
したがって $\beta$ の像は $H^2(X;\Z)\to H^2(X;\R)$ の核、すなわち $H^2(X;\Z)$ のねじれ部分である。ねじれた $c_1$ をもつ直線束は平坦な接続をもち、$F=0$ なのに非自明である（問題 \ref{ch06-prob-torsion}）。

\subsection{カップ積とコホモロジー環}

\begin{mydef}{カップ積}{ch06-cup}
可換環 $R$ 係数の $\varphi\in C^p(X;R)$, $\psi\in C^q(X;R)$ に対し、特異単体 $\sigma\colon[v_0,\dots,v_{p+q}]\to X$ 上で
\[
(\varphi\smile\psi)(\sigma)=\varphi\bigl(\sigma|_{[v_0,\dots,v_p]}\bigr)\,\psi\bigl(\sigma|_{[v_p,\dots,v_{p+q}]}\bigr)
\]
と定める（「前の $p$ 次元面」と「後の $q$ 次元面」）。$\delta(\varphi\smile\psi)=\delta\varphi\smile\psi+(-1)^p\varphi\smile\delta\psi$ が成り立つので、コホモロジー上の積
\[
\smile\colon H^p(X;R)\times H^q(X;R)\to H^{p+q}(X;R)
\]
が定まる。$H^*(X;R)=\bigoplus_kH^k(X;R)$ はこの積で環になり、\term{コホモロジー環}{cohomology ring}と呼ばれる。
\end{mydef}

\begin{mythm}{カップ積の性質}{ch06-cupprop}
\begin{enumerate}
\item 結合的で、$1\in H^0$（各点に $1$ を与える類）が単位元。
\item 自然性：$f^*(a\smile b)=f^*a\smile f^*b$。特にホモトピー同値写像はコホモロジー環の同型を誘導する。
\item \term{次数付き可換性}{graded commutativity}：$a\in H^p$, $b\in H^q$ なら $a\smile b=(-1)^{pq}\,b\smile a$。
\item 滑らかな多様体ではド・ラームの同型のもとで $[\omega]\smile[\eta]=[\omega\wedge\eta]$（$\R$ 係数）。
\end{enumerate}
\end{mythm}
(3) から、奇数次の $a$ について $2(a\smile a)=0$。$\Z$ 係数で $H^{2p}$ にねじれがなければ $a\smile a=0$ であるが、$\Z_2$ 係数では $-1=1$ なので $a\smile a\neq0$ もありうる（$\RP^n$）。

\begin{mythm}{代表的なコホモロジー環}{ch06-rings}
\begin{align*}
&H^*(S^n;\Z)=\Z[x]/(x^2),\ \lvert x\rvert=n; &&H^*(T^n;\Z)=\Lambda_\Z[x_1,\dots,x_n],\ \lvert x_i\rvert=1;\\
&H^*(\CP^n;\Z)=\Z[x]/(x^{n+1}),\ \lvert x\rvert=2; &&H^*(\RP^n;\Z_2)=\Z_2[\alpha]/(\alpha^{n+1}),\ \lvert\alpha\rvert=1.
\end{align*}
種数 $g$ の閉曲面 $\Sigma_g$ では、$H^1$ の基底 $\alpha_1,\beta_1,\dots,\alpha_g,\beta_g$ と $H^2$ の生成元 $\omega$ を適当にとると $\alpha_i\smile\beta_j=\delta_{ij}\omega=-\beta_j\smile\alpha_i$、$\alpha_i\smile\alpha_j=\beta_i\smile\beta_j=0$。
\end{mythm}
$T^n$ では $x_i=[dx^i]$（$x^i$ は周期 $1$ の座標）とみればよく、外積代数は $dx^i\wedge dx^j=-dx^j\wedge dx^i$ そのものである。$\CP^n$ の $x$ は超平面 $\CP^{n-1}$ のポアンカレ双対（次節）で、$x^n$ が「$n$ 枚の一般の超平面は1点で交わる」ことを表す。

\begin{mythm}{キュネットの公式}{ch06-kunneth}
$\mathbb F$ を体とすると、$a\otimes b\mapsto p_1^*a\smile p_2^*b$ は環の同型
\[
H^*(X;\mathbb F)\otimes H^*(Y;\mathbb F)\xrightarrow{\ \cong\ }H^*(X\times Y;\mathbb F)
\]
を与える（テンソル積の積は $(a\otimes b)(a'\otimes b')=(-1)^{\lvert b\rvert\lvert a'\rvert}aa'\otimes bb'$）。$H^*(Y;\Z)$ が有限生成自由なら $\Z$ 係数でも成り立つ。
\end{mythm}
例：$H^*(S^m\times S^n;\Z)$ は $1,x,y,xy$（$x=p_1^*\sigma_m$, $y=p_2^*\sigma_n$）を基底とし、$x^2=y^2=0$、$yx=(-1)^{mn}xy$。$T^n=(S^1)^n$ に繰り返し使えば $T^n$ の外積代数を得る。

\subsection{ポアンカレ双対と交叉数}

\begin{mydef}{基本類・キャップ積}{ch06-cap}
連結な $n$ 次元閉多様体 $M$ が向き付け可能なら、$H_n(M;\Z)\cong\Z$ であり、向きを選ぶと生成元 $[M]$（\term{基本類}{fundamental class}）が定まる。向き付け不能でも $H_n(M;\Z_2)\cong\Z_2$ で $\Z_2$ 係数の基本類がある。\term{キャップ積}{cap product} $\frown\colon H_n(M;R)\times H^k(M;R)\to H_{n-k}(M;R)$ は鎖のレベルで $\sigma\frown\varphi=\varphi(\sigma|_{[v_0,\dots,v_k]})\,\sigma|_{[v_k,\dots,v_n]}$ で定義され、$\langle\psi,\,x\frown\varphi\rangle=\langle\varphi\smile\psi,\,x\rangle$ をみたす。
\end{mydef}

\begin{mythm}{ポアンカレ双対}{ch06-pd}
$M$ を向き付けられた $n$ 次元閉多様体とする。
\begin{enumerate}
\item $D=[M]\frown\cdot\colon H^k(M;\Z)\to H_{n-k}(M;\Z)$ は同型。$\Z_2$ 係数なら向き付け不能でも成り立つ。
\item 体 $\mathbb F$ 係数で、双線形形式 $H^k(M;\mathbb F)\times H^{n-k}(M;\mathbb F)\to\mathbb F$, $(a,b)\mapsto\langle a\smile b,[M]\rangle$ は非退化。$\Z$ 係数ではねじれを除いた $H^k/T\times H^{n-k}/T\to\Z$ が\term{ユニモジュラー}{unimodular}（行列式 $\pm1$）である。
\item 特に $b_k=b_{n-k}$。
\end{enumerate}
\end{mythm}

\begin{mydef}{ポアンカレ双対形式・交叉数}{ch06-eta}
$M$ を向き付けられた $n$ 次元閉多様体、$C\subset M$ を向き付けられた $k$ 次元閉部分多様体とする。閉 $(n-k)$ 形式 $\eta_C$ で、任意の閉 $k$ 形式 $\omega$ に対し
\[
\int_C\omega=\int_M\omega\wedge\eta_C
\]
となるものの類を $C$ の\term{ポアンカレ双対}{Poincaré dual}という（$\eta_C$ は $C$ の近くに集中した「デルタ関数的な形式」にとれる）。次元の和が $n$ の部分多様体 $A,B$ が横断的に交わるとき、各交点 $p$ で $T_pA\oplus T_pB$ の向きが $T_pM$ の向きと一致すれば $+1$、逆なら $-1$ を与え、その和を\term{交叉数}{intersection number} $A\cdot B$ とする。このとき
\[
A\cdot B=\int_M\eta_A\wedge\eta_B=\langle PD(A)\smile PD(B),[M]\rangle
\]
が成り立ち、ホモロジー類だけで決まる。$A\cdot B=(-1)^{\dim A\dim B}B\cdot A$。
\end{mydef}
\begin{myex}{トーラス}{ch06-t2pd}
$T^2=\R^2/\Z^2$（向き $dx\wedge dy$）で $a$ を $x$ 方向、$b$ を $y$ 方向の円周とする。$\int_a dx=1=\int_{T^2}dx\wedge\eta_a$、$\int_a dy=0$ より $\eta_a=dy$。同様に $\int_b dy=1=\int dy\wedge\eta_b$ より $\eta_b=-dx$。よって $a\cdot b=\int dy\wedge(-dx)=\int dx\wedge dy=+1$、$b\cdot a=-1$。
\end{myex}

\begin{center}
\begin{tikzpicture}[scale=1.15]
  \draw[thick] (-3,0) .. controls (-3,1.4) and (-1.2,1.3) .. (0,0.8) .. controls (1.2,1.3) and (3,1.4) .. (3,0)
     .. controls (3,-1.4) and (1.2,-1.3) .. (0,-0.8) .. controls (-1.2,-1.3) and (-3,-1.4) .. (-3,0);
  \foreach \c in {-1.5,1.5}{
    \draw[thick] (\c-0.6,0.05) .. controls (\c-0.3,-0.25) and (\c+0.3,-0.25) .. (\c+0.6,0.05);
    \draw[thick] (\c-0.45,-0.05) .. controls (\c-0.2,0.15) and (\c+0.2,0.15) .. (\c+0.45,-0.05);
    \draw[very thick,blue!70!black,mid arrow] (\c,-0.5) arc (-90:270:1.0 and 0.5);
    \draw[very thick,red!70!black,mid arrow] (\c,-0.2) .. controls (\c-0.22,-0.45) and (\c-0.22,-0.8) .. (\c,-1.07);
    \draw[very thick,red!70!black,dashed] (\c,-0.2) .. controls (\c+0.22,-0.45) and (\c+0.22,-0.8) .. (\c,-1.07);
  }
  \node[blue!70!black] at (-1.5,0.72) {$a_1$}; \node[blue!70!black] at (1.5,0.72) {$a_2$};
  \node[red!70!black] at (-1.95,-0.95) {$b_1$}; \node[red!70!black] at (1.05,-0.95) {$b_2$};
\end{tikzpicture}
\zucap{種数2の曲面 $\Sigma_2$ の標準的なサイクル。$a_i$ と $b_i$ はちょうど1点で交わり（$a_i\cdot b_j=\delta_{ij}$）、$a_i\cdot a_j=b_i\cdot b_j=0$。$H_1(\Sigma_g)$ 上の交叉形式は標準シンプレクティック形式 $\begin{pmatrix}0&1\\-1&0\end{pmatrix}^{\oplus g}$ である。}
\end{center}

\begin{physbox}[：ブレーンの電荷とディラック対]
$p$ 次元のブレーン（$(p+1)$ 次元世界体積）は $(p+1)$ 形式ゲージ場 $C_{p+1}$ と $\int_{\text{世界体積}}C_{p+1}$ で結合し、ブレーンのつくる場の強さは $dG=\eta_{\text{ブレーン}}$（ポアンカレ双対形式）をみたす。内部空間 $Y$ のサイクルに巻き付いたブレーンは、低エネルギーで点粒子にみえ、その電荷はホモロジー類で決まる。例えば IIB 理論を6次元の $Y$ にコンパクト化すると、3サイクル $\gamma$ に巻き付いた D3 ブレーンは4次元のダイオンになり、2つのダイオンの\textbf{ディラック・シュウィンガー・ツヴァンツィガー対} $e_1g_2-e_2g_1$ は（規格化を除いて）交叉数 $\gamma_1\cdot\gamma_2$ で与えられる。3サイクルの交叉形式は反対称（$\gamma_1\cdot\gamma_2=-\gamma_2\cdot\gamma_1$）でユニモジュラーであり、電荷格子がシンプレクティック基底 $(A^I,B_I)$ をもつことが電磁双対性の数学的な表現である（第10章で扱う）。
\end{physbox}

\subsection{4次元多様体の交叉形式}

\begin{mydef}{交叉形式}{ch06-qform}
$X$ を向き付けられた4次元閉多様体とし、$L=H^2(X;\Z)/(\text{ねじれ})\cong\Z^{b_2}$ とする。
\[
Q_X\colon L\times L\to\Z,\qquad Q_X(a,b)=\langle a\smile b,[X]\rangle\ \Bigl(=\int_X\alpha\wedge\beta\Bigr)
\]
を\term{交叉形式}{intersection form}という（$\alpha,\beta$ は $a,b$ を代表する閉2形式）。ポアンカレ双対で $H_2$ に移せば、2次元曲面の交叉数 $Q(A,B)=A\cdot B$ である。
\end{mydef}
\begin{myprop}{交叉形式の基本性質}{ch06-qprop}
\begin{enumerate}
\item 対称（次数2は偶数）かつユニモジュラー（$\det Q=\pm1$、ポアンカレ双対）。
\item 実数上で対角化したときの正・負の固有値の個数を $b_+$, $b_-$ とし、$b_2=b_++b_-$、\term{符号数}{signature} $\sigma(X)=b_+-b_-$ と定める。リーマン計量を入れると、$b_\pm$ は自己双対・反自己双対な調和2形式の空間の次元に等しい（第5章のホッジ理論）。
\item $Q(a,a)$ がすべての $a$ で偶数のとき\term{偶}{even}、そうでないとき\term{奇}{odd}という。向きを逆にした $\bar X$ では $Q_{\bar X}=-Q_X$。
\end{enumerate}
\end{myprop}

\begin{myex}{代表的な交叉形式}{ch06-qex}
$H=\begin{pmatrix}0&1\\1&0\end{pmatrix}$、$E_8$ は $E_8$ のカルタン行列（問題 \ref{ch06-prob-e8}）とする。
\begin{center}
\renewcommand{\arraystretch}{1.25}
\begin{tabular}{lcccccc}
$X$ & $Q_X$ & $b_2$ & $b_+$ & $b_-$ & $\sigma$ & 偶奇\\ \hline
$S^4$ & （空） & $0$ & $0$ & $0$ & $0$ & 偶\\
$\CP^2$ & $(1)$ & $1$ & $1$ & $0$ & $1$ & 奇\\
$\overline{\CP^2}$ & $(-1)$ & $1$ & $0$ & $1$ & $-1$ & 奇\\
$S^2\times S^2$ & $H$ & $2$ & $1$ & $1$ & $0$ & 偶\\
$\CP^2\#\overline{\CP^2}$ & $(1)\oplus(-1)$ & $2$ & $1$ & $1$ & $0$ & 奇\\
$T^4$ & $3H$ & $6$ & $3$ & $3$ & $0$ & 偶\\
K3 曲面 & $3H\oplus2(-E_8)$ & $22$ & $3$ & $19$ & $-16$ & 偶
\end{tabular}
\end{center}
連結和 $X\#Y$ の交叉形式は $Q_X\oplus Q_Y$ である。K3 のオイラー標数は $\chi=24$。
\end{myex}

\begin{mythm}{偶奇性・符号数に関する事実（証明略）}{ch06-wu}
\begin{enumerate}
\item（ウー公式）$w_2(X)\in H^2(X;\Z_2)$ を第2スティーフェル・ホイットニー類とすると、すべての $a$ で $Q(a,a)\equiv\langle\bar a\smile w_2,[X]\rangle\pmod2$（$\bar a$ は $a$ の mod 2 還元）。したがって $X$ がスピン多様体（$w_2=0$）なら $Q_X$ は偶。$H_1(X;\Z)$ が2ねじれをもたなければ（例えば単連結なら）逆も成り立つ。
\item 偶なユニモジュラー形式の符号数は $8$ の倍数（ファン・デル・ブレイ）。
\item（ロホリン）滑らかな閉スピン4次元多様体では $\sigma\equiv0\pmod{16}$。第9章で指数定理から導く。
\item（セール）不定値のユニモジュラー形式は、奇なら $m(1)\oplus n(-1)$、偶なら $aH\oplus b E_8$（$b<0$ は $\lvert b\rvert(-E_8)$）に同型。定値の場合は分類が難しい（ドナルドソンの定理：滑らかな単連結閉多様体の定値な交叉形式は対角化可能）。
\end{enumerate}
\end{mythm}

\begin{physbox}[：$U(1)$ 流束とインスタントン数]
$U(1)$ 束の第1チャーン類 $c=c_1\in H^2(X;\Z)$ の実像はド・ラーム類 $[F/2\pi]$ であり、逆に任意の $c\in H^2(X;\Z)$ はある $U(1)$ 束の $c_1$ として実現される（第7章）。$H_2(X)/(\text{ねじれ})$ の基底となる曲面 $\Sigma_i$ を貫く流束を $\int_{\Sigma_i}F=2\pi n_i$（$n_i\in\Z$）とし、曲面の交叉行列を $M_{ij}=\Sigma_i\cdot\Sigma_j$ とする。双対基底 $e^j\in H^2$（$\int_{\Sigma_i}e^j=\delta_i^j$）で $c=\sum_in_ie^i$ と書け、$PD(\Sigma_i)=\sum_jM_{ij}e^j$ から $Q(e^i,e^j)=(M^{-1})_{ij}$ なので
\[
\frac{1}{4\pi^2}\int_XF\wedge F=Q_X(c,c)=\sum_{i,j}n_i\,(M^{-1})_{ij}\,n_j\in\Z
\]
（$M$ はユニモジュラーなので $M^{-1}$ も整数行列）。$U(1)$ インスタントン数は $k=\tfrac12Q_X(c,c)$ で、$Q_X$ が偶（例えばスピン多様体）なら整数、奇なら半整数もとる。
\end{physbox}

\subsection{絡み数}

\begin{mydef}{絡み数}{ch06-link}
$S^3$（または $\R^3$）内の互いに交わらない向き付き閉曲線 $C_1,C_2$ に対し、$\partial S_2=C_2$ となる向き付き曲面（ザイフェルト曲面）$S_2$ をとり、\term{絡み数}{linking number}を $\mathrm{Lk}(C_1,C_2)=C_1\cdot S_2$（$S_2$ の向きは、境界 $C_2$ の向きと右ねじの関係にある法線で定める）と定める。これは $S_2$ の選び方によらず、$\mathrm{Lk}(C_1,C_2)=\mathrm{Lk}(C_2,C_1)$ であり、ガウスの積分公式（動機の節）で表される。ガウスの公式の被積分関数は、写像 $\Psi\colon C_1\times C_2\to S^2$, $\Psi=(\bm x_2-\bm x_1)/\lvert\bm x_2-\bm x_1\rvert$ による $S^2$ の面積要素の引き戻しを $4\pi$ で割ったものであり、$\mathrm{Lk}=\deg\Psi\in\Z$ である。
\end{mydef}

\section{問題}

\subsection{普遍係数定理と離散ゲージ場}

\begin{problem}[難易度 A]\label{ch06-prob-uct}
$\RP^n$ の整数係数ホモロジー（例 \ref{ex:ch06-uctex}）を用いて、普遍係数定理から $H^k(\RP^n;\Z)$ を求めよ。$H_k$ のねじれと $H^k$ のねじれの次数がずれることを、$\RP^2$, $\RP^3$ について確かめよ。
\end{problem}

\begin{solution}
$H^k\cong\Hom(H_k,\Z)\oplus\Ext(H_{k-1},\Z)$ で、$\Hom(\Z_2,\Z)=0$、$\Ext(\Z_2,\Z)=\Z_2$、$\Ext(\Z,\Z)=0$ を使う。
\begin{itemize}
\item $k=0$：$\Hom(\Z,\Z)=\Z$。
\item $0<k<n$、$k$ 奇：$H_k=\Z_2$ なので $\Hom=0$。$H_{k-1}$ は $k-1=0$ なら $\Z$、$k-1>0$ 偶なら $0$ で、いずれも $\Ext=0$。よって $H^k=0$。
\item $0<k\le n$、$k$ 偶：$H_k=0$（$k<n$）または $H_n=0$（$n$ 偶）なので $\Hom=0$。$H_{k-1}=\Z_2$（$k-1$ は奇で $0<k-1<n$）なので $\Ext=\Z_2$。よって $H^k=\Z_2$。
\item $k=n$ 奇：$\Hom(H_n,\Z)=\Hom(\Z,\Z)=\Z$、$H_{n-1}=0$（$n-1$ は偶で正、ただし $n=1$ なら $H_0=\Z$）より $\Ext=0$。よって $H^n=\Z$。
\end{itemize}
まとめると
\[
H^k(\RP^n;\Z)=\begin{cases}\Z&k=0,\ \text{または }k=n\text{ が奇数}\\ \Z_2&k\text{ が偶数で }0<k\le n\\ 0&\text{それ以外}\end{cases}
\]
$\RP^2$：$H_*=(\Z,\Z_2,0)$ に対し $H^*=(\Z,0,\Z_2)$。$\RP^3$：$H_*=(\Z,\Z_2,0,\Z)$ に対し $H^*=(\Z,0,\Z_2,\Z)$。どちらも $H_1$ のねじれ $\Z_2$ が $H^2$ に現れ、ねじれの次数が1つ上がっている（命題 \ref{prop:ch06-shift}）。

直感的には、$H^1(X;\Z)=\Hom(H_1,\Z)$ は「ループに整数の巻き付きを与える」ものなので、$2\gamma\sim0$ となるループ $\gamma$ には $0$ しか与えられない。一方、$\gamma$ を2回まわった曲線が張る2次元の鎖 $c$（$\partial c=2\gamma$）には、$H^2$ のねじれ元が $\pm1\pmod 2$ を与える。
\end{solution}

\begin{problem}[難易度 B]\label{ch06-prob-lens}
(1) $H^k(\RP^n;\Z_2)\cong\Z_2$（$0\le k\le n$）を普遍係数定理で示せ。\\
(2) レンズ空間 $L=L(p;q)$（$p\ge2$）について、$\pi_1(L)=\Z_p$、$L$ が向き付け可能な3次元閉多様体であることを認める。$H_1(L)$, $H_2(L)$, $H_3(L)$ をポアンカレ双対と普遍係数定理を用いて求め、$H^*(L;\Z)$ と $H^*(L;\Z_N)$ を求めよ。\\
(3) $H^*(L;\Z_N)$ についてポアンカレ双対 $H^k\cong H^{3-k}$ が成り立っていることを確かめよ。
\end{problem}

\begin{solution}
(1) $H^k(\,\cdot\,;\Z_2)=\Hom(H_k,\Z_2)\oplus\Ext(H_{k-1},\Z_2)$。$\Hom(\Z,\Z_2)=\Hom(\Z_2,\Z_2)=\Z_2$、$\Ext(\Z_2,\Z_2)=\Z_2/2\Z_2=\Z_2$、$\Ext(\Z,\cdot)=0$。
\begin{itemize}
\item $k=0$：$\Hom(\Z,\Z_2)=\Z_2$。
\item $0<k<n$：$k$ 奇なら $\Hom(\Z_2,\Z_2)=\Z_2$、$H_{k-1}$ は $\Z$ か $0$ で $\Ext=0$。$k$ 偶なら $\Hom(0)=0$、$\Ext(\Z_2,\Z_2)=\Z_2$。いずれも $\Z_2$。
\item $k=n$：$n$ 奇なら $\Hom(\Z,\Z_2)=\Z_2$ と $\Ext(H_{n-1})=0$（$H_{n-1}$ は $0$ または $n=1$ のとき $\Z$）。$n$ 偶なら $\Hom(0)=0$ と $\Ext(\Z_2,\Z_2)=\Z_2$。
\end{itemize}
よってすべての $0\le k\le n$ で $\Z_2$。$\Z_2$ 係数では「ねじれのずれ」が $\Hom$ と $\Ext$ の両方に現れて均され、$\RP^n$ は $\Z_2$ 係数では球面の列のように各次元に1つずつ類をもつ（これは胞体 $e^0,\dots,e^n$ の個数と一致する）。

(2) $L$ は連結なので $H_0=\Z$。$H_1=\pi_1^{\mathrm{ab}}=\Z_p$。向き付け可能な閉3次元多様体なので $H_3=\Z$。ポアンカレ双対と普遍係数定理から
\[
H_2(L)\cong H^1(L;\Z)\cong\Hom(H_1,\Z)\oplus\Ext(H_0,\Z)=\Hom(\Z_p,\Z)=0.
\]
よって $H_*(L)=(\Z,\Z_p,0,\Z)$。整数係数のコホモロジーは
\[
\begin{gathered}
H^0=\Z,\qquad H^1=\Hom(\Z_p,\Z)=0,\\
H^2=\Hom(0,\Z)\oplus\Ext(\Z_p,\Z)=\Z_p,\qquad H^3=\Hom(\Z,\Z)=\Z.
\end{gathered}
\]
$d=\gcd(p,N)$ とおくと、$\Z_N$ 係数では
\[
\begin{gathered}
H^0=\Z_N,\qquad H^1=\Hom(\Z_p,\Z_N)=\Z_d,\\
H^2=\Ext(\Z_p,\Z_N)=\Z_N/p\Z_N\cong\Z_d,\qquad H^3=\Hom(\Z,\Z_N)=\Z_N.
\end{gathered}
\]
\begin{note}
$\Hom(\Z_p,\Z_N)=\{g\in\Z_N:pg=0\}$。$pg\equiv0\pmod N$ は $(p/d)g\equiv0\pmod{N/d}$ と同値で、$p/d$ と $N/d$ は互いに素なので $g\equiv0\pmod{N/d}$。解は $(N/d)\Z_N$ で、位数 $d$ の巡回群。また $\Z_N$ の中で $p$ の生成する部分群は $\gcd(p,N)=d$ の生成する部分群に等しい（ベズーの等式 $d=xp+yN$）ので $p\Z_N=d\Z_N$、$\Z_N/p\Z_N=\Z_N/d\Z_N\cong\Z_d$。
\end{note}

(3) $L$ は向き付け可能なので $\Z_N$ 係数のポアンカレ双対 $H^k(L;\Z_N)\cong H_{3-k}(L;\Z_N)$ が成り立つ。ホモロジー側を計算すると（$H_k(L;\Z_N)=H_k\otimes\Z_N\oplus\Tor(H_{k-1},\Z_N)$）、$H_0=\Z_N$, $H_1=\Z_p\otimes\Z_N=\Z_d$, $H_2=\Tor(\Z_p,\Z_N)=\Z_d$, $H_3=\Z_N$ であり、$H^k\cong H_{3-k}$ が確かに成り立つ。上の結果自体も $H^0\cong H^3\cong\Z_N$、$H^1\cong H^2\cong\Z_d$ で $k\leftrightarrow3-k$ の対称性をもつ。$N$ が $p$ と互いに素なら $d=1$ で、$L$ は $\Z_N$ 係数では $S^3$ と区別できない。
\end{solution}

\begin{problem}[難易度 B]\label{ch06-prob-zn}
物理との対応「平坦な $\Z_N$ ゲージ場」（基本事項）の記号を用いる。\\
(1) 平坦な $\Z_N$ ゲージ場のゲージ同値類全体が $H^1(X;\Z_N)\cong\Hom(\pi_1(X),\Z_N)$ であることを確かめよ。\\
(2) $T^2$, $T^3$, $\Sigma_g$, $L(p;q)$ について、平坦な $\Z_N$ ゲージ場の同値類の個数を求めよ。\\
(3) 2次元空間 $\Sigma$ 上の $\Z_N$ ゲージ理論の低エネルギー極限では、ヒルベルト空間が「平坦な $\Z_N$ ゲージ場のゲージ同値類」を基底とすることを認める。$\Sigma=T^2$ および $\Sigma_g$ での基底状態の縮退度を求めよ。これは局所的な演算子では区別できない縮退であり、\textbf{位相的秩序}の指標となる。
\end{problem}

\begin{solution}
(1) 物理との対応で見たとおり、平坦性は $\delta a=0$（$a$ は単体的1コサイクル）、ゲージ変換は $a\mapsto a+\delta\lambda$（$\lambda$ は0コチェイン）である。よって同値類の集合は単体的コホモロジー $Z^1/B^1=H^1(X;\Z_N)$ にほかならず、これは特異コホモロジーと一致する。$\Ext(H_0,\Z_N)=\Ext(\Z,\Z_N)=0$ なので $H^1(X;\Z_N)\cong\Hom(H_1(X),\Z_N)\cong\Hom(\pi_1(X),\Z_N)$。後者の同型の具体形は、ループ $\gamma$ に $\Z_N$ のウィルソン・ループ（$\gamma$ に沿う $a$ の和）$\sum_{e\in\gamma}a_e$ を対応させる写像である。$\delta a=0$ よりこれはループのホモトピー（プラケットを越える変形）で不変で、$\delta\lambda$ を加えても閉じたループ上の和は変わらない。

(2) $\Hom(\Z^m\oplus T,\Z_N)\cong\Z_N^m\oplus\Hom(T,\Z_N)$ を使う。
\begin{itemize}
\item $T^2$：$H_1=\Z^2$ より $\Z_N^2$、$N^2$ 個。
\item $T^3$：$H_1=\Z^3$ より $N^3$ 個。
\item $\Sigma_g$：$H_1=\Z^{2g}$ より $N^{2g}$ 個。
\item $L(p;q)$：$H_1=\Z_p$ より $\Z_{\gcd(p,N)}$、$\gcd(p,N)$ 個。例えば $\RP^3=L(2;1)$ 上の $\Z_2$ ゲージ場は2通り（自明なものと、生成元のループのホロノミーが $-1$ のもの）。
\end{itemize}

(3) 縮退度は $\lvert H^1(\Sigma;\Z_N)\rvert$ で、$T^2$ では $N^2$、$\Sigma_g$ では $N^{2g}$。$N=2$ の $T^2$ では $4$ 重縮退で、キタエフのトーリック符号の基底状態の縮退度と一致する。どの基底状態かはトーラスの非可縮なループに沿ったホロノミー（ウィルソン・ループ）でしか区別できず、局所的な摂動はこの縮退を（系の大きさについて指数関数的に小さい量を除いて）解かない。縮退度が空間のトポロジー（$b_1$）に依存することが位相的秩序の特徴である。
\end{solution}

\begin{problem}[難易度 D]\label{ch06-prob-torsion}
$X$ を弧状連結な閉多様体とし、$H_1(X)\cong\Z^{b_1}\oplus T$（$T$ は有限）とする。定理 \ref{thm:ch06-bock} の長完全列を用いる。\\
(1) $H^1(X;U(1))=\Hom(H_1(X),U(1))\cong U(1)^{b_1}\times T$ を示せ（$\Hom(\Z_n,U(1))\cong\Z_n$ を用いよ）。\\
(2) 完全列 $0\to H^1(X;\R)/H^1(X;\Z)\to H^1(X;U(1))\xrightarrow{\beta}\operatorname{Tors}H^2(X;\Z)\to0$ を導き、(1) と整合することを確かめよ。ここで $\operatorname{Tors}$ はねじれ部分。\\
(3) $X=L(p;q)$ とする。$X$ 上のすべての $U(1)$ 束は平坦な接続をもち、$p$ 個の同型類があることを示せ。生成元のループ $\gamma$ に沿ったホロノミー $e^{i\oint_\gamma A}$ の値を求め、電荷 $q_e\in\Z$ の粒子が $\gamma$ に沿って一周したときに得る位相を述べよ。どの電荷の粒子ならこの束を検出できないか。\\
(4) ド・ラーム・コホモロジー（第5章）ではこれらの束が区別できない理由を説明せよ。
\end{problem}

\begin{solution}
(1) $U(1)$ は可除群（任意の元の $n$ 乗根がある）なので $\Ext(\cdot,U(1))=0$ であり、普遍係数定理から $H^1(X;U(1))\cong\Hom(H_1,U(1))$。$\Hom(\Z,U(1))=U(1)$。$\Hom(\Z_n,U(1))$ は生成元の像 $z$ が $z^n=1$ をみたすものなので $\{e^{2\pi im/n}\}\cong\Z_n$。有限アーベル群 $T$ は巡回群の直和なので $\Hom(T,U(1))\cong T$。よって $\Hom(H_1,U(1))\cong U(1)^{b_1}\times T$。物理的には、$U(1)^{b_1}$ は $H_1$ の自由部分の生成元に沿った連続なウィルソン・ライン（平坦接続のモジュライ）、$T$ はねじれたループに沿った離散的なホロノミーである。

(2) 定理 \ref{thm:ch06-bock} の長完全列の一部
\[
H^1(X;\Z)\xrightarrow{i}H^1(X;\R)\to H^1(X;U(1))\xrightarrow{\beta}H^2(X;\Z)\xrightarrow{i}H^2(X;\R)
\]
を考える。完全性より $\beta$ の核は $H^1(X;\R)$ の像で、それは $H^1(X;\R)/i(H^1(X;\Z))$ と同型。$H^1(X;\Z)=\Hom(H_1,\Z)=\Z^{b_1}$ は $H^1(X;\R)=\R^{b_1}$ の格子として埋め込まれるので、これはトーラス $\R^{b_1}/\Z^{b_1}=U(1)^{b_1}$。$\beta$ の像は $H^2(X;\Z)\to H^2(X;\R)$ の核である。$H^2(X;\Z)=\Z^{b_2}\oplus T'$ に対し $H^2(X;\R)=H^2(X;\Z)\otimes\R$ への写像の核はちょうどねじれ部分 $T'$ である。よって
\[
0\to U(1)^{b_1}\to H^1(X;U(1))\xrightarrow{\beta}\operatorname{Tors}H^2(X;\Z)\to0 .
\]
命題 \ref{prop:ch06-shift} より $\operatorname{Tors}H^2(X;\Z)\cong T$（$H_1$ のねじれ）であり、(1) の $U(1)^{b_1}\times T$ と整合する（連続部分 $U(1)^{b_1}$ と離散部分 $T$）。

(3) $L=L(p;q)$ では $b_1=0$、$H^2(L;\Z)=\Z_p$（すべてねじれ）。$U(1)$ 束の同型類は $c_1\in H^2(L;\Z)=\Z_p$ で分類される（第7章）ので $p$ 個。(2) より $\beta\colon H^1(L;U(1))\to\Z_p$ は全射で、核は $U(1)^0=0$ なので同型。したがって各同型類の束は平坦な接続をちょうど1つ（ゲージ同値を除いて）もつ。

$c_1$ が $m\in\Z_p$ に対応する平坦接続のホロノミーは、生成元 $\gamma$ に沿って $e^{i\oint_\gamma A}=e^{2\pi im/p}$ である（$\beta$ の符号の規約によっては $e^{-2\pi im/p}$。いずれにせよ $m\mapsto$ ホロノミーは全単射）。電荷 $q_e$ の粒子は共変微分 $d-iq_eA$ をもつので、一周で波動関数は位相 $e^{iq_e\oint_\gamma A}=e^{2\pi iq_em/p}$ を得る。これは $F=0$ の領域を動いているのに生じる位相で、アハラノフ・ボーム効果（第2, 5章）の閉多様体版である。$q_e m\equiv0\pmod p$ なら位相は $1$ なので、特に $q_e\in p\Z$ の粒子はどの $m$ も検出できない。

(4) 平坦接続では $F=0$ なので、ド・ラーム類 $[F/2\pi]=0$。実際 $H^2_{\mathrm{dR}}(L)\cong H^2(L;\R)=\Z_p\otimes\R=0$ であり、ド・ラーム・コホモロジー（実数係数）はねじれを $\otimes\R$ で消してしまう。曲率の積分 $\int_\Sigma F$ もすべて $0$ で、流束の測定では区別できない。区別できるのはホロノミー（大域的な観測量）だけであり、整数係数のコホモロジーが必要になる。弦理論の離散的トーション、離散的な $\theta$ 角なども同じ数学に基づく。
\end{solution}

\subsection{カップ積}

\begin{problem}[難易度 A]\label{ch06-prob-tn}
$H^*(T^n;\Z)=\Lambda_\Z[x_1,\dots,x_n]$（$x_i=[dx^i]$）を認める。\\
(1) $H^k(T^n;\Z)$ の階数が $\binom nk$ であることを示し、オイラー標数 $\chi(T^n)=0$（$n\ge1$）を確かめよ。\\
(2) 一般の空間 $X$ で $a\in H^1(X;\Z)$、$H^2(X;\Z)$ にねじれがなければ $a\smile a=0$ であることを示せ。\\
(3) 場の強さが閉 $r$ 形式 $G$ で、閉 $r$ 次元サイクル上で $\int G\in2\pi\Z$（ディラック量子化）をみたすゲージ場を $T^d$ 上で考える。独立な流束量子数の個数が $b_r(T^d)$ であることを説明し、(a) マクスウェル場 $F$（$r=2$）で空間が $T^3$ の場合、(b) $r=3$ の場の強さ $G_3$ で内部空間が $T^6$ の場合に求めよ。
\end{problem}

\begin{solution}
(1) $\Lambda_\Z[x_1,\dots,x_n]$ の次数 $k$ の部分は $x_{i_1}\cdots x_{i_k}$（$i_1<\dots<i_k$）を基底とする自由加群なので、階数は $\binom nk$。
\[
\chi(T^n)=\sum_{k=0}^n(-1)^k\binom nk=(1-1)^n=0\qquad(n\ge1).
\]

(2) 次数付き可換性より $a\smile a=(-1)^{1\cdot1}a\smile a=-a\smile a$ なので $2(a\smile a)=0$。$H^2(X;\Z)$ にねじれがなければ、$2y=0$ となる $y$ は $0$ のみなので $a\smile a=0$。$T^n$ では $x_i\smile x_i=0$ がこれにあたり、微分形式の $dx^i\wedge dx^i=0$ に対応する。

(3) 閉 $r$ 形式 $G$ の積分はストークスの定理によりサイクルのホモロジー類だけで決まり、$H_r(T^d)\cong\Z^{\binom dr}$ は座標 $r$ トーラス $T^r_{i_1\cdots i_r}$ で張られる。$[G/2\pi]\in H^r(T^d;\Z)=\Hom(H_r,\Z)$ はこれらの上の整数 $\int_{T^r_{i_1\cdots i_r}}G/2\pi$ で決まり、それらは独立に選べる（$G=2\pi\sum n_{i_1\cdots i_r}dx^{i_1}\wedge\cdots\wedge dx^{i_r}$）。よって独立な流束量子数は $b_r(T^d)=\binom dr$ 個。
(a) $F$ は2形式なので $b_2(T^3)=\binom32=3$ 個：3つの座標2トーラス $T^2_{12},T^2_{23},T^2_{31}$ を貫く磁束 $2\pi n_{12}$ など。
(b) $b_3(T^6)=\binom63=20$ 個。IIB 弦理論のトーラス・コンパクト化におけるフラックス量子数の数え方の出発点である（第10章で扱う）。
\end{solution}

\begin{problem}[難易度 C]\label{ch06-prob-t2cup}
トーラス $T^2$ を、正方形 $[0,1]^2$ の対辺を平行移動で同一視して得られる Δ 複体とみる（下図）。頂点は1つ $v$、辺は下辺（$(0,0)\to(1,0)$）の $a$、左辺（$(0,0)\to(0,1)$）の $b$、対角線（$(0,0)\to(1,1)$）の $c$、2単体は $U=[(0,0),(1,0),(1,1)]$ と $L=[(0,0),(0,1),(1,1)]$ である。
\begin{center}
\begin{tikzpicture}[scale=2.3]
  \fill[blue!5] (0,0) rectangle (1,1);
  \draw[thick,red!70!black,mid arrow] (0,0) -- (1,0);
  \draw[thick,red!70!black,mid arrow] (0,1) -- (1,1);
  \draw[thick,blue!70!black,mid arrows] (0,0) -- (0,1);
  \draw[thick,blue!70!black,mid arrows] (1,0) -- (1,1);
  \draw[thick,mid arrow] (0,0) -- (1,1);
  \node[below] at (0.5,0) {$a$}; \node[above] at (0.5,1) {$a$};
  \node[left] at (0,0.5) {$b$}; \node[right] at (1,0.5) {$b$};
  \node at (0.45,0.58) {$c$};
  \node at (0.72,0.25) {$U$}; \node at (0.25,0.72) {$L$};
  \foreach \p in {(0,0),(1,0),(0,1),(1,1)} \node[vtx] at \p {};
  \node[below left] at (0,0) {\footnotesize $v$};
\end{tikzpicture}
\zucap{トーラスの Δ 複体。4頂点はすべて $v$。カップ積 $(\varphi\smile\psi)(U)=\varphi(\text{前の辺})\psi(\text{後の辺})$ は、$U$ では「$a$ の次に $b$」、$L$ では「$b$ の次に $a$」を読む。}
\end{center}
(1) 単体的コチェインで $H^1(T^2;\Z)\cong\Z^2$、$H^2(T^2;\Z)\cong\Z$ を求めよ。基本類は $[T]=[U-L]$（$\R^2$ の標準的な向き）である。\\
(2) $\alpha(a)=1,\alpha(b)=0,\alpha(c)=1$、$\beta(a)=0,\beta(b)=1,\beta(c)=1$ とするとき、$\alpha\smile\beta$, $\beta\smile\alpha$, $\alpha\smile\alpha$, $\beta\smile\beta$ を $[T]$ 上で評価し、$H^*(T^2;\Z)=\Lambda_\Z[\alpha,\beta]$ を示せ。\\
(3) 結果を微分形式 $dx\wedge dy$ と比較せよ。また、$T^2$ を空間とする $U(1)$ ゲージ理論で、2つの1形式の類 $m\alpha+n\beta$ と $m'\alpha+n'\beta$ の積の積分を求めよ。
\end{problem}

\begin{solution}
(1) $U$ の頂点を $v_0=(0,0),v_1=(1,0),v_2=(1,1)$ とすると、辺は $[v_0v_1]=a$, $[v_1v_2]=b$, $[v_0v_2]=c$。$L$ は $v_0=(0,0),v_1=(0,1),v_2=(1,1)$ で $[v_0v_1]=b$, $[v_1v_2]=a$, $[v_0v_2]=c$。よって
\[
\partial U=[v_1v_2]-[v_0v_2]+[v_0v_1]=b-c+a,\qquad\partial L=a-c+b .
\]
頂点が1つなので $\delta\colon C^0\to C^1$ は $0$。1コチェイン $\varphi$ について $(\delta\varphi)(U)=\varphi(\partial U)=\varphi(a)+\varphi(b)-\varphi(c)$、$(\delta\varphi)(L)$ も同じ。コサイクル条件は $\varphi(c)=\varphi(a)+\varphi(b)$ で、$(\varphi(a),\varphi(b))\in\Z^2$ を自由に選べるので $H^1\cong\Z^2$。2コチェイン $\psi$ は $(\psi(U),\psi(L))\in\Z^2$ で決まり、コバウンダリは $(t,t)$（$t=\varphi(a)+\varphi(b)-\varphi(c)$）の形なので、$\psi\mapsto\psi(U)-\psi(L)=\langle\psi,[U-L]\rangle$ が同型 $H^2\cong\Z$ を与える。$\partial(U-L)=0$ であり、$U$ は反時計回り、$L$ は時計回りの頂点順序をもつので $U-L$ が標準的な向きの基本類である。

(2) 定義より（$U$ では前の辺 $a$・後の辺 $b$、$L$ では前の辺 $b$・後の辺 $a$）
\begin{align*}
\langle\alpha\smile\beta,[T]\rangle&=\alpha(a)\beta(b)-\alpha(b)\beta(a)=1\cdot1-0\cdot0=1,\\
\langle\beta\smile\alpha,[T]\rangle&=\beta(a)\alpha(b)-\beta(b)\alpha(a)=0-1=-1,\\
\langle\alpha\smile\alpha,[T]\rangle&=\alpha(a)\alpha(b)-\alpha(b)\alpha(a)=0,\qquad
\langle\beta\smile\beta,[T]\rangle=\beta(a)\beta(b)-\beta(b)\beta(a)=0.
\end{align*}
$\alpha,\beta$ はコサイクル（$\alpha(c)=\alpha(a)+\alpha(b)$ など）で、$H^1$ の基底をなす。$H^2\cong\Z$ は $[T]$ 上の値で同定されるので、$\alpha\smile\beta$ は $H^2$ の生成元、$\beta\smile\alpha=-\alpha\smile\beta$、$\alpha^2=\beta^2=0$。よって $H^*(T^2;\Z)$ は $1,\alpha,\beta,\alpha\beta$ を基底とし、関係 $\alpha^2=\beta^2=0$, $\beta\alpha=-\alpha\beta$ をもつ外積代数 $\Lambda_\Z[\alpha,\beta]$ である。

(3) $\alpha$ は辺 $a$（$x$ 方向）に $1$、$b$（$y$ 方向）に $0$ を与えるので $[dx]$、$\beta=[dy]$ に対応する（$c$ に対する値 $\int_c dx=\int_c dy=1$ とも一致）。$\int_{T^2}dx\wedge dy=1$、$dy\wedge dx=-dx\wedge dy$、$dx\wedge dx=0$ は上の計算と完全に一致する。単体的カップ積は「$\wedge$ の格子版」である。
$m\alpha+n\beta$ と $m'\alpha+n'\beta$ の積は
\[
(m\alpha+n\beta)\smile(m'\alpha+n'\beta)=mn'\,\alpha\beta+nm'\,\beta\alpha=(mn'-nm')\,\alpha\beta,
\]
$[T]$ 上の値は $mn'-nm'$。例 \ref{ex:ch06-t2pd} より $m\,dx+n\,dy$ のポアンカレ双対サイクルは $(n,-m)$ サイクル（$\eta_{(p,q)}=p\,dy-q\,dx$）なので、この値は問題 \ref{ch06-prob-branes} の交叉数 $(n,-m)\cdot(n',-m')=mn'-nm'$ に一致する。1形式の類の積の積分は、双対なサイクルの交叉数である。
\end{solution}

\begin{problem}[難易度 B]\label{ch06-prob-rp2}
射影平面 $\RP^2$ を、正方形の頂点を $P=(0,0)$, $Q=(1,0)$, $R=(1,1)$, $S=(0,1)$ として、辺 $PQ$ と $RS$ を（$P\mapsto R$, $Q\mapsto S$）、辺 $PS$ と $RQ$ を（$P\mapsto R$, $S\mapsto Q$）同一視した空間とみる。辺 $a=[P,Q]=[R,S]$, $b=[P,S]=[R,Q]$, $c=[Q,S]$ と2単体 $U=[P,Q,S]$, $L=[R,Q,S]$ からなる Δ 複体（下図）を用いて、$H^*(\RP^2;\Z_2)\cong\Z_2[\alpha]/(\alpha^3)$ を示せ。また、$\Z_2$ ゲージ場の位相的作用 $\exp\bigl(i\pi\int a\smile a\bigr)$ が $\RP^2$ 上で自明でないことを説明せよ。
\begin{center}
\begin{tikzpicture}[scale=2.3]
  \fill[blue!5] (0,0) rectangle (1,1);
  \draw[thick,red!70!black,mid arrow] (0,0) -- (1,0);
  \draw[thick,red!70!black,mid arrow] (1,1) -- (0,1);
  \draw[thick,blue!70!black,mid arrows] (0,0) -- (0,1);
  \draw[thick,blue!70!black,mid arrows] (1,1) -- (1,0);
  \draw[thick,mid arrow] (1,0) -- (0,1);
  \node[below] at (0.5,0) {$a$}; \node[above] at (0.5,1) {$a$};
  \node[left] at (0,0.5) {$b$}; \node[right] at (1,0.5) {$b$};
  \node at (0.58,0.58) {$c$};
  \node at (0.28,0.28) {$U$}; \node at (0.72,0.72) {$L$};
  \node[vtx] at (0,0) {}; \node[vtx] at (1,1) {};
  \node[hvtx] at (1,0) {}; \node[hvtx] at (0,1) {};
  \node[below left] at (0,0) {\footnotesize $P$}; \node[below right] at (1,0) {\footnotesize $Q$};
  \node[above right] at (1,1) {\footnotesize $R$}; \node[above left] at (0,1) {\footnotesize $S$};
\end{tikzpicture}
\zucap{射影平面の Δ 複体。対辺を逆向きに貼る（中心に関する点対称）。黒丸 $P=R$、白丸 $Q=S$ の2頂点。}
\end{center}
\end{problem}

\begin{solution}
頂点は $v=P=R$ と $w=Q=S$ の2つ。係数は $\Z_2$ なので符号は無視してよい。$U=[P,Q,S]$ の面は $[Q,S]=c$, $[P,S]=b$, $[P,Q]=a$、$L=[R,Q,S]$ の面は $[Q,S]=c$, $[R,S]=a$, $[R,Q]=b$ なので $\partial U=\partial L=a+b+c$。また $\partial a=\partial b=v+w$, $\partial c=w+w=0$。

$H^1$：1コチェイン $\varphi$ がコサイクルである条件は $\varphi(a)+\varphi(b)+\varphi(c)=0$ で、解は2次元。0コチェイン $f$ の余境界は $\delta f(a)=\delta f(b)=f(v)+f(w)$, $\delta f(c)=0$ なので、コバウンダリは $(1,1,0)$ の張る1次元。よって $H^1(\RP^2;\Z_2)\cong\Z_2$ で、生成元として $\alpha(a)=1,\alpha(b)=0,\alpha(c)=1$ がとれる（$\alpha(c)\neq0$ なのでコバウンダリでない）。
$H^2$：2コチェインは $(\psi(U),\psi(L))$ で決まり、コバウンダリは $(t,t)$ の形なので、$\psi\mapsto\psi(U)+\psi(L)$ が同型 $H^2(\RP^2;\Z_2)\cong\Z_2$ を与える。これは $\Z_2$ 係数の基本類 $[U+L]$ 上の評価である。

カップ積：$U$ の頂点順序は $P,Q,S$、$L$ は $R,Q,S$ なので
\[
\begin{aligned}
(\alpha\smile\alpha)(U)&=\alpha([P,Q])\,\alpha([Q,S])=\alpha(a)\alpha(c)=1,\\
(\alpha\smile\alpha)(L)&=\alpha([R,Q])\,\alpha([Q,S])=\alpha(b)\alpha(c)=0 .
\end{aligned}
\]
和は $1$ なので $\alpha^2$ は $H^2$ の生成元。$\alpha^3\in H^3=0$。よって $H^*(\RP^2;\Z_2)=\Z_2\{1,\alpha,\alpha^2\}\cong\Z_2[\alpha]/(\alpha^3)$。

$\Z$ 係数なら次数付き可換性から奇数次の元の平方は位数2になるが、$\Z_2$ 係数では $-1=1$ なので制約にならず、実際 $\alpha^2\neq0$ となる。

物理：$\Z_2$ ゲージ場は $a\in H^1(X;\Z_2)$（問題 \ref{ch06-prob-zn}）で、作用 $\pi\int_Xa\smile a$ は $\langle a\smile a,[X]\rangle\in\Z_2$ に $\pi$ を掛けたものなので、$e^{i\pi\int a\smile a}=\pm1$ がゲージ不変に定まる（カップ積はコホモロジー類だけで決まる）。$X=\RP^2$, $a=\alpha$ では $\langle\alpha^2,[\RP^2]\rangle=1$ なので重みは $-1$ であり、この位相項は自明でない。一方、向き付け可能な閉曲面では $H^1$ の元の平方は $\Z_2$ 係数でも $0$（$\Sigma_g$ の環の構造、あるいはウー公式 $a\smile a=w_1\smile a$）なので、この項は向き付け不能な時空でのみ見える。こうした離散ゲージ場の位相項は、時間反転対称性で守られた相（SPT 相）の分類などに現れる。
\end{solution}

\begin{problem}[難易度 B]\label{ch06-prob-cp2}
$H^*(\CP^n;\Z)\cong\Z[x]/(x^{n+1})$（$x$ は超平面 $\CP^{n-1}$ のポアンカレ双対で $\langle x^n,[\CP^n]\rangle=1$）を認める。\\
(1) $\CP^2$ と $S^2\vee S^4$ は同型なホモロジー群をもつが、ホモトピー同値でないことを示せ（$\tilde H^*(X\vee Y)\cong\tilde H^*(X)\times\tilde H^*(Y)$ が環として成り立つことを用いてよい）。\\
(2) $\CP^2$ の交叉形式が $(1)$ であることを示し、幾何学的には「2本の複素射影直線 $\CP^1\subset\CP^2$ が1点で交わる」ことの表現であることを説明せよ。\\
(3) 写像 $f\colon\CP^n\to\CP^n$ が $f^*x=kx$ をみたせば、写像度は $k^n$ であることを示せ。$n$ が偶数なら、$\CP^n$ には向きを逆にするホモトピー同値写像がないことを結論せよ（特に $\CP^2\not\simeq\overline{\CP^2}$ を向きも込めて）。
\end{problem}

\begin{solution}
(1) どちらも $H_k=\Z$（$k=0,2,4$）、それ以外 $0$ で、コホモロジー群も同じ。$\CP^2$ では $H^2$ の生成元 $x$ について $x^2$ が $H^4$ の生成元。$S^2\vee S^4$ では $H^2$ の生成元は $y=r^*\sigma_2$（$r\colon S^2\vee S^4\to S^2$ は $S^4$ を潰す写像、$\sigma_2$ は $H^2(S^2)$ の生成元）で、$y^2=r^*(\sigma_2^2)=0$（$H^4(S^2)=0$）。$H^2$ の任意の元 $my$ の平方も $0$。ホモトピー同値写像はコホモロジー環の同型を誘導し、$H^2$ の生成元を生成元に移すので、平方が $0$ でない元の有無で区別でき、両者はホモトピー同値でない。

(2) $H^2(\CP^2;\Z)=\Z x$ で $Q(x,x)=\langle x^2,[\CP^2]\rangle=1$。よって $Q=(1)$。$x$ は直線 $\ell=\CP^1$ のポアンカレ双対なので $Q(x,x)=\ell\cdot\ell'$（$\ell'$ は $\ell$ を動かした別の直線）であり、異なる2直線は1点で交わる（射影平面の公理）。複素部分多様体どうしの交点は、複素構造による向きのもとで常に $+1$ と数えられるので $\ell\cdot\ell'=1$。$S^2\vee S^4$ の $S^2$ は自分自身から離せる（交叉数 $0$）のに対し、$\CP^2$ の $\CP^1$ は自分自身と必ず交わる、というのが (1) の違いの幾何学的な意味である。

(3) $f^*$ は環準同型なので $f^*(x^n)=(f^*x)^n=k^nx^n$。$x^n$ は $H^{2n}(\CP^n;\Z)\cong\Z$ の生成元なので、$f^*$ は最高次で $k^n$ 倍であり、写像度は $k^n$（$f_*[\CP^n]=k^n[\CP^n]$ と $\langle f^*x^n,[\CP^n]\rangle=\langle x^n,f_*[\CP^n]\rangle$ による）。ホモトピー同値写像なら $f^*$ は $H^2$ で同型なので $k=\pm1$、写像度は $(\pm1)^n$。$n$ 偶数なら写像度は $+1$ で、向きを逆にするホモトピー同値写像はない。特に $n=2$ では、交叉形式 $(1)$ と $(-1)$ は向きを保つホモトピー同値では移り合わないことと整合する。なお複素共役は $x\mapsto-x$ を誘導し、写像度は $(-1)^n$。
\end{solution}

\begin{problem}[難易度 C]\label{ch06-prob-sigmag}
$\Sigma_g$ を種数 $g$ の向き付けられた閉曲面とし、標準的なサイクル $a_i,b_i$（基本事項の図）で $a_i\cdot b_j=\delta_{ij}$, $a_i\cdot a_j=b_i\cdot b_j=0$ とする。\\
(1) $\alpha_i=-[\eta_{b_i}]$, $\beta_i=[\eta_{a_i}]$ とおくと、$\oint_{a_j}\alpha_i=\delta_{ij}$, $\oint_{b_j}\alpha_i=0$, $\oint_{a_j}\beta_i=0$, $\oint_{b_j}\beta_i=\delta_{ij}$ であり、$\int_{\Sigma_g}\alpha_i\wedge\beta_j=\delta_{ij}$, $\int\alpha_i\wedge\alpha_j=\int\beta_i\wedge\beta_j=0$ であることを示せ。\\
(2) \term{リーマンの双線形関係}{Riemann bilinear relation}：閉1形式 $\omega,\eta$ に対し
\[
\int_{\Sigma_g}\omega\wedge\eta=\sum_{i=1}^g\Bigl(\oint_{a_i}\omega\oint_{b_i}\eta-\oint_{b_i}\omega\oint_{a_i}\eta\Bigr)
\]
を示せ。\\
(3) $\Sigma_g\times\R$（$\R$ は時間）上のアーベリアン・チャーン・サイモンズ理論 $S=\frac{k}{4\pi}\int A\wedge dA$（$k\in\Z_{>0}$）では、$\Sigma_g$ 上の平坦接続のホロノミー $W(\gamma)=e^{i\oint_\gamma A}$ が演算子となり、$[\oint_{\gamma}A,\oint_{\gamma'}A]=\frac{2\pi i}{k}\,(\gamma\cdot\gamma')$ をみたすことを認める（右辺の符号は向きと作用の符号の規約による；ここではこの式を採用する）。$W(a_i)W(b_i)=e^{-2\pi i/k}W(b_i)W(a_i)$ を示せ。さらに、大ゲージ変換の存在により $W(a_i)^k$, $W(b_i)^k$ がヒルベルト空間上スカラーとして作用することを認めて、基底状態（ヒルベルト空間）の次元が $k^g$ であることを説明せよ。
\end{problem}

\begin{solution}
(1) 定義 \ref{def:ch06-eta} の規約 $\oint_C\omega=\int\omega\wedge\eta_C$ と $\int\eta_A\wedge\eta_B=A\cdot B$ を使う。
\begin{align*}
\oint_{a_j}\alpha_i&=-\int\eta_{b_i}\wedge\eta_{a_j}=-(b_i\cdot a_j)=a_j\cdot b_i=\delta_{ij},&
\oint_{b_j}\alpha_i&=-\int\eta_{b_i}\wedge\eta_{b_j}=-(b_i\cdot b_j)=0,\\
\oint_{a_j}\beta_i&=\int\eta_{a_i}\wedge\eta_{a_j}=a_i\cdot a_j=0,&
\oint_{b_j}\beta_i&=\int\eta_{a_i}\wedge\eta_{b_j}=a_i\cdot b_j=\delta_{ij}.
\end{align*}
ここで1次元サイクルどうしでは $A\cdot B=-B\cdot A$ を用いた。積分は
\[
\int\alpha_i\wedge\beta_j=-\int\eta_{b_i}\wedge\eta_{a_j}=-(b_i\cdot a_j)=\delta_{ij},\quad
\int\alpha_i\wedge\alpha_j=b_i\cdot b_j=0,\quad\int\beta_i\wedge\beta_j=a_i\cdot a_j=0 .
\]
（$T^2$ では $\alpha=dx$, $\beta=dy$ で、例 \ref{ex:ch06-t2pd} と一致する。）

(2) $A_i=\oint_{a_i}\omega$, $B_i=\oint_{b_i}\omega$ とおく。閉形式 $\omega-\sum_i(A_i\alpha_i+B_i\beta_i)$ はすべての $a_j,b_j$ 上の周期が $0$ で、$H_1(\Sigma_g)$ は $a_j,b_j$ で生成されるので、ド・ラームの定理によりこれは完全形式である。よって $[\omega]=\sum(A_i\alpha_i+B_i\beta_i)$。同様に $[\eta]=\sum(A'_i\alpha_i+B'_i\beta_i)$。$\int\omega\wedge\eta$ はストークスの定理によりコホモロジー類だけで決まる（$\int d f\wedge\eta=\int d(f\eta)=0$）ので、(1) より
\[
\int\omega\wedge\eta=\sum_{i,j}\bigl(A_iB'_j\textstyle\int\alpha_i\wedge\beta_j+B_iA'_j\int\beta_i\wedge\alpha_j\bigr)=\sum_i(A_iB'_i-B_iA'_i).
\]

(3) $X=\oint_{a_i}A$, $Y=\oint_{b_i}A$ とおくと $[X,Y]=\frac{2\pi i}{k}(a_i\cdot b_i)=\frac{2\pi i}{k}$ は $c$ 数。$[X,Y]=c$ が中心的ならベイカー・キャンベル・ハウスドルフの公式により $e^{iX}e^{iY}=e^{iY}e^{iX}e^{[iX,iY]}=e^{iY}e^{iX}e^{-c}$。よって
\[
W(a_i)W(b_i)=e^{-2\pi i/k}\,W(b_i)W(a_i).
\]
$i\neq j$ の組や $a_i$ どうしは交叉数 $0$ なので可換。$U=W(a_i)$, $V=W(b_i)$ はユニタリで $UV=\omega VU$（$\omega=e^{-2\pi i/k}$ は1の原始 $k$ 乗根）、$U^k=u$, $V^k=v$ はスカラーである（クロック・シフト代数）。$U$ の固有ベクトル $\lvert0\rangle$（$U\lvert0\rangle=\lambda\lvert0\rangle$）をとると、$UV^m\lvert0\rangle=\omega^m\lambda V^m\lvert0\rangle$ なので $V^m\lvert0\rangle$（$m=0,\dots,k-1$）は相異なる固有値をもち線形独立で、$V^k\lvert0\rangle=v\lvert0\rangle$ によりこの $k$ 次元空間は $U,V$ で閉じる。逆に任意の表現は、$U$ の固有値が $\lambda\omega^m$ の形に限られ（$\lambda^k=u$）各固有空間が $V$ で巡回的に移り合うので、この $k$ 次元表現の直和である。ヒルベルト空間はこの代数の既約表現（ハンドルあたり $k$ 次元）で与えられ、各ハンドル $(a_i,b_i)$ の代数は互いに可換なので全体は $k^g$ 次元。これが $U(1)_k$ チャーン・サイモンズ理論（分数量子ホール効果の有効理論）の $\Sigma_g$ 上の基底状態縮退度である。リーマンの双線形関係にあらわれるシンプレクティック構造（交叉形式）が、正準交換関係そのものになっている。
\end{solution}

\subsection{ポアンカレ双対と交叉数}

\begin{problem}[難易度 A]\label{ch06-prob-odd}
(1) $M$ を向き付け可能な $n$ 次元連結閉多様体とする。ポアンカレ双対と普遍係数定理から $b_k=b_{n-k}$ を示せ。\\
(2) $n$ が奇数なら $\chi(M)=0$ であることを示せ。向き付け可能でない場合にも、$\Z_2$ 係数を用いて同じ結論が得られることを説明せよ（$\chi(M)=\sum_k(-1)^k\dim_{\Z_2}H_k(M;\Z_2)$ を用いてよい）。\\
(3) $L(p;q)$ と $\RP^3$ で確かめよ。
\end{problem}

\begin{solution}
(1) ポアンカレ双対 $H^k(M;\Q)\cong H_{n-k}(M;\Q)$、および体係数の普遍係数定理 $H^k(M;\Q)\cong\Hom(H_k(M),\Q)$ から $\dim H^k(M;\Q)=b_k$。また $H_{n-k}(M;\Q)=H_{n-k}(M)\otimes\Q$ の次元は $b_{n-k}$。よって $b_k=b_{n-k}$。

(2) $n=2m+1$ とすると、$\chi=\sum_{k=0}^n(-1)^kb_k$ で $k$ と $n-k$ の項は $b_k=b_{n-k}$ かつ $(-1)^{n-k}=-(-1)^k$ なので打ち消し合う（$k=n-k$ となる $k$ はない）。よって $\chi=0$。向き付け不能でも $\Z_2$ 係数のポアンカレ双対 $H^k(M;\Z_2)\cong H_{n-k}(M;\Z_2)$ と、体 $\Z_2$ 上の $\dim H^k(M;\Z_2)=\dim H_k(M;\Z_2)$ から、$\Z_2$ ベッティ数が同じ対称性をもち、同様に $\chi=0$。

(3) $L(p;q)$：$b_*=(1,0,0,1)$ なので $\chi=1-0+0-1=0$。$\RP^3$：$\Z_2$ 係数で $\dim H_k=1$（$k=0,1,2,3$）なので $\chi=1-1+1-1=0$。整数係数のベッティ数 $(1,0,0,1)$ でも同じ値になる（$\chi$ は係数の体によらない）。物理的には、3次元の閉じた空間では（ポアンカレ・ホップの定理により）至るところ零にならないベクトル場が存在し、例えば一様な「向き」をもつ秩序変数の配位に位相的な障害はない。
\end{solution}

\begin{problem}[難易度 B]\label{ch06-prob-branes}
$T^2=\R^2/\Z^2$（座標 $(x,y)$、向き $dx\wedge dy$）上で、互いに素な整数 $(p,q)$ に対し直線 $t\mapsto(pt,qt)+\bm v$（$t\in[0,1]$）の像を $(p,q)$ サイクルと呼ぶ。ホモロジー類は $p[a]+q[b]$ である（$a$, $b$ は $x$, $y$ 方向の基本サイクル）。\\
(1) $(p,q)$ サイクルのポアンカレ双対形式は $\eta_{(p,q)}=p\,dy-q\,dx$ であることを示せ。\\
(2) 交叉数 $(p,q)\cdot(r,s)=ps-qr$ を $\int\eta\wedge\eta$ で計算せよ。\\
(3) 2本の直線の交点を直接数え、交点の数が $\lvert ps-qr\rvert$ で、すべて同じ符号をもつことを示せ（下図は $(1,0)$ と $(1,2)$）。\\
(4) 交叉するDブレーン模型では、$T^6=T^2\times T^2\times T^2$ の各 $T^2$ でブレーン $\mathcal A$ が $(n^i_{\mathcal A},m^i_{\mathcal A})$ サイクルに巻き付く（3サイクル）。2つのブレーン $\mathcal A,\mathcal B$ の交点に局在するカイラル・フェルミオンの数（世代数）が交叉数の絶対値で与えられることを認めて、それを $(n,m)$ で表せ。
\begin{center}
\begin{tikzpicture}[scale=2.6]
  \fill[blue!5] (0,0) rectangle (1,1);
  \draw[gray] (0,0) rectangle (1,1);
  \draw[very thick,blue!70!black,mid arrow] (0,0.25) -- (1,0.25);
  \draw[very thick,red!70!black,mid arrow] (0,0) -- (0.5,1);
  \draw[very thick,red!70!black,mid arrow] (0.5,0) -- (1,1);
  \node[vtx] at (0.125,0.25) {}; \node[vtx] at (0.625,0.25) {};
  \node[blue!70!black,right] at (1.02,0.25) {\footnotesize $(1,0)$};
  \node[red!70!black,right] at (1.02,0.9) {\footnotesize $(1,2)$};
  \node[below left] at (0,0) {\footnotesize $(0,0)$};
  \node[below right] at (1,0) {\footnotesize $(1,0)$};
  \node[above left] at (0,1) {\footnotesize $(0,1)$};
\end{tikzpicture}
\zucap{$T^2$ 上の $(1,0)$ サイクル（青）と $(1,2)$ サイクル（赤）。正方形の対辺は同一視されている。交点（黒丸）は $\lvert1\cdot2-0\cdot1\rvert=2$ 個。}
\end{center}
\end{problem}

\begin{solution}
(1) $C=(p,q)$ サイクルとする。閉1形式のコホモロジー類は $[dx],[dy]$ で張られるので、これらで確かめれば十分。$\int_C dx=\int_0^1p\,dt=p$, $\int_C dy=q$。一方 $\eta=p\,dy-q\,dx$ について
\[
\int_{T^2}dx\wedge\eta=p\int dx\wedge dy=p,\qquad\int_{T^2}dy\wedge\eta=-q\int dy\wedge dx=q .
\]
よって $\int_C\omega=\int\omega\wedge\eta$ が成り立ち、$\eta_{(p,q)}=p\,dy-q\,dx$（$p=1,q=0$ で $\eta_a=dy$、$p=0,q=1$ で $\eta_b=-dx$、例 \ref{ex:ch06-t2pd} と一致）。

(2)
\[
(p,q)\cdot(r,s)=\int(p\,dy-q\,dx)\wedge(r\,dy-s\,dx)=\int\bigl(-ps\,dy\wedge dx-qr\,dx\wedge dy\bigr)=ps-qr .
\]

(3) 交点は $(pt,qt)\equiv(ru,su)+\bm w\pmod{\Z^2}$（$\bm w$ は2本の直線の平行移動の差）となる $(t,u)\in[0,1)^2$ である。$M=\begin{pmatrix}p&-r\\q&-s\end{pmatrix}$ とおくと条件は $M(t,u)^{\mathsf T}\equiv\bm w$ で、$M$ は $T^2\to T^2$ の写像を誘導し、$\det M=-(ps-qr)\neq0$ のときその写像は $\lvert\det M\rvert$ 重の被覆である（$\Z^2/M\Z^2$ の位数が $\lvert\det M\rvert$）。よって交点はちょうど $\lvert ps-qr\rvert$ 個（$\bm w$ が一般の位置なら横断的）。各交点での符号は、接ベクトルの組 $((p,q),(r,s))$ の向き、すなわち $\operatorname{sign}\det\begin{pmatrix}p&r\\q&s\end{pmatrix}=\operatorname{sign}(ps-qr)$ で、直線なのですべての交点で同じ。よって交叉数 $=\pm$（交点数）$=ps-qr$ で (2) と一致する。図の例では $(1,0)\cdot(1,2)=2$。

(4) 各 $T^2_i$ での交叉数は $I^i_{\mathcal{AB}}=n^i_{\mathcal A}m^i_{\mathcal B}-m^i_{\mathcal A}n^i_{\mathcal B}$ で、交点は各因子の交点の組なので、交点の総数は $\prod_i\lvert I^i_{\mathcal{AB}}\rvert$、すべて同じ符号をもつ。よって世代数は
\[
\lvert I_{\mathcal{AB}}\rvert=\prod_{i=1}^3\bigl\lvert n^i_{\mathcal A}m^i_{\mathcal B}-m^i_{\mathcal A}n^i_{\mathcal B}\bigr\rvert .
\]
\begin{note}
符号つきの交叉数は、積の向き $T^2_1\times T^2_2\times T^2_3$ のもとで $(a_1\times a_2\times a_3)\cdot(b_1\times b_2\times b_3)$ を計算すると、接ベクトルの並び $(a_1,a_2,a_3,b_1,b_2,b_3)$ を $(a_1,b_1,a_2,b_2,a_3,b_3)$ に並べ替える置換（転倒数3）の符号 $-1$ が掛かり、$I_{\mathcal{AB}}=-\prod_iI^i_{\mathcal{AB}}$ となる。文献ではカイラリティの規約にこの符号を吸収して $I_{\mathcal{AB}}=\prod_iI^i_{\mathcal{AB}}$ と書くことが多い。符号はカイラリティ（左巻きか右巻きか）を、絶対値は世代数を与える。標準模型の3世代を得るには、例えば $\prod_i\lvert I^i\rvert=3$ となる巻き付き数を選ぶ（第10章で扱う）。
\end{note}
\end{solution}

\begin{problem}[難易度 A]\label{ch06-prob-kunneth}
(1) キュネットの公式を用いて $H^*(S^2\times S^4;\Z)$ を環として求めよ。\\
(2) $S^2\times S^4$ と $\CP^3$ は同じベッティ数 $(1,0,1,0,1,0,1)$ をもつが、ホモトピー同値でないことを示せ。\\
(3) 6次元多様体 $Y$ の $H^2$ の基底 $\omega_i$ に対する\textbf{3重交叉数} $d_{ijk}=\int_Y\omega_i\wedge\omega_j\wedge\omega_k$ は、カラビ・ヤウ・コンパクト化の低エネルギー有効理論の結合定数に現れる（第10章で扱う）。$S^2\times S^4$ と $\CP^3$ の $H^2$ の生成元について $d_{111}$ を求めよ。
\end{problem}

\begin{solution}
(1) $H^*(S^4;\Z)$ は自由なので $\Z$ 係数のキュネットの公式が使える。$x=p_1^*\sigma_2$, $y=p_2^*\sigma_4$ とおくと
\[
H^*(S^2\times S^4;\Z)=\Z\{1,x,y,xy\},\quad\lvert x\rvert=2,\ \lvert y\rvert=4,\quad x^2=0,\ y^2=0,\ xy=yx .
\]
（$x^2=p_1^*(\sigma_2^2)=0$ など。$xy$ は $H^6$ の生成元。）

(2) $\CP^3$ では $H^2$ の生成元 $u$ について $u^2$ が $H^4$ の生成元、$u^3$ が $H^6$ の生成元。$S^2\times S^4$ では $H^2$ の任意の元 $mx$ の平方は $0$。ホモトピー同値写像はコホモロジー環の同型を誘導し $H^2$ の生成元を生成元に移すので、両者はホモトピー同値でない（群は同じでも環が違う）。

(3) $S^2\times S^4$：$d_{111}=\langle x^3,[S^2\times S^4]\rangle=0$。$\CP^3$：$d_{111}=\langle u^3,[\CP^3]\rangle=1$（3枚の超平面は1点で交わる）。3重交叉数は $H^2$ の元の「3つの積」に関する情報であり、ベッティ数では見えないトポロジーを有効理論の結合定数として反映する。
\end{solution}

\subsection{4次元多様体の交叉形式}

\begin{problem}[難易度 B]\label{ch06-prob-s2s2}
(1) $S^2\times S^2$ の交叉形式が $H=\begin{pmatrix}0&1\\1&0\end{pmatrix}$ であることを示し、$b_\pm$, $\sigma$, 偶奇を求めよ。\\
(2) $\CP^2\#\overline{\CP^2}$ の交叉形式は $(1)\oplus(-1)$ である。$b_\pm$, $\sigma$, 偶奇を求め、$S^2\times S^2$ とホモトピー同値でないことを示せ。\\
(3) $S^2\times S^2$ 上の $U(1)$ 束で、2つの因子 $S^2\times\{\mathrm{pt}\}$, $\{\mathrm{pt}\}\times S^2$ を貫く流束が $2\pi a$, $2\pi b$ のものについて、$\frac{1}{4\pi^2}\int F\wedge F$ を求めよ。
\end{problem}

\begin{solution}
(1) キュネットの公式で $H^2(S^2\times S^2)=\Z x\oplus\Z y$（$x=p_1^*\sigma$, $y=p_2^*\sigma$）、$x^2=y^2=0$, $xy=yx$ は $H^4$ の生成元で $\langle xy,[S^2\times S^2]\rangle=1$（積の向き）。よって $Q=\begin{pmatrix}0&1\\1&0\end{pmatrix}=H$。固有値は $\pm1$ なので $b_+=b_-=1$, $\sigma=0$。$Q(mx+ny,mx+ny)=2mn$ は常に偶数なので偶。幾何学的には、$A=S^2\times\{\mathrm{pt}\}$ と $B=\{\mathrm{pt}\}\times S^2$ は1点で交わり（$A\cdot B=1$）、$A$ は別の点の $S^2\times\{\mathrm{pt}'\}$ にずらせば自分と交わらない（$A\cdot A=0$）。

(2) $Q=\operatorname{diag}(1,-1)$ なので $b_+=b_-=1$, $\sigma=0$。$Q(e_1,e_1)=1$ は奇数なので奇。交叉形式はカップ積と基本類で決まるのでホモトピー不変（向きを込めて；向きを逆にしても $-Q$ の偶奇は変わらない）。$H$ は偶、$(1)\oplus(-1)$ は奇なので、整数上同型でなく、両者はホモトピー同値でない。実数上ではどちらも符号数 $(1,1)$ の形で同値であり、整数性（格子の構造）が本質的である。

(3) $c=[F/2\pi]=ax'+by'$ と書ける。ここで $x',y'$ は $\int_{A}x'=1$, $\int_Bx'=0$ などをみたす双対基底で、$S^2\times S^2$ では $x'=x$, $y'=y$（$\int_Ax=\langle\sigma,[S^2]\rangle=1$, $\int_Bx=0$）。よって
\[
\frac{1}{4\pi^2}\int F\wedge F=Q(c,c)=Q(ax+by,ax+by)=2ab .
\]
$U(1)$ インスタントン数 $k=\frac{1}{8\pi^2}\int F\wedge F=ab$ は常に整数である（$S^2\times S^2$ はスピン多様体で偶な形式をもつ）。$a=b=1$ で $k=1$。
\end{solution}

\begin{problem}[難易度 B]\label{ch06-prob-t4}
$T^4=\R^4/\Z^4$（座標 $x^1,\dots,x^4$、向き $dx^1\wedge dx^2\wedge dx^3\wedge dx^4$）とする。\\
(1) $H^2(T^4;\Z)$ の基底 $e_{ij}=[dx^i\wedge dx^j]$（$i<j$）について交叉形式を計算し、$3H$ と同型であることを示せ。\\
(2) $U(1)$ 流束 $F=2\pi\sum_{i<j}n_{ij}\,dx^i\wedge dx^j$（$n_{ij}\in\Z$）について
\[
\frac{1}{4\pi^2}\int_{T^4}F\wedge F=2\,(n_{12}n_{34}-n_{13}n_{24}+n_{14}n_{23})
\]
を示せ。右辺の括弧は反対称行列 $(n_{ij})$ のパフィアン $\Pf(n)$ である。\\
(3) $T^4$ 上の $U(1)$ インスタントン数は常に整数であることを結論し、これが「$T^4$ はスピン多様体（実際に平行化可能）」であることと整合することを述べよ。
\end{problem}

\begin{solution}
(1) $\mathrm{vol}=dx^1\wedge dx^2\wedge dx^3\wedge dx^4$ とする。$e_{ij}\wedge e_{kl}$ は $\{i,j\}\cap\{k,l\}\neq\emptyset$ なら $0$。補集合の組について、置換の符号から
\[
e_{12}\wedge e_{34}=\mathrm{vol},\qquad e_{13}\wedge e_{24}=dx^1dx^3dx^2dx^4=-\mathrm{vol},\qquad e_{14}\wedge e_{23}=dx^1dx^4dx^2dx^3=+\mathrm{vol}
\]
（$1324$ は転倒数1、$1423$ は転倒数2）。2形式どうしは可換なので $e_{34}\wedge e_{12}=\mathrm{vol}$ など。基底を $(e_{12},e_{34}),(e_{13},e_{24}),(e_{14},e_{23})$ の順に並べると
\[
Q=H\oplus(-H)\oplus H .
\]
$-H$ は基底 $(e_{13},-e_{24})$ に取り替えれば $H$ になるので $Q\cong3H$。$b_2=6$, $b_+=b_-=3$, $\sigma=0$、偶。

(2) $F/2\pi=\sum_{i<j}n_{ij}e_{ij}$ として
\[
\Bigl(\frac{F}{2\pi}\Bigr)^2=2\bigl(n_{12}n_{34}\,e_{12}e_{34}+n_{13}n_{24}\,e_{13}e_{24}+n_{14}n_{23}\,e_{14}e_{23}\bigr)=2(n_{12}n_{34}-n_{13}n_{24}+n_{14}n_{23})\,\mathrm{vol}.
\]
係数 $2$ は、例えば $e_{12}e_{34}$ と $e_{34}e_{12}$ の2項が現れることによる。$\int\mathrm{vol}=1$ より主張が従う。また $\int_{T^2_{ij}}F=2\pi n_{ij}$ でディラック量子化をみたしている。

(3) $k=\frac{1}{8\pi^2}\int F\wedge F=\Pf(n)\in\Z$。例えば $n_{12}=n_{34}=1$（2つの直交する2トーラスを貫く単位磁束）で $k=1$。交叉形式 $3H$ は偶であり、ウー公式（定理 \ref{thm:ch06-wu}）によりスピン多様体の交叉形式は偶であることと整合する。$T^4$ は接束が自明（平行化可能）なので $w_2=0$ でスピンである。こうした「トーラス上の流束」は、格子ゲージ理論の 't Hooft のねじれ境界条件や、弦理論のトーラス・コンパクト化における磁化ブレーンに現れる。
\end{solution}

\begin{problem}[難易度 B]\label{ch06-prob-e8}
$E_8$ のディンキン図形（下図）に対し、カルタン行列 $C_{ij}$ を $C_{ii}=2$、$i,j$ が辺で結ばれていれば $C_{ij}=-1$、それ以外 $0$ で定める。
\begin{center}
\begin{tikzpicture}[scale=1.0]
  \foreach \i in {1,...,7} { \node[circle,draw,thick,fill=white,inner sep=2pt] (n\i) at (\i,0) {}; \node[below=2pt] at (\i,0) {\footnotesize $\i$}; }
  \node[circle,draw,thick,fill=white,inner sep=2pt] (n8) at (5,1) {};
  \node[right=2pt] at (5,1) {\footnotesize $8$};
  \foreach \i/\j in {1/2,2/3,3/4,4/5,5/6,6/7} \draw[thick] (n\i) -- (n\j);
  \draw[thick] (n5) -- (n8);
\end{tikzpicture}
\zucap{$E_8$ のディンキン図形。頂点 $5$ から長さ $4,2,1$ の腕（$\{1,2,3,4\}$, $\{6,7\}$, $\{8\}$）が出ている。}
\end{center}
(1) 木の形の図形で、端点 $v$ とその隣 $w$ について $\det C=2\det C_{\setminus v}-\det C_{\setminus\{v,w\}}$ を示せ（$C_{\setminus v}$ は $v$ の行と列を除いたもの）。\\
(2) $\det A_n=n+1$, $\det D_n=4$（$n\ge3$；$D_3=A_3$）を認めて、$\det E_6=3$, $\det E_7=2$, $\det E_8=1$ を示せ。\\
(3) $E_8$ 形式は正定値・偶・ユニモジュラーであることを示し、符号数を答えよ。\\
(4) $E_8$ 形式を交叉形式とする滑らかな閉4次元多様体は存在しないことを、ロホリンの定理から説明せよ（単連結な場合、偶な交叉形式ならスピンであることを用いてよい）。
\end{problem}

\begin{solution}
(1) $v$ の行で余因子展開する。$v$ の行の非零成分は $C_{vv}=2$ と $C_{vw}=-1$ だけ。$C_{vv}$ の余因子は $\det C_{\setminus v}$。$C_{vw}$ の余因子は、$v$ 行と $w$ 列を除いた小行列で、その中で $v$ 列の非零成分は $C_{wv}=-1$（$w$ 行）だけなので、さらに $v$ 列で展開して $\pm C_{wv}\det C_{\setminus\{v,w\}}$ になる。符号を整理すると、一般に（対称行列で $v$ が葉のとき）
\[
\det C=C_{vv}\det C_{\setminus v}-C_{vw}C_{wv}\det C_{\setminus\{v,w\}}=2\det C_{\setminus v}-\det C_{\setminus\{v,w\}} .
\]
\begin{note}
符号の確認：$\det C=\sum_\pi\operatorname{sgn}\pi\prod_iC_{i\pi(i)}$ で、$\pi(v)=v$ の項の和が $2\det C_{\setminus v}$。$\pi(v)=w$ の項では、$v$ 列の非零成分は $v$ 行と $w$ 行だけで $v$ 行はすでに使ったので $\pi(w)=v$。よって $\pi$ は互換 $(v\,w)$ と $C_{\setminus\{v,w\}}$ 上の置換の積で、符号は $-1$ 倍、寄与は $-C_{vw}C_{wv}\det C_{\setminus\{v,w\}}=-\det C_{\setminus\{v,w\}}$。
\end{note}

(2) $E_8$ で長い腕の端 $v=1$, $w=2$ をとる。$E_8\setminus\{1\}$ は腕の長さ $3,2,1$ で $E_7$、$E_8\setminus\{1,2\}$ は腕 $2,2,1$ で $E_6$。同様に $E_6\setminus\{3\}$ は腕 $1,2,1$ で $D_5$、$E_6\setminus\{3,4\}$ は頂点 $5,6,7,8$ の鎖で $A_4$。$E_7\setminus\{2\}=E_6$、$E_7\setminus\{2,3\}=D_5$。したがって
\begin{align*}
\det E_6&=2\det D_5-\det A_4=2\cdot4-5=3,\\
\det E_7&=2\det E_6-\det D_5=2\cdot3-4=2,\\
\det E_8&=2\det E_7-\det E_6=2\cdot2-3=1 .
\end{align*}
（$\det D_5=4$ も同じ漸化式 $\det D_5=2\det D_4-\det D_3=8-4$ で確かめられる。）

(3) 単純リー環の単純ルート $\alpha_i$（長さの2乗 $2$ に規格化）の内積行列がカルタン行列 $C_{ij}=\alpha_i\cdot\alpha_j$ であり（単純レース型）、線形独立なベクトルのグラム行列なので正定値。$\det=1$ でユニモジュラー。$x\in\Z^8$ に対し
\[
x^{\mathsf T}Cx=\sum_i2x_i^2-2\sum_{\langle ij\rangle}x_ix_j\in2\Z
\]
（$\langle ij\rangle$ は辺）なので偶。正定値なので $b_+=8$, $b_-=0$, $\sigma=8$。これは「偶ユニモジュラーなら $\sigma\equiv0\pmod8$」の最小の例である。

(4) $X$ が交叉形式 $E_8$ の滑らかな閉4次元多様体なら、$\sigma=8$。単連結と仮定すると、偶な交叉形式なのでウー公式から $w_2=0$、すなわちスピン。ロホリンの定理は $\sigma\equiv0\pmod{16}$ を要求するが $8\not\equiv0$ で矛盾。（単連結でない場合も、$\pi_1$ の仮定を外した形のドナルドソンの定理「滑らかな閉4次元多様体の定値な交叉形式は対角化可能」により否定される。$E_8$ 形式は偶なので $\operatorname{diag}(1,\dots,1)$ と同型でない。）なお位相多様体としては $E_8$ 多様体が存在し（フリードマン）、「滑らかさ」が本質的である。弦理論では、16次元の偶ユニモジュラー正定値格子 $E_8\oplus E_8$ と $\mathrm{Spin}(32)/\Z_2$ の格子が、ヘテロティック弦の2つのゲージ群を与える（モジュラー不変性が偶ユニモジュラー性を要求する）。
\end{solution}

\begin{problem}[難易度 C]\label{ch06-prob-k3}
K3 曲面 $X$ は単連結なコンパクト複素曲面で、$\chi(X)=24$, $\sigma(X)=-16$、スピン多様体であることを認める。\\
(1) $b_0,\dots,b_4$ と $b_\pm$ を求めよ。\\
(2) セールの分類を用いて、交叉形式が $3H\oplus2(-E_8)$ であることを示せ。行列式と符号数を確かめよ。\\
(3) ロホリンの定理との整合性を確かめよ。また、ヒルツェブルフの符号数定理 $\sigma=\frac13\int p_1$ と $\Ahat$ 種数 $\int\Ahat=-\frac{1}{24}\int p_1$ から、K3 上のディラック作用素の指数を求めよ（第9章で扱う公式を認めてよい）。\\
(4) IIA 型弦理論を K3 にコンパクト化すると、6次元の $U(1)$ ゲージ場の個数は $b_0+b_2+b_4$ で与えられる（RR 場 $C_1$ と $C_3$ を偶数次のサイクルで展開することによる）。この個数を求めよ。
\end{problem}

\begin{solution}
(1) 単連結なので $H_1=0$、$b_1=0$。ポアンカレ双対で $b_3=b_1=0$。連結かつ向き付け可能（複素多様体）なので $b_0=b_4=1$。$\chi=b_0-b_1+b_2-b_3+b_4=2+b_2=24$ より $b_2=22$。$b_++b_-=22$, $b_+-b_-=-16$ を解いて $b_+=3$, $b_-=19$。

(2) 単連結なので $H^2(X;\Z)=\Hom(H_2,\Z)\oplus\Ext(H_1,\Z)$ は自由で階数22。スピンなので交叉形式は偶（ウー公式）。$b_+=3>0$, $b_-=19>0$ なので不定値。セールの分類から $Q\cong aH\oplus bE_8$ で、階数 $2a+8\lvert b\rvert=22$、符号数 $0\cdot a+8b=-16$。よって $b=-2$、$a=3$ で $Q\cong3H\oplus2(-E_8)$。
確認：$\det(3H\oplus2(-E_8))=(\det H)^3(\det(-E_8))^2=(-1)^3\cdot((-1)^8\cdot1)^2=-1$ でユニモジュラー。符号数は $H$ が $0$、$-E_8$ が $-8$ ずつで $-16$。

(3) $\sigma=-16\equiv0\pmod{16}$ でロホリンの定理と整合する。$\int p_1=3\sigma=-48$、$\int\Ahat=-\frac{-48}{24}=2$。指数定理（第9章）より、K3 上のディラック作用素の指数は $\ind\slashed D=\int\Ahat=2$。（実際 K3 はハイパーケーラーで、正カイラリティの平行なスピノルが2つある。）ロホリンの定理は、4次元スピン多様体では $\ind\slashed D=-\sigma/8$ が偶数（四元数構造によりゼロモードが対で現れる）であることから従う：$-\sigma/8\in2\Z\iff\sigma\in16\Z$。K3 は $-\sigma/8=2$ でこの最小例である。

(4) $b_0+b_2+b_4=1+22+1=24$ 個。$C_1$ から1個（$b_0$）、$C_3$ を22個の調和2形式で展開して22個（$b_2$）、$C_3$ の6次元での双対（$b_4$）から1個。電荷格子は偶数次のコホモロジー $H^0\oplus H^2\oplus H^4$ に（交叉形式で）内積を入れた格子 $(3H\oplus2(-E_8))\oplus H\cong4H\oplus2(-E_8)$ で、符号数 $(4,20)$ の偶ユニモジュラー格子 $\Gamma^{4,20}$ となり、$T^4$ 上のヘテロティック弦の格子と一致する（IIA/K3 とヘテロティック/$T^4$ の双対性の証拠の1つ；第10章で扱う）。
\end{solution}

\begin{problem}[難易度 C]\label{ch06-prob-theta}
$X$ を向き付けられた4次元閉多様体（ユークリッド時空）とし、$U(1)$ ゲージ理論に $\theta$ 項 $S_\theta=\frac{\theta}{8\pi^2}\int_XF\wedge F$ を加える。経路積分の重みは $e^{iS_\theta}$ を含み、すべての $U(1)$ 束（$c=[F/2\pi]\in H^2(X;\Z)$）について和をとるとする。\\
(1) $S_\theta=\frac\theta2Q_X(c,c)$ を示せ。\\
(2) $X$ がスピン多様体（交叉形式が偶）なら $\theta\sim\theta+2\pi$ であることを示せ。\\
(3) $X=\CP^2$ では $S_\theta=\frac\theta2n^2$（$n\in\Z$）となること、$\theta\to\theta+2\pi$ で重みが $(-1)^n$ 倍され、周期が $4\pi$ になることを示せ。\\
(4) ウー公式を用いると、一般の $X$ で $\theta\to\theta+2\pi$ による重みの変化は $(-1)^{\langle\bar c\smile w_2,[X]\rangle}$ と書けることを示し、物理的な意味を述べよ。
\end{problem}

\begin{solution}
(1) 物理との対応「$U(1)$ 流束とインスタントン数」より、$F/2\pi$ はド・ラーム類として $c$ の実像を代表し、ド・ラームの同型は積を保つので
\[
\frac{1}{4\pi^2}\int_XF\wedge F=\int_X\frac F{2\pi}\wedge\frac F{2\pi}=\langle c\smile c,[X]\rangle=Q_X(c,c).
\]
よって $S_\theta=\frac{\theta}{8\pi^2}\int F\wedge F=\frac\theta2Q_X(c,c)$。$U(1)$ のインスタントン数は $k=\frac12Q_X(c,c)$。

(2) 偶なら $Q_X(c,c)=2m$（$m\in\Z$）、$S_\theta=\theta m$。$\theta\to\theta+2\pi$ で $e^{iS_\theta}\to e^{iS_\theta}e^{2\pi im}=e^{iS_\theta}$。すべての束の重みが不変なので理論は同じであり、$\theta$ の周期は $2\pi$。$S^2\times S^2$（$Q(c,c)=2ab$）、$T^4$（$2\Pf(n)$）、K3 がこの場合である。

(3) $H^2(\CP^2;\Z)=\Z x$, $Q(x,x)=1$。$c=nx$（$n=\frac1{2\pi}\int_{\CP^1}F$）なら $Q(c,c)=n^2$、$S_\theta=\frac\theta2n^2$。$\theta\to\theta+2\pi$ で重みは $e^{i\pi n^2}=(-1)^{n^2}=(-1)^n$ 倍され、$n$ 奇数のセクターの符号が反転するので理論は変わる。$\theta\to\theta+4\pi$ では $e^{2\pi in^2}=1$ で不変。よって周期は $4\pi$。$k=n^2/2$ は半整数をとる。

(4) ウー公式 $Q(c,c)\equiv\langle\bar c\smile w_2,[X]\rangle\pmod2$ より、$\theta\to\theta+2\pi$ で重みは $e^{i\pi Q(c,c)}=(-1)^{\langle\bar c\smile w_2,[X]\rangle}$ 倍される。スピン多様体（$w_2=0$）では常に $1$ で (2) と一致し、$\CP^2$ では $w_2=\bar x$（$c_1(T\CP^2)=3x$ の mod 2 還元）なので $(-1)^n$ で (3) と一致する。
物理的意味：$\theta\to\theta+2\pi$ のずれは、ゲージ場と時空の幾何（$w_2$）の混合した位相項 $\pi\int\bar c\smile w_2$ にすぎない。この項は局所的な自由度をもたないので、「$\theta$ の周期は $2\pi$ だが、$\theta\to\theta+2\pi$ で重力との混合位相項がずれる」と言い換えることもできる。また、フェルミオンを含む理論で非スピン多様体を扱うには spin$^c$ 構造を用い、そのとき $F/2\pi$ の量子化条件自体が $w_2$ だけずれる（第9章で扱う）。いずれにせよ、$\theta$ の周期性は時空のコホモロジー環（交叉形式の偶奇）に依存する。
\end{solution}

\subsection{絡み数とチャーン・サイモンズ理論}

\begin{problem}[難易度 D]\label{ch06-prob-cs}
(1) $\R^3$ 内の交わらない閉曲線 $C_1,C_2$ について、ガウスの積分 $\mathrm{Lk}(C_1,C_2)$（動機の節）は、$C_1$ を流れる単位電流が作るビオ・サバールの磁場
\[
\bm B_1(\bm x)=\frac1{4\pi}\oint_{C_1}\frac{d\bm x_1\times(\bm x-\bm x_1)}{\lvert\bm x-\bm x_1\rvert^3}
\]
の周回積分 $\oint_{C_2}\bm B_1\cdot d\bm x_2$ に等しいことを示せ。アンペールの法則（$\nabla\times\bm B_1$ は $C_1$ に沿った単位電流）を用いて、これが $C_1\cdot S_2$（$S_2$ は $C_2$ が張る曲面）に等しく整数であることを示せ。\\
(2) $C_1(s)=(\cos s,\sin s,0)$, $C_2(t)=(1+\cos t,0,-\sin t)$（$s,t\in[0,2\pi]$）は交わらないことを確かめ、$\mathrm{Lk}(C_1,C_2)=+1$ を示せ（ホップ絡み目）。\\
(3) アーベリアン・チャーン・サイモンズ理論 $S=\frac k{4\pi}\int_{\R^3}A\wedge dA+\sum_jq_j\oint_{C_j}A$ を考える（重み $e^{iS}$、$A$ は $U(1)$ ゲージ場、$q_j\in\Z$）。ベクトル表示で $A\wedge dA=\bm A\cdot(\nabla\times\bm A)\,d^3x$、源の項を $\int\bm A\cdot\bm J\,d^3x$ と書く。運動方程式 $\frac k{2\pi}\nabla\times\bm A=-\bm J$ を導き（\textbf{流束の付着}）、作用の停留値が $S_*=\frac12\int\bm A_*\cdot\bm J\,d^3x$ であることを示せ。ガウス積分では $\langle\prod_jW_{q_j}(C_j)\rangle\propto e^{iS_*}$ であることを認めて、
\[
\bigl\langle W_{q_1}(C_1)W_{q_2}(C_2)\bigr\rangle\propto\exp\Bigl(-\frac{2\pi i}{k}q_1q_2\,\mathrm{Lk}(C_1,C_2)\Bigr)
\]
を示せ（比例係数は各ループ単独の寄与で、$C_1,C_2$ の相対位置によらない）。これがエニオンの組みひも位相にあたることを説明せよ。
\end{problem}

\begin{solution}
(1) $\oint_{C_2}\bm B_1\cdot d\bm x_2=\frac1{4\pi}\oint\oint\frac{(d\bm x_1\times(\bm x_2-\bm x_1))\cdot d\bm x_2}{\lvert\bm x_2-\bm x_1\rvert^3}$。スカラー三重積の巡回性より
\[
(d\bm x_1\times(\bm x_2-\bm x_1))\cdot d\bm x_2=(\bm x_2-\bm x_1)\cdot(d\bm x_2\times d\bm x_1)=(\bm x_1-\bm x_2)\cdot(d\bm x_1\times d\bm x_2),
\]
これはガウスの積分の被積分関数そのもの。よって $\mathrm{Lk}(C_1,C_2)=\oint_{C_2}\bm B_1\cdot d\bm x_2$。
$C_2$ が張る向き付き曲面 $S_2$ をとり（$C_1$ と横断的に交わるとする）、ストークスの定理を使うと $\oint_{C_2}\bm B_1\cdot d\bm x=\int_{S_2}(\nabla\times\bm B_1)\cdot d\bm S$。$\nabla\times\bm B_1=\bm J_1$ は $C_1$ に集中した単位電流 $\bm J_1(\bm x)=\oint_{C_1}d\bm x_1\,\delta^3(\bm x-\bm x_1)$ なので、右辺は $C_1$ が $S_2$ を貫く各点で $\pm1$（$C_1$ の向きが $S_2$ の法線と同じなら $+1$）の和、すなわち $C_1\cdot S_2\in\Z$。
\begin{note}
$\bm B_1$ は $C_1$ 上で特異だが、$C_1\cap C_2=\emptyset$ なので $C_2$ の近くでは滑らかであり、$\bm B_1$ のビオ・サバール表示から $\nabla\cdot\bm B_1=0$、$\nabla\times\bm B_1=\bm J_1$（超関数の意味）が従う。$S_2$ を別の曲面 $S_2'$ に取り替えると、差 $S_2-S_2'$ は閉曲面で、それを $C_1$ が貫く符号付き回数は $0$（閉曲線は閉曲面の内部に入ったら出る）なので、$C_1\cdot S_2$ は $S_2$ によらない。
\end{note}

(2) $C_2$ 上の点で $C_1$ に乗りうるのは $z=-\sin t=0$ の点 $(2,0,0)$, $(0,0,0)$ だけで、どちらも $\lvert\bm x\rvert\neq1$ なので交わらない。$C_2$ は $xz$ 平面内の中心 $(1,0,0)$、半径1の円で、$S_2$ をその円板とする。$t=0$ で $C_2=(2,0,0)$、接ベクトル $(0,0,-1)$。中心から外向きの $(1,0,0)$ と接ベクトルの外積 $(1,0,0)\times(0,0,-1)=(0,1,0)$ が、境界の向きと右ねじの関係にある法線である（円の向きを $\bm r\times\dot{\bm r}$ で判定）。$C_1$ は円板と $(1,0,0)$（$s=0$）でただ1回交わり、そこでの接ベクトルは $(0,1,0)$ で法線と同じ向き。よって $\mathrm{Lk}=C_1\cdot S_2=+1$。（$C_1$ は $z=0$ 平面内にあり、$xz$ 平面とは $(\pm1,0,0)$ で交わるが、$(-1,0,0)$ は円板の外。）

(3) 変分：$\delta\int\bm A\cdot(\nabla\times\bm A)=\int\bigl(\delta\bm A\cdot\nabla\times\bm A+\bm A\cdot\nabla\times\delta\bm A\bigr)=2\int\delta\bm A\cdot(\nabla\times\bm A)$（部分積分で表面項を落とす）。よって
\[
\delta S=\int\delta\bm A\cdot\Bigl(\frac k{2\pi}\nabla\times\bm A+\bm J\Bigr)d^3x=0\ \Longrightarrow\ \frac k{2\pi}\nabla\times\bm A_*=-\bm J .
\]
第3成分を時間とみると $(\nabla\times\bm A)_3$ は磁場なので、静止した電荷 $q$ の粒子の位置に磁束 $-2\pi q/k$ が付着する（粒子が電荷と磁束を同時にもつ）。停留値は
\[
S_*=\frac k{4\pi}\int\bm A_*\cdot\Bigl(-\frac{2\pi}k\bm J\Bigr)+\int\bm A_*\cdot\bm J=\frac12\int\bm A_*\cdot\bm J .
\]
$\bm J=\sum_jq_j\bm J_j$ とし、$\nabla\times\bm B_j=\bm J_j$ をみたすビオ・サバール場 $\bm B_j$ を用いると $\bm A_*=-\frac{2\pi}k\sum_jq_j\bm B_j$（ゲージ変換を除く；$\bm J$ が保存するので源の項はゲージ不変）。よって
\[
S_*=-\frac\pi k\sum_{i,j}q_iq_j\int\bm B_i\cdot\bm J_j\,d^3x=-\frac\pi k\sum_{i,j}q_iq_j\oint_{C_j}\bm B_i\cdot d\bm x .
\]
$i\neq j$ の項は (1) より $\oint_{C_j}\bm B_i\cdot d\bm x=\mathrm{Lk}(C_i,C_j)$ で、$(i,j)=(1,2),(2,1)$ の2項から
\[
S_*=-\frac{2\pi}kq_1q_2\,\mathrm{Lk}(C_1,C_2)+(\text{$i=j$ の項}).
\]
$i=j$ の項（自己絡み）は各ループ単独で決まり、定義には枠付け（フレーミング）が必要だが、$C_1,C_2$ の相対位置によらない。よって $\langle W_{q_1}W_{q_2}\rangle\propto\exp\bigl(-\frac{2\pi i}kq_1q_2\mathrm{Lk}\bigr)$。

エニオン：$2+1$ 次元時空で、電荷 $q_2$ の粒子が電荷 $q_1$ の粒子のまわりを一周する過程の世界線は、（対生成・対消滅で閉じると）$\mathrm{Lk}=\pm1$ の絡み目になる。その振幅は絡んでいない場合と比べて位相 $e^{\mp2\pi iq_1q_2/k}$ だけ異なり、これが組みひも（ブレイディング）の位相である。$q_1q_2/k\notin\Z$ なら位相は自明でなく、粒子はボソンでもフェルミオンでもない統計（エニオン）に従う。この位相は時空の計量によらず、絡み数という位相不変量だけで決まる。分数量子ホール状態（$\nu=1/k$ のラフリン状態）の準粒子の分数統計はこの仕組みによる。
\end{solution}

%================================================
\chapter{ファイバー束と主束}\label{ch:07}
% 18問
%================================================

\section{物理からの動機}

\subsection*{ディラック・モノポールのベクトル・ポテンシャル}
原点に磁荷をもつ静的な磁場を考える。第5章の規約では、原点を囲む球面 $S^2$ 上で
\[
\int_{S^2}F=2\pi n\qquad(n\in\Z)
\]
である（ディラックの量子化条件）。$n\neq0$ ならば、$S^2$ 全体で滑らかなベクトル・ポテンシャル $A$（$F=dA$）は存在しない。存在すればストークスの定理により $\int_{S^2}dA=\int_{\partial S^2}A=0$ となるからである。ディラックは、ポテンシャルが特異になる半直線（\term{ディラック・ストリング}{Dirac string}）を許すことでこの困難を回避した。

ウー（T.~T.~Wu）とヤン（C.~N.~Yang）は、より明快な描像を与えた。$S^2$ を2つの開集合
\[
U_N=S^2\setminus\{\text{南極}\},\qquad U_S=S^2\setminus\{\text{北極}\}
\]
で覆い、それぞれの上で滑らかなポテンシャル
\[
A_N=\frac n2(1-\cos\theta)\,d\varphi,\qquad A_S=-\frac n2(1+\cos\theta)\,d\varphi
\]
を用いる。重なり $U_N\cap U_S$（赤道を含む帯）では
\[
A_N-A_S=n\,d\varphi=i\,g\,dg^{-1},\qquad g=e^{in\varphi}
\]
であり、2つの記述は\textbf{ゲージ変換} $g_{NS}=e^{in\varphi}$ で結ばれている。荷電粒子の波動関数も $\psi_N=e^{in\varphi}\psi_S$ と貼り合わされる。ここで $e^{in\varphi}$ が赤道上で1価であるためには $n\in\Z$ が必要であり、ディラックの量子化条件は「貼り合わせのゲージ変換が well-defined であること」として現れる。
\begin{center}
\begin{tikzpicture}[scale=1.15]
  \shade[ball color=blue!8] (0,0) circle (1.5);
  \draw[thick] (0,0) circle (1.5);
  \fill[orange!25,opacity=0.7] (-1.5,0) arc (180:360:1.5 and 0.38) -- (1.5,0.25) arc (0:-180:1.5 and 0.38) -- cycle;
  \draw[gray,dashed] (1.5,0) arc (0:180:1.5 and 0.38);
  \draw[very thick,orange!80!black] (-1.5,0) arc (180:360:1.5 and 0.38);
  \node[vtx] at (0,1.5) {}; \node[vtx] at (0,-1.5) {};
  \node[above left] at (0,1.5) {\footnotesize 北極};
  \node[below left] at (0,-1.5) {\footnotesize 南極};
  \draw[very thick,red!70!black,densely dashed] (0,-1.55) -- (0,-2.6);
  \node[red!70!black,right] at (0.05,-2.35) {\footnotesize $A_N$ のストリング};
  \draw[very thick,blue!60!black,densely dashed] (0,1.55) -- (0,2.6);
  \node[blue!60!black,right] at (0.05,2.35) {\footnotesize $A_S$ のストリング};
  \node at (-0.45,0.75) {\footnotesize $U_N$：$A_N$};
  \node at (-0.45,-0.8) {\footnotesize $U_S$：$A_S$};
  \draw[-{Stealth[length=2mm]}] (2.6,0.1) -- (1.55,-0.2);
  \node[right,align=left] at (2.6,0.15) {\footnotesize 重なり（赤道付近）\\[-1pt]\footnotesize $A_N=A_S+i\,g\,dg^{-1}$,\ $g=e^{in\varphi}$};
\end{tikzpicture}
\zucap{ウー・ヤンの2枚の地図。$A_N$ は南極（を通る半直線）で、$A_S$ は北極で特異になる（波線：ディラック・ストリング）。ストリングは各地図を全球に延長しようとしたときにだけ現れる見かけのものである。}
\end{center}

\subsection*{ゲージ理論は「貼り合わせ」で定義される}
一般に、時空（あるいは空間）$M$ 上のゲージ理論は次のように記述される。
\begin{itemize}
\item $M$ を開集合 $\{U_i\}$ で覆い、各 $U_i$ 上でゲージ場 $A_i$ と物質場 $\psi_i$ を与える。
\item 重なり $U_i\cap U_j$ では、ゲージ変換 $g_{ij}\colon U_i\cap U_j\to G$ によって $\psi_i=g_{ij}\psi_j$、$A_i=g_{ij}A_jg_{ij}^{-1}+i\,g_{ij}\,dg_{ij}^{-1}$。
\item 3つの重なり $U_i\cap U_j\cap U_k$ で矛盾しないために $g_{ij}g_{jk}=g_{ik}$。
\end{itemize}
これが\term{ファイバー束}{fibre bundle}（より正確には\term{主束}{principal bundle}とその\term{同伴ベクトル束}{associated vector bundle}）の定義そのものである。$g_{ij}$ を連続的に変形しても（局所的なゲージ変換で）移り合わないとき、束は位相的に異なり、その違いが磁荷やインスタントン数のような\textbf{位相的量子数}になる。例えば $S^4$ 上の $SU(2)$ 束は、赤道 $S^3$ 上の貼り合わせ $g\colon S^3\to SU(2)$ の巻き付き数で分類され、これがインスタントン数である。

フェルミオンを扱うには、さらに\term{スピン構造}{spin structure}が必要である。円周上のフェルミオンには周期的（ラモン，R）と反周期的（ヌヴー・シュワルツ，NS）の2つの境界条件があるが、これは $S^1$ 上の2つのスピン構造にほかならない。弦理論の世界面の分配関数はスピン構造についての和をとり、そのモジュラー不変性から GSO 射影が導かれる。また、$SU(N)$ ゲージ理論をトーラス上で考えると、周期境界条件を中心 $\Z_N$ だけ「ねじる」ことができ（'t Hooft のねじれた境界条件）、これは $SU(N)/\Z_N$ 束の位相的な類として理解される。

\begin{physbox}[：本章の辞書]
\begin{center}
\small
\begin{tabular}{p{0.46\linewidth}p{0.44\linewidth}}
物理 & 数学\\ \hline
ゲージ群 $G$ & 構造群 $G$、主 $G$ 束 $P\to M$\\
局所的なゲージの選択 & 局所切断 $s_i\colon U_i\to P$（局所自明化）\\
重なりでのゲージ変換 $\psi_i=g_{ij}\psi_j$ & 推移関数 $g_{ij}$、コサイクル条件\\
ゲージ変換（場の再定義） & 推移関数の取り換え $h_ig_{ij}h_j^{-1}$\\
表現 $V$ に属する物質場 & 同伴ベクトル束 $P\times_GV$ の切断\\
ヒッグス場 & 同伴束の切断（零点が位相欠陥）\\
磁荷、インスタントン数、't Hooft フラックス & 束の同型類（貼り合わせ写像のホモトピー類）\\
ディラック・ストリング & 1枚の地図で済ませようとした痕跡\\
フェルミオンの境界条件（R/NS） & スピン構造\\
ゲージ場（接続）、場の強さ（曲率） & 第8章
\end{tabular}
\end{center}
\end{physbox}

\section{基本事項のまとめ}

\subsection{ファイバー束と推移関数}

\begin{mydef}{ファイバー束}{ch07-bundle}
空間 $E$, $M$, $F$ と連続な全射 $\pi\colon E\to M$ について、$M$ の開被覆 $\{U_i\}$ と同相写像
\[
\Phi_i\colon U_i\times F\xrightarrow{\ \cong\ }\pi^{-1}(U_i),\qquad \pi\circ\Phi_i(x,f)=x
\]
（\term{局所自明化}{local trivialization}）が存在するとき、$\pi\colon E\to M$ をファイバー $F$ の\textbf{ファイバー束}という。$E$ を\term{全空間}{total space}、$M$ を\term{底空間}{base space}、$E_x=\pi^{-1}(x)\cong F$ を $x$ 上の\term{ファイバー}{fibre}という。$E\cong M\times F$（ファイバーを保つ同相）のとき\term{自明束}{trivial bundle}という。

群 $G$ が $F$ に忠実に作用し、重なりで
\[
\Phi_i^{-1}\circ\Phi_j(x,f)=(x,\ g_{ij}(x)f),\qquad g_{ij}\colon U_i\cap U_j\to G\ \text{（連続）}
\]
となるとき、$G$ を\term{構造群}{structure group}、$g_{ij}$ を\term{推移関数}{transition function}という。
\end{mydef}

点 $e\in\pi^{-1}(U_i\cap U_j)$ の「$i$ 番目の地図での座標」を $f_i$ と書くと $f_i=g_{ij}f_j$ である。本書ではこれを物理の記法に合わせて
\[
\psi_i=g_{ij}\,\psi_j\qquad\text{（$U_i\cap U_j$ 上）}
\]
と書く。定義から直ちに次が従う（問題 \ref{ch07-prob-cocycle}）。

\begin{myprop}{コサイクル条件と束の同値}{ch07-cocycle}
\begin{enumerate}
\item $g_{ii}=1$、$g_{ji}=g_{ij}^{-1}$、そして $U_i\cap U_j\cap U_k$ 上で
\[
g_{ij}\,g_{jk}\,g_{ki}=1\qquad(\text{\term{コサイクル条件}{cocycle condition}}).
\]
\item 局所自明化を $\psi_i'=h_i\psi_i$（$h_i\colon U_i\to G$）と取り換えると、推移関数は
\[
g_{ij}'=h_i\,g_{ij}\,h_j^{-1}
\]
に変わる。これは物理でいう\textbf{ゲージ変換}である。
\item 逆に、コサイクル条件をみたす $\{g_{ij}\}$ が与えられれば、$\bigsqcup_iU_i\times F$ を $(x,f)_j\sim(x,g_{ij}(x)f)_i$ で貼り合わせて束が作れる。2組の推移関数が同じ束を定めるのは、ある $h_i$ により $g'_{ij}=h_ig_{ij}h_j^{-1}$ となるとき、かつそのときに限る（共通の細分をとった上で）。
\end{enumerate}
\end{myprop}

\begin{mydef}{引き戻し}{ch07-pullback}
束 $\pi\colon E\to M$ と連続写像 $f\colon N\to M$ に対し
\[
f^*E=\{(y,e)\in N\times E:\ f(y)=\pi(e)\}\longrightarrow N,\qquad (y,e)\mapsto y
\]
を\term{引き戻し束}{pullback bundle}という。推移関数は $g_{ij}\circ f$（被覆 $\{f^{-1}(U_i)\}$ 上）である。$f\simeq f'$ ならば $f^*E\cong f'^*E$（底空間がコンパクトなど、適当な条件の下で）。特に可縮な空間上の束は自明である。
\end{mydef}

\subsection{主束と同伴束}

\begin{mydef}{主束}{ch07-principal}
リー群 $G$ が空間 $P$ に右から自由に作用し（$p\mapsto pg$）、$\pi\colon P\to M=P/G$ が、$G$ 同変な局所自明化 $U_i\times G\cong\pi^{-1}(U_i)$、$(x,g)\mapsto s_i(x)g$ をもつとき、$P$ を\textbf{主 $G$ 束}という。ここで $s_i\colon U_i\to P$ は\term{局所切断}{local section}である。重なりでは、$G$ の作用がファイバー上で単純推移的なので
\[
s_j(x)=s_i(x)\,g_{ij}(x)
\]
となる $g_{ij}(x)\in G$ が一意に決まり、これが推移関数である（ファイバー $G$ に左からかけ算で作用）。
\end{mydef}

\begin{mydef}{同伴ベクトル束}{ch07-assoc}
主 $G$ 束 $P$ と表現 $\rho\colon G\to GL(V)$ に対し
\[
P\times_GV=(P\times V)/\bigl((pg,v)\sim(p,\rho(g)v)\bigr)
\]
を\textbf{同伴ベクトル束}という。$U_i$ 上の点を $[s_i(x),\psi_i]$ と表すと、$s_j=s_ig_{ij}$ から $[s_j,\psi_j]=[s_i,\rho(g_{ij})\psi_j]$、すなわち
\[
\psi_i=\rho(g_{ij})\,\psi_j .
\]
物質場（表現 $V$ に属する場）は $P\times_GV$ の\term{切断}{section}であり、局所的には $V$ 値関数の組 $\{\psi_i\}$ で、重なりで上の関係をみたすものである。局所切断を $s_i'=s_ih_i^{-1}$ に取り換えると $\psi_i'=\rho(h_i)\psi_i$ となる（ゲージ変換 $\psi\mapsto g\psi$）。
\end{mydef}

\begin{mythm}{主束の自明性}{ch07-trivsec}
主 $G$ 束 $P\to M$ が自明であるための必要十分条件は、大域的な切断 $s\colon M\to P$ が存在することである（問題 \ref{ch07-prob-trivsec}）。物理的にいえば、「全空間で1つのゲージ（地図）を選べる」ことと「束が自明」であることは同値である。
\end{mythm}

ベクトル束は零切断を必ずもつので、自明性は「至る所1次独立な $k$ 個の切断（大域的な枠）が存在すること」と同値になる。直和 $E\oplus E'$、テンソル積 $E\otimes E'$、双対 $E^*$、行列式束 $\det E=\Lambda^{\mathrm{top}}E$ は、推移関数をそれぞれ $g\oplus g'$, $g\otimes g'$, $(g^{T})^{-1}$, $\det g$ として作る。

\begin{myex}{代表的なベクトル束}{ch07-vbex}
\begin{itemize}
\item \textbf{接束} $TM$：座標 $x_i^\mu$ を地図とすると、ベクトルの成分は $v_i^\mu=\dfrac{\partial x_i^\mu}{\partial x_j^\nu}v_j^\nu$ と変換し、推移関数はヤコビ行列 $g_{ij}=\partial x_i/\partial x_j\in GL(n,\R)$。\textbf{余接束} $T^*M$ は $(g_{ij}^T)^{-1}$。
\item \textbf{枠束} $F(M)$：各点の接空間の基底全体からなる主 $GL(n,\R)$ 束。計量を入れて正規直交枠に制限すると主 $O(n)$ 束、向きも込めると主 $SO(n)$ 束。$TM=F(M)\times_{GL(n)}\R^n$。
\item \textbf{トートロジカル直線束}：$\CP^n$ の点 $[z]$ に直線 $\C z\subset\C^{n+1}$ を対応させる複素直線束 $L=\mathcal O(-1)=\{([z],v):v\in\C z\}$。その双対を $\mathcal O(1)$、$\mathcal O(m)=\mathcal O(1)^{\otimes m}$（$m<0$ は $\mathcal O(-1)^{\otimes\lvert m\rvert}$）とする。
\item \textbf{メビウス束}：$E=[0,1]\times\R/\bigl((0,t)\sim(1,-t)\bigr)\to S^1$。構造群 $O(1)=\{\pm1\}$ の実直線束で、自明でない（例 \ref{ex:ch07-mobius}、問題 \ref{ch07-prob-mobius}）。
\end{itemize}
\end{myex}

\begin{myex}{メビウス束}{ch07-mobius}
$S^1=[0,1]/(0\sim1)$ を $U_1=S^1\setminus\{s=0\}$、$U_2=S^1\setminus\{s=\tfrac12\}$ で覆うと、$U_1\cap U_2$ は2つの弧 $(0,\tfrac12)$, $(\tfrac12,1)$ からなり、推移関数は一方で $+1$、他方で $-1$ である。
\begin{center}
\begin{tikzpicture}[scale=1.1]
  \fill[blue!6] (0,-1) rectangle (4,1);
  \draw[very thick,green!50!black] (0,1) -- (4,1);
  \draw[very thick,orange!90!black] (0,-1) -- (4,-1);
  \draw[thick,red!70!black,mid arrow] (0,-1) -- (0,1);
  \draw[thick,red!70!black,mid arrow] (4,1) -- (4,-1);
  \draw[thick,gray] (1.3,-1) -- (1.3,1) node[above] {\footnotesize $\pi^{-1}(s)$};
  \draw[dashed] (0,0) -- (4,0) node[right] {\footnotesize 零切断};
  \node[hvtx,inner sep=1.8pt] at (4,1) {}; \node[hvtx,inner sep=1.8pt] at (0,-1) {};
  \node[vtx,inner sep=1.6pt] at (4,-1) {}; \node[vtx,inner sep=1.6pt] at (0,1) {};
  \draw[-{Stealth[length=2.2mm]}] (2,-1.35) -- (2,-1.95) node[midway,right] {$\pi$};
  \draw[thick] (0,-2.3) -- (4,-2.3);
  \node[vtx] at (0,-2.3) {}; \node[vtx] at (4,-2.3) {};
  \node[below] at (0,-2.3) {\footnotesize $0$}; \node[below] at (4,-2.3) {\footnotesize $1$};
  \node[right] at (4.1,-2.3) {$[0,1]/(0\sim1)\cong S^1$};
\end{tikzpicture}
\zucap{メビウス束（ファイバー $\R$ を $[-1,1]$ に切り詰めて描いた）。左右の辺を矢印の向きがそろうように、すなわち $(0,t)\sim(1,-t)$ で貼る。切断は $t(0)=-t(1)$ をみたす関数 $t(s)$ なので、中間値の定理により必ず零点をもつ。}
\end{center}
\end{myex}

\begin{physbox}[：ゲージ理論と主束]
\begin{itemize}
\item ゲージ群 $G$ のゲージ理論の「場の配位空間」は、主 $G$ 束 $P$ の選択と、その上の接続（ゲージ場、第8章）および同伴束の切断（物質場）からなる。異なる位相型の $P$ は互いに連続変形で移れない\textbf{位相的セクター}を与える。
\item 局所的には $A_i$ はエルミート行列値1形式で、重なりで $A_i=g_{ij}A_jg_{ij}^{-1}+i\,g_{ij}\,dg_{ij}^{-1}$。このとき $D_i\psi_i=g_{ij}D_j\psi_j$、$F_i=g_{ij}F_jg_{ij}^{-1}$ となり（問題 \ref{ch07-prob-cocycle}）、$\tr F^k$ のようなゲージ不変量は大域的な微分形式になる。
\item $SU(N)$ ゲージ理論で全ての場が随伴表現（$SU(N)$ の中心 $\Z_N$ が自明に作用）ならば、実際の構造群は $SU(N)/\Z_N$ であり、$SU(N)$ 束にならない束も許される（問題 \ref{ch07-prob-thooft}）。
\end{itemize}
\end{physbox}

\subsection{球面上の束：貼り合わせによる分類}
$S^n$ を北半球と南半球をすこし広げた開集合 $U_N$, $U_S$ で覆う。どちらも可縮なので束は各半球上で自明であり、束の情報は赤道付近の推移関数 $g_{NS}$ だけにある。赤道付近は $S^{n-1}$ にホモトピー同値なので、$g_{NS}$ は写像 $S^{n-1}\to G$ とみなせる（\term{貼り合わせ写像}{clutching function}）。

\begin{center}
\begin{tikzpicture}[scale=1]
  \fill[blue!8] (0,0) circle (1.1); \draw[thick] (0,0) circle (1.1);
  \node at (0,0) {$U_N\simeq D^n$};
  \fill[red!8] (5,0) circle (1.1); \draw[thick] (5,0) circle (1.1);
  \node at (5,0) {$U_S\simeq D^n$};
  \draw[-{Stealth[length=2.2mm]},thick,orange!80!black] (1.2,0.25) to[bend left=20] node[midway,above] {\footnotesize $g_{NS}\colon S^{n-1}\to G$} (3.8,0.25);
  \node[below] at (0,-1.15) {\footnotesize $\partial U_N\cong S^{n-1}$};
  \node[below] at (5,-1.15) {\footnotesize $\partial U_S\cong S^{n-1}$};
  \draw[-{Stealth[length=2.2mm]},thick] (6.4,0) -- (7.4,0);
  \shade[ball color=gray!10] (8.7,0) circle (1.1); \draw[thick] (8.7,0) circle (1.1);
  \draw[very thick,orange!80!black] (7.6,0) arc (180:360:1.1 and 0.3);
  \draw[orange!80!black,dashed] (9.8,0) arc (0:180:1.1 and 0.3);
  \node at (8.7,0.55) {\footnotesize $D^n\times G$};
  \node at (8.7,-0.6) {\footnotesize $D^n\times G$};
  \node[below] at (8.7,-1.15) {\footnotesize $S^n$ 上の束};
\end{tikzpicture}
\zucap{貼り合わせによる構成。2つの自明束 $D^n\times G$ を、境界 $S^{n-1}$ の各点 $x$ で $g_{NS}(x)\in G$ によって貼り合わせると、$S^n$ 上の主 $G$ 束が得られる。}
\end{center}

\begin{mythm}{貼り合わせによる分類}{ch07-clutch}
$G$ を弧状連結なリー群、$n\ge1$ とする。$S^n$ 上の主 $G$ 束の同型類は、貼り合わせ写像 $g_{NS}\colon S^{n-1}\to G$ のホモトピー類と1対1に対応する：
\[
\{S^n\text{ 上の主 }G\text{ 束}\}/{\cong}\ \longleftrightarrow\ [S^{n-1},G]\cong\pi_{n-1}(G).
\]
$G$ が連結でないとき（$O(1)$ など）は $[S^{n-1},G]$ を $\pi_0(G)$ の共役作用で割ったものになる。
\end{mythm}
\begin{note}
証明の要点：局所自明化を取り換えると $g_{NS}'=h_Ng_{NS}h_S^{-1}$。$h_N$, $h_S$ は円板（可縮）上の写像なので定値写像とホモトピックであり、$G$ が弧状連結なのでさらに定数 $1$ とホモトピック。したがって $g'_{NS}\simeq g_{NS}$。逆に $g_{NS}\simeq g_{NS}'$ ならホモトピーを赤道の近傍の「首」に埋め込んで束の同型が作れる。最後に、$G$ は群なので基点を $1$ に移すことができ、自由ホモトピー類と基点付きホモトピー類は一致する（$\pi_1(G)$ の $\pi_{n-1}(G)$ への作用は自明）。
\end{note}

\begin{center}
\renewcommand{\arraystretch}{1.25}\small
\begin{tabular}{llll}
底空間 & 構造群 $G$ & 分類 $\pi_{n-1}(G)$ & 物理\\ \hline
$S^1$ & 連結な $G$ & $\pi_0(G)=0$ & 束は自明（ホロノミーは残る）\\
$S^1$ & $O(1)=\Z_2$ & $\Z_2$ & 自明束とメビウス束\\
$S^2$ & $U(1)$ & $\pi_1(U(1))=\Z$ & ディラック・モノポール\\
$S^2$ & $SU(N)$ & $0$ & 自明\\
$S^2$ & $U(N)$ & $\Z$（$\det$ で検出） & GNO モノポール（問題 \ref{ch07-prob-s2bundles}）\\
$S^2$ & $SO(3)$ & $\Z_2$ & $SU(2)$ に持ち上がらない束\\
$S^3$ & コンパクト連結 $G$ & $\pi_2(G)=0$ & 自明\\
$S^4$ & $SU(N)$ ($N\ge2$) & $\pi_3(SU(N))=\Z$ & インスタントン数\\
$S^4$ & $U(1)$ & $\pi_3(U(1))=0$ & 自明
\end{tabular}
\end{center}

\begin{myex}{モノポール束}{ch07-monopole}
$S^2$ を極座標 $(\theta,\varphi)$ で表し、向きは $d\theta\wedge d\varphi$（外向き法線に関する通常の向き）とする。$U_N\ni$ 北極、$U_S\ni$ 南極とし、
\[
g_{NS}=e^{in\varphi}\qquad(\psi_N=e^{in\varphi}\psi_S)
\]
で貼り合わせた主 $U(1)$ 束を $P_n$ と書く。整数 $n$ は、赤道を $\varphi$ が増える向き（北半球の境界としての向き）に1周したときの $g_{NS}$ の巻き付き数であり、本書ではこれを束の\term{モノポール荷}{monopole charge}と呼ぶ。動機の節の $A_N$, $A_S$ はこの束の上のゲージ場で、$\dfrac1{2\pi}\int_{S^2}F=n$ となる（問題 \ref{ch07-prob-monopole}）。一般に「$g_{NS}$ の巻き付き数 $=$ 第1チャーン数 $\dfrac1{2\pi}\int F$」が成り立つ（第8章）。
\end{myex}

\begin{myex}{ホップ束}{ch07-hopf}
$S^3=\{(z_1,z_2)\in\C^2:\lvert z_1\rvert^2+\lvert z_2\rvert^2=1\}$ に $U(1)$ を $z\mapsto z\lambda$ で作用させ、ホップ写像（第1章）
\[
h\colon S^3\to S^2,\qquad h(z)=z^\dagger\bm\sigma z=\bigl(2\operatorname{Re}(\bar z_1z_2),\ 2\operatorname{Im}(\bar z_1z_2),\ \lvert z_1\rvert^2-\lvert z_2\rvert^2\bigr)
\]
を射影とすると主 $U(1)$ 束になる（\term{ホップ束}{Hopf bundle}）。$z=\bigl(\cos\tfrac\theta2\,e^{i\alpha},\ \sin\tfrac\theta2\,e^{i(\alpha+\varphi)}\bigr)$ とおくと $h(z)$ は極座標 $(\theta,\varphi)$ の点である。局所切断
\[
s_N=\bigl(\cos\tfrac\theta2,\ \sin\tfrac\theta2\,e^{i\varphi}\bigr),\qquad s_S=\bigl(\cos\tfrac\theta2\,e^{-i\varphi},\ \sin\tfrac\theta2\bigr)
\]
は $s_S=s_N\,e^{-i\varphi}$ をみたすので $g_{NS}=e^{-i\varphi}$、すなわち\textbf{ホップ束はモノポール荷 $n=-1$ の束 $P_{-1}$} である。電荷1の同伴直線束 $S^3\times_{U(1)}\C$ はトートロジカル直線束 $\mathcal O(-1)$ と同型であり、その双対 $\mathcal O(1)$（電荷 $-1$ の場）がモノポール荷 $+1$ をもつ（問題 \ref{ch07-prob-hopf}）。
\begin{center}
\begin{tikzpicture}[scale=1]
  \draw[thick,red!70!black] (1.5,0) arc (0:180:1.5 and 0.45);
  \draw[white,line width=4pt] (0,-2) -- (0,2);
  \draw[thick,blue!70!black,-{Stealth[length=2mm]}] (0,-2) -- (0,2.1);
  \draw[white,line width=4pt] (-1.5,0) arc (180:360:1.5 and 0.45);
  \draw[thick,red!70!black] (-1.5,0) arc (180:360:1.5 and 0.45);
  \node[blue!70!black,right] at (0.05,1.9) {\footnotesize $h^{-1}(\text{南極})$（$\infty$ を通る）};
  \node[red!70!black,right] at (1.5,-0.25) {\footnotesize $h^{-1}(\text{北極})$};
  \begin{scope}[xshift=6.2cm]
    \draw[thick,green!50!black] (1.6,0) arc (0:180:0.8 and 0.32);
    \draw[white,line width=4pt] (0,0) circle (0.8);
    \draw[thick,violet] (0,0) circle (0.8);
    \draw[white,line width=4pt] (0,0) arc (180:360:0.8 and 0.32);
    \draw[thick,green!50!black] (0,0) arc (180:360:0.8 and 0.32);
    \node[violet] at (-1.15,0.6) {\footnotesize $h^{-1}(x)$};
    \node[green!50!black] at (2.2,0.55) {\footnotesize $h^{-1}(y)$};
  \end{scope}
\end{tikzpicture}
\zucap{ホップ束のファイバー（立体射影で $S^3\setminus\{\text{1点}\}\cong\R^3$ に描いた模式図）。各ファイバーは円周で、異なる2本のファイバーは必ず1回絡み合う。$S^3$ は互いに絡み合う円周で埋め尽くされている。}
\end{center}
\end{myex}

\begin{myex}{$S^4$ 上の $SU(2)$ 束}{ch07-s4}
$S^4=\R^4\cup\{\infty\}$ を $U_N=\{\lvert x\rvert<R+\epsilon\}$ と $U_S=\{\lvert x\rvert>R-\epsilon\}\cup\{\infty\}$ で覆う。$SU(2)$ 束は $\pi_3(SU(2))=\Z$ で分類され、代表的な貼り合わせ写像は
\[
g_{NS}(x)=\frac{x_4+i\,\bm x\cdot\bm\sigma}{\lvert x\rvert}\qquad(x=(\bm x,x_4)),
\]
その巻き付き数は記法と規約の $w(g)=\frac{1}{24\pi^2}\int_{S^3}\tr(g^{-1}dg)^3$ で $w=1$ である。束のインスタントン数は $k=w(g_{NS})$ となる（問題 \ref{ch07-prob-s4}；第8章で再び扱う）。
\end{myex}

\subsection{ファイバー束のホモトピー完全列}
第3章で道具として用いた次の定理は、ファイバー束の\term{ホモトピー持ち上げ性質}{homotopy lifting property}から従う。

\begin{mythm}{ホモトピー長完全列}{ch07-les}
$F\to E\xrightarrow{\pi}M$ をファイバー束（$M$ は弧状連結、$e_0\in F=\pi^{-1}(x_0)$）とすると、完全列
\[
\begin{gathered}\cdots\to\pi_n(F)\xrightarrow{i_*}\pi_n(E)\xrightarrow{\pi_*}\pi_n(M)\xrightarrow{\partial}\pi_{n-1}(F)\to\cdots\\ \cdots\to\pi_1(M)\to\pi_0(F)\to\pi_0(E)\to0\end{gathered}
\]
がある。
\end{mythm}
\begin{note}
$\partial$ の作り方：$[f]\in\pi_n(M)$ を $f\colon(D^n,S^{n-1})\to(M,x_0)$ で表し、持ち上げ性質で $\tilde f\colon D^n\to E$（$\pi\circ\tilde f=f$）を作る。$\tilde f(S^{n-1})\subset F$ なので $\partial[f]=[\tilde f|_{S^{n-1}}]$。完全性の確認は定義を追うだけである。主 $U(1)$ 束 $P_n\to S^2$ では、$\partial\colon\pi_2(S^2)=\Z\to\pi_1(U(1))=\Z$ は $\pm n$ 倍である（貼り合わせ写像そのものが $\partial$ の像を与える）。
\end{note}

代表的な例は、ホップ束 $S^1\to S^3\to S^2$、$SU(N-1)\to SU(N)\to S^{2N-1}$、$SO(n-1)\to SO(n)\to S^{n-1}$、$SU(N)\to U(N)\xrightarrow{\det}U(1)$、そして被覆空間（離散ファイバー）である。

\subsection{向きとスピン構造}

\begin{mydef}{向き付け可能性}{ch07-orient}
$n$ 次元多様体 $M$ が\term{向き付け可能}{orientable}であるとは、接束の推移関数（ヤコビ行列）を $\det g_{ij}>0$ となるように選べることをいう。リーマン計量を入れて $g_{ij}\in O(n)$ とすると、$\det g_{ij}\in\{\pm1\}$ は $\Z_2$ に値をとるコサイクルであり、その類
\[
w_1(M)\in H^1(M;\Z_2)
\]
（\term{第1シュティーフェル・ホイットニー類}{first Stiefel--Whitney class}）が $0$ であることが向き付け可能性と同値である。
\end{mydef}
例えば $\RP^n$ は $n$ が奇数のとき、かつそのときに限り向き付け可能である（対蹠写像 $x\mapsto-x$ が $S^n$ の向きを保つのは $n$ が奇数のとき）。$\RP^2$ はメビウスの帯を含むので向き付け不可能である。

\begin{mydef}{スピン群とスピン構造}{ch07-spin}
$n\ge3$ のとき $SO(n)$ の普遍被覆を $\Spin(n)$ と書き、2対1の準同型 $\rho\colon\Spin(n)\to SO(n)$（核 $\{\pm1\}$）をもつ。$\Spin(3)=SU(2)$、$\Spin(4)=SU(2)\times SU(2)$。低次元では $\Spin(2)=U(1)\to SO(2)=U(1)$、$\lambda\mapsto\lambda^2$、$\Spin(1)=\{\pm1\}\to SO(1)=1$ と定める。

向きと計量をもつ $M$ の正規直交枠束（主 $SO(n)$ 束）を $P_{SO}$ とする。主 $\Spin(n)$ 束 $\tilde P$ と2対1の写像 $\tilde P\to P_{SO}$ で $\rho$ と両立するものを\textbf{スピン構造}という。推移関数でいえば、$g_{ij}\in SO(n)$ の持ち上げ $\tilde g_{ij}\in\Spin(n)$（$\rho(\tilde g_{ij})=g_{ij}$）で、コサイクル条件 $\tilde g_{ij}\tilde g_{jk}\tilde g_{ki}=1$ をみたすものを選ぶことである。
\end{mydef}

\begin{mythm}{スピン構造の存在と分類}{ch07-spinclass}
持ち上げは常に $\tilde g_{ij}\tilde g_{jk}\tilde g_{ki}=\epsilon_{ijk}\in\{\pm1\}$ をみたし、$\{\epsilon_{ijk}\}$ は $\Z_2$ 係数の2-コサイクルを定める。その類
\[
w_2(M)\in H^2(M;\Z_2)
\]
（\term{第2シュティーフェル・ホイットニー類}{second Stiefel--Whitney class}）が $0$ であることがスピン構造の存在と同値である。存在するとき、2つのスピン構造は $\tilde g_{ij}'=\eta_{ij}\tilde g_{ij}$（$\eta_{ij}\in\{\pm1\}$ は1-コサイクル）で移り合い、連結な $M$ のスピン構造は $H^1(M;\Z_2)$ と1対1に対応する。特に、その個数は $\lvert H^1(M;\Z_2)\rvert$ である。
\end{mythm}

\begin{center}
\renewcommand{\arraystretch}{1.2}\small
\begin{tabular}{lll}
$M$ & スピン構造の数 & 備考\\ \hline
$S^1$ & $2$ & R（周期的）と NS（反周期的）\\
$T^2$ & $4$ & (NS,NS), (NS,R), (R,NS), (R,R)\\
$\Sigma_g$ & $2^{2g}$ & 偶 $2^{g-1}(2^g+1)$ 個、奇 $2^{g-1}(2^g-1)$ 個（問題 \ref{ch07-prob-arf}）\\
$S^n$ ($n\ge2$) & $1$ & $H^1=0$ かつ $w_2=0$\\
$T^n$ & $2^n$ & \\
$\CP^2$ & $0$ & $w_2\neq0$：スピン多様体でない
\end{tabular}
\end{center}
$\CP^2$ がスピンでないことは、交叉形式 $(1)$ が奇であること（第6章）と、ウーの公式「$x\cdot x\equiv x\cdot w_2\pmod2$」から従う（ここでは証明しない）。このような多様体上でもフェルミオンを $U(1)$ ゲージ場と結合させれば定義できる（$\Spin^c$ 構造）。

\begin{physbox}[：スピン構造とフェルミオンの境界条件]
\begin{itemize}
\item 円周 $S^1$ 上で、接ベクトルに沿った枠を用いるとスピノルは $\psi(\sigma+2\pi)=\pm\psi(\sigma)$。$+$ が R（ラモン）、$-$ が NS（ヌヴー・シュワルツ）である。円板の境界として延長できる（\term{バウンディング}{bounding}な）のは NS である（問題 \ref{ch07-prob-spins1}）。円筒から平面への写像 $z=e^{w}$ で NS セクターが平面上の1価な場に対応するのはこのためである。
\item 有限温度の場の理論でフェルミオンが虚時間方向に反周期的なのは、$\Tr e^{-\beta H}$ が時間方向 NS のスピン構造に対応するからであり、$\Tr(-1)^Fe^{-\beta H}$ は時間方向 R に対応する。
\item 弦の世界面（種数 $g$）の経路積分はスピン構造について和をとる。その整合性（モジュラー不変性）が GSO 射影を導く（問題 \ref{ch07-prob-gso}）。
\end{itemize}
\end{physbox}

\section{問題}

\subsection{束・推移関数・切断}

\begin{problem}[難易度 A]\label{ch07-prob-cocycle}
$G$ を行列リー群とし、開被覆 $\{U_i\}$ の上で、$G$ の表現 $V$ に属する場 $\psi_i\colon U_i\to V$ とゲージ場 $A_i$（$\mathfrak g$ 値1形式）が、重なりで
\[
\psi_i=g_{ij}\psi_j,\qquad A_i=g_{ij}A_jg_{ij}^{-1}+i\,g_{ij}\,dg_{ij}^{-1}
\]
と貼り合わされているとする。ただし $\psi_j$ は任意の値をとりうるものとする。\\
(1) $g_{ii}=1$、$g_{ji}=g_{ij}^{-1}$、$g_{ij}g_{jk}g_{ki}=1$ を示せ。\\
(2) $g\ast A:=gAg^{-1}+i\,g\,dg^{-1}$ とおくとき $(gh)\ast A=g\ast(h\ast A)$ を示し、$A$ の貼り合わせ規則が3つの重なりで矛盾しないことを確かめよ。\\
(3) $D_i=d-iA_i$ とするとき $D_i\psi_i=g_{ij}D_j\psi_j$ を示せ。\\
(4) 各 $U_i$ でゲージ変換 $\psi_i'=h_i\psi_i$ を行うと、推移関数が $g_{ij}'=h_ig_{ij}h_j^{-1}$ に変わることを示せ。
\end{problem}

\begin{solution}
(1) $U_i$ 上で $\psi_i=g_{ii}\psi_i$ が任意の $\psi_i$ で成り立つので $g_{ii}=1$。$\psi_i=g_{ij}\psi_j=g_{ij}g_{ji}\psi_i$ より $g_{ij}g_{ji}=1$。$U_i\cap U_j\cap U_k$ 上で $\psi_i=g_{ij}\psi_j=g_{ij}g_{jk}\psi_k$ かつ $\psi_i=g_{ik}\psi_k=g_{ki}^{-1}\psi_k$ なので $g_{ij}g_{jk}=g_{ki}^{-1}$、すなわち $g_{ij}g_{jk}g_{ki}=1$。（$\rho$ が忠実でなければ $g$ の代わりに $\rho(g)$ についての等式になる。）

(2) $d(h^{-1}g^{-1})=(dh^{-1})g^{-1}+h^{-1}dg^{-1}$ を用いて
\[
i\,(gh)\,d(gh)^{-1}=i\,g\,h\,(dh^{-1})\,g^{-1}+i\,g\,dg^{-1}.
\]
一方
\[
g\ast(h\ast A)=g\bigl(hAh^{-1}+i\,h\,dh^{-1}\bigr)g^{-1}+i\,g\,dg^{-1}=(gh)A(gh)^{-1}+i\,gh\,(dh^{-1})\,g^{-1}+i\,g\,dg^{-1}
\]
であり、両者は一致する。よって $A_i=g_{ij}\ast A_j=g_{ij}\ast(g_{jk}\ast A_k)=(g_{ij}g_{jk})\ast A_k=g_{ik}\ast A_k$ となり、$U_i\cap U_k$ での規則 $A_i=g_{ik}\ast A_k$ と矛盾しない。ここでコサイクル条件を使った。

(3) $g=g_{ij}$ と書く。$g\,dg^{-1}\,g=-dg\,g^{-1}g=-dg$（$d(gg^{-1})=0$ より $dg^{-1}=-g^{-1}dg\,g^{-1}$）を用いると
\begin{align*}
D_i\psi_i&=d(g\psi_j)-i\bigl(gA_jg^{-1}+i\,g\,dg^{-1}\bigr)g\psi_j\\
&=dg\,\psi_j+g\,d\psi_j-i\,gA_j\psi_j+g\,dg^{-1}g\,\psi_j\\
&=dg\,\psi_j+g\,d\psi_j-i\,gA_j\psi_j-dg\,\psi_j=g(d-iA_j)\psi_j .
\end{align*}
共変微分は物質場と同じ規則で貼り合わされ、$D\psi$ も同伴束の（1形式値の）切断になる。

(4) $\psi_i'=h_i\psi_i=h_ig_{ij}\psi_j=h_ig_{ij}h_j^{-1}\psi_j'$。よって $g_{ij}'=h_ig_{ij}h_j^{-1}$。このとき $g'$ もコサイクル条件をみたす：$g'_{ij}g'_{jk}g'_{ki}=h_ig_{ij}g_{jk}g_{ki}h_i^{-1}=1$。
\end{solution}

\begin{problem}[難易度 A]\label{ch07-prob-trivsec}
(1) 主 $G$ 束 $\pi\colon P\to M$ が自明（$P\cong M\times G$、$G$ 同変かつファイバーを保つ）であるための必要十分条件は、大域切断 $s\colon M\to P$ が存在することであることを示せ。\\
(2) ランク $k$ の実ベクトル束 $E$ が自明であるための必要十分条件は、各点で1次独立な $k$ 個の大域切断が存在することであることを、枠束に (1) を適用して示せ。\\
(3) 「ゲージ固定」を局所切断の選択と解釈し、ディラック・モノポール（$n\neq0$）の背景では、全空間で特異点のないゲージを選ぶことができない理由を述べよ。束 $P_n$（$n\neq0$）が自明でないことは認めてよい。
\end{problem}

\begin{solution}
(1) 自明化 $\Phi\colon M\times G\to P$ があれば $s(x)=\Phi(x,1)$ が大域切断である。逆に大域切断 $s$ があるとき $\Phi(x,g)=s(x)g$ とおく。$\pi(\Phi(x,g))=x$ で、$\Phi(x,g)g'=\Phi(x,gg')$（同変）。各ファイバー上で $G$ は自由かつ推移的に作用するので $\Phi$ は全単射である。逆写像は $p\mapsto(\pi(p),\gamma(p))$、ただし $p=s(\pi(p))\gamma(p)$ で $\gamma(p)\in G$ を定める。$\gamma$ の連続性：$U_i$ 上の局所切断 $s_i$ を用いて $s(x)=s_i(x)a_i(x)$、$p=s_i(\pi(p))b_i(p)$ と書くと、$a_i$, $b_i$ は局所自明化から決まる連続関数で、$\gamma(p)=a_i(\pi(p))^{-1}b_i(p)$ は連続。よって $\Phi$ は同相で $P$ は自明。

(2) $E$ の枠束 $F(E)$（各ファイバーの基底全体、主 $GL(k,\R)$ 束）の切断は、各点で1次独立な $k$ 個の切断 $(e_1,\dots,e_k)$ にほかならない。$E\cong F(E)\times_{GL(k)}\R^k$ であり、$F(E)$ が自明なら $E$ も自明、逆に $E\cong M\times\R^k$ なら標準基底が大域的な枠を与える。(1) より主張が従う。

(3) $U$ 上のゲージの選択（各点で「位相の基準」を選ぶこと）は、主 $U(1)$ 束の $U$ 上の切断 $s$ を選ぶことである。そのとき場は $\psi=[s,\psi_U]$ によって通常の関数 $\psi_U$ で表され、ゲージ場も通常の1形式 $A_U$ になる。$S^2$ 全体で特異点のないゲージがあれば、それは大域切断であり (1) から $P_n$ は自明になる。$n\neq0$ ではこれは矛盾なので、少なくとも2枚の地図（ウー・ヤン）か、特異点（ディラック・ストリング）を許したゲージが必要である。ストリングの位置は切断の零点・特異点の位置であり、別のゲージに移れば動かせるが、消すことはできない。
\end{solution}

\begin{problem}[難易度 B]\label{ch07-prob-mobius}
例 \ref{ex:ch07-mobius} のメビウス束 $E=[0,1]\times\R/\bigl((0,t)\sim(1,-t)\bigr)\to S^1$ を考える。\\
(1) $U_1=\{0<s<1\}$、$U_2=\{s\neq\tfrac12\}$ 上の局所自明化を具体的に書き、推移関数を求めよ。\\
(2) $E$ の任意の切断は零点をもつことを示し、$E$ が自明束でないことを結論せよ。また、推移関数の取り換え $h_2g_{21}h_1^{-1}$（$h_i\colon U_i\to\R^\times$）で $g_{21}\equiv1$ にできないことを直接示せ。\\
(3) $E\oplus E$ は自明であることを示せ。\\
(4) 円周 $\sigma\sim\sigma+2\pi$ 上の場で反周期的境界条件 $\phi(\sigma+2\pi)=-\phi(\sigma)$ をみたすものを考える。実スカラー場1つではこれを周期的な場に「場の再定義」で直すことはできないが、複素スカラー場では $\phi=e^{i\sigma/2}\tilde\phi$ により周期的な $\tilde\phi$ に直せることを示せ。後者で、ねじれがどこに移るかを述べよ（自由場の運動項 $\lvert\partial_\sigma\phi\rvert^2$ を $\tilde\phi$ で書け）。
\end{problem}

\begin{solution}
(1) $U_1$ 上は $\Phi_1^{-1}([s,t])=(s,t)$、$0<s<1$。$U_2$ 上は
\[
\Phi_2^{-1}([s,t])=\begin{cases}(s,\ t)&(0\le s<\tfrac12)\\ (s,\ -t)&(\tfrac12<s\le1)\end{cases}
\]
とおく。同一視 $[0,t]=[1,-t]$ の両辺で第2成分はともに $t$ になるので well-defined。重なり $U_1\cap U_2$ は2つの弧 $(0,\tfrac12)$, $(\tfrac12,1)$ で、ファイバー座標の関係は $t_2=t_1$（前者）、$t_2=-t_1$（後者）。よって $\psi_2=g_{21}\psi_1$ の推移関数は
\[
g_{21}=+1\ \ \text{on }(0,\tfrac12),\qquad g_{21}=-1\ \ \text{on }(\tfrac12,1).
\]

(2) 切断は連続関数 $t\colon[0,1]\to\R$ で $t(1)=-t(0)$ をみたすもの（$[0,t(0)]=[1,t(1)]$ のため）と同じである。$t(0)\neq0$ なら $t(0)$ と $t(1)$ は異符号なので、中間値の定理により $t(s_0)=0$ となる $s_0$ がある。$t(0)=0$ なら $s=0$ が零点。自明束 $S^1\times\R$ には零点のない切断（定数 $1$）があるので、$E$ は自明でない。

推移関数による議論：$U_1$, $U_2$ はともに連結（弧）なので、連続な $h_i\colon U_i\to\R^\times$ の符号は一定である。すると $g'_{21}=h_2g_{21}h_1^{-1}$ の符号は、2つの弧の上で $\operatorname{sgn}(h_2h_1^{-1})\cdot(\pm1)$ となり、必ず一方で正、他方で負になる。よって $g'_{21}\equiv1$ にはできない。（これは $w_1(E)\neq0$ の最も簡単な例である。）

(3) $E\oplus E$ の推移関数は $g_{21}\mathbf 1_2$ で、$(\tfrac12,1)$ 上で $-\mathbf 1_2$ である。$-\mathbf 1_2$ は $SO(2)$ の角度 $\pi$ の回転なので、符号の議論は使えない。実際 $\R^2\cong\C$ と同一視すると $E\oplus E$ は推移関数 $\pm1\in U(1)$ の複素直線束であり、その切断は $\tau(1)=-\tau(0)$ をみたす連続関数 $\tau\colon[0,1]\to\C$ である。$\tau(s)=e^{i\pi s}$ はこれをみたし、どこでも零にならない。直線束（ランク1）で零点のない切断は大域的な枠なので、問題 \ref{ch07-prob-trivsec} (2) より $E\oplus E$ は自明である。（貼り合わせによる分類でいえば、$S^1$ 上の束は $\pi_0$ で分類され、$U(1)$ は連結なので複素直線束はすべて自明。）

(4) 実の反周期的な場は $E$ の切断であり、場の再定義（$\phi\mapsto h\phi$, $h$ は $\R^\times$ 値の周期関数）は推移関数を $h_2g_{21}h_1^{-1}$ に変えるだけなので、(2) によりねじれは除けない。複素場では $\tilde\phi(\sigma)=e^{-i\sigma/2}\phi(\sigma)$ とおくと
\[
\tilde\phi(\sigma+2\pi)=e^{-i\sigma/2}e^{-i\pi}\bigl(-\phi(\sigma)\bigr)=\tilde\phi(\sigma)
\]
で周期的になる。運動項は $\partial_\sigma\phi=e^{i\sigma/2}(\partial_\sigma+\tfrac i2)\tilde\phi$ より
\[
\lvert\partial_\sigma\phi\rvert^2=\bigl\lvert(\partial_\sigma-iA_\sigma)\tilde\phi\bigr\rvert^2,\qquad A=-\tfrac12\,d\sigma .
\]
すなわち束は自明になるが、ねじれは平坦な背景ゲージ場 $A$（$F=0$）に移り、そのホロノミー $e^{i\oint A}=e^{-i\pi}=-1$（ウィルソン・ライン）として残る。位相的に自明な束の上でも、平坦接続のホロノミーは物理的に意味をもつ（AB 効果と同じ構造；第2, 8章）。
\end{solution}

\begin{problem}[難易度 C]\label{ch07-prob-higgssplit}
$E\to M$ をエルミート計量をもつランク $k$ の複素ベクトル束とする。\\
(1) 零点をもたない切断 $\Phi$ が存在すれば、$E\cong\underline{\C}\oplus\Phi^\perp$（$\underline{\C}=M\times\C$、$\Phi^\perp$ は各点で $\Phi$ に直交する部分空間のなす束）であることを示せ。$\Phi^\perp$ が局所自明であることも確かめよ。\\
(2) $k=2$ で、$E$ が主 $SU(2)$ 束の基本表現の同伴束である（したがって $\det E=\Lambda^2E$ が自明）とする。零点のない切断があれば $E$ は自明であることを示せ。\\
(3) $S^4$ 上の $SU(2)$ ゲージ理論に、基本表現のヒッグス場（2重項）$\Phi$ を加える。インスタントン数 $k\neq0$ の束の上では、$\Phi$ は必ずどこかで零になることを示せ。ただし $k\neq0$ の束が自明でないことは認めてよい。
\end{problem}

\begin{solution}
(1) $U_i$ 上の局所的な枠 $(e_1,\dots,e_k)$ をとり、$\Phi$ が零にならないので（必要なら $U_i$ を小さくして並べ替え）$(\Phi,e_2,\dots,e_k)$ が各点で1次独立になるようにできる。これにグラム・シュミットの直交化を施すと、正規直交枠 $(\hat\Phi,f_2,\dots,f_k)$、$\hat\Phi=\Phi/\lvert\Phi\rvert$ が得られ、係数は計量と枠の連続関数である。$(f_2,\dots,f_k)$ は $\Phi^\perp$ の $U_i$ 上の枠なので $\Phi^\perp$ はランク $k-1$ のベクトル束になる。写像
\[
\underline{\C}\oplus\Phi^\perp\to E,\qquad (x,c)\oplus v\mapsto c\,\hat\Phi(x)+v
\]
は各ファイバーで線形同型で連続（局所的に枠で書けば行列が単位行列）なので、束の同型である。

(2) (1) より $E\cong\underline{\C}\oplus L$、$L=\Phi^\perp$ は直線束。写像
\[
L\to\Lambda^2E,\qquad v\mapsto\hat\Phi\wedge v
\]
は各ファイバーで線形かつ単射（$v\neq0$ は $\hat\Phi$ と1次独立なので $\hat\Phi\wedge v\neq0$）、1次元同士なので同型であり、$L\cong\Lambda^2E$。$\Lambda^2E$ の推移関数は $\det g_{ij}=1$ なので $\Lambda^2E$ は自明。よって $E\cong\underline{\C}\oplus\underline{\C}$ は自明。

さらに主 $SU(2)$ 束 $P$ も自明である：$\hat\Phi$ と $\hat\Phi^\perp$ の生成元 $\hat\Phi'$（$\det(\hat\Phi,\hat\Phi')=1$ となるように正規化）を並べた $(\hat\Phi,\hat\Phi')$ は各点で $SU(2)$ 枠、すなわち $P$ の大域切断を与える。

(3) 零点のない $\Phi$ があれば (2) により束は自明になり、$k\neq0$ に反する。したがって $\Phi$ は零点をもつ。物理的には、インスタントン背景でヒッグス2重項は（電弱理論の真空 $\lvert\Phi\rvert=v$ に一様に留まることができず）どこかで零になり、そこでエネルギー（作用）のコストが生じる。これは電弱理論のインスタントンが「束縛されたインスタントン」として大きさが固定されることや、スファレロンの構成とつながる。逆に、ヒッグス場が至る所で零でない配位では、束を自明化する（ユニタリ・ゲージを大域的にとる）ことができる。
\end{solution}

\subsection{球面上の束：モノポール、ホップ束、接束}

\begin{problem}[難易度 B]\label{ch07-prob-monopole}
$S^2$（向き $d\theta\wedge d\varphi$）上で
\[
A_N=\frac n2(1-\cos\theta)\,d\varphi\quad(\theta\neq\pi),\qquad A_S=-\frac n2(1+\cos\theta)\,d\varphi\quad(\theta\neq0)
\]
とする。\\
(1) $A_N$ が北極の近傍で滑らかな1形式であることを、直交座標 $x=\sin\theta\cos\varphi$, $y=\sin\theta\sin\varphi$ を用いて示せ（$A_S$ も南極で同様）。\\
(2) $F=dA_N=dA_S$ を求め、$\int_{S^2}F=2\pi n$ を示せ。\\
(3) 重なりで $A_N=A_S+i\,g\,dg^{-1}$ となる $g$ を求め、$g$ が1価であるために $n\in\Z$ が必要であることを示せ。\\
(4) $A_N$ を $\R^3\setminus\{\text{負の }z\text{ 軸}\}$ の1形式とみなすと、負の $z$ 軸を小さく囲む円周に沿った積分が半径 $\to0$ で $2\pi n$ に近づくことを示し、これが「ストリング」であることを説明せよ。また単位電荷の粒子にとってこのストリングが観測できない理由を述べよ。\\
(5) 推移関数を $g'_{NS}=h_Ng_{NS}h_S^{-1}$（$h_N\colon U_N\to U(1)$, $h_S\colon U_S\to U(1)$）に取り換えても、赤道に沿った巻き付き数 $n$ は変わらないことを示せ。
\end{problem}

\begin{solution}
(1) 北極付近では $\rho^2=x^2+y^2=\sin^2\theta$、$x\,dy-y\,dx=\rho^2d\varphi$。$1-\cos\theta=\dfrac{\sin^2\theta}{1+\cos\theta}$ なので
\[
A_N=\frac n2\,\frac{x\,dy-y\,dx}{1+\cos\theta},\qquad \cos\theta=\sqrt{1-x^2-y^2},
\]
は $(x,y)=(0,0)$ の近傍で滑らかである。一方 $\theta\to\pi$ では $\frac n2(1-\cos\theta)\to n\neq0$ の係数が $d\varphi$（極で特異）にかかるので特異になる。$A_S$ も $1+\cos\theta=\dfrac{\sin^2\theta}{1-\cos\theta}$ により南極で滑らか。

(2) $dA_N=\frac n2\sin\theta\,d\theta\wedge d\varphi=dA_S$。よって $F$ は $S^2$ 全体で定義された2形式で
\[
\int_{S^2}F=\frac n2\int_0^\pi\sin\theta\,d\theta\int_0^{2\pi}d\varphi=\frac n2\cdot2\cdot2\pi=2\pi n .
\]
（半径 $r$ の球面上の磁場 $\bm B=\frac{n}{2}\frac{\hat{\bm r}}{r^2}$ で、全磁束 $2\pi n$。）

(3) $A_N-A_S=\frac n2\bigl[(1-\cos\theta)+(1+\cos\theta)\bigr]d\varphi=n\,d\varphi$。$U(1)$ では $g=e^{i\alpha}$ に対し $i\,g\,dg^{-1}=i\,e^{i\alpha}(-i\,d\alpha)e^{-i\alpha}=d\alpha$ なので、$\alpha=n\varphi$、$g=e^{in\varphi}$ と選べる（定数倍の自由度を除く）。$\varphi\to\varphi+2\pi$ で $g$ が1価であるには $e^{2\pi in}=1$、すなわち $n\in\Z$。ゲージ場だけを見れば $A_N-A_S=n\,d\varphi$ は任意の実数 $n$ で成り立つが、電荷1の物質場 $\psi_N=g\psi_S$ を矛盾なく貼り合わせるには $g$ が1価でなければならない。これがディラックの量子化条件である。

(4) 負の $z$ 軸の近く（$\theta=\pi-\delta$）で $\varphi$ を一周する円周 $C_\delta$ をとると
\[
\oint_{C_\delta}A_N=\frac n2\bigl(1-\cos(\pi-\delta)\bigr)\cdot2\pi=\pi n(1+\cos\delta)\xrightarrow{\delta\to0}2\pi n .
\]
一方 $C_\delta$ が囲む小円板を通る $F$ の磁束は $\to0$ である。ストークスの定理と矛盾しないためには、$A_N$ が負の $z$ 軸上で特異で、そこに磁束 $2\pi n$ の無限に細いソレノイド（ストリング）があるとみなさなければならない。これは $U_N$ の地図を南極まで延長しようとしたことの痕跡である。単位電荷の粒子がストリングを一周するときの AB 位相は $e^{i\oint A}=e^{2\pi in}=1$ なので、どんな干渉実験でも観測されない。ストリングの位置は $A_S$ に移る（ゲージ変換）ことで北極側へ動かせ、物理的意味をもたない。

(5) 赤道付近の帯の写像 $S^1\to U(1)$ の巻き付き数を $\deg$ と書くと、積について $\deg(f_1f_2)=\deg f_1+\deg f_2$（偏角の増分が足し合わされる）。$h_N$ は円板 $U_N$ 全体で定義された連続写像なので、その赤道への制限は定値写像とホモトピックであり $\deg h_N|_{\text{赤道}}=0$（第1, 2章）。同様に $\deg h_S=0$。よって
\[
\deg g'_{NS}=\deg h_N+\deg g_{NS}-\deg h_S=n .
\]
$n$ は局所的なゲージ変換で変わらない、束の位相不変量である。
\end{solution}

\begin{problem}[難易度 B]\label{ch07-prob-section}
問題 \ref{ch07-prob-monopole} のモノポール背景（$g_{NS}=e^{in\varphi}$）での電荷1の粒子の波動関数は、$U_N$, $U_S$ 上の連続関数の組 $(\psi_N,\psi_S)$ で $\psi_N=e^{in\varphi}\psi_S$ をみたすもの、すなわち同伴直線束 $L_n$ の切断である。\\
(1) 確率密度 $\lvert\psi\rvert^2$ は $S^2$ 上の通常の関数であることを示せ。\\
(2) 零点 $p$ の\textbf{指数}を、$p$ を正の向きに回る小さな円周に沿った（$p$ を含む地図での）$\psi_i$ の巻き付き数と定める。零点が有限個で赤道上にないとき、指数の総和は $n$ に等しいことを示せ。特に $n\neq0$ なら、零点をもたない波動関数は存在しない。\\
(3) $n=1$ のとき
\[
\psi^{(1)}:\ (\psi_N,\psi_S)=\bigl(e^{i\varphi}\sin\tfrac\theta2,\ \sin\tfrac\theta2\bigr),\qquad \psi^{(2)}:\ (\psi_N,\psi_S)=\bigl(\cos\tfrac\theta2,\ e^{-i\varphi}\cos\tfrac\theta2\bigr)
\]
が切断であることを確かめ、それぞれの零点とその指数を求めよ。
\end{problem}

\begin{solution}
(1) 重なりで $\lvert\psi_N\rvert=\lvert e^{in\varphi}\rvert\lvert\psi_S\rvert=\lvert\psi_S\rvert$ なので、$\lvert\psi\rvert^2$ は $U_N$ と $U_S$ で一致し、$S^2$ 上の関数として well-defined。ゲージ不変な量だけが観測量であることの反映である。

(2) 必要なら赤道をわずかに動かして（$g_{NS}$ のホモトピー類は変わらない）、零点が赤道上にないとする。北半球 $D_N$ 内の零点の指数の和を $\Sigma_N$、南半球 $D_S$ 内の和を $\Sigma_S$ とする。$\psi_N$ は $D_N$ から零点のまわりの小円板を除いた領域で零にならないので、偏角の原理（各零点を囲む小円周の巻き付き数の和が外周の巻き付き数に等しい）により、赤道を $\varphi$ が増す向き（$D_N$ の境界の向き）に回ったときの $\psi_N$ の巻き付き数は $\Sigma_N$。同様に $D_S$ の境界の向きは $\varphi$ が減る向きなので、$\varphi$ が増す向きでの $\psi_S$ の巻き付き数は $-\Sigma_S$。赤道上で $\psi_N=e^{in\varphi}\psi_S$ なので巻き付き数を比べて
\[
\Sigma_N=n+(-\Sigma_S),\qquad\text{すなわち}\qquad\Sigma_N+\Sigma_S=n .
\]
$n\neq0$ なら少なくとも1つ零点がある。（零点がなければ $\psi/\lvert\psi\rvert$ が主束の大域切断になり、問題 \ref{ch07-prob-trivsec} とも整合する。）

(3) $\psi^{(1)}$：重なりで $e^{i\varphi}\sin\frac\theta2=e^{i\varphi}\cdot\psi_S$ なので関係式をみたす。$\psi_S=\sin\frac\theta2$ は南極で $1$ に近づき、南極の近くで $\theta'=\pi-\theta$ の偶関数 $\cos\frac{\theta'}2$ なので滑らか。$\psi_N=e^{i\varphi}\sin\frac\theta2\approx\frac\theta2e^{i\varphi}\approx\frac12(x+iy)$ は北極で滑らか。零点は北極のみで、北極のまわりを正の向き（$\varphi$ 増加）に回ると $\psi_N$ の巻き付き数は $+1$。指数の和 $=1=n$。

$\psi^{(2)}$：$\cos\frac\theta2=e^{i\varphi}\cdot e^{-i\varphi}\cos\frac\theta2$ で関係式をみたす。零点は南極のみ。南極のまわりの正の向き（南極での外向き法線に関する向き）は $\varphi$ が減る向きであり、$\psi_S\approx\frac{\theta'}2e^{-i\varphi}$ の巻き付き数はこの向きで $+1$。指数の和 $=1=n$。

なお $\psi^{(1)}$, $\psi^{(2)}$ はスピン $\tfrac12$ のコヒーレント状態の成分 $(\bar z_2,\bar z_1)$ に相当し、モノポール背景の最低ランダウ準位（$n+1=2$ 重縮退）をなす。縮退度が $n+1$ となることは指数定理で理解される（第9章）。
\end{solution}

\begin{problem}[難易度 B]\label{ch07-prob-hopf}
例 \ref{ex:ch07-hopf} の記号を用いる。\\
(1) $z=\bigl(\cos\tfrac\theta2e^{i\alpha},\ \sin\tfrac\theta2e^{i(\alpha+\varphi)}\bigr)$ に対し $h(z)=(\sin\theta\cos\varphi,\sin\theta\sin\varphi,\cos\theta)$ を示せ。\\
(2) $s_N$ が南極以外で、$s_S$ が北極以外で連続な切断であることを示し、$\Phi_N(x,\lambda)=s_N(x)\lambda$ が局所自明化 $U_N\times U(1)\cong h^{-1}(U_N)$ を与えることを示せ。\\
(3) $s_S=s_Ng_{NS}$ となる $g_{NS}$ を求め、ホップ束のモノポール荷が $-1$ であることを示せ。\\
(4) 電荷1の同伴直線束 $S^3\times_{U(1)}\C$（$[z\lambda,c]=[z,\lambda c]$）がトートロジカル直線束 $\mathcal O(-1)=\{(x,v):v\in\C z,\ h(z)=x\}$ と同型であることを示せ。\\
(5) $s_N$ を用いて $A_N=i\,s_N^\dagger ds_N$ を計算し、問題 \ref{ch07-prob-monopole} の $A_N$ で $n=-1$ としたものに一致することを確かめよ（これは量子力学のベリー接続であり、第8章で扱う）。
\end{problem}

\begin{solution}
(1) $\bar z_1z_2=\cos\frac\theta2\sin\frac\theta2\,e^{i\varphi}=\frac12\sin\theta\,e^{i\varphi}$ なので $2\operatorname{Re}(\bar z_1z_2)=\sin\theta\cos\varphi$、$2\operatorname{Im}(\bar z_1z_2)=\sin\theta\sin\varphi$。$\lvert z_1\rvert^2-\lvert z_2\rvert^2=\cos^2\frac\theta2-\sin^2\frac\theta2=\cos\theta$。

(2) $s_N=(\cos\frac\theta2,\sin\frac\theta2e^{i\varphi})$ は $\alpha=0$ の点なので $h(s_N)=(\theta,\varphi)$ の点。第1成分 $\cos\frac\theta2=\sqrt{(1+\cos\theta)/2}$ は $\theta<\pi$ で滑らか、第2成分は $\sin\frac\theta2e^{i\varphi}=\dfrac{x+iy}{2\cos(\theta/2)}$ と書け（$\sin\theta=2\sin\frac\theta2\cos\frac\theta2$）、$\theta<\pi$ で連続。$\theta\to\pi$ では第2成分が $e^{i\varphi}$ に近づき方向に依存するので連続に延長できない。$s_S$ も同様に北極以外で連続。

$\Phi_N(x,\lambda)=s_N(x)\lambda$ は連続で $h\circ\Phi_N(x,\lambda)=x$。逆写像は $z\mapsto\bigl(h(z),\ z_1/\lvert z_1\rvert\bigr)$：実際 $h(z)\in U_N$ は $z_1\neq0$ と同値で、$z=s_N(h(z))\cdot(z_1/\lvert z_1\rvert)$（第1成分の偏角が $\lambda$、ファイバーは $U(1)$ 軌道）。よって局所自明化である。

(3) $s_Ne^{-i\varphi}=(\cos\frac\theta2e^{-i\varphi},\sin\frac\theta2)=s_S$。よって $g_{NS}=e^{-i\varphi}$。赤道を $\varphi$ 増加の向きに回ると巻き付き数は $-1$ なので、定義（例 \ref{ex:ch07-monopole}）によりモノポール荷は $-1$。すなわちホップ束は $P_{-1}$ であり、特に自明でない（全空間 $S^3$ は $S^2\times S^1$ と同相でない：$\pi_1(S^3)=0\neq\Z$ からもわかる）。

(4) $F\colon S^3\times_{U(1)}\C\to\mathcal O(-1)$、$[z,c]\mapsto(h(z),\,c\,z)$ と定める。$[z\lambda,c]\mapsto(h(z),c\lambda z)$、$[z,\lambda c]\mapsto(h(z),\lambda c z)$ で一致するので well-defined。各ファイバーで線形同型（$c\mapsto cz$ は $\C\to\C z$ の同型）で、局所的には連続なので束の同型。

(5) $s_N^\dagger ds_N=\cos\frac\theta2\,d\cos\frac\theta2+\sin\frac\theta2e^{-i\varphi}\,d\bigl(\sin\frac\theta2e^{i\varphi}\bigr)=\tfrac12d\bigl(\cos^2\tfrac\theta2+\sin^2\tfrac\theta2\bigr)+i\sin^2\tfrac\theta2\,d\varphi=i\sin^2\tfrac\theta2\,d\varphi$。よって
\[
A_N=i\cdot i\sin^2\tfrac\theta2\,d\varphi=-\tfrac12(1-\cos\theta)\,d\varphi,
\]
これは $n=-1$ の $A_N$ である。$s_S$ についても同様に $A_S=i\,s_S^\dagger ds_S=i\cdot(-i\cos^2\frac\theta2)d\varphi=\frac12(1+\cos\theta)d\varphi$（$n=-1$ の $A_S$）で、$A_N-A_S=-d\varphi=i\,g_{NS}\,dg_{NS}^{-1}$ と整合する。
\end{solution}

\begin{problem}[難易度 B]\label{ch07-prob-adjhiggs}
$SU(2)$ ジョージ・グラショウ模型の 't Hooft--Polyakov モノポールでは、遠方の球面 $S^2_\infty$ 上でヒッグス場が $\Phi=v\,\hat\Phi\cdot\tfrac{\bm\sigma}2$、$\hat\Phi(\hat{\bm x})=\hat{\bm x}$（ヘッジホッグ）となる。$S^2_\infty$ 上の $SU(2)$ 束は自明（問題 \ref{ch07-prob-s2bundles}）なので、基本表現の束は $S^2\times\C^2$ とみなせる。\\
(1) 各点 $\hat{\bm x}$ で $\hat{\bm x}\cdot\bm\sigma$ の固有値 $\pm1$ の固有空間を $L_\pm$ とすると、$S^2\times\C^2=L_+\oplus L_-$ と直線束に分解することを示せ。\\
(2) $L_+$ がホップ束の電荷1の同伴束（$\mathcal O(-1)$）と同型であることを示し、$L_+$, $L_-$ のモノポール荷（推移関数の巻き付き数）を求めよ。\\
(3) 非破れの $U(1)$ は $T=\hat\Phi\cdot\frac{\bm\sigma}2$ で生成される。随伴表現 $\C^2\otimes(\C^2)^*$ から $T$ の電荷 $q=0,\pm1$ の成分（$W^\pm$ ボソンなど）を取り出し、それぞれの直線束のモノポール荷が $-2q$ であることを示せ。2重項の成分（$q=\pm\frac12$）についても同じ関係が成り立つことを確かめよ。
\end{problem}

\begin{solution}
(1) $\hat{\bm x}\cdot\bm\sigma$ はエルミートで $(\hat{\bm x}\cdot\bm\sigma)^2=1$、跡が $0$ なので、固有値 $\pm1$ をそれぞれ1つずつもつ。射影 $P_\pm=\frac12(1\pm\hat{\bm x}\cdot\bm\sigma)$ は $\hat{\bm x}$ について連続でランク1なので、$L_\pm=\operatorname{Im}P_\pm$ は直線束（局所的に、$P_\pm$ をある固定ベクトルに作用させたものが枠を与える）。$P_++P_-=1$、$P_+P_-=0$ より $\C^2=L_+\oplus L_-$（直交分解）。

(2) $h(z)=z^\dagger\bm\sigma z$ に対し $(h(z)\cdot\bm\sigma)z=z$ を示す。$\sum_a\sigma^a_{ij}\sigma^a_{kl}=2\delta_{il}\delta_{jk}-\delta_{ij}\delta_{kl}$ を用いると
\[
\bigl((h\cdot\bm\sigma)z\bigr)_i=\sum_a\bar z_k\sigma^a_{kl}z_l\,\sigma^a_{ij}z_j=2z_i\,\bar z_kz_k-z_i\,\bar z_kz_k=z_i\quad(\lvert z\rvert=1).
\]
よって $L_+$ の $\hat{\bm x}$ 上のファイバーは $\C z$（$h(z)=\hat{\bm x}$）であり、$L_+=\mathcal O(-1)$。問題 \ref{ch07-prob-hopf} の局所枠 $s_N$, $s_S$ を用いると $s_S=s_Ne^{-i\varphi}$ なので、成分の関係は $\psi_N=e^{-i\varphi}\psi_S$、モノポール荷は $-1$。

$L_-=L_+^\perp$ の局所枠として $s_N^\perp=(-\sin\frac\theta2e^{-i\varphi},\cos\frac\theta2)$（北極で滑らか）、$s_S^\perp=(-\sin\frac\theta2,\cos\frac\theta2e^{i\varphi})$（南極で滑らか）をとると $s_S^\perp=s_N^\perp e^{i\varphi}$。よって $L_-$ のモノポール荷は $+1$。これは $L_+\otimes L_-\cong\Lambda^2\C^2$（自明）とも整合する（テンソル積の推移関数は積なので巻き付き数は和、$-1+1=0$）。

(3) $\C^2=L_+\oplus L_-$ で、$T$ は $L_\pm$ に $\pm\frac12$ で作用するので、2重項成分の電荷は $q=\pm\frac12$、モノポール荷は $\mp1=-2q$。随伴表現（$\mathfrak{sl}(2,\C)\cong\End_0(\C^2)$、$T$ は交換子で作用）は
\[
\End(\C^2)=\C^2\otimes(\C^2)^*=(L_+\otimes L_+^*)\oplus(L_-\otimes L_-^*)\oplus(L_+\otimes L_-^*)\oplus(L_-\otimes L_+^*)
\]
と分解し、跡のない部分は $\C\,T\oplus(L_+\otimes L_-^*)\oplus(L_-\otimes L_+^*)$。$L_+\otimes L_-^*$ への $[T,\cdot\,]$ の固有値は $\frac12-(-\frac12)=1$（$W^+$）、$L_-\otimes L_+^*$ は $-1$（$W^-$）、$T$ 方向は $0$（光子と、ヒッグスの動径方向）。双対の推移関数は逆数なので巻き付き数は符号が反転し、
\[
L_+\otimes L_-^*:\ -1-(+1)=-2,\qquad L_-\otimes L_+^*:\ +2,\qquad L_\pm\otimes L_\pm^*:\ 0 .
\]
よっていずれも「モノポール荷 $=-2q$」が成り立つ。すなわち対称性の破れ $SU(2)\to U(1)$ によって構造群が $U(1)$ に\textbf{縮小}し、ヘッジホッグの位相（写像 $S^2\to S^2$ の次数 $1$）が非破れ $U(1)$ 束の磁荷になる。$q=\frac12$ の場が最小のディラック磁荷を感じるので、't Hooft--Polyakov モノポールの磁荷は、$W$ ボソンの電荷を単位としたディラック量子化の最小単位の2倍にあたる（第3, 8章）。次数 $m$ の $\hat\Phi$ では $L_+=\hat\Phi^*\mathcal O(-1)$ となり、推移関数が $g_{NS}\circ\hat\Phi$ になるのでモノポール荷は $-m$ 倍になる。
\end{solution}

\begin{problem}[難易度 B]\label{ch07-prob-ts2}
$S^2$ 上の複素座標 $w=\tan\frac\theta2\,e^{i\varphi}$（北極の近傍の地図、$S^2$ の向きと整合）と $w'=1/w$（南極の近傍）を用いる。$S^2$ の接ベクトルを複素成分 $v^w$（$v=v^w\partial_w+\overline{v^w}\partial_{\bar w}$）で表し、$TS^2$ を複素直線束とみなす。\\
(1) 重なりで $v^w=-w^2\,v^{w'}$ を示し、赤道上での推移関数 $g_{NS}$ の巻き付き数が $2$ であることを示せ。\\
(2) $TS^2$ は自明でないことを結論し、「$S^2$ 上の連続な接ベクトル場は必ず零点をもつ」（毛玉の定理）を導け。地球上の風速場について何がいえるか。\\
(3) $\mathcal O(-1)$ の推移関数が $w^{-1}$ であることを用いて $TS^2\cong\mathcal O(2)$ を示せ。
\end{problem}

\begin{solution}
(1) 正則な座標変換 $w'=1/w$ に対し、ベクトルの成分は $v^{w'}=\dfrac{\partial w'}{\partial w}v^w=-w^{-2}v^w$ と変換する（$\bar w$ 成分はその複素共役）。よって $v^w=-w^2v^{w'}$、すなわち $\psi_N=v^w$、$\psi_S=v^{w'}$ として $g_{NS}=-w^2$。赤道 $\theta=\frac\pi2$ では $w=e^{i\varphi}$ で $g_{NS}=-e^{2i\varphi}$、$\varphi$ が $0\to2\pi$ と動くとき偏角は $4\pi$ 増加するので巻き付き数 $2$（構造群は $GL(1,\C)\simeq U(1)$ だが、$g/\lvert g\rvert$ を考えればよい）。

(2) 巻き付き数 $2\neq0$ なので、問題 \ref{ch07-prob-monopole} (5) と同じ議論で $TS^2$ は自明束と同型でない（自明束は $g_{NS}=1$、巻き付き数 $0$）。零点のない接ベクトル場 $v$ があれば、複素直線束 $TS^2$ の至る所零でない切断であり、問題 \ref{ch07-prob-trivsec} (2) により $TS^2$ は自明になってしまう。よって任意の接ベクトル場は零点をもつ。地表の水平風速場は $S^2$ 上の接ベクトル場なので、どの瞬間にも少なくとも1か所、水平風速が $0$ の点（台風の目のような点）が存在する。より詳しく、零点の指数の和は $2=\chi(S^2)$ である（問題 \ref{ch07-prob-section} (2) と同じ議論；ポアンカレ・ホップの定理、第8章のガウス・ボンネ）。

(3) 例の通り $\mathcal O(-1)$ は $g_{NS}=w^{-1}$（問題 \ref{ch07-prob-hopf} で赤道上 $e^{-i\varphi}$）。$\mathcal O(2)=\mathcal O(-1)^{*\otimes2}$ の推移関数は $(w^{-1})^{-2}=w^2$。$TS^2$ の推移関数 $-w^2$ とは定数 $-1$ だけ異なるが、$h_N=-1$ で取り換えれば一致するので $TS^2\cong\mathcal O(2)$。巻き付き数は $c_1(TS^2)=2=\chi(S^2)$（第8章）。
\end{solution}

\begin{problem}[難易度 A]\label{ch07-prob-parallel}
接束が自明な多様体を\term{平行化可能}{parallelizable}という。\\
(1) $S^1\subset\C$ と $S^3\subset\HH$（単位四元数）について、$v(x)=ix$（$S^1$）、$v_1(x)=xi$, $v_2(x)=xj$, $v_3(x)=xk$（$S^3$）が各点で接空間の正規直交基底を与えることを示し、$S^1$, $S^3$ が平行化可能であることを示せ。\\
(2) 任意のリー群 $G$ は平行化可能であることを示せ。\\
(3) $S^n$ が平行化可能となるのは $n=0,1,3,7$ に限る（$S^7$ は八元数による；アダムスの定理）。$n$ が偶数の場合に平行化不可能であることを、毛玉の定理の一般化（$S^{2m}$ 上の接ベクトル場は零点をもつ）から説明せよ。
\end{problem}

\begin{solution}
(1) $\HH=\R^4$ の内積は $\langle a,b\rangle=\operatorname{Re}(a\bar b)$ で、$\lvert ab\rvert=\lvert a\rvert\lvert b\rvert$。$x\in S^3$ の接空間は $\{v:\langle x,v\rangle=0\}$。$u$ が純虚四元数（$\bar u=-u$）なら
\[
\langle x,xu\rangle=\operatorname{Re}(x\bar u\bar x)=-\operatorname{Re}(xu\bar x)=0,
\]
なぜなら $xu\bar x=xux^{-1}$ は純虚四元数 $u$ の共役で、実部が $0$ だから。よって $xi,xj,xk$ は接ベクトル。さらに $\langle xu,xu'\rangle=\operatorname{Re}(xu\bar u'\bar x)=\operatorname{Re}(u\bar u')\lvert x\rvert^2=\langle u,u'\rangle$（実部は共役で不変）なので、$i,j,k$ の正規直交性が保たれ、$xi,xj,xk$ は正規直交基底。これらは $x$ について滑らかなので、$TS^3\cong S^3\times\R^3$。$S^1$ では $\langle x,ix\rangle=\operatorname{Re}(x\,\overline{ix})=\operatorname{Re}(-i\lvert x\rvert^2)=0$、$\lvert ix\rvert=1$ で同様。

(2) リー代数 $\mathfrak g=T_1G$ の基底 $X_1,\dots,X_n$ をとり、左移動 $L_g(h)=gh$ の微分で $X_a(g)=(dL_g)_1X_a$ と定める（左不変ベクトル場）。$dL_g$ は線形同型なので、各点で基底になり、$TG\cong G\times\mathfrak g$。行列群なら $X_a(g)=gX_a$ で、(1) はこの特別な場合（$S^3=SU(2)$、$S^1=U(1)$）である。物理では、$SU(2)$ 上の左不変形式 $g^{-1}dg$（マウラー・カルタン形式）が至る所で使われる。

(3) $S^{2m}$ が平行化可能なら、その枠の1本目は零点のない接ベクトル場であり、一般化された毛玉の定理に反する。（$S^{2m}$ の接ベクトル場の零点の指数の和は $\chi(S^{2m})=2\neq0$。）奇数次元球面 $S^{2m-1}\subset\C^m$ には零点のないベクトル場 $v(x)=ix$ が常にあるが、$2m-1$ 本の独立な場が存在するのは $2m-1=1,3,7$ のときだけである。
\end{solution}

\begin{problem}[難易度 C]\label{ch07-prob-frame}
(1) $S^2$ の向きの付いた正規直交枠束 $F(S^2)=\{(x,e_1,e_2)\}$（$x\in S^2$, $e_1,e_2\in T_xS^2$ 正規直交、$e_1\times e_2=x$）が $SO(3)$ と同相であり、射影が $R\mapsto Re_3$、構造群 $SO(2)$ の作用が右からのかけ算 $R\mapsto R\begin{pmatrix}r(\alpha)&0\\0&1\end{pmatrix}$ で与えられることを示せ（$r(\alpha)$ は角度 $\alpha$ の2次元回転）。\\
(2) ファイバー束 $SO(2)\to SO(3)\to S^2$ のホモトピー完全列と $\pi_1(SO(3))=\Z_2$（第2章）、$\pi_2(SO(3))=0$ を用いて、境界写像 $\partial\colon\pi_2(S^2)\to\pi_1(SO(2))$ が $\pm2$ 倍であることを示し、問題 \ref{ch07-prob-ts2} の巻き付き数 $2$ と比べよ。\\
(3) 主 $U(1)$ 束 $P_n\to S^2$（モノポール荷 $n$）について $\partial$ が $\pm n$ 倍であることを認めて、$n\neq0$ のとき $\pi_1(P_n)\cong\Z_{\lvert n\rvert}$、$n=0$ のとき $\pi_1(P_0)\cong\Z$ を示せ。ホップ束 $P_{-1}=S^3$、$F(S^2)\cong SO(3)$ と整合することを確かめよ。
\end{problem}

\begin{solution}
(1) $(x,e_1,e_2)$ に列ベクトルを並べた行列 $R=(e_1\ e_2\ x)$ を対応させる。列が正規直交なので $R\in O(3)$、$\det R=(e_1\times e_2)\cdot x=1$ なので $R\in SO(3)$。逆に $R\in SO(3)$ の列 $(e_1,e_2,x)$ は $x\in S^2$、$e_1,e_2\perp x$ で $e_1\times e_2=x$（$\det=1$ の正規直交基底）。この対応は連続な全単射で逆も連続（コンパクト・ハウスドルフ）なので $F(S^2)\cong SO(3)$。射影は $R\mapsto x=Re_3$。枠の回転 $(e_1,e_2)\mapsto(\cos\alpha\,e_1+\sin\alpha\,e_2,\ -\sin\alpha\,e_1+\cos\alpha\,e_2)$ は $R\mapsto R\operatorname{diag}(r(\alpha),1)$ であり、$SO(2)$ はファイバー（$x$ を固定する枠全体）上で自由かつ推移的に作用する。よって $SO(3)\to SO(3)/SO(2)=S^2$ は $S^2$ の枠束そのもので、$F(S^2)\cong SO(3)\cong\RP^3$（第1章）。$e_1$ が $e_2$ を決めるので、単位接ベクトル束とも同じである。

(2) 長完全列の一部
\[
\pi_2(SO(3))\to\pi_2(S^2)\xrightarrow{\partial}\pi_1(SO(2))\to\pi_1(SO(3))\to\pi_1(S^2)
\]
は $0\to\Z\xrightarrow{\partial}\Z\to\Z_2\to0$ となる（$\pi_2(SO(3))\cong\pi_2(S^3)=0$ は2重被覆 $S^3\to SO(3)$ による）。よって $\partial$ は単射で、その余核が $\Z_2$。$\Z\to\Z$ の単射準同型は $m$ 倍（$m\neq0$）で余核は $\Z_{\lvert m\rvert}$ なので $m=\pm2$。一方 $F(S^2)$ は $TS^2$（$U(1)=SO(2)$ 構造の複素直線束）の単位円束なので、推移関数は $TS^2$ と同じく巻き付き数 $2$ をもつ。$\partial$ の値が貼り合わせ写像の巻き付き数（符号を除く）と一致することが確かめられた。

(3) 長完全列
\[
\pi_2(S^1)=0\to\pi_2(P_n)\to\pi_2(S^2)\xrightarrow{\partial=\pm n}\pi_1(S^1)\to\pi_1(P_n)\to\pi_1(S^2)=0
\]
より $\pi_1(P_n)\cong\coker(\pm n\colon\Z\to\Z)=\Z/n\Z$。$n\neq0$ なら $\Z_{\lvert n\rvert}$（このとき $\partial$ は単射なので $\pi_2(P_n)=0$）。$n=0$ なら $\Z$ で、実際 $P_0=S^2\times S^1$。ホップ束 $n=-1$：$\pi_1=0$ で $S^3$ と整合。$F(S^2)$ は $n=\pm2$（向きの取り方による）：$\pi_1=\Z_2$ で $SO(3)\cong\RP^3$ と整合。一般に $P_{-m}\cong S^3/\Z_m$（$\Z_m\subset U(1)$ をホップのファイバーに沿って作用させたレンズ空間）であり、$U(1)/\Z_m\cong U(1)$（$\lambda\mapsto\lambda^m$）によって推移関数 $e^{-i\varphi}$ が $e^{-im\varphi}$ になることからも確かめられる。
\end{solution}

\begin{problem}[難易度 C]\label{ch07-prob-s2bundles}
定理 \ref{thm:ch07-clutch} を用いて $S^2$ 上の主束を分類する。$\pi_1(SU(N))=0$（$N\ge2$）、$\pi_2(G)=0$（コンパクト・リー群）は認めてよい。\\
(1) $S^2$ と $S^3$ 上の主 $SU(N)$ 束はすべて自明であることを示せ。\\
(2) $\det\colon U(N)\to U(1)$ が同型 $\pi_1(U(N))\cong\Z$ を誘導することを示し、$S^2$ 上の $U(N)$ 束は整数 $\deg(\det g_{NS})$ で分類されることを示せ。\\
(3) $U(N)$ ゲージ理論の GNO 型モノポールで $g_{NS}=\operatorname{diag}(e^{in_1\varphi},\dots,e^{in_N\varphi})$ とする。束の位相型は $\sum_kn_k$ だけで決まることを示し、$N=2$、$(n_1,n_2)=(1,-1)$ の場合に $g_{NS}$ の自明化（定値写像へのホモトピー）を具体的に作れ。\\
(4) $S^2$ 上の $SO(3)$ 束は2種類あることを示せ。$U(1)$ の電荷 $n$ のモノポール束を $U(1)\cong SO(2)\subset SO(3)$ で $SO(3)$ 束とみたとき、その類は $n\bmod2$ であることを示し、$n$ が奇数のとき $SU(2)$ 束に持ち上がらない（2重項を含む理論では許されない）ことを説明せよ。
\end{problem}

\begin{solution}
(1) $S^2$ 上：$\pi_1(SU(N))=0$ より束は1種類、すなわち自明束のみ。$S^3$ 上：$\pi_2(SU(N))=0$ より同様。物理的には、$SU(N)$ ゲージ理論では3次元空間の遠方球面 $S^2_\infty$ 上の束に位相的な制約がなく、磁荷はゲージ群が $U(1)$ に破れて初めて現れる（問題 \ref{ch07-prob-adjhiggs}）。

(2) $SU(N)\to U(N)\xrightarrow{\det}U(1)$ はファイバー束（$U(N)\cong SU(N)\times U(1)$、第1章）。完全列 $\pi_1(SU(N))=0\to\pi_1(U(N))\xrightarrow{\det_*}\pi_1(U(1))\to\pi_0(SU(N))=0$ より $\det_*$ は同型。よって $S^2$ 上の $U(N)$ 束は $[g_{NS}]\in\pi_1(U(N))\cong\Z$、すなわち $\deg(\det g_{NS})$ で分類される。これは行列式束 $\det E$ のモノポール荷である。

(3) $\det g_{NS}=e^{i(\sum_kn_k)\varphi}$ なので $\deg=\sum_kn_k$。(2) より束の型はこれだけで決まり、例えば $(n_1,\dots,n_N)$ は $(\sum n_k,0,\dots,0)$ と同じ束を与える。個々の $n_k$（GNO 電荷）は位相的には保存されず、安定性は力学で決まる。

$N=2$、$g(\varphi)=\operatorname{diag}(e^{i\varphi},e^{-i\varphi})=\cos\varphi+i\sin\varphi\,\sigma_3\in SU(2)$。$SU(2)\cong S^3$ の点 $a_0+i\bm a\cdot\bm\sigma$ を $(a_0,\bm a)\in\R^4$ と同一視すると $g$ は大円 $(\cos\varphi,0,0,\sin\varphi)$。
\[
H(\varphi,t)=\frac{(1-t)(\cos\varphi,0,0,\sin\varphi)+t\,(0,1,0,0)}{\bigl\lvert(1-t)(\cos\varphi,0,0,\sin\varphi)+t\,(0,1,0,0)\bigr\rvert}
\]
とおくと、分子の第2成分は $t$、他の成分の2乗和は $(1-t)^2$ なので分母は $\sqrt{(1-t)^2+t^2}>0$。$H$ は $S^1\times I\to S^3$ の連続写像で、$H(\cdot,0)=g$、$H(\cdot,1)=i\sigma_1$（定値）。よって $g\simeq$ 定値で、束は自明。

(4) $\pi_1(SO(3))=\Z_2$（第2章）なので2種類。$e^{i\alpha}\mapsto R_3(\alpha)$（第3軸まわりの回転）で $g_{NS}=e^{in\varphi}$ は $R_3(n\varphi)$ になる。$R_3(\varphi)$、$\varphi\in[0,2\pi]$ は $\pi_1(SO(3))$ の生成元である：$SU(2)$ への持ち上げ $\tilde R(\varphi)=e^{-i\varphi\sigma_3/2}$ は $\tilde R(2\pi)=-1$ に終わる閉じない道だからである（第2章）。よって $R_3(n\varphi)$ の類は $n\bmod2$。$n$ が奇数のとき、$g_{NS}$ を $SU(2)$ に連続に持ち上げると赤道を一周して $-1$ が掛かり、1価な $\tilde g_{NS}\colon S^1\to SU(2)$ は作れない。すなわちこの $SO(3)$ 束は $SU(2)$ 束に持ち上がらず、2重項の場（$-1\in SU(2)$ が $-1$ で作用）はこの背景では定義できない。随伴表現の場だけの理論では許され、そのモノポールの磁荷は2重項が存在する場合に許される最小値の半分になる。持ち上げの障害は $w_2\in H^2(S^2;\Z_2)=\Z_2$ であり、スピン構造の障害と同じ構造である。
\end{solution}

\begin{problem}[難易度 C]\label{ch07-prob-s4}
例 \ref{ex:ch07-s4} の状況を考える。$x=(\bm x,x_4)\in\R^4$、$\R^4$ の向きは $dx^1\wedge dx^2\wedge dx^3\wedge dx^4$ とする。\\
(1) $g_1(x)=x_4+i\,\bm x\cdot\bm\sigma$（$\lvert x\rvert=1$）について、点 $p=(0,0,0,1)$ で $\tr(g_1^{-1}dg_1)^3=12\,dx^1\wedge dx^2\wedge dx^3$ を示せ。さらに、この3形式が左不変であることを用いて、$S^3$ を $p$ で $dx^1\wedge dx^2\wedge dx^3$ が正となる向きにすれば $w(g_1)=1$ であることを示せ（$\mathrm{Vol}(S^3)=2\pi^2$）。\\
(2) 球 $B=\{\lvert x\rvert\le R\}$ の境界としての $S^3$ の向き（外向き法線を先頭に置く）は、(1) の向きと逆であることを示せ。\\
(3) $U_S$ 上でゲージ場が $A_S=0$（赤道付近）であり、$U_N$ 上のゲージ場が赤道付近で $A_N=i\,g\,dg^{-1}$、$g=g_{NS}$ であるとする。記法と規約の $CS(A)$ を用いて、インスタントン数 $k=\frac1{8\pi^2}\int\tr F\wedge F$ が $k=w(g_{NS})$ となることを示せ。\\
(4) $w(gh)=w(g)+w(h)$ を認めて、任意の $m\in\Z$ に対しインスタントン数 $m$ の $SU(2)$ 束が存在することを示せ。
\end{problem}

\begin{solution}
(1) $p$ で $g_1=1$、また $T_pS^3$ 上で $dx^4=0$ なので $g_1^{-1}dg_1=i\,\sigma_a\,dx^a$（$a=1,2,3$）。
\[
(g_1^{-1}dg_1)^3=i^3\sigma_a\sigma_b\sigma_c\,dx^a\wedge dx^b\wedge dx^c=-i\,\epsilon_{abc}\sigma_a\sigma_b\sigma_c\,dx^1\wedge dx^2\wedge dx^3 .
\]
\begin{note}
$\sigma_a\sigma_b=\delta_{ab}+i\epsilon_{abd}\sigma_d$ より $\epsilon_{abc}\sigma_a\sigma_b=i\epsilon_{abc}\epsilon_{abd}\sigma_d=2i\sigma_c$、よって $\epsilon_{abc}\sigma_a\sigma_b\sigma_c=2i\sigma_c\sigma_c=6i\cdot\mathbf 1$、跡は $12i$。
\end{note}
したがって $\tr(g_1^{-1}dg_1)^3=-i\cdot12i\,dx^{123}=12\,dx^1\wedge dx^2\wedge dx^3$。

$g\mapsto g_0g$（$g_0\in SU(2)$）で $g^{-1}dg$ は不変なので、$\omega=\tr(g^{-1}dg)^3$ は $S^3=SU(2)$ 上の左不変3形式である。左移動 $x\mapsto g_0x$ は $\R^4$ の線形写像で $\det(g_0g)=\det g$ よりノルム $\lvert x\rvert^2=\det(x_4+i\bm x\cdot\bm\sigma)$ を保ち、恒等写像と連結なので $SO(4)$ の元、すなわち $S^3$ の向きを保つ等長写像である。よって体積形式 $\mathrm{vol}$（$p$ で $dx^{123}$）も左不変で、$\omega$ と $\mathrm{vol}$ は $p$ で比 $12$ なので至る所 $\omega=12\,\mathrm{vol}$。
\[
w(g_1)=\frac{1}{24\pi^2}\cdot12\cdot2\pi^2=1 .
\]
これは記法と規約の「恒等写像 $S^3\cong SU(2)$ が $w=1$」という向きの約束を、座標で表したものである。

(2) $p$ での外向き法線は $\partial_4$。境界の向きでは $(v_1,v_2,v_3)$ が正 $\iff$ $(\partial_4,v_1,v_2,v_3)$ が $\R^4$ で正。$dx^{1234}(\partial_4,\partial_1,\partial_2,\partial_3)$ は置換 $(4,1,2,3)$ の符号（互換3回）で $-1$。よって境界の向きは $p$ で $-dx^1\wedge dx^2\wedge dx^3$、(1) の向きと逆である。

(3) 平坦 $F=dA-iA\wedge A=0$ のとき $dA=iA\wedge A$ なので
\[
\tr\Bigl(A\wedge dA-\frac{2i}3A^3\Bigr)=\Bigl(i-\frac{2i}3\Bigr)\tr A^3=\frac i3\tr A^3 .
\]
$A=i\,g\,dg^{-1}=-i\,dg\,g^{-1}$ なので $A^3=(-i)^3(dg\,g^{-1})^3=i\,(dg\,g^{-1})^3$、跡の巡回性より $\tr(dg\,g^{-1})^3=\tr(g^{-1}dg)^3$。よって
\[
CS(A_N)=\frac1{8\pi^2}\cdot\frac i3\cdot i\,\tr(g^{-1}dg)^3=-\frac{1}{24\pi^2}\tr(g^{-1}dg)^3 .
\]
$U_S$ 側（球の外）では $A_S=0$ として $F=0$ なので寄与がなく、
\[
k=\frac1{8\pi^2}\int_B\tr F\wedge F=\int_BdCS(A_N)=\int_{\partial B}CS(A_N)=-\frac1{24\pi^2}\int_{\partial B}\tr(g^{-1}dg)^3 .
\]
(2) により $\partial B$ の向きは (1) の向きの逆なので、$\int_{\partial B}=-\int_{S^3}$ となり $k=w(g_{NS})$。特に $g_{NS}=g_1$ の束はインスタントン数 $k=1$ をもつ。（$A_S=0$ でない一般の場合も、$CS(A_N)-CS(A_S)$ が $-\frac1{24\pi^2}\tr(g^{-1}dg)^3$ と完全形式の差しかないことから同じ結論になる。第8章。）

(4) $g_1^m$（$m<0$ では $(g_1^{-1})^{\lvert m\rvert}$）は $w(g_1^m)=mw(g_1)=m$（$w(g^{-1})=-w(g)$ は $w(gg^{-1})=w(1)=0$ から）。$g_{NS}=g_1^m$ で貼り合わせた束はインスタントン数 $m$ をもつ。定理 \ref{thm:ch07-clutch} と $\pi_3(SU(2))=\Z$ により、これで $S^4$ 上の $SU(2)$ 束はすべて尽くされる。ユークリッド時空 $\R^4\cup\{\infty\}$ の作用有限な配位はこれらの束のどれかの上の接続であり、インスタントン数の異なる配位は連続変形で移り合えない。
\end{solution}

\begin{problem}[難易度 B]\label{ch07-prob-les}
定理 \ref{thm:ch07-les} を用いる。$i<k$ で $\pi_i(S^k)=0$、$\pi_k(S^k)=\Z$、$k\ge2$ で $\pi_k(S^1)=0$ は認めてよい。\\
(1) ホップ束 $S^1\to S^3\to S^2$ から $\pi_2(S^2)\cong\Z$ と、$n\ge3$ で $\pi_n(S^3)\cong\pi_n(S^2)$ を示せ。特に $\pi_3(S^2)\cong\Z$ の生成元は何か。\\
(2) (1) の結果を、$O(3)$ 非線形シグマ模型（標的 $S^2$）の位相的ソリトンに関連づけて述べよ（2次元空間のベラヴィン・ポリャコフ・インスタントンと、3次元空間のホップフィオン）。\\
(3) $SU(N-1)\to SU(N)\to S^{2N-1}$（$SU(N)/SU(N-1)\cong S^{2N-1}$、第1章）を用いて、$N\ge2$ で $\pi_3(SU(N))\cong\Z$ を示せ（$SU(2)\cong S^3$ から始めよ）。また、$SU(2)\subset SU(N)$ の左上への埋め込みが $\pi_3$ の同型を誘導することを示せ。
\end{problem}

\begin{solution}
(1) 完全列
\[
\pi_n(S^1)\to\pi_n(S^3)\xrightarrow{h_*}\pi_n(S^2)\xrightarrow{\partial}\pi_{n-1}(S^1)\to\pi_{n-1}(S^3)
\]
で $n=2$ とすると $0=\pi_2(S^3)\to\pi_2(S^2)\to\pi_1(S^1)=\Z\to\pi_1(S^3)=0$ なので $\pi_2(S^2)\cong\Z$。$n\ge3$ では $\pi_n(S^1)=0=\pi_{n-1}(S^1)$ なので $h_*$ は同型。$n=3$ で $\pi_3(S^2)\cong\pi_3(S^3)=\Z$ であり、生成元 $[\id_{S^3}]$ の像 $h_*[\id]=[h]$、すなわちホップ写像自身が $\pi_3(S^2)$ の生成元である。

(2) 2次元ユークリッド空間（$\R^2\cup\{\infty\}=S^2$）上の $O(3)$ 模型の作用有限な配位 $\bm n\colon S^2\to S^2$ は $\pi_2(S^2)=\Z$ の整数（位相的電荷、スカーミオン数の2次元版）で分類され、ベラヴィン・ポリャコフ・インスタントンはその最小作用配位である。3次元空間（$\R^3\cup\{\infty\}=S^3$）の静的な配位 $\bm n\colon S^3\to S^2$ は $\pi_3(S^2)=\Z$ のホップ不変量で分類され、ファデーエフ模型ではこれが結び目状のソリトン（ホップフィオン）を安定化させる。ホップ不変量は2つの点の逆像（閉曲線）の絡み数として理解できる（第3, 6章）。

(3) 完全列
\[
\pi_4(S^{2N-1})\to\pi_3(SU(N-1))\xrightarrow{i_*}\pi_3(SU(N))\to\pi_3(S^{2N-1})
\]
で $N\ge3$ なら $2N-1\ge5>4$ なので両端が $0$ で、$i_*$ は同型。$\pi_3(SU(2))=\pi_3(S^3)=\Z$ から帰納的に $\pi_3(SU(N))\cong\Z$、しかも同型は包含 $SU(2)\subset SU(3)\subset\cdots\subset SU(N)$（左上に埋め込む）で誘導される。したがって $SU(N)$ インスタントンも整数 $k$ で分類され、$k=1$ の配位は $SU(2)$ の BPST インスタントンを $SU(N)$ に埋め込んで得られる（第8章）。
\end{solution}

\subsection{ねじれた境界条件}

\begin{problem}[難易度 D]\label{ch07-prob-thooft}
2次元トーラス $T^2=\R^2/(L\Z)^2$ 上の $SU(N)$ ゲージ理論を、$\R^2$ 上の場に「ゲージ変換を除いて周期的」という境界条件を課して記述する：$\mu=1,2$ について、$\Omega_\mu\colon\R^2\to SU(N)$ が与えられ
\[
A(x+Le_\mu)=\Omega_\mu(x)\ast A(x),\qquad \phi(x+Le_\mu)=\Omega_\mu(x)\,\phi(x)\,\Omega_\mu(x)^{-1}
\]
（$\ast$ は問題 \ref{ch07-prob-cocycle} の記号、$\phi$ は随伴表現の場）。
\begin{center}
\begin{tikzpicture}[scale=1]
  \draw[thick,fill=blue!5] (0,0) rectangle (2.4,2.4);
  \node[left] at (0,1.2) {\footnotesize $x$};
  \node[right] at (2.4,1.2) {\footnotesize $x+Le_1$：$\Omega_1$};
  \node[below] at (1.2,0) {\footnotesize $x$};
  \node[above] at (1.2,2.4) {\footnotesize $x+Le_2$：$\Omega_2$};
  \node[vtx] at (2.4,2.4) {};
  \node[above right] at (2.4,2.4) {\footnotesize $x+L(e_1+e_2)$};
\end{tikzpicture}
\zucap{基本領域。右の辺と上の辺の値は、左・下の辺の値から $\Omega_1$, $\Omega_2$ で決まる。右上の角には2通りの経路で到達する。}
\end{center}
(1) 点 $x+L(e_1+e_2)$ の値を2通りに計算し、随伴表現の場が任意の値をとりうるための条件が
\[
\Omega_1(x+Le_2)\,\Omega_2(x)=z\,\Omega_2(x+Le_1)\,\Omega_1(x),\qquad z=e^{2\pi im/N}\in\Z_N
\]
（$m\in\Z_N$ は $x$ によらない）であることを示せ。\\
(2) $\R^2$ 上のゲージ変換 $h(x)$ で場を $A\mapsto h\ast A$ と変換すると、境界条件は $\Omega_\mu'(x)=h(x+Le_\mu)\Omega_\mu(x)h(x)^{-1}$ に変わる。$m$ がこの下で不変であることを示せ。\\
(3) $N\times N$ の時計行列 $P=\operatorname{diag}(1,\omega,\dots,\omega^{N-1})$（$\omega=e^{2\pi i/N}$）とシフト行列 $Q$（$Qe_k=e_{k+1}$、添字は $\bmod N$）について $PQ=\omega QP$ を示し、定数の $\Omega_1$, $\Omega_2$ で $m=1$ となるものを作れ（$SU(N)$ に入るよう位相を調整せよ）。$N=2$ では $\Omega_1=i\sigma_3$, $\Omega_2=i\sigma_1$ でよいことを確かめよ。\\
(4) 基本表現の場 $\psi(x+Le_\mu)=\Omega_\mu(x)\psi(x)$ が存在すれば $z=1$ でなければならないことを示せ。$m$ が $SU(N)/\Z_N$ 束の位相的な類（$\pi_1(SU(N)/\Z_N)=\Z_N$）であることを説明せよ。
\end{problem}

\begin{solution}
(1) $A(x+L(e_1+e_2))$ を、まず $e_2$ 方向、次に $e_1$ 方向にずらして計算すると、問題 \ref{ch07-prob-cocycle} (2) の合成則より
\[
A(x+Le_2+Le_1)=\Omega_1(x+Le_2)\ast\bigl(\Omega_2(x)\ast A(x)\bigr)=\bigl(\Omega_1(x+Le_2)\Omega_2(x)\bigr)\ast A(x).
\]
逆の順では $\bigl(\Omega_2(x+Le_1)\Omega_1(x)\bigr)\ast A(x)$。随伴場についても同様に $\phi\mapsto U\phi U^{-1}$ で $U=\Omega_1(x+Le_2)\Omega_2(x)$ と $U'=\Omega_2(x+Le_1)\Omega_1(x)$ が現れる。任意の $\phi(x)\in\mathfrak{su}(N)$ に対し $U\phi U^{-1}=U'\phi U'^{-1}$ となるには、$c=U'^{-1}U$ が全ての $\mathfrak{su}(N)$ の元と可換、すなわち（シューアの補題より）スカラー行列でなければならない。$c\in SU(N)$ のスカラー行列は $\det c=1$ より $c=z\mathbf 1$、$z^N=1$。このとき $z\ast A=A$（$z$ は定数で中心）なので $A$ についても矛盾しない。$c(x)$ は連続で離散集合 $\Z_N$ に値をとるので定数。よって $U=z\,U'$、すなわち与式（$U'$ と $U$ の順序から $z$ の位置が決まる）。

(2) 新しい境界条件で $U$ を作ると
\[
\begin{aligned}\Omega_1'(x+Le_2)\Omega_2'(x)&=h(x+Le_1+Le_2)\,\Omega_1(x+Le_2)\,h(x+Le_2)^{-1}\,h(x+Le_2)\,\Omega_2(x)\,h(x)^{-1}\\&=h(x+L(e_1+e_2))\,U\,h(x)^{-1}.\end{aligned}
\]
同様に $U'_{\rm new}=h(x+L(e_1+e_2))\,U'\,h(x)^{-1}$。$U=zU'$ で $z$ は中心なので $U_{\rm new}=zU'_{\rm new}$、$z$ は不変。したがって $m$ は「周期的なゲージ変換」はもちろん、周期的でない（境界条件を変える）ゲージ変換でも変わらない、位相的な量（'t Hooft フラックス）である。

(3) $P e_k=\omega^{k-1}e_k$ とすると、$PQe_k=Pe_{k+1}=\omega^ke_{k+1}$、$QPe_k=\omega^{k-1}e_{k+1}$。よって $PQ=\omega QP$（$k=N$ でも $\omega^N=1$ により成立）。$\det P=\omega^{N(N-1)/2}=e^{i\pi(N-1)}=(-1)^{N-1}$、$\det Q=(-1)^{N-1}$（長さ $N$ の巡回置換）。$\zeta=e^{i\pi(N-1)/N}$ とおくと $\zeta^N=(-1)^{N-1}$ なので、$\Omega_1=\zeta^{-1}Q$, $\Omega_2=\zeta^{-1}P$ は $SU(N)$ の定数行列である。定数なので
\[
\Omega_1\Omega_2=\zeta^{-2}QP,\qquad \Omega_2\Omega_1=\zeta^{-2}PQ=\omega\,\zeta^{-2}QP,
\]
すなわち $\Omega_1\Omega_2=\omega^{-1}\cdot\Omega_2\Omega_1$。(1) の式 $\Omega_1\Omega_2=z\,\Omega_2\Omega_1$ と比べて $z=\omega^{-1}$、$m=-1$。$\Omega_1$ と $\Omega_2$ を入れ替える（$\Omega_1=\zeta^{-1}P$, $\Omega_2=\zeta^{-1}Q$）と $\Omega_1\Omega_2=\zeta^{-2}PQ=\omega\,\Omega_2\Omega_1$ で $m=1$。一般の $m$ には $\Omega_1=\zeta^{-m}P^m$ とすればよい。

$N=2$：$\Omega_1\Omega_2=(i\sigma_3)(i\sigma_1)=-\sigma_3\sigma_1=-i\sigma_2$、$\Omega_2\Omega_1=-\sigma_1\sigma_3=i\sigma_2$。よって $\Omega_1\Omega_2=-\Omega_2\Omega_1$、$z=-1=e^{2\pi i\cdot1/2}$、$m=1$。このとき $A=0$ は境界条件をみたす（定数の $\Omega$ では $\Omega\ast0=0$）ので、ねじれたセクターにも平坦な配位が存在する。

(4) 基本表現では $\psi(x+L(e_1+e_2))=U\psi(x)=U'\psi(x)$ が必要だが、$U=zU'$ なので $(z-1)U'\psi(x)=0$。$z\neq1$ なら $\psi\equiv0$ しか許されない。つまり $m\neq0$ のセクターは、中心 $\Z_N$ が自明に作用する場だけからなる理論（純粋ヤン・ミルズや随伴物質）でのみ存在する。

このとき真の構造群は $PSU(N)=SU(N)/\Z_N$ であり、$\Omega_\mu$ を $PSU(N)$ に射影したものはコサイクル条件をみたす（$z$ が消える）ので $T^2$ 上の $PSU(N)$ 束を定める。その $SU(N)$ への持ち上げのコサイクル条件が $z$ だけずれていることが、$PSU(N)$ 束が $SU(N)$ 束に持ち上がらない障害であり、その値 $m\in\Z_N=\pi_1(PSU(N))\cong H^2(T^2;\Z_N)$ が束の類である（問題 \ref{ch07-prob-s2bundles} (4) の $SO(3)=SU(2)/\Z_2$ と同じ構造）。$T^4$ 上では各平面 $(\mu\nu)$ ごとにねじれ $n_{\mu\nu}\in\Z_N$ を入れられ、インスタントン数が一般に $1/N$ の整数倍となる（分数インスタントン；第8章）。これは閉じ込めの議論（'t Hooft ループ、中心対称性）で重要である。
\end{solution}

\subsection{スピン構造}

\begin{problem}[難易度 B]\label{ch07-prob-spins1}
(1) $S^1$ のスピン構造は、正規直交枠束（接ベクトルに沿った枠による自明化 $P_{SO}\cong S^1$）の2重被覆であることを用いて、ちょうど2つあることを示せ。それぞれに対しスピノル $\psi(\sigma)$（$\sigma\sim\sigma+2\pi$）の境界条件を述べよ。\\
(2) $S^1$ 上のディラック作用素 $-i\,d/d\sigma$ の固有値を、R（周期的）と NS（反周期的）それぞれで求め、ゼロモードが存在するのはどちらか答えよ。弦理論の R セクターと NS セクターの基底状態の違いと関連づけよ。\\
(3) 円板 $D^2$ の境界 $S^1$ に誘導されるスピン構造は NS であることを、$\Spin(2)=U(1)\to SO(2)$, $\lambda\mapsto\lambda^2$ を用いて示せ。\\
(4) $T^2$ のスピン構造の数を求め、平坦な $T^2$ でディラック作用素のゼロモードをもつのはどれか答えよ。
\end{problem}

\begin{solution}
(1) $S^1$ では $SO(1)=1$、$\Spin(1)=\{\pm1\}$ で、枠束は $S^1$ 自身（各点で単位接ベクトル $\partial_\sigma$ が唯一の正の枠）。スピン構造は $S^1$ の2重被覆（主 $\Z_2$ 束）であり、それは $\pi_1(S^1)=\Z\to\Z_2$ の準同型（第2章の被覆の分類）、すなわち $\Z_2$ の2元に対応する：自明な2重被覆 $S^1\sqcup S^1$ と、連結な2重被覆 $z\mapsto z^2$。これは $\lvert H^1(S^1;\Z_2)\rvert=2$ と一致する。スピノルは2重被覆上の関数で、被覆変換で符号を変えるものなので、自明な被覆では $\psi(\sigma+2\pi)=\psi(\sigma)$（R）、連結な被覆では一周すると他方のシートに移るので $\psi(\sigma+2\pi)=-\psi(\sigma)$（NS）。

(2) $-i\psi'=\lambda\psi$ の解は $e^{i\lambda\sigma}$。R では $e^{2\pi i\lambda}=1$ より $\lambda\in\Z$、NS では $e^{2\pi i\lambda}=-1$ より $\lambda\in\Z+\tfrac12$。ゼロモード（$\lambda=0$）は R にのみ存在する。弦の世界面フェルミオン $\psi^\mu$ のモード展開は R で $\psi^\mu_n$（$n\in\Z$）、NS で $\psi^\mu_r$（$r\in\Z+\frac12$）となり、R セクターのゼロモード $\psi^\mu_0$ はクリフォード代数 $\{\psi_0^\mu,\psi_0^\nu\}=\eta^{\mu\nu}$ をなすので基底状態は時空のスピノル（時空フェルミオン）になる。NS セクターにはゼロモードがなく基底状態は一意（時空ボソン；GSO 射影前はタキオン）。

(3) $D^2$ は可縮なので、枠束は定数枠 $(\partial_x,\partial_y)$ で自明化され、スピン構造は自明なもの1つだけである（$H^1(D^2;\Z_2)=0$）。その上のスピノルは定数枠に関して1価な関数である。境界 $\sigma=\varphi$ 上の接ベクトル枠 $t(\varphi)=(-\sin\varphi,\cos\varphi)$ は定数枠を角度 $\varphi+\frac\pi2$ 回転したもので、$\varphi$ が $0\to2\pi$ と動くとき回転角は $2\pi$ 増える。$\Spin(2)\to SO(2)$ は $\lambda\mapsto\lambda^2$ なので、角度 $\beta$ の回転の持ち上げは $e^{i\beta/2}$ であり、スピノル成分を定数枠から接ベクトル枠に書き直すと因子 $e^{-i(\varphi+\pi/2)/2}$ が掛かる。これは一周で $e^{-i\pi}=-1$ を生むので、接ベクトル枠に関して反周期的、すなわち NS である。R は円板の境界として延長できない（バウンディングでない）。共形場理論で円筒 $w=\tau+i\sigma$ を平面 $z=e^w$ に写すと、重み $\frac12$ の場は $\psi_{\rm cyl}=(dz/dw)^{1/2}\psi_{\rm pl}=e^{w/2}\psi_{\rm pl}$ と変換し、平面上で1価な場が円筒上で反周期的（NS）になるのは同じ事実の表れである。

(4) $\lvert H^1(T^2;\Z_2)\rvert=\lvert\Z_2^2\rvert=4$。平坦な $T^2$ では接束が定数枠で自明化されるので、2つの基本周期の方向それぞれについて独立に周期的（R）・反周期的（NS）を選べ、(NS,NS), (NS,R), (R,NS), (R,R) の4つ。ディラック作用素の固有値は運動量 $\bm k\in(\Z+\epsilon_1)\times(\Z+\epsilon_2)$（を基本格子で読み替えたもの、$\epsilon=0$ が R、$\frac12$ が NS）の大きさに比例するので、ゼロモード（$\bm k=0$）があるのは (R,R) だけで、それは定数スピノルである。
\end{solution}

\begin{center}
\begin{tikzpicture}[scale=0.95]
  \foreach \i/\sa/\sb/\lab/\col in {0/-/-/{(NS,NS)}/blue!6, 1/-/+/{(NS,R)}/blue!6, 2/+/-/{(R,NS)}/blue!6, 3/+/+/{(R,R)}/red!10}{
    \begin{scope}[xshift=\i*3.3cm]
      \fill[\col] (0,0) rectangle (2,2);
      \draw[thick,mid arrow] (0,0) -- (2,0);
      \draw[thick,mid arrow] (0,2) -- (2,2);
      \draw[thick,mid arrows] (0,0) -- (0,2);
      \draw[thick,mid arrows] (2,0) -- (2,2);
      \node[below] at (1,-0.05) {\footnotesize $a$：$\sa$};
      \node[right] at (2.02,1) {\footnotesize $b$：$\sb$};
      \node at (1,1) {\small \lab};
    \end{scope}}
  \node[red!60!black] at (10.9,-0.75) {\footnotesize 奇（ゼロモードあり）};
  \node at (3.3,-0.75) {\footnotesize 偶（モジュラー変換で互いに移る）};
\end{tikzpicture}
\zucap{$T^2$ の4つのスピン構造。$a$ 方向（空間 $\sigma$）と $b$ 方向（虚時間）それぞれについて、スピノルが周期的（$+$，R）か反周期的（$-$，NS）かを指定する。(R,R) のみが奇スピン構造で、定数スピノルのゼロモードをもつ。}
\end{center}

\begin{problem}[難易度 C]\label{ch07-prob-gso}
平坦なトーラス $T^2=\C/(\Z+\tau\Z)$ の4つのスピン構造を、$a$ サイクル（$z\to z+1$）と $b$ サイクル（$z\to z+\tau$）に沿った符号 $(\epsilon_a,\epsilon_b)\in\{\pm1\}^2$ で表す（$+1$ が R、$-1$ が NS；定数枠に関する周期性）。\\
(1) 原始的なサイクル $pa+qb$（$\gcd(p,q)=1$）に沿った直線上の符号は $\epsilon_a^p\epsilon_b^q$ であることを説明せよ。\\
(2) モジュラー変換 $T\colon\tau\mapsto\tau+1$（新しい基底 $a'=a$, $b'=a+b$）と $S\colon\tau\mapsto-1/\tau$（$a'=b$, $b'=-a$）によるスピン構造の移り方を求め、$SL(2,\Z)$ の作用の軌道が $\{(\text{NS,NS}),(\text{NS,R}),(\text{R,NS})\}$ と $\{(\text{R,R})\}$ であることを示せ。\\
(3) $a$ を空間、$b$ を虚時間の方向とすると、自由フェルミオンの分配関数で $\Tr_{\rm NS}q^{H}$, $\Tr_{\rm NS}(-1)^Fq^{H}$, $\Tr_{\rm R}q^H$, $\Tr_{\rm R}(-1)^Fq^H$ はそれぞれどのスピン構造に対応するか。各スピン構造の分配関数 $Z_{(\epsilon_a,\epsilon_b)}$ がモジュラー変換で（位相を除き）スピン構造の移り方どおりに置換されるとして、NS セクターだけからなる理論はモジュラー不変になりえず、R セクターと $(-1)^F$ の挿入（GSO 射影）が必要であることを説明せよ。
\end{problem}

\begin{solution}
(1) 平坦なトーラスでは接束が定数枠で大域的に自明化され、スピン構造はこの枠に関して、格子のベクトル $\gamma\in\Z+\tau\Z$ だけ平行移動したときにスピノルに掛かる符号 $\epsilon(\gamma)$ で決まる。2回続けて移動すれば符号は掛け算されるので $\epsilon(\gamma+\gamma')=\epsilon(\gamma)\epsilon(\gamma')$（準同型）。よって $\epsilon(pa+qb)=\epsilon_a^p\epsilon_b^q$。サイクル $pa+qb$ を表す直線の接ベクトルは定数なので、これはそのサイクル自身の接ベクトル枠に関する周期性（R/NS）でもある。

(2) $SL(2,\Z)$ の元は線形写像として定数枠を定数枠に写すので、新しい基底に沿った符号を読めばよい。$T$：$\epsilon_{a'}=\epsilon_a$、$\epsilon_{b'}=\epsilon_a\epsilon_b$。
\begin{gather*}
(\text{NS,NS})=(-,-)\mapsto(-,+)=(\text{NS,R}),\\ (\text{NS,R})\mapsto(\text{NS,NS}),\quad(\text{R,NS})\mapsto(\text{R,NS}),\quad(\text{R,R})\mapsto(\text{R,R}).
\end{gather*}
$S$：$\epsilon_{a'}=\epsilon_b$、$\epsilon_{b'}=\epsilon_a^{-1}=\epsilon_a$ なので成分の入れ替え：$(\text{NS,R})\leftrightarrow(\text{R,NS})$、$(\text{NS,NS})$, $(\text{R,R})$ は不変。$T$, $S$ は $SL(2,\Z)$ を生成するので、軌道は $T$ と $S$ で移り合うものの集まりで、$\{(\text{NS,NS}),(\text{NS,R}),(\text{R,NS})\}$（$T$ で前2者、$S$ で後2者が結ばれる）と $\{(\text{R,R})\}$。(R,R) は $\epsilon(\gamma)=+1$（全ての $\gamma$）で明らかに不変である。

(3) 虚時間方向の境界条件は、$\Tr\,q^H$ ではフェルミオンに対して反周期的（経路積分でのフェルミオンのトレースの標準的な性質）、$\Tr(-1)^Fq^H$ では周期的。したがって
\begin{gather*}
\Tr_{\rm NS}q^H\leftrightarrow(\text{NS,NS}),\quad\Tr_{\rm NS}(-1)^Fq^H\leftrightarrow(\text{NS,R}),\\ \Tr_{\rm R}q^H\leftrightarrow(\text{R,NS}),\quad\Tr_{\rm R}(-1)^Fq^H\leftrightarrow(\text{R,R}).
\end{gather*}
トーラス上の理論の分配関数は、世界面の座標の取り方（$\tau$ の $SL(2,\Z)$ による取り換え）によらない量でなければならない。$Z_{(\text{NS,NS})}$ だけでは $T$ によって $Z_{(\text{NS,R})}$ に、さらに $S$ によって $Z_{(\text{R,NS})}$ に移るので不変にならない。不変な組合せを作るには偶の軌道の3つを（位相を除いて）等しい重みで足す必要がある：
\[
Z\propto\Tr_{\rm NS}\frac{1\pm(-1)^F}{2}q^H-\Tr_{\rm R}\frac{1\pm(-1)^F}{2}q^H
\]
の形（符号は位相因子とスピン統計で決まる）。すなわち NS セクターに $(-1)^F$ による射影を施し、R セクターを加えることが必要で、これが GSO 射影である。(R,R) は単独で不変なので、その係数の符号は自由に選べ、型 IIA と IIB（あるいは R セクターの時空カイラリティ）の選択に対応する。高種数の世界面でも同様に、$2^{2g}$ 個のスピン構造の和をとる（問題 \ref{ch07-prob-arf}）。
\end{solution}

\begin{problem}[難易度 D]\label{ch07-prob-arf}
種数 $g$ の閉曲面 $\Sigma_g$ のスピン構造は、二次形式
\[
q\colon H_1(\Sigma_g;\Z_2)\to\Z_2,\qquad q(x+y)=q(x)+q(y)+x\cdot y
\]
（$x\cdot y$ は $\bmod2$ の交叉数、第6章）と1対1に対応することが知られている：埋め込まれた円周 $\gamma$ に沿ってスピノルが（$\gamma$ の接ベクトル枠に関して）NS なら $q([\gamma])=0$、R なら $1$。これを認めて以下に答えよ。$a_i,b_i$（$i=1,\dots,g$）を $a_i\cdot b_j=\delta_{ij}$, $a_i\cdot a_j=b_i\cdot b_j=0$ をみたす $H_1(\Sigma_g;\Z_2)$ の基底とする。\\
(1) $q$ は $(q(a_i),q(b_i))\in\Z_2^{2g}$ で決まり、任意の値をとりうることを示し、スピン構造の数 $2^{2g}=\lvert H^1(\Sigma_g;\Z_2)\rvert$ を確かめよ。$g=1$ のとき、問題 \ref{ch07-prob-gso} の乗法的な記述 $\epsilon(\gamma)$ と $q$ の関係 $\epsilon=-(-1)^q$ が二次形式の性質と整合することを確かめよ。\\
(2) $\displaystyle\sum_{x\in H_1}(-1)^{q(x)}=2^g(-1)^{\mathrm{Arf}(q)}$、$\mathrm{Arf}(q)=\sum_iq(a_i)q(b_i)\bmod2$ を示せ。これより $\mathrm{Arf}(q)$ が基底の取り方によらないこと、および $q$ が多数派としてとる値が $\mathrm{Arf}(q)$ であることを示せ。\\
(3) $\mathrm{Arf}(q)=0$（偶）のスピン構造は $2^{g-1}(2^g+1)$ 個、$\mathrm{Arf}(q)=1$（奇）は $2^{g-1}(2^g-1)$ 個であることを示せ。$g=1,2$ で具体的な数を答えよ。
\end{problem}

\begin{solution}
(1) $x=\sum_i(u_ia_i+v_ib_i)$ に対し、二次形式の性質を繰り返し用いると
\[
q(x)=\sum_i\bigl(u_iq(a_i)+v_iq(b_i)\bigr)+\sum_{\text{ペア}}(\text{交叉数})=\sum_i\bigl(u_iq(a_i)+v_iq(b_i)+u_iv_i\bigr),
\]
ここで交叉数で生き残るのは $a_i\cdot b_i=1$ の項だけである（$u_i^2=u_i$ にも注意）。よって $q$ は $2g$ 個の値 $(q(a_i),q(b_i))$ で決まる。逆に任意の値 $(x_i,y_i)$ に対して右辺の式で $q$ を定義すると、$q(x+x')-q(x)-q(x')=\sum_i(u_iv_i'+u_i'v_i)=x\cdot x'$ なので二次形式になる。したがって $2^{2g}$ 個で、$H^1(\Sigma_g;\Z_2)\cong\Z_2^{2g}$ の位数と一致する（2つの二次形式の差 $q-q'$ は線形形式、すなわち $H^1$ の元であり、定理 \ref{thm:ch07-spinclass} の「$H^1$ による移り合い」に対応する）。

$g=1$：問題 \ref{ch07-prob-gso} では $\epsilon(\gamma+\gamma')=\epsilon(\gamma)\epsilon(\gamma')$。$\epsilon=-(-1)^q$ を代入すると $(-1)^{q(\gamma+\gamma')+1}=(-1)^{q(\gamma)+q(\gamma')}$、すなわち $q(\gamma+\gamma')=q(\gamma)+q(\gamma')+1$ となり、$a\cdot b=1$ のときの二次形式の性質と一致する。例えば (R,R) は $q(a)=q(b)=1$、$q(a+b)=1+1+1=1$（R）で、$a+b$ 方向の直線でも周期的であることと整合する。

(2) $x_i=q(a_i)$, $y_i=q(b_i)$ とおくと、(1) の式より $(-1)^{q(x)}=\prod_i(-1)^{u_ix_i+v_iy_i+u_iv_i}$ で、和は各 $i$ について独立に分解する：
\[
\sum_{x}(-1)^{q(x)}=\prod_{i=1}^g\ \sum_{u,v\in\Z_2}(-1)^{ux_i+vy_i+uv}.
\]
\begin{note}
1つのハンドルの和 $s(x_i,y_i)=\sum_{u,v}(-1)^{ux_i+vy_i+uv}$：$(0,0)$ では $1+1+1-1=2$。$(1,0)$ では $(u,v)=(0,0),(1,0),(0,1),(1,1)$ の指数が $0,1,0,1+1=0$ で $1-1+1+1=2$。$(0,1)$ も対称に $2$。$(1,1)$ では指数が $0,1,1,1+1+1=1$ で $1-1-1-1=-2$。まとめて $s=2(-1)^{x_iy_i}$。
\end{note}
よって $\sum_x(-1)^{q(x)}=2^g(-1)^{\sum_ix_iy_i}=2^g(-1)^{\mathrm{Arf}(q)}$。左辺は基底によらない量なので $\mathrm{Arf}(q)$ も基底によらない。また $q(x)=0$ となる $x$ の個数を $N_0$、$=1$ を $N_1$ とすると $N_0+N_1=4^g$、$N_0-N_1=2^g(-1)^{\mathrm{Arf}}$ なので
\[
N_{\mathrm{Arf}(q)}=\frac{4^g+2^g}{2}>\frac{4^g-2^g}{2}=N_{1-\mathrm{Arf}(q)},
\]
すなわち $q$ は多数派として値 $\mathrm{Arf}(q)$ をとる。

(3) スピン構造は $(x,y)\in\Z_2^g\times\Z_2^g$ と1対1で、$\mathrm{Arf}=\sum_ix_iy_i$。偶の個数を $E_g$、奇を $O_g$ とすると、(2) と全く同じ計算で
\[
E_g-O_g=\sum_{x,y}(-1)^{\sum_ix_iy_i}=\prod_{i=1}^g\sum_{x_i,y_i}(-1)^{x_iy_i}=2^g,\qquad E_g+O_g=4^g .
\]
よって $E_g=\frac{4^g+2^g}2=2^{g-1}(2^g+1)$、$O_g=2^{g-1}(2^g-1)$。$g=1$：偶3個、奇1個（(R,R) のみ、問題 \ref{ch07-prob-gso} と一致）。$g=2$：偶10個、奇6個。

物理的には、奇スピン構造の上では（カイラルな）ディラック作用素のゼロモードの数が奇数になる（アティヤの mod 2 指数：$\dim\ker\equiv\mathrm{Arf}(q)\pmod2$）。超弦理論の種数 $g$ の振幅はスピン構造の和をとり、偶スピン構造と奇スピン構造の寄与は性質が異なる（奇スピン構造ではゼロモードの吸収が必要で、パリティを破る項 $\epsilon^{\mu\nu\cdots}$ を生む）。$g=1$ の (R,R) が単独のモジュラー軌道をなしたことは、偶・奇がモジュラー群（写像類群）で保たれることの最も簡単な例である。
\end{solution}

%================================================
\chapter{接続・曲率・特性類}\label{ch:08}
% 18問
%================================================

\section{物理からの動機}

\subsection*{ゲージ場は接続である}
ゲージ理論の出発点は、場 $\psi(x)$ の「内部の向き」を点ごとに独立に回してよい、という要請であった。異なる点 $x$, $x'$ の $\psi(x)$ と $\psi(x')$ は別々のベクトル空間（ファイバー）に属するので、そのままでは比較できない。比べるには、道に沿って一方を他方の点まで運ぶ規則が要る。それを与えるのがゲージ場 $A$ であり、数学ではファイバー束の\term{接続}{connection}と呼ぶ。共変微分 $D=d-iA$ は「運んでから比べた」微分であり、$D\psi=0$ をみたすように運ぶことを\term{平行移動}{parallel transport}という。

閉じた道に沿って1周運ぶと、一般にはもとに戻らない。そのずれ
\[
U(C)=P\exp\Bigl(i\oint_CA\Bigr)
\]
が\term{ホロノミー}{holonomy}であり、そのトレース $W(C)=\tr U(C)$ が\term{ウィルソン・ループ}{Wilson loop}である。微小なループのホロノミーのずれが場の強さ $F$、すなわち\term{曲率}{curvature}である。ウィルソン・ループはゲージ不変であり、格子ゲージ理論の基本変数であり、その期待値の面積則がクォークの閉じ込めの判定条件になる（下図）。アハラノフ・ボーム効果（第2, 5章）は、$F=0$ の領域でもホロノミーが自明でない（\textbf{平坦だが自明でない接続}）ことの観測であった。

\begin{center}
\begin{tikzpicture}[scale=1.05]
  % 左：平行移動
  \draw[thick,blue!70!black,mid arrow] (0,0) .. controls (0.8,1.4) and (2.0,-0.4) .. (3,1) node[pos=0.5,below=3pt] {$\gamma$};
  \node[vtx] at (0,0) {}; \node[below left] at (0,0) {$x_0$};
  \node[vtx] at (3,1) {}; \node[above right] at (3,1) {$x_1$};
  \draw[-{Stealth[length=2mm]},thick,red!70!black] (0,0) -- (0.1,0.7) node[left] {\footnotesize $\psi(x_0)$};
  \draw[-{Stealth[length=2mm]},thick,red!70!black] (3,1) -- (3.55,1.45) node[right] {\footnotesize $U(\gamma)\psi(x_0)$};
  \node at (1.5,-1.0) {\footnotesize $U(\gamma)=P\exp\bigl(i\int_\gamma A\bigr)$};
  % 右：ウィルソン・ループ
  \begin{scope}[xshift=6.3cm,yshift=-0.9cm]
    \draw[-{Stealth[length=2mm]}] (-0.4,0) -- (3.6,0) node[right] {\footnotesize $x$};
    \draw[-{Stealth[length=2mm]}] (0,-0.3) -- (0,2.9) node[above] {\footnotesize $t$};
    \fill[blue!8] (0.62,0.42) rectangle (2.78,2.38);
    \draw[very thick,blue!70!black,mid arrow] (0.6,0.4) -- (2.8,0.4);
    \draw[very thick,blue!70!black,mid arrow] (2.8,0.4) -- (2.8,2.4);
    \draw[very thick,blue!70!black,mid arrow] (2.8,2.4) -- (0.6,2.4);
    \draw[very thick,blue!70!black,mid arrow] (0.6,2.4) -- (0.6,0.4);
    \node at (1.7,1.4) {\footnotesize 面積 $RT$};
    \draw[decorate,decoration={brace,amplitude=4pt,mirror}] (0.6,0.25) -- (2.8,0.25) node[midway,below=4pt] {\footnotesize $R$};
    \draw[decorate,decoration={brace,amplitude=4pt}] (3.0,2.4) -- (3.0,0.4) node[midway,right=4pt] {\footnotesize $T$};
    \node[left] at (0.6,1.9) {\footnotesize $q$};
    \node[right] at (2.8,1.9) {\footnotesize $\bar q$};
  \end{scope}
\end{tikzpicture}
\zucap{左：道 $\gamma$ に沿った平行移動。右：静的なクォーク・反クォーク対の世界線がつくる長方形のウィルソン・ループ。$\langle W\rangle\sim e^{-V(R)T}$ で、面積則 $\langle W\rangle\sim e^{-\sigma RT}$ なら線形ポテンシャル $V(R)=\sigma R$（閉じ込め）。}
\end{center}

\subsection*{位相的な量は曲率の積分で書ける}
物理には、場の強さの積分で書かれるのに、場の連続的な変形で変わらない量がいくつも現れる。
\begin{itemize}
\item 磁気単極子の磁荷：$\displaystyle\frac{1}{2\pi}\int_{S^2}F\in\Z$（ディラックの量子化、第5, 7章）。
\item インスタントン数：$\displaystyle k=\frac{1}{8\pi^2}\int_{\R^4}\tr F\wedge F\in\Z$。
\item 2次元重力：$\displaystyle\frac{1}{4\pi}\int_\Sigma\sqrt h\,R\,d^2\sigma=\chi(\Sigma)$（ガウス・ボンネ）。
\end{itemize}
これらが整数になり、計量やゲージ場の局所的な変形で変わらない理由を一般的に説明するのが\term{チャーン・ヴェイユ理論}{Chern--Weil theory}である：曲率の不変多項式は閉形式であり、接続を変えても完全形式しか変わらないので、そのコホモロジー類（\term{特性類}{characteristic class}）は束だけで決まる。整数性はさらに、特性類が整数係数のコホモロジー類であることから来る。

\subsection*{チャーン・サイモンズ形式と θ 真空}
$\tr F\wedge F$ は局所的には完全形式 $8\pi^2\,dCS(A)$ である。インスタントン数は、ユークリッド時間の過去と未来の真空のチャーン・サイモンズ数の差として書ける。チャーン・サイモンズ数は大きなゲージ変換で整数だけずれるので、ヤン・ミルズ理論の古典的真空は整数でラベルされた無限個の点に分かれ、インスタントンはその間のトンネル効果を表す。ゲージ不変な真空は重ね合わせ（θ 真空）であり、経路積分には重み $e^{i\theta k}$ が現れる。これが強い CP 問題の θ 角である。$k\in\Z$ だから $\theta\sim\theta+2\pi$。3次元では、チャーン・サイモンズ作用 $\tfrac{\kappa}{4\pi}\int\tr(A\,dA-\tfrac{2i}{3}A^3)$ のレベル $\kappa$ が同じ理由で整数に量子化される。

\subsection*{ベリー位相}
量子力学でハミルトニアン $H(\lambda)$ のパラメータ $\lambda$ をゆっくり1周させると、固有状態は動的位相のほかに幾何学的な位相を獲得する（\term{ベリー位相}{Berry phase}）。これはパラメータ空間上の $U(1)$ 接続（ベリー接続）のホロノミーそのものであり、磁場中のスピン $1/2$ では、パラメータ空間 $\bm B$ の原点（縮退点）に置かれた磁荷 $-1/2$ の単極子の「AB 位相」になる。

\begin{physbox}[：本章の辞書]
\begin{center}
\small
\begin{tabular}{p{0.45\linewidth}p{0.45\linewidth}}
物理 & 数学\\ \hline
ゲージ場 $A_\mu$ & 主束（ベクトル束）の接続\\
場の強さ $F_{\mu\nu}$ & 曲率\\
ゲージ変換 & 束の局所自明化の取りかえ\\
ウィルソン線・ウィルソン・ループ & 平行移動・ホロノミー（のトレース）\\
AB 効果、真空の Wilson line & 平坦接続 $\leftrightarrow\Hom(\pi_1,G)/G$\\
磁荷・インスタントン数・θ 項 & チャーン類・チャーン指標の積分\\
重力の曲率 $R^a{}_b$ & レヴィ・チヴィタ接続（スピン接続）の曲率\\
2次元・4次元のオイラー密度 & オイラー類（ガウス・ボンネ）\\
ベリー接続・ベリー位相 & パラメータ空間上の $U(1)$ 接続・ホロノミー\\
大きなゲージ変換によるチャーン・サイモンズ数のずれ & $\pi_3(G)=\Z$ の巻き付き数
\end{tabular}
\end{center}
\end{physbox}

\section{基本事項のまとめ}

本章では記法と規約のゲージ場・特性類の規約を厳密に用いる。特に
\[
D=d-iA,\qquad F=dA-iA\wedge A,\qquad A\mapsto gAg^{-1}+i\,g\,dg^{-1}.
\]
束の大域的な構成（推移関数、主束、同伴束）は第7章で扱ったものとし、ここでは局所的な記述を中心にまとめる。

\subsection{接続と曲率}

\begin{mydef}{接続（局所表示）}{ch08-conn}
$G$ をコンパクトなリー群（$U(N)$ の部分群）、$E\to M$ を階数 $N$ の複素ベクトル束（構造群 $G$）とする。開被覆 $\{U_i\}$ と推移関数 $g_{ij}\colon U_i\cap U_j\to G$（$\psi_i=g_{ij}\psi_j$）をとる。各 $U_i$ 上のエルミート行列（$\mathfrak g$ 値）1形式 $A_i$ の族で
\[
A_i=g_{ij}A_jg_{ij}^{-1}+i\,g_{ij}\,dg_{ij}^{-1}\qquad(U_i\cap U_j\text{ 上})
\]
をみたすものを $E$ 上の\term{接続}{connection}という。共変微分 $D\psi_i=d\psi_i-iA_i\psi_i$ はこの貼り合わせ規則のもとで $D\psi_i=g_{ij}D\psi_j$ をみたし、切断の微分として大域的に意味をもつ。$\mathfrak g$ 値 $p$ 形式 $\eta$（随伴表現に値をとるもの）には $D\eta=d\eta-i[A,\eta]$（$[A,\eta]=A\wedge\eta-(-1)^p\eta\wedge A$）と作用する。
\end{mydef}
推移関数は「パッチ間のゲージ変換」であり、1つのパッチの中でのゲージ変換 $g\colon U\to G$ は $A\mapsto A^g=gAg^{-1}+i\,g\,dg^{-1}$ である。無限小 $g=1+i\epsilon$ では $\delta A=d\epsilon-i[A,\epsilon]=D\epsilon$。

\begin{mythm}{曲率の性質}{ch08-curv}
$F=dA-iA\wedge A$（成分では $F_{\mu\nu}=\partial_\mu A_\nu-\partial_\nu A_\mu-i[A_\mu,A_\nu]$）について次が成り立つ。
\begin{enumerate}
\item $F$ は $D^2\psi=-iF\psi$ をみたす。ゲージ変換で $F\mapsto gFg^{-1}$（パッチ間では $F_i=g_{ij}F_jg_{ij}^{-1}$）。
\item \term{ビアンキ恒等式}{Bianchi identity}：$DF=dF-i[A,F]=0$。
\item $U(1)$ の場合 $F=dA$ は大域的な閉2形式で、ビアンキ恒等式 $dF=0$ はマクスウェル方程式の斉次部分である。
\end{enumerate}
\end{mythm}
\begin{note}
$D^2\psi=(d-iA)(d\psi-iA\psi)=-i\,dA\,\psi+iA\wedge d\psi-iA\wedge d\psi-A\wedge A\psi=-i(dA-iA\wedge A)\psi$。1形式 $A$ について $d(A\psi)=dA\,\psi-A\wedge d\psi$ に注意。性質 1 は問題 \ref{ch08-prob-bianchi} で直接確かめる。
\end{note}

\subsection{平行移動・ホロノミー・平坦接続}

\begin{mydef}{平行移動とホロノミー}{ch08-hol}
道 $\gamma\colon[0,1]\to M$ に沿って $D\psi/ds=0$、すなわち
\[
\frac{d\psi}{ds}=i\,A_\mu(\gamma(s))\dot\gamma^\mu(s)\,\psi
\]
を解いて得られる線形写像 $\psi(0)\mapsto\psi(1)$ を
\[
U(\gamma)=P\exp\Bigl(i\int_\gamma A\Bigr)=\sum_{n=0}^\infty i^n\!\!\int_{1\ge s_1\ge\cdots\ge s_n\ge0}\!\!\mathcal A(s_1)\cdots\mathcal A(s_n)\,ds_1\cdots ds_n
\]
と書く（$\mathcal A(s)=A_\mu\dot\gamma^\mu$。\term{経路順序積}{path ordering}：後の点ほど左）。閉じた道 $C$ に対する $U(C)$ を\term{ホロノミー}{holonomy}、$W(C)=\tr U(C)$ を\term{ウィルソン・ループ}{Wilson loop}という。ゲージ変換で $U(\gamma)\mapsto g(\gamma(1))\,U(\gamma)\,g(\gamma(0))^{-1}$、したがって $W(C)$ はゲージ不変である（問題 \ref{ch08-prob-transport}）。道の積については $U(\gamma_1\cdot\gamma_2)=U(\gamma_2)U(\gamma_1)$。
\end{mydef}

$U(1)$ では順序は無関係で $U(C)=\exp(iq\oint_CA)$（電荷 $q$）。$C=\partial S$ ならストークスの定理により $U(C)=\exp(iq\int_SF)$ で、ホロノミーは囲んだ磁束で決まる。

\begin{mythm}{平坦接続とホロノミー}{ch08-flat}
$M$ を連結多様体とする。$F=0$ の接続（\term{平坦接続}{flat connection}）では、ホロノミー $U(C)$ は基点 $x_0$ を固定した閉曲線 $C$ のホモトピー類だけで決まり、準同型 $\rho\colon\pi_1(M,x_0)\to G$ を与える（$\pi_1$ の積 $[\gamma_1][\gamma_2]$ は $\gamma_1$ を先にたどるので、正確には $\rho([\gamma])=U(\gamma)^{-1}$ が準同型である。以下この区別は重要でない）。$x_0$ でのゲージ変換は $\rho$ を共役で変える。逆に任意の準同型から平坦接続（を持つ束）が構成でき、
\[
\{\text{平坦 }G\text{ 接続}\}/\text{ゲージ変換}\;\cong\;\Hom(\pi_1(M),G)/G\quad(G\text{ は共役で作用}).
\]
特に単連結な $M$ 上の平坦接続は純ゲージ $A=i\,g\,dg^{-1}$ である。
\end{mythm}
例えば $M=\R^2\setminus\{0\}$ 上の $U(1)$ 平坦接続は $\Hom(\Z,U(1))=U(1)$ でラベルされ、それが AB 位相 $e^{iq\Phi}$ である（問題 \ref{ch08-prob-ab}）。トーラス上の平坦接続は問題 \ref{ch08-prob-torus} で扱う。

\subsection{ベリー接続}

\begin{mydef}{ベリー接続}{ch08-berry}
パラメータ $\lambda\in\mathcal P$ に滑らかに依存するハミルトニアン $H(\lambda)$ が、縮退しない固有状態 $\lvert n(\lambda)\rangle$（$\langle n|n\rangle=1$）をもつとする。固有状態の位相の取り方は任意なので、$\{\lvert n(\lambda)\rangle\}$ は $\mathcal P$ 上の複素直線束 $L_n$ を定める。本書では
\[
\mathcal A=i\langle n|d n\rangle\quad(\text{実の1形式}),\qquad\mathcal F=d\mathcal A
\]
を\term{ベリー接続}{Berry connection}・\term{ベリー曲率}{Berry curvature}と呼ぶ。断熱的に $\lambda$ を閉曲線 $C$ に沿って1周させると、状態は
\[
\lvert\psi(T)\rangle=e^{-i\int_0^TE_n\,dt}\,e^{i\gamma_n(C)}\lvert n(\lambda(0))\rangle,\qquad\gamma_n(C)=\oint_C\mathcal A
\]
となる。すなわち $e^{i\gamma_n}$ は記法と規約の意味での $U(1)$ 接続 $\mathcal A$ のホロノミーである（問題 \ref{ch08-prob-berry}）。位相の取りかえ $\lvert n\rangle\mapsto e^{i\alpha}\lvert n\rangle$ で $\mathcal A\mapsto\mathcal A-d\alpha$ となり、$\gamma_n$ は $2\pi$ を法としてゲージ不変である。
\end{mydef}

\begin{center}
\begin{tikzpicture}[scale=1.25]
  \shade[ball color=blue!8,opacity=0.5] (0,0) circle (1.6);
  \fill[orange!30,opacity=0.8] (-1.226,1.028) arc[start angle=140,end angle=40,radius=1.6]
       arc[start angle=0,end angle=-180,x radius=1.226,y radius=0.25];
  \draw[thick] (0,0) circle (1.6);
  \draw[gray,dashed] (1.6,0) arc[start angle=0,end angle=180,x radius=1.6,y radius=0.35];
  \draw[gray] (-1.6,0) arc[start angle=180,end angle=360,x radius=1.6,y radius=0.35];
  \draw[red!70!black,dashed] (1.226,1.028) arc[start angle=0,end angle=180,x radius=1.226,y radius=0.25];
  \draw[red!70!black,very thick,mid arrow] (-1.226,1.028) arc[start angle=180,end angle=360,x radius=1.226,y radius=0.25];
  \draw[-{Stealth[length=2mm]},thick] (0,0) -- (0.75,0.87) node[right] {\footnotesize $\hat{\bm n}=\bm B/\lvert\bm B\rvert$};
  \draw[-{Stealth[length=1.6mm]}] (0,0) -- (0,2.1) node[above] {\footnotesize $B_z$};
  \draw (0,0.45) arc[start angle=90,end angle=49,radius=0.45];
  \node at (0.2,0.62) {\footnotesize $\theta$};
  \node at (0,1.35) {\footnotesize $\Omega$};
  \node[red!70!black] at (1.75,0.55) {\footnotesize $C$};
  \node[vtx] at (0,0) {};
  \node[below] at (0,-0.05) {\footnotesize 縮退点 $\bm B=0$};
\end{tikzpicture}
\zucap{ブロッホ球（磁場の向きの空間）。磁場の向き $\hat{\bm n}$ が閉曲線 $C$ を1周すると、$\hat{\bm n}$ 方向を向いたスピン $1/2$ の状態は、$C$ が張る立体角 $\Omega$ の半分のベリー位相 $\gamma=-\Omega/2$ を得る（問題 \ref{ch08-prob-berry}）。ベリー曲率は縮退点 $\bm B=0$ に置かれた磁荷 $-1/2$ の単極子の磁場である。}
\end{center}

\subsection{チャーン・ヴェイユ理論と特性類}

\begin{mydef}{不変多項式}{ch08-invpoly}
$\mathfrak g$ 上の多項式 $P(X)$ が\term{不変多項式}{invariant polynomial}であるとは、$P(gXg^{-1})=P(X)$（$g\in G$）をみたすことをいう。例：$\tr X^k$、$\det(1+X)$。曲率 $F$ は偶数次の形式で互いに可換なので $P(F)$ は意味をもち、パッチ間で $F_i=g_{ij}F_jg_{ij}^{-1}$ だから $P(F)$ は $M$ 上の大域的な微分形式である。
\end{mydef}

\begin{mythm}{チャーン・ヴェイユの定理}{ch08-cw}
$P$ を不変多項式とする。
\begin{enumerate}
\item $P(F)$ は閉形式：$dP(F)=0$。
\item 同じ束上の2つの接続 $A_0$, $A_1$ について、$\eta=A_1-A_0$（大域的な $\mathfrak g$ 値1形式）、$A_t=A_0+t\eta$ とすると、$P$ が $k$ 次斉次のとき
\[
P(F_1)-P(F_0)=d\,Q,\qquad Q=k\int_0^1P(\eta,F_t,\dots,F_t)\,dt
\]
（$P(\cdot,\dots,\cdot)$ は $P$ の対称多重線形化）。$Q$ を\term{トランスグレッション形式}{transgression form}という。
\end{enumerate}
したがってド・ラーム類 $[P(F)]\in H^{2k}_{dR}(M)$ は接続によらず、束だけで決まる。
\end{mythm}
\begin{note}
1 は $dP(F)=kP(DF,F,\dots,F)=0$（ビアンキ）から。2 は $\dot F_t=D_t\eta$ と $P(D_t\eta,F_t,\dots)=dP(\eta,F_t,\dots)$ による。$P=\tr X^2$ の場合を問題 \ref{ch08-prob-transg} で具体的に計算する。
\end{note}

\begin{mydef}{チャーン類・チャーン指標}{ch08-chern}
エルミートな曲率 $F$ に対し（記法と規約）
\[
c(F)=\det\Bigl(1+\frac{F}{2\pi}\Bigr)=1+c_1+c_2+\cdots,\]\[\ch(F)=\tr\exp\Bigl(\frac{F}{2\pi}\Bigr)=\operatorname{rank}E+\ch_1+\ch_2+\cdots.
\]
低次では
\[
c_1=\ch_1=\frac{\tr F}{2\pi},\qquad c_2=\frac{1}{8\pi^2}\bigl(\tr F\wedge\tr F-\tr F\wedge F\bigr),\]\[\ch_2=\frac{1}{8\pi^2}\tr F\wedge F=\frac{c_1^2}{2}-c_2 .
\]
$SU(N)$ 束では $\tr F=0$ なので $c_1=0$、$c_2=-\ch_2$。$c_k(E)$ を\term{チャーン類}{Chern class}、$\ch(E)$ を\term{チャーン指標}{Chern character}という。
\end{mydef}

\begin{mythm}{整数性と基本性質}{ch08-int}
\begin{enumerate}
\item $c_k(E)$ は整数係数コホモロジー $H^{2k}(M;\Z)$ の元の実像である。特に閉曲面 $\Sigma$ 上 $\displaystyle\int_\Sigma c_1\in\Z$、閉4次元多様体上 $\displaystyle\int_Mc_2\in\Z$。
\item ホイットニー和公式 $c(E\oplus E')=c(E)c(E')$、$\ch(E\oplus E')=\ch E+\ch E'$、$\ch(E\otimes E')=\ch E\cdot\ch E'$。
\item 直線束 $L$ では $c(L)=1+c_1(L)$、$\ch L=e^{c_1(L)}$、$c_1(L^*)=-c_1(L)$、$c_1(L\otimes L')=c_1(L)+c_1(L')$。$S^2$ 上の直線束は $\int c_1=n\in\Z$ で完全に分類される（第7章の貼り合わせによる分類）。
\end{enumerate}
\end{mythm}
整数性の直観：$S^2$ 上の $U(1)$ 束を北・南のパッチで書くと、$\int_{S^2}F=\oint_{\text{赤道}}(A_N-A_S)=\oint i\,g\,dg^{-1}$ で、推移関数 $g\colon S^1\to U(1)$ の巻き付き数の $2\pi$ 倍になる（問題 \ref{ch08-prob-monopole}）。同様に $S^4$ 上の $SU(2)$ 束のインスタントン数は、赤道 $S^3$ 上の推移関数の巻き付き数である。

\begin{myex}{代表的な束の特性類}{ch08-exchar}
\begin{itemize}
\item 磁荷 $n$ の単極子束（$S^2$ 上）：$\int c_1=n$。
\item $\CP^1$ 上の\term{トートロジカル直線束}{tautological line bundle} $L=\mathcal O(-1)$（ホップ束に同伴）：$\int c_1=-1$。その双対 $\mathcal O(1)$：$+1$（問題 \ref{ch08-prob-hopf}）。
\item $\CP^n$ の接束：$c(T\CP^n)=(1+x)^{n+1}$。$x=c_1(\mathcal O(1))$、$\int_{\CP^n}x^n=1$（問題 \ref{ch08-prob-cpn}）。
\item $\R^4\cup\{\infty\}=S^4$ 上の $k$ インスタントンの $SU(2)$ 束：$\int\ch_2=k$、$\int c_2=-k$。
\end{itemize}
\end{myex}

\begin{mydef}{実ベクトル束の特性類}{ch08-pont}
実ベクトル束 $E$（構造群 $O(n)$）の接続の曲率を実反対称行列値2形式 $R$ とする（$D=d+\omega$、$R=d\omega+\omega\wedge\omega$。エルミート表示では $F=iR$）。
\begin{itemize}
\item \term{ポントリャーギン類}{Pontryagin class}：$p(E)=\det\bigl(1+\tfrac{R}{2\pi}\bigr)=1+p_1+p_2+\cdots$、$p_1=-\dfrac{1}{8\pi^2}\tr R\wedge R$。複素化を用いて $p_k(E)=(-1)^kc_{2k}(E\otimes\C)$。
\item \term{オイラー類}{Euler class}：向きのついた階数 $2m$ の束に対し $e(E)=\Pf\bigl(\tfrac{R}{2\pi}\bigr)$。階数 2 なら $e=\dfrac{R_{12}}{2\pi}$、階数 4 なら $e=\dfrac{1}{32\pi^2}\epsilon_{abcd}R_{ab}\wedge R_{cd}$。$e^2=p_m$。
\item \term{シュティーフェル・ホイットニー類}{Stiefel--Whitney class} $w_k(E)\in H^k(M;\Z_2)$：$\Z_2$ 係数なので曲率では書けない。$w_1=0$ が向き付け可能性、$w_2=0$ がスピン構造の存在条件である（第7章）。ここでは述べるだけにする。
\end{itemize}
向き付けられた階数 2 の実束は複素直線束とみなせ、$e(E)=c_1(E)$ である。
\end{mydef}

\subsection{リーマン幾何とガウス・ボンネの定理}
計量 $g$ をもつ $n$ 次元多様体で、正規直交枠 $e^a$（$g=\sum_ae^a\otimes e^a$）をとると、ねじれなしの条件 $de^a+\omega^a{}_b\wedge e^b=0$（$\omega_{ab}=-\omega_{ba}$）が\term{レヴィ・チヴィタ接続}{Levi-Civita connection}（スピン接続）$\omega$ を一意に定める。曲率 $R^a{}_b=d\omega^a{}_b+\omega^a{}_c\wedge\omega^c{}_b$。2次元では $R^1{}_2=K\,e^1\wedge e^2$（$K$ はガウス曲率、スカラー曲率は $R=2K$）。

\begin{mythm}{ガウス・ボンネ・チャーンの定理}{ch08-gb}
向き付けられた閉多様体 $M^{2m}$ について $\displaystyle\chi(M)=\int_Me(TM)$。特に
\[
\chi(\Sigma)=\frac{1}{2\pi}\int_\Sigma K\,dS=\frac{1}{4\pi}\int_\Sigma\sqrt h\,R\,d^2\sigma,\qquad\chi(M^4)=\frac{1}{32\pi^2}\int_M\epsilon_{abcd}R^{ab}\wedge R^{cd}.
\]
\end{mythm}
2次元の\textbf{アインシュタイン・ヒルベルト作用は位相不変量}であり、局所的な力学をもたない。弦理論では世界面の作用にこの項が入り、種数 $g$ の世界面の寄与は $g_s^{-\chi}=g_s^{2g-2}$ で重みづけられる（問題 \ref{ch08-prob-gb}、第10章）。4次元ではオイラー密度は（ガウス・ボンネ項として）作用に加えても運動方程式を変えない。

曲率の幾何学的な意味は、平行移動のずれである。$S^2$ 上で閉曲線に沿って接ベクトルを平行移動すると、囲んだ面積上の $\int K\,dS$ だけ回転する（下図、問題 \ref{ch08-prob-lc}）。
\begin{center}
\begin{tikzpicture}[scale=2.0]
  \shade[ball color=blue!6,opacity=0.5] (0,0) circle (1);
  \draw[thick] (0,0) circle (1);
  \draw[gray,dashed] (1,0) arc[start angle=0,end angle=180,x radius=1,y radius=0.342];
  \draw[gray] (-1,0) arc[start angle=180,end angle=360,x radius=1,y radius=0.342];
  \fill[orange!25,opacity=0.8] plot[domain=0:90,samples=30] ({-0.707*sin(\x)},{0.94*cos(\x)-0.242*sin(\x)})
     -- plot[domain=0:90,samples=30] ({-0.707*cos(\x)+0.707*sin(\x)},{-0.242*(cos(\x)+sin(\x))})
     -- plot[domain=90:0,samples=30] ({0.707*sin(\x)},{0.94*cos(\x)-0.242*sin(\x)}) -- cycle;
  \draw[very thick,blue!70!black,mid arrow] plot[domain=0:90,samples=30] ({-0.707*sin(\x)},{0.94*cos(\x)-0.242*sin(\x)});
  \draw[very thick,blue!70!black,mid arrow] plot[domain=0:90,samples=30] ({-0.707*cos(\x)+0.707*sin(\x)},{-0.242*(cos(\x)+sin(\x))});
  \draw[very thick,blue!70!black,mid arrow] plot[domain=90:0,samples=30] ({0.707*sin(\x)},{0.94*cos(\x)-0.242*sin(\x)});
  \node[vtx] at (0,0.94) {}; \node[above] at (0,0.97) {\footnotesize $N$};
  \node[vtx] at (-0.707,-0.242) {}; \node[left] at (-0.72,-0.25) {\footnotesize $A$};
  \node[vtx] at (0.707,-0.242) {}; \node[right] at (0.72,-0.25) {\footnotesize $B$};
  \draw[-{Stealth[length=1.8mm]},thick,red!70!black] (0,0.94) -- (-0.32,0.83);
  \draw[-{Stealth[length=1.8mm]},thick,green!45!black] (0,0.94) -- (0.32,0.83);
  \draw[-{Stealth[length=1.8mm]},thick,red!70!black] (-0.707,-0.242) -- (-0.707,-0.56);
  \draw[-{Stealth[length=1.8mm]},thick,red!70!black] (0.707,-0.242) -- (0.707,-0.56);
  \node[red!70!black] at (-0.45,1.0) {\footnotesize $v_0$};
  \node[green!45!black] at (0.48,1.0) {\footnotesize $v_1$};
\end{tikzpicture}
\zucap{単位球面の8分の1の測地三角形 $N\to A\to B\to N$（3つの角はすべて直角、面積 $\pi/2$）。北極の接ベクトル $v_0$ を辺に沿って平行移動して戻ると $v_1$ になり、$\int K\,dS=\pi/2$ だけ回転している。}
\end{center}

\subsection{チャーン・サイモンズ形式と大きなゲージ変換}

\begin{mydef}{チャーン・サイモンズ形式}{ch08-cs}
\[
CS(A)=\frac{1}{8\pi^2}\tr\Bigl(A\wedge dA-\frac{2i}{3}A\wedge A\wedge A\Bigr),\qquad dCS(A)=\frac{1}{8\pi^2}\tr F\wedge F=\ch_2(F).
\]
これは $A_0=0$、$A_1=A$ とした $\tr F^2$ のトランスグレッション形式（を $8\pi^2$ で割ったもの）である（問題 \ref{ch08-prob-transg}）。閉3次元多様体 $M^3$ 上の積分 $N_{CS}[A]=\int_{M^3}CS(A)$ を\term{チャーン・サイモンズ数}{Chern--Simons number}という。
\end{mydef}

\begin{mythm}{チャーン・サイモンズ形式のゲージ変換}{ch08-csgauge}
$A^g=gAg^{-1}+i\,g\,dg^{-1}$ に対し
\[
CS(A^g)-CS(A)=-\frac{1}{24\pi^2}\tr\bigl(g^{-1}dg\bigr)^3+\frac{i}{8\pi^2}\,d\,\tr\bigl(g^{-1}dg\wedge A\bigr).
\]
したがって閉3次元多様体上で
\[
\int_{M^3}CS(A^g)=\int_{M^3}CS(A)-w(g),\qquad w(g)=\frac{1}{24\pi^2}\int_{M^3}\tr(g^{-1}dg)^3\in\Z .
\]
特に純ゲージ $A=i\,g\,dg^{-1}$ について $\int CS(i\,g\,dg^{-1})=-w(g)$。
\end{mythm}
証明は問題 \ref{ch08-prob-csgauge} で行う。符号は規約 $A\mapsto gAg^{-1}+i\,g\,dg^{-1}$ から決まり、「大きなゲージ変換 $w(g)=1$ はチャーン・サイモンズ数を $1$ だけ\textbf{下げる}」。

\paragraph{向きについての注意.}
$\R^4$ の向きを $d^4x=dx^1\wedge dx^2\wedge dx^3\wedge dx^4$ とし、無限遠の球面 $S^3_\infty$ には境界としての向き（外向き法線を先頭に置いたとき $\R^4$ の向きになるもの）を入れてストークスの定理 $\int_{B}d\alpha=\int_{\partial B}\alpha$ を用いる。写像
\[
g(x)=\frac{x^4\mathbf 1+i\,x^a\sigma^a}{\lvert x\rvert}\colon S^3_\infty\to SU(2)
\]
はこの向きに関して $w(g)=-1$ である（北極 $x=(0,0,0,1)$ で $\tr(g^{-1}dg)^3=12\,dx^1\wedge dx^2\wedge dx^3$、一方、境界の向きは $-dx^1\wedge dx^2\wedge dx^3$）。記法と規約の「恒等写像 $S^3\cong SU(2)$ が $w=1$」は、この向きの $S^3$ を $x\mapsto x^4-i\,x^a\sigma^a$ で $SU(2)$ と同一視したものと読めばよい。

\begin{myex}{BPST インスタントン（正則ゲージ）}{ch08-bpst}
$SU(2)$、$T^a=\sigma^a/2$ とし、\term{トフーフト記号}{'t Hooft symbol} $\eta^a_{\mu\nu}$ を
\[
\eta^a_{bc}=\epsilon_{abc},\quad\eta^a_{b4}=\delta_{ab},\quad\eta^a_{4b}=-\delta_{ab},\quad\eta^a_{44}=0\qquad(a,b,c=1,2,3)
\]
で定める（自己双対：$\tfrac12\epsilon_{\mu\nu\rho\sigma}\eta^a_{\rho\sigma}=\eta^a_{\mu\nu}$、$\epsilon_{1234}=1$）。中心 $0$、大きさ $\rho>0$ の\term{BPST インスタントン}{BPST instanton}は
\[
A_\mu=\frac{2\eta^a_{\mu\nu}x^\nu}{x^2+\rho^2}\,T^a=-i\,\frac{\operatorname{Im}(dX\,X^\dagger)_\mu}{x^2+\rho^2},\qquad X=x^4\mathbf 1+i\,x^a\sigma^a,
\]
で与えられ、
\[
F^a_{\mu\nu}=-\frac{4\rho^2\,\eta^a_{\mu\nu}}{(x^2+\rho^2)^2},\qquad F=*F,\qquad\tr F\wedge F=\frac{48\rho^4}{(x^2+\rho^2)^4}\,d^4x,\qquad k=1
\]
である（問題 \ref{ch08-prob-bpst1}, \ref{ch08-prob-bpst2}）。遠方で $A\to i\,g\,dg^{-1}$（$g=X/\lvert x\rvert$、$w(g)=-1$）となり、$k=\int_{S^3_\infty}CS=-w(g)=1$ と整合する。作用は $S_E=8\pi^2/g_{\mathrm{YM}}^2$。
\end{myex}

\begin{center}
\begin{tikzpicture}[xscale=1.1,yscale=2.4]
  \draw[-{Stealth[length=2mm]}] (-3.3,0) -- (3.5,0) node[right] {\footnotesize $x^1$};
  \draw[-{Stealth[length=2mm]}] (0,-0.05) -- (0,1.25);
  \draw[very thick,blue!70!black,domain=-3:3,samples=120,smooth] plot (\x,{1/(1+\x*\x)^4});
  \draw[very thick,red!70!black,dashed,domain=-3:3,samples=120,smooth] plot (\x,{1/(1+\x*\x/2.56)^4});
  \node[blue!70!black] at (1.3,0.85) {\footnotesize $\rho=1$};
  \node[red!70!black] at (2.2,0.35) {\footnotesize $\rho=1.6$};
  \node[left] at (0,1.05) {\footnotesize $\dfrac{\rho^8}{(x^2+\rho^2)^4}$};
  \draw (1,0.03) -- (1,-0.03) node[below] {\footnotesize $1$};
  \draw (-1,0.03) -- (-1,-0.03) node[below] {\footnotesize $-1$};
\end{tikzpicture}
\zucap{BPST インスタントンの作用密度 $\tr F_{\mu\nu}F_{\mu\nu}\propto\rho^4/(x^2+\rho^2)^4$ を原点での値で規格化したもの。大きさ $\rho$ は任意（古典的なスケール不変性）で、どの $\rho$ でも全作用は $8\pi^2/g_{\mathrm{YM}}^2$。}
\end{center}

\subsection{θ 真空とチャーン・サイモンズ・レベル}
ユークリッド時間 $\tau$ をもつ $\R_\tau\times\R^3$ で、テンポラル・ゲージ $A_\tau=0$ をとり、各時刻の空間 $\R^3\cup\{\infty\}\cong S^3$ のチャーン・サイモンズ数を $N_{CS}(\tau)$ とすると（向きは $d\tau\wedge(\R^3\text{ の向き})$ が $\R^4$ の向きになるようにとる）
\[
k=\frac{1}{8\pi^2}\int\tr F\wedge F=N_{CS}(+\infty)-N_{CS}(-\infty).
\]
古典的真空 $A=i\,g\,dg^{-1}$ は $N_{CS}=-w(g)\in\Z$ をもち、インスタントンは $N_{CS}$ を $k$ だけ変えるトンネル過程である（下図）。大きなゲージ変換 $\mathcal T$（$w=1$）は $\lvert n\rangle\mapsto\lvert n-1\rangle$ と作用するので、その固有状態
\[
\lvert\theta\rangle=\sum_{n\in\Z}e^{-in\theta}\lvert n\rangle,\qquad\mathcal T\lvert\theta\rangle=e^{-i\theta}\lvert\theta\rangle
\]
が物理的な真空であり（\term{θ 真空}{θ vacuum}）、経路積分では $k$ セクターに重み $e^{i\theta k}$ がつく。これは作用に $\theta$ 項 $-i\theta\,\frac{1}{8\pi^2}\int\tr F\wedge F$ を加えるのと同じで（ユークリッド作用 $S_E$、重み $e^{-S_E}$）、$k\in\Z$ だから $\theta\sim\theta+2\pi$ である（問題 \ref{ch08-prob-theta}）。

\begin{center}
\begin{tikzpicture}[xscale=1.25,yscale=1.3]
  \draw[-{Stealth[length=2mm]}] (-2.8,0) -- (2.9,0) node[right] {\footnotesize $N_{CS}$};
  \draw[-{Stealth[length=2mm]}] (0,-0.1) -- (0,1.6) node[above] {\footnotesize $V$};
  \draw[very thick,blue!70!black,domain=-2.5:2.5,samples=150,smooth] plot (\x,{sin(180*\x)*sin(180*\x)});
  \foreach \n in {-2,-1,0,1,2} {\node[vtx] at (\n,0) {}; \node[below] at (\n,-0.05) {\footnotesize $\n$};}
  \node[hvtx] at (0.5,1) {}; \node[above] at (0.5,1.03) {\footnotesize スファレロン};
  \draw[-{Stealth[length=2mm]},red!70!black,thick,dashed] (-1,0.12) .. controls (-0.7,0.3) and (-0.3,0.3) .. (0,0.12);
  \draw[-{Stealth[length=2mm]},red!70!black,thick,dashed] (1,0.12) .. controls (1.3,0.3) and (1.7,0.3) .. (2,0.12);
  \node[red!70!black] at (-0.5,0.45) {\footnotesize $k=1$};
  \node[red!70!black] at (1.5,0.45) {\footnotesize $k=1$};
\end{tikzpicture}
\zucap{ゲージ場の配位空間を $N_{CS}$ に沿って見た模式図。古典的真空は整数 $N_{CS}=n$ に並び、障壁の頂上（半整数）は電弱理論ではスファレロンに対応する。インスタントン（$k=1$）は障壁の下をくぐって $n\to n+1$ と移るトンネル過程。大きなゲージ変換は真空を互いに移すので、物理的な真空はその重ね合わせ $\lvert\theta\rangle$ である。}
\end{center}

\begin{myprop}{3次元チャーン・サイモンズ理論のレベル量子化}{ch08-cslevel}
$SU(N)$ ゲージ場の3次元作用（$\tr$ は基本表現、$\tr T^aT^b=\tfrac12\delta^{ab}$）
\[
S_\kappa[A]=\frac{\kappa}{4\pi}\int_{M^3}\tr\Bigl(A\wedge dA-\frac{2i}{3}A\wedge A\wedge A\Bigr)=2\pi\kappa\int_{M^3}CS(A)
\]
は、大きなゲージ変換で $S_\kappa\mapsto S_\kappa-2\pi\kappa\,w(g)$ と変わる。$e^{iS_\kappa}$ がゲージ不変であるためには $\kappa\in\Z$ が必要である（問題 \ref{ch08-prob-cslevel}）。
\end{myprop}

\begin{physbox}[：本章の公式の使いどころ]
\begin{itemize}
\item 指数定理（第9章）は、ディラック作用素のゼロモードの数 $n_+-n_-$ を $\int\ch(F)\Ahat(TM)$ で与える。$\R^4$ 上では $\ind\slashed D=k$ で、これが 't Hooft 頂点・$U(1)_A$ 問題の解決につながる。
\item アノマリー多項式と降下方程式（第9章）は、$\ch_2=dCS$ と本章のチャーン・サイモンズ形式のゲージ変換則を高次元に拡張したものである。
\item コンパクト化（第10章）では、内部空間の $c_1$ や $\chi$ が世代数やフラックス量子化を決める。
\end{itemize}
\end{physbox}

\section{問題}

\subsection{接続・曲率・ホロノミー}

\begin{problem}[難易度 A]\label{ch08-prob-bianchi}
$A$ をエルミート行列値1形式、$F=dA-iA\wedge A$ とする。\\
(1) $A^g=gAg^{-1}+i\,g\,dg^{-1}$ に対し $F^g:=dA^g-iA^g\wedge A^g=gFg^{-1}$ を示せ。\\
(2) ビアンキ恒等式 $DF:=dF-i(A\wedge F-F\wedge A)=0$ を示せ。\\
(3) 無限小ゲージ変換 $g=1+i\epsilon$ について $\delta A=d\epsilon-i[A,\epsilon]$ を示せ。\\
(4) $U(1)$ の場合、(2) が $\nabla\cdot\bm B=0$、$\nabla\times\bm E+\partial_t\bm B=0$ を与えることを説明せよ（$F=\tfrac12F_{\mu\nu}dx^\mu\wedge dx^\nu$、$F_{0i}=-E_i$、$F_{ij}=\epsilon_{ijk}B_k$ とする）。
\end{problem}

\begin{solution}
$\theta=g^{-1}dg$ とおく。$d(g^{-1})=-g^{-1}dg\,g^{-1}$ より $i\,g\,dg^{-1}=-i\,dg\,g^{-1}=g(-i\theta)g^{-1}$。よって $A^g=gBg^{-1}$、$B=A-i\theta$ と書ける。また $d\theta=d(g^{-1})\wedge dg=-\theta\wedge\theta$。

(1) $dg=g\theta$、$dg^{-1}=-\theta g^{-1}$ を用いると
\[
d(gBg^{-1})=dg\wedge Bg^{-1}+g\,dB\,g^{-1}-gB\wedge dg^{-1}=g(\theta\wedge B+dB+B\wedge\theta)g^{-1}.
\]
したがって $F^g=g(dB+\theta B+B\theta-iB\wedge B)g^{-1}$（$\wedge$ を省略）。ここで
\[
dB=dA+i\theta^2,\qquad -iB^2=-iA^2-(A\theta+\theta A)+i\theta^2,\qquad\theta B+B\theta=\theta A+A\theta-2i\theta^2 .
\]
和をとると $\theta$ を含む項はすべて消え、$F^g=g(dA-iA^2)g^{-1}=gFg^{-1}$。

(2) $dA=F+iA^2$ を微分すると $0=d(dA)=dF+i(dA\,A-A\,dA)$。$dA$ に再び $F+iA^2$ を代入すると、$A^3$ の項は打ち消し合い $dF=-i(FA-AF)=i(AF-FA)$。よって $DF=dF-i(AF-FA)=0$。

(3) $g=1+i\epsilon$、$g^{-1}=1-i\epsilon$ として1次まで：$gAg^{-1}=A+i(\epsilon A-A\epsilon)$、$i\,g\,dg^{-1}=i(-i\,d\epsilon)=d\epsilon$。よって $\delta A=d\epsilon+i[\epsilon,A]=d\epsilon-i[A,\epsilon]=D\epsilon$。

(4) $U(1)$ では $A\wedge F-F\wedge A=0$ なので $dF=0$。$F=F_{0i}dt\wedge dx^i+\tfrac12F_{ij}dx^i\wedge dx^j=-E_i\,dt\wedge dx^i+B_1dx^2\wedge dx^3+B_2dx^3\wedge dx^1+B_3dx^1\wedge dx^2$。$dF$ の $dx^1\wedge dx^2\wedge dx^3$ 成分は $\nabla\cdot\bm B$、$dt\wedge dx^j\wedge dx^k$ 成分は $\partial_tB_i+(\nabla\times\bm E)_i$（$(i,j,k)$ は巡回）であり、$dF=0$ はこれらがゼロであることを意味する。$F=dA$ と書けること（局所的なポテンシャルの存在）と同値な条件であり、ゲージ場を接続とみなす立場では斉次マクスウェル方程式は「方程式」ではなく恒等式である。
\end{solution}

\begin{problem}[難易度 B]\label{ch08-prob-transport}
道 $\gamma\colon[0,1]\to M$ に沿う平行移動 $U(s)$ を、$D\psi=0$ の解 $\psi(s)=U(s)\psi(0)$ により定める。\\
(1) $U(s)$ がみたす微分方程式を書き、それが $U(1)=P\exp\bigl(i\int_\gamma A\bigr)$（後の点が左）で解かれることを示せ。$U(1)$ ゲージ群なら $U=\exp(i\int_\gamma A)$ であることを確認せよ。\\
(2) ゲージ変換 $A\mapsto A^g$ のもとで $U\mapsto g(\gamma(1))\,U\,g(\gamma(0))^{-1}$ を示せ。\\
(3) ウィルソン・ループ $W(C)=\tr P\exp(i\oint_CA)$ がゲージ不変で、基点の取り方にもよらないことを示せ。また、$\bar\psi(x_1)U(\gamma)\psi(x_0)$ がゲージ不変な（非局所）演算子であることを説明せよ。
\end{problem}

\begin{solution}
(1) $\mathcal A(s)=A_\mu(\gamma(s))\dot\gamma^\mu(s)$ とおく。$\gamma$ に沿う $D\psi=0$ は $\dot\psi=i\mathcal A(s)\psi$ なので
\[
\frac{dU}{ds}=i\mathcal A(s)U(s),\qquad U(0)=\mathbf 1 .
\]
積分方程式 $U(s)=\mathbf 1+i\int_0^s\mathcal A(s_1)U(s_1)ds_1$ を反復代入すると
\[
U(s)=\mathbf 1+i\int_0^s\!\mathcal A(s_1)ds_1+i^2\!\int_0^s\!ds_1\!\int_0^{s_1}\!ds_2\,\mathcal A(s_1)\mathcal A(s_2)+\cdots
\]
となり、各項は $s_1\ge s_2\ge\cdots$ の順に左から並ぶ。これが経路順序積 $P\exp(i\int\mathcal A\,ds)$ の定義である（級数は $\lVert\mathcal A\rVert$ が有界なら $\sum(\lVert\mathcal A\rVert s)^n/n!$ で抑えられ収束する）。$U(1)$ では $\mathcal A(s)$ は数で互いに可換なので、$n$ 次の項は $\frac{i^n}{n!}(\int_0^s\mathcal A)^n$ となり $U=\exp(i\int_\gamma A)$。

(2) $\tilde U(s)=g(s)U(s)g(0)^{-1}$（$g(s)=g(\gamma(s))$）とおくと
\[
\frac{d\tilde U}{ds}=\dot g\,Ug(0)^{-1}+g\,i\mathcal AUg(0)^{-1}=\bigl(\dot g g^{-1}+i\,g\mathcal Ag^{-1}\bigr)\tilde U .
\]
一方 $A^g$ の引き戻しは $\mathcal A^g=g\mathcal Ag^{-1}+i\,g\,\tfrac{d}{ds}g^{-1}=g\mathcal Ag^{-1}-i\dot gg^{-1}$ なので $i\mathcal A^g=i\,g\mathcal Ag^{-1}+\dot gg^{-1}$。よって $\tilde U$ は $A^g$ に対する方程式 $d\tilde U/ds=i\mathcal A^g\tilde U$ と初期条件 $\tilde U(0)=\mathbf 1$ をみたし、常微分方程式の解の一意性から $U[A^g]=g(\gamma(1))U[A]g(\gamma(0))^{-1}$。

(3) 閉曲線では $\gamma(1)=\gamma(0)=x_0$ なので $U\mapsto g(x_0)Ug(x_0)^{-1}$、トレースの巡回性から $W$ は不変。基点を $C$ 上の別の点 $x_0'$ に移すと、$C$ を $x_0\to x_0'$ の部分 $\alpha$ と $x_0'\to x_0$ の部分 $\beta$ に分けて $U_{x_0}=U(\beta)U(\alpha)$、$U_{x_0'}=U(\alpha)U(\beta)$ となり、トレースは等しい。

$\psi\mapsto g\psi$、$\bar\psi\mapsto\bar\psi g^{-1}$（$\bar\psi=\psi^\dagger\gamma^0$、$g$ はユニタリ）なので
\[
\bar\psi(x_1)U(\gamma)\psi(x_0)\mapsto\bar\psi(x_1)g(x_1)^{-1}\,g(x_1)U(\gamma)g(x_0)^{-1}\,g(x_0)\psi(x_0)=\bar\psi(x_1)U(\gamma)\psi(x_0).
\]
これは離れた点のクォーク・反クォークを「ゲージひも」で結んだ演算子であり、格子ゲージ理論ではリンク変数の積として現れる。
\end{solution}

\begin{problem}[難易度 A]\label{ch08-prob-ab}
半径 $a$ の無限に長いソレノイドの外 $r>a$ で $U(1)$ ゲージ場が $A=\dfrac{\Phi}{2\pi}d\varphi$（円柱座標）であるとする。\\
(1) $r>a$ で $F=0$ であることを確かめ、ソレノイドを1周する閉曲線 $C$ に沿う電荷 $q$ の粒子のホロノミー $e^{iq\oint_CA}$ を求めよ。\\
(2) $q\Phi\in2\pi\Z$ のとき、$r>a$ の領域で $A$ は（電荷 $q$ の場に対する）$A=d\beta$（$e^{iq\beta}$ は一価）と純ゲージの形に書けることを示せ。$q\Phi\notin2\pi\Z$ のときは書けないことを示せ。\\
(3) 定理 \ref{thm:ch08-flat} の言葉で (1), (2) を言いかえ、AB 効果の干渉縞のずれが何で決まるかを述べよ。
\end{problem}

\begin{solution}
(1) $d\varphi$ は閉形式（局所的には関数 $\varphi$ の微分）なので $F=dA=\frac{\Phi}{2\pi}d(d\varphi)=0$。$\oint_Cd\varphi=2\pi$（$C$ が $z$ 軸を正の向きに1周）より $\oint_CA=\Phi$、ホロノミーは $e^{iq\Phi}$。ソレノイド内部の磁束が $\Phi$ であることと、ストークスの定理 $\oint_CA=\int_SF$ は整合している（$S$ はソレノイドを横切る）。

(2) 電荷 $q$ の場 $\psi$ に対するゲージ変換 $\psi\mapsto e^{iq\beta}\psi$ のもとで $A\mapsto A+d\beta$（記法と規約で $g=e^{iq\beta}$、$qA\mapsto qA+i\,g\,dg^{-1}=qA+q\,d\beta$）。$A=d\beta$ となるのは $\beta=\frac{\Phi}{2\pi}\varphi+\text{定数}$ のときに限る（$r>a$ は連結で、$d(\beta-\frac{\Phi}{2\pi}\varphi)=0$ だから）。$e^{iq\beta}=e^{iq\Phi\varphi/2\pi}$ が $\varphi\to\varphi+2\pi$ で一価であることは $q\Phi\in2\pi\Z$ と同値。したがって $q\Phi\in2\pi\Z$ なら $A$ は（電荷 $q$ の粒子にとって）ゲージ変換で $0$ にでき、そうでなければできない。

(3) $r>a$ の領域は $\R^2\setminus D^2$（$\times\R$）で、$\pi_1=\Z$。平坦接続のゲージ同値類は $\Hom(\Z,U(1))=U(1)$、すなわち生成元のホロノミー $e^{iq\Phi}$ で完全に分類される。$e^{iq\Phi}=1$ なら自明な類（純ゲージ）で、そうでなければ $F=0$ でも自明でない。二重スリット実験で、ソレノイドの両側を通る2つの経路の振幅の相対位相は $e^{iq\oint_CA}=e^{iq\Phi}$ だけずれ、干渉縞は $q\Phi$ を $2\pi$ を法として決まる量だけずれる。観測されるのはゲージ不変なホロノミーであり、局所的な場の強さ（$F=0$）ではない。
\end{solution}

\begin{problem}[難易度 B]\label{ch08-prob-monopole}
$S^2$ を $U_N=S^2\setminus\{\text{南極}\}$、$U_S=S^2\setminus\{\text{北極}\}$ で覆い、実数 $n$ に対して
\[
A_N=\frac n2(1-\cos\theta)\,d\varphi,\qquad A_S=-\frac n2(1+\cos\theta)\,d\varphi
\]
とおく（$\theta$ は北極からの角、$S^2$ の向きは $d\theta\wedge d\varphi$、すなわち外向き法線）。\\
(1) $A_N$ が北極で、$A_S$ が南極で滑らかであることを説明し、$F=dA_N=dA_S$ と $\int_{S^2}F$ を求めよ。\\
(2) 緯度 $\theta_0$ の円 $C$（$\varphi$ 増加の向き）に沿うホロノミー（電荷 $1$）を $A_N$ と $A_S$ で別々に計算し、北側の帽子の立体角 $\Omega=2\pi(1-\cos\theta_0)$ で表せ。両者が一致するための条件を求めよ。\\
(3) $U_N\cap U_S$ 上で $A_N=A_S+i\,g\,dg^{-1}$ となる推移関数 $g$ を求め、$n\in\Z$ のとき $c_1=\frac1{2\pi}\int F$ が $g\colon S^1\to U(1)$ の巻き付き数に等しいことを確かめよ。
\end{problem}

\begin{solution}
(1) 北極付近で $1-\cos\theta\approx\theta^2/2$ なので $A_N\approx\frac n4\theta^2d\varphi=\frac n4(x\,dy-y\,dx)$（$x=\theta\cos\varphi$, $y=\theta\sin\varphi$ は北極付近の直交座標）で滑らか。$A_S$ も南極付近で $1+\cos\theta\approx(\pi-\theta)^2/2$ により同様。一方 $A_N$ は南極で $-n\,d\varphi$ に近づき、$d\varphi$ が特異なので滑らかでない（ディラック・ストリング）。
\[
F=dA_N=\frac n2\sin\theta\,d\theta\wedge d\varphi=dA_S,\qquad\int_{S^2}F=\frac n2\cdot4\pi=2\pi n .
\]

(2) $\oint_CA_N=\frac n2(1-\cos\theta_0)\cdot2\pi=\frac n2\Omega$。$\oint_CA_S=-\frac n2(1+\cos\theta_0)\cdot2\pi=\frac n2\Omega-2\pi n$。これはストークスの定理とも整合する：$C$ は北側の帽子 $D_N$ の正の向きの境界、南側の帽子 $D_S$ の負の向きの境界なので、$\oint_CA_N=\int_{D_N}F=\frac n2\Omega$、$\oint_CA_S=-\int_{D_S}F=-\frac n2(4\pi-\Omega)$。ホロノミーは
\[
e^{i\oint A_N}=e^{in\Omega/2},\qquad e^{i\oint A_S}=e^{in\Omega/2}e^{-2\pi in}
\]
であり、一致する（粒子がどちらのパッチで記述されても物理が同じ）ための条件は $e^{2\pi in}=1$、すなわち $n\in\Z$。これがホロノミーから見たディラックの量子化条件である。

(3) $A_N-A_S=n\,d\varphi$。$g=e^{in\varphi}$ とすると $i\,g\,dg^{-1}=i\cdot(-in)\,d\varphi=n\,d\varphi$。よって $g_{NS}=e^{in\varphi}$（$\psi_N=e^{in\varphi}\psi_S$）で、$n\in\Z$ のときに限り一価。赤道 $C$ に対し
\[
\int_{S^2}F=\int_{D_N}dA_N+\int_{D_S}dA_S=\oint_CA_N-\oint_CA_S=\oint_Cn\,d\varphi=2\pi n,
\]
すなわち $c_1=n$ は $g_{NS}\colon S^1\to U(1)$, $\varphi\mapsto e^{in\varphi}$ の巻き付き数 $\frac1{2\pi i}\oint g^{-1}dg=n$ に等しい。第7章の「$S^2$ 上の $U(1)$ 束は $\pi_1(U(1))=\Z$ で分類される」ことの微分形式による表現である。
\end{solution}

\begin{problem}[難易度 B]\label{ch08-prob-torus}
2次元トーラス $T^2=\R^2/\Z^2$（座標 $(x,y)\sim(x+1,y)\sim(x,y+1)$）上の平坦接続を考える。定理 \ref{thm:ch08-flat} を用いてよい。\\
(1) $G=U(1)$ のとき、平坦接続のゲージ同値類が $A=a\,dx+b\,dy$（$a,b$ 定数）で代表され、$(a,b)\sim(a+2\pi m,b+2\pi n)$（$m,n\in\Z$）が同値類を与えること、したがってモジュライ空間がトーラス $T^2$ であることを示せ。どのようなゲージ変換が $(a,b)$ をずらすか。\\
(2) $G=SU(2)$ のとき、互いに可換な $U,V\in SU(2)$ は同時に $\exp(i\alpha\sigma^3)$ の形に共役できることを示し、平坦接続のモジュライ空間が $(\alpha,\beta)\in T^2$ を $(\alpha,\beta)\sim(-\alpha,-\beta)$ で割った空間（「枕カバー」）であることを示せ。固定点はいくつあり、どのようなホロノミーに対応するか。\\
(3) 物理的な意味：コンパクト化された次元に沿ったウィルソン線の真空期待値（ホソトニ機構）や、トーラス上のゲージ理論の古典的真空として、(1), (2) を解釈せよ。
\end{problem}

\begin{solution}
$\pi_1(T^2)=\Z^2$ は生成元 $\gamma_x$（$x$ 方向に1周）と $\gamma_y$ をもつ可換群であり、$\Hom(\Z^2,G)=\{(U,V)\in G^2:UV=VU\}$。

(1) $U(1)$ では $\Hom(\Z^2,U(1))=U(1)\times U(1)$、共役作用は自明。$A=a\,dx+b\,dy$ は平坦で、ホロノミーは $(e^{ia},e^{ib})$。定理 \ref{thm:ch08-flat} により、平坦接続の類はホロノミーの組で決まるので、任意の平坦接続はある $(a,b)$ の $A$ とゲージ同値であり、$(a,b)$ と $(a',b')$ が同値 $\iff$ $e^{ia}=e^{ia'}$ かつ $e^{ib}=e^{ib'}$ $\iff$ $(a'-a,b'-b)\in2\pi\Z^2$。モジュライ空間は $\R^2/2\pi\Z^2\cong T^2$。

（直接の議論：$F=dA=0$ で、$a=\oint_{\gamma_x}A$, $b=\oint_{\gamma_y}A$ とおけば $A-a\,dx-b\,dy$ は周期がすべて $0$ の閉1形式なので完全形式 $d\beta$ であり（第5章）、ゲージ変換 $e^{-i\beta}$ で消せる。）

$(a,b)$ をずらすのは「大きなゲージ変換」$g=e^{2\pi i(mx+ny)}$ で、$T^2$ 上一価であり、$i\,g\,dg^{-1}=2\pi(m\,dx+n\,dy)$。これは $T^2\to U(1)$ の写像のホモトピー類 $[T^2,S^1]\cong\Z^2$ の非自明な元であり、連続的に恒等変換に変形できない。

(2) $U$ はユニタリなので対角化でき、固有値は $e^{\pm i\alpha}$。(i) $U\ne\pm\mathbf 1$ なら固有値は異なり、$UV=VU$ から $V$ は $U$ の各固有空間を保つので、$U$ を対角化する基底で $V$ も対角。(ii) $U=\pm\mathbf 1$ なら $V$ を対角化すればよい（$U$ はどの基底でも対角）。いずれも基底変換は $SU(2)$ でとれる（ユニタリ行列に位相をかけて行列式を $1$ にする）。よって $(U,V)=h(e^{i\alpha\sigma^3},e^{i\beta\sigma^3})h^{-1}$。

対角形どうしの同値性：$\mathrm{diag}$ の組が $h$ で共役なら、固有値の組 $\{e^{\pm i\alpha}\}$ が保たれ、同時固有ベクトルの対応から $h$ は対角行列（$(\alpha,\beta)$ 不変）か $w=i\sigma^2$ 倍の対角行列（$w\sigma^3w^{-1}=-\sigma^3$ により $(\alpha,\beta)\mapsto(-\alpha,-\beta)$）に限られる（ワイル群 $\Z_2$）。$\alpha$, $\beta$ は $2\pi$ 周期なので
\[
\mathcal M_{\text{flat}}(T^2,SU(2))=\bigl(\R^2/2\pi\Z^2\bigr)\big/\bigl((\alpha,\beta)\sim(-\alpha,-\beta)\bigr).
\]
対応する接続は $A=(\alpha\,dx+\beta\,dy)\sigma^3$（ホロノミー $P\exp(i\oint A)=e^{i\alpha\sigma^3}$ など）。$\Z_2$ 作用の固定点は $(\alpha,\beta)\equiv(-\alpha,-\beta)$、すなわち $\alpha,\beta\in\{0,\pi\}$ の4点で、ホロノミーが中心 $\{\pm\mathbf 1\}$ に値をとる4つの平坦接続である（$U,V\in\{\pm\mathbf 1\}$）。商空間は4つの角をもつ球面（枕カバー）と同相である。

(3) 時空 $\R^{d}\times T^2$ でコンパクトな方向に沿う定数ゲージ場 $A_x$, $A_y$ は、$F=0$ なので古典的エネルギーをもたず、その値はホロノミー（ウィルソン線）を通してのみ物理に効く。(1) の周期性は、大きなゲージ変換が $A_x$ を $2\pi$ ずらすことを表し、ウィルソン線の期待値の空間がトーラスになる（弦理論の T 双対性ではこの $(a,b)$ が双対トーラス上の D ブレーンの位置になる）。(2) では、$e^{i\alpha\sigma^3}$ が $\pm\mathbf 1$ でないときゲージ群は $U(1)$ に破れる（ウィルソン線による対称性の破れ、ホソトニ機構）。固定点ではホロノミーが中心にあり $SU(2)$ が破れない。
\end{solution}

\subsection{ベリー位相}

\begin{problem}[難易度 B]\label{ch08-prob-berry}
(1) 時間に依存するハミルトニアン $H(\lambda(t))$ のもとで、断熱近似の状態を $\lvert\psi(t)\rangle=c(t)\lvert n(\lambda(t))\rangle$ と書く。シュレーディンガー方程式から $c(t)=e^{-i\int E_n dt}e^{i\int\mathcal A}$（$\mathcal A=i\langle n|dn\rangle$）を導き、$\mathcal A$ が実であること、$c$ のうち幾何学的な部分が $(d-i\mathcal A)c=0$（記法と規約の平行移動）をみたすことを示せ。\\
(2) 磁場中のスピン $1/2$、$H=-\mu\bm B\cdot\bm\sigma$（$\mu>0$）の基底状態は $\hat{\bm n}\cdot\bm\sigma\lvert+\rangle=\lvert+\rangle$（$\hat{\bm n}=\bm B/\lvert\bm B\rvert=(\sin\theta\cos\varphi,\sin\theta\sin\varphi,\cos\theta)$）をみたす。
\[
\lvert+\rangle_N=\begin{pmatrix}\cos\frac\theta2\\ e^{i\varphi}\sin\frac\theta2\end{pmatrix}
\]
についてベリー接続 $\mathcal A_N$ と曲率 $\mathcal F$ を求め、$\hat{\bm n}$ が緯度 $\theta_0$ の円を $\varphi$ の増える向きに1周したときのベリー位相を、囲む立体角 $\Omega$ で表せ。\\
(3) $\mathcal F$ を $\bm B$ 空間の2形式として書き、それが原点に置かれた磁荷 $-\frac12$（$\frac1{2\pi}\int_{S^2}\mathcal F=-1$）の単極子の場であることを示せ。南極で滑らかなゲージ $\lvert+\rangle_S$ を求め、問題 \ref{ch08-prob-monopole} と比較せよ。\\
(4) 励起状態 $\lvert-\rangle$ のベリー位相はどうなるか。2つの準位のベリー曲率の和について述べよ。
\end{problem}

\begin{solution}
(1) $i\partial_t\lvert\psi\rangle=H\lvert\psi\rangle$ に代入し、断熱近似（他の準位への遷移を無視）で $\langle n|$ を作用させると
\[
i\dot c+i c\langle n|\dot n\rangle=E_nc\quad\Longrightarrow\quad\dot c=-iE_nc+i\bigl(i\langle n|\dot n\rangle\bigr)c .
\]
よって $c(t)=e^{-i\int_0^tE_ndt'}e^{i\int\mathcal A}c(0)$、$\mathcal A=i\langle n|dn\rangle$。$\langle n|n\rangle=1$ を微分して $\langle dn|n\rangle+\langle n|dn\rangle=0$、つまり $\langle n|dn\rangle$ は純虚数なので $\mathcal A$ は実。動的位相を除いた $\tilde c=e^{i\int\mathcal A}$ は $d\tilde c=i\mathcal A\tilde c$、すなわち $(d-i\mathcal A)\tilde c=0$ をみたし、記法と規約の $D=d-iA$ による平行移動である。閉曲線を1周したときの位相 $\gamma=\oint\mathcal A$ はホロノミー $e^{i\oint\mathcal A}$ を与える。

(2) $\lvert+\rangle_N$ を微分すると $d\lvert+\rangle_N=\bigl(-\tfrac12\sin\tfrac\theta2\,d\theta,\ e^{i\varphi}(\tfrac12\cos\tfrac\theta2\,d\theta+i\sin\tfrac\theta2\,d\varphi)\bigr)^T$。よって
\[
\langle+|d+\rangle=-\tfrac12\sin\tfrac\theta2\cos\tfrac\theta2\,d\theta+\tfrac12\sin\tfrac\theta2\cos\tfrac\theta2\,d\theta+i\sin^2\tfrac\theta2\,d\varphi=i\sin^2\tfrac\theta2\,d\varphi,
\]
\[
\mathcal A_N=-\sin^2\tfrac\theta2\,d\varphi=-\tfrac12(1-\cos\theta)\,d\varphi,\qquad\mathcal F=d\mathcal A_N=-\tfrac12\sin\theta\,d\theta\wedge d\varphi .
\]
$\theta=0$ では $\lvert+\rangle_N=(1,0)^T$ で $\varphi$ によらず、$\mathcal A_N$ も滑らか。緯度 $\theta_0$ の円では
\[
\gamma=\oint\mathcal A_N=-\pi(1-\cos\theta_0)=-\frac{\Omega}{2},\qquad\Omega=2\pi(1-\cos\theta_0).
\]
（ストークスにより $\gamma=\int_{\text{北の帽子}}\mathcal F=-\frac12\Omega$ と一般の閉曲線に拡張される。）例えば赤道を1周すると $\gamma=-\pi$ で、状態の符号が反転する。

(3) $\bm B$ 空間（原点を除く）では $\hat{\bm n}=\bm B/\lvert\bm B\rvert$ を通して
\[
\mathcal F=-\frac12\,\frac{\epsilon_{ijk}B_i\,dB_j\wedge dB_k}{2\lvert\bm B\rvert^3},
\]
（単位球面に制限すると $\epsilon_{ijk}n_idn_j\wedge dn_k=2\sin\theta\,d\theta\wedge d\varphi$ だから一致）。これは $\mathcal F=\frac12\mathcal F_{ij}dB_i\wedge dB_j$、$\mathcal F_{ij}=\epsilon_{ijk}\mathcal B_k$ と書いたとき「磁場」$\bm{\mathcal B}=-\frac12\bm B/\lvert\bm B\rvert^3$、すなわち原点の磁荷 $-\frac12$ の単極子場であり、全束 $\int_{S^2}\mathcal F=-2\pi$。単極子の位置はエネルギー準位の縮退点 $\bm B=0$ である。

$\lvert+\rangle_N$ は南極で $(0,e^{i\varphi})^T$ となり $\varphi$ に依存するので特異。$\lvert+\rangle_S=e^{-i\varphi}\lvert+\rangle_N=(e^{-i\varphi}\cos\frac\theta2,\sin\frac\theta2)^T$ は南極で $(0,1)^T$ と滑らか。$\lvert n\rangle\mapsto e^{i\alpha}\lvert n\rangle$ で $\mathcal A\mapsto\mathcal A-d\alpha$ なので（$\alpha=-\varphi$）
\[
\mathcal A_S=\mathcal A_N+d\varphi=\tfrac12(1+\cos\theta)\,d\varphi .
\]
これは問題 \ref{ch08-prob-monopole} の $A_N$, $A_S$ で $n=-1$ としたものに一致し、$c_1=-1$。

(4) $\lvert-\rangle_N=(-e^{-i\varphi}\sin\frac\theta2,\cos\frac\theta2)^T$（北極で滑らか）について同様に $\langle-|d-\rangle=-i\sin^2\frac\theta2\,d\varphi$、$\mathcal A^{(-)}_N=+\frac12(1-\cos\theta)d\varphi$、ベリー位相 $+\Omega/2$、チャーン数 $+1$。2準位のベリー曲率の和は $0$ である。これは $\lvert+\rangle$, $\lvert-\rangle$ の張る束が自明な束 $S^2\times\C^2$ であり、$c_1$ の和（ホイットニー和公式）が $c_1(\text{自明束})=0$ になることの反映である。
\end{solution}

\begin{problem}[難易度 C]\label{ch08-prob-spinj}
スピン $j$ の表現で、$\hat{\bm n}(\theta,\varphi)$ 方向を向いた状態（スピン・コヒーレント状態）を
\[
\lvert\hat{\bm n};m\rangle=e^{-i\varphi J_z}e^{-i\theta J_y}\lvert j,m\rangle\qquad(m=-j,\dots,j)
\]
とおく。$\hat{\bm n}\cdot\bm J\lvert\hat{\bm n};m\rangle=m\lvert\hat{\bm n};m\rangle$ を認めてよい。\\
(1) $\mathcal A=i\langle\hat{\bm n};m|d|\hat{\bm n};m\rangle=m\cos\theta\,d\varphi$ を示せ。\\
(2) $\lvert\hat{\bm n};m\rangle$ は半整数 $j$ では $\varphi\mapsto\varphi+2\pi$ で一価でなく、また北極でも滑らかでない。$\lvert\hat{\bm n};m\rangle_N=e^{im\varphi}\lvert\hat{\bm n};m\rangle$ が北極の近傍で一価かつ滑らかであることを示し、$\mathcal A_N$ を求めよ。\\
(3) 状態 $\lvert\hat{\bm n};m\rangle$ のなす直線束のチャーン数が $-2m$、特に最高ウェイト状態 $m=j$ では $-2j$ であることを示し、全ての $m$ について和をとると $0$ になる理由を述べよ。\\
(4) スピンの経路積分では、作用にベリー位相 $-j\,\Omega[\hat{\bm n}(t)]$ が現れる（$\Omega$ は $\hat{\bm n}(t)$ の描く閉曲線が囲む立体角）。立体角は北側の領域と南側の領域のどちらで測るかで $4\pi$ だけ異なる。$e^{-ij\Omega}$ が well-defined であるための条件を求め、ディラックの量子化条件と比較せよ。
\end{problem}

\begin{solution}
(1) $\lvert\hat{\bm n};m\rangle=R\lvert j,m\rangle$、$R=e^{-i\varphi J_z}e^{-i\theta J_y}$ とすると
\[
d\lvert\hat{\bm n};m\rangle=-i\,d\varphi\,J_zR\lvert j,m\rangle+e^{-i\varphi J_z}(-i\,d\theta\,J_y)e^{-i\theta J_y}\lvert j,m\rangle .
\]
$\langle\hat{\bm n};m|$ をかけると
\[
\langle\hat{\bm n};m|d|\hat{\bm n};m\rangle=-i\,d\varphi\,\langle j,m|e^{i\theta J_y}J_ze^{-i\theta J_y}|j,m\rangle-i\,d\theta\,\langle j,m|J_y|j,m\rangle .
\]
$\langle J_y\rangle=0$（$J_y=(J_+-J_-)/2i$ は $m$ を変える）。$e^{i\theta J_y}J_ze^{-i\theta J_y}=J_z\cos\theta-J_x\sin\theta$（$y$ 軸まわりの回転、符号は $\langle J_x\rangle=0$ なので結果に効かない）より期待値は $m\cos\theta$。よって $\langle d\rangle=-im\cos\theta\,d\varphi$、$\mathcal A=m\cos\theta\,d\varphi$。

(2) $e^{-2\pi iJ_z}\lvert j,m'\rangle=e^{-2\pi im'}\lvert j,m'\rangle=(-1)^{2j}\lvert j,m'\rangle$（$m'-j\in\Z$）なので、$\varphi\mapsto\varphi+2\pi$ で $\lvert\hat{\bm n};m\rangle\mapsto(-1)^{2j}\lvert\hat{\bm n};m\rangle$。さらに $\theta=0$ で $\lvert\hat{\bm n};m\rangle=e^{-im\varphi}\lvert j,m\rangle$ は $\varphi$ に依存し、北極（$\varphi$ が定義されない点）で連続でない。$e^{im\varphi}$ をかけると、$\theta=0$ で $\lvert j,m\rangle$ となり $\varphi$ によらず、また $e^{2\pi im}(-1)^{2j}=e^{2\pi i(m-j)}e^{2\pi ij}(-1)^{2j}=1$（$e^{2\pi ij}=(-1)^{2j}$）なので一価。このとき $\mathcal A\mapsto\mathcal A-d(m\varphi)$ だから
\[
\mathcal A_N=-m(1-\cos\theta)\,d\varphi .
\]
同様に $\lvert\hat{\bm n};m\rangle_S=e^{-im\varphi}\lvert\hat{\bm n};m\rangle$ は南極付近で一価かつ滑らかである。実際 $\theta=\pi$ では $e^{-i\pi J_y}\lvert j,m\rangle=c\lvert j,-m\rangle$（$c$ は定数の位相）、$e^{-i\varphi J_z}\lvert j,-m\rangle=e^{im\varphi}\lvert j,-m\rangle$ なので $\lvert\hat{\bm n};m\rangle_S=c\lvert j,-m\rangle$ は $\varphi$ によらない。このとき $\mathcal A_S=\mathcal A+d(m\varphi)=m(1+\cos\theta)\,d\varphi$。

(3) $\mathcal F=d\mathcal A_N=-m\sin\theta\,d\theta\wedge d\varphi$、$\frac1{2\pi}\int_{S^2}\mathcal F=-\frac{m}{2\pi}\cdot4\pi=-2m$。貼り合わせは $\lvert\cdot\rangle_N=e^{2im\varphi}\lvert\cdot\rangle_S$ で、状態の係数（断面の成分）は $c_N=e^{-2im\varphi}c_S$ と変換する。推移関数 $e^{-2im\varphi}$ の巻き付き数 $-2m$ がチャーン数に一致する（$\mathcal A_N-\mathcal A_S=-2m\,d\varphi$、問題 \ref{ch08-prob-monopole} で $n=-2m$ の場合）。$m=j$ では $-2j$。$j=\frac12$ では問題 \ref{ch08-prob-berry} の $-1$ に一致する。

$2j+1$ 個の直線束の直和は、各点で $\{\lvert\hat{\bm n};m\rangle\}_m$ が張る空間、すなわちスピン $j$ 表現空間全体 $\C^{2j+1}$ であり、$S^2\times\C^{2j+1}$ という自明な束である。ホイットニー和公式より $\sum_mc_1=c_1(\text{自明束})=0$。実際 $\sum_{m=-j}^j(-2m)=0$。

(4) 2つの測り方の差は $\Omega_N-(-\Omega_S)=\Omega_N+\Omega_S=4\pi$（南側の領域は逆向きに囲まれる）なので、$e^{-ij\Omega}$ の不定性は $e^{-4\pi ij}$。これが $1$ となる条件は $2j\in\Z$。すなわちスピンの大きさの量子化は、「磁荷 $-j$（全磁束 $-4\pi j$）の単極子の場の中を運動する電荷 $1$ の粒子」に対するディラックの量子化条件 $\int F\in2\pi\Z$、すなわち $2j\in\Z$ と同じものである（(3) のチャーン数 $-2j$ が整数であること）。
\end{solution}

\begin{problem}[難易度 B]\label{ch08-prob-hopf}
$\CP^1$ の点 $[z_1:z_2]$ に複素直線 $\ell=\C(z_1,z_2)\subset\C^2$ を対応させる直線束 $L=\{([z],v):v\in\ell\}$ を\textbf{トートロジカル直線束}という。$\C^2$ の内積から誘導されるエルミート計量のもとで、単位長さの局所切断 $s$ に対し $\mathcal A=i\,s^\dagger ds$ とおく（$\C^2$ の自明な接続を $L$ に射影したもの）。\\
(1) $w=z_2/z_1$ を座標とするパッチで $s_N(w)=(1,w)/\sqrt{1+\lvert w\rvert^2}$ とする。$\mathcal A_N$ と $\mathcal F=d\mathcal A_N$ を $w=\varrho e^{i\varphi}$ で表し、$\CP^1$ の複素構造による向き（$w$ 平面で $d(\mathrm{Re}\,w)\wedge d(\mathrm{Im}\,w)$）に関して $\int_{\CP^1}\mathcal F/2\pi=-1$ を示せ。\\
(2) $L$ の単位ベクトル全体が $S^3$ であり、$S^3\to\CP^1\cong S^2$ がホップ写像（ホップ束、第7章）であることを説明せよ。また $w=\tan\frac\theta2e^{i\varphi}$ とおくと、$s_N$ が問題 \ref{ch08-prob-berry} の $\lvert+\rangle_N$ と一致することを確かめよ。\\
(3) 双対束 $L^*=\mathcal O(1)$ と $L^{\otimes n}$ の $c_1$ を求め、磁荷 $n$ の単極子束をこれらで表せ。
\end{problem}

\begin{solution}
(1) $N=1+\lvert w\rvert^2$、$s=N^{-1/2}(1,w)$ とする。
\[
s^\dagger ds=N^{-1/2}(1,\bar w)\cdot\bigl[N^{-1/2}(0,dw)+(1,w)\,d(N^{-1/2})\bigr]=\frac{\bar w\,dw}{N}+N^{1/2}d(N^{-1/2})=\frac{\bar w\,dw}{N}-\frac{dN}{2N}.
\]
$dN=\bar w\,dw+w\,d\bar w$ なので $s^\dagger ds=\dfrac{\bar w\,dw-w\,d\bar w}{2N}$。$w=\varrho e^{i\varphi}$ では $\bar w\,dw-w\,d\bar w=2i\varrho^2d\varphi$ だから
\[
\mathcal A_N=i\cdot\frac{2i\varrho^2d\varphi}{2(1+\varrho^2)}=-\frac{\varrho^2}{1+\varrho^2}\,d\varphi,\qquad\mathcal F=-\frac{2\varrho\,d\varrho\wedge d\varphi}{(1+\varrho^2)^2}.
\]
$d(\mathrm{Re}\,w)\wedge d(\mathrm{Im}\,w)=\varrho\,d\varrho\wedge d\varphi$ なので、この向きで
\[
\int_{\CP^1}\mathcal F=-2\pi\int_0^\infty\frac{2\varrho\,d\varrho}{(1+\varrho^2)^2}=-2\pi\Bigl[-\frac{1}{1+\varrho^2}\Bigr]_0^\infty=-2\pi .
\]
（$w$ のパッチは $\CP^1$ から1点 $[0:1]$ を除いたもので、積分には影響しない。）よって $c_1(L)=-1$。

(2) $L$ の単位ベクトルは $(z_1,z_2)\in\C^2$、$\lvert z_1\rvert^2+\lvert z_2\rvert^2=1$、すなわち $S^3$ の点であり、各ファイバー（直線 $\ell$）の単位円 $\{e^{i\alpha}z\}$ が $S^1$。射影 $S^3\to\CP^1$, $z\mapsto[z]$ を $\CP^1\cong S^2$（ブロッホ球、$[z]\mapsto z^\dagger\bm\sigma z$）と合成したものがホップ写像であり、$L$ は主 $U(1)$ 束 $S^3\to S^2$ に同伴する直線束である。$w=\tan\frac\theta2e^{i\varphi}$ とおくと $N=1/\cos^2\frac\theta2$、$s_N=\cos\frac\theta2(1,\tan\frac\theta2e^{i\varphi})=(\cos\frac\theta2,e^{i\varphi}\sin\frac\theta2)$ で $\lvert+\rangle_N$ に一致する。さらに $\varrho=\tan\frac\theta2$ は $\theta$ の増加関数なので $\varrho\,d\varrho\wedge d\varphi$ と $d\theta\wedge d\varphi$ は同じ向きを与え、$c_1=-1$ はベリー位相の計算（問題 \ref{ch08-prob-berry}）と一致する。$\mathcal A_N$ 自体も $\varrho^2/(1+\varrho^2)=\sin^2\frac\theta2$ により一致している。

(3) 双対束の接続は $-\mathcal A$ で、$c_1(L^*)=+1$。テンソル積では接続が足されるので $c_1(L^{\otimes n})=-n$、$c_1((L^*)^{\otimes n})=n$。したがって磁荷 $n$（$\int F/2\pi=n$）の単極子束は $\mathcal O(n)=(L^*)^{\otimes n}$ と同型であり（$S^2$ 上の直線束は $c_1$ で分類される）、問題 \ref{ch08-prob-monopole} の $A_N$ は $-n\mathcal A_N=n\sin^2\frac\theta2\,d\varphi$ に一致する。ホップ束の主 $U(1)$ 束としての「チャーン数 $\pm1$」は、どちらの表現（電荷 $\pm1$）で同伴束をつくるかによる符号の違いである。
\end{solution}

\subsection{チャーン・ヴェイユ理論と特性類}

\begin{problem}[難易度 B]\label{ch08-prob-transg}
同じ束上の接続 $A_0$, $A_1$ に対し $\eta=A_1-A_0$、$A_t=A_0+t\eta$、$F_t$ をその曲率とする。\\
(1) $\dfrac{dF_t}{dt}=D_t\eta:=d\eta-i(A_t\wedge\eta+\eta\wedge A_t)$ を示せ。\\
(2) $\tr(D_t\eta\wedge F_t)=d\,\tr(\eta\wedge F_t)$ を示せ（ビアンキ恒等式を用いよ）。\\
(3) $\tr F_1\wedge F_1-\tr F_0\wedge F_0=d\Bigl(2\displaystyle\int_0^1\tr(\eta\wedge F_t)\,dt\Bigr)$ を示し、$A_0=0$, $A_1=A$ のとき右辺が $8\pi^2\,dCS(A)$ になることを確かめよ。\\
(4) 閉4次元多様体上の同じ束の2つの接続について $\int\tr F_1^2=\int\tr F_0^2$ を結論し、インスタントン数が場の小さな揺らぎで変わらないことの意味を述べよ。
\end{problem}

\begin{solution}
(1) $F_t=dA_0+t\,d\eta-i(A_0+t\eta)\wedge(A_0+t\eta)$ を $t$ で微分すると
\[
\dot F_t=d\eta-i(A_0\eta+\eta A_0)-2it\,\eta\wedge\eta=d\eta-i(A_t\eta+\eta A_t),
\]
（$A_t\eta+\eta A_t=A_0\eta+\eta A_0+2t\eta^2$）。

(2) 1形式 $\eta$ と2形式 $F$ について $d\,\tr(\eta F)=\tr(d\eta\,F)-\tr(\eta\,dF)$。一方
\[
\tr(D\eta\,F)-\tr(\eta\,DF)=\tr(d\eta F)-\tr(\eta dF)-i\tr(A\eta F+\eta AF)+i\tr(\eta AF-\eta FA).
\]
$\tr(\eta FA)$ では1形式 $A$ を3形式 $\eta F$ の前へ移すと符号 $(-1)^{1\cdot3}=-1$ がつくので $\tr(\eta FA)=-\tr(A\eta F)$。よって余分な項は $-i\tr(A\eta F)-i\tr(\eta FA)=-i\tr(A\eta F)+i\tr(A\eta F)=0$ となり、$\tr(D\eta F)-\tr(\eta DF)=d\tr(\eta F)$。ビアンキ $D_tF_t=0$ より $\tr(D_t\eta\wedge F_t)=d\tr(\eta\wedge F_t)$。

(3) 2形式どうしは可換なので $\frac{d}{dt}\tr F_t^2=\tr(\dot F_tF_t+F_t\dot F_t)=2\tr(D_t\eta\,F_t)=2\,d\tr(\eta F_t)$。$t$ で $0$ から $1$ まで積分すれば主張を得る。$A_0=0$、$\eta=A$ なら $F_t=t\,dA-it^2A^2$ で
\[
2\int_0^1\tr\bigl(A(t\,dA-it^2A^2)\bigr)dt=2\Bigl(\frac12\tr(A\,dA)-\frac i3\tr A^3\Bigr)=\tr\Bigl(A\,dA-\frac{2i}{3}A^3\Bigr)=8\pi^2CS(A).
\]
特に $\tr F^2=8\pi^2dCS(A)$。

(4) $\eta$ は大域的な1形式なので $Q=2\int_0^1\tr(\eta F_t)dt$ も大域的な3形式であり、閉多様体上 $\int dQ=0$。よって $\int\tr F\wedge F$ は接続によらない。場の小さな揺らぎ $A\to A+\delta A$（同じ束の上の変形）でインスタントン数 $k$ は不変であり、摂動論の範囲では異なる $k$ のセクターは移り合わない。$k$ を変えるには束そのもの（貼り合わせ）を変える必要があり、それは連続変形ではできない。なお $CS(A)$ は $A$ が大域的に定義されているときにだけ大域的な形式であり、自明でない束では $\tr F^2=8\pi^2dCS$ はパッチごとにしか成り立たない（だから $\int\tr F^2\ne0$ になりうる）。
\end{solution}

\begin{problem}[難易度 A]\label{ch08-prob-chernpoly}
(1) $\det(1+X)=1+\tr X+\frac12\bigl((\tr X)^2-\tr X^2\bigr)+O(X^3)$ を用いて、$c_1$, $c_2$ を $F$ で表し、$\ch_2=\frac12c_1^2-c_2$ を示せ。\\
(2) $SU(N)$ 束では $c_2=-\ch_2$ であることを示せ。BPST インスタントン（$k=1$）の $\int c_2$ を答えよ。\\
(3) 直線束の直和 $E=L_1\oplus L_2$（$c_1(L_1)=a$, $c_1(L_2)=b$）について $c(E)$ と $\ch(E)$ を求め、(1) の関係を確かめよ。また $\ch(L^{\otimes n})=e^{nc_1(L)}$ を説明せよ。
\end{problem}

\begin{solution}
(1) $X=F/2\pi$ とすると
\[
c_1=\frac{\tr F}{2\pi},\qquad c_2=\frac{1}{8\pi^2}\bigl(\tr F\wedge\tr F-\tr F\wedge F\bigr).
\]
（$\det(1+X)$ の展開：固有値 $x_i$ で $\prod(1+x_i)=1+\sum x_i+\sum_{i<j}x_ix_j+\cdots$、$\sum_{i<j}x_ix_j=\frac12((\sum x_i)^2-\sum x_i^2)$。$F$ は偶数次の形式なので固有値の議論が形式的に使える。）一方 $\ch(F)=\tr e^{X}=N+\tr X+\frac12\tr X^2+\cdots$ で $\ch_2=\frac1{8\pi^2}\tr F^2$。よって $\frac12c_1^2-c_2=\frac1{8\pi^2}\bigl(\tr F\wedge\tr F-(\tr F\wedge\tr F-\tr F^2)\bigr)=\ch_2$。

(2) $\mathfrak{su}(N)$ はトレースゼロなので $c_1=0$、$c_2=-\ch_2$。BPST では $\int\ch_2=k=1$ なので $\int c_2=-1$。

(3) 曲率は $\mathrm{diag}(F_1,F_2)$ なので
\[
c(E)=(1+a)(1+b)=1+(a+b)+ab,\qquad\ch(E)=e^a+e^b=2+(a+b)+\tfrac12(a^2+b^2)+\cdots .
\]
$\frac12c_1^2-c_2=\frac12(a+b)^2-ab=\frac12(a^2+b^2)=\ch_2$ で (1) と一致。$L^{\otimes n}$ では電荷 $n$ の場として接続が $nA$、曲率が $nF$ なので $\ch=e^{nF/2\pi}=e^{nc_1(L)}$（$\ch$ が積を積に移すことの例）。
\end{solution}

\begin{problem}[難易度 B]\label{ch08-prob-cpn}
$\CP^n$ の接束について $c(T\CP^n)=(1+x)^{n+1}$ が成り立つことを認める。ここで $x\in H^2(\CP^n;\Z)$ は $\mathcal O(1)$ の $c_1$ で $\int_{\CP^n}x^n=1$（$x$ の冪 $1,x,\dots,x^n$ がコホモロジーの基底、第6章）。また、複素 $n$ 次元のコンパクト複素多様体 $M$ について $\chi(M)=\int_Mc_n(TM)$（ガウス・ボンネ・チャーン、$c_n(TM)=e(TM)$）を認める。\\
(1) $n=1$ のとき $\int c_1(TS^2)=2=\chi(S^2)$ を確かめよ。\\
(2) $\chi(\CP^n)=n+1$ を示し、胞体分割 $\CP^n=e^0\cup e^2\cup\cdots\cup e^{2n}$ と比較せよ。\\
(3) $\CP^2$ について $p_1=c_1^2-2c_2$ を用いて $p_1(T\CP^2)=3x^2$ を示せ。ヒルツェブルフの符号数定理 $\sigma=\frac13\int p_1$（第9章）が $\CP^2$ の交叉形式 $(+1)$（第6章）と整合することを確かめよ。\\
(4) $c_1(T\CP^n)$ がゼロにならないことから、$\CP^n$ はカラビ・ヤウ多様体（第10章）でないことを述べよ。
\end{problem}

\begin{solution}
(1) $n=1$：$c(TS^2)=(1+x)^2=1+2x$（$x^2=0$）。$\int_{\CP^1}c_1=2\int x=2=\chi(S^2)$。ガウス・ボンネ（問題 \ref{ch08-prob-lc}, \ref{ch08-prob-gb}）とも一致する。

(2) $c_n=\binom{n+1}{n}x^n=(n+1)x^n$ なので $\chi(\CP^n)=(n+1)\int x^n=n+1$。胞体分割では各偶数次元に1つずつ胞体があり $\chi=\sum(-1)^k\#(k\text{ 胞体})=n+1$（第4章）。

(3) $n=2$：$c(T\CP^2)=(1+x)^3=1+3x+3x^2$。$p_1=c_1^2-2c_2=9x^2-6x^2=3x^2$。

（$p_1=c_1^2-2c_2$ の理由：実束としての $T\CP^2$ の複素化は $T\oplus\bar T$、$c(\bar T)=1-c_1+c_2-\cdots$ なので $c(T\otimes\C)=(1+c_1+c_2)(1-c_1+c_2)=1+(2c_2-c_1^2)+\cdots$、$p_1=-c_2(T\otimes\C)=c_1^2-2c_2$。）

よって $\frac13\int p_1=1$。一方 $H^2(\CP^2)=\Z x$、$\int x^2=1$ なので交叉形式は $(+1)$、符号数 $\sigma=1$ で一致する。また $\chi=3$ は (2) と一致し、$\chi=b_0+b_2+b_4=1+1+1$。

(4) カラビ・ヤウ多様体は $c_1(TM)=0$（リッチ平坦計量をもつ）ケーラー多様体である。$c_1(T\CP^n)=(n+1)x\ne0$ なので $\CP^n$ はカラビ・ヤウではない（実際 $c_1>0$ の「ファノ多様体」で、正のリッチ曲率をもつフビニ・スタディ計量をもつ）。カラビ・ヤウ多様体は $\CP^{n}$ の中の超曲面として作られる（第10章）。
\end{solution}

\subsection{曲率と重力：ガウス・ボンネ}

\begin{problem}[難易度 C]\label{ch08-prob-lc}
単位球面 $S^2$（計量 $d\theta^2+\sin^2\theta\,d\varphi^2$）上で正規直交枠 $e^1=d\theta$, $e^2=\sin\theta\,d\varphi$（双対ベクトルは $\bm e_1=\partial_\theta$, $\bm e_2=\frac1{\sin\theta}\partial_\varphi$）をとる。\\
(1) ねじれなしの条件 $de^a+\omega^a{}_b\wedge e^b=0$ からレヴィ・チヴィタ接続 $\omega^1{}_2$ を求め、$R^1{}_2=d\omega^1{}_2=K\,e^1\wedge e^2$ により $K=1$ を確かめよ。\\
(2) 接ベクトル $v=v^1\bm e_1+v^2\bm e_2$ の平行移動 $dv^a+\omega^a{}_bv^b=0$ を複素数 $v=v^1+iv^2$ で書くと $dv=iA\,v$、$A=\omega^1{}_2$ となることを示せ（$TS^2$ を複素直線束とみなし、記法と規約の $D=d-iA$ を用いている）。\\
(3) 緯度 $\theta_0$ の円を $\varphi$ の増える向きに1周する平行移動で、接ベクトルは（枠 $\bm e_1,\bm e_2$ に対して）どれだけ回転するか。それが北側の帽子の $\int K\,dS$ に等しいことを示し、緯度 $\lambda$ の地点でのフーコーの振り子の1恒星日あたりの回転角を求めよ。\\
(4) 北極・南極で滑らかな枠への取りかえにより $A_N$, $A_S$ を求め、推移関数の巻き付き数が $2$ であること、すなわち $\int_{S^2}e(TS^2)=c_1(TS^2)=2=\chi(S^2)$ を示せ。
\end{problem}

\begin{solution}
(1) $de^1=0$、$de^2=\cos\theta\,d\theta\wedge d\varphi$。$\omega^1{}_1=\omega^2{}_2=0$、$\omega^2{}_1=-\omega^1{}_2$。条件は
\[
de^1=-\omega^1{}_2\wedge e^2,\qquad de^2=-\omega^2{}_1\wedge e^1=\omega^1{}_2\wedge d\theta .
\]
$\omega^1{}_2=-\cos\theta\,d\varphi$ とおくと、第2式は $-\cos\theta\,d\varphi\wedge d\theta=\cos\theta\,d\theta\wedge d\varphi$ ✓、第1式は $\cos\theta\sin\theta\,d\varphi\wedge d\varphi=0$ ✓。（ねじれなしの解は一意。）$R^1{}_2=d\omega^1{}_2+\omega^1{}_c\wedge\omega^c{}_2=d\omega^1{}_2=\sin\theta\,d\theta\wedge d\varphi=e^1\wedge e^2$。よって $K=1$。

(2) $dv^1=-\omega^1{}_2v^2$、$dv^2=-\omega^2{}_1v^1=\omega^1{}_2v^1$ より
\[
d(v^1+iv^2)=-\omega^1{}_2v^2+i\omega^1{}_2v^1=i\omega^1{}_2(v^1+iv^2).
\]
よって $dv=iAv$、$A=\omega^1{}_2=-\cos\theta\,d\varphi$（$i$ をかけることは $\bm e_1\mapsto\bm e_2$ の $90^\circ$ 回転で、外向き法線について反時計回り）。曲率 $F=dA=R^1{}_2=K\,dS$。

(3) $\theta=\theta_0$ 上で $dv/d\varphi=-i\cos\theta_0\,v$ なので $v(2\pi)=e^{-2\pi i\cos\theta_0}v(0)=e^{2\pi i(1-\cos\theta_0)}v(0)$。すなわち枠に対して角度
\[
\Delta=2\pi(1-\cos\theta_0)\ (\mathrm{mod}\ 2\pi)
\]
だけ反時計回り（外から見て）に回転する。$\int_{\text{北の帽子}}K\,dS=\int_0^{2\pi}\!\int_0^{\theta_0}\sin\theta\,d\theta\,d\varphi=2\pi(1-\cos\theta_0)$ に等しい。（枠 $\bm e_1,\bm e_2$ は緯線上で一価なので「枠に対する回転」は意味をもつ。）

フーコーの振り子：振動面は（地球の自転に比べ速い振動の平均として）地表に沿って平行移動される。緯度 $\lambda$（北緯）では余緯度 $\theta_0=\frac\pi2-\lambda$、1恒星日で振り子は緯線を1周するので、地表の枠（南北・東西）に対して振動面は $2\pi(1-\sin\lambda)\equiv-2\pi\sin\lambda$ だけ回る。すなわち北半球では上から見て時計回りに1日 $2\pi\sin\lambda$（北極で $2\pi$、赤道で $0$）。

(4) 枠を回す（$v\mapsto e^{i\alpha}v$）と $A\mapsto A+d\alpha$（$g=e^{i\alpha}$、$i\,g\,dg^{-1}=d\alpha$）。北極付近で滑らかな枠は、$(\bm e_1,\bm e_2)$ を $-\varphi$ だけ回したもの（北極で $\varphi$ によらない向き）で、$v_N=e^{i\varphi}v$。よって $A_N=-\cos\theta\,d\varphi+d\varphi=(1-\cos\theta)d\varphi$（$\theta\to0$ で $\theta^2d\varphi/2$、滑らか）。南極では $v_S=e^{-i\varphi}v$、$A_S=-(1+\cos\theta)d\varphi$。$v_N=e^{2i\varphi}v_S$ で推移関数 $e^{2i\varphi}$ の巻き付き数は $2$。問題 \ref{ch08-prob-monopole} で $n=2$ の場合であり
\[
\int_{S^2}\frac{F}{2\pi}=\frac1{2\pi}\int_{S^2}K\,dS=\frac{4\pi}{2\pi}=2=\chi(S^2).
\]
「毛玉の定理」（$S^2$ 上にゼロ点のないベクトル場は存在しない）は、この束が自明でない（$e\ne0$）ことの言いかえである。
\end{solution}

\begin{problem}[難易度 A]\label{ch08-prob-gb}
ガウス・ボンネの定理 $\chi(\Sigma)=\frac1{2\pi}\int_\Sigma K\,dS$ を認める。\\
(1) 半径 $a$ の球面で右辺を計算し、$a$ によらないことを確かめよ。\\
(2) 種数 $g\ge2$ の閉曲面が $K\equiv-1$ の計量をもつとき、その面積を求めよ。\\
(3) 2次元の計量 $h_{\alpha\beta}$ のスカラー曲率は $R=2K$ である。弦の世界面の作用の項 $S_\lambda=\frac{\lambda}{4\pi}\int d^2\sigma\sqrt h\,R$ が $\lambda\chi$ に等しいことを示し、経路積分の重み $e^{-S_\lambda}$ が種数 $g$ の閉世界面に $g_s^{2g-2}$（$g_s=e^\lambda$）を与えることを示せ。\\
(4) 2次元のアインシュタイン方程式が自明（アインシュタイン・テンソルが恒等的にゼロ）であることを、(3) を用いて説明せよ。
\end{problem}

\begin{solution}
(1) $K=1/a^2$、面積 $4\pi a^2$ で $\frac1{2\pi}\cdot\frac{4\pi a^2}{a^2}=2=\chi(S^2)$。

(2) $\chi=2-2g=\frac1{2\pi}\int(-1)dS=-\frac{\text{面積}}{2\pi}$ より、面積 $=4\pi(g-1)$。（双曲幾何の面積は位相で決まる。）

(3) $\sqrt h\,d^2\sigma=dS$、$R=2K$ なので $\frac{1}{4\pi}\int\sqrt hR\,d^2\sigma=\frac1{2\pi}\int K\,dS=\chi$。よって $S_\lambda=\lambda\chi$、$e^{-S_\lambda}=e^{-\lambda\chi}=g_s^{-\chi}=g_s^{2g-2}$。球面（$g=0$）が $g_s^{-2}$、トーラスが $g_s^0$、と種数ごとに $g_s^2$ ずつ抑えられ、弦の摂動展開が種数展開になる（境界・クロスキャップがあれば $\chi$ はそれも数える。第10章）。

(4) $S_\lambda$ は計量の局所的な変化 $h\to h+\delta h$ で変わらない（$\chi$ は整数で連続変化しない）ので $\delta S_\lambda=0$。一方、一般に $\delta\int\sqrt hR=\int\sqrt h\,(R_{\alpha\beta}-\frac12Rh_{\alpha\beta})\delta h^{\alpha\beta}$（境界項を除く）なので、2次元ではアインシュタイン・テンソル $G_{\alpha\beta}$ が恒等的にゼロ（実際 $R_{\alpha\beta}=\frac12Rh_{\alpha\beta}$）。2次元重力は局所的な力学をもたず、位相（オイラー数）だけを感じる。
\end{solution}

\subsection{インスタントン・チャーン・サイモンズ・θ 真空}

\begin{problem}[難易度 C]\label{ch08-prob-bpst1}
ユークリッド $\R^4$（$x^2=x^\mu x^\mu$）上の $SU(2)$ ゲージ場
\[
A_\mu=A_\mu^aT^a,\qquad A^a_\mu=\frac{2\eta^a_{\mu\nu}x^\nu}{x^2+\rho^2},\qquad T^a=\frac{\sigma^a}{2}
\]
を考える（$\eta^a_{\mu\nu}$ は例 \ref{ex:ch08-bpst} のトフーフト記号）。恒等式
\[
\epsilon^{abc}\eta^a_{\mu\lambda}\eta^b_{\nu\sigma}=\delta_{\mu\nu}\eta^c_{\lambda\sigma}-\delta_{\mu\sigma}\eta^c_{\lambda\nu}-\delta_{\lambda\nu}\eta^c_{\mu\sigma}+\delta_{\lambda\sigma}\eta^c_{\mu\nu}
\]
を認めてよい。\\
(1) 記法と規約のもとで $F^c_{\mu\nu}=\partial_\mu A^c_\nu-\partial_\nu A^c_\mu+\epsilon^{abc}A^a_\mu A^b_\nu$ となることを示せ。\\
(2) $F^c_{\mu\nu}=-\dfrac{4\rho^2\eta^c_{\mu\nu}}{(x^2+\rho^2)^2}$ を示し、$F$ が自己双対 $F=*F$ であることを結論せよ。\\
(3) $A$ が原点で滑らかであること、$\lvert x\rvert\to\infty$ で $A_\mu=O(1/\lvert x\rvert)$、$F_{\mu\nu}=O(1/\lvert x\rvert^4)$ であることを確かめよ。$A$ が遠方で純ゲージ $i\,g\,dg^{-1}$、$g=(x^4\mathbf 1+ix^a\sigma^a)/\lvert x\rvert$ に近づくことを認めて、純ゲージなのに $A=O(1/\lvert x\rvert)$ で遠方の $F$ は速く減衰する理由を述べよ。
\end{problem}

\begin{solution}
(1) $-i[A_\mu,A_\nu]=-iA^a_\mu A^b_\nu[T^a,T^b]=-iA^a_\mu A^b_\nu\,i\epsilon^{abc}T^c=\epsilon^{abc}A^a_\mu A^b_\nu T^c$。

(2) $h=2/(x^2+\rho^2)$ とおくと $A^c_\nu=h\,\eta^c_{\nu\lambda}x^\lambda$、$\partial_\mu h=-h^2x_\mu$（添字の上下は区別しない）。
\[
\partial_\mu A^c_\nu=h\,\eta^c_{\nu\mu}-h^2x_\mu\eta^c_{\nu\lambda}x_\lambda .
\]
$\eta$ の反対称性から
\[
\partial_\mu A^c_\nu-\partial_\nu A^c_\mu=-2h\,\eta^c_{\mu\nu}-h^2\bigl(x_\mu\eta^c_{\nu\lambda}-x_\nu\eta^c_{\mu\lambda}\bigr)x_\lambda .
\]
非可換項は与えられた恒等式を $x_\lambda x_\sigma$ で縮約して（$\eta^c_{\lambda\sigma}x_\lambda x_\sigma=0$）
\[
\epsilon^{abc}A^a_\mu A^b_\nu=h^2\bigl(-x_\mu\eta^c_{\lambda\nu}x_\lambda-x_\nu\eta^c_{\mu\sigma}x_\sigma+x^2\eta^c_{\mu\nu}\bigr)=h^2\bigl(x_\mu\eta^c_{\nu\lambda}-x_\nu\eta^c_{\mu\lambda}\bigr)x_\lambda+h^2x^2\eta^c_{\mu\nu}.
\]
和をとると $x_\mu x_\lambda$ を含む項はちょうど打ち消し合い
\[
F^c_{\mu\nu}=\eta^c_{\mu\nu}\bigl(-2h+h^2x^2\bigr)=\eta^c_{\mu\nu}\Bigl(-\frac{4}{x^2+\rho^2}+\frac{4x^2}{(x^2+\rho^2)^2}\Bigr)=-\frac{4\rho^2\eta^c_{\mu\nu}}{(x^2+\rho^2)^2}.
\]
$\eta^c$ は自己双対（$\frac12\epsilon_{\mu\nu\rho\sigma}\eta^c_{\rho\sigma}=\eta^c_{\mu\nu}$）なので $(*F)_{\mu\nu}=\frac12\epsilon_{\mu\nu\rho\sigma}F_{\rho\sigma}=F_{\mu\nu}$。

\begin{note}
自己双対性の確認：例えば $\eta^a_{12}=\epsilon_{a12}=\delta_{a3}$、$\eta^a_{34}=\delta_{a3}$ で $\frac12(\epsilon_{1234}\eta_{34}+\epsilon_{1243}\eta_{43})=\eta_{34}=\eta_{12}$。他の成分も同様。$\eta^a_{\mu\nu}\eta^a_{\mu\nu}=12$（各 $a$ について非ゼロ成分が4つ）。
\end{note}

(3) 分母 $x^2+\rho^2\ge\rho^2>0$ なので $A$ は $\R^4$ 全体で滑らか（原点付近で $A^a_\mu\approx2\eta^a_{\mu\nu}x^\nu/\rho^2$）。遠方で $A\sim2\eta x/x^2=O(1/\lvert x\rvert)$、$F\sim4\rho^2/x^4$。純ゲージ $i\,g\,dg^{-1}$ は $g$ が角度方向だけに依存するので $dg=O(1/\lvert x\rvert)$ であり、$A$ がこの形に $O(1/\lvert x\rvert^3)$ の速さで近づく（$\frac{1}{x^2+\rho^2}=\frac1{x^2}-\frac{\rho^2}{x^4}+\cdots$）。純ゲージの曲率はゼロなので、$F$ はずれ $O(1/\lvert x\rvert^3)$ の微分程度、$O(1/\lvert x\rvert^4)$ になる。作用 $\int F^2d^4x$ の有限性に必要なのは $F$ の減衰だけで、$A$ 自身はゼロに近づく必要はなく、遠方の $g\colon S^3_\infty\to SU(2)$ の巻き付き数が位相的電荷になる（第3章）。
\end{solution}

\begin{problem}[難易度 C]\label{ch08-prob-bpst2}
問題 \ref{ch08-prob-bpst1} の BPST インスタントンについて考える。ヤン・ミルズ作用は、結合定数を $A$ に吸収した記法で
\[
S_E=\frac{1}{2g_{\mathrm{YM}}^2}\int d^4x\,\tr F_{\mu\nu}F_{\mu\nu}=\frac{1}{g_{\mathrm{YM}}^2}\int\tr F\wedge*F
\]
である。\\
(1) 正準規格化した場 $A^{\mathrm{can}}=A/g_{\mathrm{YM}}$ で書くと $S_E=\frac14\int(F^{\mathrm{can},a}_{\mu\nu})^2d^4x$ となることを確かめよ。\\
(2) ボゴモルニー不等式 $S_E\ge\dfrac{8\pi^2}{g_{\mathrm{YM}}^2}\lvert k\rvert$ を示し、等号が（反）自己双対で成り立つことを示せ。\\
(3) BPST インスタントンの $k$ を $\tr F\wedge F$ の積分で計算し、$k=1$、$S_E=8\pi^2/g_{\mathrm{YM}}^2$ を示せ。\\
(4) 同じ $k$ を、定理 \ref{thm:ch08-csgauge} と向きの注意（基本事項のまとめ）を用いて $S^3_\infty$ 上の積分 $k=\int_{S^3_\infty}CS(A)$ から求め、一致を確かめよ。\\
(5) 解は $\rho$ と中心 $x_0$ と大域的なゲージの向き（$SU(2)$）をパラメータとしてもつ。$k=1$ の解のモジュライ空間の次元を数えよ（$SU(2)$ の中心 $\pm\mathbf 1$ は $A$ に作用しないことに注意）。
\end{problem}

\begin{solution}
(1) $F=dA-iA\wedge A=g_{\mathrm{YM}}(dA^{\mathrm{can}}-ig_{\mathrm{YM}}A^{\mathrm{can}}\wedge A^{\mathrm{can}})=g_{\mathrm{YM}}F^{\mathrm{can}}$。$\tr T^aT^b=\frac12\delta^{ab}$ より $\tr F_{\mu\nu}F_{\mu\nu}=\frac{g_{\mathrm{YM}}^2}{2}F^{\mathrm{can},a}_{\mu\nu}F^{\mathrm{can},a}_{\mu\nu}$、よって $S_E=\frac14\int(F^{\mathrm{can},a}_{\mu\nu})^2$。また $F\wedge*F=\frac12F_{\mu\nu}F_{\mu\nu}d^4x$ から2つ目の表示が従う。

(2) ユークリッド4次元の2形式では $**=1$ なので
\[
0\le\int\tr(F\mp*F)\wedge*(F\mp*F)=2\int\tr F\wedge*F\mp2\int\tr F\wedge F
\]
（$*F\wedge**F=*F\wedge F=F\wedge*F$、$F\wedge*(*F)=F\wedge F$）。左辺の非負性は $\tr(G\wedge*G)=\frac12\tr G_{\mu\nu}G_{\mu\nu}d^4x$、$G$ エルミートで $\tr G_{\mu\nu}^2\ge0$ による。よって $\int\tr F\wedge*F\ge\pm8\pi^2k$、すなわち $S_E\ge8\pi^2\lvert k\rvert/g_{\mathrm{YM}}^2$。等号は $F=\pm*F$ のとき。

(3) $\tr F\wedge F=\frac14\epsilon_{\mu\nu\rho\sigma}\tr(F_{\mu\nu}F_{\rho\sigma})d^4x$。自己双対なら $\frac12\epsilon_{\mu\nu\rho\sigma}F_{\rho\sigma}=F_{\mu\nu}$ なので $=\frac12\tr F_{\mu\nu}F_{\mu\nu}\,d^4x$。
\[
\tr F_{\mu\nu}F_{\mu\nu}=\frac12F^c_{\mu\nu}F^c_{\mu\nu}=\frac12\cdot\frac{16\rho^4}{(x^2+\rho^2)^4}\cdot12=\frac{96\rho^4}{(x^2+\rho^2)^4}.
\]
$d^4x=2\pi^2r^3dr$（$S^3$ の体積 $2\pi^2$）と
\[
\int_0^\infty\frac{r^3dr}{(r^2+\rho^2)^4}=\frac12\int_0^\infty\frac{u\,du}{(u+\rho^2)^4}=\frac12\Bigl(\frac{1}{2\rho^4}-\frac{\rho^2}{3\rho^6}\Bigr)=\frac{1}{12\rho^4}
\]
より
\[
\int\tr F\wedge F=\frac12\cdot96\rho^4\cdot2\pi^2\cdot\frac1{12\rho^4}=8\pi^2,\qquad k=1 .
\]
$\rho$ によらない。自己双対なので (2) の等号が成り立ち $S_E=8\pi^2/g_{\mathrm{YM}}^2$。トンネル振幅は $e^{-8\pi^2/g_{\mathrm{YM}}^2}$ で、摂動論では見えない。

(4) $A$ は $\R^4$ 全体で滑らかなので $CS(A)$ は大域的な3形式で、半径 $R$ の球 $B_R$ にストークスを用いると $\frac1{8\pi^2}\int_{B_R}\tr F^2=\int_{S^3_R}CS(A)$。$R\to\infty$ で $A\to i\,g\,dg^{-1}$（ずれは $O(R^{-3})$ で、$S^3_R$ 上の積分への寄与は $R\to\infty$ で消える）なので、定理 \ref{thm:ch08-csgauge} より $k=\int_{S^3_\infty}CS(i\,g\,dg^{-1})=-w(g)$。向きの注意により $g=(x^4+ix^a\sigma^a)/\lvert x\rvert$ は境界の向きで $w(g)=-1$（北極で $g^{-1}dg=i\sigma^adx^a$、$\tr(i\sigma^adx^a)^3=-i\cdot2i\epsilon^{abc}dx^adx^bdx^c=12\,dx^{123}$、左不変性から $\frac{1}{24\pi^2}\cdot12\cdot(-2\pi^2)=-1$）。よって $k=1$ で (3) と一致する。

(5) 中心4、大きさ1、ゲージの向き $SU(2)/\{\pm\mathbf1\}=SO(3)$ の3、合計 $4+1+3=8$。（指数定理（第9章）により、$SU(2)$ の $k$ インスタントンのモジュライ空間は $8k$ 次元であることが示される。）大きさ $\rho$ が任意であるのは古典的なスケール不変性の反映で、量子論では結合定数の走りにより $\rho$ 積分が赤外発散する。
\end{solution}

\begin{problem}[難易度 D]\label{ch08-prob-csgauge}
定理 \ref{thm:ch08-csgauge} を証明する。$\theta=g^{-1}dg$、$B=A-i\theta$ とおく。\\
(1) $A^g=gBg^{-1}$、$d\theta=-\theta\wedge\theta$、$dA^g=g(\theta B+dB+B\theta)g^{-1}$ を示せ。\\
(2) $8\pi^2\,CS(A^g)=\tr(B\,dB)+2\tr(B^2\theta)-\frac{2i}{3}\tr B^3$ を示せ。\\
(3) $B=A-i\theta$ を代入して
\[
8\pi^2\bigl(CS(A^g)-CS(A)\bigr)=-\frac13\tr\theta^3+i\,d\,\tr(\theta\wedge A)
\]
を示せ。\\
(4) 閉3次元多様体上で $\int CS(A^g)=\int CS(A)-w(g)$ を結論し、純ゲージ $i\,g\,dg^{-1}$ のチャーン・サイモンズ数が $-w(g)$ であることを直接（$F=0$ を用いて）確かめよ。
\end{problem}

\begin{solution}
(1) $i\,g\,dg^{-1}=-i\,dg\,g^{-1}=g(-i\theta)g^{-1}$ より $A^g=g(A-i\theta)g^{-1}=gBg^{-1}$。$d\theta=d(g^{-1})\wedge dg=-g^{-1}dg\,g^{-1}\wedge dg=-\theta\wedge\theta$。$dg=g\theta$、$dg^{-1}=-\theta g^{-1}$ を用いて
\[
d(gBg^{-1})=dg\wedge Bg^{-1}+g\,dB\,g^{-1}-gB\wedge dg^{-1}=g(\theta B+dB+B\theta)g^{-1}.
\]

(2) $\tr(A^g\,dA^g)=\tr\bigl(B(\theta B+dB+B\theta)\bigr)=\tr(B\theta B)+\tr(B\,dB)+\tr(B^2\theta)$。$\tr(B\theta B)$ で最後の1形式 $B$ を2形式 $B\theta$ の前に移しても符号は $(-1)^{1\cdot2}=+1$ なので $\tr(B\theta B)=\tr(B^2\theta)$。$\tr(A^g)^3=\tr B^3$。よって主張を得る。

(3) 以下 $\wedge$ を省略し、奇数次の形式のトレースの巡回性（1形式を2形式の反対側へ移すときは符号なし）を用いる。$dB=dA-i\,d\theta=dA+i\theta^2$。
\begin{align*}
\tr(B\,dB)&=\tr(A\,dA)+i\tr(A\theta^2)-i\tr(\theta\,dA)+\tr\theta^3,\\
2\tr(B^2\theta)&=2\tr(A^2\theta)-4i\tr(A\theta^2)-2\tr\theta^3,\\
-\tfrac{2i}{3}\tr B^3&=-\tfrac{2i}{3}\tr A^3-2\tr(A^2\theta)+2i\tr(A\theta^2)+\tfrac23\tr\theta^3 .
\end{align*}
\begin{note}
2行目：$B^2=A^2-i(A\theta+\theta A)-\theta^2$、$\tr((A\theta+\theta A)\theta)=2\tr(A\theta^2)$。3行目：$\tr B^3=\tr A^3-3i\tr(A^2\theta)-3\tr(A\theta^2)+i\tr\theta^3$（$\theta$ を1個含む3つの項 $\tr(AA\theta)$, $\tr(A\theta A)$, $\tr(\theta AA)$ は巡回性で等しい。2個含む項も同様。$(-i)^2=-1$, $(-i)^3=i$）。
\end{note}
和をとると、$\tr(A\,dA)-\frac{2i}3\tr A^3=8\pi^2CS(A)$ のほかに、$\tr\theta^3$ の係数 $1-2+\frac23=-\frac13$、$\tr(A^2\theta)$ の係数 $0$、$\tr(A\theta^2)$ の係数 $i-4i+2i=-i$ が残り
\[
8\pi^2\bigl(CS(A^g)-CS(A)\bigr)=-\tfrac13\tr\theta^3-i\tr(\theta\,dA)-i\tr(A\theta^2).
\]
一方 $d\tr(\theta A)=\tr(d\theta\,A)-\tr(\theta\,dA)=-\tr(\theta^2A)-\tr(\theta\,dA)=-\tr(A\theta^2)-\tr(\theta\,dA)$。よって最後の2項は $i\,d\tr(\theta A)$ に等しく、主張を得る。

(4) 閉多様体上では完全形式の積分はゼロなので
\[
\int CS(A^g)-\int CS(A)=-\frac{1}{24\pi^2}\int\tr(g^{-1}dg)^3=-w(g).
\]
直接の確認：$A=i\,g\,dg^{-1}=-i\,dg\,g^{-1}$ は $F=0$、すなわち $dA=iA^2$ をみたすので
\[
8\pi^2CS(A)=\tr\bigl(iA^3-\tfrac{2i}3A^3\bigr)=\tfrac i3\tr A^3,\qquad A^3=(-i)^3(dg\,g^{-1})^3=i(dg\,g^{-1})^3 .
\]
$\tr(dg\,g^{-1})^3=\tr(g^{-1}dg)^3$（0形式 $g^{-1}$ を先頭へ巡回）なので $8\pi^2CS=\frac i3\cdot i\tr\theta^3=-\frac13\tr\theta^3$、$\int CS(i\,g\,dg^{-1})=-w(g)$。これは $A=0$ に対して (3) を適用したものと一致する。
\end{solution}

\begin{problem}[難易度 C]\label{ch08-prob-theta}
$\R^4=\R_\tau\times\R^3$ で $A_\tau=0$（テンポラル・ゲージ）をとり、各時刻で $\lvert\bm x\rvert\to\infty$ のとき $A\to0$ が十分速いとする（空間を $S^3$ とみなせる）。$\R^3$ の向きは $d\tau\wedge(\R^3\text{ の向き})$ が $\R^4$ の向きになるようにとり、$N_{CS}(\tau)=\int_{\R^3}CS(A(\tau))$ とおく。\\
(1) $\tau\to\pm\infty$ で $A$ が真空（純ゲージ）に近づく有限作用の配位について $k=N_{CS}(+\infty)-N_{CS}(-\infty)$ を示せ。\\
(2) 古典的真空 $A=i\,g\,dg^{-1}$（$g\colon S^3\to G$）を $N_{CS}=n$ で分類し、量子状態 $\lvert n\rangle$ とする。$w(h)=1$ の大きなゲージ変換 $h$ を表す演算子 $\mathcal T$ が $\mathcal T\lvert n\rangle=\lvert n-1\rangle$ と作用することを示し、$\mathcal T$ の固有状態 $\lvert\theta\rangle=\sum_ne^{-in\theta}\lvert n\rangle$ の固有値を求めよ。\\
(3) $\langle\theta'|e^{-HT}|\theta\rangle$ を $k$ セクターの経路積分 $Z_k$ で表し、θ 真空では各 $k$ に重み $e^{i\theta k}$ がつくこと、これはユークリッド作用に $-i\theta k=-\frac{i\theta}{8\pi^2}\int\tr F\wedge F$ を加えることと同じであることを示せ。$\theta\sim\theta+2\pi$ の理由を述べよ。\\
(4) 2次元 $U(1)$ ゲージ理論の θ 項 $\frac{\theta}{2\pi}\int_\Sigma F$ と、円周上の粒子 $L=\frac I2\dot\varphi^2+\frac{\theta}{2\pi}\dot\varphi$ について、同じ周期性 $\theta\sim\theta+2\pi$ を説明せよ。後者ではエネルギー準位を求めよ。
\end{problem}

\begin{solution}
(1) $[\tau_1,\tau_2]\times\R^3$ の向き $d\tau\wedge(\R^3)$ に関して、ストークスの定理の境界は $\{\tau_2\}\times\R^3-\{\tau_1\}\times\R^3$（外向き法線 $\pm\partial_\tau$ を先頭に置く規約）であり、空間無限遠からの寄与は仮定によりない。$\tr F^2=8\pi^2dCS(A)$ で、テンポラル・ゲージでは $CS(A)$ の時間一定面への制限は空間成分だけのチャーン・サイモンズ形式なので
\[
\frac1{8\pi^2}\int_{[\tau_1,\tau_2]\times\R^3}\tr F\wedge F=N_{CS}(\tau_2)-N_{CS}(\tau_1).
\]
$\tau_{1,2}\to\mp\infty$ として $k=N_{CS}(+\infty)-N_{CS}(-\infty)$。両端は真空なので $N_{CS}(\pm\infty)$ は整数（(2)）、その差が $k$ である。

(2) 定理 \ref{thm:ch08-csgauge} より $N_{CS}[i\,g\,dg^{-1}]=-w(g)\in\Z$。ゲージ変換 $h$ は $A\mapsto A^h$ で、$i\,g\,dg^{-1}\mapsto i(hg)d(hg)^{-1}$、$N_{CS}[A^h]=N_{CS}[A]-w(h)$（$w(hg)=w(h)+w(g)$ と整合）。よって $w(h)=1$ なら $\mathcal T\lvert n\rangle=\lvert n-1\rangle$。
\[
\mathcal T\lvert\theta\rangle=\sum_ne^{-in\theta}\lvert n-1\rangle=\sum_me^{-i(m+1)\theta}\lvert m\rangle=e^{-i\theta}\lvert\theta\rangle .
\]
小さなゲージ変換（ガウスの法則で生成される）は物理状態に自明に作用するが、大きなゲージ変換は位相 $e^{-i\theta}$ までしか固定されない。$\mathcal T$ はハミルトニアンと可換なので $\theta$ は超選択則のラベルであり、異なる $\theta$ は別の理論とみなせる。

(3) $\langle m|e^{-HT}|n\rangle=Z_{m-n}$（境界条件 $n\to m$ の経路積分、(1) により $k=m-n$ のセクター）。
\[
\langle\theta'|e^{-HT}|\theta\rangle=\sum_{m,n}e^{im\theta'}e^{-in\theta}Z_{m-n}=\sum_ne^{in(\theta'-\theta)}\sum_ke^{ik\theta'}Z_k\ \propto\ \delta(\theta-\theta')\sum_ke^{ik\theta}Z_k .
\]
θ 真空の間では各 $k$ セクターに重み $e^{i\theta k}$ がつき、これは $e^{-S_E}$ の $S_E$ に $-i\theta k=-\frac{i\theta}{8\pi^2}\int\tr F\wedge F$ を加えることと同じ。$k\in\Z$ なので $e^{i(\theta+2\pi)k}=e^{i\theta k}$ で $\theta\sim\theta+2\pi$。（QCD では $\theta$ は CP を破り、中性子の電気双極子モーメントの上限から $\lvert\theta\rvert\lesssim10^{-10}$。これが強い CP 問題である。）

(4) 2次元：閉曲面上 $\int_\Sigma F\in2\pi\Z$（ディラックの量子化、記法と規約）なので $\frac{\theta}{2\pi}\int F=\theta n$、$e^{i\theta n}$ は $\theta\mapsto\theta+2\pi$ で不変。円周上の粒子：$\frac{\theta}{2\pi}\int\dot\varphi\,dt$ は周期的な（ユークリッド）時間の経路の巻き付き数 $\nu$ の $\theta$ 倍で、同じ理由で周期的。正準運動量 $p=I\dot\varphi+\frac{\theta}{2\pi}$、波動関数 $e^{ip\varphi}$ の一価性から $p=\ell\in\Z$、
\[
H=\frac{1}{2I}\Bigl(p-\frac{\theta}{2\pi}\Bigr)^2,\qquad E_\ell=\frac{1}{2I}\Bigl(\ell-\frac{\theta}{2\pi}\Bigr)^2 .
\]
$\theta\mapsto\theta+2\pi$ は $\ell\mapsto\ell+1$ と付け替えれば同じスペクトル。これは AB 磁束 $\Phi=\theta$ の円環上の粒子と同じであり（問題 \ref{ch08-prob-ab}）、θ 項は「配位空間の非自明なループを貫く磁束」とみなせる。ヤン・ミルズでは配位空間（ゲージ同値類）の非自明なループが大きなゲージ変換であり、その $\pi_1$ が $\pi_3(G)=\Z$ に対応する。
\end{solution}

\begin{problem}[難易度 D]\label{ch08-prob-cslevel}
閉3次元多様体 $M$ 上の $SU(N)$ チャーン・サイモンズ理論
\[
S_\kappa[A]=\frac{\kappa}{4\pi}\int_M\tr\Bigl(A\wedge dA-\frac{2i}{3}A\wedge A\wedge A\Bigr)
\]
を考える（$\tr$ は基本表現、$\tr T^aT^b=\frac12\delta^{ab}$、$A$ は大域的に定義されているとする）。\\
(1) $S_\kappa=2\pi\kappa\int_MCS(A)$ を確かめ、ゲージ変換 $g\colon M\to SU(N)$ のもとで $S_\kappa\mapsto S_\kappa-2\pi\kappa\,w(g)$ を示せ。\\
(2) $M=S^3$ で $w(g)=1$ となる $g$ が存在することを示し、$e^{iS_\kappa}$ のゲージ不変性から $\kappa\in\Z$ を導け。\\
(3) 4次元多様体 $X$ で $\partial X=M$、$A$ が $X$ 上の接続に拡張されるとき、$S_\kappa=2\pi\kappa\int_X\ch_2(F)$ と書けることを示せ。2つの拡張 $X$, $X'$ による値の差が $2\pi\kappa\int_{Y}\ch_2$（$Y=X\cup_M(-X')$ は閉4次元多様体）であることから、この定義で $e^{iS_\kappa}$ が well-defined であるための条件を求め、(2) と比較せよ。\\
(4) $U(1)$ チャーン・サイモンズ $S=\frac{\kappa}{4\pi}\int A\wedge dA$ について (3) と同じ議論を行うと、$\kappa$ に対する条件はどうなるか。閉4次元多様体の交叉形式が偶（スピン多様体）の場合とそうでない場合に分けて答えよ。
\end{problem}

\begin{solution}
(1) $\frac{\kappa}{4\pi}\cdot8\pi^2=2\pi\kappa$ より $S_\kappa=2\pi\kappa\int CS(A)$。定理 \ref{thm:ch08-csgauge}（問題 \ref{ch08-prob-csgauge}）より、閉多様体上 $\int CS(A^g)=\int CS(A)-w(g)$ なので $S_\kappa[A^g]=S_\kappa[A]-2\pi\kappa\,w(g)$。$w(g)$ は整数である（$g^*$ による $H^3(SU(N);\Z)\cong\Z$ の生成元の引き戻しの積分、$M=S^3$ なら写像度）。

(2) $SU(2)\subset SU(N)$ を左上の $2\times2$ ブロックに埋め込み、$S^3\cong SU(2)$（向きの規約により $w=1$ となる同一視）を $g$ とすれば $w(g)=1$。（$\tr$ は基本表現で、埋め込んだ $SU(2)$ の基本表現と $N-2$ 個の自明表現の和に分解するので、$\tr(g^{-1}dg)^3$ は $SU(2)$ のものと同じ。）すると $e^{iS_\kappa}\mapsto e^{iS_\kappa}e^{-2\pi i\kappa}$。経路積分の被積分関数がゲージ不変であるには $e^{-2\pi i\kappa}=1$、すなわち $\kappa\in\Z$。

(3) $A$ が $X$ 上に（大域的に）拡張されるならストークスと $dCS=\ch_2$ により $\int_X\ch_2(F)=\int_MCS(A)$、よって $S_\kappa=2\pi\kappa\int_X\ch_2$。この右辺は束が自明でなくても意味をもつので、一般にはこれを定義とする（$SU(N)$ の場合、3次元多様体上の束は常に自明で、拡張 $X$ も常に存在する（$\Omega_3^{SO}(BSU(N))=0$）という事実を認める）。2つの拡張の差は、$X$ と向きを逆にした $X'$ を $M$ で貼り合わせた閉多様体 $Y$ 上の積分
\[
2\pi\kappa\int_X\ch_2-2\pi\kappa\int_{X'}\ch_2=2\pi\kappa\int_Y\ch_2(F)=-2\pi\kappa\int_Yc_2 .
\]
$c_2$ は整数類なので $\int_Yc_2\in\Z$。逆に $Y=S^4$、$k=1$ の束（インスタントン）で $\int_Y\ch_2=1$ となる。したがって $e^{iS_\kappa}$ が拡張によらない条件は $\kappa\in\Z$ で、(2) と同じ結論を、ゲージ変換を経由せず、自明でない束も含めて得る。（(2) とのつながり：$M=S^3$ 上の $A$ と $A^g$ を $S^3\times I$ の両端に置き、補間する接続を考えて両端を $g$ で貼り合わせると、$S^3\times S^1$ 上の束ができ、その $\int\ch_2$ が $\pm w(g)$ である。）

(4) $U(1)$：$S=\frac{\kappa}{4\pi}\int_XF\wedge F$。$F/2\pi=c_1$ を用いると
\[
\frac{\kappa}{4\pi}\int_YF\wedge F=\frac{\kappa}{4\pi}(2\pi)^2\int_Yc_1^2=\pi\kappa\int_Yc_1^2 .
\]
$\int_Yc_1^2=Q(c_1,c_1)$ は交叉形式（第6章）の値で整数。$e^{iS}$ が well-defined であるには、すべての閉 $Y$ と束について $\kappa\,Q(c_1,c_1)\in2\mathbb Z$。一般の $Y$（例えば $\CP^2$、$Q=(+1)$ で $c_1^2=1$ となる束がある）を許すと $\kappa\in2\Z$。$Y$ をスピン多様体に限る（理論がフェルミオンを含みスピン構造に依存する）と、交叉形式は偶で $Q(c_1,c_1)\in2\Z$ なので $\kappa\in\Z$ でよい。$U(1)$ では $\kappa$ 奇数のチャーン・サイモンズ理論は「スピン TQFT」であり、これは $\ch_2=\frac12c_1^2$ の $\frac12$ が整数性を壊すことに由来する。$SU(N)$ では $c_1=0$ で $\ch_2=-c_2$ が常に整数だったので、この区別は生じない。
\end{solution}

%================================================
\chapter{指数定理とアノマリー}\label{ch:09}
% 18問
%================================================

\section{物理からの動機}

\subsection*{ゼロモードは数えられる}
インスタントンやモノポールのような位相的に非自明な背景場の中にフェルミオンを置くと、ディラック方程式 $\slashed{D}\psi=0$ の規格化可能な解（\term{ゼロモード}{zero mode}）が現れる。ゼロモードの数そのものは背景場を少し変えると変わりうるが、その\textbf{カイラリティごとの差}
\[
\ind\slashed D=n_+-n_-
\]
は連続変形で変わらない。しかもそれは背景場のトポロジー（第8章の特性類）だけで決まる：
\[
\ind\slashed D=\int_M\ch(F)\,\Ahat(TM)\qquad(\text{アティヤ・シンガーの指数定理}).
\]
左辺は解析的（微分方程式の解の数）、右辺は位相的（曲率の積分）な量である。この等式が、本章で扱う物理の多くの現象の数学的な核心になる。

\subsection*{量子異常（アノマリー）}
古典的な対称性が量子論で破れる現象を\term{アノマリー}{anomaly}という。代表例は質量ゼロのフェルミオンの軸性 $U(1)$ 対称性 $\psi\mapsto e^{i\alpha\gamma_5}\psi$ である。古典的には軸性電荷 $Q_5$ は保存するが、量子論では
\[
\Delta Q_5=2N_f\,k\qquad(N_f\text{ 個のディラック・フェルミオン、インスタントン数 }k)
\]
だけ変化する（\term{ABJアノマリー}{Adler--Bell--Jackiw anomaly}）。藤川の方法では、これは経路積分の測度のヤコビアンとして現れ、ヤコビアンの指数部分は\textbf{指数密度}そのものである。

1+1次元の例が直観を与える。長さ $L$ の円周上の質量ゼロのフェルミオンに一様な電場をかけ、$A_1$ を $0$ から $2\pi/L$ まで断熱的に変えると、右巻きの準位はすべて1段下がり、左巻きの準位は1段上がる（下図左）。ディラックの海から右巻き（または左巻き）の粒子が1個ずつ押し出され、$N_R-N_L$ が $2$ だけ変化する。このとき時空を貫く磁束 $\int F=\int E\,dt\,dx=2\pi$ は1単位である。準位がゼロを横切る現象を\term{スペクトル流}{spectral flow}といい、スペクトル流の正味の数が指数に等しい。
\begin{center}
\begin{tikzpicture}[scale=1.0]
  \begin{scope}
    \draw[-{Stealth[length=2mm]}] (0,-1.9) -- (0,1.9) node[above] {\footnotesize $E$};
    \draw[-{Stealth[length=2mm]}] (0,0) -- (3.4,0) node[right] {\footnotesize $t$};
    \draw[dashed,gray] (3,-1.8) -- (3,1.8);
    \node[below] at (3,-1.8) {\footnotesize $A_1=2\pi/L$};
    \foreach \m in {-3,-2,-1,0,1,2}{
      \draw[thick,blue!70!black] (0,{0.5*(\m+0.5)}) -- (3,{0.5*(\m-0.5)});
      \draw[thick,red!70!black,densely dashed] (0,{0.5*(\m+0.5)+0.06}) -- (3,{0.5*(\m+1.5)+0.06});
    }
    \foreach \m in {-3,-2,-1}{ \node[vtx,blue!70!black] at (0,{0.5*(\m+0.5)}) {}; }
    \node[blue!70!black] at (1.5,-2.2) {\footnotesize 実線：右巻き　破線：左巻き};
  \end{scope}
  \begin{scope}[xshift=6.3cm]
    \draw[-{Stealth[length=2mm]}] (0,-1.9) -- (0,1.9) node[above] {\footnotesize $\lambda$};
    \draw[-{Stealth[length=2mm]}] (0,0) -- (3.4,0) node[right] {\footnotesize $t$};
    \draw[dashed,gray] (3,-1.8) -- (3,1.8);
    \draw[very thick,red!70!black,domain=0:3,samples=60] plot (\x,{0.75*cos(60*\x)});
    \draw[very thick,red!70!black,densely dashed,domain=0:3,samples=60] plot (\x,{-0.75*cos(60*\x)});
    \foreach \s in {1,-1}{
      \draw[thick,gray,domain=0:3,samples=40] plot (\x,{\s*(1.25+0.15*sin(120*\x))});
      \draw[thick,gray,domain=0:3,samples=40] plot (\x,{\s*(1.65-0.1*sin(60*\x))});
    }
    \node[below] at (0,-1.8) {\footnotesize $A$};
    \node[below] at (3,-1.8) {\footnotesize $A^U$};
    \node[red!70!black] at (1.5,-2.2) {\footnotesize 1対の固有値 $\pm\lambda$ が入れ替わる};
  \end{scope}
\end{tikzpicture}
\zucap{左：1+1次元の円周上で $A_1$ を1量子分ずらすと、右巻きの準位（実線）は1段下がり、左巻きの準位（破線）は1段上がる。黒丸ははじめに占有された準位。右：$SU(2)$ の大きなゲージ変換 $U$（$\pi_4(SU(2))=\Z_2$ の非自明元）に沿ったスペクトル流。奇数対の固有値が交差すると、パフィアンの符号が反転する（ウィッテン・アノマリー）。}
\end{center}

\subsection*{トフーフト頂点}
4次元の $SU(N_c)$ ゲージ理論で $k=1$ のインスタントン背景を考えると、基本表現の質量ゼロのディラック・フェルミオンは1フレーバーあたりちょうど1個のゼロモード（正カイラリティ）をもつ。グラスマン積分ではゼロモードの積分変数が作用に現れないため、ゼロモードを「吸収」する演算子を挿入しない限り経路積分はゼロになる。結果として、インスタントンは $2N_f$ 本のフェルミオン線をもつ実効的な相互作用（\term{トフーフト頂点}{'t Hooft vertex}）を生み、軸性電荷を $2N_f$ だけ破る。
\begin{center}
\begin{tikzpicture}[scale=1.0]
  \shade[ball color=orange!30] (0,0) circle (0.55);
  \node at (0,0) {\footnotesize $k=1$};
  \foreach \a/\lab in {90/u,150/d,210/s}{
    \draw[thick,mid arrow] (\a:1.9) -- (\a:0.58);
    \node at (\a:2.25) {\footnotesize $\psi_{\lab}$};
  }
  \foreach \a/\lab in {30/u,330/d,270/s}{
    \draw[thick,mid arrow] (\a:0.58) -- (\a:1.9);
    \node at (\a:2.25) {\footnotesize $\bar\psi_{\lab}$};
  }
  \begin{scope}[xshift=6.2cm]
    \shade[ball color=orange!30] (0,0) circle (0.55);
    \node at (0,0) {\footnotesize $SU(2)_L$};
    \foreach \i in {0,...,11}{
      \draw[thick,mid arrow] ({30*\i+15}:1.8) -- ({30*\i+15}:0.58);
    }
    \foreach \i/\lab in {0/q,1/q,2/q,3/\ell,4/q,5/q,6/q,7/\ell,8/q,9/q,10/q,11/\ell}{
      \node at ({30*\i+15}:2.1) {\footnotesize $\lab$};
    }
  \end{scope}
\end{tikzpicture}
\zucap{左：$N_f=3$ の QCD のトフーフト頂点。各フレーバーにつき $\psi$ と $\bar\psi$ のゼロモードが1本ずつ現れ、6フェルミオン相互作用になる（軸性電荷を $6$ 破る）。右：電弱 $SU(2)_L$ インスタントンは各左巻き二重項（世代あたりクォーク3色＋レプトン1個）に1本ずつ線を出し、3世代で12本：$\Delta B=\Delta L=3$。}
\end{center}

\subsection*{ゲージ・アノマリーと理論の整合性}
大域的対称性のアノマリーは物理的な効果（$\pi^0\to\gamma\gamma$、$\eta'$ の質量、$B+L$ の破れ）をもたらすが、\textbf{ゲージ}対称性のアノマリーは理論を壊す（ユニタリ性や繰り込み可能性が失われる）。したがってカイラルなゲージ理論では、フェルミオンの表現は\term{アノマリー相殺条件}{anomaly cancellation condition}をみたさなければならない。標準模型の一見でたらめなハイパーチャージ $\frac16,-\frac23,\frac13,-\frac12,1$ は、この条件を見事にみたしている。さらに、摂動論では見えない大域的なアノマリーもある。$SU(2)$ ゲージ理論で左巻きの二重項が奇数個あると理論は矛盾する（ウィッテンの $SU(2)$ アノマリー）。その起源はホモトピー群 $\pi_4(SU(2))=\Z_2$（第3章）である。

アノマリーは「$d$ 次元の理論の問題を $d+2$ 次元の特性類で記述する」という構造をもつ（\term{アノマリー多項式}{anomaly polynomial}と\term{降下方程式}{descent equations}）。これは第8章のチャーン・サイモンズ形式の議論を一段進めたものであり、弦理論のグリーン・シュワルツ機構（第10章）にもつながる。

\begin{physbox}[：本章の辞書]
\begin{center}
\small
\begin{tabular}{p{0.46\linewidth}p{0.44\linewidth}}
物理 & 数学\\ \hline
ゼロモードの数（カイラリティ付き） & フレドホルム指数 $\ind\slashed D=n_+-n_-$\\
ゼロモード数の安定性 & 指数の変形不変性\\
軸性電荷の非保存 $\Delta Q_5$ & $2N_f\times$（指数）\\
藤川のヤコビアン & 熱核の超トレース $\operatorname{Tr}\gamma_5e^{-t\slashed D^2}$\\
最低ランダウ準位の縮退度 & $T^2$ や $S^2$ 上のディラック作用素の指数\\
ゲージ・アノマリー & アノマリー多項式 $I_{d+2}=[\Ahat\ch(F)]_{d+2}$ と降下方程式\\
ウィッテンの $SU(2)$ アノマリー & $\pi_4(SU(2))=\Z_2$、mod 2 のスペクトル流\\
アノマリー流入 & 境界のある多様体上のチャーン・サイモンズ項
\end{tabular}
\end{center}
\end{physbox}

\section{基本事項のまとめ}

\subsection{フレドホルム指数とディラック作用素}

\begin{mydef}{フレドホルム指数}{ch09-fredholm}
ヒルベルト空間の間の（有界または閉）線形作用素 $T\colon H_1\to H_2$ が、$\ker T$ と $\coker T=H_2/\operatorname{im}T$ がともに有限次元で $\operatorname{im}T$ が閉であるとき、\term{フレドホルム作用素}{Fredholm operator}といい、
\[
\ind T=\dim\ker T-\dim\coker T=\dim\ker T-\dim\ker T^\dagger
\]
をその\term{指数}{index}という。
\end{mydef}

\begin{mythm}{指数の変形不変性}{ch09-deform}
フレドホルム作用素の連続な族 $T_s$（$s\in[0,1]$）について、$\ind T_s$ は $s$ によらない。$\dim\ker T_s$ 自体は跳びうるが、跳ぶときは $\ker T_s$ と $\ker T_s^\dagger$ が同じだけ増減する。
\end{mythm}
有限次元では $\ind T=\dim H_1-\dim H_2$（階数・退化次数の定理）であり、$T$ の詳細によらない。指数は無限次元でもこの「次元の差」の頑健さを保つ量である。

\begin{mydef}{ディラック作用素とカイラリティ}{ch09-dirac}
偶数次元 $d=2m$ の向き付けられたコンパクトなスピン多様体 $M$（第7章）上で、エルミートなガンマ行列 $\{\gamma^\mu,\gamma^\nu\}=2g^{\mu\nu}$ と、ゲージ場とスピン接続を含む共変微分 $D_\mu$ を用いて
\[
\slashed D=i\gamma^\mu D_\mu
\]
とおく（エルミート）。カイラリティ演算子 $\gamma_5$（$\gamma_5^2=1$, $\gamma_5^\dagger=\gamma_5$）は $\{\gamma_5,\slashed D\}=0$ をみたす。スピノル束を $S=S^+\oplus S^-$（$\gamma_5=\pm1$）に分けると
\[
\slashed D=\begin{pmatrix}0&\slashed D_-\\ \slashed D_+&0\end{pmatrix},\qquad \slashed D_\pm\colon\Gamma(S^\pm\otimes E)\to\Gamma(S^\mp\otimes E),\quad \slashed D_-=\slashed D_+^\dagger
\]
であり、$n_\pm=\dim\ker\slashed D_\pm$ として
\[
\ind\slashed D:=\ind\slashed D_+=n_+-n_- .
\]
$\gamma_5$ の符号は、指数定理が $\int\ch(F)\Ahat$ の形になるように選ぶ（記法と規約）。具体的には、$D_\mu=\partial_\mu-iA_\mu$、向き $dx^1\wedge dx^2(\wedge dx^3\wedge dx^4)$ のもとで、平坦な2次元では $\gamma^1=\sigma_1$, $\gamma^2=\sigma_2$, $\gamma_5=\sigma_3=-i\gamma^1\gamma^2$、4次元では $\gamma_5=-\gamma^1\gamma^2\gamma^3\gamma^4$ がこれにあたる（問題 \ref{ch09-prob-fujikawa2d}）。
\end{mydef}

\begin{myprop}{非ゼロモードの対と超トレース}{ch09-pair}
$\slashed D\varphi=\lambda\varphi$ なら $\slashed D(\gamma_5\varphi)=-\lambda\gamma_5\varphi$。$\lambda\neq0$ の固有空間では、$\slashed D^2$ の固有空間 $E_{\lambda^2}=E^+_{\lambda^2}\oplus E^-_{\lambda^2}$ に対し $\slashed D/\lvert\lambda\rvert\colon E^+_{\lambda^2}\to E^-_{\lambda^2}$ が同型になる。したがって任意の $t>0$ で
\[
\ind\slashed D=\operatorname{Tr}\bigl(\gamma_5e^{-t\slashed D^2}\bigr)\qquad(\text{マッキーン・シンガーの公式})
\]
が成り立ち、$t\to0$ の極限から指数が局所的な量（曲率の多項式）の積分で書けることがわかる。
\end{myprop}

\subsection{アティヤ・シンガーの指数定理}

\begin{mythm}{アティヤ・シンガーの指数定理（ディラック作用素）}{ch09-as}
$M$ を偶数次元の向き付けられた閉スピン多様体、$E\to M$ を接続 $A$（曲率 $F$）をもつエルミート・ベクトル束とする。$E$ に結合したディラック作用素について
\[
\ind\slashed D_E=\int_M\ch(F)\,\Ahat(TM),
\]
\[
\ch(F)=\tr\exp\frac{F}{2\pi},\quad \Ahat(TM)=1-\frac{p_1}{24}+\frac{7p_1^2-4p_2}{5760}-\cdots,
\]
ただし $p_1=-\dfrac{1}{8\pi^2}\tr R\wedge R$（$R$ は $\mathfrak{so}(d)$ 値の曲率2形式）。積分では $d$ 次の部分だけが寄与する。
\end{mythm}

\begin{myex}{低次元での形}{ch09-ascases}
\begin{itemize}
\item \textbf{2次元}：$\Ahat$ に2次の項はないので $\displaystyle\ind\slashed D=\frac{1}{2\pi}\int_M\tr F$。$U(1)$ で電荷1なら磁束量子数 $n=\int F/2\pi$ に等しい。
\item \textbf{4次元}：表現 $R$ のディラック・フェルミオンで $\tr_R(T^aT^b)=T(R)\delta^{ab}$（基本表現で $T=\tfrac12$）とすると
\[
\ind\slashed D=\frac{1}{8\pi^2}\int\tr_RF\wedge F-\frac{\dim R}{24}\int p_1=2T(R)\,k-\frac{\dim R}{24}\int_Mp_1 ,
\]
$k$ は基本表現で測ったインスタントン数（記法と規約）。平坦な $S^4$ 上の基本表現なら $\ind=k$、随伴表現の $SU(N)$ なら $2Nk$。
\end{itemize}
\end{myex}

指数定理は、他の楕円型作用素にも同じ形で拡張される。特に次の2つは位相不変量を「指数」として表す。

\begin{mythm}{ガウス・ボンネとヒルツェブルフの符号数定理}{ch09-gbsign}
$M$ を $d$ 次元の向き付けられた閉リーマン多様体とする。
\begin{enumerate}
\item $d+d^\dagger\colon\Omega^{\text{even}}\to\Omega^{\text{odd}}$ の指数はオイラー標数 $\chi(M)=\sum(-1)^pb_p$ に等しく、$\chi(M)=\int_Me(TM)$。2次元では $\displaystyle\chi=\frac{1}{2\pi}\int_MK\,dA$（$K$ はガウス曲率）。
\item $d=4$ で、$*$ の固有値 $\pm1$ に対応する作用素（符号数作用素）の指数は符号数 $\sigma(M)=b_2^+-b_2^-$（第6章の交叉形式の符号数）に等しく
\[
\sigma(M)=\frac13\int_Mp_1 .
\]
\end{enumerate}
\end{mythm}

4次元のスピン多様体では、$\ind\slashed D=-\frac1{24}\int p_1=-\frac{\sigma}{8}$ である。

\begin{myex}{いくつかの4次元多様体}{ch09-4mfd}
\begin{center}
\renewcommand{\arraystretch}{1.2}\small
\begin{tabular}{lccccl}
$M$ & $\chi$ & $\sigma$ & $\int p_1$ & $\Ahat[M]$ & スピン？\\ \hline
$S^4$ & 2 & 0 & 0 & 0 & ○\\
$S^2\times S^2$ & 4 & 0 & 0 & 0 & ○\\
$\CP^2$ & 3 & 1 & 3 & $-\tfrac18$ & ×\\
$K3$ & 24 & $-16$ & $-48$ & 2 & ○
\end{tabular}
\end{center}
$\CP^2$ では $\Ahat[M]$ が整数でないので、指数定理より $\CP^2$ はスピン多様体でない（問題 \ref{ch09-prob-cp2}）。$K3$ はリッチ平坦で、2個の正カイラリティの平行スピノルをもつ（問題 \ref{ch09-prob-k3}）。
\end{myex}

\begin{mythm}{ロホリンの定理}{ch09-rokhlin}
閉じた滑らかな4次元スピン多様体の符号数は16の倍数である。
\end{mythm}
\begin{note}
指数定理から $\sigma=-8\ind\slashed D$ なので $\sigma\equiv0\pmod 8$ はすぐわかる。さらに4次元の半スピノル表現 $S^\pm$ は四元数的（擬実）で、$\ker\slashed D_\pm$ には四元数構造 $J$（$J^2=-1$、$\slashed D$ と可換な反線形写像）が作用するので $n_\pm$ は偶数。よって $\ind\slashed D$ は偶数で $\sigma\equiv0\pmod{16}$。
\end{note}

\begin{myex}{ゼロモードの例}{ch09-zmex}
\begin{itemize}
\item \textbf{$S^2$ とモノポール}：磁束 $\int F=2\pi n$（$n>0$）の背景では、正カイラリティのゼロモードが $n$ 個、負カイラリティのゼロモードはない（問題 \ref{ch09-prob-s2}）。ゼロモードは正則関数 $1,z,\dots,z^{n-1}$ で与えられる。
\item \textbf{トーラスのランダウ準位}：$T^2$ 上で一様な磁場の磁束が $2\pi N$ なら、最低ランダウ準位の縮退度 $N$ が指数に等しい（問題 \ref{ch09-prob-landau}）。
\item \textbf{BPSTインスタントン}：$k=1$ の $SU(2)$ インスタントン（大きさ $\rho$、正則ゲージ）中の基本表現のフェルミオンのゼロモードは
\[
\psi_0(x)\propto\frac{\rho}{(x^2+\rho^2)^{3/2}}\,\chi,\qquad \lvert\psi_0\rvert^2=\frac{2}{\pi^2}\frac{\rho^2}{(x^2+\rho^2)^3}
\]
（$\chi$ はスピンと色の添字をからめた定数スピノル）で、インスタントンのまわり半径 $\sim\rho$ に局在する。$\int d^4x\,\rho^2/(x^2+\rho^2)^3=\pi^2/2$ を用いた。
\end{itemize}
\end{myex}

\begin{center}
\begin{tikzpicture}[scale=1.0]
  \draw[-{Stealth[length=2mm]}] (-3.6,0) -- (3.8,0) node[right] {\footnotesize $\lvert x\rvert/\rho$};
  \draw[-{Stealth[length=2mm]}] (0,-0.2) -- (0,2.8);
  \draw[very thick,blue!70!black,domain=-3.4:3.4,samples=120,smooth] plot (\x,{2.4/(1+\x*\x)/(1+\x*\x)/(1+\x*\x)});
  \draw[thick,red!70!black,densely dashed,domain=-3.4:3.4,samples=120,smooth] plot (\x,{2.4/(1+\x*\x)/(1+\x*\x)/(1+\x*\x)/(1+\x*\x)});
  \draw[thick,orange!70!black,domain=-3.4:3.4,samples=120,smooth] plot (\x,{2.4/(1+\x*\x)^1.5});
  \foreach \t in {-3,-2,-1,1,2,3}{ \draw (\t,0.05) -- (\t,-0.05) node[below] {\scriptsize $\t$}; }
  \node[blue!70!black,right] at (1.1,1.5) {\footnotesize $\lvert\psi_0\rvert^2\propto(1+x^2/\rho^2)^{-3}$};
  \node[red!70!black,right] at (1.1,1.0) {\footnotesize $\tr F_{\mu\nu}F_{\mu\nu}\propto(1+x^2/\rho^2)^{-4}$};
  \node[orange!70!black,right] at (1.6,2.2) {\footnotesize $\lvert\psi_0\rvert\propto(1+x^2/\rho^2)^{-3/2}$};
\end{tikzpicture}
\zucap{BPSTインスタントン中のフェルミオンのゼロモード（原点での値で規格化）。ゼロモードは作用密度と同程度に局在するが、遠方では $\lvert\psi_0\rvert\sim\rho/\lvert x\rvert^3$ とべき的に減衰する（質量ゼロの伝播関数と同じふるまい）。}
\end{center}

\subsection{藤川の方法とABJアノマリー}
ユークリッド時空 $M$（4次元、コンパクト化したもの）上で、$SU(N_c)$ の基本表現の質量ゼロのディラック・フェルミオン $N_f$ 個を考える。作用は $S=\int\bar\psi\slashed D\psi$、重みは $e^{-S}$ で、$\psi$ と $\bar\psi$ は独立なグラスマン変数である。$\slashed D\varphi_n=\lambda_n\varphi_n$ の正規直交固有関数で
\[
\psi=\sum_na_n\varphi_n,\qquad \bar\psi=\sum_n\bar b_n\varphi_n^\dagger,\qquad \mathcal D\psi\mathcal D\bar\psi=\prod_n da_n\,d\bar b_n
\]
と展開する（フレーバー添字は省略）。

\begin{mythm}{藤川のヤコビアン}{ch09-fujikawa}
軸性変換 $\psi'=e^{i\alpha(x)\gamma_5}\psi$、$\bar\psi'=\bar\psi e^{i\alpha(x)\gamma_5}$ のもとで
\[
\mathcal D\psi'\mathcal D\bar\psi'=\exp\Bigl(-2i\int_Md^dx\,\alpha(x)\,\mathcal A(x)\Bigr)\mathcal D\psi\mathcal D\bar\psi,
\]
\[
\mathcal A(x)=\lim_{M_{\mathrm{reg}}\to\infty}\sum_{f}\sum_n\varphi_n^\dagger(x)\gamma_5e^{-\lambda_n^2/M_{\mathrm{reg}}^2}\varphi_n(x).
\]
正則化した $\mathcal A(x)$ は指数密度に等しく、$\mathcal A\,d^dx=N_f\,[\ch(F)\Ahat]_d$。特に $\alpha$ が定数なら、ヤコビアンは $\exp(-2i\alpha N_f\ind\slashed D)=\exp(-2i\alpha N_fk)$。
\end{mythm}
\begin{note}
グラスマン変数では $a'_n=\sum_mC_{nm}a_m$ のとき $\prod da'_n=(\det C)^{-1}\prod da_n$（通常の積分と逆）。$C_{nm}=\int\varphi_n^\dagger e^{i\alpha\gamma_5}\varphi_m$ で、$\bar\psi$ についても同じ行列が現れるので、ヤコビアンは $(\det C)^{-2}=\exp(-2\operatorname{Tr}\log C)=\exp(-2i\int\alpha\sum_n\varphi_n^\dagger\gamma_5\varphi_n)+O(\alpha^2)$。和 $\sum_n\varphi_n^\dagger\gamma_5\varphi_n=\operatorname{tr}\gamma_5\,\delta(0)$ は定義されないので、ゲージ不変な正則化 $e^{-\slashed D^2/M_{\mathrm{reg}}^2}$ を入れる。$\alpha$ が定数なら、マッキーン・シンガーの公式（命題 \ref{prop:ch09-pair}）により正則化によらず $\ind\slashed D$ に等しい。
\end{note}

\begin{myprop}{選択則と軸性電荷の非保存}{ch09-selection}
演算子 $\mathcal O$ が定数の軸性変換で $\mathcal O\mapsto e^{iq\alpha}\mathcal O$ と変換するとき、インスタントン数 $k$ のセクターで $\langle\mathcal O\rangle_k\neq0$ となるのは
\[
q=2N_fk
\]
のときに限る。ミンコフスキー時空の言葉では、インスタントン数 $k$ のトンネル過程で軸性電荷 $Q_5=\int j_5^0$ は $\Delta Q_5=2N_fk$ だけ変化する。局所的には
\[
\partial_\mu j_5^\mu=2N_f\,q(x),\qquad q(x)=\frac{1}{32\pi^2}\epsilon^{\mu\nu\rho\sigma}\tr F_{\mu\nu}F_{\rho\sigma},\quad\int q\,d^4x=k
\]
（全体の符号は $\epsilon$ と $\gamma_5$ の規約による）。$U(1)$ で電荷1のときは $\partial_\mu j_5^\mu=\dfrac{N_f}{2\pi^2}\bm E\cdot\bm B$。2次元では同様に $\Delta Q_5=2N_f\cdot\frac{1}{2\pi}\int F$。
\end{myprop}
\begin{note}
$\langle\mathcal O\rangle=\int\mathcal D\psi\mathcal D\bar\psi\,e^{-S}\mathcal O$ で積分変数を $\psi'$ に取りかえると、$S$ は（$\alpha$ 定数で）不変、測度は $e^{-2i\alpha N_fk}$、$\mathcal O$ は $e^{iq\alpha}$ を出すので $\langle\mathcal O\rangle_k=e^{i\alpha(q-2N_fk)}\langle\mathcal O\rangle_k$。
\end{note}

\begin{myprop}{$\theta$ 角の移動}{ch09-theta}
ユークリッドの重みを $e^{-S+i\theta k}$、質量項を $\bar\psi\,(MP_-+M^\dagger P_+)\psi$（$P_\pm=\frac12(1\pm\gamma_5)$、$M$ はフレーバー空間の複素行列）とする。全フレーバーの軸性変換 $\psi\to e^{i\alpha\gamma_5}\psi$ によって、理論は
\[
(\theta,\ M)\ \longrightarrow\ (\theta-2N_f\alpha,\ e^{-2i\alpha}M)
\]
のものと等価になる。したがって物理的なのは
\[
\bar\theta=\theta-\arg\det M
\]
だけである（この符号は質量項の書き方に依存する）。どれか1つのクォークの質量がゼロなら $\bar\theta$ は消去でき、強いCP問題は生じない。
\end{myprop}

\begin{physbox}[：ABJアノマリーの帰結]
\begin{itemize}
\item $\pi^0\to\gamma\gamma$：軸性アイソスピン流の電磁アノマリー（三角形図、下図）が崩壊振幅を決める。
\item $U(1)_A$ 問題：QCD で $U(1)_A$ はインスタントンによって $\Z_{2N_f}$ まで陽に破れるので、$\eta'$ は南部・ゴールドストン粒子にならず重い（$958\,\mathrm{MeV}$）。
\item 電弱理論の $B+L$ の破れ：$SU(2)_L$ インスタントンは $\Delta B=\Delta L=N_{\mathrm{gen}}k$ を引き起こす。$B-L$ は保存。
\item $\theta$ 真空（第8章）：$\bar\theta=\theta-\arg\det M$ が物理的な CP の破れのパラメータ。
\end{itemize}
\end{physbox}

\subsection{ゲージ・アノマリー、アノマリー多項式、降下方程式}
左巻きのワイル・フェルミオン（反粒子を荷電共役で左巻きに直したもの）をすべて並べ、表現 $R=\bigoplus_iR_i$ でゲージ場と結合させる。ゲージ流の非保存は1ループの三角形図で決まり、その係数は
\[
\mathcal A^{abc}=\sum_{\text{左巻きワイル}}\tr_R\bigl(T^a\{T^b,T^c\}\bigr)
\]
に比例する。

\begin{center}
\begin{tikzpicture}[scale=1.1]
  \coordinate (A) at (90:1);
  \coordinate (B) at (210:1);
  \coordinate (C) at (330:1);
  \draw[thick,mid arrow] (A) -- (B);
  \draw[thick,mid arrow] (B) -- (C);
  \draw[thick,mid arrow] (C) -- (A);
  \node[vtx] at (A) {}; \node[vtx] at (B) {}; \node[vtx] at (C) {};
  \begin{scope}[shift={(A)},rotate=90]
    \draw[thick,blue!70!black,domain=0:1.2,samples=80] plot (\x,{0.08*sin(\x*1200)});
  \end{scope}
  \begin{scope}[shift={(B)},rotate=210]
    \draw[thick,blue!70!black,domain=0:1.2,samples=80] plot (\x,{0.08*sin(\x*1200)});
  \end{scope}
  \begin{scope}[shift={(C)},rotate=330]
    \draw[thick,blue!70!black,domain=0:1.2,samples=80] plot (\x,{0.08*sin(\x*1200)});
  \end{scope}
  \node at (90:2.45) {\footnotesize $A^a_\mu$};
  \node at (210:2.45) {\footnotesize $A^b_\nu$};
  \node at (330:2.45) {\footnotesize $A^c_\rho$};
  \node at (0,0) {\footnotesize $\psi_L$};
  \node[right] at (2.8,0.3) {\small $\propto\ \tr_R\bigl(T^a\{T^b,T^c\}\bigr)$};
  \node[right] at (2.8,-0.3) {\small （交差図との和で対称化）};
\end{tikzpicture}
\zucap{三角形図。左巻きワイル・フェルミオンのループに3本のゲージ場がつく。ゲージ流の発散はこの図の係数 $\tr_R T^a\{T^b,T^c\}$ に比例し、ゲージ・アノマリーが相殺するにはその総和がゼロでなければならない。}
\end{center}

\begin{mythm}{摂動論的アノマリー相殺条件（4次元）}{ch09-pertcond}
4次元のゲージ理論が摂動論的に整合的であるための必要十分条件は
\[
\sum_{\text{左巻きワイル}}\tr_R\bigl(T^a\{T^b,T^c\}\bigr)=0\quad(\text{すべての }a,b,c).
\]
重力との混合アノマリー（三角形の2本が重力子）が消える条件は、$U(1)$ 因子の生成子 $Y$ に対して $\sum\tr_RY=0$。
\end{mythm}
$SU(N)$ の基本表現 $\mathbf N$ で $\tr(T^a\{T^b,T^c\})=\frac12d^{abc}$ とおくと、$\overline{\mathbf N}$ では $-\frac12d^{abc}$。$SU(2)$ では $d^{abc}=0$ であり、実または擬実な表現（$R\cong\bar R$）ではアノマリー係数は自動的に消える（問題 \ref{ch09-prob-witten}）。

\begin{mydef}{アノマリー多項式}{ch09-anpoly}
$d$ 次元の（正カイラリティの）ワイル・フェルミオンのアノマリー多項式を
\[
I_{d+2}=\bigl[\Ahat(TM)\ch(F)\bigr]_{d+2}
\]
で定める（負カイラリティでは符号が反転）。具体的には
\[
I_4=\frac{1}{2}\frac{1}{(2\pi)^2}\tr F^2\ (d=2),\qquad
I_6=\frac{1}{3!}\frac{1}{(2\pi)^3}\tr F^3-\frac{p_1}{24}\frac{1}{2\pi}\tr F\ (d=4).
\]
\end{mydef}
$\tr F^3=F^aF^bF^c\tr(T^aT^bT^c)=\frac12F^aF^bF^c\tr(T^a\{T^b,T^c\})$（2形式は互いに可換）なので、$I_6$ の第1項が消えることが定理 \ref{thm:ch09-pertcond} の条件、第2項が重力との混合アノマリーの条件にあたる。

\begin{mythm}{降下方程式}{ch09-descent}
$I_{d+2}$ は閉形式でゲージ不変なので、局所的に $I_{d+2}=dI^{(0)}_{d+1}$ と書ける（$I^{(0)}_{d+1}$ はチャーン・サイモンズ形式）。$I^{(0)}_{d+1}$ の無限小ゲージ変換 $\delta_\lambda$ は閉形式なので、再び局所的に
\[
\delta_\lambda I^{(0)}_{d+1}=dI^{(1)}_d
\]
と書ける。$d$ 次元のワイル・フェルミオンの有効作用 $W[A]=-\log Z[A]$ のゲージ変換は
\[
\delta_\lambda W=\pm2\pi i\int_MI^{(1)}_d
\]
で与えられる（\term{整合的アノマリー}{consistent anomaly}。符号はカイラリティと向きの規約による）。
\end{mythm}
\begin{center}
\begin{tikzcd}[column sep=large]
I_{d+2}\ (d+2\text{ 次元}) \arrow[r,"d^{-1}"] & I^{(0)}_{d+1}\ (d+1\text{ 次元}) \arrow[r,"\delta_\lambda\text{ のち }d^{-1}"] & I^{(1)}_d\ (d\text{ 次元のアノマリー})
\end{tikzcd}
\zucap{降下方程式。$d$ 次元のアノマリーは、2つ上の次元の特性類から「微分の逆」と「ゲージ変換」を交互に施して得られる。}
\end{center}
$U(1)$ で電荷1の場合を例にとると（問題 \ref{ch09-prob-descent}）
\[
d=2:\ I_4=\frac{F^2}{8\pi^2},\ I_3^{(0)}=\frac{A\,dA}{8\pi^2},\ I_2^{(1)}=\frac{\lambda\,dA}{8\pi^2};
\]
\[
d=4:\ I_6=\frac{F^3}{48\pi^3},\ I_5^{(0)}=\frac{A\,(dA)^2}{48\pi^3},\ I_4^{(1)}=\frac{\lambda\,(dA)^2}{48\pi^3}.
\]
非可換の場合、$F=dA-iA^2$ の規約で $I_5^{(0)}=\frac{1}{6(2\pi)^3}\tr\bigl(A(dA)^2-\frac{3i}{2}A^3dA-\frac35A^5\bigr)$ である。

\begin{myex}{標準模型の1世代}{ch09-sm}
1世代のフェルミオンをすべて左巻きのワイル場で書き、$SU(3)\times SU(2)\times U(1)_Y$ の表現を $(\mathbf r_3,\mathbf r_2)_Y$ で表す：
\begin{center}
\renewcommand{\arraystretch}{1.2}\small
\begin{tabular}{lccccc}
場 & $Q$ & $u^c$ & $d^c$ & $L$ & $e^c$\\ \hline
表現 & $(\mathbf 3,\mathbf 2)_{1/6}$ & $(\bar{\mathbf 3},\mathbf 1)_{-2/3}$ & $(\bar{\mathbf 3},\mathbf 1)_{1/3}$ & $(\mathbf 1,\mathbf 2)_{-1/2}$ & $(\mathbf 1,\mathbf 1)_{1}$\\
成分数 & 6 & 3 & 3 & 2 & 1
\end{tabular}
\end{center}
電荷は $Q_{\mathrm{em}}=T_3+Y$。非自明な条件は $SU(3)^3$, $SU(3)^2U(1)$, $SU(2)^2U(1)$, $U(1)^3$, 重力$^2U(1)$ の5つで、すべてゼロになる（問題 \ref{ch09-prob-sm}）。$SU(2)^3$ と $SU(3)^2SU(2)$ 等は群論的に自明に消える。さらに $SU(2)$ 二重項の数は世代あたり $3+1=4$（偶数）で、ウィッテン・アノマリーもない。
\end{myex}

\paragraph{重力アノマリーとグリーン・シュワルツ機構（概略）}
$d=4k+2$ 次元ではカイラルなフェルミオンや自己双対形式場が純粋な重力アノマリーをもちうる。10次元の $N=1$ 超重力＋超対称ヤン・ミルズ理論では、アノマリー多項式 $I_{12}$ が $I_{12}=X_4\wedge X_8$ と因数分解できれば、2形式場 $B$ の変換と相互作用 $\int B\wedge X_8$ でアノマリーを打ち消せる（\term{グリーン・シュワルツ機構}{Green--Schwarz mechanism}）。これが可能なのはゲージ群の次元が496のとき（$SO(32)$, $E_8\times E_8$）であり、詳しくは第10章で扱う。

\paragraph{アノマリー流入（概略）}
$d$ 次元のアノマリーは、$d+1$ 次元のチャーン・サイモンズ項 $2\pi\int I^{(0)}_{d+1}$ のゲージ変換が境界に残す項 $2\pi\int_{\partial}I^{(1)}_d$ で打ち消せる。すなわちアノマリーをもつ理論は「より高い次元のバルクの境界」として整合的に実現できる（\term{アノマリー流入}{anomaly inflow}、カラン・ハーヴェイ）。量子ホール効果のエッジ状態がその典型例である（問題 \ref{ch09-prob-inflow}）。

\subsection{ウィッテンの $SU(2)$ アノマリー}
$SU(2)$ の表現は実または擬実なので、摂動論的なアノマリーは存在しない。しかし、左巻きの二重項1個の分配関数は、ディラック作用素の行列式の「平方根」（パフィアン）$Z[A]=\Pf(\cdots)$ で定義され、その符号はあいまいである。

\begin{mythm}{ウィッテンの $SU(2)$ アノマリー}{ch09-witten}
$\pi_4(SU(2))=\Z_2$ の非自明な元を代表する大きなゲージ変換 $U\colon S^4\to SU(2)$ をとる。$A$ から $A^U$ へゲージ場の空間の中の道 $A_t$ でつなぐと、表現 $R$ のワイル・フェルミオンについて、ゼロを横切る固有値の対の数は $\bmod 2$ で $2T(R)$（$k=1$ インスタントン中のゼロモードの数）に等しい。したがって $2T(R)$ が奇数となるワイル・フェルミオン（例えば二重項）が奇数個あると
\[
Z[A^U]=-Z[A]
\]
となり、ゲージ不変性から $Z\equiv0$：理論は矛盾する。二重項が偶数個なら矛盾はない。
\end{mythm}
標準模型では世代あたり $SU(2)_L$ 二重項が $Q$ の3色と $L$ の計4個あるので、この条件をみたす。

\section{問題}

\subsection{指数の基本性質}

\begin{problem}[難易度 A]\label{ch09-prob-pair}
$\slashed D$ をコンパクト多様体上のエルミートなディラック作用素、$\gamma_5$ を $\{\gamma_5,\slashed D\}=0$ をみたすカイラリティ演算子とする。\\
(1) $\slashed D\varphi=\lambda\varphi$ ならば $\slashed D(\gamma_5\varphi)=-\lambda\gamma_5\varphi$ であることを示せ。$\lambda\neq0$ のとき、$\varphi$ と $\gamma_5\varphi$ は直交し、$\varphi$ はカイラリティの固有状態でないことを示せ。\\
(2) $\lambda\neq0$ のとき、$\slashed D^2$ の固有値 $\lambda^2$ の固有空間 $E_{\lambda^2}$ は $\gamma_5$ と可換で $E_{\lambda^2}=E^+\oplus E^-$ と分かれ、$\slashed D\colon E^+\to E^-$ が同型であることを示せ。\\
(3) $2$ 次元の空間 $\C\oplus\C$（第1成分が正カイラリティ）上の $\slashed D(s)=\begin{pmatrix}0&s\\ s&0\end{pmatrix}$（$s\in\R$）について、$n_\pm(s)$ と指数を求め、指数の変形不変性を確かめよ。物理的に、なぜ「ゼロモードの総数」ではなく「$n_+-n_-$」が背景場の位相だけで決まると期待されるかを述べよ。
\end{problem}

\begin{solution}
(1) $\slashed D\gamma_5\varphi=-\gamma_5\slashed D\varphi=-\lambda\gamma_5\varphi$。$\gamma_5\varphi$ は固有値 $-\lambda$ の固有ベクトルで、$\slashed D$ はエルミートなので異なる固有値の固有ベクトルは直交する：$\lambda\neq-\lambda$（$\lambda\neq0$）より $\langle\varphi,\gamma_5\varphi\rangle=0$。もし $\gamma_5\varphi=\pm\varphi$ なら $\langle\varphi,\gamma_5\varphi\rangle=\pm\lVert\varphi\rVert^2\neq0$ となって矛盾する。

(2) $\gamma_5\slashed D^2=-\slashed D\gamma_5\slashed D=\slashed D^2\gamma_5$ なので $\gamma_5$ は $E_{\lambda^2}$ を保ち、$\gamma_5^2=1$ より $E_{\lambda^2}=E^+\oplus E^-$（$\gamma_5=\pm1$）。$\slashed D$ は $\gamma_5$ と反可換なので $E^+$ を $E^-$ に写す（$\gamma_5\slashed D\varphi=-\slashed D\gamma_5\varphi=-\slashed D\varphi$）。また $\slashed D^2=\lambda^2\ne0$ が $E_{\lambda^2}$ 上で可逆なので、$\slashed D|_{E^+}$ は単射、同様に $\slashed D|_{E^-}\colon E^-\to E^+$ も単射。有限次元（楕円型作用素の固有空間）なので $\dim E^+\le\dim E^-\le\dim E^+$、よって同型。

(3) $s\neq0$ では $\slashed D(s)$ は可逆（$\det=-s^2$）なので $n_+=n_-=0$。$s=0$ では $\slashed D=0$ で $n_+=n_-=1$。いずれも $\ind=n_+-n_-=0$ で、$s$ によらない。ゼロモードの総数は $0\to2$ と跳ぶが、跳ぶときは正と負のカイラリティが対になって現れ、対になって消える。

物理的には、(2) により非ゼロの固有値は必ず $(+,-)$ の対をなして現れる。背景場を連続に変えると、固有値は連続に動き、ゼロに落ちてくるときもゼロから離れるときも正負カイラリティの対としてしか動けない。したがって $n_+$ と $n_-$ は個別には背景場の詳細（偶然の縮退）に依存しうるが、差 $n_+-n_-$ は連続変形で不変であり、背景場の位相的なデータだけで決まると期待される。それを具体的に与えるのが指数定理である。
\end{solution}

\begin{problem}[難易度 B]\label{ch09-prob-mckean}
コンパクト多様体上のエルミートなディラック作用素 $\slashed D$ の固有値は離散的で、各固有空間は有限次元とする。\\
(1) $t>0$ に対し $\operatorname{Tr}\bigl(\gamma_5e^{-t\slashed D^2}\bigr)=n_+-n_-$ を示せ（マッキーン・シンガーの公式）。\\
(2) 熱核の漸近展開により、$t\to0$ で $\operatorname{Tr}(\gamma_5e^{-t\slashed D^2})=\int_M\bigl(a_0(x)t^{-d/2}+a_1(x)t^{-d/2+1}+\cdots\bigr)$（$a_j$ は曲率と場の強さの局所的な多項式）と書けることを認める。$\ind\slashed D=\int_Ma_{d/2}(x)\,d^dx$ であることを説明し、ゲージ場 $A$ を連続に変形しても指数が変わらないことを結論せよ。\\
(3) 平坦な $\R^4$ 上の $SU(2)$ ゲージ場で、作用有限のものを連続変形で $k=0$ から $k=1$ に移すことはできない。これを指数の言葉で説明せよ（$\R^4\cup\{\infty\}=S^4$ 上の束の位相型が変わることに注意）。
\end{problem}

\begin{solution}
(1) $\slashed D^2$ の固有空間分解 $L^2=\bigoplus_{\mu}E_\mu$（$\mu=\lambda^2\ge0$）で
\[
\operatorname{Tr}\bigl(\gamma_5e^{-t\slashed D^2}\bigr)=\sum_\mu e^{-t\mu}\operatorname{tr}\bigl(\gamma_5|_{E_\mu}\bigr)=\sum_\mu e^{-t\mu}\bigl(\dim E_\mu^+-\dim E_\mu^-\bigr).
\]
前問 (2) より $\mu>0$ では $\dim E^+_\mu=\dim E^-_\mu$ なので寄与しない。$\mu=0$ では $E_0=\ker\slashed D^2=\ker\slashed D$（$\lVert\slashed D\varphi\rVert^2=\langle\varphi,\slashed D^2\varphi\rangle$）で、$\dim E_0^\pm=n_\pm$。よって $t$ によらず $n_+-n_-$。

(2) 左辺は $t$ によらない（(1)）ので、$t\to0$ の漸近展開で $t$ の負べきの係数はすべて積分するとゼロ、$t^0$ の係数 $\int a_{d/2}$ が指数に等しい（正べきの係数もゼロ）。$a_{d/2}(x)$ は $A$ とその微分の多項式（ゲージ不変な曲率の多項式）なので、$A$ を連続に動かすと $\int a_{d/2}$ は連続に変わる。一方、それは整数 $n_+-n_-$ に等しいので、連続な整数値関数として一定である。

(3) $\R^4$ 上の作用有限な配位は $S^4$ 上の束と接続に延長され、$k$ はその束の第2チャーン類（で決まる整数）である（第7, 8章）。同じ束の上の接続の空間はアフィン空間で連結なので、その中では指数 $=k$ は一定である。逆に、連続変形（作用有限のまま）で $k=0$ から $k=1$ に移れたとすると、指数 $=k$ が連続変形で変わることになり (2) に反する。物理的には、$k$ が変わる途中では、作用無限大の配位（インスタントンが点に縮む特異な配位など）を通らなければならない。
\end{solution}

\subsection{具体的なゼロモードと指数定理}

\begin{problem}[難易度 C]\label{ch09-prob-s2}
単位球面 $S^2$ を北極からの立体射影座標 $z=x+iy$ で表すと、計量は $g=e^{2\sigma}(dx^2+dy^2)$、$e^{\sigma}=\dfrac{2}{1+\lvert z\rvert^2}$ である。電荷1のフェルミオンが、一様な磁場 $F=\dfrac n2\,\mathrm{vol}_{S^2}$（$n\in\Z$、$\int_{S^2}F=2\pi n$）の中にあるとする。次の事実を認めてよい。
\begin{itemize}
\item[(i)] $S^2$ 上のディラック作用素のゼロモードは $\psi=e^{-\sigma/2}\psi_0$ で、$\psi_0$ は平坦な作用素 $\slashed D_0=i(\sigma_1D_1+\sigma_2D_2)$（$D_j=\partial_j-iA_j$）のゼロモード。ノルムは $\int e^{2\sigma}\lvert\psi\rvert^2d^2x=\int e^{\sigma}\lvert\psi_0\rvert^2d^2x$。カイラリティは $\gamma_5=\sigma_3$。
\item[(ii)] 以下で現れる形のゼロモードが北極でも滑らかであることは、ノルムが有限であることと同値。
\end{itemize}
(1) $A_j=-\epsilon_{jk}\partial_k\chi$（$\epsilon_{12}=1$）とおくと $F_{12}=\nabla^2\chi$ であることを示し、$\chi=\dfrac n2\log(1+\lvert z\rvert^2)$ が上の磁場を与えることを確かめよ（$\mathrm{vol}_{S^2}=e^{2\sigma}dx\wedge dy$）。\\
(2) 正カイラリティ $\psi_0=(u,0)^T$ のゼロモード方程式が $(\partial_{\bar z}+\partial_{\bar z}\chi)u=0$ となることを示し、$u=f(z)e^{-\chi}$（$f$ は整関数）を得よ。$n>0$ のとき、ノルム有限な解は $f$ が $n-1$ 次以下の多項式であるものに限り、$n_+=n$ であることを示せ。\\
(3) 負カイラリティのゼロモードは $v=g(\bar z)e^{\chi}$ の形であることを示し、$n>0$ では $n_-=0$ であることを示せ。$n<0$ の場合も論じ、$\ind\slashed D=n$ が指数定理と一致することを確かめよ。\\
(4) $n=0$ ではゼロモードがないことを、(2), (3) から確かめよ（曲率 $R>0$ のためのリヒネロヴィッツの消滅定理と整合する）。
\end{problem}

\begin{solution}
(1) $A_1=-\partial_2\chi$、$A_2=\partial_1\chi$ なので $F_{12}=\partial_1A_2-\partial_2A_1=\partial_1^2\chi+\partial_2^2\chi=\nabla^2\chi$。$r=\lvert z\rvert$ として
\[
\nabla^2\log(1+r^2)=\frac1r\frac{d}{dr}\Bigl(r\cdot\frac{2r}{1+r^2}\Bigr)=\frac1r\cdot\frac{4r}{(1+r^2)^2}=\frac{4}{(1+r^2)^2}.
\]
よって $F_{12}=\dfrac{2n}{(1+r^2)^2}=\dfrac n2e^{2\sigma}$、すなわち $F=\frac n2\mathrm{vol}_{S^2}$。全磁束は $\int\frac{2n}{(1+r^2)^2}d^2x=2n\cdot\pi\int_0^\infty\frac{dr^2}{(1+r^2)^2}=2\pi n$。

(2) $\sigma_1(u,0)^T=(0,u)^T$、$\sigma_2(u,0)^T=(0,iu)^T$ より $\slashed D_0(u,0)^T=i\bigl(0,(D_1+iD_2)u\bigr)^T$。$\partial_{\bar z}=\frac12(\partial_1+i\partial_2)$ として
\[
D_1+iD_2=2\partial_{\bar z}-i(A_1+iA_2)=2\partial_{\bar z}-i(-\partial_2\chi+i\partial_1\chi)=2\partial_{\bar z}+(\partial_1+i\partial_2)\chi=2(\partial_{\bar z}+\partial_{\bar z}\chi).
\]
$u=e^{-\chi}f$ とおくと $(\partial_{\bar z}+\partial_{\bar z}\chi)u=e^{-\chi}\partial_{\bar z}f$ なので、方程式は $\partial_{\bar z}f=0$、すなわち $f$ は $\C$ 上の正則関数（整関数）。

ノルム：$e^{-\chi}=(1+r^2)^{-n/2}$ なので
\[
\int e^{\sigma}\lvert u\rvert^2d^2x=\int\frac{2\lvert f(z)\rvert^2}{(1+r^2)^{n+1}}d^2x .
\]
$f=\sum_mc_mz^m$ と展開すると、重みは回転対称なので $\int_0^{2\pi}z^m\bar z^{m'}d\phi=2\pi r^{2m}\delta_{mm'}$ により
\[
\lVert u\rVert^2=\sum_m\lvert c_m\rvert^2N_m,\qquad N_m=4\pi\int_0^\infty\frac{r^{2m+1}}{(1+r^2)^{n+1}}dr .
\]
被積分関数は $r\to\infty$ で $r^{2m-2n-1}$ のようにふるまうので、$N_m<\infty\iff2m-2n-1<-1\iff m\le n-1$。よってノルム有限 $\iff$ $m\ge n$ の係数がすべてゼロ、すなわち $f$ は $n-1$ 次以下の多項式。線形独立な解は $z^m(1+\lvert z\rvert^2)^{-n/2}$（$m=0,\dots,n-1$）で $n_+=n$。

(3) $\sigma_1(0,v)^T=(v,0)^T$、$\sigma_2(0,v)^T=(-iv,0)^T$ より $\slashed D_0(0,v)^T=i\bigl((D_1-iD_2)v,0\bigr)^T$。$\partial_z=\frac12(\partial_1-i\partial_2)$ として
\[
D_1-iD_2=2\partial_z-i(A_1-iA_2)=2\partial_z-i(-\partial_2\chi-i\partial_1\chi)=2\partial_z-(\partial_1-i\partial_2)\chi=2(\partial_z-\partial_z\chi),
\]
よって $v=e^{\chi}g$、$\partial_zg=0$（$g$ は $\bar z$ の整関数）。ノルムは $\int\frac{2\lvert g\rvert^2(1+r^2)^{n}}{1+r^2}d^2x$。$g=\sum_m d_m\bar z^m$ とすると各項は $r^{2m+2n-1}$ のようにふるまい、$n>0$ では $m\ge0$ のどの項も $\int^\infty r^{2m+2n-1}dr=\infty$。よって $n_-=0$。

$n<0$ では役割が入れかわる：$u=f(z)(1+r^2)^{\lvert n\rvert/2}$ はどれもノルム無限で $n_+=0$、$v=g(\bar z)(1+r^2)^{-\lvert n\rvert/2}$ は $g$ が $\lvert n\rvert-1$ 次以下の多項式のとき有限で $n_-=\lvert n\rvert$。いずれの場合も
\[
\ind\slashed D=n_+-n_-=n=\frac1{2\pi}\int_{S^2}F ,
\]
指数定理（2次元では $\Ahat$ の寄与なし）と一致する。また $n>0$ ならゼロモードはすべて正カイラリティで、カイラリティは磁束の符号で決まる。

(4) $n=0$ では $\chi=0$。(2) で $N_m=4\pi\int_0^\infty r^{2m+1}(1+r^2)^{-1}dr=\infty$（すべての $m\ge0$）、(3) も同様で、ゼロモードはない。これは $\slashed D^2=\nabla^\dagger\nabla+R/4$（$R=2>0$）から $\slashed D\psi=0\Rightarrow\psi=0$ となるリヒネロヴィッツの議論と整合する。物理的には、(2) のゼロモードは「ハルデーン球面」上の最低ランダウ準位で、磁気単極子 $n$ のまわりの縮退度が $n$ であることを表す。
\end{solution}

\begin{problem}[難易度 B]\label{ch09-prob-landau}
平坦なトーラス $T^2$（面積 $\mathcal A$）上に一様な磁場 $F_{12}=B>0$ があり、全磁束は $B\mathcal A=2\pi N$（$N\in\Z_{>0}$）とする。$\slashed D=i(\sigma_1D_1+\sigma_2D_2)$、$D_j=\partial_j-iA_j$ とする。\\
(1) $\slashed D^2=-D_1^2-D_2^2-\sigma_3F_{12}$ を示せ。\\
(2) 負カイラリティ（$\sigma_3=-1$）のゼロモードがないことを示し、指数定理から正カイラリティのゼロモードがちょうど $N$ 個あることを結論せよ。\\
(3) 電子（スピン $\tfrac12$、$g=2$）の非相対論的なハミルトニアン $H=\dfrac{1}{2m}(\bm\sigma\cdot\bm\Pi)^2$（$\Pi_j=-iD_j$）が $H=\dfrac{1}{2m}\slashed D^2$ であることを示し、エネルギー $0$ の状態（スピン上向きの最低ランダウ準位）の縮退度が $N=B\mathcal A/2\pi$ であることを導け。さらにスピンのない粒子のハミルトニアン $\frac{1}{2m}\bm\Pi^2$ の最低準位がエネルギー $B/2m$ で縮退度 $N$ であることを示せ。\\
(4) 磁場が一様でなくても $B(x)>0$ がいたるところで成り立ち全磁束が $2\pi N$ なら、エネルギー $0$ の状態はやはり $N$ 個であることを示せ（アハロノフ・キャッシャーの定理の一形）。
\end{problem}

\begin{solution}
(1) $\slashed D^2=-(\sigma_jD_j)^2=-\sigma_j\sigma_kD_jD_k$。$\sigma_j\sigma_k=\delta_{jk}+i\epsilon_{jkl}\sigma_l$（2次元なので $l=3$）より
\[
\sigma_j\sigma_kD_jD_k=D_jD_j+i\sigma_3(D_1D_2-D_2D_1)=D_jD_j+i\sigma_3[D_1,D_2].
\]
$[D_1,D_2]=-i(\partial_1A_2-\partial_2A_1)=-iF_{12}$ なので $\sigma_j\sigma_kD_jD_k=D^2+\sigma_3F_{12}$、よって $\slashed D^2=-D^2-\sigma_3F_{12}$。

(2) 負カイラリティでは $\slashed D^2=-D^2+B$。$-D^2=\sum_j(iD_j)^\dagger(iD_j)\ge0$ なので $\langle v,\slashed D^2v\rangle\ge B\lVert v\rVert^2>0$（$v\ne0$）であり、$\slashed D v=0$ の解はない：$n_-=0$。指数定理より $n_+=n_+-n_-=\frac1{2\pi}\int F=N$。

(3) $\Pi_j=-iD_j$ なので $\bm\sigma\cdot\bm\Pi=-i\sigma_jD_j=-\slashed D$、よって $(\bm\sigma\cdot\bm\Pi)^2=\slashed D^2$ で $H=\slashed D^2/2m\ge0$。$H\psi=0\iff\slashed D\psi=0$ なので、エネルギー $0$ の状態はゼロモードであり、(2) より縮退度は $N=B\mathcal A/2\pi$（スピン上向き $\sigma_3=+1$）。これはランダウ準位の縮退度「単位面積あたり $B/2\pi$（$=eB/h$）」にほかならない。

正カイラリティの空間では $\slashed D^2=-D^2-B$、すなわち $\bm\Pi^2=\slashed D^2|_{+}+B$。$\slashed D^2|_+\ge0$ で、その最小固有値 $0$ の固有空間の次元は $N$ なので、$\bm\Pi^2$ の最小固有値は $B$、縮退度 $N$。よってスピンのない粒子の最低準位は $B/2m$（$=\omega_c/2$）で縮退度 $N$。

(4) 負カイラリティでは $\slashed D^2=-D^2+B(x)\ge\min_xB>0$（$T^2$ はコンパクトなので最小値が存在）。よって $n_-=0$、指数定理（局所的な磁場の形によらず $\frac1{2\pi}\int F$ だけで決まる）から $n_+=N$。ゼロエネルギー状態の数は磁場の詳細によらない。これは指数の頑健さの典型例である。
\end{solution}

\begin{problem}[難易度 A]\label{ch09-prob-gb}
向き付けられた閉曲面 $\Sigma$ 上で、作用素 $\mathcal D=d+d^\dagger\colon\Omega^0\oplus\Omega^2\to\Omega^1$ を考える。\\
(1) $\omega\in\Omega^0\oplus\Omega^2$ について $\lVert\mathcal D\omega\rVert^2=\lVert d\omega\rVert^2+\lVert d^\dagger\omega\rVert^2$ を示し、$\ker\mathcal D$ が調和0形式と調和2形式からなることを示せ。同様に $\ker\mathcal D^\dagger$（$\mathcal D^\dagger=d+d^\dagger\colon\Omega^1\to\Omega^0\oplus\Omega^2$）は調和1形式からなる。\\
(2) ホッジの定理（第5章）を用いて $\ind\mathcal D=b_0-b_1+b_2=\chi(\Sigma)$ を示せ。\\
(3) 指数定理の帰結であるガウス・ボンネの定理 $\chi=\frac1{2\pi}\int_\Sigma K\,dA$ を、半径 $a$ の球面と種数 $g$ の曲面について確かめよ（後者は $\chi=2-2g$ を用いてよい）。
\end{problem}

\begin{solution}
(1) $\omega=f+\eta$（$f\in\Omega^0$, $\eta\in\Omega^2$）とすると $\mathcal D\omega=df+d^\dagger\eta\in\Omega^1$。
\[
\lVert df+d^\dagger\eta\rVert^2=\lVert df\rVert^2+\lVert d^\dagger\eta\rVert^2+2\,\mathrm{Re}\langle df,d^\dagger\eta\rangle,\qquad\langle df,d^\dagger\eta\rangle=\langle d^2f,\eta\rangle=0 .
\]
$d\omega=df$（$d\eta=0$）、$d^\dagger\omega=d^\dagger\eta$（$d^\dagger f=0$）なので主張が従う。よって $\mathcal D\omega=0\iff df=0,\ d^\dagger\eta=0$。$f$ については $d^\dagger f=0$、$\eta$ については $d\eta=0$ が自動的なので、$f$, $\eta$ はともに調和形式。1形式についても同様に $\lVert(d+d^\dagger)\alpha\rVert^2=\lVert d\alpha\rVert^2+\lVert d^\dagger\alpha\rVert^2$（$d\alpha\in\Omega^2$, $d^\dagger\alpha\in\Omega^0$ は次数が異なり直交）なので、$\ker\mathcal D^\dagger$ は調和1形式。

(2) ホッジの定理より調和 $p$ 形式の空間の次元は $b_p$。よって $\ind\mathcal D=\dim\ker\mathcal D-\dim\ker\mathcal D^\dagger=(b_0+b_2)-b_1=\chi(\Sigma)$。

(3) 半径 $a$ の球面：$K=1/a^2$、面積 $4\pi a^2$ なので $\frac1{2\pi}\int K\,dA=2=\chi(S^2)$。半径によらないことが位相不変性を表す。種数 $g$：$\chi=2-2g$ なので $\int K\,dA=2\pi(2-2g)$。例えば $g\ge2$ の定曲率 $K=-1$ の計量では面積が $4\pi(g-1)$ と決まる。$g=1$ では平坦な計量が可能で $\int K=0$。
\end{solution}

\begin{problem}[難易度 B]\label{ch09-prob-cp2}
$\CP^2$ の全ポントリャーギン類は $p(\CP^2)=1+3x^2$（$x\in H^2(\CP^2;\Z)$ は生成元で $\int_{\CP^2}x^2=1$）であることを認める（第8章）。\\
(1) ヒルツェブルフの符号数定理から $\sigma(\CP^2)$ を求め、交叉形式（第6章）と比べよ。$S^2\times S^2$ についても同様に論ぜよ（$p_1(S^2\times S^2)=0$ を用いてよい）。\\
(2) $\CP^2$ がスピン構造をもたないことを、指数定理を用いて示せ。\\
(3) $\CP^2$ はスピン$^c$構造をもち、そのディラック作用素の指数は $\int_{\CP^2}e^{c}\Ahat$（$c=(m+\tfrac12)x$、$m\in\Z$）で与えられる。この指数を計算し、つねに整数であることを確かめよ。物理的には、$\CP^2$ 上の電荷をもつフェルミオンは「半整数の磁束」の背景にあるときに限り定義できることを意味する。
\end{problem}

\begin{solution}
(1) $\sigma=\frac13\int p_1=\frac13\cdot3=1$。$\CP^2$ の交叉形式は $H^2=\Z x$ 上の $(1)$（$x\cdot x=1$）で、符号数 $b_2^+-b_2^-=1-0=1$ と一致する。$S^2\times S^2$ では $\sigma=0$。交叉形式は $\begin{pmatrix}0&1\\1&0\end{pmatrix}$ で固有値 $\pm1$、符号数 $0$ と一致する。

(2) $\CP^2$ がスピンなら、ディラック作用素の指数は $\int\Ahat=-\frac1{24}\int p_1=-\frac{3}{24}=-\frac18$ となる。指数は整数なので矛盾。よって $\CP^2$ はスピンでない（$w_2(\CP^2)=x\bmod2\ne0$ と整合する）。

(3) 4次の部分を取り出すと
\[
\bigl[e^c\Ahat\bigr]_4=\frac{c^2}{2}-\frac{p_1}{24}=\Bigl(\frac{(2m+1)^2}{8}-\frac{3}{24}\Bigr)x^2=\frac{4m^2+4m}{8}x^2=\frac{m(m+1)}{2}x^2 .
\]
よって $\ind=\frac{m(m+1)}2\in\Z$（$m(m+1)$ は偶数）。例えば $m=0,1,2$ で $0,1,3$。$m=1$ のとき $2c=3x=c_1(\CP^2)$ で、スピン$^c$ディラック作用素はドルボー作用素 $\bar\partial+\bar\partial^\dagger$ に一致し、指数はトッド種数 $1$ になる。

物理的解釈：$c=(m+\frac12)x$ は「スピノル束の分の $\frac12x$」と「電荷を担う線束の分」の和である。スピン構造がないため $c$ を整数類（$m$ だけ）にすると (2) と同様に指数が整数にならず、フェルミオンが定義できない。半整数ずれた磁束 $\int_{\CP^1}F/2\pi=m+\frac12$ をとって初めて整合的になる（曲がった空間上の量子化条件のずれ）。
\end{solution}

\begin{problem}[難易度 B]\label{ch09-prob-k3}
$K3$ 曲面は単連結な4次元閉スピン多様体で、$\chi=24$、$\sigma=-16$ であり、ホロノミーが $SU(2)$ のリッチ平坦計量をもつ。\\
(1) $\int_{K3}p_1$ と $\Ahat[K3]$ を求めよ。\\
(2) リヒネロヴィッツの公式 $\slashed D^2=\nabla^\dagger\nabla+\frac R4$ を用いて、リッチ平坦（$R=0$）の閉多様体上のディラック作用素のゼロモードは平行スピノル（$\nabla\psi=0$）であることを示せ。$\Spin(4)=SU(2)_+\times SU(2)_-$ で $S^\pm$ が $(\mathbf 2,\mathbf 1)$, $(\mathbf 1,\mathbf 2)$ であり、ホロノミー群 $SU(2)$ がどちらか一方の因子に含まれることを用いて、$(n_+,n_-)=(2,0)$ であることを指数と整合的に説明せよ。\\
(3) ロホリンの定理を、(i) 指数定理 $\sigma=-8\ind\slashed D$、(ii) $n_\pm$ が偶数（四元数構造）から導け。$K3$ がこれをみたすことを確かめ、交叉形式が $E_8$ 形式（偶、ユニモジュラー、符号数8）となる単連結な滑らかな閉4次元多様体が存在しないことを結論せよ（単連結な4次元多様体では、交叉形式が偶であることとスピンであることは同値であることを用いてよい）。\\
(4) 弦理論で10次元の理論を $K3$ 上にコンパクト化すると、6次元のカイラルなフェルミオンが現れる。(1), (2) のどの事実がそれを保証するか、一言で述べよ（詳細は第10章）。
\end{problem}

\begin{solution}
(1) 符号数定理より $\int p_1=3\sigma=-48$。$\Ahat[K3]=-\frac1{24}\int p_1=2$。

(2) $\slashed D\psi=0$ なら $0=\langle\psi,\slashed D^2\psi\rangle=\lVert\nabla\psi\rVert^2+\frac14\int R\lvert\psi\rvert^2=\lVert\nabla\psi\rVert^2$ なので $\nabla\psi=0$。逆に平行スピノルは $\slashed D\psi=i\gamma^\mu\nabla_\mu\psi=0$。平行スピノルは1点での値で決まり、その値はホロノミー群で不変でなければならない（平行移動で一周しても戻る）。ホロノミー $SU(2)$ が例えば $SU(2)_-$ に含まれるとすると、$S^+=(\mathbf 2,\mathbf 1)$ には自明に作用するので $S^+$ の2次元すべてが不変：平行スピノルが2個。$S^-=(\mathbf 1,\mathbf 2)$ には基本表現として作用し不変ベクトルはない。よって $(n_+,n_-)=(2,0)$ で $\ind=2=\Ahat[K3]$ と整合する（逆に、指数が $+2$ であることが、規約のもとでホロノミーが $SU(2)_-$ 側にあることを決めている）。

(3) (i) より $\sigma=-8(n_+-n_-)$、(ii) より $n_+-n_-$ は偶数なので $\sigma\in16\Z$。$K3$：$\sigma=-16$ でみたす。$E_8$ 形式をもつ単連結な滑らかな閉4次元多様体があれば、交叉形式が偶なのでスピンであり、ロホリンの定理から $\sigma\equiv0\pmod{16}$ でなければならないが $\sigma=8$ なので矛盾。よって存在しない（位相多様体としては存在する：フリードマン。滑らかさが本質的）。

(4) $\ind\slashed D=2\ne0$ で、ゼロモードが一方のカイラリティにだけ偏っていること（$(n_+,n_-)=(2,0)$）。10次元のワイル・スピノルを $K3$ 上のゼロモードで展開すると、6次元のスピノルのカイラリティは $K3$ 方向のカイラリティと連動するので、6次元に正味のカイラルなフェルミオンが残る。
\end{solution}

\begin{problem}[難易度 B]\label{ch09-prob-reps}
平坦な $S^4$（$\int p_1=0$）上で $SU(N)$ のインスタントン数 $k$ の背景を考える。\\
(1) $\tr_R(T^aT^b)=T(R)\delta^{ab}$ とするとき、表現 $R$ のディラック・フェルミオンの指数が $2T(R)k$ であることを示せ。\\
(2) 随伴表現では $T(\mathrm{adj})=N$ である。$\mathcal N=1$ 超対称ヤン・ミルズ理論のゲージーノ（随伴表現のマヨラナ＝ワイル・フェルミオン）は、$k=1$ の背景で何個のゼロモードをもつか。$U(1)_R$（ゲージーノの位相回転）がどの離散部分群に破れるかを答えよ（ワイル・フェルミオン1個の測度は、ディラックの場合の半分の位相 $e^{-i\alpha\cdot2Nk}$ を出すとしてよい）。\\
(3) $SU(2)$ のスピン $j$ 表現で $T(j)=\frac13j(j+1)(2j+1)$ を示し、$j=\frac12,1,\frac32,\frac52$ の指数を $k=1$ で求めよ。\\
(4) $K3$ 上で $SU(N)$ の基本表現のディラック・フェルミオンが、インスタントン数 $k$ の背景にあるときの指数を求めよ。
\end{problem}

\begin{solution}
(1) $\ch(F)$ の4次の部分は $\frac{1}{8\pi^2}\tr_RF^2=\frac{1}{8\pi^2}F^aF^b\tr_R(T^aT^b)=\frac{T(R)}{8\pi^2}F^aF^a$。基本表現（$T=\frac12$）で $k=\frac1{8\pi^2}\int\frac12F^aF^a$ なので $\frac1{8\pi^2}\int\tr_RF^2=2T(R)k$。$\int p_1=0$ なので指数は $2T(R)k$。

(2) 随伴表現の指数は $2Nk$。ゲージーノはワイル（かつマヨラナ）なので、ディラックの場合の正カイラリティ側だけを数え、$k=1$ でゼロモードは $2N$ 個。ゲージーノの回転 $\lambda\to e^{i\alpha}\lambda$ のもとで測度は $e^{-2iN\alpha k}$ を出すので（題意）、これが $1$ となる $\alpha\in\frac{2\pi}{2N}\Z$ だけが対称性として残る：$U(1)_R\to\Z_{2N}$。$2N$ 個のゼロモードを吸収する演算子は $\lambda^{2N}$ 型で、例えばゲージーノ凝縮 $\langle\lambda\lambda\rangle$ の議論に使われる。

(3) $T_3$ の固有値 $m=-j,\dots,j$ で $\tr(T_3^2)=\sum_mm^2=\frac13j(j+1)(2j+1)$。$\tr(T^aT^b)=T\delta^{ab}$ の $a=b=3$ 成分から $T(j)=\frac13j(j+1)(2j+1)$。$k=1$ の指数 $2T(j)$ は
\[
j=\tfrac12:\ 1,\qquad j=1:\ 4,\qquad j=\tfrac32:\ 10,\qquad j=\tfrac52:\ 35 .
\]
（$j=1$ は (2) の $2N=4$ と一致。）

(4) $\ch(F)=N+\ch_2$、$\Ahat=1-\frac{p_1}{24}$ より 4次の部分は $\ch_2-\frac{N}{24}p_1$。$\int\ch_2=k$、$\int_{K3}p_1=-48$ なので $\ind=k+2N$。
\end{solution}

\subsection{藤川の方法とABJアノマリー}

\begin{problem}[難易度 A]\label{ch09-prob-schwinger}
2次元の $U(1)$ ゲージ理論で、電荷1の質量ゼロのディラック・フェルミオン1個を考える（シュウィンガー模型）。\\
(1) コンパクト化したユークリッド時空 $M$ 上で磁束 $\int_MF=2\pi n$ の背景があるとき、命題 \ref{prop:ch09-selection} を用いて、非ゼロとなりうる演算子の軸性電荷を求めよ。\\
(2) ミンコフスキー時空のアノマリー方程式 $\partial_\mu j_5^\mu=\frac{1}{2\pi}\epsilon^{\mu\nu}F_{\mu\nu}$（$\epsilon^{01}=1$）から、長さ $L$ の空間円周上で一様な電場 $E=F_{01}$ を時間 $T$ だけかけたときの $\Delta Q_5$ を求めよ。$A_1$ が $2\pi/L$ だけ変化するごとに $\Delta Q_5=2$ となることを確かめ、本章冒頭のスペクトル流の図と比べよ。
\end{problem}

\begin{solution}
(1) 2次元の指数は $n$（例 \ref{ex:ch09-ascases}）。$N_f=1$ の選択則より $q=2n$。例えば $n=1$ では軸性電荷 $2$ の演算子 $\bar\psi P_+\psi$ が期待値をもちうる（実際、シュウィンガー模型ではカイラル凝縮 $\langle\bar\psi\psi\rangle\ne0$ がこのセクターから生じる）。

(2) $\epsilon^{\mu\nu}F_{\mu\nu}=F_{01}-F_{10}=2E$ なので $\partial_\mu j_5^\mu=E/\pi$。空間について積分すると（円周上で全微分は落ちる）$\dot Q_5=EL/\pi$、よって
\[
\Delta Q_5=\frac{ELT}{\pi}=2\cdot\frac{1}{2\pi}\int F_{01}\,dt\,dx .
\]
$E=\partial_0A_1$（$A_0=0$ のゲージ）なら $\int E\,dt=\Delta A_1$、$\Delta Q_5=L\Delta A_1/\pi$。$\Delta A_1=2\pi/L$ で $\Delta Q_5=2$。図では右巻き粒子が1個減り（または増え）、左巻きが1個増え（または減り）、$N_R-N_L$ が2変化する。ベクトル電荷 $N_R+N_L$ は保存し、軸性電荷だけが変化する。
\end{solution}

\begin{problem}[難易度 C]\label{ch09-prob-fujikawa2d}
平坦な2次元ユークリッド空間で $\slashed D=i(\gamma^1D_1+\gamma^2D_2)$、$\gamma^1=\sigma_1$, $\gamma^2=\sigma_2$, $\gamma_5=\sigma_3$、$D_\mu=\partial_\mu-iA_\mu$ とする（$U(1)$, 電荷1）。\\
(1) 局所的な軸性変換 $\psi\to e^{i\alpha(x)\gamma_5}\psi$, $\bar\psi\to\bar\psi e^{i\alpha(x)\gamma_5}$ のもとで、ヤコビアンが $\exp\bigl(-2i\int\alpha(x)\mathcal A(x)d^2x\bigr)$、$\mathcal A(x)=\sum_n\varphi_n^\dagger(x)\gamma_5\varphi_n(x)$ となることを示せ（グラスマン積分の変数変換の規則は既知としてよい）。\\
(2) 正則化 $\mathcal A(x)=\lim_{M\to\infty}\operatorname{tr}\bigl[\gamma_5\,e^{-\slashed D^2/M^2}\bigr](x,x)$ を、平面波基底
\[
\operatorname{tr}\bigl[\gamma_5e^{-\slashed D^2/M^2}\bigr](x,x)=\int\frac{d^2p}{(2\pi)^2}\operatorname{tr}\Bigl[\gamma_5\,e^{-ip\cdot x}e^{-\slashed D^2/M^2}e^{ip\cdot x}\Bigr]
\]
で計算し、$\mathcal A(x)=\dfrac{F_{12}}{2\pi}$ を示せ。$\slashed D^2=-D^2-\sigma_3F_{12}$（問題 \ref{ch09-prob-landau}）を用いてよい。\\
(3) 結果が指数定理（$\int\mathcal A=n$）と一致すること、$\gamma_5$ を $-\sigma_3$ にとると符号が反転すること（規約の確認）を述べよ。また、正則化に $\slashed D^2$ のかわりにゲージ不変でない $-\partial^2$ を用いると何が問題か述べよ。
\end{problem}

\begin{solution}
(1) $\psi'=\sum_na'_n\varphi_n$ とすると $a'_n=\sum_mC_{nm}a_m$、$C_{nm}=\int\varphi_n^\dagger e^{i\alpha\gamma_5}\varphi_m$。グラスマン変数では $\prod_nda'_n=(\det C)^{-1}\prod_nda_n$。$\bar\psi'=\bar\psi e^{i\alpha\gamma_5}$ の係数 $\bar b'_m=\sum_n\bar b_nC_{nm}$ も同じ行列なので $\prod d\bar b'=(\det C)^{-1}\prod d\bar b$。無限小 $\alpha$ で $C=1+iK$、$K_{nm}=\int\alpha\varphi_n^\dagger\gamma_5\varphi_m$ なので
\[
(\det C)^{-2}=\exp(-2\operatorname{Tr}\log C)=\exp\Bigl(-2i\sum_nK_{nn}\Bigr)=\exp\Bigl(-2i\int\alpha(x)\sum_n\varphi_n^\dagger\gamma_5\varphi_n\,d^2x\Bigr).
\]
有限の $\alpha$ はこれを積み重ねて得られる（指数の肩は $\alpha$ について線形のまま）。

(2) $e^{-ip\cdot x}D_\mu e^{ip\cdot x}=D_\mu+ip_\mu$ なので
\[
e^{-ip\cdot x}\slashed D^2e^{ip\cdot x}=-(D+ip)^2-\sigma_3F_{12}=p^2-2ip\cdot D-D^2-\sigma_3F_{12}.
\]
$p=Mq$ と変数変換すると $d^2p=M^2d^2q$ で
\[
\operatorname{tr}[\gamma_5e^{-\slashed D^2/M^2}](x,x)=M^2\int\frac{d^2q}{(2\pi)^2}e^{-q^2}\operatorname{tr}\Bigl[\gamma_5\exp\Bigl(\frac{2iq\cdot D}{M}+\frac{D^2+\sigma_3F_{12}}{M^2}\Bigr)\Bigr]\cdot1 .
\]
（最後の $1$ は作用素を定数関数1に作用させることを表す。）指数関数を $1/M$ で展開する。$\operatorname{tr}\gamma_5=\operatorname{tr}\sigma_3=0$ なので、トレースが残るのは $\sigma_3$ を1個含む項（$\operatorname{tr}(\sigma_3\sigma_3)=2$）だけである。$\sigma_3$ は $\sigma_3F_{12}/M^2$ からしか来ないので、$M^2$ をかけて $M\to\infty$ で生き残るのは $\sigma_3F_{12}/M^2$ の1次の項だけ（他の因子をかけると $1/M$ 以上の負べきが余分につく）。よって
\[
\mathcal A(x)=M^2\cdot\frac{F_{12}}{M^2}\operatorname{tr}(\sigma_3\sigma_3)\int\frac{d^2q}{(2\pi)^2}e^{-q^2}=2F_{12}\cdot\frac{\pi}{4\pi^2}=\frac{F_{12}}{2\pi}.
\]

(3) $\int\mathcal A\,d^2x=\frac1{2\pi}\int F_{12}d^2x=n$ で、マッキーン・シンガーの公式と指数定理に一致する。$\gamma_5=-\sigma_3$ なら $\operatorname{tr}(\gamma_5\sigma_3)=-2$ で $\mathcal A=-F_{12}/2\pi$、指数が $-n$ になる。記法と規約（$S^2$ で $\int F/2\pi=n$ なら $\ind=n$）は $\gamma_5=\sigma_3=-i\gamma^1\gamma^2$ を選ぶことにあたる。

ゲージ不変でない正則化 $e^{\partial^2/M^2}$ を使うと、正則化された量がゲージ場 $A_\mu$ そのもの（$F$ でなく）に依存しうるので、ゲージ不変でない項（例えばベクトルの変換に対して $M^2$ や $\partial_\mu A^\mu$ に比例する項）が現れ、ベクトル電流の保存まで破れてしまう。$\slashed D^2$ を用いると、ベクトル $U(1)$ のゲージ不変性を保ったまま軸性の部分にアノマリーを押し付けることができ、結果は一意的にゲージ不変な指数密度になる。
\end{solution}

\begin{problem}[難易度 B]\label{ch09-prob-thooft}
$SU(N_c)$ ゲージ理論に、基本表現の質量ゼロのディラック・フェルミオンが $N_f$ 個あるとする。$k=1$ の自己双対なインスタントン背景では、各フレーバーについて $\slashed D$ のゼロモードは正カイラリティの1個 $\varphi_0$ だけであることを認める。\\
(1) 指数定理と整合することを確かめよ。\\
(2) $\psi_f=a_{f,0}\varphi_0+\cdots$、$\bar\psi_f=\bar b_{f,0}\varphi_0^\dagger+\cdots$ と展開するとき、$\int\prod_fda_{f,0}d\bar b_{f,0}$ のために、$\langle\mathcal O\rangle_{k=1}\ne0$ となる $\mathcal O$ は各フレーバーの $\psi_f$ と $\bar\psi_f$ を少なくとも1個ずつ含まなければならないことを説明せよ。最小の演算子として $\det_{fg}\bigl(\bar\psi_fP_+\psi_g\bigr)$ が現れることを述べ、その軸性電荷が命題 \ref{prop:ch09-selection} と一致することを確かめよ。\\
(3) $N_f=3$（$u,d,s$）のとき、トフーフト頂点は何フェルミオン相互作用か。それが $SU(3)_L\times SU(3)_R\times U(1)_V$ を保ち、$U(1)_A$ を $\Z_6$ に破ることを示せ。
\end{problem}

\begin{solution}
(1) $\ind\slashed D=n_+-n_-=1-0=1=k$（基本表現、$\int p_1=0$）で整合する。

(2) $\varphi_0$ は $\slashed D\varphi_0=0$ なので作用 $S=\int\bar\psi\slashed D\psi=\sum_n\lambda_n\bar b_na_n$ に $a_{f,0}$, $\bar b_{f,0}$ は現れない。グラスマン積分では $\int da=0$、$\int da\,a=1$ なので、被積分関数に $a_{f,0}$ と $\bar b_{f,0}$ がちょうど1個ずつ含まれていない限り積分はゼロになる。$a_{f,0}$ は $\psi_f$ から、$\bar b_{f,0}$ は $\bar\psi_f$ からしか来ないので、$\mathcal O$ は全フレーバーの $\psi_f$ と $\bar\psi_f$ を含まなければならない。$\varphi_0$ は正カイラリティなので、$\psi_f$ の中で効くのは $P_+\psi_f$、$\bar\psi_f$ では $\varphi_0^\dagger=\varphi_0^\dagger P_+$ より $\bar\psi_fP_+$ の部分。フレーバーについて反対称化（グラスマン性）すると、最小の演算子は $\det_{fg}(\bar\psi_fP_+\psi_g)$（$2N_f$ フェルミオン）。

軸性変換で $\bar\psi P_+\psi\to\bar\psi e^{i\alpha\gamma_5}P_+e^{i\alpha\gamma_5}\psi=e^{2i\alpha}\bar\psi P_+\psi$ なので、$\det$ は $e^{2iN_f\alpha}$：$q=2N_f=2N_fk$ で選択則と一致する。

(3) $2N_f=6$ フェルミオン相互作用（図）。カイラル対称性 $SU(3)_+\times SU(3)_-$（$=SU(3)_L\times SU(3)_R$ のどちらかの対応）は $P_\pm\psi\to U_\pm P_\pm\psi$ と作用し、$\bar\psi P_+$ は $P_-\psi$ の共役にあたるので $\bar\psi P_+\to\bar\psi P_+U_-^\dagger$ と変換する。すると$\det(\bar\psi P_+\psi)\to\det(U_-^\dagger)\det(U_+)\det(\bar\psi P_+\psi)=\det(\bar\psi P_+\psi)$（$SU(3)$ では $\det=1$）。$U(1)_V$（$\psi\to e^{i\beta}\psi$, $\bar\psi\to\bar\psi e^{-i\beta}$）では各因子が不変。$U(1)_A$ では $e^{6i\alpha}$ がかかるので、不変なのは $\alpha\in\frac{2\pi}{6}\Z$：$U(1)_A\to\Z_6$（$\alpha=\pi$ はフェルミオン数 $(-1)^F$ に等しく、$\Z_6$ に含まれる）。
\end{solution}

\begin{problem}[難易度 C]\label{ch09-prob-theta}
命題 \ref{prop:ch09-theta} の設定（重み $e^{-S+i\theta k}$、質量項 $\bar\psi(MP_-+M^\dagger P_+)\psi$）で考える。\\
(1) フレーバーに依存する軸性変換 $\psi\to e^{i\beta\gamma_5}\psi$, $\bar\psi\to\bar\psi e^{i\beta\gamma_5}$（$\beta$ はフレーバー空間のエルミート行列で定数）により、理論が $(\theta,M)\to(\theta-2\tr\beta,\ e^{-i\beta}Me^{-i\beta})$ と等価になることを示し、$\bar\theta=\theta-\arg\det M$ が不変であることを示せ。\\
(2) $k=1$ セクターで、質量について1次の寄与が $e^{i\theta}\det M^\dagger\propto e^{i\bar\theta}$ の位相をもつことを前問の結果から説明せよ。\\
(3) $u$ クォークの質量がゼロなら $\bar\theta$ は物理的でないことを示せ。\\
(4) 仮に $U(1)_A$ にアノマリーがなく自発的にのみ破れていたら、$N_f=3$ のQCDには9個目の擬南部・ゴールドストン粒子が現れるはずである（ワインバーグはその質量が $\lesssim\sqrt3m_\pi$ であることを示した）。実際の $\eta'(958)$ がなぜ重くてよいのか、本章の言葉で説明せよ。
\end{problem}

\begin{solution}
(1) $\beta$ を対角化するフレーバー基底で考えれば十分で、各フレーバー $f$ が独立に角度 $\beta_f$ で回る。各フレーバーの測度は $e^{-2i\beta_fk}$ を出すので、全体で $e^{-2i\tr\beta\,k}$。$Z(\theta,M)=\sum_ke^{i\theta k}\int\mathcal D\psi\mathcal D\bar\psi e^{-S_M[\psi]}$ で $\psi=e^{i\beta\gamma_5}\chi$ と変数変換すると $\mathcal D\psi\mathcal D\bar\psi=e^{-2i\tr\beta\,k}\mathcal D\chi\mathcal D\bar\chi$。運動項は不変（$\beta$ 定数）、質量項は
\[
\bar\chi e^{i\beta\gamma_5}(MP_-+M^\dagger P_+)e^{i\beta\gamma_5}\chi=\bar\chi\bigl(e^{-i\beta}Me^{-i\beta}P_-+e^{i\beta}M^\dagger e^{i\beta}P_+\bigr)\chi
\]
（$e^{i\beta\gamma_5}P_\mp=e^{\mp i\beta}P_\mp$ と $P_\pm$ が $\beta$ と可換なことを用いた）。よって $Z(\theta,M)=Z(\theta-2\tr\beta,\ e^{-i\beta}Me^{-i\beta})$。$\det(e^{-i\beta}Me^{-i\beta})=e^{-2i\tr\beta}\det M$ なので $\arg\det M\to\arg\det M-2\tr\beta$、$\theta-\arg\det M$ は不変。

(2) $k=1$ では各フレーバーの正カイラリティのゼロモードを吸収する必要がある（前問）。$e^{-S}$ を質量で展開すると、$-\bar\psi M^\dagger P_+\psi$ の項がちょうど $\bar\psi P_+\psi$ 型なので、最低次は $\det M^\dagger$ に比例し、$\theta$ 項と合わせて $e^{i\theta}\det M^\dagger=\lvert\det M\rvert e^{i(\theta-\arg\det M)}=\lvert\det M\rvert e^{i\bar\theta}$。物理量は $\bar\theta$ だけに依存する。

(3) $m_u=0$ なら $M$ の固有値の1つがゼロで $\det M=0$、$\arg\det M$ は定義されない。より具体的には、$u$ クォークだけを回す $\beta=\operatorname{diag}(\beta_u,0,0)$ は質量項を変えずに $\theta\to\theta-2\beta_u$ と任意にずらせるので、$\theta$ は物理的でない。

(4) $U(1)_A$ はアノマリーにより、インスタントン（より一般にトポロジカル電荷の揺らぎ）のある真空では陽に破れている：測度が $e^{-2iN_f\alpha k}$ を出し、$U(1)_A$ は $\Z_{2N_f}$ にまで破れる（前問）。陽に破れた連続対称性に南部・ゴールドストンの定理は適用されないので、$\eta'$ はゴールドストン粒子である必要はない。トフーフト頂点は $\eta'$ の方向（$\det$ の位相）にだけ質量を与え、$SU(3)_L\times SU(3)_R$ の8個のゴールドストン粒子には影響しない。ワインバーグの上限はアノマリーを無視した議論から来ており、アノマリーを考慮すると解消される（ウィッテン・ヴェネツィアーノ公式は $m_{\eta'}^2$ をトポロジカル感受率で与える）。
\end{solution}

\begin{problem}[難易度 B]\label{ch09-prob-bl}
標準模型の $SU(2)_L$ インスタントン（$k=\pm1$）を考える。$SU(2)_L$ に結合するのは左巻きの二重項 $Q$（各世代3色）と $L$ だけであり、各二重項（ワイル・フェルミオン）は $k=\pm1$ の背景でちょうど1個のゼロモードをもつことを認める。\\
(1) 3世代の場合のトフーフト頂点のフェルミオン線の数と、そのフレーバー構造を答えよ。\\
(2) バリオン数 $B$（クォーク $\frac13$）とレプトン数 $L$ の変化 $\Delta B$, $\Delta L$ を求め、$B-L$ が保存することを示せ。\\
(3) 同じ結論を、$B$ 流と $L$ 流の $SU(2)^2$ アノマリー係数 $\sum_{\text{二重項}}(\text{電荷})\times T(\mathbf 2)$ から導け。さらに $SU(3)$ のインスタントンは $B$ を破らないことを示せ（$Q$ は $B=\frac13$、$u^c$, $d^c$ は $B=-\frac13$）。\\
(4) 零温度でのこの過程の振幅は $e^{-8\pi^2/g^2}$（$g$ は $SU(2)$ 結合定数、$\alpha_W=g^2/4\pi\approx1/30$）程度に抑えられる。おおよその桁を見積もれ。
\end{problem}

\begin{solution}
(1) 二重項は世代あたり $Q$ 3個（色）と $L$ 1個の計4個、3世代で12個。各二重項から1本の線が出るので12本：世代ごとに $qqq\ell$、全体で $\prod_{i=1}^3(q_iq_iq_i\ell_i)$ の形（色は $\epsilon_{abc}$ で反対称化）。

(2) 各世代で3個のクォーク（$B=\frac13$ ずつ）と1個のレプトンが生成（または消滅）するので、$\Delta B=\pm3\cdot\frac13\cdot3=\pm3$、$\Delta L=\pm3$（3世代）。一般に $\Delta B=\Delta L=N_{\mathrm{gen}}k$。$\Delta(B-L)=0$、$\Delta(B+L)=2N_{\mathrm{gen}}k$。

(3) 左巻きワイル・フェルミオン1個は $k$ の背景で $\pm2T(R)k$ だけ数が変わる（ゼロモード数＝$2T(R)k$）。二重項は $2T=1$。$B$ 流：$Q$ が世代あたり3個、各 $B=\frac13$ なので係数 $3\cdot\frac13=1$、3世代で $\Delta B=3k$。$L$ 流：$L$ が世代あたり1個で $\Delta L=3k$。$U(1)_Y$ の $F_Y\wedge F_Y$ の寄与は、$\R^4$ 上の作用有限な配位では $\int F_Y\wedge F_Y=0$ なので効かない。

$SU(3)$：$B$ 流の $SU(3)^2$ 係数は $\sum B\cdot T(R)\cdot(\text{多重度})=\frac12\bigl(2\cdot\frac13-\frac13-\frac13\bigr)=0$（$Q$ は $SU(2)$ 二重項なので3重項が2個）。よって QCD のインスタントンは $B$ を変えない（$u,d$ についてカイラリティは変えるが、クォーク数は保つ）。

(4) $8\pi^2/g^2=2\pi/\alpha_W\approx2\pi\cdot30\approx190$。振幅は $e^{-190}\sim10^{-82}$、確率（振幅の2乗）は $10^{-164}$ 程度で、実験的には全く観測されない。ただし高温（電弱相転移付近）ではスファレロンによる熱的な遷移で $B+L$ の破れが速くなり、バリオン数生成の議論で重要になる。
\end{solution}

\subsection{ゲージ・アノマリーの相殺}

\begin{problem}[難易度 B]\label{ch09-prob-sm}
例 \ref{ex:ch09-sm} の標準模型1世代（左巻きワイル）について、\\
(1) $SU(3)^3$, $SU(3)^2U(1)_Y$, $SU(2)^2U(1)_Y$, $U(1)_Y^3$, 重力$^2U(1)_Y$ の5つのアノマリー係数を表にしてすべてゼロであることを確かめよ。$SU(3)$ の3次の係数は $A(\mathbf3)=1$, $A(\bar{\mathbf 3})=-1$、$SU(N)$ の2次の係数は $T=\frac12$ とせよ。\\
(2) $SU(2)^3$, $SU(2)^2SU(3)$, $SU(3)^2SU(2)$, $SU(2)U(1)^2$ のような係数が自動的に消える理由を述べよ。\\
(3) 右巻きニュートリノ $\nu^c\sim(\mathbf 1,\mathbf 1)_0$ を加えると、$U(1)_{B-L}$（$Q:\frac13$, $u^c,d^c:-\frac13$, $L:-1$, $e^c,\nu^c:+1$）もゲージ化できることを、$U(1)_{B-L}$ を含む6つの条件 $SU(3)^2U(1)_{B-L}$, $SU(2)^2U(1)_{B-L}$, $U(1)_{B-L}^3$, 重力$^2U(1)_{B-L}$, $U(1)_Y^2U(1)_{B-L}$, $U(1)_YU(1)_{B-L}^2$ を調べて確かめよ。$\nu^c$ がないとどの条件が破れるか。
\end{problem}

\begin{solution}
(1) 多重度（色・アイソスピンの成分数）を考慮して計算する。
\begin{center}
\renewcommand{\arraystretch}{1.3}\small
\resizebox{\linewidth}{!}{\begin{tabular}{l|ccccc|c}
 & $Q$ & $u^c$ & $d^c$ & $L$ & $e^c$ & 和\\ \hline
$SU(3)^3$：$A\times(SU(2)\text{ 多重度})$ & $2$ & $-1$ & $-1$ & -- & -- & $0$\\
$SU(3)^2U(1)$：$T\times Y\times(\cdots)$ & $\frac12\cdot2\cdot\frac16$ & $\frac12\cdot(-\frac23)$ & $\frac12\cdot\frac13$ & -- & -- & $0$\\
$SU(2)^2U(1)$：$T\times Y\times(\text{色})$ & $\frac12\cdot3\cdot\frac16$ & -- & -- & $\frac12\cdot(-\frac12)$ & -- & $0$\\
$U(1)^3$：$(\text{成分数})\times Y^3$ & $6\cdot\frac1{216}$ & $3\cdot(-\frac8{27})$ & $3\cdot\frac1{27}$ & $2\cdot(-\frac18)$ & $1$ & $0$\\
重力$^2U(1)$：$(\text{成分数})\times Y$ & $6\cdot\frac16$ & $3\cdot(-\frac23)$ & $3\cdot\frac13$ & $2\cdot(-\frac12)$ & $1$ & $0$
\end{tabular}}
\end{center}
各行の和：$SU(3)^2U(1)$ は $\frac16-\frac13+\frac16=0$。$SU(2)^2U(1)$ は $\frac14-\frac14=0$。$U(1)^3$ は
\[
\frac1{36}-\frac89+\frac19-\frac14+1=\frac{1-32+4-9+36}{36}=0 .
\]
重力は $1-2+1-1+1=0$。

(2) $SU(2)^3$：$SU(2)$ の表現は実または擬実で $d^{abc}=0$（問題 \ref{ch09-prob-witten}）。$G_1^2G_2$（$G_2$ が非可換単純群）：係数は $\tr(T_1^aT_1^b)\tr(T_2^c)$ に比例し、単純リー群の生成子はトレースゼロなので消える。$SU(2)U(1)^2$ も $\tr T_2^a=0$ で消える。$SU(3)^2SU(2)$ なども同様。

(3) $X=B-L$ とする。
\begin{itemize}
\item $SU(3)^2X$：$\frac12\bigl(2\cdot\frac13-\frac13-\frac13\bigr)=0$。
\item $SU(2)^2X$：$\frac12\bigl(3\cdot\frac13-1\bigr)=0$。
\item $X^3$：$6\cdot\frac1{27}+3\cdot(-\frac1{27})+3\cdot(-\frac1{27})+2\cdot(-1)+1+1=0$。
\item 重力$^2X$：$6\cdot\frac13-1-1-2+1+1=0$。
\item $Y^2X$：$6\cdot\frac1{36}\cdot\frac13+3\cdot\frac49\cdot(-\frac13)+3\cdot\frac19\cdot(-\frac13)+2\cdot\frac14\cdot(-1)+1\cdot1\cdot1+0=\frac1{18}-\frac49-\frac19-\frac12+1=0$。
\item $YX^2$：$6\cdot\frac16\cdot\frac19+3\cdot(-\frac23)\cdot\frac19+3\cdot\frac13\cdot\frac19+2\cdot(-\frac12)\cdot1+1\cdot1+0=\frac19-\frac29+\frac19-1+1=0$。
\end{itemize}
$\nu^c$ は $Y=0$ なので $Y^2X$, $YX^2$ には寄与しないが、$X^3$ と重力$^2X$ にはそれぞれ $+1$ 寄与する。$\nu^c$ がないとこの2つが $-1$ となり破れる。したがって $B-L$ をゲージ化するには（世代ごとに）右巻きニュートリノが必要である。
\end{solution}

\begin{problem}[難易度 D]\label{ch09-prob-charge}
標準模型の1世代の場の $SU(3)\times SU(2)$ 表現（$Q,u^c,d^c,L,e^c$）を固定し、ハイパーチャージ $Y_Q,Y_u,Y_d,Y_L,Y_e$ を未知数とする。\\
(1) $SU(3)^2U(1)$, $SU(2)^2U(1)$, 重力$^2U(1)$, $U(1)^3$ の4条件を書き、$Y_Q=a$ とおいて解け。$a\ne0$ のとき $\{Y_u,Y_d\}=\{-4a,2a\}$、$Y_L=-3a$, $Y_e=6a$ となることを示せ。\\
(2) $a=0$ の解も求めよ。\\
(3) ヒッグス二重項 $H\sim(\mathbf 1,\mathbf 2)_{Y_H}$ との湯川結合 $Qu^cH$, $Qd^cH^\dagger$, $Le^cH^\dagger$ がゲージ不変であることを要請すると、(1) の解のうちどれが選ばれ、(2) の解が排除されることを示せ。$Y_H=\frac12$ の規格化で標準模型の値を再現せよ。\\
(4) 電荷 $Q_{\mathrm{em}}=T_3+Y$ を各粒子について求め、水素原子（陽子 $uud$ と電子）が電気的に中性であることがアノマリー相殺から導かれることを説明せよ。
\end{problem}

\begin{solution}
(1) 条件は（$Q$ は2成分×3色、$L$ は2成分、$u^c,d^c$ は3色）
\begin{align*}
SU(3)^2U(1):&\quad 2Y_Q+Y_u+Y_d=0,\\
SU(2)^2U(1):&\quad 3Y_Q+Y_L=0,\\
\text{重力}^2U(1):&\quad 6Y_Q+3Y_u+3Y_d+2Y_L+Y_e=0,\\
U(1)^3:&\quad 6Y_Q^3+3Y_u^3+3Y_d^3+2Y_L^3+Y_e^3=0 .
\end{align*}
$Y_Q=a$ とおくと、第2式より $Y_L=-3a$、第1式より $Y_u+Y_d=-2a$。第3式は $6a-6a-6a+Y_e=0$ で $Y_e=6a$。第4式に代入して
\[
6a^3+3(Y_u^3+Y_d^3)-54a^3+216a^3=0\quad\Longrightarrow\quad Y_u^3+Y_d^3=-56a^3 .
\]
$s=Y_u+Y_d=-2a$、$p=Y_uY_d$ とおくと $Y_u^3+Y_d^3=s^3-3sp=-8a^3+6ap$。よって $6ap=-48a^3$、$a\ne0$ なら $p=-8a^2$。$Y_u,Y_d$ は $t^2+2at-8a^2=(t+4a)(t-2a)=0$ の根で、$\{Y_u,Y_d\}=\{-4a,2a\}$。

(2) $a=0$ なら $Y_L=Y_e=0$、$Y_u=-Y_d=b$（任意）で、第4式は $b^3-b^3=0$ と自動的にみたされる。これはクォークだけがベクトル的な電荷をもつ解である。

(3) 湯川結合の不変性は
\[
Y_Q+Y_u+Y_H=0,\qquad Y_Q+Y_d-Y_H=0,\qquad Y_L+Y_e-Y_H=0 .
\]
(1) の解で $Y_u=-4a$, $Y_d=2a$ とすると第1式から $Y_H=3a$、第2式 $a+2a-3a=0$、第3式 $-3a+6a-3a=0$ ですべてみたされる。逆の割り当て $Y_u=2a$, $Y_d=-4a$ では第1式から $Y_H=-3a$、第2式は $a-4a+3a=0$ でみたされるが、第3式は $-3a+6a+3a=6a\ne0$ となり不適。(2) の解では第1式から $Y_H=-b$ で、第2式 $0+Y_d-Y_H=-b+b=0$ はみたされるが、第3式 $0+0-Y_H=0$ から $Y_H=0$、よって $b=0$：自明な $U(1)$ しか残らない。$Y_H=3a=\frac12$ とすれば $a=\frac16$ で
\[
Y_Q=\tfrac16,\ Y_u=-\tfrac23,\ Y_d=\tfrac13,\ Y_L=-\tfrac12,\ Y_e=1,
\]
標準模型の値が再現される。全体の規格化以外にハイパーチャージの自由度はない。

(4) $Q=(u,d)$：$Q_{\mathrm{em}}(u)=\frac12+\frac16=\frac23$、$Q_{\mathrm{em}}(d)=-\frac12+\frac16=-\frac13$。$L=(\nu,e)$：$0$, $-1$。$u^c$, $d^c$, $e^c$ は $-\frac23$, $\frac13$, $1$（反粒子の電荷）。陽子 $uud$ は $\frac23+\frac23-\frac13=1$、電子 $-1$ で、水素原子は厳密に中性。

これが規格化 $Y_H=\frac12$ によらないことを見るため、一般の $a$ で考える。破れずに残る電荷は、ヒッグスの真空期待値をもつ成分（$T_3=-\frac12$）を中性にする組合せ $Q_{\mathrm{em}}=T_3+cY$ で、$-\frac12+c\cdot3a=0$ より $c=\frac{1}{6a}$。すると
\begin{gather*}
Q_{\mathrm{em}}(u)=\tfrac12+\tfrac{a}{6a}=\tfrac23,\quad Q_{\mathrm{em}}(d)=-\tfrac12+\tfrac16=-\tfrac13,\\
Q_{\mathrm{em}}(e)=-\tfrac12+\tfrac{-3a}{6a}=-1,\quad Q_{\mathrm{em}}(\nu)=0
\end{gather*}
となり、$a$ によらない。したがって $Q_{\mathrm{em}}(p)+Q_{\mathrm{em}}(e)=0$ は、ゲージ・アノマリーの相殺（と湯川結合の存在）だけから導かれる。電荷の量子化と原子の中性は、大統一理論を仮定しなくてもアノマリー相殺の帰結として理解できる。
\end{solution}

\begin{problem}[難易度 C]\label{ch09-prob-descent}
電荷1の $U(1)$ ゲージ場 $A$（$F=dA$）、ゲージ変換 $\delta_\lambda A=d\lambda$ とする。\\
(1) 2次元：$I_4=\frac{1}{8\pi^2}F^2$ から $I_3^{(0)}=\frac1{8\pi^2}A\,dA$、$I_2^{(1)}=\frac1{8\pi^2}\lambda\,dA$ を導け。$2\pi\int I_2^{(1)}$ を、問題 \ref{ch09-prob-fujikawa2d} の指数密度 $\frac{1}{2\pi}F$（ワイル・フェルミオン1個の\term{共変的アノマリー}{covariant anomaly}。ディラック・フェルミオンの軸性アノマリーはその2倍）と比べ、ワイル・フェルミオン1個の整合的アノマリーは「共変的アノマリー $\frac1{2\pi}F$」の $\frac12$ であることを確かめよ。\\
(2) 4次元：$I_6=\frac1{48\pi^3}F^3$ について同様に $I_5^{(0)}$, $I_4^{(1)}$ を求め、$2\pi I_4^{(1)}=\frac{1}{24\pi^2}\lambda F\wedge F$ を示せ。共変的アノマリー（指数密度 $\frac1{8\pi^2}F\wedge F$）との比が $\frac13$ であることを確かめよ。\\
(3) チャーン・サイモンズ形式の取り方の不定性 $I^{(0)}_{d+1}\to I^{(0)}_{d+1}+dX_d$ のもとで、$I^{(1)}_d$ が $\delta_\lambda X_d$（＋完全形式）だけ変わることを示し、これが有効作用への局所的な counterterm $\int X_d$ の付加に対応することを説明せよ。\\
(4) $U(1)$ でワイル・フェルミオン（電荷 $q_i$）が複数あるとき、2次元・4次元のゲージ・アノマリー相殺条件がそれぞれ $\sum_i\pm q_i^2=0$, $\sum_iq_i^3=0$ であることを説明せよ（2次元の符号はカイラリティ）。
\end{problem}

\begin{solution}
(1) $d(A\,dA)=dA\,dA=F^2$ なので $I_3^{(0)}=\frac1{8\pi^2}AdA$。$\delta_\lambda(AdA)=d\lambda\,dA$（$\delta F=0$）$=d(\lambda\,dA)$ なので $I_2^{(1)}=\frac1{8\pi^2}\lambda\,dA$。よって
\[
2\pi\int I_2^{(1)}=\frac1{4\pi}\int\lambda F .
\]
ディラック・フェルミオン（正・負カイラリティ各1）の軸性アノマリー $\frac1{2\pi}F\cdot2$（$\Delta Q_5=2n$）は、各カイラリティについて $\pm\frac{1}{2\pi}F$ ずつの寄与の差である。これが共変的アノマリーで、ワイル・フェルミオン1個の整合的アノマリー $\frac1{4\pi}F$ はその $\frac12$。

(2) $d(A(dA)^2)=(dA)^3=F^3$ なので $I_5^{(0)}=\frac1{48\pi^3}A\,F^2$。$\delta_\lambda I_5^{(0)}=\frac1{48\pi^3}d\lambda\,F^2=d\bigl(\frac1{48\pi^3}\lambda F^2\bigr)$（$dF=0$）なので $I_4^{(1)}=\frac{1}{48\pi^3}\lambda F^2$、$2\pi I_4^{(1)}=\frac1{24\pi^2}\lambda F\wedge F$。指数密度 $\frac1{8\pi^2}F\wedge F$（正カイラリティ1個のゼロモードの数に対応）との比は $\frac{8\pi^2}{24\pi^2}=\frac13$。成分では、整合的アノマリー $\frac1{96\pi^2}\epsilon^{\mu\nu\rho\sigma}F_{\mu\nu}F_{\rho\sigma}$ と共変的アノマリー $\frac1{32\pi^2}\epsilon^{\mu\nu\rho\sigma}F_{\mu\nu}F_{\rho\sigma}$。三角形図の3つの頂点を対称に扱う（ボーズ対称性）と、アノマリーが3つの頂点に $\frac13$ ずつ分配されることにあたる。

(3) $\delta_\lambda(I^{(0)}_{d+1}+dX_d)=dI^{(1)}_d+d\,\delta_\lambda X_d$ なので、新しい $I^{(1)}_d$ は $I^{(1)}_d+\delta_\lambda X_d$（完全形式の不定性を除く）。有効作用を $W\to W\mp2\pi i\int X_d$ と局所項で修正すると、そのゲージ変換は $\mp2\pi i\int\delta_\lambda X_d$ だけ変わるので、アノマリーの表式の変化と一致する。したがってアノマリーとは「局所的な counterterm で消せない部分」であり、それは $I_{d+2}$（ゲージ不変な特性類）だけで決まる。例えば4次元のディラック・フェルミオン（ベクトル場 $V$ と軸性場 $A$ に結合）では、バーディーンの counterterm によりアノマリーをすべて軸性の側に移し、ベクトル電流を保存させることができる。

(4) アノマリー多項式は加法的で、電荷 $q$ では $F\to qF$。2次元：$I_4=\sum_i(\pm)\frac{q_i^2}{8\pi^2}F^2$（符号はカイラリティ）なので条件は $\sum_\pm q_i^2=0$（右巻きと左巻きの電荷の2乗和が等しい）。4次元：すべて左巻きにそろえると $I_6=\sum_i\frac{q_i^3}{48\pi^3}F^3$ で条件は $\sum q_i^3=0$。これは定理 \ref{thm:ch09-pertcond} の $U(1)^3$ 条件に一致する（重力との混合アノマリーからは $\sum q_i=0$）。
\end{solution}

\begin{problem}[難易度 C]\label{ch09-prob-witten}
(1) 表現 $R$ が自分自身の複素共役と同値（$R\cong\bar R$）なら、$\tr_R(T^a\{T^b,T^c\})=0$ であることを示せ。$\bar R$ の生成子は $-(T^a)^*$ であること、$\tr(T^a\{T^b,T^c\})$ が実数であることを用いよ。特に $SU(2)$ には摂動論的なゲージ・アノマリーがないことを結論せよ。\\
(2) $SU(2)$ ゲージ理論で、左巻きワイル二重項が $n_D$ 個ある。$\pi_4(SU(2))=\Z_2$（第3章）を認め、定理 \ref{thm:ch09-witten} の論理を次の順に説明せよ：(i) 大きなゲージ変換 $U$ の存在、(ii) $A$ と $A^U$ を結ぶ道に沿ったスペクトル流、(iii) $Z=(\Pf)^{n_D}$ の符号、(iv) $n_D$ が奇数なら $Z\equiv0$。\\
(3) 標準模型、および「$SU(2)$ に二重項を1個だけ加えた理論」について判定せよ。\\
(4) スピン $j$ 表現のワイル・フェルミオン1個がウィッテン・アノマリーをもつのはどのような $j$ か。問題 \ref{ch09-prob-reps} の結果を用いて $j=\frac12,\frac32,\frac52$ を判定せよ。\\
(5) 「トフーフト頂点のフェルミオン線の数が奇数だと $(-1)^F$ が破れる」という観点から、(2) の結論を直観的に説明せよ。
\end{problem}

\begin{solution}
(1) $R\cong\bar R$ なら、ユニタリな $U$ があって $-(T^a)^*=UT^aU^{-1}$。よって
\[
\tr\bigl(T^a\{T^b,T^c\}\bigr)=\tr\bigl(UT^aU^{-1}\{UT^bU^{-1},UT^cU^{-1}\}\bigr)=-\tr\bigl((T^a)^*\{(T^b)^*,(T^c)^*\}\bigr)
\]
\[
=-\overline{\tr\bigl(T^a\{T^b,T^c\}\bigr)} .
\]
一方 $T$ はエルミートなので $\overline{\tr(T^aT^bT^c)}=\tr((T^aT^bT^c)^\dagger)=\tr(T^cT^bT^a)$ で、$\overline{\tr(T^a\{T^b,T^c\})}=\tr(T^cT^bT^a+T^bT^cT^a)=\tr(T^aT^cT^b+T^aT^bT^c)$（巡回性）$=\tr(T^a\{T^b,T^c\})$：実数。したがって $x=-x$ で $x=0$。$SU(2)$ の既約表現（スピン $j$）はすべて実（$j$ 整数）または擬実（$j$ 半整数）なので、$SU(2)$ には摂動論的なアノマリーがない。

(2) (i) $\pi_4(SU(2))=\Z_2$ なので、$S^4$（コンパクト化した時空）から $SU(2)$ への写像で、定値写像に連続変形できないもの $U$ がある。$A^U=UAU^{-1}+iUdU^{-1}$ は $A$ とゲージ同値。

(ii) ゲージ場の空間はアフィン空間なので、$A_t=(1-t)A+tA^U$ で結べる。$t$ に沿った（4次元の）ディラック作用素の固有値は $\pm\lambda_n(t)$ の対で現れ、$t=0$ と $t=1$ でスペクトルは一致する（ゲージ同値）。定理 \ref{thm:ch09-witten} によれば、二重項1個につき奇数対の固有値が入れ替わる（$\lambda_1(0)>0$ が $\lambda_1(1)<0$ へ；冒頭の図右）。これは5次元の時空 $S^4\times S^1$（$t$ 方向を周期化し $U$ で貼り合わせた写像トーラス）上の mod 2 指数として位相的に定まる量である。

(iii) ワイル二重項は擬実なので、その経路積分は（ディラック二重項の行列式 $\prod_n\lambda_n^2$ の平方根である）パフィアン $\prod_n\lambda_n$（各対から1つ）で定義される。これを $t$ に沿って連続に追跡すると、奇数個の $\lambda_n$ が符号を変えるので $\Pf(A^U)=-\Pf(A)$。$n_D$ 個なら $Z[A^U]=(-1)^{n_D}Z[A]$。

(iv) ゲージ不変性は $Z[A^U]=Z[A]$ を要求するので、$n_D$ 奇数では $Z[A]=-Z[A]$、すなわち $Z\equiv0$。期待値 $\langle\mathcal O\rangle=Z_{\mathcal O}/Z$ が $0/0$ となり理論は定義されない。

(3) 標準模型：世代あたり二重項 $Q$（3色）＋$L$ で4個、偶数なので問題なし（第3世代までで12個）。$SU(2)$ に二重項1個だけの理論は $n_D=1$ で矛盾する。例えば標準模型にベクトル的でない二重項を1個追加すると、ハイパーチャージのアノマリーとは独立にこの条件で排除される。

(4) $k=1$ のゼロモード数 $2T(j)$ が奇数のとき、スペクトル流が奇数となりアノマリーをもつ。$2T(\frac12)=1$（奇数：アノマリーあり）、$2T(\frac32)=10$（偶数：なし）、$2T(\frac52)=35$（奇数：あり）。一般に $2T(j)=\frac23j(j+1)(2j+1)$ で、$r\in\Z_{\ge0}$ として
\[
j=2r+\tfrac12:\ 2T=\frac{(4r+1)(4r+3)(2r+1)}{3},
\]
\[
j=2r+\tfrac32:\ 2T=2\cdot\frac{(4r+3)(4r+5)(r+1)}{3}.
\]
前者は奇数の積を3で割った整数なので奇数、後者は偶数。よって $j=\frac12,\frac52,\frac92,\dots$ がアノマリーをもつ。整数 $j$ では $2T=4\cdot\frac{j(j+1)(2j+1)}{6}=4\sum_{m=1}^jm^2$ で偶数（表現は実）。

(5) $k=1$ のインスタントンは各二重項から1本ずつ、計 $n_D$ 本のフェルミオン線をもつトフーフト頂点を生む。$n_D$ が奇数だと、真空からフェルミオン数が奇数だけ変化する振幅が現れる。しかし $2\pi$ 回転 $(-1)^F$ はローレンツ群（$SU(2)$ ゲージ変換の $-1$ とも一致する）の元であり、奇数個のフェルミオンを含む演算子の真空期待値はゼロでなければならない。この矛盾が「$Z=0$」の形で現れていると理解できる。
\end{solution}

\begin{problem}[難易度 D]\label{ch09-prob-inflow}
3次元の境界付き多様体 $M$（例えば円板 $\times$ 時間）上の $U(1)$ チャーン・サイモンズ作用
\[
S_{CS}=\frac{k}{4\pi}\int_MA\wedge dA\qquad(k\in\Z)
\]
を考える（量子ホール流体の低エネルギー有効作用）。\\
(1) $A\to A+d\lambda$ のもとで $\delta S_{CS}=\frac{k}{4\pi}\int_{\partial M}\lambda\,dA$ となることを示せ。\\
(2) 境界 $\partial M$ 上の電荷1のカイラル（ワイル）フェルミオン1個のゲージ変換は、問題 \ref{ch09-prob-descent}(1) より大きさ $\frac{1}{4\pi}\int\lambda F$ をもつ。全体がゲージ不変となるには、境界に正味 $k$ 個の（適切なカイラリティの）カイラル・フェルミオンが必要であることを説明せよ。\\
(3) $S_{CS}$ から電流 $j^\mu=\delta S_{CS}/\delta A_\mu$ を求め、$j^0=\frac{k}{2\pi}B$、$j^x=\frac{k}{2\pi}E_y$ を示せ（成分は $A\wedge dA=\epsilon^{\mu\nu\rho}A_\mu\partial_\nu A_\rho\,d^3x$、$\epsilon^{txy}=1$、$E_i=\partial_iA_t-\partial_tA_i$ とせよ）。ホール伝導度 $\sigma_{xy}=k/2\pi$（$\hbar$, $e$ を戻すと $ke^2/h$）を結論せよ。\\
(4) 円環（コルビノ円板）の穴に磁束を $0$ から $2\pi$ まで断熱的に挿入するとき（ラフリンの議論）、ちょうど $k$ 単位の電荷が内側の境界から外側の境界へ運ばれることを (3) から示し、これを境界のカイラル・フェルミオンのスペクトル流（本章冒頭の図）と関係づけよ。
\end{problem}

\begin{solution}
(1) $\delta(A\wedge dA)=d\lambda\wedge dA+A\wedge d(d\lambda)=d\lambda\wedge dA=d(\lambda\,dA)$。ストークスの定理より $\delta S_{CS}=\frac k{4\pi}\int_M d(\lambda dA)=\frac{k}{4\pi}\int_{\partial M}\lambda\,dA$。閉じた $M$ では不変だが、境界があるとゲージ不変性が破れる。

(2) 境界のカイラル・フェルミオン1個の有効作用は $\delta_\lambda W=\pm\frac{i}{4\pi}\int\lambda F$（ユークリッド。符号はカイラリティによる）、すなわち作用の言葉で大きさ $\frac1{4\pi}\int\lambda F$ だけ変化する。これは $k=1$ のときの (1) の境界項と同じ形である。したがって (1) の変化と打ち消し合うカイラリティのフェルミオンが（右巻きと左巻きの差として）正味 $k$ 個あれば、バルク＋境界の全体はゲージ不変になる。バルクのチャーン・サイモンズ項のアノマリーが境界に「流れ込み」、境界のアノマリーを相殺している（アノマリー流入）。境界のカイラル・モード単独ではゲージ・アノマリーをもつので、2次元の理論としては存在できず、3次元のバルクの境界としてのみ存在できる。これは整数量子ホール効果のエッジ状態がカイラルで「後方散乱で局在しない」理由でもある。

(3) 2形式は1形式と可換なので $A\wedge d\delta A=d\delta A\wedge A=d(\delta A\wedge A)+\delta A\wedge dA$。よって
\[
\delta S_{CS}=\frac{k}{4\pi}\int(\delta A\wedge dA+A\wedge d\delta A)=\frac{k}{2\pi}\int\delta A\wedge dA+(\text{境界項}),
\]
すなわち成分で $j^\mu=\frac{k}{2\pi}\epsilon^{\mu\nu\rho}\partial_\nu A_\rho$。
\[
j^t=\frac{k}{2\pi}(\partial_xA_y-\partial_yA_x)=\frac{k}{2\pi}B,
\]
\[
j^x=\frac{k}{2\pi}\bigl(\epsilon^{xty}\partial_tA_y+\epsilon^{xyt}\partial_yA_t\bigr)=\frac{k}{2\pi}(\partial_yA_t-\partial_tA_y)=\frac{k}{2\pi}E_y .
\]
（$\epsilon^{xty}=-1$, $\epsilon^{xyt}=+1$。）電場に垂直な電流が流れ、$\sigma_{xy}=k/2\pi$。単位を戻すと $\sigma_{xy}=k\,e^2/2\pi\hbar=k\,e^2/h$。

(4) 穴を貫く磁束 $\Phi(t)$ を変えると、ファラデーの法則により円環の周方向に電場 $\oint E_\phi\,d\ell=-\dot\Phi$ が生じる。(3) より動径方向の電流（ホール電流）が流れ、内側の境界から外側へ運ばれる電荷は（向きの符号を除き）
\[
\Delta q=\int dt\oint j^r\,d\ell=\frac{k}{2\pi}\int dt\oint E_\phi\,d\ell=\frac{k}{2\pi}\lvert\Delta\Phi\rvert=k .
\]
境界から見ると、磁束 $2\pi$ の挿入は境界の円周上で $A_\phi$ を1量子ずらすことにあたる。冒頭の図のように、各カイラル・フェルミオンの準位が1段ずれ、ディラックの海から1個の粒子が押し出される（スペクトル流）。外側の境界のカイラル・モード $k$ 個はそれぞれ電荷 $1$ を得、内側の境界（逆向きのカイラリティ）はそれぞれ $1$ を失う。各境界の電荷の非保存（2次元のアノマリー）は、バルクを流れるホール電流（流入）によって正確に補われている。
\end{solution}

%================================================
\chapter{弦理論・コンパクト化のトポロジー}\label{ch:10}
% 17問
%================================================

\section{物理からの動機}

\subsection*{世界面の位相による摂動展開}
弦理論の摂動論では、ファインマン図の代わりに弦の\term{世界面}{worldsheet}についての和をとる。ユークリッド化した世界面 $\Sigma$ 上のポリャコフ作用には、定数のディラトン $\Phi$ が
\[
S_\Phi=\frac{\Phi}{4\pi}\int_\Sigma d^2\sigma\sqrt{h}\,R+\frac{\Phi}{2\pi}\oint_{\partial\Sigma}ds\,k_g=\Phi\,\chi(\Sigma)
\]
の形で入る（$R$ は世界面のスカラー曲率、$k_g$ は境界の測地的曲率）。2番目の等号は境界付きのガウス・ボンネの定理（第8章）である。したがって振幅は $e^{-S_\Phi}=g_s^{-\chi(\Sigma)}$（$g_s=e^\Phi$ は\term{弦の結合定数}{string coupling}）で重みづけられ、閉弦の真空振幅は
\[
Z=\sum_{g=0}^\infty g_s^{2g-2}Z_g
\]
と\term{種数展開}{genus expansion}される。ループの数を数えていたファインマン図のオイラー標数（第4章）の役割を、世界面のオイラー標数が果たすのである。
\begin{center}
\begin{tikzpicture}[scale=0.9]
  \shade[ball color=blue!8] (0,0) circle (0.75);
  \draw[thick] (0,0) circle (0.75);
  \draw[gray] (-0.75,0) arc (180:360:0.75 and 0.2);
  \draw[gray,dashed] (0.75,0) arc (0:180:0.75 and 0.2);
  \node at (0,-1.2) {\footnotesize $g=0$：$g_s^{-2}$};
  \node at (1.45,0) {$+$};
  \begin{scope}[shift={(3.3,0)}]
    \shade[ball color=blue!8] (0,0) ellipse (1.25 and 0.75);
    \draw[thick] (0,0) ellipse (1.25 and 0.75);
    \draw[thick] (-0.45,0.05) to[bend right=35] (0.45,0.05);
    \draw[thick,fill=white] (-0.34,-0.01) to[bend left=35] (0.34,-0.01) to[bend left=20] (-0.34,-0.01);
    \node at (0,-1.2) {\footnotesize $g=1$：$g_s^{0}$};
  \end{scope}
  \node at (5.15,0) {$+$};
  \begin{scope}[shift={(7.6,0)}]
    \shade[ball color=blue!8] (0,0) ellipse (1.9 and 0.8);
    \draw[thick] (0,0) ellipse (1.9 and 0.8);
    \foreach \x in {-0.8,0.8}{
      \draw[thick] (\x-0.42,0.05) to[bend right=35] (\x+0.42,0.05);
      \draw[thick,fill=white] (\x-0.32,-0.01) to[bend left=35] (\x+0.32,-0.01) to[bend left=20] (\x-0.32,-0.01);
    }
    \node at (0,-1.2) {\footnotesize $g=2$：$g_s^{2}$};
  \end{scope}
  \node at (10.4,0) {$+\ \cdots$};
\end{tikzpicture}
\zucap{閉弦の真空振幅の種数展開。取っ手が1つ増えるごとに $\chi$ は $2$ 減り、$g_s^2$ が掛かる。開弦・向き付け不可能な弦では、境界や交差帽を持つ世界面も現れる（問題 \ref{ch10-prob-ws}）。}
\end{center}

\subsection*{コンパクト化：4次元の物理を内部空間のトポロジーが決める}
超弦理論は10次元時空で定義されるので、時空を $M_4\times K$（$K$ はコンパクトな6次元多様体）とし、$K$ の大きさより長い波長で見た4次元の有効理論を考える。このとき
\begin{itemize}
\item 10次元の $p$ 形式場を $K$ 上の\textbf{調和形式}で展開すると、4次元の質量ゼロ場が得られる。その個数はベッチ数 $b_k(K)$、$K$ が複素多様体ならホッジ数 $h^{p,q}(K)$ で決まる（ホッジの定理、第5章）。
\item 4次元で何個の超対称性が残るかは $K$ の\textbf{ホロノミー群}で決まる。$\mathcal N=1$ が残る代表例が\term{カラビ・ヤウ多様体}{Calabi--Yau manifold}である。
\item ヘテロティック弦をカラビ・ヤウ多様体 $X$ にコンパクト化すると、カイラルなクォーク・レプトンの\textbf{世代数}が $\lvert\chi(X)\rvert/2$ になる。これは指数定理（第9章）の帰結である。
\item 弦の真空の連続パラメータ（\term{モジュライ}{moduli}）の数は $h^{1,1}$, $h^{2,1}$ で数えられる。
\item フラックスやDブレーンの電荷は、一般化されたディラックの量子化により\textbf{整数係数コホモロジー}の元になる。コンパクトな空間ではガウスの法則により総電荷がゼロでなければならず（\term{タドポール相殺}{tadpole cancellation}）、その条件に $\chi/24$ のような特性数が現れる。
\item 10次元のアノマリーの相殺（グリーン・シュワルツ機構）がゲージ群を $SO(32)$ か $E_8\times E_8$ に限る。
\end{itemize}
本章では、これらを支える代数トポロジー・複素幾何の道具を、弦理論の具体例（トーラス、K3、5次超曲面、オービフォールド）で計算する。

\begin{physbox}[：本章の辞書]
\begin{center}
\small
\begin{tabular}{p{0.46\linewidth}p{0.44\linewidth}}
物理 & 数学\\ \hline
弦の摂動の次数 & 世界面のオイラー標数 $\chi=2-2g-b-c$\\
4次元の質量ゼロ場の数 & ベッチ数、ホッジ数 $h^{p,q}$\\
残る超対称性 & ホロノミー群、共変定数スピノル\\
世代数 & 指数 $\int\ch(V)\Td(X)=\chi/2$\\
ケーラー・モジュライ／複素構造モジュライ & $h^{1,1}$ ／ $h^{2,1}$\\
オービフォールドの「ねじれたセクター」 & 固定点、特異点の解消\\
フラックス量子・ブレーン電荷 & $H^{k}(K;\Z)$ の元\\
タドポール相殺 & コンパクト空間上のガウス則、$\chi/24$\\
グリーン・シュワルツ機構 & $I_{12}$ の因数分解、$\dim G=496$\\
ミラー対称性 & $h^{1,1}\leftrightarrow h^{2,1}$
\end{tabular}
\end{center}
\end{physbox}

\section{基本事項のまとめ}

\subsection{曲面の分類と世界面}

\begin{mythm}{コンパクト曲面の分類}{ch10-surf}
コンパクトで連結な曲面（境界があってもよい）は、球面 $S^2$ に $g$ 個の取っ手（\term{ハンドル}{handle}）を付け、$b$ 個の円板を取り除き（\term{境界}{boundary}）、$c$ 個の\term{交差帽}{crosscap}（円板を取り除いて境界の対蹠点を同一視したもの）を付けて得られる。そのオイラー標数は
\[
\chi=2-2g-b-c
\]
である。$c\ge1$ なら向き付け不可能。取っ手1つと交差帽1つの組は交差帽3つと同じ曲面を与えるので、$c\ge1$ のときは $g=0$ としてよい。
\end{mythm}
\begin{note}
$\chi$ の公式は、円板を取り除くと（面が1つ減るので）$\chi$ が $1$ 減ること、交差帽は円板を取り除いて $\RP^2$ から円板を除いたもの（メビウスの帯、$\chi=0$）を貼ることで $\chi$ を $1$ 減らすこと、取っ手は $\chi$ を $2$ 減らすことから従う。一般に $\chi(A\cup B)=\chi(A)+\chi(B)-\chi(A\cap B)$（第4章のMV完全列）と $\chi(S^1)=0$ を使えばよい。
\end{note}

\begin{physbox}[：弦の種類と世界面]
\begin{itemize}
\item 閉じた向き付けられた弦（II型、ヘテロティック）：閉じた向き付け可能曲面 $\Sigma_g$ のみ。重み $g_s^{2g-2}$。
\item 開弦を含む理論：境界 $b\ge1$ が現れる。境界は開弦の端点の軌跡で、Dブレーン上にある。
\item 向き付けられていない弦（I型）：交差帽 $c\ge1$ も現れる。交差帽はオリエンティフォールド面（O面）との相互作用を表す。
\end{itemize}
$\chi$ を固定すると $2g+b+c=2-\chi$ なので、各次数に寄与する世界面の位相型は有限個である。
\end{physbox}

\subsection{複素多様体とホッジ数}

\begin{mydef}{複素多様体とドルボー・コホモロジー}{ch10-dolb}
$\C^n$ の開集合を正則な座標変換で貼り合わせたものを複素 $n$ 次元の\term{複素多様体}{complex manifold}という。局所座標 $z^i$ で微分形式は $dz^{i_1}\wedge\dots\wedge dz^{i_p}\wedge d\bar z^{j_1}\wedge\dots\wedge d\bar z^{j_q}$ の線形結合として $(p,q)$ 型に分解され、外微分は $d=\partial+\bar\partial$ と分かれる。$\bar\partial^2=0$ なので
\[
H^{p,q}_{\bar\partial}(X)=\frac{\{\bar\partial\text{-閉な }(p,q)\text{ 形式}\}}{\bar\partial\{(p,q-1)\text{ 形式}\}}
\]
が定まる（\term{ドルボー・コホモロジー}{Dolbeault cohomology}）。$h^{p,q}=\dim_\C H^{p,q}_{\bar\partial}(X)$ を\term{ホッジ数}{Hodge number}という。
\end{mydef}

\begin{mydef}{ケーラー多様体}{ch10-kahler}
エルミート計量 $g_{i\bar j}$ をもつ複素多様体で、2形式 $J=ig_{i\bar j}\,dz^i\wedge d\bar z^j$ が $dJ=0$ をみたすものを\term{ケーラー多様体}{K\"ahler manifold}、$J$ を\term{ケーラー形式}{K\"ahler form}という。コンパクトなケーラー多様体では $\int_XJ^n>0$ なので $[J]\neq0\in H^2(X;\R)$、特に $b_{2k}\ge1$（$0\le k\le n$）。
\end{mydef}
$\C^n$、トーラス $\C^n/\Lambda$、$\CP^n$（フビニ・スタディ計量）、$\CP^n$ の複素部分多様体はすべてケーラー多様体である。

\begin{mythm}{ホッジ分解}{ch10-hodge}
コンパクトなケーラー多様体 $X$（複素 $n$ 次元）について
\[
H^k(X;\C)\cong\bigoplus_{p+q=k}H^{p,q}(X),\qquad b_k=\sum_{p+q=k}h^{p,q},
\]
\[
h^{p,q}=h^{q,p}\ (\text{複素共役}),\qquad h^{p,q}=h^{n-p,n-q}\ (\text{セール双対性}).
\]
各類はラプラシアン $\Delta_d=2\Delta_{\bar\partial}$ について調和な $(p,q)$ 形式を一意に代表元にもつ。特に奇数次のベッチ数は偶数。
\end{mythm}
\begin{note}
ホッジの定理（第5章）により $H^k(X;\R)$ は調和 $k$ 形式の空間と同型である。ケーラー条件のもとでは $\Delta_d$ が $(p,q)$ 型を保つので、調和形式が型ごとに分解する。$h^{p,q}=h^{q,p}$ は調和形式の複素共役が調和であることから、$h^{p,q}=h^{n-p,n-q}$ はホッジの $*$ と複素共役の組み合わせ（あるいはポアンカレ双対、第6章）から従う。
\end{note}

ホッジ数は\term{ホッジ・ダイヤモンド}{Hodge diamond}と呼ばれる菱形に並べる。$k$ 段目（上から $0$ 段目）に $p+q=k$ の数を並べ、横の和がベッチ数 $b_k$ になる。

\begin{center}
\begin{tikzpicture}[x=0.62cm,y=0.5cm]
  \foreach \p/\q/\t in {0/0/1, 1/0/0, 0/1/0, 2/0/0, 1/1/h^{1,1}, 0/2/0,
      3/0/1, 2/1/h^{2,1}, 1/2/h^{2,1}, 0/3/1, 3/1/0, 2/2/h^{1,1}, 1/3/0,
      3/2/0, 2/3/0, 3/3/1}{
    \node at ({\q-\p},{-(\p+\q)}) {\small $\t$};
  }
  \node[blue!60!black] at (-0.2,1.2) {\footnotesize 一般のCY3};
  \begin{scope}[xshift=6.2cm]
  \foreach \p/\q/\t in {0/0/1, 1/0/0, 0/1/0, 2/0/0, 1/1/1, 0/2/0,
      3/0/1, 2/1/101, 1/2/101, 0/3/1, 3/1/0, 2/2/1, 1/3/0,
      3/2/0, 2/3/0, 3/3/1}{
    \node at ({\q-\p},{-(\p+\q)}) {\small $\t$};
  }
  \node[blue!60!black] at (0,1.2) {\footnotesize 5次超曲面 $X_5\subset\CP^4$};
  \end{scope}
\end{tikzpicture}
\zucap{カラビ・ヤウ3次元多様体（CY3）のホッジ・ダイヤモンド。$h^{p,q}$ を座標 $(q-p,\,-(p+q))$ に置いた。左右対称は $h^{p,q}=h^{q,p}$、上下対称はセール双対性 $h^{p,q}=h^{3-p,3-q}$。}
\end{center}

\begin{myex}{トーラスと射影空間のホッジ数}{ch10-hodgeex}
\begin{itemize}
\item 平坦なトーラス $T^{2n}=\C^n/\Lambda$：調和形式は定数係数の形式なので $h^{p,q}=\binom{n}{p}\binom{n}{q}$（問題 \ref{ch10-prob-torus}）。$T^6$ では $h^{1,0}=3$, $h^{2,0}=3$, $h^{1,1}=9$, $h^{2,1}=9$, $h^{3,0}=1$。
\item $\CP^n$：$h^{p,p}=1$（$0\le p\le n$）、それ以外は $0$。$H^{2p}$ はケーラー形式のべき $J^p$ で生成される。
\end{itemize}
\end{myex}

\subsection{カラビ・ヤウ多様体}

\begin{mydef}{カラビ・ヤウ多様体}{ch10-cy}
コンパクトなケーラー多様体 $X$ で第1チャーン類 $c_1(X)=c_1(TX)$ が（実係数で）ゼロであるものを\term{カラビ・ヤウ多様体}{Calabi--Yau manifold}という。本書では特に、単連結で、ホロノミーがちょうど $SU(n)$ のものを考える。複素 $3$ 次元のものをCY3と略す。
\end{mydef}

\begin{mythm}{カラビ予想（ヤウの定理）と帰結}{ch10-yau}
$c_1(X)=0$ のコンパクトなケーラー多様体の各ケーラー類 $[J]$ には、リッチ平坦なケーラー計量がただ1つ存在する。その計量のホロノミーは $SU(n)$ に含まれ、次が成り立つ。
\begin{enumerate}
\item 至るところゼロにならない正則 $n$ 形式 $\Omega$ が存在する（標準束 $K_X=\Lambda^nT^*X$ が自明）。ホロノミーがちょうど $SU(n)$ なら $h^{n,0}=1$、$h^{p,0}=0$（$0<p<n$）。
\item \term{共変定数スピノル}{covariantly constant spinor} $\nabla\eta=0$ が存在する（$U(n)$ ホロノミーではスピノルは $\Lambda^{0,*}$ のように分解し、$SU(n)$ に縮小すると $\Lambda^{0,0}$ と $\Lambda^{0,n}$ の成分が不変になる）。
\end{enumerate}
\end{mythm}

\begin{myprop}{CY3のホッジ数とモジュライ}{ch10-cy3}
ホロノミー $SU(3)$ のCY3では、ホッジ・ダイヤモンドは上図の形になり、独立なホッジ数は $h^{1,1}$ と $h^{2,1}$ の2つだけである。
\[
b_0=b_6=1,\quad b_1=b_5=0,\quad b_2=b_4=h^{1,1},\quad b_3=2+2h^{2,1},
\]
\[
\chi(X)=2\bigl(h^{1,1}-h^{2,1}\bigr).
\]
リッチ平坦計量の変形は、$h^{1,1}$ 個の実パラメータ（\term{ケーラー・モジュライ}{K\"ahler moduli}：$J$ の変形）と $h^{2,1}$ 個の複素パラメータ（\term{複素構造モジュライ}{complex structure moduli}）からなる。後者は $H^1(X,TX)\cong H^{2,1}(X)$（$\Omega$ との縮約による同型）で数えられる。
\end{myprop}

\begin{myex}{代表的なカラビ・ヤウ多様体}{ch10-cyex}
\begin{center}
\small
\renewcommand{\arraystretch}{1.15}
\begin{tabular}{lll}
多様体 & ホッジ数など & $\chi$\\ \hline
楕円曲線 $T^2$（$\CP^2$ の3次曲線） & $h^{1,0}=1$ & $0$\\
K3曲面（$\CP^3$ の4次曲面など） & $h^{2,0}=1$, $h^{1,1}=20$, $b_2=22$ & $24$\\
5次超曲面 $X_5\subset\CP^4$ & $(h^{1,1},h^{2,1})=(1,101)$ & $-200$\\
ミラー5次超曲面 & $(101,1)$ & $200$\\
$T^6/\Z_3$ の解消 & $(36,0)$ & $72$
\end{tabular}
\end{center}
$T^6$ は $c_1=0$ だがホロノミーが自明で、$h^{1,0}=3\ne0$ なので上の意味のCY3ではない。
\end{myex}

\subsection{チャーン類と添加公式}
チャーン類・ガウス・ボンネ・チャーンの定理の一般論は第8章で扱った。本章では次を用いる。複素 $n$ 次元のコンパクト複素多様体について
\[
\chi(X)=\int_Xc_n(TX).
\]

\begin{myprop}{射影空間と超曲面の全チャーン類}{ch10-adj}
$H\in H^2(\CP^N;\Z)$ を超平面の類（ポアンカレ双対）とすると $\int_{\CP^N}H^N=1$ であり
\[
c(T\CP^N)=(1+H)^{N+1}.
\]
$d$ 次の斉次多項式の零点として定まる滑らかな超曲面 $X\subset\CP^N$（複素 $N-1$ 次元）について、\term{添加公式}{adjunction formula}
\[
c(TX)=\frac{(1+H)^{N+1}}{1+dH}\Big|_X,\qquad \int_XH^{N-1}=d
\]
が成り立つ。特に $c_1(TX)=(N+1-d)H$ なので、$X$ がカラビ・ヤウになるのは $d=N+1$ のときである。
\end{myprop}
\begin{note}
$T\CP^N$ の式はオイラー完全列 $0\to\underline{\C}\to\mathcal O(1)^{\oplus(N+1)}\to T\CP^N\to0$ から、超曲面の式は $0\to TX\to T\CP^N|_X\to N_X\to0$ と法束 $N_X=\mathcal O(d)|_X$（$c(N_X)=1+dH$）から、チャーン類の積公式 $c(E)=c(E')c(E'')$ によって従う。$\int_XH^{N-1}=d$ は、$X$ と一般の位置にある $N-1$ 枚の超平面の交点（＝$d$ 次方程式1つと1次方程式 $N-1$ 個の共通解）が $d$ 個であることを表す。
\end{note}

\begin{mythm}{レフシェッツの超平面定理}{ch10-lef}
$X\subset\CP^N$ を滑らかな超曲面（複素 $n=N-1$ 次元）とし、$i\colon X\hookrightarrow\CP^N$ を包含とする。$i^*\colon H^k(\CP^N;\Z)\to H^k(X;\Z)$ は $k<n$ で同型、$k=n$ で単射である。
\end{mythm}
例えば5次超曲面（$n=3$）では $H^1(X)=0$、$H^2(X)\cong\Z$（$H$ で生成）。

\subsection{オービフォールド}

\begin{mydef}{大域的オービフォールド}{ch10-orb}
多様体 $X$ に有限群 $G$ が作用しているとき、商空間 $X/G$ を\term{オービフォールド}{orbifold}という。$g\in G$ の固定点集合を $X^g=\{x:gx=x\}$ と書く。$G$ が自由に作用する（$g\ne1$ なら $X^g=\emptyset$）とき $X/G$ は多様体で、$X\to X/G$ は $\lvert G\rvert$ 重の被覆である（第2章）。固定点があると、その像は一般に特異点（例：$\C^2/\Z_2$ は $\RP^3$ 上の錐）になる。
\end{mydef}

\begin{mythm}{オービフォールドのオイラー標数}{ch10-orbchi}
\begin{enumerate}
\item（位相的な商空間）$\displaystyle\chi(X/G)=\frac{1}{\lvert G\rvert}\sum_{g\in G}\chi(X^g)$。自由な作用なら $\chi(X/G)=\chi(X)/\lvert G\rvert$。
\item（弦理論のオービフォールド；ディクソン・ハーヴェイ・ヴァッファ・ウィッテン）弦の真空の「オイラー標数」は
\[
\chi_{\rm str}(X,G)=\frac{1}{\lvert G\rvert}\sum_{\substack{g,h\in G\\ gh=hg}}\chi\bigl(X^g\cap X^h\bigr).
\]
$X$ がCYで $G$ が正則体積形式を保つとき（複素次元 $\le3$）、これは特異点をカラビ・ヤウ性を保って解消した多様体（\term{クレパント解消}{crepant resolution}）のオイラー標数に等しい（ここでは認める）。
\end{enumerate}
\end{mythm}
\begin{note}
(1) は $G$ と両立する三角形分割をとり、各単体を「その単体を固定する元の数」で数え直すことで示される（バーンサイドの補題の位相版）。(2) の和で $h=1$ の項は (1) を再現する（\term{ねじれていないセクター}{untwisted sector}）。$h\ne1$ の項は、$h$ でねじれた境界条件 $X(\sigma+2\pi)=hX(\sigma)$ をもつ閉弦（\term{ねじれたセクター}{twisted sector}）からの寄与で、固定点に局在している。
\end{note}

\begin{center}
\begin{tikzpicture}[scale=1.6]
  \fill[blue!6] (0,0) rectangle (2,1);
  \draw[thick] (0,0) rectangle (2,2);
  \draw[densely dashed,gray] (0,1)--(2,1);
  \foreach \x/\y in {0/0,1/0,0/1,1/1,2/0,2/1,0/2,1/2,2/2}{\node[vtx,red!70!black] at (\x,\y) {};}
  \node[below] at (0,0) {\footnotesize $0$};
  \node[below] at (1,0) {\footnotesize $\tfrac12$};
  \node[below] at (2,0) {\footnotesize $1$};
  \node[left] at (0,1) {\footnotesize $\tfrac{\tau}{2}$};
  \node[left] at (0,2) {\footnotesize $\tau$};
  \node[right] at (1,1.15) {\footnotesize $\tfrac{1+\tau}{2}$};
  \node at (1,-0.5) {\footnotesize $T^2$ と $z\mapsto-z$ の固定点};
  \begin{scope}[xshift=3.4cm,yshift=0.2cm]
    \fill[blue!6] (0,0) .. controls (0.6,-0.12) and (1.2,-0.12) .. (1.8,0)
       .. controls (1.92,0.5) and (1.92,1.1) .. (1.8,1.6)
       .. controls (1.2,1.72) and (0.6,1.72) .. (0,1.6)
       .. controls (-0.12,1.1) and (-0.12,0.5) .. (0,0);
    \draw[thick] (0,0) .. controls (0.6,-0.12) and (1.2,-0.12) .. (1.8,0)
       .. controls (1.92,0.5) and (1.92,1.1) .. (1.8,1.6)
       .. controls (1.2,1.72) and (0.6,1.72) .. (0,1.6)
       .. controls (-0.12,1.1) and (-0.12,0.5) .. (0,0);
    \draw[gray,densely dashed] (0,0) .. controls (0.5,0.35) and (1.3,0.35) .. (1.8,0);
    \foreach \x/\y in {0/0,1.8/0,1.8/1.6,0/1.6}{\node[vtx,red!70!black] at (\x,\y) {};}
    \node at (0.9,-0.7) {\footnotesize $T^2/\Z_2$（枕）$\cong S^2$};
  \end{scope}
\end{tikzpicture}
\zucap{左：$T^2=\C/(\Z+\tau\Z)$（図は $\tau=i$）の $\Z_2$ 作用 $z\mapsto-z$ の固定点（赤、格子で同一視すると4点）。影をつけた下半分が基本領域。右：商空間は4つの角（錐角 $\pi$ の円錐特異点）をもつ「枕」で、位相的には球面である（問題 \ref{ch10-prob-pillow}）。}
\end{center}

\begin{myex}{特異点の解消}{ch10-resol}
\begin{itemize}
\item $\C^2/\Z_2$（$z\mapsto-z$）：原点を $\CP^1\cong S^2$ に置き換えた $T^*S^2$（\term{江口・ハンソン空間}{Eguchi--Hanson space}）で解消できる。解消空間は $S^2$ にホモトピー同値で、例外曲線の自己交叉数は $-2$。
\item $\C^3/\Z_3$（$z\mapsto\omega z$, $\omega=e^{2\pi i/3}$）：原点を $\CP^2$ に置き換えた空間（$\CP^2$ 上の直線束 $\mathcal O(-3)$ の全空間）で解消できる。$\CP^2$ にホモトピー同値。
\end{itemize}
1点を $E$ に置き換えると、オイラー標数は $\chi(E)-1$ だけ増える（$\chi(A\cup B)=\chi(A)+\chi(B)-\chi(A\cap B)$ を用いる）。
\end{myex}

\subsection{ヘテロティック弦のコンパクト化と世代数}
ヘテロティック弦 $E_8\times E_8$ を、$\mathcal N=1$ の超対称性を保つためにCY3 $X$ にコンパクト化する（問題 \ref{ch10-prob-susy}）。ゲージ場の真空期待値として、$X$ のスピン接続（$SU(3)$ に値をとる）を一方の $E_8$ の部分群 $SU(3)$ に埋め込む（\term{標準埋め込み}{standard embedding}）。これはベクトル束 $V=TX$ の上の接続を選ぶことにあたり、後で述べるビアンキ恒等式の条件を自動的にみたす。$E_8\supset E_6\times SU(3)$ の分解
\[
\mathbf{248}=(\mathbf{78},\mathbf 1)\oplus(\mathbf 1,\mathbf 8)\oplus(\mathbf{27},\mathbf 3)\oplus(\overline{\mathbf{27}},\overline{\mathbf 3})
\]
により、4次元には $E_6$ の大統一理論が残る。

\begin{mythm}{世代数（候補・ホロヴィッツ・ストロミンジャー・ウィッテン）}{ch10-gen}
標準埋め込みでは、4次元の質量ゼロのカイラル多重項として $\mathbf{27}$ が $h^{2,1}(X)$ 個、$\overline{\mathbf{27}}$ が $h^{1,1}(X)$ 個現れる。正味の世代数は
\[
n_{\rm gen}=\bigl\lvert h^{2,1}-h^{1,1}\bigr\rvert=\frac{\lvert\chi(X)\rvert}{2}.
\]
これは $X$ 上の $V$ に結合したディラック作用素の指数 $\int_X\ch(V)\Td(X)=\tfrac12\int_Xc_3(V)$ に等しい（第9章、問題 \ref{ch10-prob-gen}）。
\end{mythm}

\subsection{$p$ 形式場、フラックス、ブレーン電荷}
弦理論には電磁場の一般化である $p$ 形式ゲージ場が現れる。$D$ 次元時空で空間 $p$ 次元に広がった物体（\term{$p$ ブレーン}{$p$-brane}、世界体積 $p+1$ 次元）は、$(p+1)$ 形式ゲージ場 $C_{p+1}$ に
\[
\exp\Bigl(i\mu\int_{W_{p+1}}C_{p+1}\Bigr)
\]
の形で結合する（$\mu$ は電荷）。場の強さ $F_{p+2}=dC_{p+1}$ はゲージ変換 $C\to C+d\lambda_p$ で不変である。$p=0$ が通常の荷電粒子と $A=C_1$ の結合 $\exp(iq\int A)$ である。

\begin{mythm}{一般化されたディラックの量子化}{ch10-dirac}
$D$ 次元時空で、電気的 $p$ ブレーン（電荷 $\mu_e$）と、それを囲む $(p+2)$ 次元球面 $S^{p+2}$ を貫く磁気的フラックス $\int_{S^{p+2}}F_{p+2}=\mu_m$ をもつ物体が共存するなら
\[
\mu_e\,\mu_m\in2\pi\Z .
\]
磁気的な物体は $(D-p-4)$ ブレーンである（$S^{p+2}$ で囲まれる物体の横方向は $p+3$ 次元なので、世界体積は $D-p-3$ 次元）。電荷の単位を $\mu_e=1$ に選ぶと、任意の閉 $(p+2)$ 次元部分多様体 $\Sigma$ について $\int_\Sigma F_{p+2}\in2\pi\Z$。
\end{mythm}
証明は第5章のディラックの議論と同じである（問題 \ref{ch10-prob-dirac}）。

\begin{center}
\begin{tikzpicture}[scale=1.1]
  \fill[orange!15] (-2.4,-0.5) -- (1.6,-0.5) -- (2.4,0.5) -- (-1.6,0.5) -- cycle;
  \draw[thick,orange!60!black] (-2.4,-0.5) -- (1.6,-0.5) -- (2.4,0.5) -- (-1.6,0.5) -- cycle;
  \node[orange!50!black] at (-1.6,-0.25) {\footnotesize D$p$};
  \shade[ball color=blue!15,opacity=0.55] (0.2,0) circle (1.05);
  \draw[blue!60!black,thick] (0.2,0) circle (1.05);
  \draw[blue!60!black,dashed] (1.25,0) arc (0:180:1.05 and 0.28);
  \draw[blue!60!black] (-0.85,0) arc (180:360:1.05 and 0.28);
  \node[vtx] at (0.2,0) {};
  \foreach \a in {30,90,150,210,270,330}{
    \draw[-{Stealth[length=1.6mm]},red!70!black] ($(0.2,0)+(\a:0.25)$) -- ($(0.2,0)+(\a:1.55)$);
  }
  \node[blue!60!black] at (2.2,1.1) {\footnotesize $S^{8-p}$};
  \node[red!70!black] at (-1.6,1.2) {\footnotesize $F_{8-p}={*F_{p+2}}$};
\end{tikzpicture}
\zucap{Dブレーンの電荷の測り方（模式図）。D$p$ ブレーンは10次元のうち $p+1$ 次元に広がり、横方向 $\R^{9-p}$ の1点に見える。その点を囲む球面 $S^{8-p}$ 上で $\int_{S^{8-p}}F_{8-p}=2\pi N$ が、ブレーンの枚数 $N$ を与える。}
\end{center}

\begin{physbox}[：Dブレーンと RR 場]
II型超弦理論には\term{ラモン・ラモン場}{Ramond--Ramond field}（RR場）$C_{p+1}$ があり、D$p$ ブレーンがその電気的な源になる。IIA型では $p$ が偶数（D0, D2, D4, D6, D8）、IIB型では奇数（D$(-1)$, D1, D3, D5, D7, D9）。$D=10$ では D$p$ の磁気的双対は D$(6-p)$ ブレーンで、電荷は
\[
N=\frac{1}{2\pi}\int_{S^{8-p}}F_{8-p},\qquad F_{8-p}={*F_{p+2}}
\]
で測られる（適当な規格化のもと）。コンパクトな内部空間 $K$ 上のフラックスは $[F/2\pi]\in H^{p+2}(K;\Z)$ の像に値をとり、独立なフラックス量子の数は $b_{p+2}(K)$ である。より精密には、RR電荷はコホモロジーではなく\textbf{K理論}で分類される（第11章で扱う）。
\end{physbox}

\begin{myprop}{タドポール相殺（ガウスの法則）}{ch10-tad}
コンパクトな空間 $K$ 上で $dF=\rho$（$\rho$ は源の電荷密度を表す閉形式）なら、ストークスの定理より $\int_K\rho=\int_KdF=0$：\textbf{総電荷はゼロでなければならない}。弦理論では、ブレーンの電荷に加えて、フラックスや曲率（オリエンティフォールド面、高次微分補正）も電荷を担う。代表的な結果（ここでは認める）：
\begin{itemize}
\item M理論をCY4 $Y$ にコンパクト化すると、M2ブレーンの数 $N_{M2}$ と4形式フラックス $G$ について
$\displaystyle N_{M2}+\frac12\int_Y\frac{G}{2\pi}\wedge\frac{G}{2\pi}=\frac{\chi(Y)}{24}$。
F理論（楕円ファイバー付きCY4）では $N_{M2}$ が D3ブレーンの数 $N_{D3}$ に置き換わる。
\item ヘテロティック弦では $dH=\frac{\alpha'}{4}\bigl(\tr R\wedge R-\tr F\wedge F\bigr)$（$\tr$ は適当に規格化したトレース）。したがって $[\tr R\wedge R]=[\tr F\wedge F]$、すなわち $SU(3)$ 束 $V$ について $c_2(V)=c_2(TX)$ が必要（5ブレーンがないとき）。標準埋め込み $V=TX$ はこれを自動的にみたす。
\end{itemize}
\end{myprop}

\subsection{カルツァ・クライン還元と質量ゼロ場の数}

\begin{myprop}{調和形式による展開}{ch10-kk}
$D=4+n$ 次元の $p$ 形式ゲージ場 $C_p$ を $M_4\times K$（$K$ はコンパクトなリーマン多様体）上で考える。$K$ 上の調和 $k$ 形式の基底 $\omega^{(\alpha)}_k$（$\alpha=1,\dots,b_k(K)$）で
\[
C_p=\sum_{k=0}^{p}\sum_{\alpha=1}^{b_k(K)}c^{(\alpha)}_{p-k}(x)\wedge\omega_k^{(\alpha)}(y)+(\text{質量をもつモード})
\]
と展開すると、4次元の質量ゼロ場は、各 $k$ について $b_k(K)$ 個の $(p-k)$ 形式場 $c^{(\alpha)}_{p-k}$ である。
\end{myprop}
\begin{note}
ゲージ $d^\dagger C=0$ で運動方程式は $\Delta_{M_4\times K}C=(\Delta_{M_4}+\Delta_K)C=0$。$\Delta_K$ の固有値 $m^2>0$ のモードは4次元で質量 $m$ をもつので、質量ゼロは $\Delta_K\omega=0$（調和形式）のモードに限る。その数はホッジの定理（第5章）により $b_k(K)$。4次元では2形式は（質量ゼロなら）スカラーと双対、3形式は伝播自由度をもたない。
\end{note}

\subsection{アノマリー相殺とミラー対称性}

\begin{mythm}{グリーン・シュワルツ機構（主張）}{ch10-gs}
10次元 $\mathcal N=1$ 超重力＋超対称ヤン・ミルズ理論（ゲージ群 $G$）のアノマリー多項式 $I_{12}$（第9章）は、$\tr R^6$（重力の既約部分）を係数 $\dim G-496$ で含み、ゲージ部分は $\operatorname{Tr}F^6$（随伴表現のトレース）の既約部分を含む。これらが消えるとき、
\[
I_{12}\propto\Bigl(\tr R^2-\tfrac{1}{30}\operatorname{Tr}F^2\Bigr)\wedge X_8
\]
と因数分解し、$B$ 場の相互作用 $\int B\wedge X_8$ と修正されたビアンキ恒等式 $dH\propto\tr R^2-\tfrac{1}{30}\operatorname{Tr}F^2$ によってアノマリーが相殺される。異常のない弦理論のゲージ群は $SO(32)$ と $E_8\times E_8$ である。
\end{mythm}

\begin{physbox}[：ミラー対称性]
CY3 の組 $(X,W)$ で $h^{1,1}(X)=h^{2,1}(W)$, $h^{2,1}(X)=h^{1,1}(W)$ をみたし、IIA型弦理論を $X$ にコンパクト化したものと IIB型を $W$ にコンパクト化したものが同じ物理を与えるものを\term{ミラー対}{mirror pair}という。ホッジ・ダイヤモンドは $90^\circ$ 回転し、$\chi(W)=-\chi(X)$。$X$ のケーラー・モジュライ空間（の量子補正を含む幾何、世界面インスタントンの数え上げ）が $W$ の複素構造モジュライ空間の古典幾何で計算できる。
\end{physbox}

\section{問題}

\subsection{世界面とトーラス}

\begin{problem}[難易度 A]\label{ch10-prob-ws}
次の世界面のオイラー標数 $\chi$ と、振幅に掛かる $g_s$ のべき $g_s^{-\chi}$ を求めよ：球面 $S^2$、円板 $D^2$、射影平面 $\RP^2$、トーラス $T^2$、円環（アニュラス）、メビウスの帯、クラインの壺。また、閉じた向き付けられた弦の理論（II型、ヘテロティック）とI型弦理論で、それぞれどれが現れるかを答えよ。さらに、I型弦理論の1ループの真空振幅には $T^2$ 以外にどの世界面が同じ次数で寄与するか。
\begin{center}
\begin{tikzpicture}[scale=0.85]
  \foreach \x/\nm in {0/{$S^2$},2.1/{$D^2$},4.2/{$\RP^2$},6.3/{$T^2$},8.4/{円環},10.5/{メビウス},12.6/{クライン}}{
    \shade[ball color=blue!8] (\x,0) circle (0.75);
    \draw[thick] (\x,0) circle (0.75);
    \node at (\x,-1.1) {\footnotesize \nm};
  }
  % disk: one boundary
  \draw[very thick,blue!70!black,fill=white] (2.1,0) circle (0.25);
  % RP2: one crosscap
  \draw[red!70!black,fill=white] (4.2,0) circle (0.22);
  \draw[red!70!black] (4.05,-0.15)--(4.35,0.15); \draw[red!70!black] (4.05,0.15)--(4.35,-0.15);
  % torus: handle
  \draw[thick] (6.0,0.05) to[bend right=35] (6.6,0.05);
  \draw[thick,fill=white] (6.08,-0.01) to[bend left=35] (6.52,-0.01) to[bend left=20] (6.08,-0.01);
  % annulus: two boundaries
  \draw[very thick,blue!70!black,fill=white] (8.1,0) circle (0.2);
  \draw[very thick,blue!70!black,fill=white] (8.7,0) circle (0.2);
  % Mobius: boundary + crosscap
  \draw[very thick,blue!70!black,fill=white] (10.2,0) circle (0.2);
  \draw[red!70!black,fill=white] (10.8,0) circle (0.2);
  \draw[red!70!black] (10.66,-0.14)--(10.94,0.14); \draw[red!70!black] (10.66,0.14)--(10.94,-0.14);
  % Klein: two crosscaps
  \foreach \c in {12.3,12.9}{
    \draw[red!70!black,fill=white] (\c,0) circle (0.2);
    \draw[red!70!black] (\c-0.14,-0.14)--(\c+0.14,0.14); \draw[red!70!black] (\c-0.14,0.14)--(\c+0.14,-0.14);
  }
\end{tikzpicture}
\zucap{球面に取っ手（$T^2$）、穴＝境界（青の太線）、交差帽（$\times$ 印）を付けて表した世界面。}
\end{center}
\end{problem}

\begin{solution}
定理 \ref{thm:ch10-surf} の $\chi=2-2g-b-c$ を用いる。
\begin{center}
\begin{tabular}{lcccccc}
世界面 & $g$ & $b$ & $c$ & $\chi$ & $g_s^{-\chi}$ & 現れる理論\\ \hline
$S^2$ & 0&0&0& $2$ & $g_s^{-2}$ & 閉弦すべて\\
$D^2$ & 0&1&0& $1$ & $g_s^{-1}$ & I型\\
$\RP^2$ & 0&0&1& $1$ & $g_s^{-1}$ & I型\\
$T^2$ & 1&0&0& $0$ & $g_s^{0}$ & 閉弦すべて\\
円環 & 0&2&0& $0$ & $g_s^{0}$ & I型\\
メビウスの帯 & 0&1&1& $0$ & $g_s^{0}$ & I型\\
クラインの壺 & 0&0&2& $0$ & $g_s^{0}$ & I型
\end{tabular}
\end{center}
胞体分割でも確かめられる：$\RP^2=e^0\cup e^1\cup e^2$ で $\chi=1-1+1=1$。クラインの壺は正方形の辺を同一視したもので頂点1、辺2、面1なので $\chi=0$。円板は可縮で $\chi=1$、円環とメビウスの帯は円周に変位レトラクトするので $\chi=\chi(S^1)=0$。

閉じた向き付けられた弦では、世界面は境界も交差帽ももたない向き付け可能な閉曲面 $\Sigma_g$ に限られるので、上の表では $S^2$（木）と $T^2$（1ループ）だけが現れる。I型弦理論は開弦（境界）を含み、世界面の向きを区別しない（交差帽を許す）ので、すべてが現れる。

1ループ（$\chi=0$, $g_s^0$）の真空振幅には $T^2$ に加えて、円環・メビウスの帯・クラインの壺が寄与する。$2g+b+c=2$ の解は $(g,b,c)=(1,0,0),(0,2,0),(0,1,1),(0,0,2)$ のちょうど4つで、これで尽きている。I型弦理論では、後の3つの和に現れる発散（RR場のタドポール）が相殺することを要求すると、ゲージ群が $SO(32)$ に決まる（グリーン・シュワルツ。円環が $N^2$、メビウスの帯が $N$ に比例し、$N=32$ でクラインの壺と相殺する）。
\end{solution}

\begin{problem}[難易度 A]\label{ch10-prob-torus}
$T^{2n}=\C^n/\Lambda$（$\Lambda$ は $\R^{2n}$ の格子）に平坦な計量を入れる。調和形式が定数係数の形式であること（第5章）を用いて次を示せ。\\
(1) $h^{p,q}(T^{2n})=\binom np\binom nq$。\\
(2) $b_k(T^{2n})=\binom{2n}{k}$ と $\chi(T^{2n})=0$。\\
(3) $T^2$, $T^4$, $T^6$ のホッジ・ダイヤモンドを書け。$T^6$ が（ホロノミー $SU(3)$ の意味の）CY3 でないのはホッジ数のどこに現れているか。
\end{problem}

\begin{solution}
(1) 定数係数の $(p,q)$ 形式の基底は $dz^{i_1}\wedge\dots\wedge dz^{i_p}\wedge d\bar z^{j_1}\wedge\dots\wedge d\bar z^{j_q}$（$i_1<\dots<i_p$, $j_1<\dots<j_q$）で、$i$ の選び方が $\binom np$、$j$ の選び方が $\binom nq$ 通り。よって $h^{p,q}=\binom np\binom nq$。

(2) ホッジ分解とファンデルモンドの恒等式より
\[
b_k=\sum_{p+q=k}\binom np\binom nq=\binom{2n}{k}.
\]
（実際、定数係数の実 $k$ 形式は $dx^1,\dots,dx^{2n}$ から $k$ 個選ぶので直接 $\binom{2n}k$ である。）$\chi=\sum_k(-1)^k\binom{2n}{k}=(1-1)^{2n}=0$。

(3) $n=1,2,3$ で
\[
T^2:\ \begin{matrix}&1&\\1&&1\\&1&\end{matrix}\qquad
T^4:\ \begin{matrix}&&1&&\\&2&&2&\\1&&4&&1\\&2&&2&\\&&1&&\end{matrix}\qquad
T^6:\ \begin{matrix}&&&1&&&\\&&3&&3&&\\&3&&9&&3&\\1&&9&&9&&1\\&3&&9&&3&\\&&3&&3&&\\&&&1&&&\end{matrix}
\]
$T^6$ のベッチ数は $1,6,15,20,15,6,1$（和は $2^6=64$）。$T^6$ も $c_1=0$（$\Omega=dz^1\wedge dz^2\wedge dz^3$ がある）で $h^{3,0}=1$ だが、$h^{1,0}=3$, $h^{2,0}=3$ が $0$ でない。ホロノミー $SU(3)$ なら $h^{1,0}=h^{2,0}=0$ なので、$T^6$ はそうでない。実際 $T^6$ のホロノミーは自明で、平行な1形式 $dz^i$ が存在する。これが問題 \ref{ch10-prob-susy} で見るように超対称性が多く残ること（$\mathcal N=4$）に対応する。
\end{solution}

\subsection{ホッジ数・チャーン類・カラビ・ヤウ多様体}

\begin{problem}[難易度 A]\label{ch10-prob-cpn}
$c(T\CP^n)=(1+H)^{n+1}$、$\int_{\CP^n}H^n=1$ を用いて次を示せ。\\
(1) $c_k(T\CP^n)=\binom{n+1}{k}H^k$、$\chi(\CP^n)=n+1$。これを胞体分割 $\CP^n=e^0\cup e^2\cup\dots\cup e^{2n}$ から得られる値と比べよ。\\
(2) $n=1$ のとき、(1) をガウス・ボンネの定理と比べよ。\\
(3) $\CP^n$ はカラビ・ヤウ多様体でないことを示せ。\\
(4) $\CP^2$ の第1ポントリャーギン類 $p_1=c_1^2-2c_2$ を計算し、ヒルツェブルフの符号数定理 $\sigma=\frac13\int p_1$ が $\sigma(\CP^2)=1$ を与えることを確かめよ。
\end{problem}

\begin{solution}
(1) 2項展開より $c_k=\binom{n+1}{k}H^k$。$\chi=\int c_n=\binom{n+1}{n}\int H^n=n+1$。胞体分割では偶数次元の胞体が $n+1$ 個なので $\chi=\sum_k(-1)^k(\text{$k$ 胞体の数})=n+1$ で一致する。

(2) $n=1$：$c_1(T\CP^1)=2H$、$\int c_1=2=\chi(S^2)$。ガウス・ボンネでは $\frac{1}{2\pi}\int_{S^2}K\,dA=2$ で、$c_1(TS^2)$ がちょうど $\frac{1}{2\pi}K\,dA$ で代表される（第8章）ことと合う。

(3) $c_1(T\CP^n)=(n+1)H$ であり、$\int_{\CP^1}H=1$（$\CP^1\subset\CP^n$ は直線）なので $\int_{\CP^1}c_1=n+1\ne0$。よって $c_1\ne0$ で、カラビ・ヤウではない（$\CP^n$ は正のリッチ曲率をもつ）。

(4) $n=2$：$c_1=3H$, $c_2=3H^2$。$p_1=c_1^2-2c_2=9H^2-6H^2=3H^2$、$\int p_1=3$ で $\sigma=1$。一方 $H^2(\CP^2;\Z)=\Z H$ の交叉形式は $\int H\wedge H=1$ の $(1)$ で、符号数 $1$ に一致する。
\end{solution}

\begin{problem}[難易度 B]\label{ch10-prob-adj}
命題 \ref{prop:ch10-adj} の添加公式を用いる。\\
(1) $\CP^2$ の滑らかな $d$ 次曲線 $C$ の $\chi(C)$ と種数 $g(C)$ を求め、$g=\frac{(d-1)(d-2)}{2}$ を示せ。$c_1(C)=0$ となるのはどの $d$ か。\\
(2) $\CP^3$ の滑らかな4次曲面 $S$ について $c(TS)=1+6H^2$ を示し、$\chi(S)=24$ を導け。レフシェッツの超平面定理を用いて $b_1(S)=0$, $b_2(S)=22$ を示せ。\\
(3) 複素1次元・2次元・3次元のカラビ・ヤウ超曲面を、射影空間の次元と次数で答えよ。
\end{problem}

\begin{solution}
(1) $N=2$：$c(TC)=(1+H)^3(1+dH)^{-1}=(1+3H)(1-dH)=1+(3-d)H$（$C$ 上で $H^2=0$）。よって $c_1(C)=(3-d)H$、$\chi(C)=\int_Cc_1=(3-d)\int_CH=d(3-d)$。$\chi=2-2g$ より
\[
g=1-\frac{d(3-d)}{2}=\frac{d^2-3d+2}{2}=\frac{(d-1)(d-2)}{2}.
\]
$d=1,2$ で $g=0$（$\CP^1$）、$d=3$ で $g=1$、$d=4$ で $g=3$。$c_1=0$ は $d=3$ のみで、これは楕円曲線（$T^2$）である。

(2) $N=3$, $d=4$：$(1+H)^4=1+4H+6H^2+\dots$、$(1+4H)^{-1}=1-4H+16H^2-\dots$。
\[
c(TS)=1+(4-4)H+(6-16+16)H^2=1+6H^2.
\]
$\chi=\int_Sc_2=6\int_SH^2=6\cdot4=24$。レフシェッツの定理（$n=2$）より $H^1(S)\cong H^1(\CP^3)=0$ なので $b_1=0$、ポアンカレ双対で $b_3=0$。$\chi=b_0-b_1+b_2-b_3+b_4=2+b_2=24$ より $b_2=22$。$c_1=0$ なのでこれはK3曲面である（問題 \ref{ch10-prob-k3}）。

(3) $c_1=(N+1-d)H=0$ より $d=N+1$。複素1次元：$\CP^2$ の3次曲線（楕円曲線）。2次元：$\CP^3$ の4次曲面（K3）。3次元：$\CP^4$ の5次超曲面（\term{5次超曲面}{quintic}）。
\end{solution}

\begin{problem}[難易度 B]\label{ch10-prob-k3}
\term{K3曲面}{K3 surface} $S$ は、単連結なコンパクト・ケーラー曲面で $c_1(S)=0$、$h^{2,0}=1$、$\chi(S)=24$ をもつものである。\\
(1) $S$ のホッジ・ダイヤモンドを決定せよ。\\
(2) コンパクト・ケーラー曲面では $b_2^+=2h^{2,0}+1$、$b_2^-=h^{1,1}-1$ であること（ホッジの指数定理）を認めて、交叉形式の符号数 $\sigma(S)=-16$ を示せ。ヒルツェブルフの符号数定理 $\sigma=\frac13\int p_1$ と $p_1=c_1^2-2c_2$ からも同じ値が得られることを確かめよ。\\
(3) ネーターの公式 $\sum_q(-1)^qh^{0,q}=\frac{1}{12}\int(c_1^2+c_2)$ が成り立つことを確かめよ。\\
(4) $H^0\oplus H^2\oplus H^4$ に（$H^0\oplus H^4$ には双曲型の形式 $\bigl(\begin{smallmatrix}0&1\\1&0\end{smallmatrix}\bigr)$ を入れて）内積を入れると、符号 $(4,20)$ の格子になることを示せ。これはヘテロティック弦を $T^4$ にコンパクト化したときのナラン格子 $\Gamma^{4,20}$ と同じ符号であり、IIA/K3 とヘテロティック/$T^4$ の双対性の手がかりとなった。
\end{problem}

\begin{solution}
(1) 単連結なので $b_1=0$、よって $h^{1,0}=h^{0,1}=0$、セール双対性で $h^{2,1}=h^{1,2}=0$。$h^{0,0}=h^{2,2}=1$、$h^{2,0}=h^{0,2}=1$。$\chi=b_0+b_2+b_4=2+b_2=24$ より $b_2=22=h^{2,0}+h^{1,1}+h^{0,2}$、したがって $h^{1,1}=20$。
\[
\begin{matrix}&&1&&\\&0&&0&\\1&&20&&1\\&0&&0&\\&&1&&\end{matrix}
\]

(2) $b_2^+=2\cdot1+1=3$、$b_2^-=20-1=19$、$\sigma=3-19=-16$。一方 $c_1=0$ なので $p_1=-2c_2$、$\int p_1=-2\chi=-48$、$\sigma=-48/3=-16$ で一致する。

(3) 左辺 $=h^{0,0}-h^{0,1}+h^{0,2}=1-0+1=2$。右辺 $=\frac{1}{12}(0+24)=2$。

(4) $H^2$ の交叉形式は $(b_2^+,b_2^-)=(3,19)$。$H^0\oplus H^4$ の双曲型形式は固有値 $\pm1$ で符号 $(1,1)$。和は $(4,20)$、階数は $24=\chi(S)$ に等しい。ヘテロティック弦の $T^4$ コンパクト化では、左右の運動量・巻き付き数がなす偶自己双対格子が、右向き4次元・左向き $4+16=20$ 次元の $\Gamma^{4,20}$ になり、符号が一致する。なお $H^2(S;\Z)$ 自体は偶ユニモジュラー格子 $(-E_8)^{\oplus2}\oplus U^{\oplus3}$（$U$ は双曲型）であることが知られている（第6章の交叉形式）。
\end{solution}

\begin{problem}[難易度 C]\label{ch10-prob-quintic}
(1) 5次超曲面 $X\subset\CP^4$ について
\[
c(TX)=1+10H^2-40H^3
\]
を示し、$\int_XH^3=5$ を用いて $\chi(X)=-200$ を導け。\\
(2) 一般に、$\CP^N$ の滑らかな $d$ 次超曲面 $X$ のオイラー標数が
\[
\chi(X)=\frac{(1-d)^{N+1}-1}{d}+N+1
\]
であることを示せ。\\
(3) (2) を用いて $d=N+1$（カラビ・ヤウ）の場合の $\chi$ を $N=2,3,4,5$ で計算し、問題 \ref{ch10-prob-adj} と (1) に一致することを確かめよ。
\end{problem}

\begin{solution}
(1) $(1+H)^5=1+5H+10H^2+10H^3+\dots$、$(1+5H)^{-1}=1-5H+25H^2-125H^3+\dots$。積の各次数は
\begin{align*}
H^1&:\ 5-5=0,\\
H^2&:\ 10-5\cdot5+25=10,\\
H^3&:\ 10-10\cdot5+5\cdot25-125=-40.
\end{align*}
よって $c(TX)=1+10H^2-40H^3$（$X$ は複素3次元なので $H^4$ 以上は不要）。$c_1=0$ でカラビ・ヤウ。$\chi=\int_Xc_3=-40\cdot5=-200$。

(2) $\chi(X)=\int_Xc_{N-1}(TX)=d\cdot a_{N-1}$ で、$a_{N-1}$ は $(1+H)^{N+1}(1+dH)^{-1}$ の $H^{N-1}$ の係数：
\[
a_{N-1}=\sum_{j=0}^{N-1}\binom{N+1}{j}(-d)^{N-1-j}.
\]
$d\,(-d)^{N-1-j}=\frac1d(-d)^{N+1-j}$ だから
\begin{align*}
\chi&=\frac1d\sum_{j=0}^{N-1}\binom{N+1}{j}(-d)^{N+1-j}
=\frac1d\Bigl[(1-d)^{N+1}-\binom{N+1}{N}(-d)-\binom{N+1}{N+1}\Bigr]\\
&=\frac{(1-d)^{N+1}-1}{d}+N+1 .
\end{align*}
（2項定理 $(1-d)^{N+1}=\sum_{j=0}^{N+1}\binom{N+1}{j}(-d)^{N+1-j}$ から $j=N,N+1$ の項を除いた。）

(3) $d=N+1$ で $\chi=\dfrac{(-N)^{N+1}-1}{N+1}+N+1$。
\begin{center}
\begin{tabular}{cccl}
$N$ & $d$ & $\chi$ & \\ \hline
$2$ & $3$ & $\frac{-8-1}{3}+3=0$ & 楕円曲線\\
$3$ & $4$ & $\frac{81-1}{4}+4=24$ & K3\\
$4$ & $5$ & $\frac{-1024-1}{5}+5=-200$ & 5次超曲面\\
$5$ & $6$ & $\frac{15625-1}{6}+6=2610$ & 6次超曲面（CY4）
\end{tabular}
\end{center}
$N=2,3,4$ は問題 \ref{ch10-prob-adj} と (1) に一致する。$N=5$ の値は問題 \ref{ch10-prob-dirac} で使う。
\end{solution}

\begin{problem}[難易度 C]\label{ch10-prob-h21}
5次超曲面 $X\subset\CP^4$ のホッジ数を3通りの方法で求める。$X$ はホロノミー $SU(3)$ のCY3であり、$\chi(X)=-200$ とする。\\
(1) レフシェッツの超平面定理から $h^{1,1}$ を求め、$\chi=2(h^{1,1}-h^{2,1})$ から $h^{2,1}$ と $b_3$ を求めよ。\\
(2) 複素構造の変形はすべて定義多項式 $P$ の係数を変えることで得られる（認めてよい）。5変数の5次斉次多項式の空間の次元から、$GL(5,\C)$ による座標変換で移り合う分を引いて $h^{2,1}=101$ を確かめよ。\\
(3) フェルマー型 $P_0=\sum_{i=1}^5z_i^5$ の近くでは、無限小変形 $P_0+\epsilon Q$ のうち $GL(5,\C)$ の無限小変換で得られるのは $Q\in\operatorname{span}\{z_j\,\partial P_0/\partial z_i\}$ である。各指数が $3$ 以下の5次単項式の数を数えることで、再び $101$ を得よ。
\end{problem}

\begin{solution}
(1) レフシェッツの定理（$n=3$）より $H^1(X)\cong H^1(\CP^4)=0$、$H^2(X;\Z)\cong H^2(\CP^4;\Z)=\Z$。よって $b_1=0$, $b_2=1$。CY3なので $h^{2,0}=0$ で $h^{1,1}=b_2=1$。
\[
-200=2(1-h^{2,1})\quad\Longrightarrow\quad h^{2,1}=101,\qquad b_3=2+2h^{2,1}=204.
\]
（$\chi=b_0-b_1+b_2-b_3+b_4-b_5+b_6=1+1-b_3+1+1=4-b_3$ からも $b_3=204$。）

(2) 5変数の5次単項式の数は、重複組合せで $\binom{5+5-1}{5}=\binom94=126$。$GL(5,\C)$ は $25$ 次元で、一般の5次式の固定部分群は有限群なので、軌道は $25$ 次元。全体のスカラー倍 $P\to\lambda P$ は同じ超曲面を与えるが、これは $z\to\lambda^{1/5}z$ として $GL(5,\C)$ に含まれている。よって超曲面の変形の次元は
\[
126-25=101=h^{2,1}.
\]

(3) $z_j\partial_iP_0=5z_i^4z_j$（$i,j=1,\dots,5$）は25個の相異なる単項式 $z_i^5$（5個）と $z_i^4z_j$, $j\ne i$（20個）で、ちょうど「ある指数が $4$ 以上の5次単項式」の全体である。したがって真の変形は、これらを除いた単項式、すなわち各指数が $3$ 以下の5次単項式で張られ、その数は $126-25=101$。例えば $z_1z_2z_3z_4z_5$（よく知られた1パラメータ族 $P_0-5\psi z_1z_2z_3z_4z_5$ を与える）、$z_1^3z_2z_3$、$z_1^2z_2^2z_3$ などである。（内訳：指数の型 $(1,1,1,1,1)$ が $1$ 個、$(2,1,1,1,0)$ が $5\cdot4=20$ 個、$(2,2,1,0,0)$ が $\binom52\cdot3=30$ 個、$(3,1,1,0,0)$ が $5\cdot\binom42=30$ 個、$(3,2,0,0,0)$ が $5\cdot4=20$ 個で、合計 $101$。）

3つの方法が一致し、ホッジ・ダイヤモンドは本文の図の右のとおりである。
\end{solution}

\subsection{ヘテロティック弦：超対称性と世代数}

\begin{problem}[難易度 B]\label{ch10-prob-susy}
ヘテロティック弦の10次元の超対称性は、マヨラナ・ワイル・スピノル（$\mathbf{16}$、実16成分）の超電荷で与えられる。$M_4\times K$（$K$ は6次元リーマン多様体）で、$\mathrm{Spin}(9,1)\supset\mathrm{Spin}(3,1)\times\mathrm{Spin}(6)$、$\mathrm{Spin}(6)\cong SU(4)$ のもとで
\[
\mathbf{16}=(\mathbf 2,\mathbf 4)\oplus(\bar{\mathbf 2},\bar{\mathbf 4})
\]
と分解する。4次元に残る超電荷は $K$ 上の共変定数スピノルに対応し、その数は $\mathbf 4$ の中でホロノミー群 $\mathrm{Hol}(K)\subset SU(4)$ の自明表現の個数 $n_0$ で決まる（4次元の $\mathcal N=n_0$、超電荷は $4n_0$ 個）。\\
(1) $K=T^6$, $K=\mathrm{K3}\times T^2$, $K=$ CY3（ホロノミー $SU(3)$）, ホロノミーが $SO(6)$ 全体の一般の $K$ について、$\mathcal N$ を求めよ。\\
(2) II型弦理論（超電荷32個）では、それぞれ何個の4次元超対称性が残るか。\\
(3) 4次元 $\mathcal N=1$ が、カイラルなフェルミオン（標準模型）を許す唯一の可能性である理由を簡単に述べよ。
\end{problem}

\begin{solution}
(1) \textbf{$T^6$}：ホロノミーは自明なので $\mathbf 4$ の4成分すべてが不変、$n_0=4$、$\mathcal N=4$（超電荷16個すべてが残る）。

\textbf{K3$\times T^2$}：$\mathrm{Spin}(6)\supset\mathrm{Spin}(4)\times\mathrm{Spin}(2)=SU(2)_L\times SU(2)_R\times U(1)$ のもとで $\mathbf 4=(\mathbf 2,\mathbf 1)_{+1/2}\oplus(\mathbf 1,\mathbf 2)_{-1/2}$。K3 のホロノミーは $SU(2)\subset SO(4)$ で、$SU(2)_L\times SU(2)_R$ の一方（例えば $SU(2)_L$）である。したがって $\mathbf 4\to\mathbf 2\oplus\mathbf 1\oplus\mathbf 1$ で $n_0=2$、$\mathcal N=2$。

\textbf{CY3}：$SU(3)\subset SU(4)$ で $\mathbf 4=\mathbf 3\oplus\mathbf 1$、$n_0=1$、$\mathcal N=1$。

\textbf{一般の $K$}：ホロノミー $SO(6)$ はスピン群として $SU(4)$ 全体で、$\mathbf 4$ は既約。$n_0=0$、超対称性は残らない。

（逆に、$\mathcal N\ge1$ を要求すると $K$ は共変定数スピノルをもち、ホロノミーが $SU(3)$ に含まれる。これからケーラーかつリッチ平坦、すなわちカラビ・ヤウであることが従う。これが弦理論でカラビ・ヤウ多様体が現れる理由である。）

(2) II型は2つの $\mathbf{16}$（IIA は逆カイラリティ、IIB は同カイラリティ）をもつので、それぞれ2倍：$T^6$ で $\mathcal N=8$、K3$\times T^2$ で $\mathcal N=4$、CY3 で $\mathcal N=2$。

(3) $\mathcal N\ge2$ の超対称多重項では、ゲージ場とフェルミオンが同じ表現に属するか、ハイパー多重項のフェルミオンがベクトル的（左右対称な表現）になり、ゲージ群の複素表現にカイラルなフェルミオンを置けない。$\mathcal N=1$ のカイラル多重項はワイル・フェルミオンを任意の表現にもてるので、カイラルな標準模型（およびその超対称拡張）を実現できる。ヘテロティック弦では CY3 がちょうどこれを与える。
\end{solution}

\begin{problem}[難易度 B]\label{ch10-prob-gen}
ヘテロティック弦を CY3 $X$ に標準埋め込みでコンパクト化する（定理 \ref{thm:ch10-gen}）。\\
(1) 5次超曲面では $\mathbf{27}$ と $\overline{\mathbf{27}}$ がそれぞれいくつ現れるか。正味の世代数を求めよ。\\
(2) 階数3の複素ベクトル束 $V$ で $c_1(V)=0$ のもの、および $c_1(X)=0$ の複素3次元多様体について、$\Td(X)=1+\frac{1}{12}c_2(X)+\cdots$（6形式部分は $\frac{1}{24}c_1c_2=0$）と $\ch(V)=3-c_2(V)+\frac12c_3(V)$ を用いて
\[
\int_X\ch(V)\Td(X)=\frac12\int_Xc_3(V)
\]
を示せ。$V=TX$ ならこれは $\chi(X)/2$ である。また $V$ を自明な直線束にした $\int_X\Td(X)$ がドルボー・コホモロジーの交代和 $\sum_q(-1)^qh^{0,q}$ と一致することを確かめよ。\\
(3) 有限群 $G$ が $X$ に自由に作用するとき $\chi(X/G)=\chi(X)/\lvert G\rvert$ を示せ。5次超曲面に自由に作用する $\Z_5\times\Z_5$ による商、および $\chi=-18$ の田・ヤウ多様体に自由に作用する $\Z_3$ による商の世代数を求めよ。3世代を得るには $\chi$ はいくつでなければならないか。
\end{problem}

\begin{solution}
(1) $(h^{1,1},h^{2,1})=(1,101)$ なので $\mathbf{27}$ が $101$ 個、$\overline{\mathbf{27}}$ が $1$ 個。$\mathbf{27}$ と $\overline{\mathbf{27}}$ の1対は質量項 $\mathbf{27}\cdot\overline{\mathbf{27}}$ を組めるので、低エネルギーに残る正味の世代数は $101-1=100=\lvert-200\rvert/2$。

(2) $\ch(V)\Td(X)$ の6形式部分は
\[
\ch_3(V)+\ch_2(V)\Td_1+\ch_1(V)\Td_2+\ch_0(V)\Td_3=\tfrac12c_3(V)+(-c_2(V))\cdot0+0\cdot\tfrac{c_2}{12}+3\cdot0
\]
（$\Td_1=\frac12c_1(X)=0$, $\ch_1(V)=c_1(V)=0$, $\Td_3=\frac1{24}c_1c_2=0$）。よって $\int\ch(V)\Td(X)=\frac12\int c_3(V)$。$\ch_3=\frac16(c_1^3-3c_1c_2+3c_3)=\frac12c_3$ を用いた。$V=TX$ なら $\int c_3(TX)=\chi(X)$ なので指数は $\chi/2$。5次超曲面で $-100$ で、絶対値が (1) の世代数に一致する（符号はどちらを「世代」と呼ぶかの規約による）。

自明な直線束では $\int_X\Td(X)=\int\Td_3=0$。一方 CY3 では $h^{0,0}=1$, $h^{0,1}=0$, $h^{0,2}=h^{2,0}=0$, $h^{0,3}=h^{3,0}=1$ なので $\sum_q(-1)^qh^{0,q}=1-0+0-1=0$ で、ヒルツェブルフ・リーマン・ロッホの定理と整合する。

(3) $X\to X/G$ は $\lvert G\rvert$ 重の被覆。$X/G$ の三角形分割を（十分細かくとって各単体が被覆で自明になるようにし）$X$ に持ち上げると、各単体がちょうど $\lvert G\rvert$ 個の単体に持ち上がる。よって各次元の単体数が $\lvert G\rvert$ 倍で、$\chi(X)=\lvert G\rvert\chi(X/G)$。

5次超曲面の $\Z_5\times\Z_5$ 商：$\chi=-200/25=-8$、世代数 $4$。田・ヤウ多様体の $\Z_3$ 商：$\chi=-18/3=-6$、世代数 $3$。3世代には $\chi=\pm6$ が必要である。

なお商空間は $\pi_1(X/G)=G$ をもつので、$E_6$ を標準模型に近い群に破るウィルソン線（平坦接続のホロノミー、第2, 8章）を入れることができる。これもトポロジーが現象論に効く例である。
\end{solution}

\subsection{オービフォールド}

\begin{problem}[難易度 B]\label{ch10-prob-pillow}
$T^2=\C/(\Z+\tau\Z)$（$\operatorname{Im}\tau>0$）に $\Z_2$ が $z\mapsto-z$ で作用する。\\
(1) 固定点が $0,\ \frac12,\ \frac\tau2,\ \frac{1+\tau}{2}$ の4点であることを示せ。\\
(2) 定理 \ref{thm:ch10-orbchi}(1) で $\chi(T^2/\Z_2)$ を求め、$T^2/\Z_2\cong S^2$ であることを説明せよ（「枕」の図を参照）。\\
(3) 平坦な計量は商空間に下りて、固定点の像以外では平坦、固定点の像では錐角 $\pi$ の円錐特異点になることを示せ。円錐の頂点の曲率を欠損角 $2\pi-(\text{錐角})$ とみなして、ガウス・ボンネの定理 $\sum(\text{欠損角})=2\pi\chi$ が成り立つことを確かめよ。\\
(4) 複素多様体としては $T^2/\Z_2$ は滑らかで、$\CP^1$ と同型であることを、局所座標 $w=z^2$ を用いて説明せよ。
\end{problem}

\begin{solution}
(1) $-z\equiv z$ $\Leftrightarrow$ $2z\in\Z+\tau\Z$ $\Leftrightarrow$ $z\in\frac12(\Z+\tau\Z)$。$\frac12(\Z+\tau\Z)/(\Z+\tau\Z)\cong(\Z_2)^2$ なので $T^2$ 上では $0,\frac12,\frac\tau2,\frac{1+\tau}{2}$ の4点。

(2) $\chi(T^2)=0$、$\chi(X^g)=4$（4点）：
\[
\chi(T^2/\Z_2)=\frac12(0+4)=2 .
\]
基本領域は $\{s+t\tau:0\le s\le1,\ 0\le t\le\frac12\}$ で、左右の辺は平行移動 $z\to z+1$ で同一視され（円筒）、下の辺 $t=0$ は $s\mapsto-s\equiv1-s$ で中点 $\frac12$ を折り目として半分に折られ、上の辺 $t=\frac12$ も $s+\frac\tau2\mapsto-s-\frac\tau2\equiv(1-s)+\frac\tau2$ で同様に折られる。円筒の両端の円周をそれぞれ半分に折って閉じると、2つの円板でふさいだのと同じで球面になる。角が4つの固定点の像で、形は枕である。閉じた向き付け可能曲面（$z\mapsto-z$ は向きを保つ）で $\chi=2$ だから、分類定理からも $S^2$。

(3) $z\mapsto-z$ は等長変換なので、計量は商に下りる。固定点以外では作用は自由で、局所的に等長なので平坦。固定点 $z_0$ の近傍では、極座標 $z-z_0=re^{i\theta}$ で作用は $\theta\mapsto\theta+\pi$ なので、商は角度 $\theta\in[0,\pi]$ の扇形の両辺を貼ったもの、すなわち錐角 $\pi$ の円錐である。欠損角は $2\pi-\pi=\pi$ で、4点で計 $4\pi=2\pi\cdot2=2\pi\chi(S^2)$。ガウス・ボンネが成り立つ（平坦な部分は寄与しない）。

(4) 固定点の近くで $w=(z-z_0)^2$ は $\Z_2$ 不変な正則関数で、商の近傍を $w$ 平面の近傍と正則に同一視する（$\C/\Z_2\cong\C$, $z\mapsto z^2$）。よって $T^2/\Z_2$ は特異点のないコンパクトなリーマン面で、種数 $0$ なので $\CP^1$。具体的には、ワイエルシュトラスの $\wp$ 関数（偶関数）が同型 $T^2/\Z_2\to\CP^1$ を与える。計量は特異でも、複素構造は滑らかであることに注意。複素1次元では固定点は複素構造の特異点にならないが、複素2次元以上の $\C^n/\Z_k$ は本当の特異点になる（問題 \ref{ch10-prob-t4}）。
\end{solution}

\begin{problem}[難易度 C]\label{ch10-prob-t4}
$T^4=\C^2/\Lambda$（$\Lambda\cong\Z^4$）に $\Z_2$ が $(z_1,z_2)\mapsto(-z_1,-z_2)$ で作用する。\\
(1) 固定点の数を求めよ。各固定点の近傍は $\C^2/\Z_2$ で、その「境界」は $S^3/\Z_2=\RP^3$ である。\\
(2) 商空間 $T^4/\Z_2$ のオイラー標数と、$\Z_2$ 不変な調和形式から得られるベッチ数を求めよ。\\
(3) 各特異点を江口・ハンソン空間で解消した多様体 $\widetilde X$ の $\chi$ と $b_2$ を求めよ。また $h^{1,1}(\widetilde X)$, $h^{2,0}(\widetilde X)$ を求め、$\widetilde X$ がK3曲面のホッジ数をもつことを確かめよ。\\
(4) 弦理論のオービフォールド公式 $\chi_{\rm str}$ を計算し、(3) と比べよ。
\end{problem}

\begin{solution}
(1) $2z\in\Lambda$、すなわち $z\in\frac12\Lambda/\Lambda\cong(\Z_2)^4$ で、固定点は $2^4=16$ 個。

(2) 定理 \ref{thm:ch10-orbchi}(1)：$\chi(T^4/\Z_2)=\frac12(0+16)=8$。

$T^4$ の調和形式（定数係数）のうち不変なもの：$dx^i\mapsto-dx^i$ なので、$k$ 形式は $(-1)^k$ 倍される。不変なのは偶数次だけで、$b_0=1$, $b_1=0$, $b_2=\binom42=6$, $b_3=0$, $b_4=1$。商空間の（実係数の）コホモロジーは不変部分で与えられ、$\chi=1+6+1=8$ と一致する。

(3) 各固定点の像の近傍（$\C^2/\Z_2$ の球近傍、錐 $C(\RP^3)$、可縮で $\chi=1$）を江口・ハンソン空間の対応する部分（$S^2$ にホモトピー同値、$\chi=2$）に取り替える。境界 $\RP^3$ は $\chi=0$ なので、1つあたり $\chi$ が $2-1=+1$ 変わる：
\[
\chi(\widetilde X)=8+16\cdot1=24 .
\]
各例外曲線 $\CP^1$ が新しい2次ホモロジー類を与えるので $b_2=6+16=22$（$b_1=b_3=0$ は変わらず、$\chi=1+22+1=24$ と整合）。

ホッジ数：$T^4$ 上の不変な2形式は $dz^1\wedge dz^2$（$(2,0)$）、$d\bar z^1\wedge d\bar z^2$（$(0,2)$）、$dz^i\wedge d\bar z^j$（$(1,1)$, 4個）。例外曲線は複素曲線なので、そのポアンカレ双対は $(1,1)$ 類。よって
\[
h^{2,0}=1,\qquad h^{1,1}=4+16=20,\qquad h^{1,0}=0 .
\]
これは問題 \ref{ch10-prob-k3} のK3のホッジ・ダイヤモンドに一致する（実際 $\widetilde X$ はK3曲面で、\term{クンマー曲面}{Kummer surface}と呼ばれる）。正則2形式 $dz^1\wedge dz^2$ が不変なので $\Z_2$ は $SU(2)$ に含まれ、解消後も $c_1=0$ が保たれる。

(4) $G=\Z_2=\{1,g\}$ は可換で、可換な組 $(h_1,h_2)$ は4つ。$X^1\cap X^1=T^4$（$\chi=0$）、他の3組 $(1,g),(g,1),(g,g)$ ではいずれも16点（$\chi=16$）：
\[
\chi_{\rm str}=\frac12\bigl(0+3\cdot16\bigr)=24 .
\]
(3) の解消多様体のオイラー標数に一致する。素朴な商空間の値 $8$ との差 $16$ は、ねじれたセクター（固定点に局在した閉弦の状態）からの寄与で、幾何学的には特異点の解消で現れる16個の2サイクルに対応する。
\end{solution}

\begin{problem}[難易度 D]\label{ch10-prob-t6z3}
$\omega=e^{2\pi i/3}$、$\Lambda_0=\Z+\Z\omega$（六角格子）とし、$T^2_0=\C/\Lambda_0$ に $\Z_3$ が $z\mapsto\omega z$ で作用する。\\
(1) $T^2_0$ 上の固定点が、格子座標 $z=a+b\omega$ で $(a,b)=(0,0),(\frac23,\frac13),(\frac13,\frac23)$ の3点であることを示せ（図）。\\
(2) $T^6=(T^2_0)^3$ に $\Z_3$ が $(z_1,z_2,z_3)\mapsto(\omega z_1,\omega z_2,\omega z_3)$ で作用する。固定点の数を求め、この作用が正則3形式 $\Omega=dz^1\wedge dz^2\wedge dz^3$ を保つことを確かめよ。\\
(3) $\Z_3$ 不変な $(p,q)$ 形式を数え、ねじれていないセクターのホッジ数 $h^{1,1}_{\rm untw}$, $h^{2,1}_{\rm untw}$ を求めよ。\\
(4) 各固定点を $\mathcal O(-3)\to\CP^2$ で解消すると $h^{1,1}$ と $h^{2,2}$ が1ずつ増える（例 \ref{ex:ch10-resol}）。解消した CY3 $\widetilde X$ の $(h^{1,1},h^{2,1})$ と $\chi$ を求め、$\chi_{\rm str}$ および「素朴な $\chi$＋例外因子の寄与」と一致することを確かめよ。\\
(5) ヘテロティック弦を $\widetilde X$ にコンパクト化したときの世代数を求めよ。$h^{2,1}=0$（\term{剛的}{rigid}）であることから、ミラー多様体が通常の CY3 として存在しえないことを説明せよ。
\begin{center}
\begin{tikzpicture}[scale=2.0]
  \coordinate (w) at (-0.5,0.8660);
  \fill[blue!6] (0,0) -- (1,0) -- (0.5,0.866) -- (-0.5,0.866) -- cycle;
  \foreach \xx/\yy in {-1/0,0/0,1/0,2/0,-1.5/0.866,-0.5/0.866,0.5/0.866,1.5/0.866,
      -0.5/-0.866,0.5/-0.866,1.5/-0.866}{
    \node[circle,fill=gray,inner sep=0.8pt] at (\xx,\yy) {};
  }
  \draw[gray!60] (-1,0)--(2,0) (-1.5,0.866)--(1.5,0.866) (-0.5,-0.866)--(1.5,-0.866);
  \draw[gray!60] (-1,0)--(-0.5,0.866) (-0.5,-0.866)--(0.5,0.866) (0.5,-0.866)--(1.5,0.866) (1.5,-0.866)--(2,0);
  \draw[gray!60] (-1.5,0.866)--(-0.5,-0.866) (-0.5,0.866)--(0.5,-0.866) (0.5,0.866)--(1.5,-0.866) (1.5,0.866)--(2,0);
  \draw[thick] (0,0) -- (1,0) -- (0.5,0.866) -- (-0.5,0.866) -- cycle;
  \draw[gray,densely dashed] (0,0) -- (0.5,0.866);
  \node[vtx,red!70!black] at (0,0) {};
  \node[red!70!black,below left] at (0,0) {\footnotesize $0$};
  \node[vtx,red!70!black] at (0.5,0.2887) {};
  \node[red!70!black,below] at (0.5,0.24) {\footnotesize $\tfrac{2+\omega}{3}$};
  \node[vtx,red!70!black] at (0,0.5774) {};
  \node[red!70!black,left] at (-0.04,0.5774) {\footnotesize $\tfrac{1+2\omega}{3}$};
  \node[below] at (1,0) {\footnotesize $1$};
  \node[above left] at (-0.5,0.866) {\footnotesize $\omega$};
  \node[above right] at (0.5,0.866) {\footnotesize $1+\omega$};
\end{tikzpicture}
\zucap{六角格子 $\Lambda_0=\Z+\Z\omega$ と基本領域（影）。$z\mapsto\omega z$（$120^\circ$ 回転）の固定点は、原点と2つの正三角形の重心の3点（赤）。}
\end{center}
\end{problem}

\begin{solution}
(1) $\omega z\equiv z$ $\Leftrightarrow$ $(1-\omega)z\in\Lambda_0$。$\lvert1-\omega\rvert^2=3$ なので $(1-\omega)\Lambda_0$ は指数 $3$ の部分格子で、解 $z\in(1-\omega)^{-1}\Lambda_0$ は $\Lambda_0$ を法として3個。直接確かめると
\[
\omega\cdot\frac{2+\omega}{3}=\frac{2\omega+\omega^2}{3}=\frac{2\omega-1-\omega}{3}=\frac{\omega-1}{3}=\frac{2+\omega}{3}-1,
\]
\[
\omega\cdot\frac{1+2\omega}{3}=\frac{\omega+2\omega^2}{3}=\frac{\omega-2-2\omega}{3}=\frac{1+2\omega}{3}-(1+\omega)
\]
（$\omega^2=-1-\omega$ を用いた）で、いずれも格子ベクトルの差しかない。$\frac{2+\omega}{3}=\frac23+\frac13\omega$, $\frac{1+2\omega}{3}=\frac13+\frac23\omega$ で、それぞれ三角形 $(0,1,1+\omega)$ と $(0,1+\omega,\omega)$ の重心である。

(2) 固定点は各因子の固定点の組なので $3^3=27$ 個。$\Omega\mapsto\omega^3\Omega=\Omega$ で不変（作用は $SU(3)$ に含まれる：$\operatorname{diag}(\omega,\omega,\omega)$, $\det=1$）。

(3) $dz^i\mapsto\omega\,dz^i$, $d\bar z^j\mapsto\bar\omega\,d\bar z^j=\omega^2d\bar z^j$ なので、$(p,q)$ 形式は $\omega^{p+2q}=\omega^{p-q}$ 倍される。不変 $\Leftrightarrow$ $p\equiv q\pmod3$。
\begin{itemize}
\item $(1,1)$：$p-q=0$ で不変。$h^{1,1}_{\rm untw}=3\cdot3=9$。
\item $(2,1)$：$p-q=1$ で不変でない。$h^{2,1}_{\rm untw}=0$。
\item $(1,0)$, $(2,0)$：不変でない（$h^{1,0}=h^{2,0}=0$）。$(3,0)$：$p-q=3\equiv0$ で不変、$h^{3,0}=1$。
\end{itemize}
よって不変部分のホッジ・ダイヤモンドはCY3の形で、$(h^{1,1},h^{2,1})_{\rm untw}=(9,0)$。

(4) 27個の固定点の解消でそれぞれ $h^{1,1}$ が1増え、$h^{2,1}$ は変わらない：
\[
h^{1,1}(\widetilde X)=9+27=36,\qquad h^{2,1}(\widetilde X)=0,\qquad \chi(\widetilde X)=2(36-0)=72 .
\]
\textbf{弦理論の公式}：$\Z_3$ の可換な組は $3\times3=9$ 個。$(1,1)$ は $\chi(T^6)=0$。それ以外の8組では、$g\ne1$ の固定点集合は $g=\omega$ でも $g=\omega^2$ でも同じ27点で、$X^g\cap X^h$ は27点（$g$ または $h$ が $1$ ならもう一方の固定点集合）：
\[
\chi_{\rm str}=\frac13(0+8\cdot27)=72 .
\]
\textbf{素朴な $\chi$＋解消}：$\chi(T^6/\Z_3)=\frac13(0+27+27)=18$。各点を $\CP^2$ にホモトピー同値な空間に置き換えると $\chi$ は $\chi(\CP^2)-1=2$ 増えるので $18+27\cdot2=72$。3つとも一致する。

(5) 世代数 $\lvert\chi\rvert/2=36$（$\overline{\mathbf{27}}$ が $h^{1,1}=36$ 個、$\mathbf{27}$ が $h^{2,1}=0$ 個）。ミラー $W$ は $h^{1,1}(W)=h^{2,1}(\widetilde X)=0$ をもたねばならないが、コンパクト・ケーラー多様体ではケーラー形式の類が $[J]\ne0\in H^{1,1}$ なので $h^{1,1}\ge1$。よって通常の幾何学的な CY3 としてのミラーは存在しない。（ミラーはランダウ・ギンツブルグ模型などの非幾何学的な $\mathcal N=2$ 共形場理論として存在する。）
\end{solution}

\subsection{カルツァ・クライン還元とミラー対称性}

\begin{problem}[難易度 B]\label{ch10-prob-kk}
命題 \ref{prop:ch10-kk} を用いる。\\
(1) 10次元の2形式場 $B_{MN}$ を $M_4\times T^6$ に還元したときの4次元の質量ゼロ場（2形式、ベクトル、スカラー）の数を求めよ。10次元の質量ゼロの2形式の物理的自由度（光円錐ゲージで横方向 $SO(8)$ の反対称テンソル）と、4次元の自由度の和が一致することを確かめよ（4次元の質量ゼロの2形式は1自由度、ベクトルは2自由度）。\\
(2) 10次元の計量 $g_{MN}$ についても同様に、4次元の計量・ベクトル・スカラーの数を数え、自由度（10次元の重力子は $SO(8)$ の対称トレースレス・テンソル）を確かめよ。ベクトル $g_{\mu i}$ は $T^6$ の平行なベクトル場（$b_1=6$ 個）に対応する。\\
(3) $B$ を CY3 $X$ に還元すると、4次元で何個のスカラー（\term{アクシオン}{axion}）とベクトルが得られるか。5次超曲面ではいくつか。CY3 のリッチ平坦計量の変形から得られる実スカラーの数も答えよ。
\end{problem}

\begin{solution}
(1) $B=\frac12B_{\mu\nu}dx^\mu dx^\nu+B_{\mu i}dx^\mu dy^i+\frac12B_{ij}dy^idy^j$ と分け、$b_0=1$, $b_1=6$, $b_2=15$ を用いる：
\begin{center}
\begin{tabular}{lll}
成分 & 4次元の場 & 数\\ \hline
$B_{\mu\nu}$ & 2形式（スカラーと双対） & $b_0=1$\\
$B_{\mu i}$ & ベクトル & $b_1=6$\\
$B_{ij}$ & スカラー & $b_2=15$
\end{tabular}
\end{center}
自由度：10次元では $\binom82=28$。4次元では $1\cdot1+6\cdot2+15\cdot1=28$ で一致する。

(2) $g_{\mu\nu}$：4次元の計量1個（2自由度）。$g_{\mu i}$：ベクトル $6$ 個（12自由度）。$g_{ij}$：$T^6$ の平坦計量の変形、対称行列の成分で $\frac{6\cdot7}{2}=21$ 個のスカラー（21自由度）。合計 $2+12+21=35$。10次元の重力子は $8\times8$ 対称トレースレスで $\frac{8\cdot9}{2}-1=35$。一致する。

(3) CY3 では $b_1=0$, $b_2=h^{1,1}$（$h^{2,0}=0$）なので、$B$ からはベクトルは $0$ 個、スカラー（アクシオン）は $h^{1,1}$ 個、それに4次元の $B_{\mu\nu}$（もう1個のアクシオンと双対）が1個。5次超曲面ではアクシオン $1$ 個と $B_{\mu\nu}$ が $1$ 個。

計量の変形は命題 \ref{prop:ch10-cy3} より実 $h^{1,1}+2h^{2,1}$ 個：5次超曲面で $1+2\cdot101=203$。（$T^6$ の場合のように $b_1\neq0$ だと $g_{\mu i}$ からベクトルが得られるが、CY3 では平行なベクトル場がないので得られない。）
\end{solution}

\begin{problem}[難易度 B]\label{ch10-prob-mirror}
IIA型弦理論を CY3 $X$ にコンパクト化すると4次元 $\mathcal N=2$ 超重力で、ベクトル多重項が $h^{1,1}(X)$ 個、ハイパー多重項が $h^{2,1}(X)+1$ 個現れる。IIB型では、ベクトル多重項が $h^{2,1}(X)$ 個、ハイパー多重項が $h^{1,1}(X)+1$ 個である（認めてよい）。\\
(1) 5次超曲面 $X_5$ について、IIA型とIIB型の多重項の数を求めよ。\\
(2) ミラー5次超曲面 $W$ のホッジ・ダイヤモンドを書き、$\chi(W)$ を求めよ。IIA/$X_5$ と IIB/$W$ の多重項の数が一致することを確かめよ。\\
(3) 一般に、ミラー対 $(X,W)$ について $\chi(W)=-\chi(X)$ と $b_3(W)=2b_2(X)+2$ を示せ。\\
(4) IIB型では、10次元の自己双対4形式 $C_4$ から $b_3(X)$ 個の4次元ベクトルの「半分」、すなわち $\frac12b_3=h^{2,1}+1$ 個の独立なベクトル場が得られる。これが（ベクトル多重項の数）＋（重力多重項のグラビフォトン1個）と一致することを確かめよ。
\end{problem}

\begin{solution}
(1) $(h^{1,1},h^{2,1})=(1,101)$。IIA/$X_5$：ベクトル多重項 $1$、ハイパー多重項 $102$。IIB/$X_5$：ベクトル多重項 $101$、ハイパー多重項 $2$。

(2) $W$ は $(h^{1,1},h^{2,1})=(101,1)$：
\[
\begin{matrix}
&&&1&&&\\
&&0&&0&&\\
&0&&101&&0&\\
1&&1&&1&&1\\
&0&&101&&0&\\
&&0&&0&&\\
&&&1&&&
\end{matrix}
\]
$\chi(W)=2(101-1)=200$。IIB/$W$：ベクトル多重項 $h^{2,1}(W)=1$、ハイパー多重項 $h^{1,1}(W)+1=102$。IIA/$X_5$ と一致する。

(3) $\chi(W)=2(h^{1,1}(W)-h^{2,1}(W))=2(h^{2,1}(X)-h^{1,1}(X))=-\chi(X)$。$b_3(W)=2+2h^{2,1}(W)=2+2h^{1,1}(X)=2+2b_2(X)$。すなわちミラー対称性は $X$ の偶数次コホモロジー $H^0\oplus H^2\oplus H^4\oplus H^6$（次元 $2+2h^{1,1}$）と $W$ の $H^3$ を入れ替える。偶数次元のDブレーン（IIA）が $X$ の偶数次サイクルに、奇数次元のDブレーン（IIB）が $W$ の3サイクルに巻き付くことと対応している。

(4) $C_4=\sum_{\alpha=1}^{b_3}A^{(\alpha)}\wedge\omega_3^{(\alpha)}+\cdots$ から $b_3$ 個の4次元1形式が得られるが、自己双対性 $F_5={*F_5}$ により、4次元の電気的・磁気的な場が対になって同一視され、独立なのは半分の $\frac12b_3=h^{2,1}+1$ 個。ベクトル多重項 $h^{2,1}$ 個のベクトルと、重力多重項のグラビフォトン1個の和に等しい。5次超曲面では $102$。
\end{solution}

\subsection{ブレーン電荷とフラックスの量子化}

\begin{problem}[難易度 A]\label{ch10-prob-dbrane}
(1) 10次元時空で D$p$ ブレーンを囲む球面の次元を求め、そこで電荷を測る場の強さ $F_{8-p}={*F_{p+2}}$ の次数と一致することを確かめよ。\\
(2) D$p$ ブレーンの磁気的双対が D$(6-p)$ ブレーンであることを示し、IIA型・IIB型それぞれで電気・磁気の対を列挙せよ。自分自身と双対なものはどれか。\\
(3) 一般の $D$ 次元で、$p$ ブレーンの磁気的双対は $(D-p-4)$ ブレーンである。$D=4$ の粒子、$D=6$ のひも、$D=11$ の M2 ブレーンについて双対を答えよ。
\end{problem}

\begin{solution}
(1) D$p$ の世界体積は $p+1$ 次元、横方向は $10-(p+1)=9-p$ 次元で、ブレーンは横方向の1点に見える。その点を囲む球面は $S^{8-p}$。$F_{p+2}$ の10次元のホッジ双対は $10-(p+2)=8-p$ 形式で、球面の次元と一致するので $\int_{S^{8-p}}{*F_{p+2}}$ が定義でき、ガウスの法則によりこれが電荷を測る。

(2) $F_{8-p}=dC_{7-p}$ に電気的に結合するのは、世界体積 $7-p$ 次元、すなわち $(6-p)$ ブレーンである。$C_{7-p}$ の電気的な源は $F_{8-p}$ の、つまり $F_{p+2}$ の磁気的な源なので、D$p$ と D$(6-p)$ が電気・磁気の対になる。$p$ と $6-p$ の偶奇は同じなので、同じ理論の中で対になる。
\begin{itemize}
\item IIA：D0–D6、D2–D4（D8 は $F_{10}$ に対応する特別な場合で、双対の $F_0$ は「質量」パラメータ）。
\item IIB：D$(-1)$–D7、D1–D5、D3–D3。D3 は自己双対で、$F_5={*F_5}$ に対応する。$N$ 枚の D3 の電荷は $\int_{S^5}F_5=2\pi N$ で、$\mathrm{AdS}_5\times S^5$ の $S^5$ を貫くフラックスになる。
\end{itemize}

(3) $D=4$, $p=0$：$4-0-4=0$、粒子の双対は粒子（電荷と磁気単極子、第5, 8章）。$D=6$, $p=1$：$6-1-4=1$、ひもは自分自身と双対（自己双対な3形式場 $H={*H}$ のひも）。$D=11$, $p=2$：$11-2-4=5$、M2 の双対は M5 ブレーン（4形式 $G_4$ と $G_7={*G_4}$）。
\end{solution}

\begin{problem}[難易度 D]\label{ch10-prob-dirac}
(1) 定理 \ref{thm:ch10-dirac} を示せ。（ヒント：磁気的な物体を囲む $S^{p+2}$ の上で、電気的 $p$ ブレーンの閉じた世界体積 $\Sigma_{p+1}$（例えば赤道 $S^{p+1}$）に沿った位相因子 $\exp(i\mu_e\oint_\Sigma C)$ を考える。$C$ は $S^{p+2}$ 全体では定義できないので、$\Sigma$ を境界とする北半球 $B_+$ と南半球 $B_-$ を用いて $\exp(i\mu_e\int_{B_\pm}F)$ と書き直し、両者の一致を要求せよ。）\\
(2) 電荷の単位を $\mu_e=1$ にとると、コンパクトな内部空間 $K$ 上のフラックス $F_{p+2}$ は格子 $H^{p+2}(K;\Z)/\text{ねじれ}$ に値をとる。IIB型弦理論を5次超曲面にコンパクト化したとき、3形式フラックス $F_3$, $H_3$ を指定する整数は全部で何個あるか。\\
(3) M理論を CY4 $Y$ にコンパクト化するときのタドポール条件（命題 \ref{prop:ch10-tad}）
\[
N_{M2}+\frac12\int_Y\frac{G}{2\pi}\wedge\frac{G}{2\pi}=\frac{\chi(Y)}{24}
\]
を考える。(a) $Y=\mathrm{K3}\times\mathrm{K3}$、$G=0$ のとき $N_{M2}$ を求めよ。(b) $Y$ を $\CP^5$ の6次超曲面とすると $\chi(Y)=2610$（問題 \ref{ch10-prob-quintic}）で $\chi/24$ は整数でない。ウィッテンのフラックス量子化条件「$\frac{G}{2\pi}-\frac12c_2(Y)\in H^4(Y;\Z)$」と $c_2(Y)=15H^2$, $\int_YH^4=6$ を用いて、$G=0$ が許されないことを示し、$\frac{G}{2\pi}=\frac12H^2$ のとき $N_{M2}$ が整数になることを確かめよ。
\end{problem}

\begin{solution}
(1) 磁気的な物体は $S^{p+2}$ の内部にあり、$S^{p+2}$ 上で $F=F_{p+2}$ は閉形式で $\int_{S^{p+2}}F=\mu_m$。$\mu_m\ne0$ なら $F$ は $S^{p+2}$ 上で完全形式でなく、$C$ は大域的に定義できない（第5章のディラック・モノポールと同じ）。しかし北半球 $B_+$、南半球 $B_-$ はそれぞれ可縮なので、その上では $F=dC_\pm$ と書ける。

電気的 $p$ ブレーンの世界体積が赤道 $\Sigma=S^{p+1}$ を掃くとき、その経路積分の位相は $\exp(i\mu_e\oint_\Sigma C)$ で、ストークスの定理より北側の表示では $\exp\bigl(i\mu_e\int_{B_+}F\bigr)$、南側では $\exp\bigl(-i\mu_e\int_{B_-}F\bigr)$ と書ける（$\partial B_+=\Sigma$、$\partial B_-=-\Sigma$ と向きを付ける。$S^{p+2}=B_+\cup B_-$ は外向きの向き）。物理的な位相は表示によらないので
\[
\exp\Bigl(i\mu_e\int_{B_+}F\Bigr)=\exp\Bigl(-i\mu_e\int_{B_-}F\Bigr)
\iff \exp\Bigl(i\mu_e\int_{S^{p+2}}F\Bigr)=1
\iff \mu_e\mu_m\in2\pi\Z .
\]
（$C_+$ と $C_-$ は赤道の近傍でゲージ変換 $C_+-C_-=d\lambda_p$ で結ばれ、$\exp(i\mu_e\oint_\Sigma d\lambda)$ に相当する「大きなゲージ変換」の一価性と言ってもよい。）磁気的な物体の横方向は $S^{p+2}$ を含む $p+3$ 次元、世界体積は $D-(p+3)=D-p-3$ 次元なので $(D-p-4)$ ブレーンである。$\mu_e=1$ なら $\mu_m\in2\pi\Z$ で、任意の閉 $(p+2)$ 次元部分多様体 $\Sigma$ について同じ議論（$\Sigma$ を2つに分けて境界に沿ってブレーンを動かす）から $\int_\Sigma F\in2\pi\Z$。

(2) $[F/2\pi]$ は $H^{p+2}(K;\R)$ の中の格子（$H^{p+2}(K;\Z)$ の像）に属し、その階数は $b_{p+2}(K)$。5次超曲面では $b_3=204$（問題 \ref{ch10-prob-h21}）なので、$F_3$ と $H_3$ でそれぞれ204個、計408個の整数がフラックスを指定する。（これらの整数の組み合わせの膨大さが、いわゆる弦理論の「ランドスケープ」の起源の一つである。）

(3)(a) キュネットの公式で $\chi(\mathrm{K3}\times\mathrm{K3})=24\cdot24=576$。$N_{M2}=576/24=24$。

(b) $\chi/24=2610/24=108.75$。$G=0$ なら量子化条件は $-\frac12c_2=-\frac{15}{2}H^2\in H^4(Y;\Z)$ を要求する。整数係数の類 $\alpha$ について $\int_Y\alpha\wedge\alpha\in\Z$ だが、$\alpha=-\frac{15}{2}H^2$ なら $\int\alpha^2=\frac{225}{4}\cdot6=337.5\notin\Z$。よって $G=0$ は許されない（タドポール条件の右辺が整数でないことと整合）。

$\frac{G}{2\pi}=\frac12H^2$ とすると $\frac{G}{2\pi}-\frac12c_2=\frac12H^2-\frac{15}{2}H^2=-7H^2$ で、$H$ は整数係数なので条件をみたす。このとき
\[
\frac12\int_Y\Bigl(\frac{H^2}{2}\Bigr)^2=\frac18\int_YH^4=\frac68=\frac34,\qquad N_{M2}=108.75-0.75=108\in\Z .
\]
特性類に由来する半整数のずれと、タドポール条件の右辺の分数部分がちょうど打ち消し合う。（超対称性の条件など物理的な制約はここでは考えていない。）
\end{solution}

\subsection{グリーン・シュワルツ機構}

\begin{problem}[難易度 C]\label{ch10-prob-gs}
10次元 $\mathcal N=1$ 超重力＋超対称ヤン・ミルズ理論（ゲージ群 $G$）のアノマリーが、グリーン・シュワルツ機構で相殺されるための条件は次の2つである（定理 \ref{thm:ch10-gs}）。(i) $\dim G=496$。(ii) 随伴表現のトレース $\operatorname{Tr}F^6$ が、より低次のトレースの積で書ける（既約な6次のカシミールの寄与がない）。\\
(1) $SO(N)$ では、ベクトル表現のトレース $\tr$ と随伴表現のトレース $\operatorname{Tr}$ の間に
\begin{align*}
\operatorname{Tr}F^2&=(N-2)\tr F^2,\\
\operatorname{Tr}F^4&=(N-8)\tr F^4+3(\tr F^2)^2,\\
\operatorname{Tr}F^6&=(N-32)\tr F^6+15\tr F^2\tr F^4
\end{align*}
が成り立つ。(i), (ii) がともに $N=32$ を与えることを示せ。$\frac1{30}\operatorname{Tr}F^2=\tr F^2$ となることも確かめよ。\\
(2) $E_8$ は $SO(16)$ の随伴表現 $\mathbf{120}$ と、一方のカイラリティのスピノル表現から作られる。スピノルの次元を求めて $\dim E_8=248$ を確かめ、$E_8\times E_8$ が (i) をみたすことを示せ。$E_8$ では $\operatorname{Tr}F^4=\frac1{100}(\operatorname{Tr}F^2)^2$, $\operatorname{Tr}F^6=\frac1{7200}(\operatorname{Tr}F^2)^3$ が成り立つ（既約な4次・6次のカシミールがない）ことを認めて、(ii) もみたすことを説明せよ。\\
(3) $SU(N)$ が (i) をみたしえないことを示せ。\\
(4) 修正されたビアンキ恒等式 $dH\propto\tr R\wedge R-\frac1{30}\operatorname{Tr}F\wedge F$ が CY3 上で課す位相的条件を述べ、標準埋め込みがそれをみたす理由を説明せよ。
\end{problem}

\begin{solution}
(1) $\dim SO(N)=\frac{N(N-1)}{2}=496$ $\iff$ $N^2-N-992=0$ $\iff$ $(N-32)(N+31)=0$ で、$N=32$。一方 (ii) は $\operatorname{Tr}F^6$ の既約部分 $(N-32)\tr F^6$ が消えることなので、やはり $N=32$。2つの独立な条件が同じ $N$ を選ぶのは注目すべき一致である。$N=32$ で $\operatorname{Tr}F^2=30\tr F^2$、すなわち $\frac1{30}\operatorname{Tr}F^2=\tr F^2$。さらに
\[
\operatorname{Tr}F^6=15\tr F^2\tr F^4,\qquad \operatorname{Tr}F^4=24\tr F^4+3(\tr F^2)^2
\]
となって、$I_{12}$ が $\tr F^2$ と $\tr F^4$ の多項式で書け、因数分解の可能性が開ける。

(2) $SO(16)$ のディラック・スピノルは $2^{16/2}=256$ 次元、一方のカイラリティは $128$ 次元。$\dim E_8=120+128=248$。$\dim(E_8\times E_8)=2\cdot248=496$ で (i) をみたす。$F=F_1+F_2$（各因子の場の強さ）とすると $\operatorname{Tr}F^6=\operatorname{Tr}F_1^6+\operatorname{Tr}F_2^6=\frac1{7200}\bigl[(\operatorname{Tr}F_1^2)^3+(\operatorname{Tr}F_2^2)^3\bigr]$ で、すべて $\operatorname{Tr}F^2$ の積で書けるので (ii) もみたす。

(3) $\dim SU(N)=N^2-1=496$ なら $N^2=497$。$22^2=484<497<529=23^2$ なので整数解がない。

(4) $H$ は大域的な3形式なので $\tr R\wedge R-\frac1{30}\operatorname{Tr}F\wedge F$ は完全形式で、ド・ラーム・コホモロジー類としてゼロ：$[\tr R\wedge R]=[\tfrac1{30}\operatorname{Tr}F\wedge F]\in H^4(X;\R)$。これは（適当な規格化で）ゲージ束 $V$ と接束 $TX$ について $\ch_2(V)=\ch_2(TX)$、$SU(3)$ 束では $c_2(V)=c_2(TX)$ を意味する（命題 \ref{prop:ch10-tad}）。標準埋め込みではゲージ接続をスピン接続そのものに等しくとるので、$\tr F\wedge F$ と $\tr R\wedge R$ が形式として等しく、$dH=0$ が局所的に成り立ち、条件は自明にみたされる。$V\ne TX$ の一般のベクトル束を使う場合（非標準埋め込み）は、この位相的条件（5ブレーンがあればその寄与も含めたタドポール条件）を確かめる必要がある。
\end{solution}

%================================================
\chapter{モース理論・超対称量子力学・K理論への入口}\label{ch:11}
% 15問
%================================================

\section{物理からの動機}

\subsection*{超対称量子力学の基底状態を数える}
超対称性が自発的に破れるかどうかは、エネルギー $E=0$ の状態（超対称な真空）が存在するかどうかで決まる。超対称性を持つ理論では $H=\{Q,Q^\dagger\}\ge0$ なので、$E=0$ の状態は $Q\psi=Q^\dagger\psi=0$ をみたす状態にほかならない。ところが、摂動論で $E=0$ の状態が見つかっても、トンネル効果（インスタントン）によってエネルギーがわずかに持ち上がり、超対称性が\textbf{非摂動的に}破れることがある。破れるかどうかを摂動論に頼らず判定する量として、ウィッテンは
\[
\Tr(-1)^F e^{-\beta H}=n_B^{E=0}-n_F^{E=0}
\]
（\term{ウィッテン指数}{Witten index}）を導入した。$E>0$ の状態はボソンとフェルミオンが必ず対になるので、この量は $\beta$ によらず、パラメータを連続的に変えても変わらない。指数が $0$ でなければ超対称性は破れない。

1次元の超対称量子力学で、超ポテンシャル $W(x)=\lambda(x^3/3-a^2x)$ をとると、ボソン的なポテンシャルは $V=W'^2=\lambda^2(x^2-a^2)^2$ の二重井戸になる（下図）。調和近似では2つの井戸の底にそれぞれ $E=0$ の状態ができるが、一方はボソン的、もう一方はフェルミオン的であり、指数は $1-1=0$。実際には井戸の間のトンネル効果がこの2状態を混ぜて $E\sim e^{-2\Delta W}>0$ に持ち上げ、超対称性は破れる（問題 \ref{ch11-prob-dw}）。

\begin{center}
\begin{tikzpicture}[scale=1.15]
  \draw[-{Stealth[length=2mm]}] (-2.6,0) -- (2.8,0) node[right] {$x$};
  \draw[-{Stealth[length=2mm]}] (0,-0.2) -- (0,2.5) node[above] {$V=W'^2$};
  \draw[very thick,blue!70!black,domain=-1.6:1.6,samples=100,smooth]
     plot (\x,{0.9*(\x*\x-1)*(\x*\x-1)});
  \node[vtx] at (1,0) {}; \node[vtx] at (-1,0) {};
  \node[below] at (1,-0.05) {\footnotesize $+a$};
  \node[below] at (-1,-0.05) {\footnotesize $-a$};
  \draw[red!70!black,thick] (0.75,0.32) -- (1.25,0.32);
  \draw[green!40!black,thick] (-1.25,0.32) -- (-0.75,0.32);
  \node[red!70!black,right] at (1.3,0.42) {\footnotesize ボソン（$W''>0$）};
  \node[green!40!black,left] at (-1.3,0.42) {\footnotesize フェルミオン（$W''<0$）};
  \draw[-{Stealth[length=2mm]},thick,dashed,orange!80!black] (-0.75,0.5) .. controls (-0.3,1.3) and (0.3,1.3) .. (0.75,0.5);
  \draw[-{Stealth[length=2mm]},thick,dashed,orange!80!black] (0.75,0.45) .. controls (0.3,1.1) and (-0.3,1.1) .. (-0.75,0.45);
  \node[orange!80!black] at (0,1.55) {\footnotesize インスタントン};
\end{tikzpicture}
\zucap{超ポテンシャル $W=\lambda(x^3/3-a^2x)$ の超対称量子力学。調和近似では各井戸に $E=0$ の状態が1つずつあるが、統計が逆である。トンネル効果で両者が結びつき、どちらも $E>0$ に持ち上がる。}
\end{center}

\subsection*{超対称量子力学 $=$ ド・ラーム理論}
ウィッテン（1982）は、リーマン多様体 $M$ 上を動く粒子の超対称量子力学が、$M$ 上の微分形式の理論そのものであることに気づいた。
\begin{itemize}
\item 状態空間は $M$ 上の微分形式全体、フェルミオン数 $F$ は形式の次数。
\item 超電荷は外微分 $Q=d$、ハミルトニアンはラプラシアン $H=dd^\dagger+d^\dagger d$。
\item $E=0$ の状態は調和形式であり、ホッジの定理（第5章）によりド・ラーム・コホモロジーと1対1。よってウィッテン指数はオイラー標数 $\chi(M)=\sum(-1)^pb_p$ である。
\end{itemize}
さらに $M$ 上の関数 $h$ を超ポテンシャルとして加え（$d_t=e^{-th}de^{th}$）、結合 $t$ を大きくすると、低エネルギーの状態は $h$ の臨界点の近くに局在し、各臨界点に「フェルミオン数 $=$ その臨界点の指数」の摂動的真空が1つずつできる。真の基底状態はこのうちの一部であり、どれが生き残るかは臨界点を結ぶインスタントン（$h$ の勾配流）が決める。こうして\textbf{臨界点の数}と\textbf{トポロジー（ベッチ数）}を結ぶ古典的な\term{モース理論}{Morse theory}が、超対称量子力学の言葉で自然に導かれる。

\subsection*{Dブレーンの電荷とK理論}
タイプII弦理論の Dp ブレーンは RR 場の電荷をもち、ド・ラーム・コホモロジー（RR 形式の束縛）で数えるのが素朴な考えである。しかし、BPS でない（超対称性を保たない）にもかかわらず安定なブレーンがあり、その電荷は $\Z_2$ に値をとる（タイプI 弦理論の D0 ブレーンなど）。これはコホモロジーでは捉えられない。

センは、D9 ブレーンと反 D9 ブレーンの対の上のタキオン場が凝縮すると両者が対消滅し、タキオンが位相的に巻き付いた渦の芯に低次元の Dブレーンが残ることを示した。ブレーンと反ブレーンの上のゲージ束を $E,F$ とすると、「対消滅で消えるもの」を同一視した電荷は形式的差 $[E]-[F]$ であり、これが\term{K理論}{K-theory}の元である。球面の K 群は\term{ボット周期性}{Bott periodicity}によって完全に計算でき、その結果は既知の安定な Dブレーンの一覧と見事に一致する。

\begin{physbox}[：本章の辞書]
\begin{center}
\small
\begin{tabular}{p{0.47\linewidth}p{0.43\linewidth}}
物理 & 数学\\ \hline
超対称量子力学の状態・フェルミオン数 & 微分形式・形式の次数\\
超電荷 $Q$、ハミルトニアン $H$ & 外微分 $d$、ラプラシアン $\Delta$\\
超対称な真空（$E=0$ の状態） & 調和形式 $\cong$ ド・ラーム・コホモロジー\\
ウィッテン指数 $\Tr(-1)^F$ & オイラー標数 $\chi(M)$\\
超ポテンシャル $t\,h$ & モース関数 $h$\\
摂動的真空 & 臨界点（フェルミオン数 $=$ 指数）\\
インスタントン（トンネル効果） & 勾配流線、モース・ウィッテン微分\\
トンネル効果による超対称性の破れ & モース複体の微分が $0$ でない\\
ブレーン・反ブレーン対（Chan--Paton 束 $E,F$） & K理論の元 $[E]-[F]$\\
対消滅 & $[E]-[E]=0$\\
Dブレーン電荷（$\Z$ や $\Z_2$） & $\tilde K(S^n)$, $\widetilde{KO}(S^n)$
\end{tabular}
\end{center}
\end{physbox}

\section{基本事項のまとめ}

\subsection{モース関数}
以下 $M$ は $n$ 次元の滑らかな多様体、$h\colon M\to\R$ は滑らかな関数とする。

\begin{mydef}{臨界点・モース関数・指数}{ch11-morse}
$dh_p=0$ となる点 $p$ を $h$ の\term{臨界点}{critical point}という。局所座標でのヘッセ行列 $\bigl(\partial_i\partial_jh(p)\bigr)$ が正則なとき $p$ は\term{非退化}{nondegenerate}であるといい、その負の固有値の個数（重複込み）を $p$ の\term{指数}{index}（モース指数、$\ind p$）という（座標によらない）。すべての臨界点が非退化な関数を\term{モース関数}{Morse function}という。指数 $k$ の臨界点の個数を $c_k$ と書く。
\end{mydef}

\begin{mythm}{モースの補題}{ch11-lemma}
$p$ が指数 $k$ の非退化臨界点ならば、$p$ のまわりの局所座標 $(x_1,\dots,x_n)$ で
\[
h=h(p)-x_1^2-\dots-x_k^2+x_{k+1}^2+\dots+x_n^2
\]
となるものがとれる。特に非退化臨界点は孤立している。
\end{mythm}
物理の言葉では、臨界点の近くで $h$ は（座標を選べば）厳密に2次であり、指数は「不安定方向の数」である。極小は指数 $0$、極大は指数 $n$、その間が鞍点である。コンパクト多様体上の「一般の」関数はモース関数である（モース関数全体は稠密）。

\begin{myex}{立てたトーラスの高さ関数}{ch11-torus}
トーラス $T^2$ を鉛直面内に立て、高さ $h$ を関数とする。臨界点は下図の4点で、下から順に極小（指数 $0$）、鞍点（指数 $1$）、鞍点（指数 $1$）、極大（指数 $2$）である。$(c_0,c_1,c_2)=(1,2,1)$ はベッチ数 $(b_0,b_1,b_2)=(1,2,1)$ と一致する（問題 \ref{ch11-prob-torus}）。
\begin{center}
\begin{tikzpicture}[scale=0.95]
  \fill[blue!8,even odd rule] (0,0) circle (2) (0,0) circle (1);
  \draw[thick] (0,0) circle (2);
  \draw[thick] (0,0) circle (1);
  % tube cross sections
  \draw[gray] (0,1.5) ellipse (0.18 and 0.5);
  \draw[gray] (0,-1.5) ellipse (0.18 and 0.5);
  \draw[gray] (1.5,0) ellipse (0.5 and 0.18);
  \draw[gray] (-1.5,0) ellipse (0.5 and 0.18);
  \node[vtx,red!70!black] at (0,-2) {}; \node[below] at (0,-2) {\footnotesize 極小（指数 0）};
  \node[vtx,orange!80!black] at (0,-1) {}; \node[above] at (0,-1) {\footnotesize 鞍点（1）};
  \node[vtx,orange!80!black] at (0,1) {}; \node[below] at (0,1) {\footnotesize 鞍点（1）};
  \node[vtx,blue!70!black] at (0,2) {}; \node[above] at (0,2) {\footnotesize 極大（指数 2）};
  \draw[-{Stealth[length=2mm]}] (3.3,-2.3) -- (3.3,2.5) node[above] {$h$};
  \foreach \y/\l in {-2/-2,-1/-1,1/1,2/2}{
    \draw[densely dotted] (0.1,\y) -- (3.3,\y);
    \node[right] at (3.35,\y) {\footnotesize $\l$};
  }
\end{tikzpicture}
\zucap{鉛直に立てたトーラスの正面図（外側の輪郭の半径 $R+r=2$、穴の輪郭の半径 $R-r=1$）。灰色の楕円は管の断面。高さ関数 $h$ の臨界点は4つ。}
\end{center}
\end{myex}

\subsection{胞体分解とモース不等式}
部分レベル集合 $M^a=\{x\in M:h(x)\le a\}$ が $a$ とともにどう変わるかを見る。

\begin{mythm}{臨界点を通過すると胞体が付く}{ch11-handle}
$M$ はコンパクト、$h$ はモース関数とする。
\begin{enumerate}
\item $[a,b]$ に臨界値がなければ、$M^a$ は $M^b$ の変位レトラクトで、$M^a\cong M^b$。
\item 臨界値 $c$ に指数 $k$ の臨界点がちょうど1つあるなら、十分小さい $\varepsilon>0$ で $M^{c+\varepsilon}\simeq M^{c-\varepsilon}\cup e^k$（$k$ 胞体を1つ貼り付けたもの）。
\end{enumerate}
したがって $M$ は、指数 $k$ の臨界点1つにつき $k$ 胞体1つをもつ CW 複体とホモトピー同値である。
\end{mythm}
\begin{note}
(2) の直観：モースの補題の座標で、臨界点を通過すると「不安定方向 $x_1,\dots,x_k$ の $k$ 次元円板」が新たに $M^c$ に入ってくる。安定方向 $x_{k+1},\dots,x_n$ は縮められるので、ホモトピー的には $k$ 胞体を貼るだけである。
\end{note}

胞体ホモロジー（第4章）で $\dim C_k=c_k$ の鎖複体から $H_k$ が計算されることから、次が従う。$b_k=\dim H_k(M;\mathbb{F})$ は任意の体 $\mathbb{F}$ 係数のベッチ数（$\mathbb{F}=\Q$ か $\Z_2$）である。

\begin{mythm}{モース不等式}{ch11-ineq}
$M$ をコンパクト、$h$ をモース関数とする。
\begin{enumerate}
\item（弱モース不等式）$c_k\ge b_k$ \ （$k=0,1,\dots,n$）。
\item（強モース不等式）各 $k$ について
\[
c_k-c_{k-1}+\dots+(-1)^kc_0\ \ge\ b_k-b_{k-1}+\dots+(-1)^kb_0 ,
\]
かつ $k=n$ では等号：$\displaystyle\sum_k(-1)^kc_k=\sum_k(-1)^kb_k=\chi(M)$。
\item（多項式形）$M_t=\sum c_kt^k$、$P_t=\sum b_kt^k$ とおくと、非負整数係数の多項式 $R(t)$ があって
\[
M_t-P_t=(1+t)R(t).
\]
\end{enumerate}
\end{mythm}
\begin{note}
(3) の導出：鎖複体 $0\to C_n\to\dots\to C_0\to0$ で $Z_k=\ker\partial_k$、$B_k=\operatorname{im}\partial_{k+1}$、$r_k=\dim B_k$ とおくと $c_k=\dim Z_k+r_{k-1}$、$b_k=\dim Z_k-r_k$ なので $c_k-b_k=r_k+r_{k-1}$。よって $M_t-P_t=\sum_k(r_k+r_{k-1})t^k=(1+t)\sum_kr_kt^k$。(1)(2) は (3) の係数比較から従う（問題 \ref{ch11-prob-ineq}）。
\end{note}

$c_k=b_k$（すべての $k$）となるモース関数を\term{完全モース関数}{perfect Morse function}という。指数がすべて偶数なら（より一般に、隣り合う指数の臨界点がなければ）完全である（\term{欠落原理}{lacunary principle}、問題 \ref{ch11-prob-ineq}）。

\subsection{モース・ホモロジー（ウィッテン複体）}
$M$ にリーマン計量を入れ、勾配流
\[
\frac{dx}{d\tau}=-\nabla h(x)
\]
を考える。$h$ は流れに沿って減少する。臨界点 $p$ から出発して（$\tau\to-\infty$ で $p$ に近づき）$q$ に流れ込む（$\tau\to+\infty$ で $q$ に近づく）解を、$p$ から $q$ への\term{勾配流線}{gradient flow line}という（時間の平行移動は同一視する）。$p$ から出る流線全体は $p$ の\term{不安定多様体}{unstable manifold} $W^u(p)$ をなし、その次元は $\ind p$ である。$q$ に流れ込む流線全体 $W^s(q)$ の次元は $n-\ind q$。

すべての $p,q$ で $W^u(p)$ と $W^s(q)$ が横断的に交わるとき（\term{モース・スメール条件}{Morse--Smale condition}；計量を少し動かせば常に成り立つ）、$\ind p=\ind q+1$ なら $p$ から $q$ への流線は有限個である。各 $W^u(p)$ に向きを選んでおくと、各流線 $\gamma$ に符号 $\epsilon(\gamma)=\pm1$ が定まる。

\begin{mythm}{モース・ホモロジー}{ch11-morsehom}
$C_k$ を指数 $k$ の臨界点で生成される自由加群とし、
\[
\partial p=\sum_{\ind q=\ind p-1}n(p,q)\,q,\qquad n(p,q)=\sum_{\gamma:\,p\to q}\epsilon(\gamma)
\]
と定めると $\partial^2=0$ であり、$H_k(C_\bullet,\partial)\cong H_k(M;\Z)$。$\Z_2$ 係数では符号を無視して流線の数を $\bmod 2$ で数えればよい。
\end{mythm}
符号の規則は、1次元の場合が典型である：$\ind p=1$ の臨界点の不安定多様体は向きのついた弧で、その2つの端に流れ着く極小 $q_\pm$（向きの正の側を $q_+$）に対し $\partial p=q_+-q_-$ である。これは胞体の境界 $\partial e^1=(\text{終点})-(\text{始点})$ と同じである。
ウィッテンの立場では、余境界 $\delta=\partial^T$（指数を1上げる）が変形した外微分 $d_t$ の行列要素（インスタントン振幅）として現れる（後述）。

\subsection{超対称量子力学とウィッテン指数}

\begin{mydef}{超対称量子力学}{ch11-sqm}
ヒルベルト空間 $\mathcal{H}=\mathcal{H}_B\oplus\mathcal{H}_F$ 上に、\term{フェルミオン数演算子}{fermion number} $(-1)^F$（$\mathcal H_B$ 上 $+1$、$\mathcal H_F$ 上 $-1$）と、$(-1)^F$ と反可換な\term{超電荷}{supercharge} $Q$ があって
\[
Q^2=0,\qquad H=\{Q,Q^\dagger\}=QQ^\dagger+Q^\dagger Q
\]
となる系を（$\mathcal N=2$ の）超対称量子力学という。$[H,Q]=0$、$H\ge0$ である。
\end{mydef}

\begin{myprop}{エネルギー正の状態の対合とウィッテン指数}{ch11-pair}
\begin{enumerate}
\item $H\psi=0\iff Q\psi=Q^\dagger\psi=0$。
\item $E>0$ の固有空間 $\mathcal H_E$ では、$Q_1=(Q+Q^\dagger)/\sqrt E$ が $\mathcal H_E\cap\mathcal H_B$ と $\mathcal H_E\cap\mathcal H_F$ の同型を与える。
\item したがって\term{ウィッテン指数}{Witten index}
\[
I_W=\Tr(-1)^Fe^{-\beta H}=n_B^0-n_F^0
\]
は $\beta$ によらない。$n_{B,F}^0$ は $E=0$ のボソン・フェルミオン状態の数。
\item $I_W\neq0$ ならば $E=0$ の状態が存在し、超対称性は自発的に破れない。
\end{enumerate}
\end{myprop}
証明は問題 \ref{ch11-prob-pair}。$I_W$ はパラメータの連続変形で変わらない：状態は $E>0$ の対として出入りするだけだからである。ただし、非コンパクトな空間で漸近的振る舞いが変わる変形（基底状態が無限遠に逃げる）では変わりうる（問題 \ref{ch11-prob-W}）。

\begin{myex}{多様体上の超対称量子力学}{ch11-sqmM}
コンパクトなリーマン多様体 $M$ 上の超対称シグマ模型（符号・規格化は文献により異なるが、ラグランジアン $L=\tfrac12g_{ij}\dot x^i\dot x^j+i\bar\psi_i\tfrac{D\psi^i}{dt}-\tfrac14R_{ijkl}\psi^i\bar\psi^j\psi^k\bar\psi^l$）を量子化すると、フェルミオンは $\psi^i\to dx^i\wedge$、$\bar\psi_i\to\iota_{\partial_i}$ のように微分形式への作用として実現され
\[
\mathcal H=\bigoplus_{p=0}^n\Omega^p(M)\ (L^2),\qquad F=p,\qquad Q=d,\quad Q^\dagger=d^\dagger,\quad H=dd^\dagger+d^\dagger d=\Delta .
\]
$E=0$ の状態は調和形式であり、ホッジの定理（第5章）により $\dim\ker\Delta|_{\Omega^p}=b_p(M)$。よって
\[
\Tr(-1)^Fe^{-\beta H}=\sum_p(-1)^pb_p(M)=\chi(M).
\]
\end{myex}

\subsection{ウィッテンの変形}
$h$ を $M$ 上の関数、$t\in\R$ とし
\[
d_t=e^{-th}\,d\,e^{th}=d+t\,dh\wedge,\qquad d_t^\dagger=e^{th}\,d^\dagger e^{-th},\qquad H_t=d_td_t^\dagger+d_t^\dagger d_t
\]
とおく。$d_t^2=0$ であり、$d_t$ のコホモロジーは $e^{th}$ との共役で $d$ のコホモロジーと同型なので、$H_t$ の $E=0$ 状態の数は各次数で $b_p$ のままである（問題 \ref{ch11-prob-pair}）。これは超ポテンシャル $th$ をもつ超対称量子力学である。

\begin{myprop}{変形したハミルトニアン}{ch11-Ht}
局所的な正規直交座標（平坦近似）で $a^{i\dagger}=dx^i\wedge$、$a^i=\iota_{\partial_i}$（$\{a^i,a^{j\dagger}\}=\delta^{ij}$）とすると
\[
H_t=\Delta+t^2\lvert\nabla h\rvert^2+t\sum_{i,j}\frac{\partial^2h}{\partial x^i\partial x^j}\,[a^{i\dagger},a^j].
\]
\end{myprop}
\begin{note}
1次元 $M=\R$ では、$0$ 形式 $f$ と $1$ 形式 $g\,dx$ に対し $d_tf=(f'+th'f)\,dx$、$d_t^\dagger(g\,dx)=-g'+th'g$ で
\[
H_t\big|_{\Omega^0}=-\frac{d^2}{dx^2}+t^2h'^2-th'',\qquad H_t\big|_{\Omega^1}=-\frac{d^2}{dx^2}+t^2h'^2+th''.
\]
$[a^\dagger,a]$ は $0$ 形式上 $-1$、$1$ 形式上 $+1$ であり、一般式と一致する。$E=0$ の解は $0$ 形式 $e^{-th}$、$1$ 形式 $e^{th}dx$ である。
\end{note}

$t\to\infty$ では $t^2\lvert\nabla h\rvert^2$ が大きくなるので、低エネルギーの状態は臨界点の近くに局在する。臨界点 $p$ でヘッセ行列を対角化して固有値を $\lambda_1,\dots,\lambda_n$ とすると、調和近似で $H_t$ は独立な超対称調和振動子の和になり、その固有値は
\[
E=t\sum_{i=1}^n\bigl(\lvert\lambda_i\rvert(2N_i+1)+\lambda_i\varepsilon_i\bigr),\qquad N_i=0,1,2,\dots,\quad \varepsilon_i=\begin{cases}+1&(dx^i\ \text{を含む})\\-1&(\text{含まない})\end{cases}
\]
である。$E=0$ となるのは $N_i=0$ かつ「$\lambda_i<0$ の方向にだけ $dx^i$ を含む」ときで、ちょうど1状態、その次数は $\ind p$ である（問題 \ref{ch11-prob-witten}）。

\begin{mythm}{ウィッテンによるモース理論}{ch11-witten}
$t\to\infty$ で、$H_t|_{\Omega^k}$ の固有値のうち $O(t)$ より小さいもの（低エネルギー状態）の個数は $c_k$ に等しく、それらは指数 $k$ の臨界点に局在した摂動的真空で近似される。
\begin{enumerate}
\item 真の $E=0$ 状態は低エネルギー状態の一部なので $b_k\le c_k$（弱モース不等式）。
\item 低エネルギー状態の空間に制限した $d_t$ は、指数 $k$ の臨界点の状態を指数 $k+1$ の臨界点の状態へ写し、その行列要素は2点を結ぶ勾配流線（インスタントン）の寄与の和 $\sim\pm e^{-t(h(q)-h(p))}$ で与えられる。この複体のコホモロジーが $E=0$ 状態を与え、モース・ホモロジー（定理 \ref{thm:ch11-morsehom}）の双対版になる。
\item 指数：$\displaystyle\chi(M)=\sum_k(-1)^kc_k$。
\end{enumerate}
\end{mythm}

\subsection{ボット周期性}
$U=\bigcup_NU(N)$、$O=\bigcup_NO(N)$、$Sp=\bigcup_NSp(N)$ を包含 $U(N)\subset U(N+1)$ などによる帰納極限とする（\term{安定}{stable}な古典群）。

\begin{mythm}{ボット周期性}{ch11-bott}
\[
\pi_k(U)\cong\pi_{k+2}(U),\qquad \pi_k(O)\cong\pi_{k+8}(O),\qquad \pi_k(Sp)\cong\pi_{k+4}(O).
\]
具体的には
\begin{center}
\begin{tabular}{c|cccccccc}
$k\bmod 8$ & 0 & 1 & 2 & 3 & 4 & 5 & 6 & 7\\ \hline
$\pi_k(U)$ & 0 & $\Z$ & 0 & $\Z$ & 0 & $\Z$ & 0 & $\Z$\\
$\pi_k(O)$ & $\Z_2$ & $\Z_2$ & 0 & $\Z$ & 0 & 0 & 0 & $\Z$\\
$\pi_k(Sp)$ & 0 & 0 & 0 & $\Z$ & $\Z_2$ & $\Z_2$ & 0 & $\Z$
\end{tabular}
\end{center}
\end{mythm}

\begin{myprop}{安定域}{ch11-stable}
ファイブレーション $U(N-1)\to U(N)\to S^{2N-1}$ などのホモトピー完全列（第3章）から、包含は次の範囲で $\pi_k$ の同型を与える：
\begin{gather*}
\pi_k(U(N))\cong\pi_k(U)\ \ (k\le2N-1),\qquad \pi_k(O(N))\cong\pi_k(O)\ \ (k\le N-2),\\
\pi_k(Sp(N))\cong\pi_k(Sp)\ \ (k\le4N+1).
\end{gather*}
また $SU(N)\to U(N)$ は $k\ge2$ で $\pi_k$ の同型を与える（$U(N)\cong SU(N)\times U(1)$）。
\end{myprop}
例えば $\pi_3(SU(N))=\Z$（$N\ge2$、インスタントン数、第3, 8章）、$\pi_1(SO(N))=\Z_2$（$N\ge3$、スピン、第2章）、$\pi_4(SU(N))=0$（$N\ge3$）だが $\pi_4(SU(2))=\pi_4(Sp(1))=\Z_2$（安定域内なので $Sp$ の値；ウィッテンの $SU(2)$ アノマリー、第9章）。

\begin{center}
\begin{tikzpicture}[scale=1.0]
  \draw[thick,gray] (0,0) circle (2.2);
  \foreach \k/\g in {0/\Z_2,1/\Z_2,2/0,3/\Z,4/0,5/0,6/0,7/\Z}{
    \pgfmathsetmacro{\ang}{90-45*\k}
    \node[circle,draw,fill=white,inner sep=1.5pt,minimum size=7mm] at (\ang:2.2) {\footnotesize $\g$};
    \node at (\ang:2.95) {\footnotesize $k\equiv\k$};
  }
  \draw[-{Stealth[length=2mm]},thick,blue!60!black] (80:1.75) arc (80:-250:1.75);
  \node at (0,0.25) {\footnotesize $\pi_k(O)$};
  \node at (0,-0.25) {\footnotesize 周期 8};
  \begin{scope}[xshift=6.4cm]
    \draw[thick,gray] (0,0) circle (1.2);
    \node[circle,draw,fill=white,inner sep=1.5pt,minimum size=7mm] at (90:1.2) {\footnotesize $0$};
    \node[circle,draw,fill=white,inner sep=1.5pt,minimum size=7mm] at (-90:1.2) {\footnotesize $\Z$};
    \node at (90:1.85) {\footnotesize $k$ 偶数};
    \node at (-90:1.85) {\footnotesize $k$ 奇数};
    \node at (0,0.2) {\footnotesize $\pi_k(U)$};
    \node at (0,-0.2) {\footnotesize 周期 2};
  \end{scope}
\end{tikzpicture}
\zucap{ボット周期性の「時計」。左：$\pi_k(O)$ は $k\bmod8$ で決まる（$k=0$ を真上に置き時計回り）。右：$\pi_k(U)$ は $k\bmod2$ で決まる。}
\end{center}

\subsection{K理論}

\begin{mydef}{K群}{ch11-K}
$X$ をコンパクト空間とする。$X$ 上の複素ベクトル束の同型類全体は直和 $\oplus$ について可換モノイドをなす。その\term{グロタンディーク群}{Grothendieck group}、すなわち形式的差 $[E]-[F]$ 全体に
\[
[E]-[F]=[E']-[F']\iff E\oplus F'\oplus\C^N\cong E'\oplus F\oplus\C^N\ (\text{ある } N)
\]
で同値関係を入れた群を $K(X)$ と書く。テンソル積で環になる。$X$ が基点 $x_0$ をもつとき、階数 $[E]-[F]\mapsto\dim E_{x_0}-\dim F_{x_0}$ の核を\term{簡約K群}{reduced K-group} $\tilde K(X)$ とし、$K(X)=\Z\oplus\tilde K(X)$。実ベクトル束で同様に作ったものを $KO(X)$、$\widetilde{KO}(X)$ と書く。
\end{mydef}
コンパクトな $X$ 上の任意の束 $F$ には $F\oplus F'\cong\C^N$ となる $F'$ があるので、$K(X)$ の元は常に $[E]-[\C^N]$ の形に書ける。$[E]=[E']$ は「安定同型」（自明束を足せば同型）を意味する。

\begin{mythm}{球面の K 群}{ch11-Ksphere}
球面上の階数 $N$ の束は、赤道 $S^{n-1}$ 上の貼り合わせ関数のホモトピー類で分類される（第7章）：$\mathrm{Vect}_N(S^n)\cong[S^{n-1},U(N)]$。安定化して
\[
\tilde K(S^n)\cong\pi_{n-1}(U)=\begin{cases}\Z&(n\ \text{偶数})\\0&(n\ \text{奇数})\end{cases},\qquad
\widetilde{KO}(S^n)\cong\pi_{n-1}(O).
\]
$\widetilde{KO}(S^n)$ は $n\bmod 8=0,1,\dots,7$ に応じて $\Z,\Z_2,\Z_2,0,\Z,0,0,0$。
\end{mythm}
$n=2m$ のとき同型 $\tilde K(S^{2m})\to\Z$ は $[E]-[\C^N]\mapsto\int_{S^{2m}}\ch_m(E)$ で与えられる（$\ch=\tr\exp(F/2\pi)$）。$\tilde K(S^2)$ の生成元は単位モノポール束 $L_1$ を用いた $[L_1]-1$（問題 \ref{ch11-prob-KS2}）、$\tilde K(S^4)$ の生成元はインスタントン数 $\pm1$ の $SU(2)$ 束である。

高次の K 群は懸垂で定める：$\tilde K^{-j}(X)=\tilde K(\Sigma^jX)$。特に $\tilde K^{-1}(S^n)=\tilde K(S^{n+1})$。ボット周期性は $\tilde K^{-j}\cong\tilde K^{-j-2}$、$\widetilde{KO}^{-j}\cong\widetilde{KO}^{-j-8}$ と表される。チャーン指標 $\ch\colon K(X)\otimes\Q\to H^{\mathrm{even}}(X;\Q)$ は環同型だが、$K(X)$ のねじれ（$KO$ の $\Z_2$ など）はコホモロジーでは見えない。

\begin{mythm}{アティヤ・ボット・シャピロ構成}{ch11-ABS}
$\R^{2k}$ のクリフォード代数 $\{\Gamma_\mu,\Gamma_\nu\}=2\delta_{\mu\nu}$ の既約表現をカイラリティで分けた $S^\pm\cong\C^{2^{k-1}}$ に対し、$T(x)=x^\mu\Gamma_\mu\colon S^+\to S^-$ は $x\neq0$ で同型である。三つ組 $(S^+,S^-,T)$ は「$\R^{2k}$ 上でコンパクト台をもつ K 理論の元」すなわち $\tilde K(S^{2k})$ の元を定め、これは生成元（ボット元）である。実クリフォード代数を用いると $\widetilde{KO}$ の生成元も同様に得られる。
\end{mythm}

\subsection{Dブレーン電荷とK理論}

\begin{physbox}[：センのタキオン凝縮とウィッテンの K 理論的分類]
タイプIIB 弦理論で、$N$ 枚の D9 ブレーンと $N$ 枚の反 D9 ブレーンを重ねる。それぞれの上のチャン・パトン束を $E,F$（階数 $N$）とすると、両者を結ぶ開弦の基底状態はタキオン $T\in\Gamma(\Hom(E,F))$ である。タキオンのポテンシャルの最小値 $\lvert T\rvert=T_0$ ではブレーン対は消滅して閉弦の真空になる（センの予想）。
\begin{itemize}
\item $T$ がいたるところ同型なら $E\cong F$ で、全体が真空に崩壊する：$[E]-[F]=0$。
\item $T$ が $\R^{2k}$ 方向に巻き付いた配位（アティヤ・ボット・シャピロ型 $T=x^\mu\Gamma_\mu$、$N=2^{k-1}$）は、$x=0$ の芯に D$(9-2k)$ ブレーンを残す。
\end{itemize}
よって Dp ブレーンの電荷は横方向 $\R^{9-p}$ のコンパクト台 K 群 $\tilde K(S^{9-p})$ に値をとる。
\begin{center}
\small
\begin{tabular}{l|l|l}
理論 & 電荷の群 & 安定な Dp ブレーン\\ \hline
IIB & $\tilde K(S^{9-p})$ & $p$ 奇数で $\Z$\\
IIA & $\tilde K^{-1}(S^{9-p})=\tilde K(S^{10-p})$ & $p$ 偶数で $\Z$\\
タイプI & $\widetilde{KO}(S^{9-p})$ & $p=1,5,9$ で $\Z$；$p=-1,0,7,8$ で $\Z_2$
\end{tabular}
\end{center}
タイプIの $\Z_2$ 電荷のブレーンは BPS ではないが安定である（問題 \ref{ch11-prob-dbrane}）。
\end{physbox}

\begin{center}
\begin{tikzpicture}[scale=0.95]
  % left: D9 and anti-D9
  \fill[blue!10] (0,1.4) -- (3,1.4) -- (3.8,2.2) -- (0.8,2.2) -- cycle;
  \draw[thick,blue!70!black] (0,1.4) -- (3,1.4) -- (3.8,2.2) -- (0.8,2.2) -- cycle;
  \fill[red!10] (0,0) -- (3,0) -- (3.8,0.8) -- (0.8,0.8) -- cycle;
  \draw[thick,red!70!black] (0,0) -- (3,0) -- (3.8,0.8) -- (0.8,0.8) -- cycle;
  \node[blue!70!black,left] at (0,1.8) {\footnotesize D9 ($E$)};
  \node[red!70!black,left] at (0,0.4) {\footnotesize $\overline{\mathrm{D9}}$ ($F$)};
  \draw[thick,domain=0.45:1.75,samples=60,smooth] plot ({1.9+0.06*sin(\x*900)},\x);
  \node[right] at (1.95,1.1) {\footnotesize $T$};
  \draw[-{Stealth[length=2.5mm]},very thick] (4.3,1.1) -- (5.6,1.1);
  \node[above] at (4.95,1.15) {\footnotesize 凝縮};
  % right: vortex leaving D7
  \begin{scope}[xshift=6.2cm,yshift=0.2cm]
    \draw[dashed,gray] (0,0) -- (3.4,0) -- (4.2,1.6) -- (0.8,1.6) -- cycle;
    \foreach \a in {0,45,...,315}{
      \draw[-{Stealth[length=1.6mm]},orange!80!black,thick] ($(2.1,0.8)+(\a:0.55 and 0.3)$) -- ++(\a+90:0.3 and 0.17);
    }
    \draw[very thick,green!40!black] (2.1,-0.5) -- (2.1,2.1);
    \node[green!40!black,right] at (2.15,2.0) {\footnotesize D7（芯）};
    \node[orange!80!black] at (3.5,0.25) {\footnotesize $T\propto x^1+ix^2$};
  \end{scope}
\end{tikzpicture}
\zucap{ブレーン・反ブレーン対消滅の模式図。左：D9 と反 D9 の間の開弦がタキオン $T$ を与える。右：$T$ が横方向の平面 $(x^1,x^2)$ で1回巻き付いた渦配位で凝縮すると、芯（余次元 2）に D7 ブレーンが残る（縦の線は芯の残りの方向を表す）。}
\end{center}

\section{問題}

\subsection{モース関数とモース不等式}

\begin{problem}[難易度 A]\label{ch11-prob-torus}
例 \ref{ex:ch11-torus} の立てたトーラスを、次のように $z$ 軸のまわりの回転面として与える：$T^2=\{((R+r\cos\phi)\cos\theta,\ (R+r\cos\phi)\sin\theta,\ r\sin\phi)\}$（$R>r>0$）。高さ関数として $h=x$（$x$ 座標）をとる（$x$ 軸を鉛直とみれば、トーラスは鉛直面 $xy$ 内に立っている）。\\
(1) $h(\theta,\phi)=(R+r\cos\phi)\cos\theta$ の臨界点を求め、各点のヘッセ行列から指数を求めよ。\\
(2) $c_k$ とトーラスのベッチ数 $b_k$ を比べ、弱・強モース不等式と $\sum(-1)^kc_k=\chi(T^2)$ を確かめよ。\\
(3) 定理 \ref{thm:ch11-handle} によるトーラスの胞体分解を述べよ。
\end{problem}

\begin{solution}
(1) $\partial_\theta h=-(R+r\cos\phi)\sin\theta$、$\partial_\phi h=-r\sin\phi\cos\theta$。$R+r\cos\phi>0$ なので $\sin\theta=0$、すなわち $\theta=0,\pi$。このとき $\cos\theta=\pm1\neq0$ なので $\sin\phi=0$、$\phi=0,\pi$。臨界点は4つ。2階微分は
\[
\partial_\theta^2h=-(R+r\cos\phi)\cos\theta,\quad \partial_\phi^2h=-r\cos\phi\cos\theta,\quad \partial_\theta\partial_\phi h=r\sin\phi\sin\theta=0\ (\text{臨界点で}).
\]
\begin{center}
\begin{tabular}{c|c|cc|c}
$(\theta,\phi)$ & $h$ & $\partial_\theta^2h$ & $\partial_\phi^2h$ & 指数\\ \hline
$(\pi,0)$ & $-(R+r)$ & $R+r$ & $r$ & 0\\
$(\pi,\pi)$ & $-(R-r)$ & $R-r$ & $-r$ & 1\\
$(0,\pi)$ & $R-r$ & $-(R-r)$ & $r$ & 1\\
$(0,0)$ & $R+r$ & $-(R+r)$ & $-r$ & 2
\end{tabular}
\end{center}
いずれもヘッセ行列は対角で正則なので $h$ はモース関数であり、$R=\tfrac32$、$r=\tfrac12$ とすれば例の図の4点（高さ $-2,-1,1,2$）に一致する。

(2) $(c_0,c_1,c_2)=(1,2,1)$、$(b_0,b_1,b_2)=(1,2,1)$（第4章）。弱不等式 $c_k\ge b_k$ は等号で成り立つ。強不等式は $k=0$：$1\ge1$、$k=1$：$c_1-c_0=1\ge b_1-b_0=1$、$k=2$：$c_2-c_1+c_0=0=b_2-b_1+b_0$。よって $\chi(T^2)=1-2+1=0$ で、$h$ は完全モース関数である。

(3) 下から順に：極小で $e^0$（円板 $\simeq$ 点）、第1の鞍点で $e^1$ を付けて $S^1$ と同じ型（円筒）、第2の鞍点で $e^1$ を付けて $S^1\vee S^1$ の型（穴のあいたトーラス）、極大で $e^2$ を付けて閉じる。$T^2\simeq e^0\cup e^1\cup e^1\cup e^2$ で、2胞体の接着写像は交換子 $aba^{-1}b^{-1}$ である（第2, 4章）。
\end{solution}

\begin{problem}[難易度 B]\label{ch11-prob-cpn}
(1) $S^n\subset\R^{n+1}$ 上の高さ関数 $h=x_{n+1}$ の臨界点は北極・南極の2点で、指数はそれぞれ $n$, $0$ であることを局所座標で示せ。$n\ge2$ のとき、モース不等式から $S^n$ のベッチ数を求めよ。$n=1$ ではどうか。\\
(2) 実数 $\lambda_0<\lambda_1<\dots<\lambda_n$ に対し、$\CP^n$ 上の関数
\[
f([z])=\frac{\sum_{j=0}^n\lambda_j\lvert z_j\rvert^2}{\sum_{j=0}^n\lvert z_j\rvert^2}
\]
を考える。量子力学の言葉では、ハミルトニアン $\hat H=\operatorname{diag}(\lambda_0,\dots,\lambda_n)$ をもつ $(n+1)$ 準位系の、状態 $[z]$（射線）でのエネルギー期待値である。$f$ の臨界点が固有状態 $[e_j]$ であり、その指数が $2j$ であることを示せ。\\
(3) $\CP^n$ のベッチ数（$\Z$ 係数ホモロジー）を求め、$\chi(\CP^n)$ を計算せよ。
\end{problem}

\begin{solution}
(1) 南極 $S=(0,\dots,0,-1)$ の近くでは $(x_1,\dots,x_n)$ が局所座標で $h=-\sqrt{1-\lvert x\rvert^2}=-1+\tfrac12\lvert x\rvert^2+O(\lvert x\rvert^4)$：指数 $0$ の非退化臨界点。北極では $h=1-\tfrac12\lvert x\rvert^2+\dots$：指数 $n$。それ以外の点では、$dh$ は $S^n$ への $e_{n+1}$ の射影 $e_{n+1}-x_{n+1}x$ に対応し、これは $x\neq\pm e_{n+1}$ で $0$ でないので臨界点ではない。

$n\ge2$ では $c_0=c_n=1$、他は $0$ で、隣り合う指数がないので完全（欠落原理）：$b_0=b_n=1$、他は $0$。直接には、弱不等式で $b_k=0$（$k\neq0,n$）、$b_0=1$（連結）、$\chi=1+(-1)^n=b_0+(-1)^nb_n$ から $b_n=1$。
$n=1$ では指数 $0,1$ が隣り合うので不等式だけでは $b_1\in\{0,1\}$ の可能性が残るが、$\chi=c_0-c_1=0=b_0-b_1$ より $b_1=1$。

(2) $[e_j]$ の近傍では $z_j\neq0$ なので $z_j=1$ と規格化し、$w_k=z_k/z_j$（$k\neq j$）を複素座標にとる。
\[
f=\frac{\lambda_j+\sum_{k\neq j}\lambda_k\lvert w_k\rvert^2}{1+\sum_{k\neq j}\lvert w_k\rvert^2}=\lambda_j+\sum_{k\neq j}(\lambda_k-\lambda_j)\lvert w_k\rvert^2+O(\lvert w\rvert^4).
\]
よって $w=0$ は臨界点で、ヘッセ行列は実座標 $(\mathrm{Re}\,w_k,\mathrm{Im}\,w_k)$ について対角 $2(\lambda_k-\lambda_j)$（各2重）。$\lambda_k-\lambda_j<0$ となるのは $k<j$ の $j$ 個で、各々が実2次元を与えるので指数は $2j$。
他に臨界点がないこと：射線上の変分 $\delta f=0$ は、$\lvert z\rvert=1$ に規格化し $\delta z\perp z$ をとると $\delta f=2\,\mathrm{Re}\langle\delta z,(\hat H-f)z\rangle$ なので、すべての $\delta z\perp z$ について $0$ になることは $(\hat H-f)z\parallel z$、すなわち $\hat Hz=fz$（固有ベクトル）と同値。固有値が相異なるので $[z]=[e_j]$。これは「エネルギー期待値の停留点は固有状態」というレイリー・リッツの変分原理であり、指数 $2j$ は「より低い固有状態の数 $\times2$」（その方向へ混ぜるとエネルギーが下がる）である。

(3) 指数は $0,2,\dots,2n$ でどれも隣り合わないので完全：$b_{2j}=1$（$0\le j\le n$）、奇数次は $0$。モース・ホモロジーで見ても $C_k$ は奇数次で $0$ なので $\partial=0$、$H_{2j}(\CP^n;\Z)=\Z$、$H_{\text{odd}}=0$ が $\Z$ 係数で成り立つ。$\chi(\CP^n)=n+1$。$n=1$ で $\CP^1\cong S^2$ と (1) に一致する。
\end{solution}

\begin{problem}[難易度 C]\label{ch11-prob-rpn}
$\lambda_0<\lambda_1<\dots<\lambda_n$ とし、$S^n$ 上の関数 $f(x)=\sum_{i=0}^n\lambda_ix_i^2$（$x=(x_0,\dots,x_n)$）を考える。$f(-x)=f(x)$ なので $f$ は $\RP^n$ 上の関数 $\bar f$ を定める。\\
(1) $f$ の臨界点は $\pm e_k$ で指数 $k$ であること、$\bar f$ の臨界点は $[e_k]$（$k=0,\dots,n$）で指数 $k$ であることを示せ。\\
(2) 丸い計量に関する勾配流 $\dot x=-\nabla f$ の解が $x(\tau)=e^{-2\Lambda\tau}c/\lvert e^{-2\Lambda\tau}c\rvert$（$\Lambda=\operatorname{diag}(\lambda_i)$、$c\in\R^{n+1}\setminus0$）であることを示し、$e_k$ から $\pm e_{k-1}$ への流線をすべて求めよ。$\RP^n$ では $[e_k]$ から $[e_{k-1}]$ への流線は何本か。\\
(3) 2本の流線の符号の比は $(-1)^k$ である（$[e_k]$ の不安定多様体の閉包は $\RP^k$ で、その接着写像 $S^{k-1}\to\RP^{k-1}$ の、$[e_{k-1}]$ の2つの逆像 $\pm e_{k-1}$ での局所写像度が対蹠写像の写像度 $(-1)^k$ だけ異なる）ことを認めて、モース複体 $\partial[e_k]=(1+(-1)^k)[e_{k-1}]$ を導き、$H_*(\RP^n;\Z)$ と $H_*(\RP^n;\Z_2)$ を求めよ。$\Z_2$ 係数では $\bar f$ が完全であることを確かめよ。\\
(4) $SO(3)\cong\RP^3$ と、ネマティック液晶の秩序変数空間 $\RP^2$ について、(3) の結果を述べよ。
\end{problem}

\begin{solution}
(1) ラグランジュ乗数法：$\lvert x\rvert^2=1$ の下で $f$ の停留点は $\Lambda x=\mu x$、すなわち固有ベクトル $x=\pm e_k$（固有値が相異なる）。$e_k$ の近くで $x_k=\sqrt{1-\sum_{i\neq k}x_i^2}$ とし $(x_i)_{i\neq k}$ を座標にとると
\[
f=\lambda_k+\sum_{i\neq k}(\lambda_i-\lambda_k)x_i^2+O(\lvert x\rvert^4).
\]
負の係数は $i<k$ の $k$ 個なので指数 $k$。$-e_k$ でも同じ。$\RP^n$ への射影 $S^n\to\RP^n$ は局所微分同相なので、$\bar f$ の臨界点は $[e_k]$ の $n+1$ 個で指数 $k$。

(2) $f$ の $\R^{n+1}$ での勾配は $2\Lambda x$ で、球面への射影は $\nabla f=2\Lambda x-2(x\cdot\Lambda x)x=2\Lambda x-2fx$。$y=e^{-2\Lambda\tau}c$ とおくと $\dot y=-2\Lambda y$ で、$x=y/\lvert y\rvert$ について
\[
\dot x=\frac{\dot y}{\lvert y\rvert}-\frac{y\,(y\cdot\dot y)}{\lvert y\rvert^3}=-2\Lambda x+2(x\cdot\Lambda x)x=-\nabla f .
\]
初期値は任意にとれるので、これがすべての解である。成分は $x_i\propto c_ie^{-2\lambda_i\tau}$ なので、$\tau\to-\infty$ では $c_i\neq0$ となる最大の $i$ の成分が、$\tau\to+\infty$ では最小の $i$ の成分が支配する。$\tau\to-\infty$ で $e_k$ に近づくには「$i>k$ で $c_i=0$、$c_k>0$」、$\tau\to+\infty$ で $\pm e_{k-1}$ に近づくには「$i<k-1$ で $c_i=0$、$\pm c_{k-1}>0$」が必要十分。よって $c=(0,\dots,0,c_{k-1},c_k,0,\dots)$ で、$c$ の正の定数倍と $\tau$ の平行移動は同じ曲線を与えるので、流線は
\[
\gamma_\pm:\ x=\cos\varphi\,e_k\pm\sin\varphi\,e_{k-1}\quad(\varphi:0\to\tfrac\pi2)
\]
の2本（$e_k$ から $e_{k-1}$ へ1本、$-e_{k-1}$ へ1本）。この大円の弧上で $f=\lambda_k\cos^2\varphi+\lambda_{k-1}\sin^2\varphi$ は単調減少している。$\RP^n$ では、$\gamma_+$ と $\gamma_-$ は $[e_k]$ から $[e_{k-1}]$ への異なる2本の流線になる（$\gamma_-\neq\pm\gamma_+$ なので射影は異なる。$-e_k$ から出る流線 $-\gamma_\pm$ は $\RP^n$ では $\gamma_\pm$ と同じ曲線に射影されるので、新たには数えない）。よって $\RP^n$ では2本。

(3) 2本の符号の比が $(-1)^k$ なので、全体の符号を適当に選んで
\[
\partial[e_k]=\bigl(1+(-1)^k\bigr)[e_{k-1}]=\begin{cases}2[e_{k-1}]&(k\ \text{偶数}\ge2)\\0&(k\ \text{奇数})\end{cases},\qquad \partial[e_0]=0 .
\]
$\Z$ 係数：$C_k=\Z$（$0\le k\le n$）。$k$ が奇数なら $\partial_k=0$、偶数 $k\ge2$ なら $\partial_k=\times2$。
\begin{itemize}
\item $H_0=\Z/\operatorname{im}\partial_1=\Z$。
\item $0<k<n$、$k$ 奇数：$\ker\partial_k=\Z$、$\operatorname{im}\partial_{k+1}=2\Z$ なので $H_k=\Z_2$。
\item $0<k<n$、$k$ 偶数：$\ker\partial_k=0$ なので $H_k=0$。
\item $k=n$：$n$ 奇数なら $\ker\partial_n=\Z$ で上からの像はないので $H_n=\Z$；$n$ 偶数なら $H_n=0$。
\end{itemize}
$\Z_2$ 係数では $2=0$ なので $\partial=0$、$H_k(\RP^n;\Z_2)=\Z_2$（$0\le k\le n$）。よって $b_k(\Z_2)=1=c_k$ で $\bar f$ は $\Z_2$ 係数で完全。$\Q$ 係数では $b_0=1$、$b_n=1$（$n$ 奇数）以外は $0$ で、完全ではない（流線が「2本で打ち消さずに足し合う」ことが $\Z_2$ ねじれの原因）。$\chi(\RP^n)=\sum(-1)^kc_k=(1+(-1)^n)/2$ もわかる。

(4) $SO(3)\cong\RP^3$：$H_0=\Z$、$H_1=\Z_2$、$H_2=0$、$H_3=\Z$。$H_1=\Z_2$ は $\pi_1(SO(3))=\Z_2$ のアーベル化であり、$2\pi$ 回転が $-1$ を与えるスピン $\tfrac12$ の起源（第2章）。$\RP^2$：$H_0=\Z$、$H_1=\Z_2$、$H_2=0$；$\pi_1(\RP^2)=\Z_2$ はネマティック液晶の線欠陥（$\tfrac12$ 回位）が $\Z_2$ 電荷をもつこと（2本が合体すると消える）に対応する（第3章）。
\end{solution}

\begin{problem}[難易度 B]\label{ch11-prob-ineq}
定理 \ref{thm:ch11-ineq}(3) の多項式形 $M_t-P_t=(1+t)R(t)$（$R$ は非負整数係数）を認める。\\
(1) $t^k$ の係数を比べて弱モース不等式を、$(1+t)^{-1}=\sum(-t)^j$ を用いて強モース不等式を導け。$t=-1$ を代入して $\sum(-1)^kc_k=\chi(M)$ を示せ。\\
(2)（欠落原理）すべての $k$ で $c_kc_{k+1}=0$ ならば $h$ は完全である（$c_k=b_k$）ことを示せ。\\
(3) 種数 $g$ の閉曲面 $\Sigma_g$ 上のモース関数は少なくとも $2g+2$ 個の臨界点をもつことを示せ。\\
(4) $n$ 次元閉多様体上で、$-h$ の指数 $k$ の臨界点は $h$ の指数 $n-k$ の臨界点であることを用いて、$n$ が奇数なら $\chi(M)=0$ であることを示せ。
\end{problem}

\begin{solution}
$R(t)=\sum_kr_kt^k$（$r_k\ge0$、$r_{-1}=r_n=0$）とする。

(1) $t^k$ の係数：$c_k-b_k=r_k+r_{k-1}\ge0$。これが弱不等式。また $(M_t-P_t)/(1+t)=R(t)$ を形式的べき級数で展開すると
\[
r_k=\sum_{j=0}^k(-1)^{k-j}(c_j-b_j)\ge0 ,
\]
これが強不等式。さらに $R$ は多項式で $\deg R\le n-1$ なので $r_n=0$、すなわち $k=n$ で等号。$t=-1$ を代入すると右辺は $0$ で $\sum(-1)^kc_k=\sum(-1)^kb_k=\chi(M)$。

(2) $R\neq0$ とし、$r_j>0$ となる $j$ をとる。$t^j$ の係数より $c_j\ge b_j+r_j>0$、$t^{j+1}$ の係数より $c_{j+1}\ge b_{j+1}+r_j>0$。これは $c_jc_{j+1}=0$ に反する。よって $R=0$、$M_t=P_t$。特に指数がすべて偶数なら完全である（$\CP^n$、問題 \ref{ch11-prob-cpn}）。

(3) $b_0=b_2=1$、$b_1=2g$ なので、弱不等式より $c_0+c_1+c_2\ge b_0+b_1+b_2=2g+2$。立てた $\Sigma_g$（$g$ 個の穴が縦に並んだ面）の高さ関数は極小1、鞍点 $2g$、極大1で、この下限が実現される。

(4) $-h$ のヘッセ行列は $h$ のヘッセ行列の $-1$ 倍なので、負の固有値の数は $n-\ind p$。よって $c_k(-h)=c_{n-k}(h)$ で
\[
\chi(M)=\sum_k(-1)^kc_k(-h)=\sum_k(-1)^kc_{n-k}(h)=(-1)^n\sum_j(-1)^jc_j(h)=(-1)^n\chi(M).
\]
$n$ 奇数なら $\chi=-\chi$、$\chi(M)=0$。（ポアンカレ双対 $b_k=b_{n-k}$ の、モース理論による表れである。）
\end{solution}

\begin{problem}[難易度 B]\label{ch11-prob-sphaleron}
$M$ をコンパクトで連結な多様体、$h$ をモース関数とする。\\
(1) $h$ が2つ以上の極小（指数 $0$ の臨界点）をもつならば、指数 $1$ の臨界点が存在すること、より詳しく $c_1\ge c_0-1+b_1$ を、強モース不等式から示せ。\\
(2) $h$ の極小がただ1つでも、$b_1(M)\ge1$ ならば指数 $1$ の臨界点があることを示せ。\\
(3) 場の理論で $h$ をエネルギー汎関数、$M$ を配位空間とみなす。電弱理論では、チャーン・サイモンズ数の異なる古典的真空どうしがエネルギー障壁で隔てられており、その最低点にある不安定な静的解を\term{スファレロン}{sphaleron}という（クリンカマー・マントン）。(1)(2) が示唆するスファレロンの性質（指数）を述べ、「山越えの補題（最小の山越え経路は指数 $1$ の鞍点を通る）」との関係を説明せよ。（配位空間は無限次元なので、ここでは類推として述べればよい。）
\end{problem}

\begin{solution}
(1) 強不等式の $k=1$：$c_1-c_0\ge b_1-b_0$。連結なので $b_0=1$。よって $c_1\ge c_0-1+b_1\ge c_0-1$。$c_0\ge2$ なら $c_1\ge1$。
直観的には、2つの谷 $A,B$ を結ぶ道の最高点をできるだけ低くしたもの（峠）が指数1の鞍点である。$h$ の部分レベル集合 $M^a$ で $A,B$ が別の連結成分にある $a$ から、同じ成分になる $a$ へ移るところで、定理 \ref{thm:ch11-handle} により 1 胞体が付かなければならない（$0$ 胞体は成分を増やし、$k\ge2$ 胞体は成分をつながない）。

(2) $c_0=1$ でも $c_1\ge b_1\ge1$（弱不等式）。モース・ホモロジーで言えば、$H_1\neq0$ の生成元を代表するには1次の鎖、すなわち指数1の臨界点が要る。マントンの「ほどけないループ」の議論（配位空間の非自明なループの上でエネルギーを最大化し、それを最小化すると鞍点が得られる）はこの類推である。

(3) チャーン・サイモンズ数 $N_{CS}=n$ と $n+1$ の真空（$h$ の極小）は異なり、両者は連続な配位の道で結べるが、途中で必ずエネルギー障壁を越える。(1) の類推から、両者の間に指数1の臨界点があり、それがスファレロンである：エネルギー汎関数のヘッセ行列（揺らぎ演算子）がちょうど1つの負の固有値をもち、その方向が「真空 $n$ へ転がり落ちる方向／真空 $n+1$ へ転がり落ちる方向」を与える。最小山越え経路（2つの真空を結ぶ道の中で最高エネルギーが最小のもの）の最高点が指数1の鞍点であることは (1) の直観的説明そのものであり、スファレロンのエネルギー $E_{\mathrm{sph}}$（電弱理論で約 $9\ \mathrm{TeV}$）は、高温で $B+L$ を破る遷移率 $\sim e^{-E_{\mathrm{sph}}/T}$ を決める。
\end{solution}

\subsection{超対称量子力学とウィッテン指数}

\begin{problem}[難易度 A]\label{ch11-prob-sho}
ボソンの消滅演算子 $a$（$[a,a^\dagger]=1$）とフェルミオンの消滅演算子 $b$（$\{b,b^\dagger\}=1$、$b^2=0$）、$\omega>0$ に対し
\[
Q=\sqrt\omega\,a\,b^\dagger,\qquad F=b^\dagger b
\]
とする（\term{超対称調和振動子}{supersymmetric harmonic oscillator}）。\\
(1) $Q^2=0$ と $H=\{Q,Q^\dagger\}=\omega(a^\dagger a+b^\dagger b)$ を示せ。\\
(2) $H$ のスペクトルと縮退度を求め、$E>0$ の各準位でボソン状態とフェルミオン状態が $Q$ で対になることを確かめよ。\\
(3) $\Tr(-1)^Fe^{-\beta H}$ を直接計算し、$\beta$ によらず $1$ であることを示せ。
\end{problem}

\begin{solution}
(1) $Q^2=\omega\,a^2(b^\dagger)^2=0$（$a$ と $b^\dagger$ は可換、$(b^\dagger)^2=0$）。$Q^\dagger=\sqrt\omega\,a^\dagger b$ で
\[
\begin{aligned}
\{Q,Q^\dagger\}&=\omega\bigl(ab^\dagger a^\dagger b+a^\dagger b\,ab^\dagger\bigr)=\omega\bigl(aa^\dagger\,b^\dagger b+a^\dagger a\,bb^\dagger\bigr)\\
&=\omega\bigl((a^\dagger a+1)b^\dagger b+a^\dagger a(1-b^\dagger b)\bigr)=\omega(a^\dagger a+b^\dagger b).
\end{aligned}
\]

(2) 基底 $\lvert n,n_F\rangle$（$n=0,1,\dots$、$n_F=0,1$）で $E=\omega(n+n_F)$。$E=0$ は $\lvert0,0\rangle$（ボソン）のみ。$E=m\omega$（$m\ge1$）は $\lvert m,0\rangle$（ボソン）と $\lvert m-1,1\rangle$（フェルミオン）の2重縮退。$Q\lvert m,0\rangle=\sqrt\omega\,a\lvert m\rangle\otimes b^\dagger\lvert0\rangle=\sqrt{m\omega}\,\lvert m-1,1\rangle$、$Q^\dagger\lvert m-1,1\rangle=\sqrt{m\omega}\,\lvert m,0\rangle$ で、両者は $Q$ で対になる。

(3)
\[
\Tr(-1)^Fe^{-\beta H}=\sum_{n\ge0}e^{-\beta\omega n}-\sum_{n\ge0}e^{-\beta\omega(n+1)}=\frac{1-e^{-\beta\omega}}{1-e^{-\beta\omega}}=1 .
\]
$E>0$ の寄与は各準位で $e^{-\beta m\omega}-e^{-\beta m\omega}=0$ と打ち消し、$E=0$ のボソン1個だけが残る。座標表示では $a=(\sqrt\omega x+ip/\sqrt\omega)/\sqrt2$ で $H=\tfrac12(p^2+\omega^2x^2)+\tfrac\omega2[b^\dagger,b]$ となり、超ポテンシャル $W=\tfrac12\omega x^2$ の模型（問題 \ref{ch11-prob-W} の規格化の $\tfrac12$ 倍）である。$\omega<0$ にあたる $W=-\tfrac12\lvert\omega\rvert x^2$ では基底状態がフェルミオン的になり、指数は $-1$ となる。
\end{solution}

\begin{problem}[難易度 B]\label{ch11-prob-pair}
定義 \ref{def:ch11-sqm} の超対称量子力学を考える。\\
(1) $\langle\psi|H|\psi\rangle=\lVert Q\psi\rVert^2+\lVert Q^\dagger\psi\rVert^2$ を示し、$H\psi=0\iff Q\psi=Q^\dagger\psi=0$ を導け。\\
(2) $E>0$ の固有空間 $\mathcal H_E$ で $Q_1=(Q+Q^\dagger)/\sqrt E$ が $Q_1^2=1$ をみたし、$(-1)^F$ と反可換であることを示し、$\dim(\mathcal H_E\cap\mathcal H_B)=\dim(\mathcal H_E\cap\mathcal H_F)$ を結論せよ。これより $\Tr(-1)^Fe^{-\beta H}=n_B^0-n_F^0$ を示せ。\\
(3) 各固有空間が有限次元で、$\mathcal H$ が $\ker H\oplus\operatorname{im}Q\oplus\operatorname{im}Q^\dagger$ と直交分解されることを認めて、$\ker H\cong\ker Q/\operatorname{im}Q$（$Q$ コホモロジー）を示せ。\\
(4) $W$ をボソン的なエルミート演算子とし $Q_t=e^{-tW}Qe^{tW}$ とする。$Q_t^2=0$ を示し、$Q_t$ コホモロジーが $Q$ コホモロジーと同型であることから、（(3) の仮定のもとで）$H_t=\{Q_t,Q_t^\dagger\}$ の $E=0$ 状態の数がボソン・フェルミオン別に $t$ によらないことを示せ。$H_t$ のスペクトル自体は $t$ によって変わることに注意せよ。
\end{problem}

\begin{solution}
(1) $\langle\psi|QQ^\dagger+Q^\dagger Q|\psi\rangle=\langle Q^\dagger\psi|Q^\dagger\psi\rangle+\langle Q\psi|Q\psi\rangle$。右辺は非負で、$0$ になるのは $Q\psi=Q^\dagger\psi=0$ のとき。逆にそのとき $H\psi=0$ は明らか。

(2) $Q^2=(Q^\dagger)^2=0$ より $(Q+Q^\dagger)^2=QQ^\dagger+Q^\dagger Q=H$。$[H,Q]=[H,Q^\dagger]=0$ なので $Q_1$ は $\mathcal H_E$ を保ち、そこで $Q_1^2=H/E=1$。$(-1)^F$ と $Q$, $Q^\dagger$ は反可換なので $Q_1$ も反可換で、$Q_1$ は $\mathcal H_E\cap\mathcal H_B$ を $\mathcal H_E\cap\mathcal H_F$ に写し、その逆写像も $Q_1$ である。よって両者は同型で次元が等しい。したがって
\[
\Tr(-1)^Fe^{-\beta H}=\sum_Ee^{-\beta E}\bigl(\dim\mathcal H_E^B-\dim\mathcal H_E^F\bigr)=n_B^0-n_F^0 .
\]

(3) $\ker H\subset\ker Q$ は (1) より。$\operatorname{im}Q\subset\ker Q$（$Q^2=0$）。$\ker Q\cap\operatorname{im}Q^\dagger=0$：$Q Q^\dagger\chi=0$ なら $\lVert Q^\dagger\chi\rVert^2=\langle\chi|QQ^\dagger\chi\rangle=0$。よって直交分解から $\ker Q=\ker H\oplus\operatorname{im}Q$ となり、$\ker Q/\operatorname{im}Q\cong\ker H$。$Q$ が $F$ を保つ次数（$F\to F\pm1$）をもつ場合は各次数ごとに成り立つ。これはホッジ分解（第5章）の抽象版である。

(4) $Q_t^2=e^{-tW}Q^2e^{tW}=0$。$e^{tW}$ は可逆で $(-1)^F$ と可換なので、$\psi\mapsto e^{-tW}\psi$ は $\ker Q\to\ker Q_t$、$\operatorname{im}Q\to\operatorname{im}Q_t$ の同型を与え、$Q$ コホモロジーと $Q_t$ コホモロジーは次数ごとに同型。$H_t=\{Q_t,Q_t^\dagger\}$ に (3) を適用すると、$\ker H_t$ の次元は次数ごとに $t$ によらない。ただし $e^{tW}$ はユニタリでないので $H_t$ は $H$ とユニタリ同値ではなく、$E>0$ のスペクトルや $E=0$ 状態の波動関数は $t$ とともに変わる（ウィッテン変形で $t\to\infty$ とすると波動関数が臨界点に局在するのはこのため）。非コンパクトな空間では $e^{tW}$ が $L^2$ を保たず結論が破れうる（問題 \ref{ch11-prob-W}）。
\end{solution}

\begin{problem}[難易度 B]\label{ch11-prob-torusQM}
平坦なトーラス $T^n=\R^n/2\pi\Z^n$ 上の超対称量子力学（$Q=d$、$H=\Delta$、$F=$ 形式の次数）を考える。\\
(1) $n=1$（$S^1$）のとき、$\Omega^0$ と $\Omega^1$ での $H$ の固有関数と固有値を求め、$Q$ による対を具体的に示せ。$E=0$ の状態を求め、ウィッテン指数が $\chi(S^1)=0$ であることを確かめよ。\\
(2) 一般の $n$ で、$H$ の固有値と各次数の縮退度を求め、$\Tr(-1)^Fe^{-\beta H}=0=\chi(T^n)$ を示せ。$E=0$ の状態の数を次数ごとに求め、ベッチ数 $b_p(T^n)=\binom np$ と比べよ。\\
(3) $S^2$ 上の超対称量子力学のウィッテン指数はいくつか。$S^2$ 上で $E=0$ の状態をすべて書け。
\end{problem}

\begin{solution}
(1) $\Omega^0$：$f_m=e^{im\theta}$、$H=d^\dagger d=-\partial_\theta^2$ で $E=m^2$。$\Omega^1$：$g_m=e^{im\theta}d\theta$、$H=dd^\dagger=-\partial_\theta^2$ で $E=m^2$（$d^\dagger(g\,d\theta)=-g'$）。$Q f_m=d f_m=im\,e^{im\theta}d\theta=im\,g_m$ なので、$m\neq0$ では $f_m\leftrightarrow g_m$ が対になる（$Q^\dagger g_m=-im f_m$）。$E=0$ は $f_0=1$（ボソン）と $g_0=d\theta$（フェルミオン）。
\[
\Tr(-1)^Fe^{-\beta H}=\sum_{m\in\Z}e^{-\beta m^2}(1-1)=0=\chi(S^1).
\]
基底状態は2つあるが統計が逆なので指数は $0$。指数 $0$ は超対称性の破れを意味しない（実際破れていない）ことに注意。

(2) 平坦計量では $\Delta(f\,d\theta^{i_1}\wedge\dots\wedge d\theta^{i_p})=(-\sum_j\partial_j^2f)\,d\theta^{i_1}\wedge\dots\wedge d\theta^{i_p}$。固有形式は $e^{i\bm m\cdot\bm\theta}d\theta^I$（$\bm m\in\Z^n$、$\lvert I\rvert=p$）で $E=\lvert\bm m\rvert^2$。エネルギー $E$ の $p$ 形式の数は $N(E)\binom np$（$N(E)=\#\{\bm m:\lvert\bm m\rvert^2=E\}$）なので
\[
\Tr(-1)^Fe^{-\beta H}=\sum_{\bm m}e^{-\beta\lvert\bm m\rvert^2}\sum_{p=0}^n(-1)^p\binom np=\Bigl(\sum_{\bm m}e^{-\beta\lvert\bm m\rvert^2}\Bigr)(1-1)^n=0 .
\]
$E=0$ は $\bm m=0$ の定数係数形式 $d\theta^I$ で、$p$ 形式は $\binom np$ 個、合計 $2^n$ 個。これは $b_p(T^n)=\binom np$ に一致し（第5章）、$\chi(T^n)=(1-1)^n=0$。フェルミオン零モード $\psi^i$（$n$ 個）の量子化で $2^n$ 状態ができる、と理解することもできる。

(3) 例 \ref{ex:ch11-sqmM} より $\chi(S^2)=2$。$E=0$ 状態は調和形式で、$b_0=b_2=1$、$b_1=0$ より、定数関数 $1$（ボソン）と体積形式 $\mathrm{vol}_{S^2}$（$F=2$ でボソン）の2つ。指数 $2-0=2\neq0$ なので、どんな変形（計量を変える、超ポテンシャルを加える）をしても超対称性は破れない。
\end{solution}

\begin{problem}[難易度 B]\label{ch11-prob-W}
1次元の超対称量子力学を、$\R$ 上の $0$ 形式 $f(x)$（ボソン）と $1$ 形式 $g(x)\,dx$（フェルミオン）の空間で $Q=d_W=e^{-W}de^{W}$ により定める（$W(x)$ は実多項式、\term{超ポテンシャル}{superpotential}）。\\
(1) $H=\{Q,Q^\dagger\}$ が $0$ 形式上 $-\partial_x^2+W'^2-W''$、$1$ 形式上 $-\partial_x^2+W'^2+W''$ であることを示せ。\\
(2) $E=0$ の状態を求め、ウィッテン指数が次で与えられることを示せ：
\[
I_W=\begin{cases}+1&(W\to+\infty\ (x\to\pm\infty))\\-1&(W\to-\infty\ (x\to\pm\infty))\\0&(\text{それ以外：}x\to\pm\infty\text{ で符号が逆})\end{cases}
\]
特に $W$ の次数を $N$、最高次係数を $w_N$ とすると $I_W$ は $N$ 偶数なら $\operatorname{sign}w_N$、奇数なら $0$。\\
(3) $W=\tfrac14x^4-\tfrac12a^2x^2$（$a>0$）で、調和近似による摂動的真空を数え、(2) の厳密な結果と比べよ。\\
(4) $W_\epsilon=\epsilon x^3+x^2$ の指数を $\epsilon=0$ と $\epsilon\neq0$ で比べ、「指数は連続変形で不変」と矛盾しない理由を説明せよ。
\end{problem}

\begin{solution}
(1) $Qf=e^{-W}(e^Wf)'dx=(f'+W'f)dx$、$Q^\dagger(g\,dx)=e^{W}d^\dagger(e^{-W}g\,dx)=-e^W(e^{-W}g)'=-g'+W'g$。$0$ 形式上 $H=Q^\dagger Q$：
\[
Q^\dagger Qf=-(f'+W'f)'+W'(f'+W'f)=-f''-W''f+W'^2f .
\]
$1$ 形式上 $H=QQ^\dagger$：$(\partial+W')(-g'+W'g)=-g''+W''g+W'g'-W'g'+W'^2g=-g''+W'^2g+W''g$。

(2) 問題 \ref{ch11-prob-pair}(1) より $E=0\iff Q\psi=Q^\dagger\psi=0$。$0$ 形式では $Q^\dagger f=0$ は自動的で、$Qf=0$ より $f=Ce^{-W}$。$1$ 形式では $Q(g\,dx)=0$ は自動的で、$Q^\dagger(g\,dx)=0$ より $g=Ce^{W}$。$e^{-W}$ が $L^2$ に入るのは $W\to+\infty$（両側）のとき、$e^W$ が $L^2$ に入るのは $W\to-\infty$（両側）のときで、両方同時には起こらない。多項式なので $\lvert W\rvert\to\infty$ は少なくとも1次で起こり、指数関数的に減衰するので確かに $L^2$。よって $I_W=+1,-1,0$ が表のとおり。$N$ 偶数では $x\to\pm\infty$ で $W\sim w_Nx^N$ の符号が両側同じで $\operatorname{sign}w_N$、奇数では逆符号で $I_W=0$。

(3) $W'=x^3-a^2x=x(x-a)(x+a)$、$W''=3x^2-a^2$。臨界点 $x=\pm a$ で $W''=2a^2>0$（$W$ の極小）、$x=0$ で $W''=-a^2<0$（極大）。臨界点 $x_0$ の近くで $W'^2\approx W''(x_0)^2(x-x_0)^2$ なので、$\omega=\lvert W''(x_0)\rvert$ として $-\partial^2+\omega^2y^2$ の固有値は $\omega(2n+1)$。
\begin{itemize}
\item $W''>0$：$0$ 形式で $\omega(2n+1)-\omega=2n\omega$ → $n=0$ で $E=0$（ボソン）。$1$ 形式は $2(n+1)\omega>0$。
\item $W''<0$：$1$ 形式で $\omega(2n+1)-\omega=2n\omega$ → $E=0$（フェルミオン）。
\end{itemize}
よって摂動的真空はボソン2個（$x=\pm a$）、フェルミオン1個（$x=0$）。指数は $2-1=1$ で、(2) の厳密な値 $I_W=+1$（$N=4$、$w_4>0$）と一致する。厳密には $E=0$ 状態は $e^{-W}$ のただ1つ（両井戸に広がったボソン状態）で、残りのボソン1個とフェルミオン1個は、$0$ から $\pm a$ への勾配流線（インスタントン）を介して混ざり、$E>0$ の対に持ち上がる。モース複体で言えば $W$ の極大 $0$（指数1）から極小 $\pm a$ へ流線が1本ずつあり、$\partial(0)=(+a)-(-a)$、$H_0=\Z$、$H_1=0$（$\R$ のホモロジー）。

(4) $\epsilon=0$：$W=x^2$ は両側で $+\infty$、$I_W=1$（超対称調和振動子）。$\epsilon\neq0$：$N=3$ で $I_W=0$。$\epsilon\to0$ で、$W_\epsilon$ には $x\approx-2/(3\epsilon)$ に新しい臨界点が無限遠から現れる（$W'_\epsilon=3\epsilon x^2+2x=0$）。$\epsilon\neq0$ で摂動的真空は $x=0$（ボソン）と $x=-2/(3\epsilon)$（フェルミオン、$W''=-2$）の2つで、トンネル効果で持ち上がる。$\epsilon\to0$ でフェルミオン的真空が無限遠へ逃げ、ボソン的真空とのトンネル振幅が $0$ になるので、状態の対が「無限遠から出入り」して指数が変わる。指数の不変性は、漸近的振る舞い（$W$ の最高次）を変えない変形に対してのみ成り立つ。
\end{solution}

\begin{problem}[難易度 C]\label{ch11-prob-dw}
問題 \ref{ch11-prob-W} の設定で $W(x)=\lambda\bigl(\tfrac13x^3-a^2x\bigr)$（$\lambda,a>0$）とする。ボソン的ポテンシャル $W'^2=\lambda^2(x^2-a^2)^2$ は二重井戸である（第1節の図）。\\
(1) 調和近似で、$x=+a$ にボソン的な、$x=-a$ にフェルミオン的な $E=0$ の摂動的真空が1つずつあることを示し、ウィッテン指数を求めよ。\\
(2) 厳密な $E=0$ 状態が存在しないことを示せ。すなわち超対称性は自発的に破れる。\\
(3) $H$ の規格化 $H=p^2+W'^2\pm W''$ に対応するユークリッド作用 $S_E=\int d\tau\bigl(\tfrac14\dot x^2+W'^2\bigr)$（ボソン部分）について、$-a$ から $+a$ へ移る任意の道で $S_E\ge\lvert W(-a)-W(a)\rvert$ を示し、等号が $\dot x=\pm2W'$（勾配流）で成り立つことを示せ。$\Delta W=W(-a)-W(a)$ を計算せよ。\\
(4) 基底エネルギーが $E_0\sim e^{-2\Delta W}$（前因子を除く）のように振る舞うことを、「$\langle\psi_F\lvert Q\rvert\psi_B\rangle\sim e^{-\Delta W}$ と $E_0\approx\lvert\langle\psi_F|Q|\psi_B\rangle\rvert^2$」から説明し、これが $\lambda$ について摂動展開のどの次数にも現れない理由を述べよ。モース複体の言葉ではどう表されるか。
\end{problem}

\begin{solution}
(1) $W'=\lambda(x^2-a^2)$、$W''=2\lambda x$。$x=+a$：$W''=2\lambda a>0$、$x=-a$：$W''=-2\lambda a<0$。問題 \ref{ch11-prob-W}(3) と同じ調和近似（$\omega=2\lambda a$）で、$x=+a$ に $0$ 形式（ボソン）、$x=-a$ に $1$ 形式（フェルミオン）の $E=0$ 状態が1つずつ。指数は $1-1=0$。これは $W$ が3次（奇数次）であることと整合する。

(2) $E=0$ 状態の候補は $e^{-W}$ と $e^{W}dx$ だが、$W\sim\lambda x^3/3$ は $x\to+\infty$ で $+\infty$、$x\to-\infty$ で $-\infty$ なので、$e^{-W}$ は $x\to-\infty$ で、$e^W$ は $x\to+\infty$ で発散し、どちらも規格化できない。よって $E=0$ 状態はなく、基底エネルギー $E_0>0$：超対称性は破れる。指数 $0$ はこれと矛盾しない（指数 $\neq0$ なら破れないが、$0$ なら破れも破れないもありうる）。

(3) 平方完成：
\[
\tfrac14\dot x^2+W'^2=\Bigl(\tfrac12\dot x\mp W'\Bigr)^2\pm\dot xW'=\Bigl(\tfrac12\dot x\mp W'\Bigr)^2\pm\frac{dW}{d\tau}.
\]
積分して $S_E\ge\pm\bigl(W(x(+\infty))-W(x(-\infty))\bigr)$、符号を選んで $S_E\ge\lvert W(a)-W(-a)\rvert$。等号は $\dot x=\pm2W'$、すなわち $W$ の（時間を2倍に伸縮した）勾配流のとき。$(-a,a)$ では $W'<0$ なので、$-a$ から $a$ への流線は $\dot x=-2W'>0$（$W$ を減らす向き）。
\[
W(-a)=\lambda\Bigl(-\frac{a^3}3+a^3\Bigr)=\frac{2\lambda a^3}{3},\quad W(a)=-\frac{2\lambda a^3}3,\quad \Delta W=\frac{4\lambda a^3}{3}.
\]
これがインスタントン作用であり、インスタントンは $W$ の勾配流線（BPS 方程式 $\dot x=-2W'$ の解）である。

(4) 摂動的真空 $\psi_B$（$+a$ に局在）と $\psi_F$（$-a$ に局在）の空間に制限すると、$Q$ は $\psi_B\mapsto c\,\psi_F$ の形の $1\times1$ 行列で、低エネルギーの $H=\{Q,Q^\dagger\}$ の固有値は $\lvert c\rvert^2$。$c$ は $\psi_B\sim e^{-(W-W(a))}$ と $\psi_F\sim e^{W-W(-a)}$ の重なり（WKB の裾）で決まり、流線に沿った半古典評価で $c\sim e^{-\Delta W}$、したがって $E_0\sim e^{-2\Delta W}=e^{-8\lambda a^3/3}$。結合を $W\to W/g$ のように入れれば $E_0\sim e^{-8\lambda a^3/(3g)}$ であり、$g$ のべき級数展開ではすべての次数で $0$ になる。これは「摂動論の各次数で超対称性は破れない」（非繰り込み定理）と「非摂動的には破れる」ことの典型例である。
モース複体では、$W$ の極大 $-a$（指数1）から極小 $+a$（指数0）へ流線が1本だけあり（$-a$ から左へ出る流線は $W\to-\infty$ へ逃げて臨界点に着かない）、$\partial(-a)=\pm(+a)$。複体 $\Z\xrightarrow{\pm1}\Z$ のホモロジーは $0$ で、$E=0$ 状態がないこと（$\R$ 上で「無限遠を相対的に扱った」ホモロジーが $0$）に対応する。
\end{solution}

\begin{problem}[難易度 B]\label{ch11-prob-s1complex}
円周 $S^1$（反時計回りに向きをつける）上の関数 $h(\theta)=\cos2\theta$ を考える。臨界点は極大 $M_1$（$\theta=0$）、$M_2$（$\theta=\pi$）と極小 $m_1$（$\theta=\pi/2$）、$m_2$（$\theta=3\pi/2$）である（下図）。
\begin{center}
\begin{tikzpicture}[scale=1.25]
  \draw[thick,gray!60] (0,0) circle (1.5);
  \draw[very thick,blue!70!black,mid arrow] (5:1.5) arc (5:85:1.5);
  \draw[very thick,blue!70!black,mid arrow] (175:1.5) arc (175:95:1.5);
  \draw[very thick,blue!70!black,mid arrow] (185:1.5) arc (185:265:1.5);
  \draw[very thick,blue!70!black,mid arrow] (355:1.5) arc (355:275:1.5);
  \node[vtx,red!70!black] at (0:1.5) {}; \node[right] at (0:1.5) {$M_1$ ($\theta=0$)};
  \node[vtx,red!70!black] at (180:1.5) {}; \node[left] at (180:1.5) {$M_2$ ($\theta=\pi$)};
  \node[hvtx] at (90:1.5) {}; \node[above] at (90:1.5) {$m_1$ ($\theta=\pi/2$)};
  \node[hvtx] at (270:1.5) {}; \node[below] at (270:1.5) {$m_2$ ($\theta=3\pi/2$)};
  \draw[-{Stealth[length=2mm]},gray] (30:0.6) arc (30:150:0.6);
  \node[gray] at (0,0.1) {\footnotesize 向き};
  \node[blue!70!black] at (45:1.95) {$\gamma_1$};
  \node[blue!70!black] at (135:1.95) {$\gamma_2$};
  \node[blue!70!black] at (225:1.95) {$\gamma_3$};
  \node[blue!70!black] at (315:1.95) {$\gamma_4$};
\end{tikzpicture}
\zucap{$h=\cos2\theta$ の勾配流線（$-\nabla h$ の向き）。赤い点は極大（指数1）、白丸は極小（指数0）。}
\end{center}
(1) 4本の流線 $\gamma_1,\dots,\gamma_4$ の符号を、節の規則（$\ind p=1$ の点から、向きの正の側へ出る流線を $+1$、負の側へ出る流線を $-1$）で定め、モース複体の $\partial M_1$、$\partial M_2$ を求めよ。\\
(2) モース・ホモロジーを計算し $H_*(S^1;\Z)$ と一致することを確かめよ。$\Z_2$ 係数ではどうか。\\
(3) ウィッテン変形した超対称量子力学（$d_t=e^{-th}de^{th}$、$t\gg1$）で、摂動的真空はいくつあり、それぞれボソンかフェルミオンか。余境界 $\delta=\partial^T$ を用いて、厳密な $E=0$ 状態として生き残る組合せを求めよ。これらが厳密解 $e^{-th}$、$e^{th}d\theta$ の局在の様子と一致することを確かめよ。
\end{problem}

\begin{solution}
(1) $\gamma_1$：$M_1$ から正の向き（$\theta$ 増加）に出て $m_1$ へ、$+1$。$\gamma_4$：$M_1$ から負の向きに出て $m_2$ へ、$-1$。$\gamma_3$：$M_2$ から正の向きに出て $m_2$ へ、$+1$。$\gamma_2$：$M_2$ から負の向きに出て $m_1$ へ、$-1$。よって
\[
\partial M_1=m_1-m_2,\qquad \partial M_2=m_2-m_1,\qquad \partial m_i=0 .
\]
（$M_1$ の不安定多様体は弧 $(3\pi/2-2\pi,\pi/2)$ を正の向きに向きづけたもので、その境界が終点 $-$ 始点 $=m_1-m_2$ であることと一致する。）

(2) $C_1=\Z M_1\oplus\Z M_2$、$C_0=\Z m_1\oplus\Z m_2$、$\partial=\begin{pmatrix}1&-1\\-1&1\end{pmatrix}$（列が $M_1,M_2$）。$\ker\partial_1=\Z(M_1+M_2)$ で $H_1=\Z$（生成元は円周全体 $M_1+M_2$）。$\operatorname{im}\partial_1=\Z(m_1-m_2)$ で、$H_0=\Z^2/\Z(m_1-m_2)\cong\Z$（$m_1-m_2$ は直和因子を生成するので商はねじれなし）。$H_*(S^1)$ と一致。$\Z_2$ 係数では $\partial M_1=\partial M_2=m_1+m_2$、$\ker=\Z_2(M_1+M_2)$、$H_1=\Z_2$、$H_0=\Z_2^2/\Z_2=\Z_2$。完全ではない（$c=(2,2)$、$b=(1,1)$）が、$\sum(-1)^kc_k=0=\chi(S^1)$。

(3) 摂動的真空は4つ：$m_1,m_2$ に $0$ 形式（ボソン、$F=0$）、$M_1,M_2$ に $1$ 形式（フェルミオン、$F=1$）。指数は $2-2=0$。$\delta=\partial^T$：
\[
\delta m_1=M_1-M_2,\qquad \delta m_2=-M_1+M_2 .
\]
$\ker\delta|_{C^0}=\Z(m_1+m_2)$：ボソン的基底状態は2つの極小の\textbf{対称な}組合せ。$C^1/\operatorname{im}\delta=\Z^2/\Z(M_1-M_2)$ の生成元は $M_1$（$\equiv M_2$）で、直交補空間で代表すれば $M_1+M_2$：フェルミオン的基底状態は2つの極大の対称な組合せ。反対称な組合せ $m_1-m_2$ と $M_1-M_2$ はインスタントン（流線）で結ばれて $E>0$ に持ち上がる。
厳密解：$e^{-th}=e^{-t\cos2\theta}$ は $\cos2\theta=-1$ の点 $\theta=\pi/2,3\pi/2$（$m_1,m_2$）に等しい重みで局在し、$e^{th}d\theta$ は $\theta=0,\pi$（$M_1,M_2$）に等しい重みで局在する。いずれも $S^1$ 上で一価で、対称な組合せに対応する。これらは $b_0=b_1=1$ に対応する。
\end{solution}

\begin{problem}[難易度 D]\label{ch11-prob-witten}
$n$ 次元コンパクト・リーマン多様体 $M$ 上のウィッテン変形 $d_t=e^{-th}de^{th}$ を考える。\\
(1) 平坦な座標の近傍で $d=\sum_ia^{i\dagger}\partial_i$、$d^\dagger=-\sum_ia^i\partial_i$（$a^{i\dagger}=dx^i\wedge$、$a^i=\iota_{\partial_i}$）と書けることを用いて、命題 \ref{prop:ch11-Ht} の式
\[
H_t=-\sum_i\partial_i^2+t^2\lvert\nabla h\rvert^2+t\sum_{i,j}(\partial_i\partial_jh)[a^{i\dagger},a^j]
\]
を導け。\\
(2) 臨界点 $p$ でヘッセ行列が $\operatorname{diag}(\lambda_1,\dots,\lambda_n)$（すべて $\neq0$）となる座標で、調和近似 $h\approx h(p)+\tfrac12\sum\lambda_ix_i^2$ の下での $H_t$ のスペクトルを求めよ。$E=0$ の状態がちょうど1つで、その次数が $\ind p$ であることを示し、その波動関数を書け。\\
(3)「$t\to\infty$ で、$H_t|_{\Omega^k}$ の $O(t)$ より小さい固有値の数は、調和近似で $E=0$ となる状態の数に一致する」ことを認めて、弱モース不等式 $b_k\le c_k$ と $\chi(M)=\sum_k(-1)^kc_k$ を導け。\\
(4) ベクトル場 $V$ の孤立零点 $p$ の指数を $\operatorname{sign}\det(\partial_iV^j(p))$（非退化の場合）とする。勾配ベクトル場 $V=\nabla h$ の零点 $p$ の指数が $(-1)^{\ind p}$ であることを示し、(3) から $\sum_{V(p)=0}(\text{指数})=\chi(M)$（勾配場に対するポアンカレ・ホップの定理）を導け。$S^2$ 上で高さ関数を用いて確かめ、$S^2$ 上のいたるところ $0$ でないベクトル場は存在しないことを、「任意のベクトル場の零点の指数和は $\chi$ に等しい」（一般のポアンカレ・ホップ）を認めて結論せよ。
\end{problem}

\begin{solution}
(1) $\{a^i,a^{j\dagger}\}=\delta^{ij}$、$\{a^i,a^j\}=\{a^{i\dagger},a^{j\dagger}\}=0$。$d_t=d+t\sum_i(\partial_ih)a^{i\dagger}$、$d_t^\dagger=d^\dagger+t\sum_j(\partial_jh)a^j$。
\[
H_t=\{d,d^\dagger\}+t^2\{\partial_ih\,a^{i\dagger},\partial_jh\,a^j\}+t\bigl(\{d,\partial_jh\,a^j\}+\{\partial_ih\,a^{i\dagger},d^\dagger\}\bigr)
\]
（$\{d,d\}$ 型の項はない。同じ添字は和をとる。）第1項は $-\sum\partial_i^2$（平坦）、第2項は $(\partial_ih)(\partial_jh)\delta^{ij}=\lvert\nabla h\rvert^2$。第3項の計算は次の【注】。結果は $t(\partial_i\partial_jh)[a^{i\dagger},a^j]$。
\begin{note}
形式 $\omega$ に作用させる。微分は係数にのみ作用し、$a$ は定数行列とみなせる。
\[
\{a^{i\dagger}\partial_i,\ \partial_jh\,a^j\}\omega=a^{i\dagger}a^j(\partial_i\partial_jh)\omega+\partial_jh\,(a^{i\dagger}a^j+a^ja^{i\dagger})\partial_i\omega=(\partial_i\partial_jh)a^{i\dagger}a^j\omega+(\partial_ih)\partial_i\omega,
\]
\[
\{\partial_ih\,a^{i\dagger},\ -a^j\partial_j\}\omega=-\partial_ih\,a^{i\dagger}a^j\partial_j\omega-a^ja^{i\dagger}(\partial_j\partial_ih)\omega-\partial_ih\,a^ja^{i\dagger}\partial_j\omega=-(\partial_ih)\partial_i\omega-(\partial_i\partial_jh)a^ja^{i\dagger}\omega .
\]
和は $(\partial_i\partial_jh)(a^{i\dagger}a^j-a^ja^{i\dagger})\omega$。1階微分の項は打ち消す。
\end{note}

(2) 調和近似で $\partial_ih=\lambda_ix_i$、$\partial_i\partial_jh=\lambda_i\delta_{ij}$ なので
\[
H_t=\sum_{i=1}^n\Bigl(-\partial_i^2+t^2\lambda_i^2x_i^2+t\lambda_i[a^{i\dagger},a^i]\Bigr),
\]
互いに可換な $n$ 個の1次元の項の和。$[a^{i\dagger},a^i]=2a^{i\dagger}a^i-1$ は $dx^i$ を含む形式に $+1$、含まない形式に $-1$ をかける（$\varepsilon_i$）。$-\partial^2+t^2\lambda^2x^2$ の固有値は $t\lvert\lambda\rvert(2N+1)$ なので
\[
E=t\sum_i\bigl(\lvert\lambda_i\rvert(2N_i+1)+\lambda_i\varepsilon_i\bigr).
\]
各項は $\ge t(\lvert\lambda_i\rvert-\lvert\lambda_i\rvert)=0$ で、$0$ になるのは $N_i=0$ かつ $\varepsilon_i=-\operatorname{sign}\lambda_i$、すなわち「$\lambda_i<0$ なら $dx^i$ を含み、$\lambda_i>0$ なら含まない」とき。よって $E=0$ の状態はただ1つで、その次数は負の固有値の数 $=\ind p$。波動関数は、負の方向の添字を $i_1<\dots<i_k$ として
\[
\psi_p=\exp\Bigl(-\frac t2\sum_i\lvert\lambda_i\rvert x_i^2\Bigr)\,dx^{i_1}\wedge\dots\wedge dx^{i_k}.
\]
（1次元の $e^{-th}$, $e^{th}dx$ の積になっている。）他の状態は $E\ge2t\min\lvert\lambda_i\rvert$ で、$t\to\infty$ で持ち上がる。

(3) 認めた事実により、$\Omega^k$ の低エネルギー状態の数は「調和近似で $E=0$ の状態が次数 $k$ である臨界点」の数、すなわち $c_k$ である。真の $E=0$ 状態の数は（問題 \ref{ch11-prob-pair}(4) より $t$ によらず）$b_k$ で、それは低エネルギー状態の一部なので $b_k\le c_k$。また $\Tr(-1)^Fe^{-\beta H_t}$ は $t$ によらず $\chi(M)$ で、$t\to\infty$ で低エネルギー状態のうち $E>0$ のものは $d_t$ で対になる（低エネルギー部分空間は $d_t$ で保たれる）ので、$\chi(M)=\sum(-1)^kc_k$。

(4) $V^j=g^{jk}\partial_kh$ の零点は臨界点で、正規直交座標で $\partial_iV^j(p)=\partial_i\partial_jh(p)$。よって指数は $\operatorname{sign}\prod_i\lambda_i=(-1)^{\#\{\lambda_i<0\}}=(-1)^{\ind p}$。(3) より
\[
\sum_{V(p)=0}(-1)^{\ind p}=\sum_k(-1)^kc_k=\chi(M).
\]
$S^2$ の高さ関数では、南極（指数0）と北極（指数2）の勾配場の指数はどちらも $+1$ で、和は $2=\chi(S^2)$。北極では勾配は「流れ出し」、南極では「流れ込み」だが、2次元では湧き出しも吸い込みも指数 $+1$ である。一般のポアンカレ・ホップより、零点のないベクトル場があれば指数和は $0$ で $\chi(S^2)=2$ に反する。よって $S^2$ 上のベクトル場は必ず零点をもつ（「毛玉の定理」）。$T^2$（$\chi=0$）にはいたるところ $0$ でないベクトル場 $\partial_\theta$ がある。超対称量子力学の言葉では、ウィッテン指数はベクトル場（超ポテンシャルの勾配）の零点での局所的寄与の和に局在する。これは第9章の指数定理の局在公式の原型である。
\end{solution}

\subsection{ボット周期性とK理論}

\begin{problem}[難易度 A]\label{ch11-prob-bott}
定理 \ref{thm:ch11-bott} の表と命題 \ref{prop:ch11-stable} の安定域を用いて次を答えよ。\\
(1) $\pi_3(SU(N))$（$N\ge2$）、$\pi_3(SO(N))$（$N\ge5$）、$\pi_3(Sp(N))$。4次元ゲージ理論での意味は何か。\\
(2) $\pi_4(SU(2))$、$\pi_4(SU(N))$（$N\ge3$）、$\pi_4(Sp(N))$。4次元の $Sp(N)$ ゲージ理論にウィッテン型の大域的アノマリー（第9章）が起こりうるか。\\
(3) $\pi_1(SO(N))$（$N\ge3$）と $\pi_7(SO(N))$（$N\ge9$）。\\
(4) $\tilde K(S^n)$ と $\widetilde{KO}(S^n)$ を $n=1,\dots,8$ について表にせよ。
\end{problem}

\begin{solution}
(1) $\pi_3(SU(N))=\pi_3(U)=\Z$（$3\le2N-1$、$N\ge2$）。$\pi_3(SO(N))=\pi_3(O)=\Z$（$3\le N-2$、$N\ge5$）。$\pi_3(Sp(N))=\pi_3(Sp)=\Z$（$3\le4N+1$、すべての $N\ge1$）。いずれも、ユークリッド4次元で遠方の純ゲージ $g\colon S^3_\infty\to G$ の巻き付き数、すなわちインスタントン数が整数であることを示す（第3, 8章）。$SO(4)$ は安定域外で、実際 $\pi_3(SO(4))=\Z\oplus\Z$。

(2) $SU(2)=Sp(1)$ で $4\le4\cdot1+1$ なので $\pi_4(SU(2))=\pi_4(Sp)=\Z_2$。$N\ge3$ では $4\le2N-1$ なので $\pi_4(SU(N))=\pi_4(U)=0$。$\pi_4(Sp(N))=\pi_4(Sp)=\Z_2$（すべての $N$）。ウィッテンの $SU(2)$ アノマリーは $\pi_4(G)=\Z_2$ の生成元に沿った大きなゲージ変換の下でのフェルミオン行列式の符号反転なので、$Sp(N)$ 理論でも、擬実表現（基本表現など）のワイル・フェルミオンが奇数個あれば起こる。$SU(N\ge3)$ では $\pi_4=0$ なのでこの型の大域的アノマリーはない（その代わり摂動的アノマリーの相殺が必要）。

(3) $\pi_1(SO(N))=\pi_1(O)=\Z_2$（$1\le N-2$）：$2\pi$ 回転の非可縮性とスピノル（第2章）。$\pi_7(SO(N))=\pi_7(O)=\Z$（$7\le N-2$）。

(4) $\tilde K(S^n)=\pi_{n-1}(U)$、$\widetilde{KO}(S^n)=\pi_{n-1}(O)$ より
\begin{center}
\begin{tabular}{c|cccccccc}
$n$ & 1 & 2 & 3 & 4 & 5 & 6 & 7 & 8\\ \hline
$\tilde K(S^n)$ & 0 & $\Z$ & 0 & $\Z$ & 0 & $\Z$ & 0 & $\Z$\\
$\widetilde{KO}(S^n)$ & $\Z_2$ & $\Z_2$ & 0 & $\Z$ & 0 & 0 & 0 & $\Z$
\end{tabular}
\end{center}
例えば $\widetilde{KO}(S^1)=\Z_2$ は円周上の実直線束がメビウスの帯と自明束の2種であること、$\widetilde{KO}(S^2)=\Z_2$ は $S^2$ 上の実束のスピン構造の障害 $w_2$ に対応する。
\end{solution}

\begin{problem}[難易度 C]\label{ch11-prob-KS2}
$S^2$ を北半球 $D_N$ と南半球 $D_S$ に分け、赤道 $S^1$（座標 $\varphi$）上の貼り合わせ関数 $g_n(\varphi)=e^{in\varphi}\in U(1)$ で作った複素直線束を $L_n$ とする（モノポール電荷 $n$ の束；向きを適当に選んで $\int_{S^2}c_1(L_n)=n$、第7, 8章）。階数 $r$ の束は貼り合わせ関数のホモトピー類 $[S^1,U(r)]=\pi_1(U(r))\cong\Z$（$\det$ の巻き付き数）で分類されることを認める。\\
(1) $L_n\oplus L_m\cong L_{n+m}\oplus L_0$ を、貼り合わせ関数のホモトピーを具体的に作って示せ。\\
(2) $S^2$ 上の任意の階数 $r$ の束 $E$ は $L_n\oplus\C^{r-1}$（$n=\int c_1(E)$）と同型であることを示し、$\tilde K(S^2)\cong\Z$ で生成元が $x=[L_1]-1$ であることを示せ。\\
(3) $L_n\otimes L_m\cong L_{n+m}$ を用いて、環としての関係 $x^2=0$、$K(S^2)\cong\Z[x]/(x^2)$、$[L_n]=1+nx$ を示せ。\\
(4) チャーン指標 $\ch(L_n)=1+c_1(L_n)$ を用いて、同型 $\tilde K(S^2)\to\Z$ が $\xi\mapsto\int_{S^2}\ch(\xi)$ で与えられることを確かめよ。\\
(5) $SU(2)\to U(1)$ と破れたゲージ理論（ジョージ・グラショウ模型）の 't Hooft--Polyakov モノポールでは、遠方の $U(1)$ 束は $L_1$ で非自明だが、$SU(2)$ の随伴的な見方では $L_1\oplus L_{-1}$ の形の階数2の束が現れる。$L_1\oplus L_{-1}$ が自明であることを示し、その物理的意味を述べよ。
\end{problem}

\begin{solution}
(1) $L_n\oplus L_m$ の貼り合わせ関数は $\operatorname{diag}(e^{in\varphi},e^{im\varphi})=\operatorname{diag}(e^{i(n+m)\varphi},1)\cdot u(\varphi)$、$u(\varphi)=\operatorname{diag}(e^{-im\varphi},e^{im\varphi})\in SU(2)$。$u$ が $SU(2)$ 内で定値ループ $\mathbf 1$ にホモトピックであることを示せばよい。$SU(2)$ の元を $a_0\mathbf 1+i\bm a\cdot\bm\sigma$（$a_0^2+\lvert\bm a\rvert^2=1$）と書くと（第1章）
\[
u(\varphi)=\cos m\varphi\,\mathbf 1-i\sin m\varphi\,\sigma_3,\qquad (a_0,a_1,a_2,a_3)=(\cos m\varphi,0,0,-\sin m\varphi)
\]
で、$u$ は $a_1=a_2=0$ の大円上にある。特に点 $P=i\sigma_1$（$(a_0,\bm a)=(0,1,0,0)$）を通らない。$S^3\setminus\{P\}$ は $P$ からの立体射影 $\sigma_P$ で $\R^3$ と同相なので、
\[
u_s(\varphi)=\sigma_P^{-1}\bigl((1-s)\,\sigma_P(u(\varphi))+s\,\sigma_P(\mathbf 1)\bigr)\qquad(0\le s\le1)
\]
は $SU(2)\setminus\{P\}$ 内のループの連続な族で、$u_0=u$、$u_1=\mathbf 1$（定値）。よって $u\simeq\mathbf 1$。

したがって貼り合わせ関数は $U(2)$ 内で $\operatorname{diag}(e^{i(n+m)\varphi},1)$ にホモトピックで、これは $L_{n+m}\oplus L_0$ の貼り合わせ関数である。貼り合わせ関数がホモトピックなら束は同型（第7章）なので $L_n\oplus L_m\cong L_{n+m}\oplus\C$。（$\det$ の巻き付き数がどちらも $n+m$ であることからも、認めた分類より直ちに従う。）

(2) $E$ の貼り合わせ関数 $g\colon S^1\to U(r)$ の $\det$ の巻き付き数を $n$ とすると、$\operatorname{diag}(e^{in\varphi},1,\dots,1)$ も巻き付き数 $n$ なので、分類より $E\cong L_n\oplus\C^{r-1}$。$c_1$ は直和で加法的で $c_1(\C)=0$ なので $n=\int c_1(E)$。
$\tilde K(S^2)$ の元は $[E]-[\C^r]=[L_n]-1$ と書ける。(1) を繰り返すと $L_1^{\oplus n}\cong L_n\oplus\C^{n-1}$（$n\ge1$）なので $[L_n]-1=n([L_1]-1)=nx$。$n<0$ でも $L_n\oplus L_{-n}\cong\C^2$ より $[L_n]-1=-([L_{-n}]-1)=nx$。準同型 $\tilde K(S^2)\to\Z$、$[E]-[\C^r]\mapsto\int c_1(E)$ は well-defined（$c_1$ は安定同型で不変）で、$x\mapsto1$ なので全射、$nx\mapsto n$ より単射。よって $\tilde K(S^2)=\Z x$。

(3) テンソル積の貼り合わせ関数は積 $e^{in\varphi}e^{im\varphi}=e^{i(n+m)\varphi}$ なので $L_n\otimes L_m\cong L_{n+m}$。
\[
x^2=[L_1\otimes L_1]-2[L_1]+1=[L_2]+[L_0]-2[L_1]=0,
\]
最後は (1) で $n=m=1$ とした $L_1\oplus L_1\cong L_2\oplus L_0$。よって $K(S^2)=\Z\oplus\Z x$、$x^2=0$、すなわち $K(S^2)\cong\Z[x]/(x^2)$。$[L_n]=1+nx$ は (2) から。（$[L_1]^n=(1+x)^n=1+nx$ とも整合する。）

(4) $S^2$ 上では $4$ 形式以上は $0$ なので $\ch(L_n)=1+c_1(L_n)$ で、$\int_{S^2}\ch([L_n]-1)=\int c_1(L_n)=n$。よって $\int\ch$ は (2) の同型 $nx\mapsto n$ と一致する。一般に $\tilde K(S^{2m})\cong\Z$ は $\int\ch_m$ で与えられる（定理 \ref{thm:ch11-Ksphere} の後の注意）。

(5) (1) で $n=1$、$m=-1$ とすれば $L_1\oplus L_{-1}\cong L_0\oplus\C=\C^2$。$SU(2)$ 主束（あるいは基本表現の階数2束）は $\pi_1(SU(2))=0$ によって $S^2$ 上ではすべて自明である。物理的意味：'t Hooft--Polyakov モノポールは $SU(2)$ ゲージ場としては特異点のない滑らかな配位で、$SU(2)$ 束としては自明なまま、ヒッグス場が $U(1)$ 部分群を選ぶことで「$U(1)$ 部分束」が $L_1$（電荷を持つ成分）と $L_{-1}$ に分かれて非自明になる。モノポール電荷は全体の束のトポロジーではなく、ヒッグス場による分解 $\C^2=L_1\oplus L_{-1}$ のトポロジー、すなわち $\pi_2(SU(2)/U(1))=\pi_2(S^2)=\Z$ によって決まる（第3章）。K 理論では $[L_1]+[L_{-1}]-2=x-x=0$。
\end{solution}

\begin{problem}[難易度 D]\label{ch11-prob-dbrane}
Dブレーンの電荷が K 理論に値をとるという主張（第2節の physbox）を認める。\\
(1) 定理 \ref{thm:ch11-Ksphere} とボット周期性を用いて、IIB では $p$ が奇数、IIA では $p$ が偶数の Dp ブレーンだけが $\Z$ に値をとる保存電荷をもつことを示せ（$-1\le p\le9$）。\\
(2) タイプI弦理論では電荷が $\widetilde{KO}(S^{9-p})$ に値をとる。$p=-1,0,1,\dots,9$ について電荷の群を求め、$\Z$ 電荷をもつもの（BPS の D1, D5, D9）と $\Z_2$ 電荷をもつもの（非 BPS）を答えよ。\\
(3) IIB で D9 と反 D9 を1枚ずつ重ね、タキオンを $T(x)=T_0\,\phi(\rho)\,e^{i\theta}$（$x^1+ix^2=\rho e^{i\theta}$、$\phi(0)=0$、$\phi(\infty)=1$）とする。遠方の円周上の写像 $T/\lvert T\rvert\colon S^1\to U(1)$ の巻き付き数を求め、これが $\tilde K(S^2)=\pi_1(U)$ の生成元であること、芯に D7 ブレーンが1枚残ることを説明せよ。\\
(4) 余次元4では D9 と反 D9 を2枚ずつ（$E,F$ の階数2）とし、$x=(x^1,\dots,x^4)$ に対して
\[
T(x)=x^4\,\mathbf 1_2+i\sum_{a=1}^3x^a\sigma^a
\]
とする（アティヤ・ボット・シャピロ構成）。$T/\lvert T\rvert$ を $S^3\to SU(2)$ とみたときの巻き付き数（$w(g)=\frac1{24\pi^2}\int\tr(g^{-1}dg)^3$）が $\pm1$ であることを示し、芯に D5 ブレーンが1枚残ることを説明せよ。なぜ1枚ずつでは足りないか。\\
(5) タイプIの D0 ブレーンの電荷が $\Z_2$ であることの物理的意味を述べよ。
\end{problem}

\begin{solution}
(1) Dp ブレーンの横方向は $\R^{9-p}$ で、コンパクト台の K 群は $\tilde K(S^{9-p})$。IIB：$\tilde K(S^{9-p})=\pi_{8-p}(U)$ は $9-p$ 偶数、すなわち $p$ 奇数（$p=-1,1,3,5,7,9$）で $\Z$、$p$ 偶数で $0$。IIA：$\tilde K^{-1}(S^{9-p})=\tilde K(S^{10-p})=\pi_{9-p}(U)$ は $p$ 偶数（$p=0,2,4,6,8$）で $\Z$、$p$ 奇数で $0$。これは既知の BPS Dブレーン（IIB は RR 形式 $C_0,C_2,C_4,\dots$ に、IIA は $C_1,C_3,\dots$ に結合する）と一致する。逆に IIB の $p$ 偶数（IIA の $p$ 奇数）の Dp ブレーンは存在するがタキオンをもつ不安定な非 BPS ブレーンであり、保存電荷をもたないこと（$\tilde K=0$）と整合する。

(2) $\widetilde{KO}(S^{9-p})$ は $9-p\bmod8$ で決まる：
\begin{center}
\begin{tabular}{c|ccccccccccc}
$p$ & 9 & 8 & 7 & 6 & 5 & 4 & 3 & 2 & 1 & 0 & $-1$\\ \hline
$9-p$ & 0 & 1 & 2 & 3 & 4 & 5 & 6 & 7 & 8 & 9 & 10\\
$9-p\bmod 8$ & 0 & 1 & 2 & 3 & 4 & 5 & 6 & 7 & 0 & 1 & 2\\
電荷 & $\Z$ & $\Z_2$ & $\Z_2$ & 0 & $\Z$ & 0 & 0 & 0 & $\Z$ & $\Z_2$ & $\Z_2$
\end{tabular}
\end{center}
（$p=9$ では横方向が点で、$\widetilde{KO}(S^0)=\Z$ とみなす：D9 の枚数。）$\Z$：D9, D5, D1（BPS；D1 は $SO(32)$ 異型弦理論の基本弦と S 双対、D5 は $SO(32)$ のインスタントン）。$\Z_2$：D8, D7, D0, D$(-1)$（非 BPS だが安定）。

(3) $\rho\to\infty$ で $T/\lvert T\rvert=e^{i\theta}$、巻き付き数 $1$。$\tilde K(S^2)$ は横方向 $\R^2$ のコンパクト台の K 群で、$(E,F,T)$ の三つ組は「遠方で $T$ が同型」なので、遠方の円周上の貼り合わせ関数 $T/\lvert T\rvert\in[S^1,U(1)]$ を定め、これが $\pi_1(U)=\Z$ の元である（$S^2=D^2\cup D^2_\infty$ を $E$ と $F$ で貼り合わせた束が $L_{\pm1}$）。巻き付き数 $1$ はその生成元で、$\tilde K(S^2)$ の生成元 $x=[L_1]-1$（問題 \ref{ch11-prob-KS2}）に対応する（符号は向きの規約による）。遠方ではタキオンは真空 $\lvert T\rvert=T_0$ にあり、ブレーン対は閉弦真空へ消滅しているが、$\rho=0$ では $T=0$ で凝縮できない。この余次元2の渦の芯（$x^1=x^2=0$ の8次元）に、1単位の D7 ブレーン電荷が残る。

(4) $\lvert x\rvert=1$ では $T=x^4+i\bm x\cdot\bm\sigma\in SU(2)$（第1章の $a_0+i\bm a\cdot\bm\sigma$ 表示）で、$T|_{S^3}\colon S^3\to SU(2)\cong S^3$ は恒等写像（座標の同一視）である。記法と規約により $S^3$ の向きは恒等写像 $S^3\cong SU(2)$ が $w=1$ となるように選んであるので $w(T/\lvert T\rvert)=1$（逆の向きなら $-1$）。$\pi_3(U(2))\cong\pi_3(U)=\Z$（安定域 $3\le2\cdot2-1$）の生成元なので、$\tilde K(S^4)$ の生成元であり、対応する $S^4$ 上の階数2の束はインスタントン数 $\pm1$ の $SU(2)$ 束である（第8章）。芯 $x=0$（6次元）に D5 ブレーンが1枚残る。
1枚ずつ（$U(1)$）では足りない理由：$\pi_3(U(1))=0$ なので、$T\colon S^3\to U(1)$ は必ず可縮で位相的な巻き付きをもてない。$\pi_3(U(N))$ が $\Z$ になるのは $N\ge2$ からで、ABS 構成ではスピノル $S^\pm\cong\C^{2^{k-1}}$（$2k=4$ で $2$）が必要枚数を与える。一般に余次元 $2k$ には $2^{k-1}$ 対の D9・反 D9 で足りる。

(5) $\Z_2$ 電荷とは、D0 ブレーン2つは位相的には何もない状態と同じ（$1+1=0$）で、対消滅しうるが、1つだけでは崩壊できないことを意味する。BPS 条件による質量の下限で守られているのではなく、K 理論の位相的な電荷で守られた安定な非 BPS 状態である。実際、タイプI弦理論の D0 ブレーンは $SO(32)$ のスピノル表現の量子数をもつ最も軽い状態で、S 双対な $SO(32)$ 異型弦理論のスピノル状態（摂動的スペクトルに存在する）に対応する（セン）。この $\Z_2$ はコホモロジー（チャーン指標）では見えず、RR 場の電荷として測ることもできない。これが、Dブレーン電荷はコホモロジーではなく K 理論で分類されるべきだという主張の最も鮮やかな証拠である。
\end{solution}


\printsolutions

\backmatter
\chapter*{参考文献}
\addcontentsline{toc}{chapter}{参考文献}
\begin{enumerate}
  \item M. Nakahara, \textit{Geometry, Topology and Physics}, 2nd ed., IOP (2003).
  \item T. Frankel, \textit{The Geometry of Physics}, 3rd ed., Cambridge (2011).
  \item T. Eguchi, P. B. Gilkey, A. J. Hanson, ``Gravitation, gauge theories and differential geometry,'' Phys.\ Rept.\ \textbf{66} (1980) 213.
  \item R. Bott, L. W. Tu, \textit{Differential Forms in Algebraic Topology}, Springer (1982).
  \item A. Hatcher, \textit{Algebraic Topology}, Cambridge (2002).
  \item S. Coleman, \textit{Aspects of Symmetry}, Cambridge (1985)（特に ``The uses of instantons'', ``Classical lumps and their quantum descendants''）.
  \item N. Manton, P. Sutcliffe, \textit{Topological Solitons}, Cambridge (2004).
  \item S. Weinberg, \textit{The Quantum Theory of Fields, Vol.~II}, Cambridge (1996)（第22, 23章）.
  \item K. Fujikawa, H. Suzuki, \textit{Path Integrals and Quantum Anomalies}, Oxford (2004).
  \item L. Alvarez-Gaum\'e, P. Ginsparg, ``The structure of gauge and gravitational anomalies,'' Ann.\ Phys.\ \textbf{161} (1985) 423.
  \item E. Witten, ``Supersymmetry and Morse theory,'' J.\ Diff.\ Geom.\ \textbf{17} (1982) 661.
  \item E. Witten, ``An SU(2) anomaly,'' Phys.\ Lett.\ B \textbf{117} (1982) 324.
  \item M. B. Green, J. H. Schwarz, E. Witten, \textit{Superstring Theory, Vol.~2}, Cambridge (1987).
  \item J. Polchinski, \textit{String Theory, Vol.~II}, Cambridge (1998).
  \item K. Hori \textit{et al.}, \textit{Mirror Symmetry}, Clay Math.\ Monographs (2003).
\end{enumerate}

\end{document}
